% !TEX root = nc-cubic-threefold-moduli.tex
\documentclass[11pt,a4paper]{ltjsarticle}
\usepackage{amsmath,amssymb,amsthm,mathtools,bm,mathrsfs}
\usepackage{luatexja-fontspec}
\usepackage{enumitem,booktabs,array,needspace}
\usepackage{xcolor}
\usepackage[colorlinks=true,linkcolor=blue!45!black,citecolor=blue!45!black,urlcolor=blue!45!black]{hyperref}
\IfFileExists{/usr/share/fonts/opentype/noto/NotoSerifCJK-Regular.ttc}{%
  \setmainjfont[Path=/usr/share/fonts/opentype/noto/]{NotoSerifCJK-Regular.ttc}%
}{\setmainjfont{HaranoAjiMincho-Regular}}
\setlength{\textwidth}{15.6cm}
\setlength{\oddsidemargin}{0.25cm}
\setlength{\evensidemargin}{0.25cm}
\setlength{\textheight}{22.6cm}
\setlength{\topmargin}{-0.2cm}
\setlength{\parindent}{1em}
\setlength{\parskip}{1pt}
\numberwithin{equation}{section}
\newtheorem{theorem}{定理}[section]
\newtheorem{maintheorem}{主定理}
\renewcommand{\themaintheorem}{\Alph{maintheorem}}
\newtheorem{lemma}[theorem]{補題}
\newtheorem{proposition}[theorem]{命題}
\newtheorem{corollary}[theorem]{系}
\newtheorem{remark}[theorem]{注意}
\newtheorem{definition}[theorem]{定義}
\newtheorem{example}[theorem]{例}
\renewcommand{\proofname}{証明}
\newcommand{\C}{\mathbb C}
\providecommand{\CC}{\C}
\newcommand{\PP}{\mathbb P}
\newcommand{\Gr}{\operatorname{Gr}}
\newcommand{\A}{\mathbb A}
\newcommand{\OO}{\mathcal O}
\newcommand{\cB}{\mathcal B}
\newcommand{\cA}{\mathcal A}
\newcommand{\cI}{\mathcal I}
\newcommand{\cQ}{\mathcal Q}
\newcommand{\cZ}{\mathcal Z}
\newcommand{\Sym}{\operatorname{Sym}}
\newcommand{\gr}{\operatorname{gr}}
\newcommand{\qgr}{\operatorname{qgr}}
\newcommand{\gldim}{\operatorname{gldim}}
\newcommand{\rank}{\operatorname{rank}}
\newcommand{\Span}{\operatorname{span}}
\newcommand{\Spec}{\operatorname{Spec}}
\newcommand{\Aut}{\operatorname{Aut}}
\newcommand{\Hom}{\operatorname{Hom}}
\newcommand{\End}{\operatorname{End}}
\newcommand{\Ext}{\operatorname{Ext}}
\newcommand{\Proj}{\operatorname{Proj}}
\newcommand{\Frac}{\operatorname{Frac}}
\newcommand{\trdeg}{\operatorname{trdeg}}
\newcommand{\GL}{\mathrm{GL}}
\newcommand{\PGL}{\mathrm{PGL}}
\newcommand{\Lie}{\operatorname{Lie}}
\newcommand{\Z}{\operatorname{Z}}
\newcommand{\Tens}{\operatorname{T}}
\newcommand{\CY}{\mathrm{CY}}
\newcommand{\AS}{\mathrm{AS}}
\newcommand{\nc}{\mathrm{nc}}
\newcommand{\Symm}{\operatorname{sym}}
\title{非可換三次三様体のAbel圏モジュライ\\二次包絡とKuznetsov残余射}
\author{}
\date{2026年11月19日・公表準備稿 v21}
\hypersetup{pdftitle={非可換三次三様体のAbel圏モジュライ},pdfsubject={二次包絡とKuznetsov残余射}}
\begin{document}
\maketitle
\begin{abstract}
構造対象 $\mathcal O$ を保つ全ての完全 Abel 圏同値を射とする、
非可換三次三様体のモジュライを考える。
二つの中心一次変数を持つ Jacobian 型中心三次商から、
滑らかな有限型十七次元 Deligne--Mumford 部分 stack $\mathfrak C_{\mathrm J}$ を構成する。
二次包絡への射は、十三次元の周囲代数の stack 上の階数四ベクトル束の torsor である。
一方、$\mathfrak C_{\mathrm J}$ の非空開上では、Kuznetsov 残余圏が
滑らかな可換三次三様体の Kuznetsov 成分となる。
これにより可換三次の十次元モジュライへの支配的な射を得る。
適切な開に制限すれば、この射は相対次元七の滑らかな全射であり、
残余圏と全 compact tails 圏は相対 dg smooth proper である。
証明では、全 Abel 同値に対する二次包絡と残余操作の自然性、
圏変形の形式的認識、係数環上 perfect な同値核の効果化、
および三次成分の Serre 安定部分圏の生成性を組み合わせる。

さらに、滑らかな可換三次点の framed Abel 形式近傍を、
十次元形式円板と $\operatorname{Gr}(2,5)$ のアフィン錐の頂点完備化との積として記述する。
この同型は、残余操作が誘導する形式射に両立するよう選べる。
有限係数の Jacobian 模型の適切な非空開上では、
埋込み点スキームを楕円曲線と長さ二十四の完全交叉に分解し、
そのデータから ambient Poisson 直線と五次元の三次 Casimir 空間を復元する。
点対応の shift を含めると二次包絡と Poisson 構造の尺度も復元できる。
量子中心から特定の古典三次を選ぶ問題は残り、
単純な対称化および点スキームと Poisson 零点の直接同一視に対する障害を示す。
\end{abstract}

\setcounter{tocdepth}{1}
\tableofcontents

\section{序論}\label{sec:intro}

滑らかな三次三様体 $X\subset\PP^4$ の導来圏には半直交分解
\[
 \mathrm D^b(X)=\langle\operatorname{Ku}(X),\mathcal O_X,\mathcal O_X(1)\rangle
\]
があり、Kuznetsov 成分 $\operatorname{Ku}(X)$ は $X$ の同型類を決定する
\cite[Theorem~1.1]{BMMSCubicTorelli}。
本稿では、構造対象付き Abel 圏として三次三様体を非可換変形したとき、
この残余圏がどの可換幾何を保つかを調べる。
主な結果は、明示的な十七次元の非可換族の一般点で残余圏が再び可換三次の成分となり、
これが十次元の可換モジュライへの実際の代数的な射を与えることである。
同じ族には二次包絡への十三次元の射もあり、二つの射は異なる情報を取り出す。

対象は $(\mathcal C,\mathcal O)$ という構造対象付き線型 Abel 圏である。
三次の Hilbert 級数と相対復元条件を満たす許容豊富列の\emph{存在}を要求するが、
その選択は標識に含めない。射には $\mathcal O$ を保つ全ての完全線型同値を含め、
標識を保つ自然同型で割る。
この fpqc モジュライ問題を $\mathfrak M^{\mathcal O}_{\mathrm{cub}}$ と書く
（定義は第\ref{sec:abel-marked-moduli}節）。
非 Noetherian な係数環では大きい tails 圏 $\operatorname{QGr}$ を用い、
導来圏の内在的な有限部分として $\mathrm D(\operatorname{QGr})^c$ を取る。
以下の具体的な部分族の中心条件は、このモジュライ全体の定義には課さない。
従って可換な三次も同じモジュライ問題に含まれる。

\subsection{二つのモジュライ射}

五次元空間 $V$ 上の二次代数と中心三次直線の対
\begin{equation}\label{eq:data}
 B=\Tens(V)/(R_2),\qquad [w]\in\PP(\Z(B)_3),\qquad A=B/(w)
\end{equation}
を考える。扱う $B$ は、二つの中心一次変数 $u,v$ と三つの非中心変数を持つ
Jacobian 型の相対 Koszul--AS 五次元代数であり、$w$ は中心正則元である。
全混合係数を許した族とその明示的な非空開を
第\ref{sec:construction}、\ref{sec:normal}、\ref{sec:general17-relative-algebra}節で構成する。
三次関係を加えても二次関係は変わらないため、$A$ の二次包絡は $B$ である。
この単純な観察を、豊富列の選択を忘れた全 Abel 同値と任意の係数変更に両立させることが
最初の課題となる。

\Needspace{7\baselineskip}
\begin{maintheorem}[代数的部分モジュライと二次包絡]\label{thm:intro-atlas}
滑らかな既約有限型 Deligne--Mumford stack
$\mathfrak C_{\mathrm J}$、$\mathfrak B_{\mathrm J}$ と射
\[
 \mathfrak C_{\mathrm J}\longrightarrow\mathfrak M^{\mathcal O}_{\mathrm{cub}},
 \qquad
 \tau:\mathfrak C_{\mathrm J}\longrightarrow\mathfrak B_{\mathrm J}
\]
が存在し、次を満たす。
\begin{enumerate}[label=\textup{(\roman*)}]
\item $\dim\mathfrak C_{\mathrm J}=17$、$\dim\mathfrak B_{\mathrm J}=13$ であり、
第一の射は任意の係数環上で完全忠実である。
その本質像は、指定した Jacobian 族へ fpqc 局所的に戻る対象からなる全部分 stack である。
\item $\tau$ は二次包絡を取る射であり、$\mathfrak B_{\mathrm J}$ 上の
階数四ベクトル束の加法群の torsor である。
従ってスキームで表現可能、アフィン、滑らかかつ全射である。
\end{enumerate}
\end{maintheorem}

正確な構成と証明は、命題\ref{prop:jacobian-ambient-etale-chart}、
命題\ref{prop:jacobian-affine-four-morphism}、
定理\ref{thm:jacobian-algebraic-abel-substack}に与える。
周囲代数を同定まで固定したファイバーが四次元になる理由は、
中心一次空間 $U$ に対する中心三次の完全列
\[
 0\longrightarrow\Sym^3 U\longrightarrow Z(B)_3\longrightarrow L\longrightarrow0,
 \qquad \dim U=2,\quad\dim L=1
\]
にある。許容中心三次直線はこの列の分裂を与える。
$\Hom(L,\Sym^3 U)$ の中心スカラー重みは零なので商 stack に降下するが、
$U$ や $Z(B)_3$ 自体の通常のベクトル束としての降下を主張する必要はない。
この torsor に大域的な原点は指定していない。

$\mathfrak C_{\mathrm J}$ 上の compact tails から、atlas 上の二つの例外対象
$\mathcal O,\mathcal O(1)$ による残余圏 $\mathcal K$ を作る。
$\mathcal O(1)$ の基底線束による捩れは mutation の評価写像で相殺されるので、
残余圏は全ての $\mathcal O$-保存 Abel 同値に自然である。

\Needspace{7\baselineskip}
\begin{maintheorem}[Kuznetsov 成分による残余射]\label{thm:intro-residual}
$\mathfrak C_{\mathrm J}$ の非空開部分 stack $\mathfrak U$ と、
自己同型が自明な滑らかな可換三次のモジュライの非空開 $V$ が存在し、
相対次元七の滑らかな全射
\[
 \rho:\mathfrak U\longrightarrow V
\]
を持つ。任意の幾何学的点 $c$ において、$\rho(c)=[X_c]$ は
$\mathcal K_c\simeq\operatorname{Ku}(X_c)$ を満たす三次の唯一の同型類である。
$\mathfrak U$ 上の残余圏と全 compact tails 圏は相対 dg smooth proper であり、
任意の導来底変更と両立する。
三次の族と相対 dg Morita 同値は $\mathfrak U$ 上で étale 局所的に存在する。
\end{maintheorem}

これは定理\ref{thm:kuznetsov-cubic-moduli-morphism}で証明する。
$\tau$ の四次元ファイバーは周囲代数を固定した中心三次の自由度であり、
$\rho$ の七次元ファイバーは残余圏を固定した自由度である。
両者の関係の具体的な記述は、本稿の結論には含めない。
また、$\rho(c)$ を、Jacobian 表示に用いた三次形式の零点集合と同一視していない。
$\mathfrak C_{\mathrm J}$ の全モジュライ内での開性・局所閉性・成分性も未解決である。

\subsection{可換点の全形式近傍と点データ}

可換点では、特定の次数付き表示に属する変形だけでなく全 Abel 変形を扱える。
以下の framed 変形では特殊ファイバーとの同定を固定する。
群oid は同値の自然同型類を射とする通常の $1$-群oid であり、
dg 圏の変形に現れる高次スカラー自己同型まで消去する主張ではない。

\Needspace{7\baselineskip}
\begin{maintheorem}[滑らかな可換三次点の形式近傍]\label{thm:intro-formal}
滑らかな三次三様体 $X/\C$ において、
$\mathfrak M^{\mathcal O}_{\mathrm{cub}}$ の framed Artin 形式近傍は、
剰余体 $\C$ の可換局所 Artin $\C$-代数上の群oid値関手として、
\begin{equation}\label{eq:intro-formal-ring}
 \operatorname{Spf} R_X,\qquad
 R_X=\C[[t_1,\ldots,t_{10},a_{ij}\ (0\le i<j\le4)]]
           /(\operatorname{Pfaff}_4(a))
\end{equation}
という離散群oidと同値である。
ここで $(a_{ij})$ は $5\times5$ 交代行列であり、
$\operatorname{Pfaff}_4(a)$ はその $4\times4$ Pfaffian が生成するイデアルである。
全ての平坦 Abel Artin 変形は、全次数平坦な相対 Koszul--AS 二次包絡と
中心正則三次商の模型を持つ。
残余操作から得る射 $r:\widehat{\mathfrak M}^{\mathcal O}_{X}\to\operatorname{Def}_X$ に対して、
\eqref{eq:intro-formal-ring}の表示を $r$ が $t$ 座標への射影になるように選べる。
従って形式近傍は次元十七・接空間次元二十であり、
$r$ の形式ファイバーは $\operatorname{Gr}(2,5)$ のアフィン錐の頂点完備化である。
\end{maintheorem}

全次数平坦性は定理\ref{thm:abel-cubic-flat-envelope}、
stack 化後の形式表示は系\ref{cor:abel-full-stack-discrete-formal}、
残余射との両立は系\ref{cor:abel-poisson-formal-lift}による。
この形式表示は非標準的であり、有限 $\Aut(X)$ で割った非 framed 形式近傍や、
選んだ表示の代数化は別問題である。

有限係数の模型からの復元には、埋込み点スキームとその shift を用いる。
$\Pi_F$ は $u,v$ を Casimir とし、非中心変数の括弧を
$\{x,y\}=F_z$ とその巡回置換で定める Jacobian Poisson 構造とする。

\Needspace{7\baselineskip}
\begin{maintheorem}[点スキームと点対応による復元]\label{thm:intro-points}
$h\in\C^\times$ と三次形式 $F(x,y,z,u,v)$ に対し、
$u,v$ を中心とし
\[
 [y,z]=h\Symm(F_x),\qquad [z,x]=h\Symm(F_y),\qquad [x,y]=h\Symm(F_z)
\]
を関係とする第\ref{sec:point-scheme-poisson}節の二次代数 $B$ と中心三次商 $A=B/(w)$ を考える。
古典境界三次 $f_3=F|_{u=v=0}$ は滑らかとする。
$B/(u,v)$ の点対応は、滑らかな楕円曲線 $E_h$ 上の固定点を持たない自己同型
$\sigma$ のグラフであり、$w$ はその点加群を消すと仮定する。
$w$ の tensor 代表の可換化を $w_{\mathrm{ab}}$ とし、
\[
 C=V(F_x,F_y,F_z)\subset\PP(V^*),\qquad
 D_w=C\cap V(w_{\mathrm{ab}})
\]
と置く。$w_{\mathrm{ab}}$ が $C$ のどの既約成分も含まないと仮定する。このとき
\[
 \operatorname{Pt}(A)=E_h\sqcup D_w,\qquad \operatorname{length}(D_w)=24,
 \qquad \Gamma_A=\Gamma_\sigma\sqcup\Delta_{D_w},
\]
であり、$D_w$ は型 $(2,2,2,3)$ の完全交叉である。
埋込み点対応は $H^0(I_{\Gamma_A}(1,1))=R_B$ により二次包絡を復元する。
さらに古典境界 $f_3$ が非 Sebastiani--Thom 型であれば、埋込み点スキームから
\[
 \begin{gathered}
 U=I(E_h)_1,\qquad Q=I(D_w)_2,\qquad W=\operatorname{ann}(U)\subset V^*,\\
 R_Q=\{G\in\Sym^3V:\partial_aG\in Q\text{ for all }a\in W\}
        =\C F\oplus\Sym^3U
 \end{gathered}
\]
と ambient Poisson 直線 $[\Pi_F]$ を復元できる。
shift を含む点対応からは、二次関係の交代部分と対称部分のグラフとして
尺度込みの構造 $\Pi_{\mathrm{fin}}=h\Pi_F$ を復元できる。
\end{maintheorem}

命題\ref{prop:central-two-point-scheme}、
命題\ref{prop:point-correspondence-quadratic-recovery}、
命題\ref{prop:point-gradient-reconstruction}と
命題\ref{prop:point-recovery-relative-descent}が正確な仮定と相対版を与える。
これらの復元は指定した Jacobian 開で全 $\mathcal O$-保存同値に自然である。
復元に用いるのは埋込みと scheme 構造であり、一般の Abel 圏における点対象の
内在的な特徴付けは主張していない。
有限成分は一般には被約となるが、Fermat 例では三点の各局所長が八となる。
また、五次元空間 $R_Q$ の復元は、その中の特定の三次直線の選択とは区別される。

この区別には具体的な理由がある。
滑らかな中心超曲面と Poisson 一次括弧を持つ変形では、
点加群が局所的に一次持上げを持つ条件は、半古典 Poisson 構造の零点条件に等しい
（命題\ref{prop:first-order-point-poisson-zero}）。
しかし有限パラメータでは、中心三次の対称化逆像は一般に Poisson 超曲面を定めず、
Hesse 族の点スキーム全体を同じ射影空間内のいずれの滑らかな三次上の
Poisson 零点にもできない開が存在する
（命題\ref{prop:poisson-naive-symmetrization-fails}、
\ref{prop:hesse-point-zero-nonrealization}）。
従って点データから $\rho$ の出力を直接選ぶには、量子中心を反映した補正が必要である。

\subsection{証明の構造と先行研究}

\paragraph{座標代数から全 Abel 同値へ.}
非可換射影空間と豊富列の枠組みは Artin--Schelter、Artin--Zhang、Polishchuk
\cite{AS,AZ,PolishchukProj}に基づく。
本稿の明示族では、Mori--Smith \cite[Theorem~4.2(1)]{MS} の
三変数 Jacobian 代数、Berger--Taillefer \cite{BT} の PBW 変形、
Etingof--Ginzburg と Rains \cite[Proposition~7.3.1(ii) and \S9.1]{EG} の
中心元を二つの中心変数上で用いる。
有限射影 Koszul 解消と相対 CY 双対を構成することで、
任意の係数環上の compact tails に自然な Serre pairing を得る。
これは体上の双対定理 \cite{YZ,RRGeneralizedAS} の無条件な底変更ではない。

偶数偏極と相対 section 復元により、全 $\mathcal O$-保存 Abel 同値の層を
$\underline{\operatorname{Isom}}_{\mathrm{gr}}/\mathbb G_m$ と同定する。
この商は fpqc 層商であり、大域点集合の商ではない。
次に有限な二次・三次関係窓、作用写像の étale locus、中心三次の分裂 torsor を使って
主定理\ref{thm:intro-atlas}を得る。
Behrend--Noohi \cite{BehrendNoohiModuli} の次数付き代数のモジュライ、
Abdelgadir--Okawa--Ueda \cite{AOUNCPlanes} の非可換射影平面のモジュライに対し、
ここでの具体的な追加部分は、中心三次商の族で全 $\mathcal O$-保存同値を含む
有限型部分問題と二次包絡射を同時に記述することにある。

\paragraph{形式的な圏認識から実際の残余射へ.}
Kuznetsov \cite[Theorem~8.9]{KuznetsovHochschild} による Hochschild 計算と
Jacovskis--Lin--Liu--Zhang \cite[Theorem~4.4(2)]{JLLZInfinitesimal} による
実際の Kodaira--Spencer 写像の同型を、
Lurie の tame な compactly generated 圏の変形比較
\cite[Theorems~5.1.9,5.3.16 and Proposition~5.3.21]{LurieDAGX}と組み合わせる。
これにより通常の局所 Artin 基底上の残余圏の変形の同型類を、
可換三次の変形として認識する。
次に Bhatt \cite[Lemma~4.2]{BhattAlgebraization} を
\emph{係数環上の} perfect 複体に適用して同値核とその作用を効果化し、
smooth な幾何的 target への完全忠実な埋込みを得る。
相対 Serre trace が像を Serre 安定にする。

Stevenson \cite[Theorems~3.1(3),6.13]{StevensonSingularityActions} の孤立 hypersurface の
特異部分圏分類を冪等完備化まで用い、Orlov の ungraded・graded の比較
\cite[Theorem~3.9]{OrlovLandauGinzburg}
\cite[Theorems~2.5,3.10--3.11]{OrlovSingularities}、graded 化、および
Neeman \cite[Lemma~2.2]{NeemanCompactLocalization} の compact 生成を用いると、
非零 Serre 安定 thick 部分圏は三次成分全体となる。
この段階が埋込みを同値にする。
最後に、明示的な量子化退化、perfect 核の有限表示降下
\cite[Lemma~2.10]{ToenVaquieObjects}、categorical Torelli
\cite{BMMSCubicTorelli}を使って射を代数的な非空開へ広げる。
支配性は、古典三次への係数写像の微分に対する整数行列証明書で示す。

Lin--Zhang \cite[Theorem~1.7 and \S2]{LinZhangDeformingEquivalences} は、
滑らかな既約基底上の滑らかな射影族、底変更に両立する perfect な射影核を持つ
相対 admissible 圏、および特殊点での既存の同値を考える。
幾何的 Kodaira--Spencer 写像の全射性と負次数 Hochschild コホモロジーの消滅の下で、
両基底へ滑らかに写る共通基底上に同値核を変形する定理である。
本稿で必要な、非可換族をまず幾何的な圏として認識する段階は、この定理の仮定には含まれない。
同様に、有限位数境界上の smooth proper な十六次元族の稠密性
（第\ref{sec:dg-smooth-torsion-families}節）だけを一般点の smooth 性の根拠にはしない。
主定理\ref{thm:intro-residual}は相対 Serre 双対と三次成分固有の生成性を用いて証明する。

\paragraph{全 Artin 変形と Poisson 形式錐.}
De~Deken--Lowen \cite[Theorem~4.5]{DeDekenLowen} の既知の比較定理は、
全 twist の $H^1,H^2$ 消滅により、Abel 変形を section $\mathbb Z$-代数の変形に結び付ける。
Lowen--Van den Bergh の Abel 変形理論 \cite{LowenVdBDeformation,LowenVdBHochschild}、
Lowen の代数層との比較 \cite[Theorems~3.14,3.26]{LowenRinged}、
Yekutieli の量子化比較 \cite[Corollary~11.2]{YekutieliTwisted}を既知の入力とする。
Yekutieli の Corollary~11.2 では可換積も変動を許す。
$H^1(T_X)=0$ を追加仮定する同論文 Corollary~11.3 は三次三様体には適用しない。
Barmeier--Wang \cite[Theorems~1.2 and~3.11]{BarmeierWangSheaves} は、
有限アフィン被覆の交叉から作る quiver の関係により、
分離 Noetherian スキーム上の準連接層の Abel 変形を明示的に制御する。
本稿ではこの種の変形全体との比較を必要とする一方、
三次三様体上に tilting bundle が存在するとは仮定しない。

本稿では Cattaneo--Felder--Tomassini \cite[\S3 and Corollary~5.2(iv)]{CFT} により
中心商の量子化を構成し、実際の局所化写像と section 復元を比較して
二次包絡の\emph{全次数}平坦性を得る。
さらに非 Noetherian な fpqc 被覆を含む線対象・Hom・合成の降下と
framed 自己同値の消滅により、stack 化後の形式群oidを記述する。
全次数平坦性やこの降下は、接空間の次元一致だけから導くものではない。

滑らかな三次の射影 Poisson 空間が $\operatorname{Gr}(2,5)$ となることは
Loray--Pereira--Touzet と Mayanskiy の既知結果である
\cite[Theorem~9.1]{LPT}\cite[Theorem~1 and \S2.2]{Mayanskiy}。
本稿では Schouten--Pl\"ucker 方程式を相対化し、残余射との両立まで示す。
Pym \cite[Theorem~2.3, Proposition~2.9 and Theorem~2.10]{Pym} の
量子化と Poisson 構造の芽の比較とは、対象となるモジュライと写像を区別する。
実際の開上の非標準的 Poisson 持上げ
（命題\ref{thm:noncanonical-poisson-cubic-lift}）は、
$\rho$ と相対超越基底から従う存在命題である。
その有理座標による論法自体や、既知の量子化比較を新規結果として扱わない。

\paragraph{点スキームと復元.}
Vancliff \cite[\S6, Proposition~6.17 and Theorem~6.18]{VancliffQuantumSpaces} は、
非零パラメータで埋込み点スキームが一定という仮定の下で Poisson 零点と比較する。
本稿の Hesse 境界はパラメータとともに変動する。
Le~Bruyn--Smith--Van~den~Bergh
\cite[Theorem~4.2.2 and Propositions~4.2.5,4.3.7]{LeBruynSmithVdBCentral}の
一中心変数の場合に対し、本稿では二中心変数と中心三次商を扱い、
長さ二十四の有限成分とその双線型関係を用いる。
偏微分空間からの三次復元の原理は Ueda--Yoshinaga
\cite[\S3, Lemma~3]{UedaYoshinagaJacobian} および Hwang
\cite[Theorems~1.5--1.6]{HwangSymmetrizer}と関連する。
本稿で取り出すデータは、二次元中心空間を保つ五次元 Casimir 空間と
尺度付き ambient Poisson 構造である。
Gualtieri--Pym \cite[Theorem~7.2]{GualtieriPym} の modular residue と
指定した量子化経路の点 shift の速度の関係は
命題\ref{prop:jacobian-modular-shift-comparison}に与える。

\subsection{構成・規約と残された問題}

第\ref{sec:intrinsic-cubic}--\ref{sec:abel-marked-moduli}節でモジュライの定義と
座標表示の同値関係を整え、次に明示的な中心三次族を構成する。
第\ref{sec:relative-koszul-line-cohomology}--\ref{sec:kuznetsov-moduli-morphism}節では、
相対双対、全 Abel 同値、代数的部分 stack、残余射の順に主定理A・Bを証明する。
第\ref{sec:moduli-artin}--\ref{sec:point-scheme-poisson}節では、
全 Artin 変形と点データからの復元を扱う。
周期性、非可換射影三空間との比較、有限中心模型に関する補助結果は付録にまとめる。
これらの結果の仮定は各節で明示し、主定理に用いる範囲を参照する。

以下では原則として $\C$ 上で議論し、一般の係数環と導来底変更の使用は各命題に明記する。
$\mathfrak C_{\mathrm J}$ の滑らかさはパラメータ stack の性質であり、
圏の dg smooth 性とは区別する。

未解決の中心問題は、量子中心直線に対応する $R_Q$ 内の古典三次直線を自然に選び、
それが $\rho$ の出力する三次に一致することを示すことである。
全モジュライ内での位置、七次元ファイバー、構造対象を忘れた復元については
第\ref{sec:status}節にまとめる。
本稿は公表準備稿であり、各結果の初出・新規性の網羅調査と全証明の独立監査は未完了である。
検証範囲は研究ノートに記す。

\clearpage
\part{構造対象付き圏とJacobian中心三次族}
\section{豊富列による内在的定義と族におけるambient復元}
\label{sec:intrinsic-cubic}

本節では，超曲面埋込みや次数シフトを入力に含めず，
圏と十分な対象列からambientを構成する．
「内在的」はこの\emph{対象列付きの圏}に関する意味である．
対象列を忘れたアーベル圏，導来圏，または残余圏のみからの
一意復元は本節の結論ではない．特に，自己同値$s$と対象$E$
から初めから$E_i=s^{-i}E$と置くと1周期構造を既に選んでしまう．
一般の$\mathbb Z$-代数のモジュライを扱うには，対象列を先に固定し，
その列を移す自己同値の存在を別に問う必要がある．

\subsection{対象列と二次包絡}

基礎体を$k$とする．$k$-線型アーベル圏$\mathcal C$内の列
$\mathcal E=(E_i)_{i\in\mathbb Z}$について，
$\dim_k\Hom(E_i,M)<\infty$をすべての$i,M$に仮定する．
本節の添字規約では，列が\emph{豊富}であるとは次の二条件をいう：
任意の全射$M\twoheadrightarrow N$に対して
$\Hom(E_i,M)\to\Hom(E_i,N)$は$i\gg0$で全射となり，
任意の$M$と整数$r$に対して，ある$i_1,\ldots,i_s\geq r$を用いた
全射$\bigoplus_{\alpha=1}^s E_{i_\alpha}\twoheadrightarrow M$が存在する．
これはPolishchukのample sequenceの規約を$i\mapsto-i$で読み替えたものである
\cite[Lemma~2.1, Theorem~2.4]{PolishchukProj}．
可換射影多様体の例では$E_i=\OO_X(-i)$である．

以下では$\End(E_i)=k$，$\Hom(E_j,E_i)=0$（$j<i$）を仮定し，
\begin{equation}\label{eq:intrinsic-section-z}
 \mathcal A_{ij}=\Hom_{\mathcal C}(E_j,E_i)\quad(i\leq j),
 \qquad
 \mathcal A_{ij}\otimes\mathcal A_{j\ell}\longrightarrow\mathcal A_{i\ell},
 \quad f\otimes g\longmapsto f\circ g
\end{equation}
と置く．これが本節で用いる正の連結$\mathbb Z$-代数の積の規約である．
豊富性により，$M\mapsto\bigoplus_j\Hom(E_j,M)$は
$\mathcal C\simeq\operatorname{cohproj}\mathcal A$を与える
\cite[Theorem~2.4]{PolishchukProj}．
noetherianな場合には右辺を$\qgr\mathcal A$と書く．
ここでは$\mathcal A$を最初からHomで定義しており，
一般の斉次座標環とそのsection algebraが低次数まで一致することを
無条件に仮定してはいない．

次数一生成の正の連結$\mathbb Z$-代数$\mathcal A$に対して
\begin{align}
 V_i&=\mathcal A_{i,i+1},\\
 R_i&=\ker\bigl(V_i\otimes V_{i+1}\longrightarrow\mathcal A_{i,i+2}\bigr),
 \label{eq:intrinsic-quadratic-relations}\\
 Q(\mathcal A)&=T_{\mathbb Z}(V)/(R_i\mid i\in\mathbb Z)
 \label{eq:intrinsic-quadratic-envelope}
\end{align}
と定める．ここで
$T_{\mathbb Z}(V)_{ij}=V_i\otimes\cdots\otimes V_{j-1}$（$i<j$），
$T_{\mathbb Z}(V)_{ii}=k$であり，括弧は両側イデアルを表す．
積から標準全射$q:Q(\mathcal A)\twoheadrightarrow\mathcal A$が定まり，
これは長さ$0,1,2$の成分で同型である．
この構成にambientの選択は含まれない．
対象列付き圏の同値には，各$i$について$F(E_i)\simeq E'_i$という
両立する同定を含める．これによりHomと合成の同型，従って
$\mathcal A\simeq\mathcal A'$および$Q(\mathcal A)\simeq Q(\mathcal A')$
が誘導される．同定を変更したときの基底変更も同型として扱い，
対象列を保たない圏同値まで同じモジュライ問題に含めない．

\begin{definition}[豊富列付きの正規三次データ]
\label{def:intrinsic-cubic-datum}
組$(\mathcal C,\mathcal E)$が\emph{正規三次三様体データ}であるとは，
上の条件に加え，そのsection $\mathbb Z$-代数$\mathcal A$が次数一生成であり，
次を満たすことをいう．
\begin{enumerate}[label=\textup{(\roman*)}]
\item $\mathcal A$と$\mathcal B:=Q(\mathcal A)$は左右noetherianであり，
$\mathcal B$は左右のAS条件を満たす二次Koszul AS正則$\mathbb Z$-代数である．
その次元とGorenstein parameterはともに$5$，
最小線型解消の型は$(1,5,10,10,5,1)$とする．
\item すべての$i$と$m\geq0$について
\begin{equation}\label{eq:intrinsic-cubic-hilbert}
 \dim_k\mathcal B_{i,i+m}=b_m:=\binom{m+4}{4},\qquad
 \dim_k\mathcal A_{i,i+m}=h_m:=b_m-b_{m-3},
\end{equation}
ただし$m<0$では$b_m=0$とする．
\item $I=\ker q$の一次元空間$I_{i,i+3}$の任意の非零元$g_i$は左右正則である．
すなわち，$g_i\mathcal B_{i+3,j}$および$\mathcal B_{\ell i}g_i$を定める
乗法写像は，それぞれ$j\geq i+3$，$\ell\leq i$について単射である．
\end{enumerate}
さらに$D^b(\mathcal C)$がsmooth properな$k$-線型dg enhancementを持つとき，
このデータを\emph{滑らかな}正規三次データと呼ぶ．
この滑らかさは独立の条件であり，Hilbert級数やambientのAS正則性だけからは結論しない．
\end{definition}

この定義は，圏のHomと合成から計算できる\emph{埋込みの選択を要しない認識条件}である．
しかし，より弱い数値的・Serre関手的条件を満たす抽象圏すべてが
この定義を満たすと証明したわけではない．
また中心元による三次商は正規三次データの特殊な場合であり，
中心性は対象列に周期構造を選んだ後に初めて通常の次数付き表示で比較できる．

\begin{theorem}[内在的な正規超曲面認識]
\label{thm:intrinsic-normal-recognition}
定義\ref{def:intrinsic-cubic-datum}のデータに対して，
任意の選択$0\ne g_i\in I_{i,i+3}$は正則正規族をなし，
\begin{equation}\label{eq:intrinsic-normal-presentation}
 \mathcal A\simeq\mathcal B/(g_i\mid i\in\mathbb Z),
 \qquad \mathcal B=Q(\mathcal A)
\end{equation}
となる．$\mathcal A$への商写像と両立する二次ambientは一意である．
より正確には，二次代数$\mathcal B'$から$\mathcal A$への全射が
長さ$1,2$で同型なら，その商写像と両立する
$\mathcal B'\simeq Q(\mathcal A)$が一意に存在する．
\end{theorem}
\begin{proof}
Hilbert関数の差から
$\dim I_{i,i+3}=1$，$\dim I_{i,i+4}=5$である．
左右正則性により
\begin{equation}\label{eq:intrinsic-degree-four-normal}
 g_i\mathcal B_{i+3,i+4}
   =I_{i,i+4}
   =\mathcal B_{i,i+1}g_{i+1}
\end{equation}
が成り立つ．実際，左右の空間はともに$I_{i,i+4}$に含まれる5次元空間である．
次数一生成と帰納法により
$g_i\mathcal B_{i+3,j+3}=\mathcal B_{ij}g_j$がすべての$j\geq i$について従う．
従って$(g_i)$は正則正規族である．

$J$をこの族が生成する両側イデアルとする．正規性により
$J_{i,i+m}=g_i\mathcal B_{i+3,i+m}$（$m\geq3$）であり，
正則性からその次元は$b_{m-3}$となる．一方，
\eqref{eq:intrinsic-cubic-hilbert}から
$\dim I_{i,i+m}=b_{m-3}$である．$J\subseteq I$なので全成分で$J=I$を得る．
従って，三次での核に加えて高次の独立な方程式が残る可能性もここで排除される．

最後に，$\mathcal B'$は次数一生成の二次代数なので，
長さ一成分と，隣接する長さ一成分の積の核とで定まる．
仮定によりこれらは\eqref{eq:intrinsic-quadratic-relations}の$V_i,R_i$そのものである．
従って\eqref{eq:intrinsic-quadratic-envelope}が所望の同型を与える．
同型は生成元上で既に決まるので一意である．
\end{proof}

\begin{remark}[周期性との関係]
\label{rem:intrinsic-periodicity-scope}
第\ref{sec:moduli-ambient-periodicity}節の系\ref{cor:moduli-cubic-ambient-periodic}により，
定理\ref{thm:intrinsic-normal-recognition}で得られる体上の対$(\mathcal B,I)$は，
三次イデアルを保つ1周期同型を持つ．従って$\mathcal A$も
一つの通常の次数付き正規三次商の付随$\mathbb Z$-代数として表示できる．
これは正規三次データに対する\emph{結論}であり，定義に組み込んだ条件ではない．
この証明は正規族に沿った有限頂点orbit代数とNakayama双対性を使う．
一方，1周期構造の選択は一意とは限らず，次数3の関係元の中心性や
族全体の周期構造の降下は別途検証を要する．
\end{remark}

\subsection{任意の底変更と両立する構成}

閉点の同型類ごとの復元だけでは，モジュライ間の射はまだ得られない．
必要なのは族での構成と底変更への両立である．以下では圏の族の存在論を避け，
相対section代数が既に与えられている状況を正確に扱う．

\begin{theorem}[二次包絡と三次イデアルの相対復元]
\label{thm:intrinsic-relative-envelope}
$S$を$k$-スキームとし，$\mathcal A$を正の連結
$\OO_S$-$\mathbb Z$-代数とする．各$\mathcal A_{ij}$は有限階数局所自由，
$\mathcal A$は次数一生成と仮定する．すると
\eqref{eq:intrinsic-quadratic-relations}--\eqref{eq:intrinsic-quadratic-envelope}
は$S$上で定義され，任意の$f:T\to S$について標準同型
\begin{equation}\label{eq:intrinsic-base-change}
 f^*Q(\mathcal A)\simeq Q(f^*\mathcal A)
\end{equation}
を持つ．ここで$f$の平坦性は不要である．

さらに$\mathcal B=Q(\mathcal A)$の各成分も局所自由で，その階数と
$\mathcal A$の階数が\eqref{eq:intrinsic-cubic-hilbert}に等しく，
すべての幾何学的ファイバーにおける$I_{i,i+3}$の非零元が左右正則とする．
このとき$L_i:=I_{i,i+3}$は直線束であり，
\begin{equation}\label{eq:intrinsic-relative-ideal}
 I_{ij}=\operatorname{im}\bigl(L_i\otimes\mathcal B_{i+3,j}
                  \longrightarrow\mathcal B_{ij}\bigr)\quad(j\geq i+3).
\end{equation}
この写像は$I_{ij}$への同型である．また
\begin{equation}\label{eq:intrinsic-relative-normal}
 \operatorname{im}\bigl(L_i\otimes\mathcal B_{i+3,j+3}
                   \longrightarrow\mathcal B_{i,j+3}\bigr)
 =\operatorname{im}\bigl(\mathcal B_{ij}\otimes L_j
                   \longrightarrow\mathcal B_{i,j+3}\bigr).
\end{equation}
従って$I$はこの三次直線束族で生成され，$I$と$L_i$も任意の底変更と両立する．
\end{theorem}
\begin{proof}
次数一生成により$V_i\otimes V_{i+1}\to\mathcal A_{i,i+2}$は全射である．
終項が局所自由なので，その核$R_i$は局所自由であり，
この短完全列は局所分裂する．従って任意の底変更で核が保たれる．
長さ$m$のテンソル代数成分は$m$個の$V_i$のテンソル積である．
二次関係が生成するイデアルの同じ成分は，
その隣接する二因子を$R_j$に置き換えた有限個のテンソル積からの写像の像である．
よって$Q(\mathcal A)_{i,i+m}$はこの有限の写像の余核として表される．
テンソル積と余核はいずれも底変更と両立するので
\eqref{eq:intrinsic-base-change}を得る．

$\mathcal B\twoheadrightarrow\mathcal A$の各成分の終項も局所自由なので，
$I_{ij}$は局所自由であり，核の形成は任意の底変更と両立する．
特に$L_i$と$I_{i,i+4}$の階数はそれぞれ$1,5$である．
乗法による二写像
\begin{equation}
 L_i\otimes\mathcal B_{i+3,i+4}\longrightarrow I_{i,i+4},\qquad
 \mathcal B_{i,i+1}\otimes L_{i+1}\longrightarrow I_{i,i+4}
\end{equation}
は，すべての幾何学的ファイバーで単射となる同階数のベクトル束間の写像である．
従ってそれぞれ同型である．この事実を次数一生成と帰納法に適用すると，
\eqref{eq:intrinsic-relative-normal}を得る．

\eqref{eq:intrinsic-relative-ideal}の写像を$I_{ij}$への写像と見ると，
始項と終項の階数はいずれも$b_{j-i-3}$である．
すべての幾何学的ファイバーで正則性により単射なので，やはり同型である．
従って三次直線束族が$I$全体を生成する．
以上の構成はテンソル積・積・局所分裂した核・余核からなるため，
主張したすべての底変更同型は自然であり，底変更の合成とも両立する．
\end{proof}

\begin{remark}[族についての仮定を省略しない]
$Q(\mathcal A)$の各成分の局所自由性は上の定理の追加仮定である．
有限個の低次数のrankや個々の閉点での数値のみから，これを自動的に導いたわけではない．
圏の族から出発する場合も，相対Homが局所自由であり底変更と両立することを
別途確かめてから本定理を使う．また$L_i$は一般には非自明な直線束なので，
大域的な元$g_i$の選択は不要であり，その存在も主張しない．
ファイバーごとの1周期性から，族全体の1周期同型やその降下も自動的には従わない．
\end{remark}

\subsection{モジュライ関手と「一意」の意味}

上の相対条件と幾何学的ファイバーでのAS条件を課した
$\mathcal A$の族のファイバー群oidを$\mathfrak C^{\rm seq}_3$，
標準5次元二次AS正則$\mathbb Z$-代数の族のファイバー群oidを
$\mathfrak P^{\rm AS}_4$と書く．射は添字を保つ代数同型とする．
定理\ref{thm:intrinsic-relative-envelope}は標準的な射
\begin{equation}\label{eq:intrinsic-moduli-morphism}
 \Phi:\mathfrak C^{\rm seq}_3\longrightarrow\mathfrak P^{\rm AS}_4,
 \qquad\mathcal A\longmapsto Q(\mathcal A)
\end{equation}
を与える．必要なら両辺をfpqc stack化して同じ射を得る．
ここで\emph{代数的stackであること，局所有限表示性，有限型atlasの存在は主張しない}．
これらは無限個の成分・AS条件・noetherian性を有限の表示条件に置き換える別の課題である．
また，この群oidはまず相対座標代数のモジュライである．
圏の族のモジュライと同一視する際には，相対Homの底変更と，
座標代数から相対section代数への写像のsaturationを別途確認する．

\begin{proposition}[商データとの同値とambient忘却]
\label{prop:intrinsic-pair-groupoid}
$\mathfrak C^{\rm seq}_3$は，条件を満たす二次ambientと三次直線束族による
正規商の組$(\mathcal B,I)$のファイバー群oidに同値である．
同値は$\mathcal A\mapsto(Q(\mathcal A),\ker q)$で与えられる．
特にすべての$\mathcal A$について
\begin{equation}\label{eq:intrinsic-automorphisms}
 \Aut(\mathcal A)\simeq\Aut(Q(\mathcal A),\ker q)
 \lhook\joinrel\longrightarrow\Aut(Q(\mathcal A))
\end{equation}
となる．しかし$\Phi$は一般には同型類上で単射ではない．
\end{proposition}
\begin{proof}
正規商は長さ$1,2$の成分を変えないため，その二次包絡は元の二次ambientに戻る．
逆方向は定理\ref{thm:intrinsic-normal-recognition}と
定理\ref{thm:intrinsic-relative-envelope}で与えられる．
商の同型は長さ一成分上で同型を与え，テンソル代数の普遍性から
二次包絡の同型へ一意に延長する．積の核を保つので$I$も保存する．
逆に$(\mathcal B,I)$の同型は商に降りる．これらは互いに逆で，
底変更と両立するから群oidの同値となる．

最後の主張は可換例だけでも分かる．
すべての滑らかな三次超曲面$X\subset\PP^4$について，
$\mathcal A=\check{(k[x_0,\ldots,x_4]/(F))}$の二次包絡は
同じ$\check{k[x_0,\ldots,x_4]}$である．
従って，ambientの一意性は三次方程式の一意性を意味しない．
\end{proof}

具体的な$\mathcal B$への同型まで指定した$\Phi$のファイバーは，
その$\mathcal B$内の許容三次イデアルのパラメータ問題である．
ambientの同型を指定しなければさらに$\Aut(\mathcal B)$で割る．
したがって，像の次元を求めるにはソースの次元だけでなく，
一般ambient上の三次イデアルの変動と安定化群も計算しなければならない．
既存の17次元族の像が17次元である，あるいはambient moduliの既約成分である，
という結論は\eqref{eq:intrinsic-moduli-morphism}だけからは得られない．

\section{二つのモジュライstackと正規三次incidence}
\label{sec:algebraic-moduli-foundations}

本節では、まず添字を保つ代数同型を射とするモジュライ問題を定義する。
圏の同値における自然同型、対象列の同定の忘却、周期構造の忘却は、
この定義を変える操作であり、後続の節で別々に扱う。
基礎体は$k=\C$とし、stackは$k$-スキームのfpqc site上で考える。
有限型代数的stackの構成と、無限個の次数にわたるモジュライ問題の定義を区別する。

\subsection{固定Hilbert関数を持つ代数のstack}

$d=(d_m)_{m\geq0}$を$d_0=1$を満たす正整数列とする。
$\mathfrak{Alg}_{\mathbb Z}(d)(S)$の対象は、
\begin{equation}\label{eq:moduli-z-family}
 A_{ii}=\OO_S,\quad A_{ij}=0\ (j<i),\quad
 \rank A_{ij}=d_{j-i}\ (j>i),\qquad
 A_{ij}\otimes A_{j\ell}\longrightarrow A_{i\ell}
\end{equation}
という有限階数局所自由層と、単位律・結合律を満たす積の族である。
射は各$A_{ii}=\OO_S$上で恒等となる、添字を保つ代数同型とする。
プルバックは各成分の通常の層のプルバックで定める。
この時点では、次数一生成、noetherian性、AS正則性を課さない。

\begin{proposition}[有限窓による厳密な表示]
\label{prop:moduli-finite-window}
$\mathfrak{Alg}_{\mathbb Z}(d)$はfpqc stackである。
$J_N=\{-N,\ldots,N\}$への制限のstackを$\mathfrak{Alg}_{J_N}(d)$と書くと、
\begin{equation}\label{eq:moduli-window-limit}
 \mathfrak{Alg}_{\mathbb Z}(d)
 \simeq \mathop{2\!\varprojlim}_{N}\mathfrak{Alg}_{J_N}(d),
 \qquad
 \mathfrak{Alg}_{J_N}(d)=[X_N/G_N]
\end{equation}
である。ここで$X_N$は有限型affine scheme、$G_N$は有限個の一般線型群の積であり、
各有限窓のstackは有限型代数的stackである。
\end{proposition}
\begin{proof}
$i<j$に対し$W_{ij}=k^{d_{j-i}}$を取り、
\[
 M_N=\prod_{\substack{i<j<\ell\\ i,j,\ell\in J_N}}
 \Hom_k(W_{ij}\otimes W_{j\ell},W_{i\ell}),\qquad
 G_N=\prod_{\substack{i<j\\ i,j\in J_N}}\GL(W_{ij})
\]
と置く。$X_N\subset M_N$を、すべての$i<j<\ell<m$について
$\mu_{i\ell m}(\mu_{ij\ell}\otimes1)
=\mu_{ijm}(1\otimes\mu_{j\ell m})$を課した閉部分schemeとする。
長さ零の積は単位で指定済みなので、それ以外の式は不要である。
基底変更の作用は
$(g\mu)_{ij\ell}=g_{i\ell}\mu_{ij\ell}(g_{ij}^{-1}\otimes g_{j\ell}^{-1})$
である。有限階数ベクトル束のframe torsorを取る通常の降下により、
$[X_N/G_N]$は有限窓の代数族をちょうど分類する。

ベクトル束とその間の射はfpqc降下を満たし、結合律は降下後にも成り立つ。
この議論を成分ごとに行えば、無限個の添字についても降下は有効である。
さらに各成分、各積、各結合律はある有限窓に含まれる。
従って、制限同型を含む有限窓の両立する系から、全成分と積を一意に貼り合わせられる。
射についても同じなので\eqref{eq:moduli-window-limit}を得る。
\end{proof}

これは、次数付き代数に対するBehrend--Noohiの乗法・基底変更群・切断の構成
\cite[\S\S2.3, 2.5]{BehrendNoohiModuli}を、添字を持つ小さい線型圏に適用したものである。
有限窓をrelationの空間として記述する方法は、
Abdelgadir--Okawa--Uedaのquiver relationのモジュライ
\cite[\S\S3--5]{AOUNCPlanes}とも整合する。
これらの構成自体を新規の結果として主張しない。
また、\eqref{eq:moduli-window-limit}から全stackの有限型性や
有限窓一つによる決定性は従わない。

\begin{definition}[ambientと三次のモジュライ]
\label{def:moduli-two-stacks}
$b_m=\binom{m+4}{4}$、$h_m=b_m-b_{m-3}$とする。
以下の二つを本稿の\emph{厳密な座標代数のモジュライ}と定める。
\begin{enumerate}[label=\textup{(\roman*)}]
\item $\mathfrak P^{\rm AS}_4\subset\mathfrak{Alg}_{\mathbb Z}(b)$は、
次数一生成の相対二次代数で、すべての幾何学的ファイバーが
定義\ref{def:intrinsic-cubic-datum}の標準5次元Koszul AS正則条件を満たすもののfull substackである。
左右noetherian性、AS index $5$、左右のglobal dimensionとuniform AS Ext条件も含める。
相対二次代数とは$T_{\mathbb Z}(B_{i,i+1})/(R_i)\to B$が同型であることをいう。
\item $\mathfrak C^{\rm seq}_3\subset\mathfrak{Alg}_{\mathbb Z}(h)$は、
次数一生成の$A$で、$Q(A)\in\mathfrak P^{\rm AS}_4$、かつ三次核の非零元が
すべての幾何学的ファイバーで左右正則となるもののfull substackである。
\end{enumerate}
\end{definition}

各条件はfpqc局所的であり、二次包絡は任意の底変更に可換であるため、
これは実際に二つのfpqc stackを定める。
ここでnoetherian性などは明記した幾何学的ファイバー上の条件である。
一般の基底に対する相対$qgr$の存在論やcoherenceを、
ファイバーの条件だけから結論していない。
圏を用いたモジュライとの比較には第\ref{sec:categorical-moduli-comparison}節の条件を追加する。
滑らかな三次のみを扱うときも、相対的なsmooth proper enhancementの条件を別に課す。
その部分が開であることは、本定義からは従わない。

\subsection{比較射と全ファイバーの定義}

固定した$S$上の$B\in\mathfrak P^{\rm AS}_4(S)$に対し、
$\operatorname{Inc}^{\rm nor}_3(B)(T)$を、$B_T$内の直線部分束
\[
 L_i\subset (B_T)_{i,i+3}\quad(i\in\mathbb Z)
\]
で次を満たすものの集合とする。部分束の語は余商も局所自由であることを含む。
まず、左右の乗法
\begin{equation}\label{eq:moduli-normal-subbundles}
 L_i\otimes (B_T)_{i+3,j}\longrightarrow (B_T)_{ij},\qquad
 (B_T)_{\ell i}\otimes L_i\longrightarrow (B_T)_{\ell,i+3}
\end{equation}
はすべての$j\geq i+3$、$\ell\leq i$について部分束への同型とする。
さらに、長さ四で
\begin{equation}\label{eq:moduli-adjacent-normality}
 L_i(B_T)_{i+3,i+4}=(B_T)_{i,i+1}L_{i+1}
 \ \subset (B_T)_{i,i+4}
\end{equation}
を課す。直線の生成元、またはそれらの自明化はデータに含めない。

\begin{theorem}[二つのモジュライの比較とnormal incidence]
\label{thm:moduli-incidence-equivalence}
自然な射
\begin{equation}\label{eq:moduli-main-arrow}
 \Phi:\mathfrak C^{\rm seq}_3\longrightarrow\mathfrak P^{\rm AS}_4,
 \qquad A\longmapsto Q(A)
\end{equation}
があり、$B:S\to\mathfrak P^{\rm AS}_4$による2-fiber productは
\begin{equation}\label{eq:moduli-full-fibre}
 S\mathop{\times}_{\mathfrak P^{\rm AS}_4}\mathfrak C^{\rm seq}_3
 \simeq\operatorname{Inc}^{\rm nor}_3(B)
\end{equation}
というfpqc集合層である。特に、ambientとの同型を指定したファイバーに
追加の自己同型はない。
一方、代数閉体上でambientの同型類だけを固定した群oidは
$\operatorname{Inc}^{\rm nor}_3(B)(k)//\Aut_0(B)(k)$である。
ambientがfpqc局所的に$B$と同型である族を含めたstackは
$[\operatorname{Inc}^{\rm nor}_3(B)/\Aut_0(B)]$である。
ここで$\Aut_0$は添字を保つ代数自己同型を表す。
\end{theorem}
\begin{proof}
三次側からは、定理\ref{thm:intrinsic-relative-envelope}により
$L_i=\ker(Q(A)_{i,i+3}\to A_{i,i+3})$を得る。
同定$Q(A)\simeq B_T$を用いて$B_T$内の部分束と見れば、
\eqref{eq:moduli-normal-subbundles}--\eqref{eq:moduli-adjacent-normality}を満たす。

逆に、直線束族から$J_{ij}=L_i(B_T)_{i+3,j}$（$j-i\geq3$）、
その他の成分では$J_{ij}=0$と置く。
これは右イデアルである。次数一生成と\eqref{eq:moduli-adjacent-normality}を繰り返すと
$J_{ij}=(B_T)_{i,j-3}L_{j-3}$も成り立つので、左イデアルでもある。
\eqref{eq:moduli-normal-subbundles}から$J_{ij}$は階数$b_{j-i-3}$の部分束である。
従って$A=B_T/J$の各成分は階数$h_{j-i}$のベクトル束となる。
長さ一・二を変えないため、$Q(A)\simeq B_T$である。
また各ファイバーで正規正則三次商となり、noetherian性もambientから従う。

両構成は任意の底変更に可換で互いに逆である。
ambient上で恒等となる商の自己同型は、次数一生成により恒等である。
従って2-fiber productは集合層になる。
ambientの同定を忘れれば、その変更を$\Aut_0(B)$で割ることになる。
\end{proof}

\begin{proposition}[ファイバーの有限窓での方程式]
\label{prop:moduli-incidence-windows}
各$B/S$に対し、有限表示の準射影$S$-scheme $I_N(B)$があり、
\begin{equation}\label{eq:moduli-incidence-limit}
 \operatorname{Inc}^{\rm nor}_3(B)=\varprojlim_N I_N(B)
\end{equation}
がfpqc集合層として成り立つ。
\end{proposition}
\begin{proof}
まず$J_N$に$i,i+3$が含まれる$i$について、有限積
$\prod_i\PP_S(B_{i,i+3})$を取る。ここでは$\PP$は直線部分束の規約である。
この積には普遍直線束がある。
窓内に含まれる\eqref{eq:moduli-normal-subbundles}の有限個の写像について、
すべてが最大rankとなる開部分を取る。
この開部分では各像と余商はベクトル束である。
そこで\eqref{eq:moduli-adjacent-normality}を、例えば一方の像から他方の余商への
合成が零となる閉条件として課す。
同階数の二つの部分束の包含は等号となるため、これが必要十分である。
こうして得た局所閉部分schemeを$I_N(B)$とする。
窓の制限はこれらの間の射を定める。
すべての直線、正則性の条件、長さ四の条件はある有限窓に現れるので、
その逆極限は所望の集合層である。
\end{proof}

式\eqref{eq:moduli-incidence-limit}は、比較射のファイバーを具体的な方程式で調べる入口になる。
しかし遷移射がaffineであるとは示していないので、ここではこの極限がschemeであるとは主張しない。
同様に、AS条件を有限窓一つで認識できる定理、族の有界性、適切な安定性を証明するまでは、
全モジュライを有限型代数的stackまたはGIT quotientと呼ばない。

\subsection{一次変形と可換点での基準例}

\begin{proposition}[固定ambient上の接空間]
\label{prop:moduli-normal-tangent}
体上の点$I\subset B$における$\operatorname{Inc}^{\rm nor}_3(B)$の接空間は
\begin{equation}\label{eq:moduli-normal-tangent}
 T_I\operatorname{Inc}^{\rm nor}_3(B)
   =\Hom_{B\text{-}B,\mathbb Z}(I,B/I)
\end{equation}
である。右辺では両方の添字を保つ双加群準同型を取る。
\end{proposition}
\begin{proof}
$k[\epsilon]/(\epsilon^2)$上のflatなイデアル変形は、各成分の補空間を選ぶと
$a+\epsilon\widetilde\eta(a)$（$a\in I$）と$\epsilon I$で生成される。
$\eta(a)=\widetilde\eta(a)\bmod I$は補空間に依存しない。
これが両側イデアルとなる条件は
$\eta(bac)=b\eta(a)c$であり、これが右辺の双加群条件である。
逆に、そのような$\eta$のgraphは成分ごとにflatな商を与える。
三次より低い成分は変わらず、三次核は直線束、全核の階数は$b_{m-3}$のままである。
左右の正則性の乗法は特殊ファイバーで単射なので、Artin基底上で部分束を与える。
定理\ref{thm:intrinsic-relative-envelope}を適用するとnormal incidenceの点を得る。
\end{proof}

周期を選んだ次数付き問題では右辺も次数零のgraded双加群準同型に置き換える。
例えば中心正則元$G\in B_3$による商では、それは$\Z(B/(G))_3$になる。
中心三次直線だけを動かす部分問題の接空間は$\Z(B)_3/kG$であり、
両者を無条件に同じ空間とはしない。
全normal incidenceと、特定の次数付き表示でのcentral incidenceは異なる定義である。

\begin{proposition}[可換ambientのファイバー]
\label{prop:moduli-commutative-fibre}
$S_0=\Sym_k V$、$\dim V=5$とする。可換ambient$\check S_0$の厳密なファイバーは
\begin{equation}\label{eq:moduli-commutative-fibre}
 \operatorname{Inc}^{\rm nor}_3(\check S_0)
       \simeq\PP(\Sym^3 V)=\PP^{34}
\end{equation}
であり、これは非被約な基底も含む関手の同型である。
頂点スカラーを除いたambient自己同型群は$\PGL(V)$である。
従って第\ref{sec:categorical-moduli-comparison}節の一点の同定を保持するlevelでは、
$B\PGL(V)$という可換ambientの中立部分stackへの底変更は
$[\PP^{34}/\PGL(V)]$となる。
滑らかな可換三次の部分は判別式の補集合を取った通常の10次元モジュライである。
\end{proposition}
\begin{proof}
基底上で隣接する二つの三次直線を局所的に$F,G$で生成する。
normalityは$F V=G V$を意味する。
両辺はrank $5$の部分束なので、基底$x_0,\ldots,x_4$について
$F x_j=G\ell_j$、$\ell_j\in V\otimes\OO_S$となり、$x_j\mapsto\ell_j$は可逆である。
交換律と$G$の正則性から$x_i\ell_j=x_j\ell_i$を得る。
二次単項式の係数を比較すれば、すべての$j$で$\ell_j=\lambda x_j$となる。
可逆性から$\lambda$は単元であり、$F=\lambda G$である。
従ってすべての$L_i$は同じ三次直線に等しい。
逆に、任意の三次直線部分束は各ファイバーで非零多項式を与えるため、
全乗法写像が部分束を与え、必要なnormalityも満たす。
以上は任意の基底に対する議論であり、\eqref{eq:moduli-commutative-fibre}を得る。

最後に、$\check S_0$の自己同型を長さ一成分の$g_i\in\GL(V)$で記述する。
二次関係を保つ条件は$(g_i\otimes g_{i+1})\wedge^2 V=\wedge^2 V$である。
$(1\otimes h)\wedge^2 V=\wedge^2 V$を満たす$h$は、同じ二次係数比較によりスカラーである。
従ってすべての$g_i$の射影類は等しい。
各矢印のスカラーは頂点スカラーの比として表せるので、それらを除く群は$\PGL(V)$となる。
滑らかな三次では射影自己同型群は有限であり、次元は$34-24=10$である。
この有限性は$H^0(X,T_X)=0$からも従う。
\end{proof}

この基準例により、ambient復元の一意性と比較射の単射性は別の性質であることが、
モジュライの定義そのものから確認できる。
可換部分の10次元すべてが同じambientの点へ写る一方、
後述の非可換Jacobian族では一般ambientも13次元分だけ動く。

\section{豊富列付き圏と座標代数のモジュライの比較}
\label{sec:categorical-moduli-comparison}

第\ref{sec:intrinsic-cubic}節の復元をモジュライの比較に用いるには，
対象の同型類だけでなく，射とその自然同型を指定しなければならない．
本節では，対象列の各項の同定を保持する場合，その同定を忘れる場合，
および$E_0$の同定だけを保持する場合を区別する．
以下の比較は，明記するsaturationと相対復元の仮定の下での定理である．
全AS正則代数に対するこれらの仮定や，得られるstackの代数性を
ここで自動的に結論するものではない．

\subsection{全対象の同定を保持する比較}

まず体$k$上で考える．正の連結$\mathbb Z$-代数$A$に対し，
$P_i=e_iA$，$\overline P_i=\pi P_i\in\qgr A$と書く．
本節で\emph{飽和した座標代数}とは，$A$が局所有限・右noetherianで，
$(\overline P_i)$が第\ref{sec:intrinsic-cubic}節の意味で豊富であり，
$\qgr A$における必要なHom空間が有限次元，かつ標準写像
\begin{equation}\label{eq:categorical-saturation}
 A_{ij}\longrightarrow
 \Hom_{\qgr A}(\overline P_j,\overline P_i)
\end{equation}
がすべての$i,j$について同型であるものをいう．
$i>j$では左辺を$0$とする．特に，これは低次数を含む条件であり，
高次数での同型だけでは代用しない．本稿のモジュライ部分問題では
さらに$\dim A_{i,i+1}=5$を課す．

$\mathcal C$を$k$-線型Abel圏，$(E_i)$を豊富列とし，
$\End(E_i)=k$，$\Hom(E_j,E_i)=0$（$i>j$）を仮定する．
section代数は$A_{ij}=\Hom(E_j,E_i)$である．
Polishchukの復元定理\cite[Theorem~2.4]{PolishchukProj}は
この圏を$\operatorname{cohproj}A$として復元する．
以下では$A$が右noetherianである部分問題に制限し，これを$\qgr A$と書く．
逆向きに，任意の代数$A$から始めて\eqref{eq:categorical-saturation}を
無条件に結論することはしない．

\begin{definition}[全同定付きの2群oid]
\label{def:categorical-marked}
$\mathbf{Cat}^{\mathrm{mark}}(k)$の対象は，上記の条件を満たす
$(\mathcal C,(E_i))$である．1射は，$k$-線型圏同値
$F:\mathcal C\to\mathcal C'$と同型
\[
 \alpha_i:F(E_i)\xrightarrow{\sim}E'_i\qquad(i\in\mathbb Z)
\]
の組とする．$(F,\alpha)$から$(G,\beta)$への2射は自然同型
$\eta:F\Rightarrow G$であって
$\beta_i\eta_{E_i}=\alpha_i$をすべての$i$について満たすものとする．
添字の平行移動はこの定義の射に含めない．
\end{definition}

同定$\alpha_i$を保持することは，$\Hom(E_j,E_i)$の基底を選ぶこととは異なる．
ここで保持しているのは，圏同値の下でどの対象をどの対象に同定するかというデータである．

\begin{theorem}[全同定付き圏と厳密な代数同型]
\label{thm:categorical-marked-equivalence}
section代数の構成は，$\mathbf{Cat}^{\mathrm{mark}}(k)$から，
飽和した座標代数と添字を保つ代数同型の群oidへの同値を与える．
左辺では各Hom群oidの2同型類を取る．より強く，
各Hom群oidは離散群oidと同値である．1射の像は
\begin{equation}\label{eq:categorical-induced-map}
 \phi_{ij}(a)=\alpha_iF(a)\alpha_j^{-1}.
\end{equation}
従って，同じ$\phi$を誘導する二つの1射の間には，
同定と両立する2射が一意に存在する．
\end{theorem}
\begin{proof}
式\eqref{eq:categorical-induced-map}は単位と合成を保つ．
$\beta_i\eta_{E_i}=\alpha_i$と自然性を用いると，2同型な1射の像は等しい．
逆に，添字を保つ代数同型$\phi:A\xrightarrow{\sim}A'$は
加群の構造の移送を与え，torsion加群を保つため$\qgr$の同値を誘導する．
これは$\overline P_i$を$\overline P'_i$へ標準的に同定する．
式\eqref{eq:categorical-saturation}により，そのsection代数上の作用は
もとの$\phi$である．

ここで，自然変換の一意性だけでなく存在も確認する．
$\Gamma_{\mathcal E}(M)_j=\Hom(E_j,M)$と置く．
1射$(F,\alpha)$は，各$j$の写像
\[
 \Hom(E'_j,F(M))\longrightarrow\Hom(E_j,M),
 \qquad u\longmapsto F^{-1}(u\alpha_j)
\]
により，$\Gamma_{\mathcal E'}F$と
$\Gamma_{\mathcal E}$の$\phi$に沿った構造移送との自然同型を与える．
この写像は合成の定義から加群構造と両立する．
十分高い次数への切断を取り，その像を$\qgr$で考えれば，
復元同値により$F$は$\phi$が誘導する関手へ自然に同型となる．
この同型は$E_i$上で指定された$\alpha_i$と両立する．
従って同じ$\phi$を与える二つの1射の間に所望の2射が存在する．

その一意性は豊富性からも直接分かる．
二つの2射がすべての$E_i$上で一致するとき，任意の$M$への全射
$\bigoplus E_{i_a}\twoheadrightarrow M$に自然性を適用する．
圏同値は全射を保つので，二つの自然変換は$M$上でも一致する．
特に1射の2自己同型は恒等変換に限る．
最後に，飽和した$A$は$(\qgr A,(\overline P_i))$から復元され，
豊富列付き圏はPolishchukの定理で復元される．
以上で対象とHom群oidの両方について同値を得る．
\end{proof}

\subsection{頂点スカラーと一点を保持する比較}

次に，1射には圏同値$F$だけを含め，
各$i$について$F(E_i)\simeq E'_i$であることを条件とし，
具体的な同定$\alpha_i$を含めない．2射はすべての自然同型とする．
この2群oidを$\mathbf{Cat}^{\mathrm{weak}}(k)$と書く．
別に$\mathbf{Cat}^{\mathrm{pt}}(k)$では，1射に
$\alpha_0:F(E_0)\simeq E'_0$だけを含め，2射には
$\beta_0\eta_{E_0}=\alpha_0$を課す．
いずれも列自体を忘れた圏のモジュライではない．

\begin{proposition}[スカラーgaugeの正確な作用]
\label{prop:categorical-scalar-gauge}
$T(k)=(k^\times)^{\mathbb Z}$，
$T_0(k)=\{(\lambda_i)\in T(k)\mid\lambda_0=1\}$と置く．
頂点スカラーは
\begin{equation}\label{eq:categorical-vertex-gauge}
 \rho_\lambda(a_{ij})=\lambda_i\lambda_j^{-1}a_{ij}
\end{equation}
で添字を保つ代数自己同型を定める．
$\mathbf{Cat}^{\mathrm{weak}}(k)$は，対象を飽和した代数$A$，
1射を厳密な代数同型$\phi:A\to A'$，2射を
\begin{equation}\label{eq:categorical-gauge-two-arrow}
 \lambda:\phi\Rightarrow\psi,
 \qquad \psi=\rho_\lambda\phi
\end{equation}
とする2群oidと同値である．
$\mathbf{Cat}^{\mathrm{pt}}(k)$では同じ記述で$\lambda_0=1$を課す．

各$A_{i,i+1}\ne0$の下で，$\ker\rho$は対角$k^\times$であり，
$\rho(T_0)=\rho(T)$，$\rho|_{T_0}$は単射である．
従って弱い自己同値の2同型類の群と恒等自己同値の2自己同型群はそれぞれ
\begin{equation}\label{eq:categorical-weak-inertia}
 \Aut_{\mathrm{seq}}(\mathcal C)
   =\Aut_{\mathbb Z}(A)/\rho(T_0),
 \qquad
 \Aut_2(\operatorname{Id}_{\mathcal C})=k^\times.
\end{equation}
一点の同定を保持すると後者は自明となり，
前者を自己同型群とする通常の群oidを得る．
\end{proposition}
\begin{proof}
各$i$で同定$\alpha_i:F(E_i)\simeq E'_i$を選べば
式\eqref{eq:categorical-induced-map}を得る．
同定を$\lambda_i\alpha_i$に変更すると像は$\rho_\lambda\phi$になる．
二つの1射に同定$\alpha_i,\beta_i$を選び，自然同型$\eta:F\Rightarrow G$
を与えると，$\End(E'_i)=k$より
$\beta_i\eta_{E_i}\alpha_i^{-1}=\lambda_i\operatorname{Id}_{E'_i}$と書ける．
自然性はちょうど\eqref{eq:categorical-gauge-two-arrow}である．
逆にこの式が成り立つなら，$G$の同定を$\lambda_i^{-1}\beta_i$へ変更し，
定理\ref{thm:categorical-marked-equivalence}を適用すると所望の自然同型が一意に得られる．
$E_0$での条件は$\lambda_0=1$と同値である．

$\rho_\lambda=1$なら，$A_{i,i+1}\ne0$により
$\lambda_i=\lambda_{i+1}$であり，すべてのスカラーが等しい．
逆は明らかである．任意の$\lambda$を$\lambda_0$で割れば
同じ$\rho_\lambda$を与える$T_0$の元となる．
また添字を保つ$k$-線型代数同型は$\rho_\lambda$と可換なので，
この像は$\Aut_{\mathbb Z}(A)$の中心部分群である．
以上から自己同値の2同型類と2自己同型群の記述が従う．

恒等関手の自然自己同型については，別の直接的な確認もできる．
各$E_i$上の作用はスカラーであり，隣接Hom空間との自然性により
すべて同じスカラーになる．豊富列からの全射により，
このスカラーは圏全体の自然変換を一意に決める．
$E_0$上で恒等であるという条件を加えると，そのスカラーも$1$となる．
\end{proof}

従って，全同定付きモジュライにおいて頂点スカラーを既に同一視したと考えてはならない．
例えば$(\operatorname{Id},(\lambda_i\operatorname{Id}))$は，
$\lambda_i$が一定でない限り，異なる全同定付き自己同型を与える．
同定を忘れた後に初めて，これらが自然同型で結ばれる．
対角$\mathbb G_m$は弱い圏の\emph{2自己同型}として残るので，
通常の代数stackの慣性群から単に一つの$\mathbb G_m$を除く操作と混同しない．

\subsection{相対線束とfpqc降下}

底をスキーム$S$とし，$A_{ij}$が有限階数局所自由である場合を考える．
相対圏との比較には，次を明示的な仮定とする．
圏の族は，基底変更とその合成の整合性を備えた$S$-線型圏のfpqc stackであり，
直線束とのテンソル積と相対Homを持つものとする．
1射はこの構造および基底変更と両立する線型圏同値，
2射は基底変更と両立する自然同型である．
選んだ対象列の相対Homは$A_{ij}$であり任意の底変更と両立すること，
およびsection関手による復元同値がすべての底変更後に自然な形で成立することを課す．
座標代数から出発するときも，\eqref{eq:categorical-saturation}の相対版を
すべての底変更後に要求する．
相対tailsを大きなAbel圏$\operatorname{QGr}(A)$で扱い，
幾何学的点上のnoetherian部分圏を$\qgr$とすることもできる．
有限生成部分圏を基底上で直接使う場合は，その部分圏がAbelであることも必要である．
本節ではこれらの条件を満たす相対模型とそのfpqc局所形式を比較対象とし，
一般のAbel圏の族に対する相対Artin--Zhang定理を主張しない．

相対版の「列を保つ」には二通りある．大域同型
$F(E_i)\simeq E'_i$を要求する条件と，そのような同型がfpqc局所で存在する条件は
同じではない．stackとしての降下には後者を使う．
そこで，直線束の族$L=(L_i)_{i\in\mathbb Z}$に対し
\begin{equation}\label{eq:categorical-line-twist}
 A^L_{ij}=L_i\otimes_{\OO_S}A_{ij}\otimes_{\OO_S}L_j^{-1}
\end{equation}
と置く．積は中間の$L_j^{-1}\otimes L_j$を縮約して定める．
これは$(E_i\otimes L_i)$のsection代数である．

\begin{proposition}[弱い射の線束による記述]
\label{prop:categorical-relative-descent}
上の相対復元とsaturationの仮定の下で，二つの大域的な代数模型$A,A'$の間の
弱い圏同値の群oidは，次の群oidと同値である：
\begin{enumerate}[label=\textup{(\roman*)}]
\item 対象は，直線束族$L_i$と添字を保つ同型
$\phi:A\xrightarrow{\sim}(A')^L$の組である．
\item $(L,\phi)$から$(M,\psi)$への射は直線束の同型族
$\theta_i:L_i\xrightarrow{\sim}M_i$であって
$\psi=\Theta_\theta\phi$を満たすものとする．
ここで$\Theta_\theta:(A')^L\to(A')^M$は
$\theta_i$と$\theta_j$の双対の逆写像から誘導される．
\end{enumerate}
この記述で対応する圏同値は
$F(E_i)\simeq E'_i\otimes L_i$を満たす．
一点の同定を保持する場合は，さらに$L_0\simeq\OO_S$という自明化を指定し，
射はその自明化を保つものとする．
\end{proposition}
\begin{proof}
まず線束族を与えたとき，右$A$加群$M$を
\[
 (M^L)_i=M_i\otimes L_i^{-1}
\]
とする操作は，$A$加群と$A^L$加群の圏の同値を与える．
乗法は\eqref{eq:categorical-line-twist}により定義され，逆は$L_i^{-1}$を用いる操作である．
この同値はtorsionを保つためtailsに降りる．
その逆同値は$\pi e_iA^L$を$\pi e_iA\otimes L_i$へ送る．
従って$\phi:A\to(A')^L$から，主張した対象列上の作用を持つ圏同値が得られる．

逆に，相対圏同値$F$がfpqc局所で列を保つとする．
局所的に$F(E_i)\simeq E'_i$なので，相対Hom
$L_i=\mathcal Hom(E'_i,F(E_i))$は直線束である．
評価写像$E'_i\otimes L_i\to F(E_i)$は局所同型であり，大域的にも同型となる．
この同型を用いてHomと合成を取ると
$A\xrightarrow{\sim}(A')^L$を得る．
二つの圏同値の間の自然同型は，各$E_i$上で$\theta_i$を与え，
自然性は$\psi=\Theta_\theta\phi$と同値である．
逆向きの自然同型の存在と一意性は，局所的に全同定付き比較へ帰着する．
一意性によりこれらは重なり上で一致し，降下する．

すべての$L_i$を同時に自明化するfpqc被覆を使ってよいことも確認する．
それぞれのframe束の$S$上の積を取れば，これはaffineな忠実平坦射である．
実際，有限部分積はaffine忠実平坦であり，その座標代数のfiltered colimitも
平坦で，各幾何学的ファイバーが非零なので忠実平坦である．
この射はaffine，従って準コンパクトである．
無限個の線束を同時に自明化するこの議論から，有限表示の被覆が得られるとは主張しない．

最後に$F(E_0)\simeq E'_0$という同定は，評価同型を通して
$L_0$の自明化と同じである．自然同型がこの同定を保つことは，
$\theta_0$が自明化と両立することと同値である．
\end{proof}

全頂点に共通な直線束$L_i=L$を取ると$A^L\simeq A$であるが，
対応する弱い圏同値は$M\mapsto M\otimes L$となる．
従って「線束による捩れは代数式で相殺する」ことから，
相対圏同値自体が大域的に恒等であるとはいえない．
体上で現れた対角スカラーの2自己同型は，この相対現象の局所的な形である．

\begin{definition}[比較する三種類のモジュライ]
\label{def:categorical-three-moduli}
相対代数の許容クラスを底変更とfpqc降下に関して固定する．
全同定付きの相対模型をその射とともにstack化したものを
$\mathfrak M^{\mathrm{mark}}$とする．
これは上の相対比較の仮定の下で，飽和した代数と厳密な添字保存同型のstackに同値である．

式\eqref{eq:categorical-gauge-two-arrow}の2射を加えてfpqc 2-stack化したものを
$\mathfrak M^{\mathrm{weak}}$とする．
また$E_0$の同定を保持する模型について，同じ構成の各Hom群oidの
同型類を取り，fpqc stack化したものを$\mathfrak M^{\mathrm{pt}}$とする．
命題\ref{prop:categorical-scalar-gauge}により後者には非自明な2慣性がなく，
通常の1-stackとして扱える．
\end{definition}

ここで$\mathbb T=(\mathbb G_m)^{\mathbb Z}$，
$\mathbb T_0=\{\lambda_0=1\}$はfpqc群層である．
$A_{i,i+1}$が正の一定階数を持つので，相対版でも
$\ker\rho=\mathbb G_m$，$\rho|_{\mathbb T_0}$は単射である．
大域模型$A/S$において
\begin{equation}\label{eq:categorical-pointed-aut-sheaf}
 \mathcal Aut_{\mathfrak M^{\mathrm{pt}}}(A)
   \simeq \mathcal Aut_{\mathbb Z}(A)/\rho(\mathbb T_0)
\end{equation}
は\emph{fpqc層としての商}を意味する．
右辺の$S$上の切断を，常に大域自己同型の集合を大域スカラーで割った集合と
置き換えてはならない．命題\ref{prop:categorical-relative-descent}の線束が，
局所代表を貼り合わせるために必要になる．
また弱い2-stackの対象には局所模型の捩れた形式も含まれ得る．
すべての対象が，大域的な代数と大域的に選ばれた全対象列を持つとは主張しない．
これらの定義には代数的atlasや有限型性の存在を必要としない．

\subsection{二次包絡はすべての比較と両立する}

\begin{proposition}[内在的復元のgauge同変性]
\label{prop:categorical-envelope-equivariance}
定理\ref{thm:intrinsic-relative-envelope}の仮定の下で，任意の直線束族$L$について
標準同型
\begin{equation}\label{eq:categorical-envelope-twist}
 Q(A^L)\simeq Q(A)^L
\end{equation}
が存在する．三次核も$I^L$へ対応する．
従って命題\ref{prop:intrinsic-pair-groupoid}の同値は，
全同定付き代数stack，弱い2-stack，一点の同定を保持するstackの
それぞれにおいて，三次商のモジュライと許容対$(B,I)$のモジュライの同値を与える．
各levelでambientを忘れる射が定まる．
\end{proposition}
\begin{proof}
$V_i=A_{i,i+1}$，$R_i\subset V_i\otimes V_{i+1}$とする．
捩れた長さ一成分は$V_i^L=L_i\otimes V_i\otimes L_{i+1}^{-1}$である．
隣接テンソル積では中間の直線束が相殺し，二次積の核は
$R_i^L=L_i\otimes R_i\otimes L_{i+2}^{-1}$となる．
長さ$m$のテンソル積でも同様に中間因子がすべて相殺するので，
二次関係で割ると\eqref{eq:categorical-envelope-twist}を得る．
商写像と可換な同型であるため，核の成分も
$I^L_{ij}=L_i\otimes I_{ij}\otimes L_j^{-1}$へ写る．
これらの同型は線束の同型，積，底変更と自然に両立する．
従って全同定付きでの商と二次包絡の互いに逆な構成は，
命題\ref{prop:categorical-relative-descent}の射と2射にも延長し，
stack化の後も互いに逆となる．
\end{proof}

特に，先の自己同型比較\eqref{eq:intrinsic-automorphisms}は
\emph{厳密な添字保存代数同型，または全同定付き圏同値の2同型類}
についての等式である．同定を弱めた後の正確な式は
\begin{equation}\label{eq:categorical-pointed-pair-aut}
 \frac{\mathcal Aut_{\mathbb Z}(A)}{\rho(\mathbb T_0)}
 \simeq
 \frac{\mathcal Aut_{\mathbb Z}(Q(A),I)}{\rho(\mathbb T_0)}
 \lhook\joinrel\longrightarrow
 \frac{\mathcal Aut_{\mathbb Z}(Q(A))}{\rho(\mathbb T_0)}.
\end{equation}
これはfpqc群層の同型と単射である．
頂点スカラーは各斉次成分に単元倍で作用するので，すべての斉次イデアルを保つ．
また$A$と$Q(A)$の長さ一成分は同じであり，両辺の$\mathbb T_0$は同じ部分群を表す．
このため右端への単射も保たれる．
弱い2-stackでは，同じ自己同値の2同型類に加えて，
両側に対角$\mathbb G_m$の2自己同型を保持する．

具体的なambientへの同型まで指定したファイバーと，ambientの同型類だけを
固定した商との違いは，この比較の後にも残る．
可換ambientでの検算は命題\ref{prop:moduli-commutative-fibre}を参照されたい．
周期構造を選ぶ場合はさらに別のデータを加える必要があり，
その忘却射とファイバーは第\ref{sec:period-forgetting}節で扱う．
頂点スカラーの同一視と周期構造の忘却は，異なる操作である．

\subsection{既知の復元定理との適用範囲の比較}

Artin--Zhang\cite{AZ}の偏極を持つ圏の枠組みと，
Polishchuk\cite[Theorem~2.4 and Remark~2 following it]{PolishchukProj}の
一般の豊富列の枠組みは，次数シフトを入力に持つかどうかが異なる．
本節は後者を用い，射の同定データと自然変換をさらに明示したものである．
section代数の復元そのものは既知の結果であり，
相対saturationを省略する根拠にはならない．

次数付き加群圏についても，列を保つこととシフト関手に可換であることは異なる．
Sierra\cite[Theorem~3.6 and the following Remark]{SierraGAlgebras}は，
付随$\mathbb Z$-代数の同型を，標準自由加群のシフトの各同型類を保つ圏同値で特徴付ける．
同論文のCorollary~4.4はこれをZhang twistと対応させる
\cite{SierraGAlgebras,Zhang}．
これらは次数付き加群圏についての結果であり，
任意の無標識な$\qgr$の同値を付随$\mathbb Z$-代数の同型に置き換える結果ではない．

一方，3次元二次AS正則代数については，
Abdelgadir--Okawa--Ueda\cite[Theorem~7.3, Appendix~A]{AOUNCPlanes}が
無標識なAbel圏からのより強い復元を扱う．
同付録Corollary~A.7では，圏の内部で定まる次元のfiltration，
商圏での長さ，Euler形式を用いて$\{\OO(n)\}$を認識し，
その順序と1周期性から$\mathbb Z$-代数の同型類を復元する．
Corollary~A.10はこの場合の$\qgr$同値とZhang twistの関係も与える．
この対象認識を標準5次元AS正則ambientや三次商へ移す証明は本稿にはない．
従って，ここで定義した豊富列付きモジュライから列をすべて忘れる射が
同値になることは，独立の未解決問題として残す．

\section{構造対象だけを標識したAbel圏のモジュライ}
\label{sec:abel-marked-moduli}

本節では豊富列そのものをモジュライのデータから外し，構造対象だけを残す．
二次包絡を含む三次性の定義，全ての圏同値を含む射の定義，
特定の部分族での復元定理を分ける．特に中心の階数条件は，
可換三次三様体を含む全体の定義には課さない．

\subsection{対象の条件と全ての構造対象付き同値}
\label{subsec:abel-cubic-definition}

\begin{definition}[構造対象付き三次Abel圏]
\label{def:abel-cubic-pointed}
体$k=\C$上で，$(\mathcal C,\OO)$が\emph{構造対象付き三次Abel圏}であるとは，
$\mathcal C$が$k$-線型noetherian Abel圏であり，$E_0\simeq\OO$を満たす
豊富列$(E_i)_{i\in\mathbb Z}$で，次の条件を満たすものが\emph{存在する}ことをいう．
そのsection $\mathbb Z$-代数は，第\ref{sec:intrinsic-cubic}節の
三次データの条件を満たし，かつ第\ref{sec:categorical-moduli-comparison}節の
saturationと復元の条件を満たす．すなわち，標準Hilbert関数を持つ
次数一生成の代数で，二次包絡は標準Koszul AS正則次元$5$，
その核は左右正則な正規三次族である．

1射は任意の$k$-線型完全同値$F:\mathcal C\to\mathcal C'$と同型
$\alpha:F(\OO)\xrightarrow{\sim}\OO'$の組であり，
2射は$\beta\eta_{\OO}=\alpha$を満たす自然同型$\eta:F\Rightarrow G$とする．
豊富列の選択や次数一の自己同値は対象にも射にも指定しない．
\end{definition}

例えば滑らかな三次三様体$X\subset\PP^4$は
$(\operatorname{coh}X,\OO_X)$としてこのクラスに入る．
豊富列は$E_i=\OO_X(-i)$，二次包絡は五変数多項式環の
$\mathbb Z$-代数である．一方，定義\ref{def:categorical-three-moduli}の
$\mathfrak M^{\mathrm{pt}}$は対象列自体をまだ記憶している．
この違いのため，列を忘れる射の像に，列保存の射だけを入れてはならない．
定義\ref{def:abel-cubic-pointed}は対象の条件を満たす圏の間の
\emph{全て}の構造対象付き完全同値を採用する．

相対版では第\ref{sec:categorical-moduli-comparison}節と同じく，
相対Hom，底変更，線束とのテンソル積，およびその整合性を備える圏の族を使う．
アフィン上の標準Grothendieck模型の底変更には，係数環の作用を備えた対象の圏を用いる．
非noetherianな係数環には大きな$\operatorname{QGr}$を使い，
小さい$qgr$への制限はそのAbel性が保証される場合に行う．
fpqc局所で上記の相対$qgr$/$QGr$模型を持ち，sectionのsaturationと復元が
指定した全ての底変更後に成立することを要求する．このような模型の間の
全ての構造対象付き相対完全同値を射として，Hom群oidの同型類を取り，
fpqc stack化したものを$\mathfrak M^{\OO}_{\mathrm{cub}}$と書く．
この構成は相対模型の条件と底変更構造を明示したモジュライ問題であり，
任意の抽象Abel圏に底変更を無条件に定義する定理ではない．
有限型atlas，対角の代数性，smooth proper locusの開性は，定義により従うものではない．
特に非正則Artin基底上のAS正則性には，
定理\ref{thm:abel-relative-as-lift}で明示する相対解消と両側AS双対を用い，
基底を含めた絶対的大域次元の有限性を要求しているのではない．

実際，任意の許容な豊富列について$\End(E_i)=\OO_T$であり，
隣接Homは正の階数の局所自由層である．自然自己変換の$E_i$上のスカラーは
隣接Homの忠実性により全て等しく，豊富性から圏全体の変換を決定する．
従って$\OO=E_0$上で恒等である2自己同型は自明である．
このことが，上記の構造対象付き問題を通常の1-stackとして扱う理由である．
体上の同じ議論は補題\ref{lem:abel-pointed-scalar-center}に記す．

\begin{remark}[可換三次点での普遍的な相対模型]
\label{rem:abel-cubic-essential-image}
滑らかな可換三次三様体を特殊ファイバーとして固定すると，
全ての平坦Abel Artin変形は，任意の係数変更後にも上の三次条件を満たす
標準Grothendieck模型で表示できる．
第\ref{sec:abel-formal-cubic}節の定理\ref{thm:abel-cubic-flat-envelope}が
二次包絡の全次数平坦性，正規正則三次核，相対Koszul--AS性を証明し，
系\ref{cor:abel-universal-artin-model}と\ref{cor:abel-universal-relative-cubic}が
任意の係数変更後のsectionの飽和と復元，相対解消の両立を証明する．
固定した模型の対象と射のfpqc降下も同節で示す．
標準係数族からなる相対変形群oidは全Abel変形群oidに同値であり
（定理\ref{thm:abel-canonical-relative-formal}），
その古典的同型類関手は系\ref{cor:abel-cubic-formal-ring}の
17次元・接空間20次元の完備局所環で表される．

stack化後の全形式完備化との同定について残るのは，
圏自体の降下が有効で，Artin基底上に大域的な標準模型を持つかという問題である．
これは固定した圏の対象の降下とは別である．
有限型atlas，任意の非可換点での局所完全性，または
特殊ファイバーの自己同型を含めたstack全体の表示も，この結論には含めない．
\end{remark}

\subsection{中心二変数を持つ族の偶数Veronese復元}
\label{subsec:abel-relative-veronese}

本節は中心Jacobian型のうち低次中心を制御できる部分に関する相対定理である．
三次商の可換点では次数二中心の次元は$15$であり，以下の階数$3$の条件を満たさない．
従って，本節の条件が可換点の近傍を含むことは主張しない．
中心条件については下の有限rank命題が開条件を与えるが，
他の許容条件の開性まで同時に従うものではない．Abel圏との比較に必要な
相対section algebraのsaturationやSerre関手の底変更は，本節の代数的主張とは別の入力である．

\begin{definition}[普遍的な中心二変数条件]
\label{def:abel-two-center-family}
$T$を$\C$上のスキームとする（特に$2$は可逆である）．非負次数付き
$\mathcal O_T$-代数$A=\bigoplus_{n\geq0}A_n$が本節の許容条件を満たすとは，
次を満たすことをいう．
\begin{enumerate}
\item $A_0=\mathcal O_T$であり，各$A_n$は有限階数局所自由である．
\item 階数$2$のベクトル束$E\subset A_1$が中心に含まれ，積は単射
  $P:=\operatorname{Sym}_{\mathcal O_T}E\hookrightarrow A$を与える．
  $A$は$P$加群として平坦である．
\item 任意の$T'\to T$について
  \begin{equation}\label{eq:abel-universal-center-two}
    Z(A_{T'})_2=\operatorname{Sym}^2 E_{T'}
       \ \subset (A_{T'})_2
  \end{equation}
  が成り立つ．ここで左辺は底変更後の代数で中心を取り直したものとする．
\end{enumerate}
中心の形成が底変更と自動的に可換であると仮定したのではなく，
式\eqref{eq:abel-universal-center-two}を非平坦な底変更も含めて要求している．
$P$上の自由性は十分条件だが，以下の証明では平坦性だけを用いる．
\end{definition}

\begin{lemma}[中心多項式環の埋込みは普遍的である]
\label{lem:abel-universal-polynomial-embedding}
$K$を任意の可換環，$P=K[u,v]$を標準次数付き多項式環とする．
$A$は非負次数付きの単位的$P$-代数で，$A_0=K$であり，
$A$が$P$加群として平坦であるとする．このとき$P\to A$は単射であり，
任意の$K\to K'$による底変更後にも単射である．
\end{lemma}
\begin{proof}
核は斉次イデアルなので，像が零となる斉次元$p\in P$を取る．
$P$平坦性を列
\[
 0\longrightarrow\operatorname{Ann}_P(p)\longrightarrow P
          \xrightarrow{p}P
\]
に適用すると，$p\cdot1_A=0$から
$1_A\in\operatorname{Ann}_P(p)A$を得る．
$p$は斉次なので$\operatorname{Ann}_P(p)$も斉次である．
非負次数と$A_0=K$を用いてこの等式の次数零成分を取ると，
$1\in\operatorname{Ann}_P(p)_0$となる．従って$p=0$である．
底変更後も$A\otimes_KK'$は$P\otimes_KK'$上平坦であり，
非負次数と次数零の条件を保つので，同じ議論を適用できる．
階数二の束の対称代数についても，局所基底を取れば同じ結論を得る．
\end{proof}

\begin{proposition}[二次中心の普遍性を保証する有限rank条件]
\label{prop:abel-center-rank-open}
$A$は次数一生成で，$A_1,A_2,A_3$が有限階数局所自由とする．
中心に含まれる階数$2$の束$E\subset A_1$を取り，
$\operatorname{Sym}^2E\to A_2$は階数$3$の部分束と仮定する．
交換子写像は
\begin{equation}\label{eq:abel-center-commutator-map}
 \overline c:A_2/\operatorname{Sym}^2E
     \longrightarrow\mathcal Hom_{\mathcal O_T}(A_1,A_3),
 \qquad \overline z\longmapsto(a\longmapsto za-az)
\end{equation}
を定める．$\overline c$が幾何学的ファイバー上で単射となる点の集合$T^\circ$は開である．
この開集合上では\eqref{eq:abel-universal-center-two}が任意の底変更後に成り立つ．
\end{proposition}
\begin{proof}
始項と終項は有限階数局所自由なので，局所行列の最大小行列式の非消滅が
$T^\circ$を定める．その上で$\overline c$は局所分裂した単射となり，
余商も局所自由である．従って任意の底変更後も単射である．
また$\operatorname{Sym}^2E\subset A_2$は部分束なので，その核の形成も底変更と両立する．
次数一生成により，二次元が代数全体と可換であることは$A_1$と可換であることと同値である．
従って底変更後も交換子写像の核はちょうど$\operatorname{Sym}^2E$となる．
\end{proof}

この命題は，中心条件をある具体的な有限係数族のrank-minor開部分で確認する方法を与える．
その開集合の非空性，$\operatorname{Sym}E$上の平坦性，tailsの正則性，
dg smooth proper性は別の確認を要する．可換三次点では二次中心は階数$15$なので，
この最大rank開部分には含まれない．

$R=A^{(2)}$，$D=P^{(2)}$と置き，$R_n=A_{2n}$，$D_n=P_{2n}$と次数を付ける．
以下の$\underline{\operatorname{Isom}}_{\mathrm{gr},T}$は，
$T$上のfpqc層としての次数付き代数同型層を意味する．

\begin{lemma}[低次中心の内在性と普遍的な交差表示]
\label{lem:abel-universal-center-and-cech}
定義\ref{def:abel-two-center-family}の条件の下で，任意の底変更後に
\begin{equation}\label{eq:abel-intrinsic-center}
 Z(A)_1=E,\qquad Z(R)_1=P_2,\qquad
 D=\langle Z(R)_1\rangle_{\mathrm{alg}}
\end{equation}
となる．さらに$E$を局所的に$u,v$で自明化すると，
\begin{equation}\label{eq:abel-universal-intersection}
 A=A[u^{-1}]\cap A[v^{-1}]
       \quad\text{in }A[(uv)^{-1}]
\end{equation}
が成り立ち，この表示は任意の底変更後にも成立する．
\end{lemma}
\begin{proof}
$P=K[u,v]$となる基底のaffine開集合上で証明する．
$A$の$P$平坦性から，$u,v,uv$による乗法は単射である．
$z\in Z(A)_1$なら，$uz,vz\in P_2$であり，$u(vz)=v(uz)$を満たす．
二次式の係数を比較すると，ある$\ell\in P_1$について$uz=u\ell$となる．
$u$を消去すれば$z=\ell$であり，$Z(A)_1=E$を得る．

$z\in Z(R)_1$とし，$a\in A$を斉次元とする．
$a$が偶数次数なら$[z,a]=0$である．奇数次数なら$ua\in R$なので，
$u[z,a]=[z,ua]=0$となり，再び$u$を消去できる．
従って$Z(R)_1=Z(A)_2=P_2$である．$P^{(2)}$は$P_2$で生成されるので
最後の等式も従う．

Laurent単項式の係数比較により
\[
 0\longrightarrow P\longrightarrow P[u^{-1}]\oplus P[v^{-1}]
       \longrightarrow P[(uv)^{-1}]
\]
は最初の二つの非零項で完全である．最後の射の全射性は必要ない．
これを平坦な$P$加群$A$とテンソルすれば
\eqref{eq:abel-universal-intersection}を得る．
$P$平坦性は基底変更で保たれ，中心の条件も普遍的に課しているため，
同じ議論は任意の$T'\to T$上で成立する．
\end{proof}

\begin{lemma}[二次Veronese錐の同型の持上げ]
\label{lem:abel-relative-conic-lift}
階数$2$のベクトル束$E,E'$に対し，
$P=\operatorname{Sym}E$，$P'=\operatorname{Sym}E'$，
$D=P^{(2)}$，$D'=(P')^{(2)}$と置く．制限射
\begin{equation}\label{eq:abel-conic-isom-torsor}
 \underline{\operatorname{Isom}}_{\mathrm{gr},T}(P,P')
  \longrightarrow
 \underline{\operatorname{Isom}}_{\mathrm{gr},T}(D,D')
\end{equation}
は$\mu_{2,T}$-torsorである．すなわち，右辺の任意の$T'$上の同型の
持上げ層は，$T'$上の有限\'etaleな$\mu_2$-torsorで表される．
\end{lemma}
\begin{proof}
ベクトル束を局所的に自明化する．
$\operatorname{Proj}D$は$\mathbb P^1$であり，その$\mathcal O(1)$は
$\mathbb P^1$上の$\mathcal O(2)$である．$D\simeq D'$はこの偏極付き射影直線の同型を誘導する．
射影直線の自己同型層が$\operatorname{PGL}_2$であることにより，
fpqc局所でこの射影同型をある$M:E\simeq E'$へ持ち上げられる．
この事実は非被約な試験基底も含む関手の等式である．
基礎体上の任意の代数についての直接証明は
\href{https://www.jmilne.org/math/Books/iAG2022.pdf}
{Milne, \emph{Algebraic Groups}, Proposition~20.7}を参照されたい．

$\operatorname{Sym}^{\mathrm{even}}M$を取り除いた差は射影直線上で恒等である．
従って偏極$\mathcal O(2)$上の自己同型は基底のある単元$c$による乗法であり，
対応する$D_n$上の作用は$c^n$である．ここでは
$D_n\simeq p_*\OO_{\mathbb P(E)}(2n)$が任意の底変更と両立することを使う．
従って非被約基底でも，射影化で見えない低次数の作用が残ることはない．
有限\'etale被覆$z^2=c$を取って$M$を$zM$へ置き換えれば，
与えられた$D$の同型そのものを持ち上げる．平方根被覆の忠実平坦性と\'etale性は
\href{https://stacks.math.columbia.edu/tag/03PK}
{Stacks Project, Section~59.28, Lemma~59.28.1}のKummer計算である．

二つの持上げの差は射影化で恒等なのでスカラー$\lambda$であり，
偶数次数上で恒等であることは$\lambda^2=1$と同値である．
従って持上げ層は$\mu_2$の単純推移的作用を持ち，局所的に非空である．
$\mu_2$は有限\'etaleなので，降下によりこのtorsorも有限\'etaleで表される．
\end{proof}

\begin{theorem}[相対偶数Veronese復元と同型層]
\label{thm:abel-relative-veronese-isom}
$A,A'$は定義\ref{def:abel-two-center-family}の許容族とし，
対応する環を$P,P',R,R',D,D'$と書く．制限は標準同型
\begin{equation}\label{eq:abel-veronese-cartesian}
 \underline{\operatorname{Isom}}_{\mathrm{gr},T}(A,A')
 \simeq
 \underline{\operatorname{Isom}}_{\mathrm{gr},T}(R,R')
 \mathop{\times}_{\underline{\operatorname{Isom}}_{\mathrm{gr},T}(D,D')}
 \underline{\operatorname{Isom}}_{\mathrm{gr},T}(P,P')
\end{equation}
を与える．特に
\begin{equation}\label{eq:abel-veronese-mu-two}
 \underline{\operatorname{Isom}}_{\mathrm{gr},T}(A,A')
   \longrightarrow
 \underline{\operatorname{Isom}}_{\mathrm{gr},T}(R,R')
\end{equation}
は$\mu_2$-torsorである．その作用は，偶数次数で恒等，奇数次数で
$\epsilon\in\mu_2$を掛ける自己同型との合成である．
全構成は任意の底変更と両立する．
\end{theorem}
\begin{proof}
補題\ref{lem:abel-universal-center-and-cech}から，$A$の同型は$P$を，
$R$の同型は$D$を必ず保存する．従って制限射は内在的に定まる．

$\phi:R\simeq R'$と，その$D$上の制限を持ち上げる
$\beta:P\simeq P'$を与える．$E$の局所基底を$u,v$とする．
$R,D$の次数を二倍して元の次数に戻すと，積は標準同型
\begin{equation}\label{eq:abel-localized-veronese}
 R[(u^2)^{-1}]
    \otimes_{D[(u^2)^{-1}]}P[u^{-1}]
       \xrightarrow{\sim} A[u^{-1}]
\end{equation}
を与える．実際，$P[u^{-1}]$は$D[(u^2)^{-1}]$上で基底$1,u$を持ち，
$A[u^{-1}]$の偶数部分は$R[(u^2)^{-1}]$，奇数部分はその$u$倍である．
後者は奇数斉次元$a$について$a/u=(ua)/u^2$と書けば分かる．
中心性によりこの加群分解は表示した代数同型を与える．$v$についても同様である．

式\eqref{eq:abel-localized-veronese}と$\phi,\beta$から，
$A[u^{-1}]\simeq A'[\beta(u)^{-1}]$および
$A[v^{-1}]\simeq A'[\beta(v)^{-1}]$を得る．
重なり上では同じ$\phi,\beta$から作った写像なので一致する．
$\beta(u),\beta(v)$は$E'$の基底であり，
式\eqref{eq:abel-universal-intersection}を両辺に適用すれば
一意な$A\simeq A'$へ制限される．逆向きも同様である．
これで\eqref{eq:abel-veronese-cartesian}を得た．

制限射\eqref{eq:abel-veronese-mu-two}は
\eqref{eq:abel-conic-isom-torsor}の底変更なので$\mu_2$-torsorである．
中心一次元を$\epsilon$倍すると，偶数部分を固定した持上げは奇数部分を
$\epsilon$倍する．各局所化・テンソル積は底変更に可換であり，
交差の表示は普遍的に完全な列から得られたので，底変更との両立も従う．
\end{proof}

\begin{corollary}[次数スカラーを除いた同型層]
\label{cor:abel-projective-veronese-isom}
定理\ref{thm:abel-relative-veronese-isom}の条件の下で，
次数スカラーを除いたfpqc商層に標準同型
\begin{equation}\label{eq:abel-projective-isom-sheaf}
 \underline{\operatorname{Isom}}_{\mathrm{gr},T}(A,A')/\mathbb G_m
   \xrightarrow{\sim}
 \underline{\operatorname{Isom}}_{\mathrm{gr},T}(R,R')/\mathbb G_m
\end{equation}
がある．左の$\lambda$は$A_n$に$\lambda^n$を掛け，
右の$c$は$R_n$に$c^n$を掛ける．従って制限上のスカラー写像は
$\lambda\mapsto\lambda^2$である．
\end{corollary}
\begin{proof}
前定理により局所全射である．二つの$A$の同型の偶数制限が
右辺の次数スカラー$c$だけ異なるとき，fpqc局所で$\lambda^2=c$を選ぶ．
左辺で$\lambda$の次数スカラーを合成すれば偶数制限は一致し，
残る差は前定理の$\mu_2$である．これも左辺の次数スカラーに含まれる．
従ってfpqc商層では単射でもある．
\end{proof}

\begin{remark}[同型と対象の降下障害を区別する]
\label{rem:abel-veronese-descent-obstruction}
固定した$R\simeq R'$の持上げ層は$\mu_2$-torsorなので，
その大域持上げの障害は$H^1_{\mathrm{fppf}}(T,\mu_2)$の類である．
一つの大域持上げがあれば全持上げは$\mu_2(T)$のtorsorになる．
従って「持上げが正確に二つ」という表現は代数閉体上の一点では正しいが，
一般の基底では大域切断が存在しないことも，非連結性により二つより多いこともある．

一方，局所的に許容族の偶数部分として表示される$R$について，
$A$そのものを大域的に貼り合わせる問題は別である．
局所持上げを対象，$R$上で恒等となる同型を射とするfpqc stackは
$\mu_2$でbandされたgerbeであり，対象の降下障害は
$H^2_{\mathrm{fppf}}(T,\mu_2)$に置かれる．
これは三重の重なり上で持上げ同型のcocycle条件のずれが$\mu_2$に入るという，
gerbeの標準的な降下障害である
（\href{https://stacks.math.columbia.edu/tag/0CJZ}{Stacks Project, Section~21.11}参照）．
局所的に二次Veronese錐である中心が，大域的にも
ある階数$2$の束$E$の$\operatorname{Sym}^{\mathrm{even}}E$であるとは限らない．
射影二次曲線束のBrauer障害や，偏極に含まれる基底の線束の平方根の障害を
省略してはならない．
本系から得られるのは次数スカラー商の\emph{層}の同型であり，
大域代数模型の存在やAbel圏の自然同型による商との一致は別に確認する．
\end{remark}

\begin{example}[複素基底でも大域持上げがない同型]
\label{ex:abel-veronese-nonsquare}
$K=\mathbb C[t,t^{-1}]$，$A=K[u,v]$，$R=A^{(2)}$とする．
$R_n$上で$t^n$を掛ける次数付き自己同型を考える．
これを$A$へ持ち上げる行列$M\in\operatorname{GL}_2(K)$があれば，
その射影化は恒等なので$M=\lambda I$であり，$\lambda^2=t$を要する．
しかし$K^\times=\mathbb C^\times t^{\mathbb Z}$にそのような$\lambda$はない．
持上げは有限\'etale被覆$z^2=t$の後に存在する．
これは三次三様体の例ではなく，中心二変数の持上げに現れる降下障害の最小例である．
同じ問題は任意の許容な定数族の偶数部分に次数スカラー$t$を作用させても現れる．
\end{example}

\begin{remark}[相対Serre軌道環への適用範囲]
\label{rem:abel-relative-serre-line-twist}
相対Serre関手が$\mathsf S\simeq(-2)\otimes L[3]$と，
基底の直線束$L$を伴う場合には，$\mathsf H=\mathsf S^{-1}[3]$による軌道環は
$\bigoplus_{n\geq0}A_{2n}\otimes L^{-n}$となり得る．
この場合は$L$を局所自明化して上の定理を適用し，
直線束の捩れと次数スカラーの降下を別に扱う必要がある．
この式には，相対Serre双対・section algebraのsaturation・Homの底変更という
追加の仮定が必要であり，本節の代数的復元定理から自動的には従わない．
\end{remark}

\subsection{構造対象を保つAbel圏同値の全射と自然同型}
\label{subsec:abel-pointed-field-hom}

本小節では基礎体を$k=\C$とする．
対象$(\mathcal C,\OO)$の間の\emph{構造対象付きAbel圏同値}とは，
$k$-線型完全同値$F:\mathcal C\to\mathcal C'$と同型
$\alpha:F(\OO)\xrightarrow{\sim}\OO'$の組をいう．
$(F,\alpha)$から$(G,\beta)$への射は自然同型$\eta:F\Rightarrow G$で
$\beta\eta_{\OO}=\alpha$を満たすものとする．
この群oidを$\operatorname{Eq}_{\OO}(\mathcal C,\mathcal C')$と書く．
以下の仮定は比較定理の適用範囲を指定するものであり，
非可換三次三様体全体の定義に中心の次元条件を課すものではない．

\begin{lemma}[豊富な偶数列によるスカラー中心の認識]
\label{lem:abel-pointed-scalar-center}
$\mathcal C$内の豊富列$(E_i)_{i\in\mathbb Z}$が
$\End(E_i)=k$，$\Hom(E_{i+1},E_i)\ne0$を満たすとする．
このとき$\operatorname{Nat}(\operatorname{Id}_{\mathcal C},
\operatorname{Id}_{\mathcal C})=k$であり，
$\OO=E_0$とした構造対象付き圏同値の2自己同型は自明である．
\end{lemma}
\begin{proof}
自然変換$\eta$は$E_i$上でスカラー$\lambda_i$として作用する．
非零の隣接Homに自然性を適用すると$\lambda_i=\lambda_{i+1}$となる．
任意の対象への豊富列の有限直和からの全射を用いると，
$\eta$はこの共通スカラーの恒等変換に等しい．
圏同値$F$の自然自己変換にも準逆を合成して同じ議論を適用できる．
$\OO$の指定された同定を保つという条件はスカラーを$1$に固定する．
\end{proof}

正の次数付き代数$A$に対し，$\varepsilon_\lambda$で
次数$m$上に$\lambda^m$を掛ける\emph{次数スカラー}を表す．
これは構造対象を保つ恒等tails関手と自然同型な関手を誘導する．
従って次数スカラーは，厳密な代数自己同型としては非自明でも，
構造対象付きAbel圏同値の2同型類では自明となる．

\begin{theorem}[構造対象付きAbel圏のHom群oid]
\label{thm:abel-pointed-field-hom}
$A,A'$をconnected，局所有限，次数一生成の右noetherian次数付き
$k$-代数とし，$\mathcal C_A=\qgr A$，$\OO_A=\pi A$と置く．
両方について次を仮定する．
\begin{enumerate}[label=\textup{(\roman*)}]
\item 標準section写像
\[
 A_m\longrightarrow\Hom_{\mathcal C_A}(\OO_A,\OO_A(m))
\]
はすべての整数$m$で同型である．ただし$m<0$では$A_m=0$とする．
偶数列$(\OO_A(-2i))_{i\in\mathbb Z}$は豊富で，$A_2\ne0$とする．
\item $D^b(\mathcal C_A)$はHom有限でSerre双対を持ち，そのSerre関手に
三角関手の同型$\mathsf S_A\simeq(-2)[3]$が存在する．
この同型の選択をモジュライの標識には含めない．
\item $R=A^{(2)}$，$R'=(A')^{(2)}$について，制限写像
\[
 \operatorname{Isom}_{\rm gr}(A,A')
 \longrightarrow\operatorname{Isom}_{\rm gr}(R,R')
\]
は全射で，各fiberは$\varepsilon_{-1}$の作用に関する
$\mu_2(k)$-torsorである．
\end{enumerate}
すると，構造移送は標準的な全単射
\begin{equation}\label{eq:abel-pointed-hom-quotient}
 \operatorname{Isom}_{\rm gr}(A,A')/k^\times
 \xrightarrow{\sim}
 \pi_0\operatorname{Eq}_{\OO}(\mathcal C_A,\mathcal C_{A'})
\end{equation}
を与える．左辺の$k^\times$は次数スカラーの合成によって作用する．
右辺の各Hom群oidは離散群oidに同値であり，
同じ左辺の類を与える二つの構造移送の間には，
構造対象の同定を保つ自然同型が一意に存在する．
これらの全単射は同値の合成と両立する．
特に
\begin{equation}\label{eq:abel-pointed-aut-quotient}
 \operatorname{Aut}_{\OO}(\mathcal C_A)
 \simeq \operatorname{Aut}_{\rm gr}(A)/k^\times.
\end{equation}
\end{theorem}
\begin{proof}
まず$k$-線型Abel圏同値は完全なので，項ごとに適用することにより
有界導来圏の同値を誘導する．Serre双対の表現可能性とYonedaの一意性により，
この導来同値はSerre関手と整合的に可換する
（\cite[Lemma~57.3.5, Tag~0FY3]{StacksAbel}）．
ここで用いる一意性は双対同型を備えたSerre関手の一意性であり，
自己同値としての同型の任意の選択を一意と主張するものではない．

$H_A=\mathsf S_A^{-1}[3]$と置く．仮定\textup{(ii)}により
$H_A\simeq(2)$であり，$H_A$は標準heartを保つ．
同型$H_A\simeq(2)$と$H_{A'}\simeq(2)$を一時的に選ぶ．
構造対象付き同値$(F,\alpha)$は，Serre関手との可換同型とその反復から
軌道環の次数付き同型
\begin{align*}
 \phi_R:R
 &\simeq\bigoplus_{n\ge0}
       \Hom(\OO_A,H_A^n\OO_A)\\
 &\xrightarrow{\sim}\bigoplus_{n\ge0}
       \Hom(\OO_{A'},H_{A'}^n\OO_{A'})
 \simeq R'
\end{align*}
を誘導する．同じ可換同型を反復して用いるため，この写像は積を保つ．
負の冪も用いれば，偶数豊富列のすべての相対Homにおける写像が定まる．
仮定\textup{(i)}と豊富列による復元により，これらのHomと合成上の作用は
$F$と構造対象上の同定$\alpha$を自然同型まで決定する．

仮定\textup{(iii)}で$\phi_R$を$\phi_A:A\xrightarrow{\sim}A'$へ持ち上げる．
$\phi_A$の構造移送は偶数豊富列のsection代数に$\phi_R$として作用する．
従って豊富列の復元定理により，その構造移送は$(F,\alpha)$へ自然に同型となる．
これで\eqref{eq:abel-pointed-hom-quotient}の全射性を得る．

次に核を計算する．二つの代数同型を比較するには，一方の逆と合成して，
$\phi\in\operatorname{Aut}_{\rm gr}(A)$の構造移送が
構造対象付き恒等関手へ自然同型である場合を考えればよい．
構造移送による$\OO_A(-2i)$の標準同定の下で，その自然同型は
$\OO_A(-2i)$上で$\lambda_i\in k^\times$として作用し，$\lambda_0=1$である．
$\Hom(\OO_A(-2i-2),\OO_A(-2i))=A_2$上の自然性は
\[
 \lambda_i\phi(a)=\lambda_{i+1}a\qquad(a\in A_2)
\]
を与える．従って$\phi|_{A_2}=c\operatorname{Id}$となる
$c\in k^\times$が存在する．$A$は次数一生成なので$R$は次数一生成であり，
$\phi|_{A_{2n}}=c^n\operatorname{Id}$である．
$b^2=c$を満たす$b\in k^\times$を選ぶと，
$\varepsilon_b^{-1}\phi$の$R$への制限は恒等である．
仮定\textup{(iii)}により$\varepsilon_b^{-1}\phi$は
$1$または$\varepsilon_{-1}$である．従って$\phi$は次数スカラーである．
逆に次数スカラーの構造移送は，次数ごとのスカラーによる自然同型により
構造対象付き恒等関手となる．これで単射性を得る．

最後に，仮定\textup{(i)}の偶数列は
補題\ref{lem:abel-pointed-scalar-center}の仮定を満たす．
従って構造対象付き1射の2自己同型は自明であり，
存在する自然同型は一意である．合成との両立は構造移送の合成から従う．
\end{proof}

\begin{remark}[Serre同定の選択と次数スカラー]
\label{rem:abel-pointed-serre-normalization}
上の証明で$H_A\simeq(2)$の選択を変更する場合，その差は
$(2)$の自然自己同型である．補題\ref{lem:abel-pointed-scalar-center}により
それは一つのスカラーであり，$n$回の反復においてその$n$乗として作用する．
従って誘導される偶数軌道環の同型の変化は次数スカラーに限られる．
二つの平方根による持上げの差$\varepsilon_{-1}$も，
全次数スカラー$k^\times$の部分群である．
したがって\eqref{eq:abel-pointed-hom-quotient}はこれらの選択に依存しない．
これは，厳密なSerre双対と両立する軌道環同型の正規化を，
任意の次数付き同型と無条件に一致させるという主張ではない．
また，固定した偶数環上の$\mu_2$持上げtorsorを，
構造対象付きAbel圏の追加の2慣性と解釈してはならない．
\end{remark}

\begin{remark}[中心二変数を持つ三次族への適用範囲]
\label{rem:abel-pointed-central-chart}
仮定\textup{(iii)}は，中心一次元$u,v$による$k[u,v]$上の平坦性と
$Z(A)_2=\operatorname{Sym}^2(ku\oplus kv)$などを仮定した
定理\ref{thm:abel-relative-veronese-isom}の体上の特殊化から得られる．
従って本定理は，これらの仮定，saturation，Serre双対を満たす
中心Jacobian型三次商について，既存の同型類の復元を射と自然同型まで強める．
一方，滑らかな可換三次三様体の座標環では
$Z(A)_2=A_2$は15次元であり，この3次元中心の条件は成立しない．
従ってこの条件を全三次三様体の内在的定義に採用したり，
可換点を含む開近傍で自動的に成立するとしたりしない．
一般の正規三次商，標識を忘れたAbel圏，任意の基底上での相対的な同値，
代数的atlasの存在は，本定理の結論に含まれない．
\end{remark}

\subsection{構造対象付きAbel圏との相対比較：仮定を明示した形}
\label{subsec:abel-relative-categorical-hom}

相対偶数Veronese定理は代数の同型層についての結果である．
これをAbel圏の族の同値へ移すため，ここでは必要な相対圏論の入力を明記する．
この小節では有限型$\C$-スキームとfppf位相を用いる．
これは局所Artin $\C$-代数のスペクトルを含み，ファイバー積で閉じている．
第\ref{subsec:abel-relative-veronese}小節の代数的復元定理の任意の基底に対するfpqc主張を弱めるものではない．

$T$を有限型$\C$-スキーム，$A,A'$を
定義\ref{def:abel-two-center-family}の条件を満たす族とする．
以下の条件は両方の族について，任意の有限型$\C$-スキーム$U\to T$の後に課す．
局所的には$U=\operatorname{Spec}R$として記述する．
\begin{enumerate}[label=\textup{(C\arabic*)}]
\item $\mathcal C_U=\operatorname{QGr}(A_U)$は$R$-線型locally noetherian
Grothendieck圏で，構造対象$\OO_U$と次数捩れを持つ．
これらの相対模型は，基底変更の合成に関する整合性とfppf降下を備える．
完全同値・自然同型についてもこの相対構造を使う．
従って，ここでいう族の同値には，さらなる基底変更との両立を含める．
\item 任意の整数$m$について，標準写像
\[
 (A_U)_m\xrightarrow{\sim}
 \mathcal Hom_{\mathcal C_U}(\OO_U,\OO_U(m))
\]
は同型であり，任意の有限型基底変更と両立する．
偶数列$E_i=\OO_U(-2i)$は豊富で，そのsection関手は
$\mathcal C_U$と$\operatorname{QGr}(A_U^{(2)})$との復元同値を与える．
この同値は基底変更に自然とし，偶数列上のHomと合成への作用は
完全同値を自然同型まで決定すると仮定する．
最後の条件は，ファイバーごとの豊富性だけから結論せず，
第\ref{sec:categorical-moduli-comparison}節の相対復元の入力として課す．
\item 導来圏には$\mathcal C_U$の複体から得られる標準dg enhancementを用いる．
$\mathcal P_U$を，コホモロジーが有界かつnoetherianで，
$R$上のTor振幅が有限な複体からなるfull dg部分圏とする．
ここでTor振幅は，$\mathcal C_U$の標準的な$R$-加群作用を用いた
$M\otimes_R^{\mathbf L}K$のコホモロジーの消滅で定義する．
すべての$E_i$は$\mathcal P_U$に含まれ，この部分圏と相対RHomは
指定した基底変更と両立すると仮定する．
\item $\mathcal P_U$の二対象間の相対RHomは$R$-perfectであり，
整合的な相対Serre双対
\[
 \operatorname{RHom}_R
 \bigl(\operatorname{RHom}_{\mathcal C_U}(K,L),R\bigr)
 \simeq
 \operatorname{RHom}_{\mathcal C_U}(L,\mathsf S_U K)
\]
を持つ．この双対と$\mathsf S_U$は基底変更と両立する．
fppf局所で，$\mathcal P_U$上の関手として
$\mathsf S_U\simeq(-2)[3]$となる．この表示の選択を標識には含めない．
\end{enumerate}

条件\textup{(C3)}で相対perfect部分圏を内在的に指定した理由は，
すべてのAbel圏同値がこの部分圏を保存することを保証するためである．
単に何らかのenhancementが存在してその上にSerre関手がある，という条件だけでは，
任意のAbel圏同値をそのenhancementへ持ち上げる根拠にはならない．

\begin{lemma}[完全線型同値と相対perfect部分圏]
\label{lem:abel-equivalence-preserves-relative-perfect}
上の標準dg enhancementと相対perfect部分圏を用いると，
完全$R$-線型Abel圏同値$F:\mathcal C_U\to\mathcal C'_U$は
$\mathcal P_U\simeq\mathcal P'_U$を誘導し，相対RHomを保つ．
従って相対Serre双対と両立する同型
$F\mathsf S_U\simeq\mathsf S'_U F$を持つ．
\end{lemma}
\begin{proof}
圏同値は任意の直和と余核を保ち，$R$-線型性により
$F(M\otimes_R N)\simeq M\otimes_R F(N)$となる．
自由$R$-加群による解消を用いると，この同型は導来テンソル積にも延びる．
$F$は完全なのでコホモロジーを保ち，有限Tor振幅を保存する．
noetherian対象も，部分対象の昇鎖という圏論的特徴付けにより保存する．
よって$F$は$\mathcal P_U$を$\mathcal P'_U$へ送る．
準逆にも同じ議論を適用して同値を得る．

$F$を複体に項ごとに適用すると標準導来enhancementのdg同値を得る．
また$F$は単射対象を保つので，単射解消による計算から
相対RHomを$R$-複体として保つ．
従って$F\mathsf S_U$と$\mathsf S'_U F$は同じ相対双対RHom関手を表現する．
双対同型を備えた表現のYoneda一意性により所望の同型を得る．
この一意性により，同値の合成と基底変更にも整合する．
\end{proof}

$\underline{\operatorname{Eq}}_{\OO,T}(A,A')$で，
相対完全線型同値$F$と$\alpha:F(\OO)\simeq\OO'$の組を，
$\alpha$と両立する自然同型で割ったもののfppf層を表す．
これは射の同定を忘れる際に，\emph{すべての対象列の保存を条件として残した}
第\ref{sec:categorical-moduli-comparison}節の一点付き比較とは異なる．
ここでは最初に要求するのは構造対象の保存だけであり，
偶数列の保存はSerre双対から導く．

\begin{theorem}[構造対象付き同値の相対同型層]
\label{thm:abel-relative-pointed-hom}
上の代数的条件および\textup{(C1)}--\textup{(C4)}の下で，構造移送は
fppf層の標準同型
\begin{equation}\label{eq:abel-relative-pointed-hom}
 \underline{\operatorname{Isom}}_{\mathrm{gr},T}(A,A')/\mathbb G_m
 \xrightarrow{\sim}
 \underline{\operatorname{Eq}}_{\OO,T}(A,A')
\end{equation}
を与える．商はfppf層として取り，$\lambda\in\mathbb G_m$は
次数$n$上に$\lambda^n$を掛ける．
構造対象の同定を保持したHom群oidには非自明な2自己同型がなく，
同型\eqref{eq:abel-relative-pointed-hom}は同値の合成と底変更に両立する．
\end{theorem}
\begin{proof}
fppf局所で$\mathsf S_U\simeq(-2)[3]$とその相手側の表示を選ぶ．
補題\ref{lem:abel-equivalence-preserves-relative-perfect}と
$F(\OO_U)\simeq\OO'_U$により，$F$は
$E_i=\OO_U(-2i)$を$E'_i$へ送る．
相対Serre関手との可換同型を反復すると，積と両立する次数付き同型
\[
 A_U^{(2)}
 \simeq\bigoplus_{n\ge0}
 \mathcal Hom\bigl(\OO_U,(\mathsf S_U^{-1}[3])^n\OO_U\bigr)
 \xrightarrow{\sim}(A'_U)^{(2)}
\]
を得る．
定理\ref{thm:abel-relative-veronese-isom}により，これは
さらに有限\'etale局所で$A_U\simeq A'_U$へ持ち上がる．
その構造移送と$F$は偶数列のHomと合成上で一致するので，
\textup{(C2)}の相対復元により，構造対象の同定を保つ自然同型で結ばれる．
これで局所全射性を得る．

核は局所的な自然変換から計算する．
構造移送が構造対象付き恒等関手へ自然同型であるとし，
$E_i$上のスカラーを$\lambda_i\in\mathcal O_U^\times$，$\lambda_0=1$と書く．
section saturationにより$\End(E_i)=\mathcal O_U$である．
隣接Homは$(A_U)_2$であり，正の階数を持つベクトル束なので忠実である．
自然性から$\lambda_{i+1}/\lambda_i=c$は$i$によらず，
任意の$\Hom(E_{i+n},E_i)=(A_U)_{2n}$上の作用は$c^n$となる．
系\ref{cor:abel-projective-veronese-isom}によれば，このような同型は
次数スカラーを除いた同型層で自明である．逆に次数スカラーの構造移送は
次数ごとのスカラーによって構造対象付き恒等関手へ自然に同型となる．
従って\eqref{eq:abel-relative-pointed-hom}は単射でもある．

さらに，自然自己同型は各$E_i$上でスカラーとして作用し，
隣接Homの忠実性によりすべて同じスカラーとなる．
豊富列の生成性と直和・全射への自然性から圏全体でもそのスカラーである．
構造対象上の同定を保つとスカラーは$1$になる．
従って2自己同型は自明で，存在する標識付き自然同型は一意である．
Serre表示の変更は偶数軌道環の次数スカラーを変えるだけであり，
持上げの$\mu_2$も全次数スカラーに含まれる．
よって得た商層の同型は選択によらない．
構造移送の合成と基底変更との両立，および\textup{(C1)}の降下から，
局所的な比較は全体に貼り合わさる．
\end{proof}

\begin{remark}[この相対定理で解決していないこと]
\label{rem:abel-relative-pointed-scope}
本定理は，必要な相対圏論の入力を供給したとき，
体上の定理\ref{thm:abel-pointed-field-hom}を族の同型層へ拡張できることを示す．
既存の17次元全混合族で\textup{(C1)}--\textup{(C4)}が成立すること，
その成立範囲が非空な代数的開集合であること，
得られるAbel圏モジュライが有限型代数スタックであることは，なお別の課題である．
特に，ファイバーごとの有限大域次元やSerre公式から
相対RHomのperfect性・基底変更・相対Serre双対を省略しない．
具体的な有限位数楕円境界の四係数族については，
第\ref{sec:four-parameter-central-model}節で\textup{(C1),(C2)}を全係数基底上に供給する．
さらに命題\ref{prop:central-model-intrinsic-perfect-open}は既存の一般正則性入力の下で，
非空基底開上に\textup{(C3)}と\textup{(C4)}のperfect RHom・底変更を供給する．
この部分族でなお必要なのは，相対Serre双対と$(-2)[3]$という関手の同定である．

相対Serre関手が$(-2)\otimes L[3]$で表されるときは，
$L$の局所自明化により上の証明を適用する．
大域軌道環は$\bigoplus_{n\ge0}A_{2n}\otimes L^{-n}$となり得るので，
これを大域的に$A^{(2)}$と等置する必要はない．
体上の比較自体は，この相対条件一式を証明する前から
定理\ref{thm:abel-pointed-field-hom}の仮定の下で有効である．
\end{remark}

\subsection{剛体ambientの四係数族への帰結}
\label{subsec:abel-rigid-four-parameters}

\begin{corollary}[四係数族の構造対象付き圏とその自己同型]
\label{cor:abel-rigid-four-family}
命題\ref{prop:moduli-rigid-example}で指定した入力を満たす剛体ambient $B_*$と，
系\ref{rem:moduli-rigid-generic-smoothness}の非空開部分$U\subset k[u,v]_3\simeq\A^4$を取る．
$f\in U$について$A_f=B_* /(G_*+f)$と置くと，
\[
 (\qgr A_f,\OO_f)\simeq(\qgr A_{f'},\OO_{f'})
 \quad\Longleftrightarrow\quad f=f',
 \qquad
 \operatorname{Aut}_{\OO}(\qgr A_f)=1.
\]
右辺の自己同型群は完全同値の自然同型類の群である．
さらに構造対象上の同定を保つ2自己同型も自明である．
\end{corollary}
\begin{proof}
指定した入力により$A_f$は左右noetherianな中心三次CY商であり，
balanced dualizing complexは$A_f(-2)[4]$である．
局所双対\cite[Theorem~4.2(1), Corollary~4.3]{YZ}から
$H^0_{\mathfrak m}(A_f)=H^1_{\mathfrak m}(A_f)=0$である．
section写像の基本完全列により
$A_f\simeq\bigoplus_m\Hom(\OO_f,\OO_f(m))$を得る．
同じ文献による$\chi$と標準tails表示により，標準列およびその偶数部分列は豊富である．
$U$上の有限大域次元と双対複体から
$\mathsf S=(-2)[3]$となる．低次中心と$k[u,v]$上の自由性は
命題\ref{prop:moduli-rigid-example}およびその後の系の入力から得られるので，
定理\ref{thm:abel-pointed-field-hom}を適用できる．

二次包絡による復元により，$A_f$の次数付き自己同型は$B_*$の自己同型へ
一意に延びる．$\Aut_{\rm gr}(B_*)=k^\times$であり，次数スカラーは
常に$G_*+f$の直線を保つので，$\Aut_{\rm gr}(A_f)=k^\times$である．
従って次数スカラー商は自明である．また既存の同型判定から
$A_f\simeq_{\rm gr}A_{f'}$は$f=f'$と同値であり，圏の判定も従う．
\end{proof}

この系は体上の対象と射の判定である．$U$がAbel圏モジュライの
開部分や\'etale chartであること，$U$上の相対dg smooth proper性，
構造対象を忘れた圏の相異性は結論しない．


\section{相対 Jacobian 型の中心三次族}\label{sec:construction}
本節では基礎体を $k=\C$ とする．$U=ku\oplus kv$ と
$W=kx\oplus ky\oplus kz$ を置き，すべての生成元に次数１を与える．
$F\in\Sym^3(W\oplus U)$ を
\begin{equation}\label{eq:Ffull}
 F=f_3(x,y,z)+Q_2(x,y,z;u,v)+L_1(x,y,z;u,v)+g(u,v)
\end{equation}
と書く．$f_3$ は $x,y,z$ に関して次数３，$Q_2$ は次数２，
$L_1$ は次数１，$g\in\Sym^3U$ であり，それぞれ全次数３で斉次とする．

\begin{definition}\label{def:B}
$\Symm$ を可換単項式の完全対称化とする．$B_F$ を
\begin{equation}\label{eq:B}
 B_F=\frac{k\langle x,y,z,u,v\rangle}{
 \bigl([u,-],[v,-],\ [x,y]-\Symm(F_z),\
 [y,z]-\Symm(F_x),\ [z,x]-\Symm(F_y)\bigr)}
\end{equation}
で定義する．$g(u,v)$ は右辺の微分に現れない．
\end{definition}

\begin{lemma}[平面三次式の相対 PBW 性]\label{lem:pbw}
$E=V(f_3)\subset\PP(W^*)$ が滑らかなら，$B_F$ は
$k[u,v]$ 上自由であり，$x,y,z$-次数による濾過について
\begin{align}\label{eq:gr}
 \gr_{x,y,z}B_F&\simeq S_{f_3}\otimes_k k[u,v],\\
 S_{f_3}&:=k\langle x,y,z\rangle/\bigl([x,y]-\Symm((f_3)_z),\ 
\notag\\[-2pt]
 &\hspace{8em}[y,z]-\Symm((f_3)_x),\ [z,x]-\Symm((f_3)_y)\bigr).
\end{align}
さらに $S_{f_3}$ は noetherian な３次元 Koszul AS 正則 Calabi--Yau 整域で，
$B_F$ は noetherian な５次元 Koszul AS 正則整域となる．とくに
\begin{equation}\label{eq:HilbertB}
 H_{B_F}(t)=(1-t)^{-5}.
\end{equation}
\end{lemma}
\begin{proof}
Mori--Smith \cite[Theorem~4.2(1)]{MS} により，滑らかな $f_3$ から得た
$S_{f_3}$ は３次元 Calabi--Yau であり，したがって上記の
AS 正則性と Koszul 性を持つ．$u,v$ を中心的な係数とみなし，
\eqref{eq:B} の三つの関係を
$S_{f_3}$ の potential に対する次数２以下の cyclic potential の摂動と解釈する．
各 $(u,v)=(s,t)\in k^2$ への特殊化について，
Berger--Taillefer および Etingof--Ginzburg の PBW 定理
\cite[Proposition~7.3.1(ii)]{EG} により，$x,y,z$-濾過の随伴次数付き代数は
$S_{f_3}$ に等しい．

$S_{f_3}$ の斉次な語の基底を選ぶ．元の関係式を用い，語を最高 $x,y,z$-次数から
下げる帰納法により，選んだ語の持上げは $k[u,v]$ 上 $B_F$ を生成する．
この持上げの有限 $k[u,v]$-線形関係があれば，各 $(s,t)$ への特殊化で
PBW 基底間の関係を与えるので，その全係数は $k[u,v]$ の全閉点で零となり，
したがって零である．これで自由性および \eqref{eq:gr} が得られる．
特に $u,v$ は $B_F$ 上の正則中心列である．

$B_F/(u,v)\simeq S_{f_3}$ なので，中心正則元に関する
AS 正則性の上昇を二度適用すれば次元５の AS 正則性が得られる
\cite[Propositions~2.1, 2.3]{CKS}．noetherian 性は中心正則拡大の
標準定理，Koszul 性は中心一次元による change-of-rings，整域性は
\eqref{eq:gr} の整域性に従う．$u,v$ を加えた次数計数から \eqref{eq:HilbertB} を得る．
\end{proof}

\begin{lemma}[ambient の Calabi--Yau 性]\label{lem:CY}
補題\ref{lem:pbw} の仮定の下で $B_F$ の中山自己同型は恒等であり，
$B_F$ は５次元 Calabi--Yau 代数である．
\end{lemma}
\begin{proof}
AS 正則性から $B_F$ は skew Calabi--Yau であり，中山自己同型を $\nu$ とする．
中山双加群の左右の中心作用は一致するため，$\nu$ は中心 $k[u,v]$ を
各点固定する．すなわち $\nu(u)=u,\nu(v)=v$ である．
したがって各 $(s,t)\in k^2$ について $\nu$ は
$B_F/(u-s,v-t)$ に自己同型を誘導する．
Zhou--Shen--Lu の正規元による中山自己同型の特殊化公式
\cite[Propositions~2.6, 3.6]{ZSL} により，この自己同型は
特殊化代数の中山自己同型である．
一方，特殊化代数は低次 potential の摂動であり，
Etingof--Ginzburg の定理\cite[Theorem~3.3.2(i)]{EG} によって
３次元 Calabi--Yau である．特殊化代数は随伴次数付き整域
$S_{f_3}$ を持つので，単元は $k^\times$ のみである．
従って特殊化された $\nu$ は恒等である．
$B_F$ は $k[u,v]$ 上自由だから，
$\nu(x)-x,\nu(y)-y,\nu(z)-z$ が全 $(s,t)$ で零になることから
それらは零である．よって $\nu=\mathrm{id}$．
\end{proof}

\begin{lemma}[一般点での三次中心元]\label{lem:center}
$E=V(f_3)$ が滑らかな開集合のある非空開部分で
\begin{equation}\label{eq:centS}
 \Z(S_{f_3})_1=\Z(S_{f_3})_2=0,\qquad
 \dim_k\Z(S_{f_3})_3=1
\end{equation}
が成立する．その開部分をさらに縮小し，三次中心の非零元
$\Omega_{f_3}\in\Z(S_{f_3})_3$ を代数的に選ぶ．
すると\eqref{eq:B} の $Q_2,L_1$ の全係数に関して代数的に変わる
\begin{equation}\label{eq:G}
 \mathcal G_F=\Omega_{f_3}+C_2(F)+C_1(F)\in\Z(B_F)_3
\end{equation}
が存在する．$C_2$ は $x,y,z$-次数２で $u,v$ に線形，$C_1$ は
$x,y,z$-次数１で $u,v$ に二次であり，$x,y,z$-次数０の部分を零とする
規格化によって一意である．特に
\begin{equation}\label{eq:Gg}
 G_F=\mathcal G_F+g(u,v)
\end{equation}
も中心元であり，$B_F$ 上正則である．
\end{lemma}
\begin{proof}
具体的に $f_3=(x^3+y^3+z^3)/3$ では，
$S_{f_3}=k\langle x,y,z\rangle/([x,y]-z^2,[y,z]-x^2,[z,x]-y^2)$ であり，
自由語の次数１から４の関係行列を有理数上で消去すると，
次数１，２，３の中心の次元は順に $0,0,1$ になる．
この有限計算は付属 \nolinkurl{computations/independent/verify_generic17_center.py} で再現できる．
三変数 Sklyanin 代数の楕円曲線と translation の表示
\cite{ATV,EG} において，$\tau,\tau^2\ne1$ となる一般の点では
低次数中心が零であり，中心三次元の空間は一次元である．
この条件を中心性の次数４以下の有限行列の rank として書けば
非空 Zariski 開条件になり，その上では \eqref{eq:centS} の
核は階数１の代数的線束をなす．非零小行列式を反転した
chart 上で $\Omega_{f_3}$ を選べる．

$u,v$ を数値へ特殊化した三変数の非斉次 potential を考える．
\cite[Theorem~3.3.2(ii)]{EG} は $x,y,z$-濾過次数３以下に
非スカラー中心元が存在することを保証する．その最高次数の記号は
\eqref{eq:centS} により $\Omega_{f_3}$ の非零定数倍であり，
係数を規格化できる．補正項の $x,y,z$-次数２および１の係数は
それぞれ６個および３個の未知数を持ち，中心性は線形方程式
\begin{equation}\label{eq:centerlinear}
 M(\mathbf a,\mathbf b)c=d(\mathbf a,\mathbf b)
\end{equation}
で表される．ここで $\mathbf a$ は非斉次二次 potential の係数，
$\mathbf b$ は一次 potential の係数である．
零定数部分で \eqref{eq:centerlinear} の同次解が非零なら，
$x,y,z$-濾過次数２以下の非スカラー中心元を与え，
最高記号を取ると \eqref{eq:centS} に矛盾する．
よって $M$ は全特殊化で階数９を持つ．各特殊化に解が存在するので，
一定 rank のベクトル束写像の余核中で $d$ の像は全閉点で零となる．
基底を被約な多項式係数環としたからその像は零であり，
唯一の代数的な解が存在する．

$\mathbf a$ に重み１，$\mathbf b$ に重み２を与える
$\mathbb G_m$-作用と規格化の一意性から，$C_2$ は
$\mathbf a$ に一次，$C_1$ は $\mathbf a$ に二次・$\mathbf b$ に一次である．
$\mathbf a$ を $u,v$ の一次式，$\mathbf b$ を二次式に置けば
\eqref{eq:G} の斉次中心元を得る．
$\Omega_{f_3}$ は \eqref{eq:gr} の整域における非零元なので，
$G_F$ は $B_F$ の非零正則元である．
\end{proof}

\begin{theorem}[41次元の中心三次対の代数的族]\label{thm:incidence}
$U=\langle u,v\rangle$ を固定し，$F$ の35係数を動かす．
補題\ref{lem:center} の非空 Zariski 開条件を課したパラメータ空間
$\mathcal T_U$ 上，$F\mapsto(B_F,[G_F])$ は代数的な中心三次対の族となる．
商
\begin{equation}\label{eq:A}
 A_F=B_F/(G_F)
\end{equation}
は各次数で基底上局所自由であり，
\begin{equation}\label{eq:HseriesA}
 H_{A_F}(t)=\frac{1-t^3}{(1-t)^5},\qquad
 (\dim A_1,\dim A_2,\dim A_3)=(5,15,34).
\end{equation}
$U$ を $\operatorname{Gr}(2,5)$ 上で動かすと，この族のパラメータ空間は
$6+35=41$ 次元の既約な有限型空間を与える．その ambient 二次代数の
係数空間は $6+31=37$ 次元であり，一般の固定 ambient の上で
中心三次直線の選択に４次元の自由度がある．
\end{theorem}
\begin{proof}
一般の斉次三次式の係数は $\binom{7}{3}=35$ 個であり，
$F$ に含まれる $g\in\Sym^3U$ の係数は４個である．
$B_F$ の関係式は $F$ の $x,y,z$ 偏微分から決まり，$g$ には依存しない．
補題\ref{lem:center} の $\mathcal G_F$ は最高 $x,y,z$-次数成分によって規格化され，
$G_F=\mathcal G_F+g$ であるから，$g$ の異なる４係数は
同じ ambient 内の異なる中心三次直線を与える．
補題\ref{lem:pbw} および補題\ref{lem:center} により，$B_F$ は整域，$G_F$ は
中心正則である．$B_F$ の任意の次数に対して中心正則元による完全列
\[
 0\longrightarrow B_F(-3)\xrightarrow{\cdot G_F}
 B_F\longrightarrow A_F\longrightarrow0
\]
が成立し，\eqref{eq:HseriesA} が従う．
相対的には，固定次数で有限生成自由加群間の表示行列の rank が
全閉点で一定なので，次数成分は局所自由となる．
$U$ の変化は Grassmannian 上の局所 chart ごとに同じ構成で与えられ，
全体の既約次元は $\dim\operatorname{Gr}(2,5)+35=41$ である．
\end{proof}

\begin{remark}[正則性の二つの意味]\label{rem:regularity}
本定理は ambient $B_F$ の５次元 AS 正則性，$G_F$ の正則性，
$A_F$ の相対平坦性を示す．これだけで任意の閉点について
$\gldim(\qgr A_F)=3$ が従うわけではない．
量子的な楕円境界が滑らかな固定 ambient について，十分一般の
中心三次係数で tails が滑らかになるという議論には，
中心局所化と非可換 generic smoothness の結果
\cite{RSVdB,Rains} が必要である．
本稿の17次元性および復元定理は，この追加の全閉点での
Zariski 開性を仮定しない．
\end{remark}

\section{generic な17係数正規形と真の変形次元}\label{sec:normal}
ここで重要なのは，41係数から形式的に24を引くことだけではない．
実際の座標変換の一般軌道が24次元であり，17個の係数だけを持つ族が
その有理商を支配することを示す．

\subsection{中心一次平面と群作用}
\begin{lemma}[一次中心による平面の復元]\label{lem:U}
補題\ref{lem:center} の開集合では
\begin{equation}\label{eq:Z1}
 \Z(B_F)_1=U=ku\oplus kv.
\end{equation}
したがって，次数付き代数 $B_F$ の自己同型は $U$ を保存する．
\end{lemma}
\begin{proof}
$u,v$ は中心である．$w\in (B_F)_1$ が中心なら，
$x,y,z$-次数に関する最高記号は $S_{f_3}[u,v]$ の中心である．
その $x,y,z$-次数が１なら，$\Z(S_{f_3})_1=0$ に矛盾する．
よって $w\in U$ である．
\end{proof}

$P\subset\PGL(W\oplus U)$ を平面 $U$ の保存群とする．
標識された基底 $(x,y,z;u,v)$ に関して，$P$ は
\[
 \left(\begin{array}{cc}M&C\\0&N\end{array}\right),\qquad
 M\in\GL_3,\quad N\in\GL_2,\quad C\in\operatorname{Mat}_{3\times2}
\]
の同時スカラー倍を同一視した群である．したがって
\begin{equation}\label{eq:dimP}
 \dim P=9+4+6-1=18,
 \qquad\dim\PGL_5=24.
\end{equation}
Jacobian 括弧は変数変換に対し $\det M$ の重みを持つ．
したがって ambient 二次関係式の係数 $F$ の
$\Sym^3U$ 以外の31成分への Lie 作用は
\begin{equation}\label{eq:Lie}
 \delta_{M,N,C}F=(\operatorname{tr}M)F
 -\sum_{i=1}^{3}\bigl((Mx)_i+(Cs)_i\bigr)\partial_{x_i}F
 -\sum_{j=1}^{2}(Ns)_j\partial_{s_j}F,
 \qquad s=(u,v)^t,
\end{equation}
を $\Sym^3U$ で割った像である．
\emph{中心元の規格化は変数変換で binary cubic の補正を受け得る}ため，
以下の横断性は最初に ambient の31方向で確かめ，
最後に独立な中心三次の４方向を加える．
この順序なら中心元の gauge に関する余分な同変性は不要である．

\subsection{17係数の表示}
\begin{theorem}[17係数の横断的代数的族]\label{thm:17}
$M_1,M_2,M_3\in k[u,v]_2$, $g\in k[u,v]_3$ として
\begin{equation}\label{eq:F17}
\boxed{\begin{aligned}
 F_{17}={}&a(x^3+y^3+z^3)+bxyz\\
 &+ux^2+\alpha vy^2+\beta(u+v)z^2\\
 &+M_1(u,v)x+M_2(u,v)y+M_3(u,v)z+g(u,v).
\end{aligned}}
\end{equation}
この族の係数は
\begin{equation}\label{eq:count17}
 (a,b,\alpha,\beta)\quad\text{の４個},\qquad
 (M_1,M_2,M_3)\quad\text{の９個},\qquad
 g\quad\text{の４個}
\end{equation}
であり，合計17個である．
非空 Zariski 開集合 $T\subset\A^{17}$ 上で
$F=F_{17}$ を \eqref{eq:B}, \eqref{eq:Gg} に代入すると
$A_F=B_F/(G_F)$ は次数ごとに平坦な代数的族となる．
その一般点の次数付き代数同型類は17次元分だけ動く．
より正確には，この族から第\ref{sec:construction}節の
中心三次対の Rosenlicht 有理商への有理写像は支配的かつ一般に有限で，
\begin{equation}\label{eq:17mod}
 \trdeg_k k(\cI)^{\PGL_5}=41-24=17.
\end{equation}
ここで $\cI$ は構成した中心三次対の41次元既約族の像の閉包とする．
\end{theorem}
\begin{proof}
まず35係数空間での一般の cubic $F$ の $f_3$ が滑らかな
平面三次式を定めるとする．古典的 Hesse 正規形
\cite[\S2, Lemma~2.1]{AD} により，
$\GL_3$ の座標変換で $f_3=a(x^3+y^3+z^3)+bxyz$ と書ける．
$b\ne0$ とし，$x,y,z$ を $u,v$ の一次式だけ平行移動すると，
三つの交差二次項 $xy,yz,zx$ に付く二元一次式６係数を消せる．
残る３つの二元一次式が $\PP(U)$ 上に互いに異なる零点を持つ一般点では，
$\PGL(U)$ でそれらを $u,v,u+v$ に変換し，
相対スカラーで第一の係数を１に規格化すれば \eqref{eq:F17} になる．
ただし Hesse 変換の選択は有限多価であり，この議論だけで
有理商の一般次数が１であるとはいわない．

さらに代数的な横断性を有理数体で独立に検算する．
\eqref{eq:F17} の次の点を取る：
\begin{align}\label{eq:F0}
 F_0={}&3(x^3+y^3+z^3)+5xyz+ux^2+2vy^2+3(u+v)z^2\notag\\
 &+(u^2+uv+2v^2)x+(2u^2+3uv+v^2)y\\
 &+(u^2+4uv+5v^2)z+u^3+u^2v+2uv^2+3v^3.\notag
\end{align}
$\mathfrak{gl}_3\oplus\mathfrak{gl}_2\oplus\operatorname{Mat}_{3\times2}$
の19個の行列係数に対する微分 \eqref{eq:Lie} の像を取り，
これと \eqref{eq:F17} の方向のうち $g$ 以外の13方向とを
$\Sym^3(W\oplus U)/\Sym^3U$ の31個の単項式に展開する．
同時スカラー１方向を除いた $31\times31$ の行列について
\begin{equation}\label{eq:det}
 \boxed{\det M_{31}=-5538243375000\ne0}
\end{equation}
となる．行と列の取り方を付録\ref{app:rank}に固定した．
この計算は有理数体上の厳密計算であり，
付属 \nolinkurl{computations/independent/verify_generic17_slice.py} で再現できる．
したがって ambient 二次代数の31方向において，
$P$ の一般軌道は18次元で，13係数正規形は横断的である．

残りの $\Sym^3U$ の４係数は，固定 ambient の上の
$[\mathcal G_F+g]$ の４次元のアフィン chart である．
よって全中心三次対の固定 $U$ での一般軌道も18次元であり，
17係数正規形は一般軌道に横断する．
$U$ が \eqref{eq:Z1} によって intrinsic なので，
一般の同型がこれ以外の６次元の pencil 方向を
同一視してしまうことはない．
$U$ を動かすと $\cI$ の次元は41，$\PGL_5$ の一般軌道の次元は
$6+18=24$ である．Rosenlicht の有理商の定理 \cite{Rosenlicht}
と次元公式により \eqref{eq:17mod} を得る．
\end{proof}

\begin{remark}[定義域と17次元の意味]\label{rem:meaning}
この次元は $\PGL_5$ による同型軌道についての値であり，
実際に $A$ の次数付き同型類に対応することは
第\ref{sec:moduli-iso}節の復元定理\ref{thm:moduli-iso}で保証する．
$\PGL_5$ の一般 stabilizer が有限であることは示したが，
自明であることや有限次数が１であることは示していない．
また $T$ の各点で $\qgr A$ が滑らかであるとは主張しない．
\end{remark}

\section{分離型の余次元と一般混合項の不可避性}
\label{sec:split-codimension}
一楕円Sklyaninブロックと独立な中心三次式からなる分離型は、
rotation付き可換復元の明示的な成功例である。
しかし、それが本稿の17次元一般族を支配するわけではない。
その理由を中心一次平面を固定した係数空間で正確に述べる。

\begin{proposition}[分離型は余次元15]
\label{prop:split-codimension}
$W=\langle x,y,z\rangle$、$U=\langle u,v\rangle$、
$\mathcal T=\Sym^3(W\oplus U)$とする（$\dim\mathcal T=35$）。
$W$上の滑らかな三次形式$f$、$U$上の一般の三次形式$g$、
および$L\in\Hom(U,W)$からなる
\[
 F_{f,g,L}(x,u)=f(x+L(u))+g(u)
 \tag{S}\label{eq:split-family}
\]
の像$\mathcal S\subset\mathcal T$は既約な20次元構成可能集合
の閉包を持ち、$\mathcal T$内で余次元15である。
さらに、中心平面$U$を保つ本稿の放物型座標変換群
$\mathcal P$の一般軌道が18次元であるため、
$\mathcal S$の有理商は次元2である。
従って、一般17次元の中心三次対の有理商において、
分離型の像は余次元15の真の特殊部分をなす。
\end{proposition}
\begin{proof}
$f$の係数は10、$g$は4、$L$は6で、全20変数である。
$L=0$における微分は$f$と$g$の係数方向のほか、
\[
 u f_x,\ u f_y,\ u f_z,\quad
 v f_x,\ v f_y,\ v f_z
 \tag{T}\label{eq:split-tangent}
\]
を与える。$f$が滑らかなら$f_x,f_y,f_z$は定数係数で一次独立である。
実際、定数線形関係があれば$f$はある座標に依存しない錐となり、
射影平面三次曲線は特異になる。
従って\eqref{eq:split-tangent}の六方向は独立で、
$\Sym^3W$と$\Sym^3U$の14方向とも次数で分離される。
ゆえに微分の階数は20であり、$\dim\overline{\mathcal S}=20$。

群$\mathcal P$は、$W$内と$U$内の線形座標変換、
$W$への$U$のshearを含み、$F$への作用は本稿の
Poisson中心三次対の規約（式\eqref{eq:Lie}）による。
この群は$\mathcal S$を保つ。
一般の滑らかな$f$と重根のない二元三次式$g$は、
それぞれ射影的な安定化群が有限である。
$L=0$で$f$を固定するshearの無限小方向は
\eqref{eq:split-tangent}の独立性から消える。
従って$\mathcal P$の有効な一般安定化群は有限であり、
一般軌道の次元は18となる。
ゆえに有理商の次元は$20-18=2$。
全係数空間の一般有理商は第\ref{sec:normal}節により
$35-18=17$次元なので、その差は15である。
\end{proof}

\begin{corollary}[本質的な混合項の一次近似]
\label{cor:split-normal-directions}
分離点$F=f(x)+g(u)$で、上の分離型の係数空間内の法線空間は
固定した分解$W\oplus U$の下で自然に
\[
 N_{\mathcal S/\mathcal T,F}
 \simeq U\otimes\bigl(\Sym^2W/\langle f_x,f_y,f_z\rangle\bigr)
 \ \oplus\ W\otimes\Sym^2U,
 \qquad \dim N_{\mathcal S/\mathcal T,F}=6+9=15
 \tag{N}\label{eq:split-normal-directions}
\]
と同定される。前者が中心変数に一次の本質的な混合項、
後者が中心変数に二次の混合項を表す。
\end{corollary}
\begin{proof}
全三次式の混合部分は
$\Sym^2W\otimes U\oplus W\otimes\Sym^2U$である。
分離型の接空間の混合部分は、線形shearによる
$\langle f_x,f_y,f_z\rangle\otimes U$と一致する。
滑らかな$f$では三偏導関数が独立なので
その商を取れば上式を得る。
\end{proof}

\begin{proposition}[可換cubicへの忘却と$17=10+7$の接空間]
\label{prop:jacobian-forgetful-tangent}
上の分離点$F=f(x)+g(u)$で、$f$が滑らかな平面三次式、
$g$が重根のない二元三次式とする。$J(F)$を$F$のJacobianイデアルとすると、
\[
  J(F)_3^{\mathrm{quot}}:=\Sym^3(W\oplus U)/J(F)_3
   \simeq J(f)_3^{\mathrm{quot}}
     \oplus\bigl(J(f)_2^{\mathrm{quot}}\otimes J(g)_1^{\mathrm{quot}}\bigr)
     \oplus\bigl(J(f)_1^{\mathrm{quot}}\otimes J(g)_2^{\mathrm{quot}}\bigr)
\]
であり、三項の次元は順に$1,6,3$、合計10である。
ここで$J(f)_d^{\mathrm{quot}}=(\Sym W/(f_x,f_y,f_z))_d$、
$J(g)_d^{\mathrm{quot}}=(\Sym U/(g_u,g_v))_d$とする。
放物型軌道の接空間を割った17次元の変形空間から、
元の可換cubicのJacobian変形空間への自然な線形射
\[
 \Sym^3(W\oplus U)/T_F(\mathcal P\cdot F)
 \longrightarrow J(F)_3^{\mathrm{quot}}
 \tag{J}\label{eq:forgetful-Jacobian}
\]
は全射で、その核は7次元である。
分離型の2次元商に由来する方向の像は1次元、
その余次元15の法線からの像は9次元であり、
法線の核の次元は6である。
\end{proposition}
\begin{proof}
$f$と$g$の勾配は別変数の正則列だから、Jacobian環は
$J(f)^{\mathrm{quot}}\otimes J(g)^{\mathrm{quot}}$に分解し、
Hilbert級数はそれぞれ$(1+t)^3$と$(1+t)^2$となる。
従って次数3には$(3,0),(2,1),(1,2)$の三組だけが寄与し、
次元は$1+3\cdot2+3\cdot1=10$である。
$\mathcal P$の軌道の無限小作用は$F$のスカラー倍と
$X_j\partial_i F$の線形結合である。
Euler等式により$F\in J(F)_3$なので、
\eqref{eq:forgetful-Jacobian}はwell-definedであり全射である。
その次元差は$17-10=7$。
分離型の接空間の純三次方向は$\Sym^3W\oplus\Sym^3U$、
残りは$u f_i,v f_i$なので、Jacobian環への像は
$J(f)_3^{\mathrm{quot}}$の1次元のみである。
さらに系\ref{cor:split-normal-directions}の法線の第一項は
$J(f)_2^{\mathrm{quot}}\otimes J(g)_1^{\mathrm{quot}}$の6方向に写り、
第二項は$J(f)_1^{\mathrm{quot}}\otimes J(g)_2^{\mathrm{quot}}$の3方向に写る。
よって法線からの像は9次元、核は$15-9=6$である。
\end{proof}
\begin{remark}[幾何学的17次元候補との関係]
\label{rem:forgetful-geometry-not-model}
この接空間の$17=10+7$と、幾何学的候補
$(Y,\Lambda,\xi)$の$10+6+1$は整合的である。
実際、元の可換cubicを固定するPoisson方向は
超平面pencilの6方向とPoissonスカラーの1方向を持つ。
しかしこの計算だけでは、量子化後の可換cubic$Y$、平面$\Lambda$、
楕円曲線上の直線束$\xi$を関手的に構成したことにはならない。
\end{remark}

\begin{proposition}[七次元の核の幾何的分解]\label{prop:vertical-seven}
命題\ref{prop:jacobian-forgetful-tangent}の仮定の下で、
$W=\langle x,y,z\rangle$、$U=\langle u,v\rangle$を
\emph{一次形式の空間}とする。分離点$F=f(x)+g(u,v)$で、
\[
 \mathsf T_{F,U}:=\Sym^3(W\oplus U)/T_F(\mathcal P\cdot F),
 \qquad R_F=\Sym(W\oplus U)/(F_x,F_y,F_z,F_u,F_v)
\]
とおく。一次形式の部分空間$U$で指定される平面
$\Lambda=V(U)\simeq\PP^2$の変形空間は
$T_{\Lambda}\operatorname{Gr}(2,W\oplus U)
\simeq\operatorname{Hom}(U,W)$である。
$L\in\operatorname{Hom}(U,W)$の基底表示を
$L(u)=\ell_u(x)$、$L(v)=\ell_v(x)$とすれば、
次の短完全列が成立する：
\[
 0\longrightarrow
 \mathbb C[f]\oplus T_{\Lambda}\operatorname{Gr}(2,W\oplus U)
 \xrightarrow{\ \iota\ } \mathsf T_{F,U}
 \xrightarrow{\ \mathrm{Jac}\ } (R_F)_3
 \longrightarrow0,
 \tag{V}\label{eq:vertical-seven-exact}
\]
\[
 \iota(a,L)=[\,a f+\ell_u(x)g_u(u,v)+\ell_v(x)g_v(u,v)\,].
 \tag{VM}\label{eq:vertical-moving-plane}
\]
左端の一次元方向はPoisson双ベクトルの大きさを変える方向であり、
六次元方向は可換cubicそのものを変えずに指定した平面を動かす方向である。
\end{proposition}
\begin{proof}
$F$のJacobianイデアルの三次部分を変数の双次数に分けると、
\[
 J(F)_3=
 W\langle f_x,f_y,f_z\rangle\ \oplus\
 U\langle f_x,f_y,f_z\rangle\ \oplus\
 W\langle g_u,g_v\rangle\ \oplus\
 U\langle g_u,g_v\rangle.
\]
$f$が滑らかな平面三次式であることから
第一項の次元は$10-\dim(R_f)_3=9$、第二項は6である。
$g$が重根を持たない二元三次式なので
第三項は6、第四項は4であり、全体の次元は25となる。
一方、$\mathcal P$の無限小作用は、
第一項において$\mathrm{SL}(W)$の8次元方向を占め、
第二項ではshearによる全6方向を占め、
第四項では$\mathrm{GL}(U)$による全4方向を占める。
ここで$\mathrm{GL}(U)$が$\Sym^3U$に作用する微分は同型である：
重根のない三点の組の射影的安定化群が有限だからである。
従って放物型軌道の次元は$8+6+4=18$であり、
$W\langle g_u,g_v\rangle$の第三項には寄与しない。
また$\mathbb Cf\subset W\langle f_x,f_y,f_z\rangle$は
$\mathrm{SL}(W)$の8次元接方向と独立である。
実際、滑らかな平面三次式の判別式は$\mathrm{SL}(W)$不変で
$f\mapsto(1+\varepsilon)f$で非零の一次変化を持つ。
従って$J(F)_3/T_F(\mathcal P\cdot F)$は
$\mathbb C[f]\oplus W\langle g_u,g_v\rangle$に一致する。
$g_u,g_v$は二元二次式として一次独立なので、
後者は$\operatorname{Hom}(U,W)$と同一視される。
最後に$\operatorname{Gr}(2,W\oplus U)$の平面変形は
$U\mapsto\langle u+\varepsilon\ell_u(x),
 v+\varepsilon\ell_v(x)\rangle$で記述される。
$g$へこの置換を適用したときの一次変化が
\eqref{eq:vertical-moving-plane}の第二項に一致する。
$F$のJacobian商への射は全射だから、短完全列が得られる。
\end{proof}

\begin{proposition}[PoissonスカラーをPicardパラメータと直接同一視できない]\label{prop:poisson-picard-obstruction}
$X\subset\PP^4$を滑らかな可換cubic、$\Lambda\simeq\PP^2$を
$C=X\cap\Lambda$が滑らかな平面三次曲線となる平面とする。
$X$と$\Lambda$を固定し、非零Poisson構造のスカラー倍
$\lambda\pi$、$\lambda\in\mathbb G_m$だけを変化させる。
この一変数族から$\operatorname{Pic}^0(C)$への任意の
有理写像は定数である。
とくに、同じ$X$と$\Lambda$を保つような
Poissonパラメータ$\lambda$と楕円Picardパラメータ$\xi$の間には、
支配的な有理対応を与える一価の有理写像は存在しない。
\end{proposition}
\begin{proof}
$\operatorname{Pic}^0(C)$は滑らかな射影的種数1曲線である。
$\mathbb G_m\dashrightarrow\operatorname{Pic}^0(C)$という有理写像は、
領域を$\PP^1$に埋め込めば、曲線上の有理写像に関する
正規性とターゲットの固有性により
$\PP^1\to\operatorname{Pic}^0(C)$に延長される。
標数0で$\PP^1$から種数1曲線への非定数射は
Riemann--Hurwitz公式に反する。
従って有理写像は定数である。
この議論は$(X,\Lambda)$を変数とする基底の
一般幾何ファイバー上でも成立する。
\end{proof}
\begin{remark}[復元に必要な非自明な幾何変化]\label{rem:picard-obstruction-scope}
滑らかなcubic threefold上の非零Poisson構造を射影化すると
$\operatorname{Gr}(2,5)$になることは既知である
\cite{LPT,Mayanskiy}。
非零スカラーまで保持するとその錐から零切断を除いた
$\mathbb G_m$束であり、固定cubic上で$6+1$次元を持つ。
一方$(Y,\Lambda,\xi)$の$\xi$は固定$(Y,\Lambda)$上では
楕円曲線$\operatorname{Pic}^0(Y\cap\Lambda)$を動く。
従って$17=10+6+1$という次元一致を、
可換cubicと平面を両方固定したまま
$\mathbb G_m$パラメータを楕円曲線の点へ置き換える
大域的な有理同一視と解釈してはならない。
実際の量子化後の可換復元では、元の$X$と異なる$Y$を選ぶか、
平面やmarkingを非自明に変化させる必要があり得る。
一ブロックHesse模型の$\vartheta(h,t)-t\ne0$は
そのような基底の移動が起こる明示的な例である。
この命題は一般のrotation付き可換復元の存在を否定しない。
また、有限被覆を介する多価の代数的対応の不存在も主張しない。
\end{remark}

\begin{remark}[可換復元への帰結]
\label{rem:split-implication}
分離型のrotation付き可換復元が成立しても、
それだけでは17次元一般点の可換復元は得られない。
線形shearで除去できない本質的な混合項に対する
相対Orlov比較・smooth proper性・rotationの同定が必要である。
一般混合族の残余圏の幾何性とsmooth proper性は、
定理\ref{thm:kuznetsov-cubic-moduli-morphism}で別の方法により非空開上で証明する。
そこで得るMorita同値が指定されたrotationも保つことは、追加の問題である。
ここで求めた余次元は、中心平面を指定する
Poisson中心三次対の\emph{有理商}についての結論であり、
無標識導来圏やKuznetsov成分のモジュライの余次元だと
読み替えてはいけない。
\end{remark}

\section{全混合係数を保つ中心三次式と規格化}
\label{sec:pair-explicit-centers}
中心元を加える操作は同じ周囲の中でも三次対を変えるので、
その中央スカラーを含む規格化を固定する。
以下ではRainsの中心計算\cite[\S9.1]{EG}を相対Hesse chartで用い、
命題\ref{prop:pair-rains-polynomial}に可換点で正則な規格化を与える。
ここでは対の復元に必要な全係数の表示と、関係イデアル内の直接の証明を与える。

\subsection{相対Hesse chartの有限式}
$R=k[u,v]$、$L_i\in R_1$、$M_i\in R_2$、$g\in R_3$として
\[
 F=\frac{x^3+y^3+z^3}{3}+\lambda xyz
   +\frac12(L_1x^2+L_2y^2+L_3z^2)+M_1x+M_2y+M_3z+g
\]
を考える。$\lambda^3\ne-1$なら最高次平面三次式は滑らかである。
任意の滑らかな最高次部分はHesse形へ移せる。
$\lambda=0$でも、$\omega^3=1$, $\omega\ne1$, $c^3=1/3$として
$(x,y,z)=c(X+Y+Z,X+\omega Y+\omega^2Z,X+\omega^2Y+\omega Z)$
とすればHesse係数は2となるので、必要なら$\lambda\ne0$にできる。
一般の二次混合項の交差係数を$C_{12},C_{23},C_{31}\in R_1$とすると、
$x=X-C_{23}/\lambda$, $y=Y-C_{31}/\lambda$, $z=Z-C_{12}/\lambda$
で交差三項を消せる。新たな対角係数は$L_i=\widetilde L_i+2a_i$、
一次係数は$M_i=\widetilde M_i+a_i^2+\widetilde L_i a_i-\lambda a_ja_k$である。
定数は$g$へ吸収する。従ってこのchartは全21混合方向を扱い、
モジュライ上の一意な座標選択を主張しない。

$q=1+\lambda t$と置き、$R$を中心として
\[
 xy-qyx=t(z^2+L_3z+M_3),\quad
 yz-qzy=t(x^2+L_1x+M_1),\quad
 zx-qxz=t(y^2+L_2y+M_2)
\]
で$B_t$を定める。略記を
\begin{align*}
 A&=q^3-t^3,& Q&=q^2+q+1,\\
 K&=q^4+q^3+q^2-(2q+1)t^3,& J&=Q+(q+2)t^3,\\
 H_1&=(q+1)(q^3+q^2+q-t^3),\\
 H_2&=q^3+2q^2+2q+1+(q+1)t^3,\\
 H_3&=q^4+2q^3+2q^2+q-(q+1)t^3
\end{align*}
とする。
\begin{proposition}[可換点で正則なRains規格化]
\label{prop:pair-rains-polynomial}
次の有限式は$D(t)=q(q^3-t^3)(q+1)(1+t^3)\ne0$上で中心正則元となる：
\begin{align*}
 6(G_t-g)={}&(q+1)\{A(x^3+z^3)+(1+t^3)y^3\}\\
 &+q(q+1)(\lambda Q-2t^2)zyx\\
 &+K(L_1x^2+L_3z^2)+JL_2y^2\\
 &-(q+1)t^2(qL_1zy+L_2zx+qL_3yx)\\
 &+(H_1M_1-qt^3L_1^2-qt^2L_2L_3)x\\
 &+(H_2M_2+t^3L_2^2-qt^2L_1L_3)y\\
 &+(H_3M_3-qt^3L_3^2-qt^2L_1L_2)z.
\end{align*}
$D(0)=2$、$G_0=F$であり、$B_t/(G_t)$はHilbert級数
$(1-s^3)/(1-s)^5$を持つ平坦な次数付き族である。
\end{proposition}
\begin{proof}
$L_i,M_i$を独立中心変数、重みをそれぞれ1,2として扱う。
重み4の自由語339個に対する二次関係の全倍数162本を
$\mathbb Q(q,t)$上で並べるとrank144となり、表示式の三交換子はその行空間に入る。
\texttt{verify\_central\_pair\_centers.py}がこの関係行列と式を再生成して厳密還元する。
係数を数値へ代入した補間ではない。
$q+1\ne0$では関係式を$(q+1)/2$で割ると、
量子係数$h=2t/(q+1)$、最高次Hesse係数$\lambda$の対称化potentialになる。
補題\ref{lem:pbw}の一様PBWとCalabi--Yau証明をこの係数で適用できる。
この平坦性により分数体上の中心性は$D\ne0$全体へ降下する。
次数3の関係行列の17次小行列式は$-q^5(q^3-t^3)^2$であり、
整列三次語の基底で最高次成分は
\[
 (q+1)(1+t^3)(x^3+z^3)+(q+1)(q^3-t^3)y^3
 +(q+1)(\lambda Q-2t^2)xyz
\]
となる。これは非零なので、整域性から$G_t$は正則である。
商のassociated gradedは最高次中心元による商と$R$のtensor積となり、
その基底の持上げで自由性とHilbert級数を得る。
\end{proof}
\begin{remark}
EG掲載のRains式の$y$係数の$q^3b_2$は$q^3$と読む必要がある。
元のMAGMA資料\cite{RainsMagma}と上の独立な関係行列で照合する。
中心性自体は既知入力の具体化であり、全特殊パラメータのtails正則性は含まない。
\end{remark}

\subsection{Fermatの六二次係数に対する自由代数恒等式}
$(x_1,x_2,x_3)=(x,y,z)$とし、添字$(i,j,k)$を巡回させる。
$L_i,e_i$の重みは1、$m_i$の重みは2とする。
\[
 F=\frac13\sum_i x_i^3+\frac12\sum_iL_ix_i^2
       +\sum_i e_i\Symm(x_jx_k)+\sum_i m_ix_i,
 \qquad g_i=x_i^2+L_ix_i+e_jx_k+e_kx_j+m_i.
\]
$[x,y]=hg_3$, $[y,z]=hg_1$, $[z,x]=hg_2$で関係を定める。
\begin{lemma}[全Artin商での形式PBW]
全ての$N\ge1$で、$h^N=0$上のこの代数は順序付き語$x^ay^bz^c$を
係数環上の基底に持つ。
\end{lemma}
\begin{proof}
$yx\mapsto xy-hg_3$, $zx\mapsto xz+hg_2$, $zy\mapsto yz-hg_1$を用いる。
残る$h$次数と転倒数の辞書式順序で停止する。
唯一のoverlap $zyx$の二経路の差は
$h([x,g_1]+[y,g_2]+[z,g_3])=0$である。
最後は自由代数恒等式で、交差係数の交換子が対で消える。
停止・局所合流性から基底を得て、各$N$の商写像とも両立する。
有限の非零$h$における整列基底を仮定していない。
\end{proof}
\begin{theorem}[四次補正で止まる全係数中心式]
\label{thm:pair-full-fermat-center}
$\Omega_h=\sum x_i^3/3-h^2\sum_{\sigma\in S_3}\sigma(xyz)/9$とする。
$L_i=0$のとき
\[
 Z_h=\Omega_h+\sum_i e_i\Symm(x_jx_k)+\sum_i m_ix_i
       +\frac{h^2}{6}\sum_i e_i^2x_i
       +\frac{h^4}{18}\sum_i(e_ix_i^2-e_je_kx_i)
\]
は全有限$h$・全係数で中心的である。一般の$L_i$に対しては
\begin{align*}
 X_i&=x_i+L_i/2,\qquad
 n_i=m_i-L_i^2/4-(e_jL_k+e_kL_j)/2,\\
 c_{\rm cl}&=\sum_iL_i^3/12+\sum_ie_iL_jL_k/4-\sum_im_iL_i/2,\\
 W_h(x;L,e,m)&=Z_h(X;e,n)+c_{\rm cl}
\end{align*}
が中心元で、$W_0=F$となる。
\end{theorem}
\begin{proof}
$L=0$の自由代数で$r=[x,y]-hg_3$, $s=[y,z]-hg_1$, $t=[z,x]-hg_2$とする。
次の四つの式を定める：
\begin{align*}
 C_0={}&-(ry^2+yry+y^2r)/3+(tz^2+ztz+z^2t)/3+m_3t-m_2r\\
       &-e_3(rx+xr)/2+e_2(tx+xt)/2+e_1(yt-rz+[x,s]/2),\\
 C_1={}&([yz+zy,s]+e_3[t,z]+e_2[r,y])/6,\\
 C_2={}&(r[x,z]-sx^2-t(xy+yx)-xrz+xsx+xty\\
       &\hspace{6em}+2ytx-yxt+zxr)/18,\\
 C_3={}&([r,y^2]+[t,z^2]+e_1[s,x])/18.
\end{align*}
自由代数内で$[Z_h,x]=C_0+hC_1+h^2C_2+h^3C_3$である。
右辺は41個の関係式倍であり、有理数係数の展開で直接検証できる。
巡回置換が残る二つの交換子を与える。
平行移動については$X_i^2+e_jX_k+e_kX_j+n_i=g_i$であり、
$\sum X_i^3/3+\sum e_i\Symm(X_jX_k)+\sum n_iX_i+c_{\rm cl}=F$は
自由代数の恒等式なので、全ての主張が従う。
\end{proof}
\begin{remark}[規格化と不可欠な四次項]
$L=0$では$\det\operatorname{Hess}F=8xyz-2\sum e_i^2x_i+2e_1e_2e_3$である。
従って$Z_h=\Symm(F-h^2\det\operatorname{Hess}F/12)+h^2e_1e_2e_3/6
+h^4\sum(e_ix_i^2-e_je_kx_i)/18$。
最後の項は省けない。$F=\sum x_i^3/3+rxy$では省いた候補の$x$との交換子は
$-h^5(ry^2z+r^2xz)/9\pmod{h^6}$で、形式PBWにより非零である。
中央平行移動$x\mapsto x+a$に対し
$L_i'=L_i+2a_i$, $m_i'=m_i+L_ia_i+a_i^2+e_ja_k+e_ka_j$とすると
\[
 W_h(x+a;L,e,m)=W_h(x;L',e,m')+F(a).
\]
対角Rains式（$\lambda=0,q=1,t=h$）とpreferred式の差は
\[
 G_h^{\rm R}-W_h(x;L,0,m)
 =\frac{h^2}{12}L_1L_2L_3+
   \frac{h^3}{6}(L_1m_1-L_2m_2+L_3m_3).
\]
この中央三次差を捨てると、同じ周囲の中心対を誤同定する。
\end{remark}
\begin{proposition}[有限微分性と形式正則商]
順序付きPBW固定化の積$\mu_h$の$h^r$係数は、係数$L,e,m$を微分しない
$x,y,z$の双微分作用素で、入力の全微分次数は高々$2r$である。
また中心商は各$h^N=0$上で係数環上自由となる。
\end{proposition}
\begin{proof}
Barmeier--Wang\cite[Theorem 10.3, Lemma 10.9]{PairBW}の還元グラフ表示では
$h^r$に寄与する頂点数は$r$であり、各頂点から入力への微分は高々二つ。
非零補正はactive三変数の三対だけなので、中央九係数の微分は現れない。
PBWの合流性によりこの積は指定順序の実際の積である。
中心性と結合性は同じ固定化の曲率付きMaurer--Cartan方程式を与える。
商の自由性は、$F$が$x$のmonic三次の定数倍であることから直接示せる。
$x$次数二以下の単項式空間を$V$とすると、$(a,v)\mapsto Fa+v$は線形同型。
その変形$(a,v)\mapsto W_h*a+v$は元の同型と$1+hT$の合成であり、
$h^N=0$では有限幾何級数で逆を持つ。
この分裂は任意の中央係数代入後にも保たれ、$g(u,v)$を加えても同じである。
\end{proof}

\begin{corollary}[可換点を含む一様な有限係数量子化]
滑らかな最高次平面三次式$f_3$を固定すると、$\Delta(0)\ne0$なる一変数多項式
$\Delta(h)$が存在し、その非零範囲では全21混合係数を同時に保つ中心元$G_h$を代数的に取れる。
$G_0=F$で、中心商は次数ごとに平坦である。
さらに$h\ne0$なら$\dim Z(B_h)_1=2$, $\dim Z(B_h)_3=5$とできる。
\end{corollary}
\begin{proof}
Hesse座標の選択後に$t=2h/(2-\lambda h)$, $q=(2+\lambda h)/(2-\lambda h)$を用いると、
命題\ref{prop:pair-rains-polynomial}が対称化Jacobian関係の中心元を与える。
その分母と$D(t)$の分子は$h=0$で非零で、混合係数に依存しない。
中心の低次数は交換子行列を$h$で割って調べる。
$h=0$でこれは$f_3$のJacobian Poisson中心を計算し、次数1,2,3のrankは3,6,9である。
その非零minorをさらに反転する。三次の一般rankは中心元の存在から9を超えず、
従って三次核はrank1である。
低次補正の未知数9個の行列は、active次数の順に、上の次数2・1の
divided-commutator行列を対角に持つ三角blockを含む。
同じminorを反転すれば、全混合係数に関する唯一の補正を代数的に解ける。
これら有限個の因子の積を$\Delta(h)$とする。
相対PBWで最高次中心元を引く帰納法から中心の次数1,3の等式を得る。
$h=0$ではdivided-commutatorの核はPoisson中心であり、実際の可換中心は
次数1で5、次数3で35に跳ぶ。この二つを同一視していない。
\end{proof}

\section{Poisson三次incidenceと37次元の周囲成分}
\label{sec:pair-generic-component}
\subsection{\texttt{poisson-p4}の三次incidenceとの照合}
本節は関連原稿\cite{PoissonP4Incidence}の三次延長定理と
ファイバー計算を用いる。射影四空間の全成分数やその完全分類は仮定しない。
本稿ではPoisson構造のスカラー倍を\emph{同一視しない}。
同原稿の射影化した次元36に対し、ここでの錐の次元は37である。
この区別は、有限量子パラメータを持つ代数のモジュライを考えるために必要である。

$V^*=\langle x,y,z,u,v\rangle$、
$\Omega=dx\wedge dy\wedge dz\wedge du\wedge dv$とし、
\begin{equation}\label{eq:rat113-jacobian}
 \iota_{P(F,L,M)}\Omega=dF\wedge dL\wedge dM,
 \qquad \deg F=3,\quad \deg L=\deg M=1
\end{equation}
と置く。$C_{113}^{\rm aff}$はこの全像の被約閉包である。
標準発散作用素を$\delta$と書く。
閉三形式の分解可能性から$\delta P=0$かつ$[P,P]=0$である。

\begin{proposition}[滑らかな三次からの像と四次元ファイバー]
\label{prop:poisson-cubic-ambient-image}
滑らかな三次超曲面$X=V(F)\subset\PP^4$をPoisson部分多様体にする
非零ambient Poisson構造は、必ず\eqref{eq:rat113-jacobian}型で一般階数二である。
$X$上のPoisson構造を固定した延長スキームは被約一点である。
非零ambientを固定して$X$を忘れる向きのファイバーは、
尺度を規格化した後の
\begin{equation}\label{eq:poisson-cubic-fiber}
 \{F+g(u,v):g\in\Sym^3\langle u,v\rangle,\ V(F+g)\text{ は滑らか}\}
\end{equation}
という$\A^4$の開部分である。
従って三次Poisson対の像は$C_{113}^{\rm aff}$に稠密であり、
その次元は$41-4=37$である。
\end{proposition}
\begin{proof}
$\wedge^2T_{\PP^4}\simeq\Omega^2_{\PP^4}(5)$なので、三次超曲面への
制限の核は$H^0(\PP^4,\Omega^2_{\PP^4}(2))=0$である。
Loray--Pereira--TouzetおよびMayanskiyの分類\cite{LPT,Mayanskiy}により、
$X$上の非零Poisson構造は分解可能な定数二形式に対応し、
\eqref{eq:rat113-jacobian}の制限で与えられる。
線形な制限方程式の核が零であるから、延長はスキームとしても一点である。

射影bivectorの二次斉次liftはEuler場$E$との外積だけ曖昧だが、
$\delta(E\wedge Y)=5Y-(\operatorname{tr}Y)E$より
発散零の代表は一意である。その代表の一次Casimir空間は
$U=\langle L,M\rangle$に等しい。実際、独立な第三の一次Casimirがあれば
$F$は三つの一次式だけの多項式になり、$V(F)$は錐となって滑らかさに反する。
従ってambientから$U$を回復する。同じ$U$に対する二つのJacobianが
比例する条件は$F'-cF\in\Sym^3U$である。尺度を除くと
\eqref{eq:poisson-cubic-fiber}を得る。
滑らかな三次式は係数空間の稠密開であり、
Grassmannianの6次元と、非射影化した相対Jacobianの31係数を合わせて
像の次元は37である。Poisson対の次元41は、三次式の射影係数34と
固定した三次上のPoisson錐の次元7の和でもある。
\end{proof}

\subsection{全Poissonスキーム内の接空間と一般被約性}
\begin{theorem}[\texorpdfstring{$\mathrm{Rat}(1,1,3)$}{Rat(1,1,3)}の独立な成分証明]
\label{thm:rat113-smooth-component}
$C_{113}^{\rm aff}$は、五変数の発散零二次Poisson構造の
\emph{全スキーム}の37次元既約成分であり、その一般点は滑らかである。
従って射影四空間の非射影化Poissonスキームでも同じ結論が成り立つ。
\end{theorem}
\begin{proof}
二次bivectorの係数空間は$10\cdot15=150$次元である。
全ての狭義Jacobi方程式$\frac12[P,P]=0$の350係数と、
発散$\delta P=0$の25係数を用いる。
次の整数係数三次式を取る：
\begin{align}\label{eq:rat113-test-F}
 F_*={}&3(x^3+y^3+z^3)+5xyz+ux^2+2vy^2+3(u+v)z^2\notag\\
 &+(u^2+uv+2v^2)x+(2u^2+3uv+v^2)y
 +(u^2+4uv+5v^2)z.
\end{align}
$P_*=P(F_*,u,v)$における方程式の微分を$M$とする。
列は十個の辞書式添字対と15個の二次単項式、
行は十個の辞書式添字三つ組と35個の三次単項式、
続いて五つの発散成分と五つの一次単項式の順である。
例えばJacobi成分の微分は
\begin{equation}\label{eq:rat113-jacobi-differential}
 (DQ)_{ijk}=\sum_{\mathrm{cyc}(i,j,k)}\sum_{l=1}^{5}
 \left(Q_{il}\partial_l(P_*)_{jk}
       +(P_*)_{il}\partial_lQ_{jk}\right)
\end{equation}
で与えられる。添字を$1,\ldots,5$として変数順$x,y,z,u,v$を使う。

一方、$F_*$のうち純$u,v$単項式を除いた31係数、および
$u\mapsto u+s_1x+s_2y+s_3z$、
$v\mapsto v+t_1x+t_2y+t_3z$の六係数による族の微分を$S$とする。
$M$は$375\times150$、$S$は$150\times37$の整数行列であり、
Jacobian族が恒等的にPoissonかつ発散零なので$MS=0$である。
再現コードはこの積も整数上で直接確認する。

\nolinkurl{computations/verify_rat113_tangent.py}は両行列を元の公式から生成し、
\nolinkurl{computations/rat113_tangent_certificate.json}の行・列添字で
指定した小行列式について
\begin{equation}\label{eq:rat113-minor-certificates}
 \det M_{113}\equiv844194,\qquad
 \det S_{37}\equiv948131\pmod{1000003}
\end{equation}
を得る。これらは整数行列式の\emph{非零}証明である。
従って$\rank_{\mathbb Q}M\ge113$、$\rank_{\mathbb Q}S=37$であり、
整数上の恒等式$MS=0$と合わせて$\rank_{\mathbb Q}M=113$となる。
有限体での零から特性零の恒等式を推論してはいない。
よって全スキームの接空間は$150-113=37$次元である。

この点を通るJacobian族の像は、$S$の階数37から局所次元37を持つ。
したがって全スキームの局所次元と接空間次元が一致し、この点は正則点である。
Jacobian源は既約であるから、その37次元像の閉包がこの点を通る
唯一の既約成分となり、滑らかな開部分を持つ。

最後に、発散零二次bivectorの空間は
$H^0(\PP^4,\wedge^2T_{\PP^4})$に線形同型である。
発散零の$P$に対して射影Jacobiが零なら、
$[P,P]=E\wedge A$と$\delta[P,P]=0$を用いて
$0=5A-E\wedge\delta A$、従って$E\wedge A=0$を得る。
これは係数の任意の基底環上で成立するため、Jacobiイデアルも一致する。
従って上記は射影Poissonスキームそのものについての証明である。
\end{proof}

\subsection{規格化CY変形の局所成分への移行}
Pym\cite[Definition 2.8, Proposition 2.9, Theorem 2.10]{Pym}の
$\operatorname{QDef}_0(k^5)$は、\emph{可換多項式環の点における}
Koszul Calabi--Yau変形の芽である。
同文献第2節のこれらの結果は任意の$\dim V$について述べられており、
第3節の四生成元の六成分分類とは区別する。
以下の成分性は、この局所変形空間に関する主張である。

$V=k^5$とし、$V^{\otimes2}=\wedge^2V\oplus\Sym^2V$で規格化して
\[
 R_\Pi=\langle a\otimes b-b\otimes a-\Symm(\Pi(a,b)):a,b\in V\rangle,
 \qquad B_\Pi=T(V)/(R_\Pi)
\]
と置く。
\begin{proposition}[graph成分と一般被約性]\label{prop:quantum-113-component}
$\Pi\mapsto R_\Pi$は$C_{113}^{\rm aff}$を二次関係Grassmannianの
標準アフィンchartへ閉埋込みする。その像$Q_{113}$は37次元で、
可換点の規格化CY変形の芽において対応する被約既約成分をなす。
この成分の一般点における全局所環は被約かつ正則である。
これは全AS正則$\mathbb Z$-代数の大域モジュライの成分性を意味しない。
\end{proposition}
\begin{proof}
$R_\Pi$は$\wedge^2V\to\Sym^2V$のgraphなので$\Pi$を一意に回復でき、
標準liftが線形であることから閉埋込みとなる。
最高次平面三次制限が滑らかな相対Jacobian点は稠密であり、
その点では補題\ref{lem:pbw}と\ref{lem:CY}がKoszul CY代数を与える。
CY変形の条件はこの変形芽では閉条件なので、閉包の芽も含まれる。
それを含む被約既約成分の芽を$D$とする。
一次項が恒等な形式比較の下で$D$の接錐の被約台はあるPoisson既約成分$C_D$である。
graph像は斉次錐なので
$C_{113}^{\rm aff}=TC_0Q_{113}\subset TC_0D=C_D$となる。
定理\ref{thm:rat113-smooth-component}の成分の極大性から両者は等しく、次元も37である。
従って$D$と$Q_{113}$の被約な既約芽は等しい。
Pymの比較は被約台だけの比較ではなく、Poissonスキームと規格化CY変形の
形式的、従ってある解析的な芽の同型である。
定理\ref{thm:rat113-smooth-component}による一般正則性も対応する成分へ移る。
一般点では他の成分との交叉を除けるので、全局所環も正則となる。
これはKontsevich級数自体がgraph式と一致するという主張ではない。
\end{proof}

\begin{lemma}[相対Jacobian括弧の全Poisson中心]
$R=k[u,v]$とし、$F\in R[x,y,z]$の最高active次数部分$f_3$が既約とする。
$R$を中心、$\{x,y\}=F_z$などと定めると$Z_{\rm Pois}(R[x,y,z])=R[F]$である。
\end{lemma}
\begin{proof}
中心元の最高次部分$g_m$は$df_3\wedge dg_m=0$を満たす。
Euler場との縮約から$3f_3dg_m=mg_mdf_3$。
$f_3$は素元で、少なくとも一つの非零偏微分を割らないので$f_3\mid g_m$である。
これを繰り返すと$m=3j$、$g_m=cf_3^j$ ($c\in R$)となる。
$g-cF^j$へ次数帰納法を適用すればよい。逆包含はCasimir性による。
\end{proof}
\begin{theorem}[量子成分の関数体上の中心]
\label{thm:pair-generic-center}
$K=k(Q_{113})$とし、一般ファイバーを$B_\eta$とすると
\[
 Z(B_\eta)=K[a,b,c],\qquad\deg a=\deg b=1,\quad\deg c=3.
\]
$B_\eta$は整域で、中心の次数1,2,3の次元は2,3,5となる。
全次数の中心が一つのZariski開部分で一定であるとは主張しない。
\end{theorem}
\begin{proof}
中心変数を厳密に保つ相対Jacobian量子化と中心元$G_\hbar\equiv F\pmod\hbar$を取る。
次数$d$の形式中心元を$\hbar$で還元すると前補題のPoisson中心に入り、
$u,v,G_\hbar$の同じ重みの多項式を引けば差は$\hbar$で割れる。
反復と完備性から
\[
 Z(B_\hbar)_d=\bigoplus_{i+j+3l=d}k[[\hbar]]u^iv^jG_\hbar^l.
\]
独立性は関係式の最小$\hbar$次数を取ることで従う。

この計算を成分の一般点へ戻す。
標準次数付き整域$A$の頂点完備化から、支配的パラメータ表示$T\to\Spec A$を用いて
\[
 \widehat A\longrightarrow k(T)[[\hbar]],\quad
 \sum a_n\longmapsto\sum\hbar^n\phi^*a_n
\]
を作ると、最初の非零斉次係数が消えないので単射である。
これをPoisson錐と形式量子化比較に適用し、$K$を含む体
$L=k(T)((\hbar))$を得る。$T$はpencilの6座標と三次式の31係数で取れる。
各次数の中心は有限自由空間の交換子行列の核なので体拡大と可換し、
$Z(B_\eta)_d\otimes_K L=L[u,v,G_\hbar]_d$となる。
$Z(B_\eta)_1$の基底$a,b$を取り、$c\in Z(B_\eta)_3\setminus\Sym^3\langle a,b\rangle$を選ぶと、
$L$上では$c=\lambda G_\hbar+r(u,v)$ ($\lambda\ne0$)である。
従って$K[a,b,c]\to Z(B_\eta)$は忠実平坦な体拡大後に同型、ゆえに元から同型である。
整域性も形式PBW量子化への埋込みから従う。
\end{proof}
\begin{corollary}[支配的incidenceと有限次数の良い開部分]
$Q_{113}$の次数三中心直線のincidenceには一般ファイバー$\PP^4_K$を持つ
41次元の支配的既約成分$I_{113}$が一意に存在する。
中心の次数1,3とHilbert関数の次数4以下が一般値を取る開部分では、
これはrank5の中心核束の$\PP^4$束である。
一般の有理商の次元は周囲で13、対で17、忘却の一般ファイバーは4次元である。
\end{corollary}
\begin{proof}
中心の有限行列の非零minorを反転するとrank5の核束を得る。
その射影束は既約で一般ファイバー$\PP^4_K$の閉包を与え、支配的成分は他にない。
相対Jacobianパラメータ族の6+35=41次元の像がこれを支配する。
六分岐値を持つ対角族の周囲は射影stabilizerが有限であることを
第\ref{sec:pair-diagonal-classification}節で証明する。
従って$\PGL_5$の一般軌道は24次元である。
Rosenlichtの定理により37−24=13、41−24=17を得て、忘却の一般ファイバーは4次元となる。
この計算は商多様体の有理性や一般stabilizerの自明性を含まない。
\end{proof}

\section{Ambientへの射・相対incidence・一次障害}
\label{sec:ambient-incidence}
二次包絡の一意性は、三次対からambientへの\emph{射}を与える。
その射の単射性は意味しない。本節では、有限型のambient族を一つ固定し、
この射を中心核の線形代数として構成する。
全AS正則$\mathbb Z$-代数の大域的な代数的モジュライの存在は入力しない。

\subsection{中心三次直線の相対表示}
有限型$k$-スキーム$S$上に、次数一生成の連結次数付き代数族
$\mathcal B=\bigoplus_{d\geq0}\mathcal B_d$を取り、各$\mathcal B_d$は有限階数局所自由とする。
考える例では$\mathcal B_1$の階数は5、ファイバーは二次AS正則代数である。
中心性を検出するベクトル束写像を
\begin{equation}\label{eq:central-commutator-map}
 \delta:\mathcal B_3\longrightarrow
 \mathcal H:=\mathcal Hom_S(\mathcal B_1,\mathcal B_4),
 \qquad \delta(c)(x)=cx-xc
\end{equation}
と置く。次数一生成性から、$\ker\delta_s=Z(B_s)_3$である。
$\PP_S(\mathcal B_3)$は一次元部分空間を表す規約とし、
その普遍直線束を$\mathcal L\subset p^*\mathcal B_3$と書く。

\begin{proposition}[中心incidenceと基底変換]\label{prop:central-incidence-basechange}
合成$\mathcal L\to p^*\mathcal B_3\xrightarrow{p^*\delta}p^*\mathcal H$
の零スキーム
\begin{equation}\label{eq:central-incidence-definition}
 \mathcal I^{\mathrm{cen}}_3(\mathcal B/S)\subset\PP_S(\mathcal B_3)
\end{equation}
は、中心三次直線を表す有限型の射影$S$-スキームであり、任意の基底変換と可換する。
写像$\delta$の階数がスキーム論的に一定の層$S_r$上では
$\mathcal E_r=\ker(\delta|_{S_r})$は局所自由で、
\begin{equation}\label{eq:central-projective-kernel}
 \mathcal I^{\mathrm{cen}}_3(\mathcal B/S)\times_S S_r
 \simeq\PP_{S_r}(\mathcal E_r).
\end{equation}
この核束および同型も$S_r$上の任意の基底変換と可換する。
\end{proposition}
\begin{proof}
中心三次直線とは、直線部分束$L\subset B_3$で$\delta(L)=0$を満たすものであり、
普遍零スキームの関手的定義がそのまま主張を与える。
階数$r$の層は、$(r+1)$次小行列式を零とし、少なくとも一つの$r$次小行列式を
可逆にした局所閉スキームとして取る。
その上で局所的な行列の基本変形により$\delta$は
$\operatorname{diag}(I_r,0)$となる。従って核と余核は局所自由で、
核を取る操作が任意の基底変換と可換する。普遍直線の記述から
\eqref{eq:central-projective-kernel}を得る。
単に幾何点で核次元が一致することを、非被約な底上の核の基底変換と同一視していない。
\end{proof}

全幾何ファイバー$B_s$が整域なら、incidence上の非零三次直線は各ファイバーで正則である。
このとき全次数$d$について
\begin{equation}\label{eq:relative-cubic-flat-quotient}
 0\longrightarrow\mathcal L\otimes p^*\mathcal B_{d-3}
 \xrightarrow{\ \mathrm{mult}\ }p^*\mathcal B_d
 \longrightarrow\mathcal A_d\longrightarrow0
\end{equation}
は局所自由層の完全列になる（$d<3$では左項を零とする）。
実際、最初の写像は全幾何ファイバーで単射なので、局所的に階数分の小行列式が可逆となり、
局所自由な余核を持つ。これで商族の全次数の平坦性が従う。
整域条件を外す場合は、正則性を別条件として課す。
いずれの場合も、$\qgr A_s$の大域次元3または導来圏の滑らかさは別の条件である。
その条件の開性・非空性は、三次中心核の有限行列だけからは従わない。

\subsection{群作用・ファイバー・次元の規約}
$V=k^5$を固定し、$S$が二次関係空間$R\subset V\otimes V$の
$\GL(V)$不変な族であるとする。
上の構成は変数変換と可換するので、
\begin{equation}\label{eq:central-stack-morphism}
 [\mathcal I^{\mathrm{cen}}_3/\GL(V)]
 \longrightarrow[S/\GL(V)]
\end{equation}
という射を得る。これは選んだ有限型ambient族の上の代数的stackの射である。
中心スカラーは関係部分空間と三次\emph{直線}の双方に自明に作用するため、
同じ作用は$\PGL(V)$を通り、対応する$\PGL(V)$商も取れる。
このスカラー自己同型を除いた規約と、$\GL(V)$商stackそのものとは区別する。

\begin{proposition}[固定ambient上のファイバー]\label{prop:ambient-fibers-stack}
幾何点$B$と、その点へのambientの同型を指定して
\eqref{eq:central-stack-morphism}の2ファイバーを取ると、
それは$\PP(Z(B)_3)$である。
正則元だけ、または滑らかな三次圏だけに制限した場合は、対応する部分関手を取る。
ambientの\emph{同型類}だけを固定して同型の指定を忘れると、
三次対の同型類は$\PP(Z(B)_3)$に対する$\Aut_{\mathrm{gr}}(B)$の軌道である。
\end{proposition}
\begin{proof}
商stackの射の2ファイバーの対象は、中心三次対と、そのambientを指定された$B$へ送る
同型との組である。この同型を用いると中心直線は$B_3$内の直線になり、
同型を保つ自己同型は恒等である。従ってファイバーは中心直線のスキームそのものである。
同型の指定を忘れれば、残る同一視は$B$の次数付き自己同型による作用である。
\end{proof}

定理\ref{thm:pair-generic-center}の一般点では$\dim Z(B)_3=5$であり、
このファイバーは$\PP^4$となる。
中心の階数5が一定の開部分では、射は相対次元4の射影束である。
一般の射影stabilizerが有限であることを合わせると、次元は次のようになる。
\begin{center}
\begin{tabular}{lrr}
\toprule
規約&ambient&中心三次対\\
\midrule
生成元を標識した係数空間&37&41\\
$\PGL_5$商stack・一般有理商&13&17\\
$\GL_5$商stack&12&16\\
\bottomrule
\end{tabular}
\end{center}
$\GL_5$商stackの次元が一つ小さいのは、一般点にも残る
$\mathbb G_m$自己同型のためであり、同型類の変動次元が減るという意味ではない。
一般有理商でのファイバーは、有限な射影stabilizerを$\Gamma_B$として
$\PP^4/\Gamma_B$と記述され、次元は4のままである。
一般stabilizerが自明だと証明しない限り、これを$\PP^4$そのものとは呼ばない。

命題\ref{prop:quantum-113-component}の規格化CY変形の対応成分では、
中心三次incidenceが成分の稠密開部分を支配する。
従って\emph{この成分の中で}その像の一般余次元は0である。
これは像の閉包についての結論であり、全ての境界点を像に含むという結論ではない。
また、滑らかな三次圏の部分関手へ制限した射が同じ開部分を支配することには、
その部分関手の一般ファイバーでの非空性を別途要する。

\subsection{中心三次を延長する一次障害}
$s\in S(k)$と$0\ne c\in Z(B_s)_3$を固定する。
一次変形$v\in T_sS$上で$\mathcal B_3,\mathcal H$を自明化し、
$\delta$の一次変化を$\dot\delta_v$とする。
三次直線の持上げを$c+\epsilon\dot c$と書けば、中心性は
\begin{equation}\label{eq:central-tangent-equation}
 \delta_s(\dot c)+\dot\delta_v(c)=0,
 \qquad \dot c\in(B_s)_3/kc
\end{equation}
に等しい。従って一次障害は
\begin{equation}\label{eq:central-tangent-obstruction}
 \omega_c(v):=[\dot\delta_v(c)]\in\operatorname{coker}\delta_s
\end{equation}
である。自明化を変えると右辺は$\operatorname{im}\delta_s$の元だけ変わり、
$\delta_s(c)=0$なので余核中の類は変わらない。

\begin{proposition}[incidenceの接空間完全列]\label{prop:central-tangent-exact}
$L=kc$と置くと、自然な完全列
\begin{equation}\label{eq:central-tangent-exact-sequence}
\begin{split}
 0\longrightarrow\Hom_k(L,Z(B_s)_3/L)
 \longrightarrow T_{(s,[c])}\mathcal I^{\mathrm{cen}}_3
 \longrightarrow T_sS
 \xrightarrow{\ \omega\ }\Hom_k(L,\operatorname{coker}\delta_s)
\end{split}
\end{equation}
がある。最後の射の全射性は主張しない。
特にambientの一次変形が中心三次対へ延長できるための必要十分条件は
$\omega_c(v)=0$である。
\end{proposition}
\begin{proof}
\eqref{eq:central-tangent-equation}の解の存在は
$\dot\delta_v(c)\in\operatorname{im}\delta_s$と同値である。
$v=0$の場合の解は$Z(B_s)_3/kc$であり、生成元$c$の選択を除くと
$\Hom_k(L,Z(B_s)_3/L)$となる。これが全ての箇所の完全性を与える。
\end{proof}

階数一定の層に沿っては核束が存在するため、この一次障害は消える。
一方、その層に横断するambient変形には中心三次を失う方向があり得る。
全ambientモジュライ内での像の余次元を決める際は、
中心核の階数だけでなく、この障害写像と高次の延長問題を調べる必要がある。
これは全非可換変形の障害複体を中心性の有限行列で置き換えるという主張ではない。

\subsection{中心条件と正規条件の差}
正則な正規三次元$c$は一般に中心元ではなく、
一意な次数付き自己同型$\sigma$に対して$cx=\sigma(x)c$を満たす。
二次表示$B=T(V)/(R)$が与えられている場合、標識付き正規incidenceは
\begin{equation}\label{eq:normal-incidence-equations}
 \sigma\in\GL(V),\quad
 (\sigma\otimes\sigma)(R)=R,\quad
 cx-\sigma(x)c=0\ \text{in }B_4\quad(x\in V)
\end{equation}
で書ける。二次関係を保存する条件により$\sigma$は$B$の自己同型へ延びる。
最後の等式は積について帰納的に延び、$cB=Bc$を与える。
正則性があれば、この式から$\sigma$は一意に決まる。
中心incidenceは$\sigma=1$という部分問題であり、
中心核の射影束をそのまま全正規incidenceだとみなすことはできない。
また、ここでの表示は1周期化を選んだ次数付き代数についてのものであり、
任意の$\mathbb Z$-代数の正規族を一つの$\sigma$に置き換えたものではない。

\begin{proposition}[量子多項式族での正の余次元の例]
\label{prop:quantum-polynomial-central-cubic}
$q_{ij}\in\mathbb G_m$ ($1\le i<j\le5$)、$q_{ji}=q_{ij}^{-1}$、$q_{ii}=1$として
\[
 B(q)=k\langle x_1,\ldots,x_5\rangle/
 (x_ix_j-q_{ij}x_jx_i)
\]
を考える。この10次元パラメータトーラスのうち、非零の中心三次元を持つ部分は
余次元4の部分トーラスの有限個の捩れ点平行移動の和集合である。
従って中心三次incidenceの像はこのパラメータ族では真部分集合である。
一方、任意の$q$で$x_i^3$は正則正規三次元である。
\end{proposition}
\begin{proof}
順序付き単項式$x^\alpha$はPBW基底をなし、$B(q)$は反復skew多項式拡大による整域である。
そのKoszul複体は可換多項式環の複体の多重次数による捩れとして得られ、
Koszul双対は五生成元のskew外積代数である。従ってこれは標準的な
noetherianな5次元Koszul AS正則代数の族である。
各単項式について
\[
 x_i x^\alpha=\chi_i(\alpha)x^\alpha x_i,
 \qquad\chi_i(\alpha)=\prod_jq_{ij}^{\alpha_j}
\]
となる。交換子の単項式が独立なので、次数三中心が非零であるための必要十分条件は、
$|\alpha|=3$を満たす35個の指数の少なくとも一つについて
$\chi_i(\alpha)=1$ ($i=1,\ldots,5$)となることである。
このトーラス準同型の指標格子上の写像を$\mathbb Q$へ拡大すると
$\beta\mapsto\beta\wedge\alpha$であり、$\alpha\ne0$なので階数は4である。
従って各核の連結成分は次元$10-4=6$の部分トーラスの捩れ点平行移動である。
有限和集合も次元6であり、主張を得る。最後の主張は各$x_i$の正規性と整域性による。
\end{proof}

ここでの余次元4は\emph{指定した量子多項式の係数トーラス内}の値である。
座標変換・Zhang捩れ・$\mathbb Z$-代数同型をすべて商に取った空間の
余次元だとは主張しない。また、単項式による正規商の存在は、
その商の射影圏が滑らかな三次三様体であることを意味しない。
この例は、中心三次を持つ像の大きさがambient族によって変わり、
相対Jacobian成分の余次元0という結論を全ambient族へ一般化できないことを示す。

\section{二次包絡によるモジュライ同型類の判定}\label{sec:moduli-iso}
\begin{theorem}[三次商からのambient復元]\label{thm:moduli-iso}
$B=T(V)/(R_2)$を次数一生成の二次連結次数付き代数とし、
$G\in Z(B)_3$を非零正則中心元とする。$A=B/(G)$に対して
\begin{equation}\label{eq:moduli-envelope}
 B\simeq T(A_1)/
 \left\langle\ker(A_1\otimes A_1\longrightarrow A_2)\right\rangle,
 \qquad \ker(B_3\longrightarrow A_3)=kG
\end{equation}
となる。特に第\ref{sec:normal}節の17係数一般族では、
中心三次対の同型類の次元17が、実際の$A$の
次数付き同型類の次元でもある。
さらに$A$が左右noetherianなAS Gorenstein代数で
injective dimensionが2以上なら、
$(\qgr A,\OO_A,(1))$から同じ対象を復元できる。
\end{theorem}
\begin{proof}
$G$は次数三なので$A_1=B_1$, $A_2=B_2$である。
したがって$A$の二次関係の核は$R_2$であり、
二次関係のみで定義された$B$を復元する。
商写像の次数三の核は$B_0G=kG$である。
次数付き同型$A\simeq A'$はこの復元により
$(B,kG)\simeq(B',kG')$を誘導する。
偏極付き圏については、AS Gorenstein性より
$\operatorname{Ext}_A^0(k,A)=\operatorname{Ext}_A^1(k,A)=0$であり、
Artin--Zhangのsection algebraの自然写像は
\begin{equation}\label{eq:moduli-section}
 A\xrightarrow{\ \sim\ }\bigoplus_{n\ge0}
 \operatorname{Hom}_{\qgr A}(\OO_A,\OO_A(n))
\end{equation}
を与える。よって偏極付き圏の同値も次数付き代数の
同型を誘導する。
\end{proof}


\clearpage
\part{代数的部分モジュライとKuznetsov残余射}
\section{相対 Koszul 解消による三次 tails の双対}
\label{sec:relative-koszul-line-cohomology}

この節では、体上の balanced dualizing complex の存在を、
一般の係数環上へ形式的に拡張することはしない。
有限な相対 Koszul 双加群解消から torsion 関手を直接計算する。
全ての基底で大きい $\operatorname{QGr}$ と非有界導来圏を使う。

\subsection{明示的な相対入力}

$P$ を可換環とし、$B=\bigoplus_{n\geq0}B_n$ を $B_0=P$ の次数付き $P$-代数とする。
$P$ は中心に作用し、$B$ は次数一で生成されるとする。
以下を入力とする。
\begin{enumerate}
\item 各 $B_n$ は有限射影 $P$-加群である。
\item 有限な次数付き双加群解消
\begin{equation}\label{eq:relative-diagonal-koszul-input}
 K_B^{-r}=B\otimes_P W_r\otimes_P B\quad(0\leq r\leq5),
 \qquad K_B\xrightarrow{\sim}B
\end{equation}
があり、$W_r$ は内部次数 $r$ に置かれた有限射影 $P$-加群、
$W_0=P$、$W_5$ は可逆 $P$-加群である。
これは相対二次 Koszul 解消であり、augmentation を片側へ tensor した複体は
$P$ の Koszul 解消となる。
この解消は全ての可換 $P$-代数への基底変更後にも完全である。
\item 相対 CY の規格化として、真正な次数付き $B$-双加群の導来同型
\begin{equation}\label{eq:relative-cy-five-orientation}
 \Theta_{B/P}:=
 \mathbf R\underline{\operatorname{Hom}}_{B\otimes_PB^{\mathrm{op}}}
      (B,B\otimes_PB^{\mathrm{op}})
 \simeq B(5)[-5]
\end{equation}
を指定する。
この同型を有限解消 \eqref{eq:relative-diagonal-koszul-input} の双対により表し、
任意の基底変更に両立するものとする。
\item $w\in Z(B)_3$ を取り、$A=B/(w)$ と置く。
各次数で
\begin{equation}\label{eq:relative-cubic-split-sequence}
 0\longrightarrow B_{n-3}\xrightarrow{\,w\,}B_n
 \longrightarrow A_n\longrightarrow0
\end{equation}
は完全で、$A_n$ は有限射影 $P$-加群である。
従ってこの列は $P$-加群として分裂し、$w$ の正則性は任意の基底変更後にも保たれる。
\end{enumerate}
特に $B$ と $A$ は $P$ 上平坦である。
非正則な $P$ 上の $B$ の絶対的大域次元が五であるとは仮定しない。
また有限中心、左右 noetherian 性、楕円境界の有限位数性を仮定しない。
実際の十七係数族には、第\ref{sec:general17-relative-algebra}節で、
これらの入力を一つの明示した基底上で供給する。
各体値の AS/CY 性だけから入力 (2)--(3) を推測しない。

\subsection{有限双加群解消による torsion 関手}

右次数付き加群を用い、
$M_{<n}=B/B_{\geq n}$ と書く。
$\Gamma_B$ は各元が $B_{\geq n}$ のある $n$ で消える torsion 部分加群を取る関手である。

\begin{lemma}[非 noetherian 基底上の有界 torsion 模型]
\label{lem:relative-koszul-torsion-model}
入力 (1)--(2) の下で、任意の可換 $P$-代数 $S$ 上において以下が成立する。
\begin{enumerate}
\item $\mathbf R\Gamma_{B_S}$ は振幅 $[0,5]$ を持ち、filtered colimit と直和を保つ。
\item torsion 加群 $T$ に対し $\mathbf R\Gamma_{B_S}(T)\simeq T$ である。
この同型は torsion cohomology を持つ非有界複体にも成り立つ。
\item 商写像 $\pi_B:\operatorname{Gr}(B_S)\to\operatorname{QGr}(B_S)$ と
右随伴 $\omega_B$ について、非有界導来圏上に三角
\begin{equation}\label{eq:relative-torsion-localization-triangle}
 \mathbf R\Gamma_{B_S}(N)\longrightarrow N
 \longrightarrow\mathbf R\omega_B\pi_BN\longrightarrow
\end{equation}
があり、$\mathbf R\omega_B$ は直和を保つ。
\item 各 $\mathcal O_B(i)=\pi_BB_S(i)$ は compact であり、
$\bigoplus_{i=0}^4\mathcal O_B(i)$ は
$D(\operatorname{QGr}(B_S))$ の compact generator である。
\end{enumerate}
\end{lemma}

\begin{proof}
全ての入力が基底変更後にも成立するので、$S=P$ と書いて証明する。
まず tails torsion は localizing である。
実際、次数一生成と $B_n$ の有限生成性から、右イデアル $B_{\geq n}=B_nB$ は有限生成である。
拡大の商である元が $B_{\geq n}$ で消えるなら、その有限個の生成元の像を消す
共通の冪を選ぶことで、元自身も十分高い次数のイデアルで消える。
従って torsion は拡大で閉じ、部分対象、商、直和でも閉じる。
この議論は $A$ に対しても同じである。

$B$ は $P$-射影加群であるから、$K_B$ の各項は左 $B$-射影であり、
$K_B\to B$ は左 $B$-加群の有界な分裂解消である。
従って
\[
 K_n=M_{<n}\otimes_B K_B\longrightarrow M_{<n}
\]
は任意の $n$ について右 $B$-加群の解消である。
その項は $M_{<n}\otimes_PW_r\otimes_PB$ である。
$M_{<n}$ は有限射影 $P$-加群なので、これは有限生成射影右 $B$-加群である。
この解消は $M_{<n+1}\to M_{<n}$ に自然である。

従って
\begin{equation}\label{eq:relative-bounded-torsion-complex}
 C_B(N)=\varinjlim_n
 \underline{\operatorname{Hom}}_B(K_n,N)
\end{equation}
は $\mathbf R\Gamma_B(N)$ を計算する。
具体的には、各 Hom は有界な有限射影解消による導来 Hom であり、
filtered colimit は exact なので、cohomology は
$\varinjlim_n\underline{\operatorname{Ext}}^q_B(M_{<n},N)$ となる。
各項の関手は exact で、直和と filtered colimit を保つ。
横方向の長さが五であるため、この複体は非有界複体 $N$ に対しても
通常の有限幅の全複体によって定義できる。
$K$-injective 解消を取った場合にも、各 injective 項における同じ計算と
有限幅の全複体の比較により、真の非有界導来関手を計算する。

torsion の非 noetherian な問題をここで確認する。
内部次数が上に有界な加群 $T$ を取る。
固定した内部次数 $t$ について、
$\underline{\operatorname{Hom}}_B(K_n,T)_t$ は $n\gg0$ で
$\underline{\operatorname{Hom}}_B(K_B,T)_t$ に安定する。
実際、$K_B^{-r}=\bigoplus_{d\geq0}B_d\otimes_PW_r\otimes_PB$ の
生成元の内部次数は $d+r$ であり、$d+r+t$ が $T$ の上限を超える項の Hom は零である。
従って各固定次数で有限個の $d$ しか寄与しない。
$K_B\to B$ は右 $B$-射影解消でもあるので、極限の cohomology は $T$ の次数零だけである。

任意の torsion 加群は、その有限個の斉次元で生成される部分加群の filtered colimit である。
各部分加群は内部次数が上に有界である。
前段と \eqref{eq:relative-bounded-torsion-complex} の filtered colimit 保存から、
$\mathbf R\Gamma_B(T)\simeq T$ を得る。
有界な cohomological amplitude による標準の truncation 比較により、
同じ結論が torsion cohomology を持つ非有界複体に成立する。
例えば $R^p\Gamma_B(H^qN)\Rightarrow H^{p+q}(\mathbf R\Gamma_BN)$ の
spectral sequence は $0\leq p\leq5$ の有限幅なので、非有界な $N$ にも収束し、
torsion cohomology の場合は $p=0$ の行だけになる。
さらに \eqref{eq:relative-bounded-torsion-complex} の各項は torsion 加群である。
右 $B$-作用は $M_{<n}$ の左作用から来て、$B_{\geq n}$ がその項を消すからである。

これにより $D(\operatorname{Tors}B)\to D(\operatorname{Gr}B)$ は完全忠実であり、
その像は cohomologically torsion な複体である。
一方、$\pi_B$ は exact であり、$\omega_B$ は $K$-injective 複体を
$K$-injective 複体へ送る。$\pi_B\omega_B=\operatorname{id}$ なので
\[
 D(\operatorname{Gr}B)/D_{\operatorname{tors}}(\operatorname{Gr}B)
 \simeq D(\operatorname{QGr}B)
\]
となる。これらの随伴から \eqref{eq:relative-torsion-localization-triangle} を得る。
$C_B$ が直和を保つので $\mathbf R\omega_B\pi_B$ も直和を保ち、
$\pi_B$ の本質的全射性から $\mathbf R\omega_B$ 自身も直和を保つ。

従って自由対象 $B(i)$ の compact 性と随伴により、
全ての $\mathcal O_B(i)$ は compact である。
これらが全ての cohomological shift とともに対象 $N$ に直交するなら、
$\mathbf R\omega_BN$ は全ての自由 $B(i)$ に直交するため零であり、$N=0$ である。
よって全ての整数 twist は compact generators をなす。

augmentation の Koszul 解消を $\pi_B$ へ送ると
\[
 0\to\mathcal O_B(-5)\otimes_PW_5\to\cdots\to
 \mathcal O_B(-1)\otimes_PW_1\to\mathcal O_B\to0
\]
が完全になる。
$W_5$ が可逆で、全ての $W_r$ が有限射影なので、
任意の連続する五つの twist の thick 閉包は前後の twist を含む。
従って $\bigoplus_{i=0}^4\mathcal O_B(i)$ が compact generator である。
\end{proof}

\begin{remark}[smashing の具体的な意味]
上の模型は dg 双加群のレベルでも使える。
$K_n$ は有限生成射影右 $B$-加群の複体なので、項別の同型
\[
 \underline{\operatorname{Hom}}_B(K_n,N)
 \simeq N\otimes_B
       \underline{\operatorname{Hom}}_B(K_n,B)
\]
を持つ。
右辺の第二因子は左 $B$-射影の有界複体であり、その filtered colimit は
左 $B$-flat の有界複体である。
従って
\[
 \mathbf R\Gamma_B(N)
 \simeq N\otimes_B^{\mathbf L}\mathbf R\Gamma_B(B)
\]
となり、localization の三角も $N$ と tensor して得られる。
これは単なる cohomology の階数計算ではない。
\end{remark}

\subsection{相対 CY 入力からの局所コホモロジー}

有限個の次数にのみ台を持つ有限射影 $P$-加群 $M$ について
$M^\vee=\underline{\operatorname{Hom}}_P(M,P)$ と書く。
$B^\vee=\bigoplus_n\operatorname{Hom}_P(B_{-n},P)$ は次数ごとの双対である。

\begin{lemma}[有限係数加群の双対性]
\label{lem:relative-cy-finite-coefficient-dual}
入力 (1)--(3) の下で、有限射影 $P$-加群である右次数付き $B$-加群 $M$ に対して、
自然な同型
\begin{equation}\label{eq:relative-finite-module-dual}
 \mathbf R\underline{\operatorname{Hom}}_B(M,B)
 \simeq \Theta_{B/P}\otimes_B^{\mathbf L}M^\vee
 \simeq M^\vee(5)[-5]
\end{equation}
がある。
$M$ がさらに左の代数作用を持つとき、その作用も双対して保持する。
従って $M=M_{<n}$ では真正な $B$-双加群の同型である。
全ての同型は任意の可換基底変更に両立する。
\end{lemma}

\begin{proof}
$M\otimes_BK_B$ は右 $B$-有限射影解消である。
項 $B\otimes_PW_r\otimes_PB$ に対し、有限射影性による tensor--Hom の同型は
\[
 \underline{\operatorname{Hom}}_B
       (M\otimes_PW_r\otimes_PB,B)
 \simeq B\otimes_PW_r^\vee\otimes_PM^\vee
\]
となる。
これは双加群解消の双対を $M^\vee$ と tensor した項であり、
evaluation の自然性により微分および全ての残された作用と可換する。
従って最初の導来同型を得る。
第二の同型には指定した \eqref{eq:relative-cy-five-orientation} を使う。
有限射影加群の双対、有限 tensor、解消の全てが基底変更と可換するため、最後の主張も従う。
体上の同じ計算は \cite[Proposition~4.12]{RRGeneralizedAS} に現れる。
ここでは係数有限次元性を有限射影性に置き換え、必要な項別計算を直接示した。
\end{proof}

\begin{proposition}[周囲代数の相対局所コホモロジー]
\label{prop:relative-as5-local-cohomology}
全ての可換 $P$-代数 $S$ について、真正な次数付き $B_S$-双加群として
\begin{equation}\label{eq:relative-as5-local-dual}
 \mathbf R\Gamma_{B_S}(B_S)\simeq B_S^\vee(5)[-5]
\end{equation}
であり、その形成は任意の基底変更と可換する。
従って全ての整数 $n$ に対して
\[
 \mathbf R\operatorname{Hom}_{\operatorname{QGr}(B_S)}
       (\mathcal O_B,\mathcal O_B(n))
 \simeq
 \begin{cases}
  (B_S)_n[0],&n\geq0,\\
  0,&-4\leq n\leq-1,\\
  (B_S)_{-n-5}^\vee[-4],&n\leq-5.
 \end{cases}
\]
右辺は $S$-perfect である。
\end{proposition}

\begin{proof}
前補題を $M_{<n}=B_S/B_{S,\geq n}$ に適用して filtered colimit を取る。
$\varinjlim_nM_{<n}^\vee=B_S^\vee$ なので、
補題\ref{lem:relative-koszul-torsion-model}の模型から
\eqref{eq:relative-as5-local-dual} を得る。
cohomology を計算してから双加群作用を推測するのではなく、
各有限段階の真正な双加群同型を用いている。

localization の三角の内部次数 $n$ を取る。
その中央項 $(B_S)_n$ は $n\geq0$ のみに存在し、
第一項 $B_S^\vee(5)_n[-5]$ は $n\leq-5$ のみに存在する。
従って第三項は表示した通りである。
三角は canonical section map とその高次版を計算しており、
標準の乗法写像に自然である。
任意の内部次数で各係数は有限射影なので、基底変更と perfect 性も従う。
\end{proof}

\subsection{中心三次商と全 twist の計算}

\begin{lemma}[中心商の非有界 tails 随伴]
\label{lem:relative-central-quotient-derived-adjunction}
入力 (1)--(4) の下で、任意の可換 $P$-代数 $S$ に対し
\[
 \mathbf L j^*:D(\operatorname{QGr}(B_S))
 \rightleftarrows D(\operatorname{QGr}(A_S)):j_*
\]
という随伴がある。
$\mathbf L j^*$ は $-\otimes_{B_S}^{\mathbf L}A_S$ から誘導され、
$j_*$ は忘却から誘導される exact、conservative、直和を保つ関手である。
$\mathbf L j^*\mathcal O_B(i)=\mathcal O_A(i)$ であり、
\[
 j_*\mathcal O_A(n)
 \simeq\operatorname{Cone}\bigl(
  \mathcal O_B(n-3)\xrightarrow{\,w\,}\mathcal O_B(n)\bigr).
\]
\end{lemma}

\begin{proof}
次数付き加群圏の忘却 $i_*$ は exact であり、
$A_S$-torsion と忘却後の $B_S$-torsion は同じである。
従って商圏へ exact な $j_*$ を誘導し、その kernel は零である。
忘却は直和を保つので $j_*$ も直和を保つ。

中心三次元の二項解消から、
$-\otimes_{B_S}^{\mathbf L}A_S$ の忘却後の複体は
cone$(w:N(-3)\to N)$ である。
従って cohomologically torsion な複体を同種の複体へ送る。
非有界導来圏の Gabriel 商の記述から、左随伴を tails の導来圏へ降ろせる。
随伴そのものも降下する。
実際、$N$ が $A_S$ 側の右 local 複体なら、torsion 複体 $T$ について
\[
 \operatorname{Hom}_{D(\operatorname{Gr}B_S)}(T,i_*N)
 =\operatorname{Hom}_{D(\operatorname{Gr}A_S)}
      (T\otimes_{B_S}^{\mathbf L}A_S,N)=0
\]
なので $i_*N$ は $B_S$ 側でも右 local である。
最後の二つの式は自由加群と中心三次元の二項解消に適用したものに他ならない。
\end{proof}

\begin{theorem}[有限中心なしの普遍的線コホモロジー]
\label{thm:relative-cubic-universal-line-cohomology}
入力 (1)--(4) の下で、任意の可換 $P$-代数 $S$ と整数 $i,j$ に対し、
$n=j-i$ と置くと
\begin{equation}\label{eq:relative-cubic-all-line-rhom}
 \mathbf R\operatorname{Hom}_{\operatorname{QGr}(A_S)}
       (\mathcal O_A(i),\mathcal O_A(j))
 \simeq
 \begin{cases}
  (A_S)_n[0],& n\geq0,\\
  0,& n=-1,\\
  (A_S)_{-n-2}^\vee[-3],& n\leq-2.
 \end{cases}
\end{equation}
この同型は指定した CY 規格化の下で任意の基底変更に両立する。
特に、全ての線対象間の RHom は $S$-perfect であり、
標準 section 写像は全ての次数で同型である。
\end{theorem}

\begin{proof}
twist を移して $i=0$ としてよい。
前補題の随伴により左辺は
\[
 \mathbf R\operatorname{Hom}_{\operatorname{QGr}(B_S)}
       (\mathcal O_B,j_*\mathcal O_A(n))
\]
であり、これは周囲の線コホモロジーの写像
\[
 Q_B(n-3)\xrightarrow{\,w\,}Q_B(n),\qquad
 Q_B(t)=\mathbf R\operatorname{Hom}(\mathcal O_B,\mathcal O_B(t))
\]
の cone である。
$n\geq0$ の場合、関連する負の次数は $-1,-2,-3$ だけであり、
周囲の高次 cohomology は寄与しない。
従って \eqref{eq:relative-cubic-split-sequence} によって $(A_S)_n$ を得る。
$n=-1$ では周囲の二項はともに零である。

$n\leq-2$ では、cohomological degree 四の写像は
\[
 (B_S)_{-n-2}^\vee\longrightarrow (B_S)_{-n-5}^\vee
\]
である。
命題\ref{prop:relative-as5-local-cohomology}の双加群としての自然性により、
これは $w:(B_S)_{-n-5}\to(B_S)_{-n-2}$ の双対である。
分裂完全列 \eqref{eq:relative-cubic-split-sequence} の双対から、この写像は全射で、
kernel は $(A_S)_{-n-2}^\vee$ である。
従って cone は cohomological degree 三にのみ cohomology を持ち、表示式を得る。

全ての計算は有限射影係数の写像、指定した CY 同型、および canonical な随伴により構成される。
従って非平坦なものも含む任意の基底変更に両立する。
正次数の同型は標準 section map そのものである。
\end{proof}

\begin{corollary}[compact generation と相対 properness]
\label{cor:relative-cubic-compact-proper}
任意の可換 $P$-代数 $S$ に対し
\[
 T_{A,S}=\bigoplus_{i=0}^4\mathcal O_A(i)
\]
は $D(\operatorname{QGr}(A_S))$ の compact generator であり、
\[
 D(\operatorname{QGr}(A_S))^c=\operatorname{thick}(T_{A,S}).
\]
任意の二つの compact 対象の RHom は $S$-perfect である。
さらに $P$ 上の compact 対象 $F,G$ と任意の $P\to S$ に対して、
\[
 \mathbf R\operatorname{Hom}(F,G)\otimes_P^{\mathbf L}S
 \xrightarrow{\sim}
 \mathbf R\operatorname{Hom}(F_S,G_S)
\]
となる。ここで $F_S,G_S$ は導来係数延長である。
この主張は dg smoothness を含まない。
\end{corollary}

\begin{proof}
$T_{B,S}$ は補題\ref{lem:relative-koszul-torsion-model}により compact generator である。
$j_*$ は直和を保つので、その左随伴の像
$\mathbf Lj^*T_{B,S}=T_{A,S}$ は compact である。
$T_{A,S}$ とその全ての shift に直交する対象 $N$ については、
随伴により $j_*N$ が $T_{B,S}$ とその全ての shift に直交する。
従って $j_*N=0$ であり、保守性から $N=0$ である。
compact 対象がこの generator の thick 閉包であることは、
compactly generated category の標準定理
\cite[Theorem~3.4]{KellerDgCategories}による。

定理\ref{thm:relative-cubic-universal-line-cohomology}により
$\mathbf R\operatorname{Hom}(T_{A,S},T_{A,S})$ は $S$-perfect である。
第一・第二変数について有限 cone、shift、直和因子を取ると、
全 compact 対象間の perfect 性が従う。
同様に、係数変更比較は線対象の対で同型であり、
両変数について thick に閉じているので全 compact 対象へ延びる。
導来係数延長を tails へ降ろせることは、次の torsion 比較の補題からも直接分かる。
torsion 複体は torsion 加群から localizing に生成され、
係数環の flat 解消との tensor は torsion 加群を torsion 複体へ送るからである。
\end{proof}

\subsection{三次商自身の torsion と真正な双加群双対}

次の補足は、上の線コホモロジーの計算には不要だが、
相対 Serre trace の構成へ渡す局所双対入力を与える。

\begin{lemma}[中心商に対する derived torsion の比較]
\label{lem:relative-cubic-torsion-comparison}
上の入力の下で、$A$-加群 $N$ の torsion cohomology は、忘却後の $B$ 上で計算できる。
すなわち
\[
 i_*\mathbf R\Gamma_A(N)\simeq\mathbf R\Gamma_B(i_*N).
\]
これは非有界複体にも成り立つ。
$\mathbf R\Gamma_A$ は有限 cohomological dimension を持ち、直和を保ち、
torsion cohomology を持つ複体上では恒等である。
従って通常の derived torsion は Gabriel localization の kernel への右随伴と一致する。
\end{lemma}

\begin{proof}
torsion 部分加群そのものは $A$ と $B$ で同じである。
$J$ を injective な次数付き $A$-加群とする。
有限射影 $B$-解消と induction--restriction の随伴から
\[
 \underline{\operatorname{Ext}}^q_B(M_{<n},i_*J)
 \simeq
 \underline{\operatorname{Hom}}_A
       (\operatorname{Tor}^B_q(M_{<n},A),J)
\]
となる。
二項解消により $q>1$ の Tor は零である。
$q=1$ では
\[
 \operatorname{Tor}^B_1(M_{<n},A)
 =\ker\bigl(w:M_{<n}(-3)\to M_{<n}\bigr).
\]
$w$ の $B$ 上の正則性により、この kernel を表す元の元の次数は
$n-3,n-2,n-1$ に限られる。
従って $n'\geq n+3$ に対する遷移
$\operatorname{Tor}^B_1(M_{<n'},A)\to\operatorname{Tor}^B_1(M_{<n},A)$ は零である。
filtered colimit を取ると $i_*J$ の高次 $\Gamma_B$ は全て零となる。

従って $A$-injective 解消により derived torsion を比較できる。
非有界複体の場合にも、補題\ref{lem:relative-koszul-torsion-model}の有限幅の模型を
$A$-injective 解消の各項へ適用すれば同じ比較を得る。
有限 cohomological dimension と直和保存は $B$ 側から従う。
torsion 加群上の acyclicity、さらに非有界 torsion 複体上の恒等性も $B$ 側から従う。
これらにより $D(\operatorname{Tors}A)\simeq D_{\operatorname{tors}}(\operatorname{Gr}A)$ となる。
torsion 加群は、内部次数の有限な filtration と係数加群の自由解消から生成され、
任意の複体はその項から truncation と直和によって生成される。
従ってこの kernel は torsion 加群から localizing に生成される。
\end{proof}

\begin{remark}[localization 三角そのものの係数変更]
\label{rem:relative-cubic-localization-triangle-basechange}
補題\ref{lem:relative-cubic-torsion-comparison}の比較は
$\mathbf R\Gamma_A(A)\to A$ という標準写像も保つ。
任意の $P\to S$ について、模型
\eqref{eq:relative-bounded-torsion-complex} を $N=A$ に適用すると、
各 $K_n$ は有限射影 $B$-加群の複体なので、項別に
\[
 \underline{\operatorname{Hom}}_B(K_n,A)\otimes_PS
 \simeq
 \underline{\operatorname{Hom}}_{B_S}(K_{n,S},A_S)
\]
となる。左辺の各内部次数は有限射影 $P$-加群であり、
filtered colimit 後も flat である。
横方向の複体は有界なので、ここで通常の tensor は導来 tensor を計算する。
この比較は解消の augmentation と可換するため、
$\mathbf R\Gamma_B(A)\to A$ の係数変更を標準写像ごと計算する。

係数延長は torsion の導来部分圏を保つので、localization の随伴から
\[
 \mathbf R\Gamma_A(A)\otimes_P^{\mathbf L}S
 \longrightarrow\mathbf R\Gamma_{A_S}(A_S)
\]
という標準比較が定まる。
上の $B_S$ 上の計算によりこれは擬同型であり、
$A_S$ の左右の作用を保つ比較としても擬同型である。
従ってその cone である $\mathbf R\omega\pi A$、単位 $A\to\mathbf R\omega\pi A$、
および三角の境界写像も任意係数変更に両立する。
単一 cohomology 加群の抽象的な同型だけから、三角の係数変更を推測してはいない。
\end{remark}

\begin{proposition}[三次商の相対局所双対]
\label{prop:relative-cubic-local-bimodule-dual}
任意の可換 $P$-代数 $S$ について、真正な次数付き $A_S$-双加群の導来圏で
\begin{equation}\label{eq:relative-cubic-local-bimodule-dual}
 \mathbf R\Gamma_{A_S}(A_S)
 \simeq A_S^\vee(2)[-4],
 \qquad
 \mathbf R\underline{\operatorname{Hom}}_S
       (\mathbf R\Gamma_{A_S}(A_S),S)
 \simeq A_S(-2)[4].
\end{equation}
双対は内部次数ごとに取る。
これらの同型は、指定した相対 CY 規格化の下で任意の基底変更に両立する。
\end{proposition}

\begin{proof}
$S=P$ と書く。
双加群の derived torsion は、$A$-双加群としての injective 解消を用いても計算できる。
$A$ が $P$-flat なので、この解消を片側の $A$-加群へ忘れたものも injective である。
実際、その忘却の左随伴は $A\otimes_P-$ という exact な関手である。
従って以下で残す二つの作用は、通常の片側 derived torsion を計算しながら保持できる。

中心正則元の二項解消は
\[
 \mathbf R\underline{\operatorname{Hom}}_B(A,B)
 \simeq A(3)[-1]
\]
を与える。
これは左 $B$-作用と右 $A$-作用を保持した同型であり、
cohomology が一次だけに集中するため、通常の cokernel の作用として確認できる。
central 元なので normalizing automorphism は現れない。
derived coinduction の随伴から、任意の右 $A$-加群 $M$ に対して
\[
 \mathbf R\underline{\operatorname{Hom}}_B(M,B)
 \simeq
 \mathbf R\underline{\operatorname{Hom}}_A
       \bigl(M,\mathbf R\underline{\operatorname{Hom}}_B(A,B)\bigr)
 \simeq
 \mathbf R\underline{\operatorname{Hom}}_A(M,A)(3)[-1]
\]
となる。これは Rees の change-of-rings 計算である。
体上の通常の計算は \cite[arXiv v1, Lemma~4.4 の証明および Lemma~4.5 の証明]{CSBGCCKM}
に現れるが、ここで使った二項解消と coinduction は一般の係数環上で成立する。
noetherian 性や絶対的有限大域次元は不要である。

$M=A/A_{\geq n}$ と取ると、$M$ は有限射影 $P$-加群である。
補題\ref{lem:relative-cy-finite-coefficient-dual}により
\[
 \mathbf R\underline{\operatorname{Hom}}_A(M,A)
 \simeq M^\vee(2)[-4]
\]
を得る。
ここで作用の扱いを明確にする。
derived coinduction の段階で直接保持されるのは、左 $B$ と右 $A$ の作用である。
しかし左辺の cohomology は四次だけに集中し、
その cohomology はもともと真正な $A$-双加群である。
右辺も真正な $A$-双加群である。
$B\to A$ は全射なので、cohomology 上の左 $B$-線型同型は
一意に左 $A$-線型同型となる。
従って上の同型は単一次数への集中を使って、
真正な $A$-双加群の導来同型として得られる。
restriction $D(A^e)\to D(B^e)$ の完全忠実性を仮定しているわけではない。

$n$ の filtered colimit を取れば、derived torsion の定義から第一の同型を得る。
各内部次数の $A$ 成分は有限射影なので、二重双対と係数変更はともに exact であり、
第二の同型と任意基底変更の主張が従う。
\end{proof}

\begin{remark}[ここで確定した範囲と別途必要な入力]
この節で構成したのは、大きい tails 圏の compact generator、
全 compact 間の相対 perfect RHom、任意の基底変更、
および真正な双加群としての相対局所双対である。
一般族への適用には入力 (1)--(4) の検証が必要であり、
relative CY 同型の存在を各ファイバーの CY 性だけで代替してはならない。
compact 圏上の Serre 関手の構成には、最後の局所双対を
局所化された dg 双加群と trace へ移す段階を別に記す。
線対象の RHom の階数だけから cyclic trace を結論しない。
本節は smooth proper のうち proper な側を与えるが、dg smoothness を証明していない。
\end{remark}


\subsection{導来局所化を用いた相対 Serre 双対}
\label{subsec:relative-serre-dg-localization}

有限中心模型を用いない場合，線対象間の Ext の階数を計算するだけでは，
相対 Serre pairing の両変数自然性は得られない．
この小節では，局所コホモロジーの\emph{両側加群}同型から，
線対象の導来自己準同型代数の\emph{両側加群}双対へ移す．
この手順は，高次 Ext による左作用を忘れてから推測で戻すことを避ける．

全ての次数付き双対は，内部次数を反転する係数双対
\[
 (M^{\vee_S})_n
 :=R\operatorname{Hom}_S(M_{-n},S)
\]
を意味する．同時にコホモロジー次数も通常の導来双対によって反転する．
両側加群を双対化すると左右の作用を交換する．


\begin{lemma}[線対象の同時解消による dg 局所化模型]
\label{lem:serre-dg-localization-model}
$A$ の各成分 $A_n$ は $S$-flat とする．
$\mathcal C=\operatorname{QGr}A$ は Grothendieck 圏であり，
非有界導来圏上に随伴
\[
 \pi:D(\operatorname{Gr}A)\rightleftarrows D(\mathcal C):R\omega
\]
があり，$R\omega$ は完全忠実かつ任意直和を保つとする．
$\pi$ の核は torsion cohomology を持つ複体全体とする．
このとき，内部次数付き dg 代数の写像 $A\to Q$ と，標準 dg enhancement に両立する
同値 $D(\mathcal C)\simeq D(\operatorname{Gr}Q)$ を構成できる．
$Q$ の基礎 $A$-両側加群は $R\omega\pi A$ であり，
$A\to Q$ は標準随伴の単位を表す．
\end{lemma}

\begin{proof}
$E=\pi A$ とする．内部 shift は $\mathcal C$ の完全自己同値である．
$\mathcal C$ 上の加法的 monad
\[
 \mathbb T(M)=\bigoplus_{n\geq0}A_n\otimes_SM(-n)
\]
を，$A$ の積と単位によって定める．
その加群は，同時に整合する写像
$A_n\otimes_SM\to M(n)$ を持つ対象，すなわち左 $A$-作用付き対象である．
各 $A_n$ は係数平坦なので $\mathbb T$ は完全であり，任意余極限を保つ．
この加群圏は Grothendieck 圏である．実際，核・余核と filtered colimit は
忘却関手が作り，$\mathcal C$ の生成族の自由 $\mathbb T$-加群が生成族となる．
自由加群関手が完全なので，忘却関手は injective 対象を injective 対象へ送る．

$E$ には左乗法から左 $A$-作用がある．
この作用を残した加群圏内で $E\to I^\bullet$ を通常の非負次数 injective 解消として取る．
作用を忘れた $I$ とその全ての内部 shift $I(n)$ は，
$\mathcal C$ 内の有界下 injective 複体，従って K-injective である．
\[
 Q_n=\operatorname{Hom}^{\bullet}_{\mathcal C}(I,I(n))
\]
と置き，$f\in Q_m$，$g\in Q_n$ の積を
$f(n)\circ g$ と定める．
これは strict な内部次数付き dg 代数であり，
$I$ の左 $A$-作用は strict な dg 代数写像 $A\to Q$ を与える．
従って内部 shift の strictification を別途仮定する必要はない．

線対象 $E(n)$ は compact である．実際，
\[
 R\operatorname{Hom}_{\mathcal C}(E(n),-)
 \simeq R\operatorname{Hom}_{\operatorname{Gr}A}(A(n),R\omega(-))
\]
であり，右辺は任意直和を保つ．
$A(n)$ が $D(\operatorname{Gr}A)$ を生成し $\pi$ は導来商なので，
$E(n)$ は $D(\mathcal C)$ を生成する．
従って，Keller の dg Morita 定理
\cite[Theorem~3.8(c), \S3.8]{KellerDgCategories}
をこの線対象族に適用すると，
$D(\mathcal C)\simeq D(\operatorname{Gr}Q)$ を得る．
この同値は，K-injective な代表 $N$ に対して
$\bigoplus_n\operatorname{Hom}^{\bullet}_{\mathcal C}(I,N(n))$
を取る dg Yoneda 関手により与えられる．

その $A$ への restriction を確認する．
$E\to I$ は左 $A$-作用と両立する擬同型であり，$N$ は K-injective なので
\[
 \bigoplus_n\operatorname{Hom}^{\bullet}_{\mathcal C}(I,N(n))
 \xrightarrow{\sim}
 \bigoplus_n\operatorname{Hom}^{\bullet}_{\mathcal C}(E,N(n))
\]
は右 $A$-加群の擬同型である．
右辺は section 関手の定義と随伴により $R\omega(N)$ を表す．
特に $N=I$ の場合，左作用も保つため，$Q$ は $R\omega\pi A$ を
$A$-両側加群として表す．$A\to Q$ と随伴の単位との一致も左乗法の定義から従う．

従って Morita 表示の下で $R\omega$ は restriction であり，
その左随伴 $\pi$ は extension $-\otimes_A^{\mathbf L}Q$ である．
restriction の完全忠実性により extension/restriction の余単位は同型となる．
自由対象 $Q$ における余単位を取れば
$Q\otimes_A^{\mathbf L}Q\xrightarrow{\sim}Q$ を得る．
この写像は積そのものであり，左右の作用を保つ．
\end{proof}

この構成は有限中心性を使わない．
非有界圏上の導来商と $R\omega$ の直和保存は，補題の明示した入力であり，
線対象間の Ext 群の次元だけから推定するものではない．


\begin{lemma}[homological epimorphism の両側局所化]
\label{lem:serre-dg-homological-epimorphism}
$S$ を可換環，$A$ を非負次数付き $S$-代数とし，各 $A_n$ は有限射影とする．
$A\to Q$ を内部次数を持つ $S$-線型 dg 代数の写像で，
\[
 Q\otimes_A^{\mathbf L}Q\xrightarrow{\sim}Q
\]
を満たすものとする．以下を仮定する．
\begin{enumerate}
\item $-\otimes_A^{\mathbf L}Q$ は tails torsion の導来圏を核とする局所化である．
\item $A$-両側加群の導来圏における標準三角
\[
 \Gamma\longrightarrow A\longrightarrow Q\longrightarrow\Gamma[1]
\]
があり，ある同型
\[
 \Gamma^{\vee_S}\xrightarrow{\sim}A(-2)[4]
 \quad\text{in }D(\operatorname{Gr}A^e)
\]
が指定されている．
\end{enumerate}
このとき，真正な $Q$-両側加群として自然な同型
\begin{equation}\label{eq:serre-localization-dual-bimodule}
 Q^{\vee_S}\xrightarrow{\sim}Q(-2)[3]
 \quad\text{in }D(\operatorname{Gr}Q^e)
\end{equation}
を得る．ここで $Q^e$ は必要なら導来 tensor 積で定め，
係数環上 K-flat な dg 模型を用いる．
\end{lemma}

\begin{proof}
$Q^{\vee_S}$ は，$Q$ の積を双対化することによる $Q$-両側加群である．
両側局所化を
\[
 \mathcal L^e(M)=Q\otimes_A^{\mathbf L}M\otimes_A^{\mathbf L}Q
\]
と書く．homological epimorphism の恒等式により，
任意の $Q$-両側加群 $N$ について
$\mathcal L^e(N)\simeq N$ である．
実際，左側については導来 tensor の結合則によって
\[
 Q\otimes_A^{\mathbf L}N
 \simeq(Q\otimes_A^{\mathbf L}Q)\otimes_Q^{\mathbf L}N
 \xrightarrow{\sim}Q\otimes_Q^{\mathbf L}N\simeq N
\]
であり，右側にも同じ計算を適用する．
これらの同型は $N$ の残りの $Q$-作用も保つ．
従って $\mathcal L^e(Q^{\vee_S})\simeq Q^{\vee_S}$ である．

$A_n$ は $S$-射影なので $A^{\vee_S}$ は通常の次数付き双対加群に等しい．
その内部次数は非正であり，次数 $-d$ の元は $A_{\geq d+1}$ の右作用で零となる．
よって $A^{\vee_S}$ は tails torsion であり，
\[
 A^{\vee_S}\otimes_A^{\mathbf L}Q=0,
 \qquad\mathcal L^e(A^{\vee_S})=0.
\]

指定した三角を係数双対化し，その後で $\mathcal L^e$ を適用すると，
$D(\operatorname{Gr}Q^e)$ の三角
\[
 Q^{\vee_S}\longrightarrow0\longrightarrow
 Q(-2)[4]\longrightarrow Q^{\vee_S}[1]
\]
を得る．最後の非零項では
\[
 \mathcal L^e(A(-2)[4])
 \simeq(Q\otimes_A^{\mathbf L}Q)(-2)[4]
 \simeq Q(-2)[4]
\]
を使った．この三角の境界写像が所望の同型を与える．
全ての操作を両側加群の導来圏で行ったので，
この同型は $A$ の左右の作用だけでなく，$Q$ の高次コホモロジーから来る
左右の作用にも両立する．
\end{proof}

\begin{theorem}[compact tails の相対 Serre 判定]
\label{thm:relative-serre-localization-criterion}
$S$，$A$，$Q$ は前補題の仮定を満たすとする．
さらに各内部次数 $Q_n$ は $S$-perfect とし，
標準 dg enhancement と両立する同値
\[
 D(\operatorname{QGr}A)\simeq D(\operatorname{Gr}Q)
\]
が，線対象 $\pi A(n)$ を自由 dg 加群 $Q(n)$ に送るとする．
このとき任意の compact 対象 $K,L$ に対し，自然な同型
\begin{equation}\label{eq:relative-serre-localization-formula}
 R\operatorname{Hom}_S
 \bigl(R\operatorname{Hom}_{\operatorname{QGr}A}(K,L),S\bigr)
 \xrightarrow{\sim}
 R\operatorname{Hom}_{\operatorname{QGr}A}(L,K(-2)[3])
\end{equation}
がある．両辺は $S$-perfect であり，$(-2)[3]$ は compact 圏の相対 Serre 自己同値である．

ある係数環 $P$ 上でこれらの入力が構成され，指定した全ての $P\to S$ について
\[
 A_S=A_P\otimes_PS,
 \qquad Q_S\simeq Q_P\otimes_P^{\mathbf L}S
\]
および標準三角と $\Gamma^{\vee}$ の指定同型が係数変更に両立すると仮定する．
この場合，上の相対 Serre 同型と trace は，指定した任意の導来係数変更に両立する．
\end{theorem}

\begin{proof}
$D(\operatorname{Gr}Q)$ の compact 対象は，$Q(n)$ を有限直和，
コホモロジー shift，cone，直和因子について閉じて得る perfect dg 加群である．
$Q_n$ の係数 perfect 性から，二つの自由 shift の導来 Hom は $S$-perfect である．
有限な上記の操作で閉じるので，全ての compact 対象間の導来 Hom が $S$-perfect になる．

係数双対の $Q$-両側加群 $Q^{\vee_S}$ には，評価と導来 tensor--Hom 随伴による
標準的な dg 双対性
\[
 R\operatorname{Hom}_S(R\operatorname{Hom}_{\operatorname{Gr}Q}(K,L),S)
 \simeq
 R\operatorname{Hom}_{\operatorname{Gr}Q}
 (L,K\otimes_Q^{\mathbf L}Q^{\vee_S})
\]
がある．有限自由 shift の場合は係数双対の定義そのものである．
評価による写像は $Q$-両側加群の作用に両立するので，
これは両変数に自然な dg の写像であり，有限 cell 加群と直和因子へ延びる．
前補題の $Q$-両側加群同型を代入すると
$K\otimes_Q^{\mathbf L}Q^{\vee_S}\simeq K(-2)[3]$ を得る．
内部次数 shift とコホモロジー shift は可逆なので，これは相対 Serre 自己同値である．

trace は，上の双対同型の $L=K$ の場合と恒等射の双対により
\[
 \tau_K:R\operatorname{Hom}(K,K(-2)[3])
 \longrightarrow R\operatorname{Hom}_S(R\operatorname{Hom}(K,K),S)
 \xrightarrow{(\operatorname{id}_K)^*}S
\]
と定める．特に自由対象 $Q$ の場合は，
$Q^{\vee_S}_0=R\operatorname{Hom}_S(Q_0,S)$ から単位元で評価する写像である．
一般の pairing は導来合成と $\tau_K$ の合成である．
この定義から両変数自然性と cyclic な整合性が従い，
各 fibre で別々に trace を選ぶ操作はしない．

係数変更する対象を，元の係数環上の有限 cell 表示の直和因子として取る．
その表示と分裂を新しい係数環へ延長して，対象の導来係数変更を計算する．
各 $Q_n$ が係数 perfect なので，Hom の係数変更，その係数双対の変更は同型である．
両側局所化，標準三角，境界写像，評価，単位は全て指定した係数変更に自然である．
従って Serre 同型と trace の全ての整合性も係数変更に保たれる．
\end{proof}


\paragraph{係数変更に関する dg の整合性．}
任意の可換係数変更 $S\to T$ について，
$-\otimes_S^{\mathbf L}T$ は torsion 加群を torsion cohomology の複体へ送る．
実際，係数の自由解消を用いると各項は元の torsion 加群の直和である．
既に証明した導来 torsion 部分圏の記述により，torsion cohomology の複体全体は
torsion 加群から localizing に生成される．導来 tensor が cone と任意直和を保つので，
同じ結論は任意の非有界 torsion cohomology 複体にも成立する．
従って係数変更は dg 導来商にも降り，$\pi A_S(n)$ を $\pi A_T(n)$ に送る．
その導来 Hom 上の作用は，cofibrant な dg 模型を用いて
\[
 Q_S\otimes_S^{\mathbf L}T\longrightarrow Q_T
\]
という dg 代数の比較写像を与える．この写像は単位と全ての dg 合成に両立する．
injective 解消そのものが底変更で injective のままであるとは仮定していない．

定理\ref{thm:relative-cubic-universal-line-cohomology}の標準の Hom 底変更比較により，
この $Q$ の比較写像は各内部次数で擬同型となる．
$A$ の係数平坦性と標準三角の比較から，その fibre $\Gamma_A(A)$ と
$A$ への写像も同じ係数変更に両立する．
各 $Q_n$ は係数 perfect なので
\[
 Q_S^{\vee_S}\otimes_S^{\mathbf L}T
 \simeq(Q_S\otimes_S^{\mathbf L}T)^{\vee_T}
 \simeq Q_T^{\vee_T}
\]
も両側作用を含めて成立する．
相対 Serre 同型に用いた三角，その境界写像，両側局所化，
単位での評価はこれらの比較に自然である．
これが相対 trace と双対性の底変更の整合性を与える．



\begin{theorem}[中心三次商の普遍的相対 Serre 双対]
\label{thm:koszul-cubic-relative-serre}
本節の相対入力 \textup{(1)--(4)} の下で，任意の可換 $P$-代数 $S$ に対し，
$D(\operatorname{QGr}A_S)^c$ は $S$-proper であり，
$\mathsf S=(-2)[3]$ を相対 Serre 自己同値とする．
すなわち任意の compact 対象 $K,L$ について
\[
 R\operatorname{Hom}_S
 \bigl(R\operatorname{Hom}_{\operatorname{QGr}A_S}(K,L),S\bigr)
 \simeq R\operatorname{Hom}_{\operatorname{QGr}A_S}(L,K(-2)[3])
\]
であり，この同型，合成 pairing と trace は任意の可換係数変更と両立する．
相対 CY 同型の規格化は双対を構成するために用いるが，
Abel 圏のモジュライに追加の標識として要求しない．
\end{theorem}
\begin{proof}
補題\ref{lem:relative-cubic-torsion-comparison}と
命題\ref{prop:relative-cubic-local-bimodule-dual}により，
$R\omega_A$ は直和を保ち，局所化三角の fibre は
$\Gamma^{\vee_S}\simeq A_S(-2)[4]$ を真正な双加群として満たす．
補題\ref{lem:serre-dg-localization-model}により $A_S\to Q_S$ を構成する．
定理\ref{thm:relative-cubic-universal-line-cohomology}から
全ての $(Q_S)_n$ は $S$-perfect である．
従って補題\ref{lem:serre-dg-homological-epimorphism}と
定理\ref{thm:relative-serre-localization-criterion}が適用できる．
任意係数変更の整合性は直前の dg 比較と，
命題\ref{prop:relative-cubic-local-bimodule-dual}の双加群同定の自然性による．
ここで要求したのは線対象の compact 性であり，
$D^b(\operatorname{qgr}A_S)$ 全体の有限大域次元ではない．
\end{proof}

\paragraph{既知文献の適用範囲．}
Mizuno, \emph{Dualizing complexes over $\mathbb Z$-algebras},
\cite{MizunoZDuality}, arXiv:2509.13073v3（2026年9月4日）は係数体上の理論であり（\S2冒頭），
Theorem~5.8 はさらに $\operatorname{qgr}$ の有限 global dimension を仮定する．
Brown--Sridhar, \emph{Serre duality for dg-algebras},
\cite{BrownSridharSerre}, arXiv:2410.07204v2（2025年8月15日）も係数体上の理論である（\S2冒頭）．
Brown--Soto Levins--Sridhar, \emph{Existence of balanced dualizing dg-modules},
\cite{BrownSotoSridharBalanced}, arXiv:2506.02398v1 は，balanced 双対の存在と perfect tails の Serre 双対を
さらに扱うが，その Notation~1.4 と Setup~2.1 も係数体を仮定する．
従って，いずれも次数零部分が一般の係数環である本相対結論へ
そのまま適用したものではない．
一方，局所コホモロジーの双対と $R\omega\pi A$ を結び付ける考え方は，
Yekutieli--Zhang の非可換 Serre 双対\cite{YZ}と，これらの一般化に属する．
本小節の役割は，明示した相対の\emph{両側}局所双対入力が得られた後に，
trace と compact 全体の自然性を導来局所化の形式計算で供給することである．
この具体的な組合せの初出は未確認である．


\section{一般十七係数族と構造対象付き Abel 圏の比較}
\label{sec:general17-relative-algebra}

有限中心による射影模型とは別に、二つの中心一次変数と
有限な potential 複体を用いて、一般十七係数族の相対入力を構成する。
以下の基底は明示的な非空開集合であり、楕円曲線の自己同型の
有限位数性を要求しない。
大きい $\operatorname{QGr}$、切断の飽和性、相対 CY 同型、偶数列による復元は
任意の可換係数環への基底変更について述べる。
noetherian 性は最後に、noetherian な係数環について別に証明する。

\subsection{明示的な開集合と楕円境界}

定理\ref{thm:17}の式を
\begin{equation}\label{eq:general17-explicit-potential}
 F=a(x^3+y^3+z^3)+bxyz+ux^2+\alpha vy^2+\beta(u+v)z^2
   +\sum_{i=1}^3 M_i(u,v)x_i+g(u,v)
\end{equation}
と書く。ここで $(x_1,x_2,x_3)=(x,y,z)$、
$M_i=\sum_{j=0}^2m_{ij}u^{2-j}v^j$、
$g=\sum_{j=0}^3g_ju^{3-j}v^j$ である。
以下の $p,q,r$ は Rains 公式の $q$ とは異なる。
\begin{align*}
 p&=2-b,&q&=-(2+b),&r&=-6a,\\
 D_E&=(p^3+q^3+r^3)^3-27p^3q^3r^3,\\
 L_6&=p^6+q^6+r^6-p^3q^3-q^3r^3-r^3p^3,\\
 \Delta&=a(27a^3+b^3)\,pqr\,D_E
 (p^3-q^3)(q^3-r^3)(r^3-p^3)L_6.
\end{align*}
\begin{equation}\label{eq:general17-coefficient-ring}
 P_0=\C[a,b,\Delta^{-1}],\qquad
 P=P_0[\alpha,\beta,(m_{ij})_{1\leq i\leq3,\,0\leq j\leq2},g_0,g_1,g_2,g_3]
\end{equation}
と置く。これは十七次元の整域である。
$a=1/6,b=0$ では $(p,q,r)=(2,-2,-1)$ であり、
表示した七つの因子は順に
\[
 \frac1{48},\quad4,\quad-1729,\quad16,\quad-7,\quad-9,\quad193
\]
となるので $\Delta\ne0$ である。
必要なら定理\ref{thm:17}の slice の非零小行列式も反転する。
その多項式と $\Delta$ の積は零多項式ではないので、共通開集合も非空である。
表示した多項式恒等式と非消滅点は
\nolinkurl{computations/verify_general17_open.py} で有理数上に再現できる。
これは slice の横断性の条件と相対代数の条件を同じ数値点で確認したという主張ではない。

最高次の三変数代数を $\mathscr S$ と書く。その関係は
\[
 p yz+q zy+r x^2=0,\qquad
 p zx+q xz+r y^2=0,\qquad
 p xy+q yx+r z^2=0
\]
である。$a(27a^3+b^3)$ の可逆性により、古典的な最高次三次式は滑らかである。
Mori--Smith \cite[Theorem~4.2(1), Proposition~2.8 and Corollary~2.9]{MS} から、その全ての複素体値は
三次元 Koszul AS 正則 CY 整域となる。
ここで使うのは commutator を含む potential の場合であり、
同論文の対称 potential に現れる別の例外条件を混同しない。

\begin{lemma}[楕円境界の相対表示]\label{lem:general17-relative-elliptic-boundary}
$P_0$ 上の滑らかな平面三次曲線
\[
 \mathscr E:\quad pqr(X^3+Y^3+Z^3)-(p^3+q^3+r^3)XYZ=0
\]
には、三つの双線形関係で定まる $P_0$-自己同型 $\sigma$ が存在する。
全ての幾何学的ファイバーで $\sigma,\sigma^2,\sigma^3$ は恒等でない。
$\mathcal L=\mathcal O_{\mathscr E}(1)$ とし
\[
 \mathcal L_n=\mathcal L\otimes\sigma^*\mathcal L\otimes\cdots
       \otimes(\sigma^{n-1})^*\mathcal L,
 \qquad H=\bigoplus_{n\geq0}\Gamma(\mathscr E,\mathcal L_n)
\]
を twisted homogeneous coordinate algebra とする。
$H_0=P_0$、$H_n$ は $n\geq1$ について階数 $3n$ の有限射影 $P_0$-加群で、
$H$ の形成は任意の基底変更と可換する。
標準の切断写像 $\mathscr S\to H$ は次数付き全射である。
\end{lemma}

\begin{proof}
双線形関係の第一変数側の行列は
\[
 M(X,Y,Z)=
 \begin{pmatrix}
 rX&qZ&pY\\pZ&rY&qX\\qY&pX&rZ
 \end{pmatrix}.
\]
その行列式は $\mathscr E$ の方程式の負である。
$pqrD_E$ の可逆性により各ファイバーは滑らかである。
$\mathscr E$ 上で行列の階数が一以下なら、全ての二次小行列式が零となり、
行列式の全偏微分も零となるので矛盾する。
従って階数は常に二で、核の可逆部分加群が第二射影座標を定める。
第二変数側の行列でも同じ議論が成立するので、双線形関係の対応の両射影は同型である。
これにより $\sigma$ を得る。この構成は任意の基底変更に両立する。

切断 $O=(1:-1:0)$ に対して直接計算すると
\[
 \sigma O=(p:q:r),\qquad
 \sigma^2O=
 \bigl(q(r^3-p^3):p(q^3-r^3):r(p^3-q^3)\bigr).
\]
$r$ および $r(p^3-q^3)$ が可逆なので、前二者は $O$ と異なる。
$M(\sigma^2O)$ の第一・第二行の外積の第三成分は $-pqr^2L_6$ であり、
これも可逆である。従って $\sigma^3O$ の第三成分は零でなく、$\sigma^3O\ne O$ である。
これらは係数環上の多項式恒等式である。

$\mathcal L_n$ のファイバー次数は $3n$ である。
種数一の曲線の線束コホモロジーと proper flat base change から、
$n\geq1$ では $R^1f_*\mathcal L_n=0$、$f_*\mathcal L_n$ は階数 $3n$ の有限局所自由加群であり、
任意の基底変更と可換する。
$H_0=P_0$ とその基底変更は平面三次の標準完全列からも直接従う。
双線形関係は $\sigma$ のグラフ上で零なので、次数一切断から $\mathscr S\to H$ が定まる。

全射性をファイバーで確認する。
滑らかな楕円曲線上で $L$ の次数が三、$M$ の次数が三以上なら
$H^0(L)\otimes H^0(M)\to H^0(LM)$ は全射である。
実際、$H^0(L)$ の基点のない二次元部分空間の pencil の完全列を使えば、
$\deg M>3$、または $\deg M=3$ かつ $M\not\simeq L$ の場合は
$H^1(ML^{-1})=0$ から従う。
$M\simeq L$ の場合は平面三次の埋込みに二次関係がないことから従う。
よって各ファイバーで $H$ は次数一生成である。
切断加群が有限局所自由であることと Nakayama の補題から、相対的にも全射となる。
\end{proof}

\subsection{普遍 PBW 性と中心正則三次元}

\begin{proposition}[全係数にわたる射影性と正則性]
\label{prop:general17-universal-pbw}
$P$ 上で \eqref{eq:general17-explicit-potential} の関係により $B$ を定める。
命題\ref{prop:pair-rains-polynomial}の規格化で中心三次元 $w=\mathcal G+g(u,v)$ を選び、
$A=B/(w)$ と置ける。
各 $B_n,A_n$ は有限射影 $P$-加群であり、
\[
 \operatorname{rank}_P B_n=\binom{n+4}{4},\qquad
 \operatorname{rank}_P A_n=\binom{n+4}{4}-\binom{n+1}{4}
\]
である。右辺の負次数項は零と読む。
$Q=P[u,v]$ と置くと、$B$ と $A$ は $Q$-射影加群であり、active 次数の濾過について
\begin{equation}\label{eq:general17-relative-associated-graded}
 \operatorname{gr}B\simeq\mathscr S\otimes_{P_0}Q,
 \qquad \operatorname{gr}A\simeq H\otimes_{P_0}Q.
\end{equation}
これらの同型、$w$ の正則性、および各次数の完全列
\[
 0\longrightarrow B_{n-3}\xrightarrow{w}B_n\longrightarrow A_n\longrightarrow0
\]
は任意の可換 $P$-代数への基底変更後にも成立する。
\end{proposition}

\begin{proof}
まず各 $\mathscr S_n$ は有限自由な語の加群の有限表示商であり、全閉点で
次元 $\binom{n+2}{2}$ を持つ。
被約な有限型 $\C$-基底上の一定階数の行列を局所的に標準形へ直すと、
$\mathscr S_n$ は有限局所自由で、任意の基底変更に両立することが分かる。
ここで行列の必要な小行列式が全閉点で零なら、基底の被約性により恒等的に零である。
各語の加群から $\mathscr S_n$ への全射の $P_0$-線形分裂を選ぶ。

active 次数を下げる帰納法により、これらの持上げは $B$ を $Q$-加群として生成する。
全閉点への特殊化では、低次 cyclic potential に対する PBW 定理
\cite[Theorem~3.1]{BT}\cite[Proposition~7.3.1(ii)]{EG} が適用できる。
従って持上げ間の有限 $Q$-線形関係を全閉点で評価すれば全係数が零となる。
有限射影加群の切断は被約な有限型基底の全閉点で検出されるので、関係自体が零である。
よって $B\simeq\bigoplus_{n\geq0}\mathscr S_n\otimes_{P_0}Q$ が
濾過付き $Q$-加群として成り立ち、第一の associated graded 同型は普遍的である。

Rains 公式への代入を明示する。最高次三次式を $3a$ で割り、
\[
 \lambda=\frac b{3a},\quad h=3a,\quad
 t=\frac{6a}{2-b}=-\frac rp,\quad
 q_{\mathrm R}=\frac{2+b}{2-b}=-\frac qp
\]
と置く。対角係数は
$L_1=2u/(3a), L_2=2\alpha v/(3a), L_3=2\beta(u+v)/(3a)$、
一次係数は $M_i/(3a)$ である。
公式の分母は
\[
 D(t)=q_{\mathrm R}(q_{\mathrm R}^3-t^3)(q_{\mathrm R}+1)(1+t^3).
\]
ここで $q_{\mathrm R}+1=4/p$、
$q_{\mathrm R}^3-t^3=(r^3-q^3)/p^3$、
$1+t^3=(p^3-r^3)/p^3$ であるから、全て $\Delta\ne0$ 上で可逆である。
同命題の有限公式は、active 定数部分が零の $\mathcal G$ と、
全ファイバーで非零な最高記号 $\Omega\in Z(\mathscr S)_3$ を与える。
四つの自由な定数係数はこの規格化で $g$ と書く。
別の可逆な規格化を使う場合は、その係数変換も同時に行う。

各体値の $H$ は $k(\mathscr E)[t,t^{-1};\sigma]$ に埋め込まれる。
次数 $d>0$ の中心元 $ct^d\ne0$ があれば、次数一切断との可換性から
$\sigma^d$ はその比を全て固定する。
その比は $k(\mathscr E)$ を生成するので $\sigma^d=\mathrm{id}$ となる。
前補題により $d=1,2,3$ は起こらない。
従って中心元 $\Omega$ の $H_3$ における像は全閉点で零である。
$H_3$ は有限射影、$P_0$ は被約なので、像は大域的に零である。
各閉点で $\mathscr S$ は整域、$\Omega\ne0$ なので、$\Omega$ は正則である。
Hilbert 級数と $H_n$ の階数を比較すると
\begin{equation}\label{eq:general17-elliptic-quotient}
 0\longrightarrow\mathscr S(-3)\xrightarrow{\Omega}\mathscr S
   \longrightarrow H\longrightarrow0
\end{equation}
は各次数で有限局所自由加群の完全列となる。
局所的な分裂により、この列は任意の基底変更後にも完全である。

$w$ の最高記号が $\Omega$ であるから、濾過付き正則元で割る標準の
最高次数からの帰納法により第二の associated graded 同型を得る。
その各次数商 $H_n\otimes Q$ は射影なので、濾過を次数ごとに分裂できる。
従って $A\simeq\bigoplus_{n\geq0}H_n\otimes_{P_0}Q$ であり、
正則性と associated graded の計算は非平坦な基底変更後にも有効である。
通常の全次数で各成分を数えれば階数公式と最後の分裂完全列を得る。
\end{proof}

\subsection{相対 CY orientation と有限 Koszul 対角解消}

\begin{proposition}[全係数に両立する五次元 CY 入力]
\label{prop:general17-relative-cy-koszul}
命題\ref{prop:general17-universal-pbw}の $B$ には、任意の基底変更に両立する
真正な次数付き双加群の同型
\begin{equation}\label{eq:general17-relative-cy-orientation}
 \mathbf R\underline{\operatorname{Hom}}_{B\otimes_PB^{\mathrm{op}}}
       (B,B\otimes_PB^{\mathrm{op}})\simeq B(5)[-5]
\end{equation}
がある。さらに標準の二次 Koszul 双加群複体
\[
 K_B^{-i}=B\otimes_P W_i\otimes_P B\quad(0\leq i\leq5)
\]
は $B$ の解消であり、各 $W_i$ は内部次数 $i$ の有限射影加群、
$\operatorname{rank} W_i=\binom5i$、$W_0=P$、$W_5$ は可逆である。
片側を augmentation と tensor した複体は $P$ の解消となる。
全ての完全性は任意の可換基底変更後にも保たれる。
\end{proposition}

\begin{proof}
$Q=P[u,v]$ を係数環とし、三つの active 変数の cyclic potential を
\[
 \Phi=xyz-yxz-\operatorname{Symm}(F-g)
\]
と取る。$r_i=\partial_{x_i}\Phi$ は $B$ の三関係である。
$D=B\otimes_QB^{\mathrm{op}}$ と書く。
形式的な基底 $e_i,f_i,e$ の全次数をそれぞれ $1,2,3$ として
\[
 0\longrightarrow D(-3)\longrightarrow D(-2)^3
   \longrightarrow D(-1)^3\longrightarrow D\longrightarrow B\longrightarrow0
\]
を作る。微分は双加群表記で
\begin{align*}
 d_1(1\otimes e_i\otimes1)&=x_i\otimes1-1\otimes x_i,\\
 d_2(1\otimes f_i\otimes1)&=\mathcal D(r_i),\\
 d_3(1\otimes e\otimes1)&=
       \sum_i(x_i\otimes f_i\otimes1-1\otimes f_i\otimes x_i)
\end{align*}
である。$\mathcal D$ は active 語の一文字を順に中間因子へ置く普遍微分で、
$Q$ 上では零とする。
cyclic 恒等式 $\sum_i[x_i,r_i]=0$ によりこれは複体となる。
active 濾過の associated graded は $\mathscr S$ の三次元 potential 解消を
$Q$ へ延長した複体である。
この $Q$-延長前の $\mathscr S$ の複体の各次数は有限射影 $P_0$-加群の有限複体で、
全閉点で完全であるから、
局所的に分裂して普遍的に完全である。
濾過が各全次数で有限であることから、元の複体も全基底変更後に完全である。

自己双対はファイバーの CY 性から推論しない。
基底 $e_i,f_i$ を対応させ、$1,e$ を対応させる双対写像を使う。
中間微分について必要な恒等式は
\[
 \partial^{\mathrm{nc}}_{x_j}(\partial_{x_i}\Phi)
 =\operatorname{flip}\bigl(
       \partial^{\mathrm{nc}}_{x_i}(\partial_{x_j}\Phi)\bigr)
\]
である。これは各 cyclic 語における二つの文字の位置を入れ替える恒等式であり、
全ての係数環上で成立する。
両端の微分は対応する差 $x_i\otimes1-1\otimes x_i$ から一致する。
双対複体の標準符号を入れれば、真正な双加群複体としての自己双対
$K_Q^\vee\simeq K_Q(3)[-3]$ を得る。
これは Berger--Taillefer \cite[arXiv v2, Theorem~3.6 の証明, pp.~9--12]{BT} の
明示的な微分・符号・cyclic Hessian 恒等式を、上の射影性と普遍的完全性の下で
$Q$ 上にそのまま書いたものである。
従って $\Theta_{B/Q}\simeq B(3)[-3]$ である。

次に $C=B\otimes_PB^{\mathrm{op}}$ と置く。
作用の規約を固定する。$B$-双加群は左 $C$-加群として扱い、
その双対は右 $C$-加群である。
右 $C$-加群 $N$ を双加群へ戻す作用は
$a\cdot n\cdot b=n\cdot(b\otimes a^{\mathrm{op}})$ とする。
$B$ は $Q$-射影なので、$C$ は
$Q\otimes_P Q=P[u_L,v_L,u_R,v_R]$ 上平坦である。
従って $z_1=u_L-u_R,z_2=v_L-v_R$ は $C$ 上の中心正則列であり、
$D=C/(z_1,z_2)$ である。
順序 $(z_1,z_2)$ を固定した Koszul 解消とその双対は
\[
 \mathbf R\underline{\operatorname{Hom}}_C(D,C)\simeq D(2)[-2]
\]
を与える。この同型は残る両側作用を保持する。
ここで左辺の左 $D$-作用は定義域 $D$ の右掛算による前合成、
右 $C$-作用は値域の右掛算である。
唯一のコホモロジーは通常の $D$--$C$ 双加群 $D(2)$ であり、
その左作用は $d$ の任意の $C$-持上げによる右掛算を
中心 Koszul 解消の chain map へ持ち上げて計算できる。
商においてこれは左からの通常の $d$ の掛算となり、右作用は $C\to D$ を通る。
導来随伴と前段の自己双対を合わせれば
\[
 \mathbf R\underline{\operatorname{Hom}}_C(B,C)
 \simeq
 \mathbf R\underline{\operatorname{Hom}}_D
   \bigl(B,\mathbf R\underline{\operatorname{Hom}}_C(D,C)\bigr)
 \simeq B(5)[-5].
\]
この随伴は右 $C$-作用を保持している。
最後の同型は $\Theta_{B/Q}$ の右 $D$-作用を $C\to D$ で制限したものであり、
最初に固定した反転規約により通常の $B$-双加群となる。
$D$ は $C$-perfect、$B$ は $D$-perfect なので $B$ は $C$-perfect である。
この構成の規格化は active 三変数の体積と $(u,v)$ の順序によるものであり、
中山自己同型や基底上の未指定の線束を後から挿入していない。
有限射影複体と表示した恒等式により、全ての同型は任意の基底変更に両立する。

最後に二次 Koszul 解消を整える。
$V=B_1\simeq P^5$、$R_2=\ker(V\otimes_PV\to B_2)$ とし
\[
 W_i=\bigcap_{j=0}^{i-2}
 V^{\otimes j}\otimes R_2\otimes V^{\otimes(i-2-j)}
 \ \subset V^{\otimes i}\qquad(i\geq2)
\]
と置く。
各体値で、active 濾過の associated graded は
Koszul 代数 $\mathscr S\otimes k[u,v]$ である。
全次数ごとに有限な filtered bar 複体の spectral sequence により $B$ も Koszul である。
Hilbert 級数 $(1-t)^{-5}$ から、その Koszul 双対の階数は
$\binom5i$、$i>5$ で零となる。
$W_i$ は有限射影加群間の有限行列の核として表せるので、
$2\leq i\leq6$ の一定階数により $W_i$ は有限射影で基底変更と可換する。
特に $W_6=0$ であり、高次の交わりも零である。
標準の Koszul 双加群複体の各内部次数は有限射影複体で、
全閉点で完全だから相対的にも分裂完全である。
これで普遍的完全性を得る。
この複体を片側 $B$-加群として見れば、射影対象 $B$ の有界な射影解消であるから分裂する。
従って augmentation との tensor も解消となる。
先に構成した CY 同型をこの指定した解消でも表せる。
実際、$K_B^\vee$ は有界な射影右 $C$-複体なので、導来同型は chain map
$K_B^\vee\to B(5)[-5]$ により表される。
その cone は有界かつ acyclic で、各内部次数の全項は $P$-flat である。
従って cone は任意の基底変更後にも acyclic であり、
この有限 Koszul 模型で表した orientation も全係数変更に両立する。
\end{proof}

\subsection{切断の普遍的飽和性と二次中心}

\begin{lemma}[中心二変数による飽和性]
\label{lem:two-central-variables-universal-saturation}
$S$ は任意の可換環とし、$A$ は $A_0=S$ の非負次数付き、次数一生成代数で、
$A_1$ は有限生成とする。
中心一次元 $u,v$ が与える $S[u,v]$-加群として $A$ が平坦なら、
全ての $m\in\mathbb Z$ について $A(m)$ は tails の Gabriel 局所化に対して閉じている。
従って
\[
 A_n\xrightarrow{\sim}
 \operatorname{Hom}_{\operatorname{QGr}A}(\mathcal O,\mathcal O(n))
 \qquad(n\in\mathbb Z)
\]
である。$A$ の $S[u,v]$-平坦性が任意の係数変更後にも保たれる場合、
この同型も全ての係数変更、積、合成に自然である。
\end{lemma}

\begin{proof}
$S[u,v]$ の二つの変数の Čech 列を平坦な $A(m)$ と tensor すると
\[
 0\longrightarrow A(m)\longrightarrow A(m)_u\oplus A(m)_v
     \longrightarrow A(m)_{uv}
\]
は初めの二箇所で完全である。
また $u$ は $A(m)$ 上正則である。
tails-torsion 加群 $T$ は $u,v$ の局所化で零となるので
$\operatorname{Hom}_{\operatorname{Gr}A}(T,A(m))=0$ である。
拡大 $0\to A(m)\to E\to T\to0$ を取る。
局所化により $E_u\simeq A(m)_u$、$E_v\simeq A(m)_v$ となる。
二つの自然な写像を合わせると、次数付き $A$-線形な retraction
\[
 E\longrightarrow A(m)_u\mathop{\times}_{A(m)_{uv}}A(m)_v=A(m)
\]
を得る。
従って $\operatorname{Ext}^1_{\operatorname{Gr}A}(T,A(m))=0$ である。
これは Gabriel 商の右随伴 $\omega$ の本質像に属するための条件なので
$A(m)\simeq\omega\pi A(m)$ となる。
自由対象 $A$ との随伴を取れば切断同型を得る。
全ての写像が局所化と標準切断から定まるので自然性も従う。
ここでは高次コホモロジーの消滅や noetherian 性を使用していない。
\end{proof}

\begin{proposition}[任意係数上の切断と二次中心]
\label{prop:general17-universal-saturation-center}
\eqref{eq:general17-coefficient-ring} の任意の可換 $P$-代数 $S$ について
\[
 (A_S)_n\simeq
 \operatorname{Hom}_{\operatorname{QGr}(A_S)}(\mathcal O,\mathcal O(n)),
 \qquad
 Z(A_S)_2=S u^2\oplus S uv\oplus S v^2
\]
が成り立ち、全ての係数変更に両立する。
\end{proposition}

\begin{proof}
命題\ref{prop:general17-universal-pbw}により
$A_S$ は $S[u,v]$-射影なので、第一の主張は前補題である。
各複素体値で、補題\ref{lem:general17-relative-elliptic-boundary}と
\eqref{eq:general17-elliptic-quotient}から
$Z(\mathscr S)_1=Z(\mathscr S)_2=0$ となる。
最高 active 記号を取れば $Z(B)_2=\operatorname{Sym}^2\langle u,v\rangle$ であり、
補題\ref{lem:moduli-no-new-quadratic-center}によって $Z(A)_2$ も同じである。

普遍性を単なるファイバーの等式から推論しない。
PBW 分裂により $P u^2\oplus Puv\oplus Pv^2\subset A_2$ は直和因子である。
その商は階数十二の有限射影加群であり、三つの交換子による写像
\[
 A_2/(P u^2\oplus Puv\oplus Pv^2)\longrightarrow A_3^{\oplus3}
\]
は全閉点で単射である。
最大階数の小行列式が局所的に可逆となるので、この写像は局所的に分裂単射である。
従って任意の、非平坦な基底変更後にも核は零であり、表示した中心の等式を得る。
\end{proof}

\subsection{偶数 Gabriel 切断による全同値の復元}

\begin{lemma}[任意係数上の偶数 Veronese 局所化]
\label{lem:general-even-gabriel-reconstruction}
$A$ は可換環 $S$ 上の非負次数付き、次数一生成代数で、$A_1$ は有限生成とする。
$R=A^{(2)}$ と置くと、偶数成分を取る関手は同値
\[
 \operatorname{QGr}A\xrightarrow{\sim}\operatorname{QGr}R
\]
を誘導する。
さらに全ての整数 $i$ を用いた切断関手
\[
 \omega_{\mathrm{ev}}(N)_i=
 \operatorname{Hom}_{\operatorname{QGr}A}(\mathcal O(-2i),N)
\]
は、偶数 Veronese の Gabriel 右随伴そのものであり、完全忠実である。
これらは任意の係数作用の延長に両立する。
\end{lemma}

\begin{proof}
$F(M)=M^{(2)}$、$G(N)=N^{\mathrm{dbl}}\otimes_{R^{\mathrm{dbl}}}A$ と置く。
ここで dbl は次数を二倍にする操作である。
$FG=\operatorname{id}$ であり、counit $GFM\to M$ は偶数成分で同型である。
その核と余核は偶数成分が零なので、次数一生成性により $A_{\geq1}$ で消える。
従って両者は tails-torsion である。

$F$ は torsion を保つ。逆に $FM$ が torsion なら、偶数次数の斉次元は
十分大きい偶数次数の右因子で消え、次数一生成性から全ての十分高い次数で消える。
奇数次数の斉次元 $m$ については、$A_1$ の有限個の生成元 $x_j$ に対する
偶数次数の元 $mx_j$ を同時に消す bound を取ればよい。
これにより $M$ も torsion となる。
$FG=\operatorname{id}$ から $G$ も torsion を保ち、
同じ反映性により torsion 同値を torsion 同値へ送る。
従って商上で $F,G$ は互いに逆となる。

$F(\pi A(-2i))=\pi R(-i)$ である。
$\operatorname{Gr}R\to\operatorname{QGr}R$ の右随伴の第 $i$ 成分は
$\operatorname{Hom}(\pi R(-i),-)$ なので、表示した関手を得る。
Gabriel 右随伴は完全忠実である。
係数作用の延長との整合性は、局所化の係数作用比較と
\[
 \operatorname{Hom}_{\mathcal D_T}(T\otimes_SE,N)
   =\operatorname{Hom}_{\mathcal D}(E,UN)
\]
に含まれる $T$-作用から従う。
\end{proof}

\begin{remark}[同値の復元に必要なデータ]
従って、全偶数線対象とその間の全ての Hom を保つ完全同値は、
この切断関手を通して実際の関手として自然同型まで復元される。
線対象上の一致から一般の対象への射の持上げを仮定する必要はない。
相対 Serre 同型が偶数線列を内在的に識別する場合、
この補題を直接使える。
これは全整数添字の Gabriel 切断に関する主張であり、
正の次数だけの切断、あるいは coherent 部分圏の ample 条件と混同しない。
\end{remark}

\subsection{Noetherian 係数環上の模型の許容性}

\begin{proposition}[有限位数に依存しない相対 strong noetherian 性]
\label{prop:general17-strong-noetherian}
任意の可換 noetherian $P$-代数 $S$ に対し、
$H_S,\mathscr S_S,B_S,A_S$ は左右 noetherian である。
\end{proposition}

\begin{proof}
$\mathscr E_S\to\operatorname{Spec}S$ は滑らかな射影種数一曲線で、
$\deg(\mathcal L_n|_{\mathscr E_s})=3n$ である。
任意の可逆層 $\mathcal H$ のファイバー次数は局所一定であり、
noetherian な準コンパクト基底上では有限個の値を取る。
従って $n\gg0$ なら $\mathcal H\otimes\mathcal L_n$ は全てのファイバーで正次数を持つ。
proper morphism のファイバー上の ampleness 判定
\cite[Proposition~2.7]{KeelerAmpleFilters} により相対 ample であり、
基底が affine なので ample である。
Keeler \cite[Theorem~1.3, condition~(A4)]{KeelerAmpleFilters} から
$(\mathcal L_n)$ は ample sequence となり、
同 \cite[Theorem~7.14]{KeelerAmpleFilters} により $H_S$ は左右 noetherian である。
ここでは corrected arXiv v5 の定理を用いている。
$\sigma$ の有限位数性は使っていない。

中心正則元 $\Omega$ による下降濾過の associated graded は
$\operatorname{gr}_{\Omega}(\mathscr S_S)=H_S[z]$ である。
これは noetherian なので、任意の斉次片側イデアルの associated graded の有限斉次生成元を持ち上げる。
最高 $\Omega$-adic 成分を順に引く操作は、固定した通常次数では
$\deg\Omega=3$ によって有限回で止まる。
よって元の斉次イデアルも有限生成である。
通常次数の最高成分を取る帰納法により、非斉次イデアルも有限生成となる。
従って $\mathscr S_S$ は左右 noetherian である。
最後に \eqref{eq:general17-relative-associated-graded} と
正の active 濾過についての最高成分からの帰納法を適用すれば、
$B_S,A_S$ の左右 noetherian 性を得る。
この最後の議論は各元の active 次数が有限であることだけを用いる。
\end{proof}



\begin{proposition}[相対 Abel 模型の許容性]
\label{prop:general17-abel-admissible}
$P$ 上の $A$ は，定義\ref{def:abel-cubic-pointed}の相対模型を与える．
任意の noetherian $P$-代数 $S$ 上で
$\operatorname{qgr}(A_S)$ は noetherian Abel 圏であり，
$(\mathcal O(-i))_{i\in\mathbb Z}$ および偶数列は豊富である．
任意の可換 $P$-代数上では大きい $\operatorname{QGr}$ 模型を用い，
相対 AS 解消，切断の飽和，Gabriel 復元および標準の係数延長が成立する．
\end{proposition}
\begin{proof}
相対 AS 解消は命題\ref{prop:general17-relative-cy-koszul}の augmentation と
双対から得られ，三次核の左右正則性と標準 Hilbert 関数は
命題\ref{prop:general17-universal-pbw}である．
普遍的な飽和と偶数復元は，直前の二つの結果による．
標準線対象は $S$-平坦である．実際，$A_S(i)$ は $S$-平坦であり，
$S$-加群との tensor を完全な商関手 $\pi$ に通せばよい．

noetherian な $S$ については，命題\ref{prop:general17-strong-noetherian}により
$qgr$ が Abel である．豊富性の高次消滅も明示する．
有限生成次数付き $A_S$-加群 $M$ を，有限生成 $B_S$-加群とみる．
$B_S$ は noetherian なので，有限生成次数付き自由加群による最初の六段の解消
\[
 0\to K_6\to F_5\to\cdots\to F_0\to M\to0
\]
を取れる．命題\ref{prop:relative-as5-local-cohomology}から，
全ての $F_i$ の局所コホモロジーは内部次数が上に有界である．
共通の上界を超える内部次数 $n$ では，長完全列の dimension shifting により
\[
 H^q(R\Gamma_{B_S}M)_n
 \simeq H^{q+6}(R\Gamma_{B_S}K_6)_n=0\qquad(q\geq0)
\]
となる．最後には cohomological dimension が五以下であることを使った．
補題\ref{lem:relative-cubic-torsion-comparison}と局所化三角から，
$n\gg0$ では $R\omega_A\pi_A M$ の第 $n$ 成分は $M_n[0]$ に一致する．
これが必要な Serre 消滅である．

有限生成加群の高い truncation から，任意に遠い負の線対象の有限和による
全射を得る．全射に沿う Hom の持上げは，その kernel に対する上の
$\operatorname{Ext}^1$ 消滅による．従って標準列は豊富である．
偶数列についても，同じ消滅と偶数 Veronese の Gabriel 同値を適用できる．

任意係数延長の標準模型は
補題\ref{lem:abel-arbitrary-coefficient-tails}による係数作用付き対象の圏と一致する．
同補題の次数一有限生成性の条件はここで満たされている．
係数延長の合成と対象・射の fpqc 降下は既存の係数降下補題による．
このように非 noetherian な底には大きい圏を使い，
小さい $qgr$ の Abel 性を余分に仮定してはいない．
\end{proof}

\subsection{任意の係数環上の全ての構造対象付き同値}

\begin{theorem}[compact 双対と偶数 Gabriel 復元による相対 Hom]
\label{thm:general-relative-pointed-hom}
$S$ を可換 $\mathbb C$-代数とし，$A,A'$ を第\ref{sec:relative-koszul-line-cohomology}節の相対入力\textup{(1)--(4)}を満たす
中心三次商とする．さらに両者は
定義\ref{def:abel-two-center-family}の普遍的な中心二変数条件を満たすとする．
標準の係数族について，任意の可換 $S$-代数への変更に自然な fpqc 層の同型
\begin{equation}\label{eq:general-relative-pointed-hom}
 \underline{\operatorname{Isom}}_{\mathrm{gr},S}(A,A')/\mathbb G_m
 \xrightarrow{\sim}
 \underline{\operatorname{Eq}}_{\mathcal O,S}
       (\operatorname{QGr}A,\operatorname{QGr}A')
\end{equation}
がある．右辺には構造対象を保つ全ての完全線型同値を含め，
自然同型で割った同型類を取る．構造対象上の同定を保つ自然自己同型は自明である．
この結論に $S$ の noetherian 性や三次商の正則性は必要ない．
左辺は層商であり，大域的な同型の集合を $S^\times$ で割った集合と
無条件には同一視しない．
\end{theorem}

\begin{proof}
完全 $S$-線型同値 $F$ は非有界複体に項別に作用し，
標準 dg enhancement の同値を誘導する．同値は直和を保つため compact 対象を保つ．
線対象は系\ref{cor:relative-cubic-compact-proper}により compact であり，
定理\ref{thm:koszul-cubic-relative-serre}の双対は
この内在的部分圏上で成り立つ．相対 RHom とその係数双対を保つこと，
および表現対象の Yoneda 一意性から，指定した双対と両立する同型
\[
 F\mathsf S_A\simeq\mathsf S_{A'}F,
 \qquad \mathsf S_A=(-2)[3],\quad\mathsf S_{A'}=(-2)[3]
\]
を得る．これは全ての Abel 同値に適用でき，列保存を仮定していない．

$H_A=\mathsf S_A^{-1}[3]\simeq(2)$ と置く．
$F(\mathcal O_A)\simeq\mathcal O_{A'}$ と上の可換同型を反復することで，
全ての整数 $i$ に対する偶数線対象の同定と，その全ての Hom の同定を得る．
反復には一つの可換同型を一貫して用いるので，合成も保たれる．
普遍的な section の飽和から，これは次数付き代数同型
\[
 \phi_R:A^{(2)}\xrightarrow{\sim}(A')^{(2)}
\]
を与える．双対の表示を変更する自然変換の各線対象上の成分は単元である．
隣接する偶数線の Hom は正階数の有限射影 $S$-加群 $A_2$ なので忠実であり，
自然性からこれらの単元は全て等しい．従って表示の変更は $\phi_R$ を
一つの次数スカラーで変更するだけである．

補題\ref{lem:general-even-gabriel-reconstruction}の偶数 Gabriel 復元により，
\[
 N\longmapsto\bigoplus_{i\in\mathbb Z}
   \operatorname{Hom}(\mathcal O_A(-2i),N)
\]
は $\operatorname{QGr}A\simeq\operatorname{QGr}(A^{(2)})$ の右随伴に一致し，
完全忠実である．従って全偶数線とその Hom・合成に対する作用は，
$F$ を自然同型まで決定する．ここでは線対象を射影対象と仮定せず，
coherent 対象の有限表示にも依存しない．

定理\ref{thm:abel-relative-veronese-isom}により，$\phi_R$ の
$A\simeq A'$ への持上げ層は $\mu_2$-torsor であり，fpqc 局所で持上げ $\phi$ を選べる．
$\phi$ による構造移送と $F$ は偶数 Gabriel 右随伴の下で同じ作用を持つので，
構造対象の同定を保つ自然同型を持つ．これで層としての全射性が従う．

核を調べるため，$\phi\in\operatorname{Aut}_{\mathrm{gr},S}(A)$ の構造移送が
標識付き恒等関手に自然同型であるとする．その $\mathcal O_A(-2i)$ 上の成分を
$\lambda_i\in S^\times$ とし，$\lambda_0=1$ とする．
隣接 Hom 上の自然性により $\phi|_{A_2}=c\operatorname{id}$ となる
$c\in S^\times$ があり，次数一生成性から
$\phi|_{A_{2n}}=c^n\operatorname{id}$ となる．
忠実平坦な平方根被覆上で $\lambda^2=c$ を選ぶと，
$\varepsilon_\lambda^{-1}\phi$ は偶数部分上で恒等である．
偶数復元の $\mu_2$ も次数スカラーに含まれるため，$\phi$ は局所的に次数スカラーである．
逆に次数スカラーの構造移送は次数ごとのスカラーにより標識付き恒等関手に同型となる．
これで単射性を得る．

自然自己同型の各偶数線上の成分も $S^\times$ に属し，
隣接 Hom の忠実性から全て等しい．構造対象上で恒等なら全て一であり，
偶数線が生成族であることから圏全体で恒等である．
標準係数族の任意係数延長と固定模型の fpqc Hom 降下は
補題\ref{lem:abel-canonical-family-equivalences}により両立する．
従ってここで得た局所同型は指定した stack 化後にも正確に全ての射を記述する．
構造移送，Serre の反復，Gabriel 右随伴の各構成から，合成と係数変更への自然性も従う．
\end{proof}

\begin{corollary}[同型層の有限な方程式]
\label{cor:general-relative-isom-scheme}
前定理の下で，さらに $A,A'$ は標準 Hilbert 関数
$H_A(t)=(1-t^3)/(1-t)^5$ を持つとする．
$A_1\simeq(A')_1\simeq S^5$ を自明化する．
$V=S^5$ とし，$d=2,3$ に対して
\[
 I_d=\ker(V^{\otimes d}\to A_d),\qquad
 I'_d=\ker(V^{\otimes d}\to A'_d)
\]
と置く．これらはそれぞれ階数 $10,91$ の部分束である．
このとき式\eqref{eq:general-relative-pointed-hom}の層は
\begin{equation}\label{eq:general-finite-isom-equations}
 \left\{[g]\in\operatorname{PGL}_{5,S}:\
    g^{\otimes2}I_2=I'_2,\quad g^{\otimes3}I_3=I'_3\right\}
\end{equation}
で表される．従ってこれは $S$ 上アフィンかつ有限表示のスキームである．
この記述は非被約な係数環を含む．
\end{corollary}

\begin{proof}
$A$ は二次関係と一つの三次関係で表示されるので，
次数一の可逆写像が代数同型へ延びる必要十分条件は表示した二条件である．
$A_2,A_3$ は有限射影で，それぞれ階数 $15,34$ だから，$I_2,I_3$ は主張した部分束となる．
$\operatorname{GL}_5$ 上の二つの部分束の対応条件は Grassmannian への写像による閉条件であり，
次数スカラーの中心 $\mathbb G_m$ はこの条件に自明に作用する．
従って条件は $\operatorname{PGL}_5$ 上に降下し，閉部分スキームを定める．
局所的な包含条件は，両部分束の階数が等しく $g$ が可逆であることから等号を意味する．
この議論は体値の点だけを使わない．
$\operatorname{GL}_5\to\operatorname{PGL}_5$ は fpqc $\mathbb G_m$-torsor なので，
この閉部分スキームは正確に次数スカラーの層商を表す．
\end{proof}

これは二つの指定した族の間の射を有限な方程式で記述する結果である．
一般の Abel 圏の有効降下，モジュライ全体の代数性，有限型 atlas，
またはパラメータ族の開埋込みは，この同型層の表示だけからは従わない．


\begin{corollary}[一般十七次元の Abel 圏変動と二次包絡]
\label{cor:general17-abel-moduli-comparison}
$P$ を式\eqref{eq:general17-coefficient-ring}とする．
族 $A/P$ は $\operatorname{Spec}P\to\mathfrak M^{\mathcal O}_{\mathrm{cub}}$ を定める．
この族から任意の基底へ引き戻した二対象の間には，
定理\ref{thm:general-relative-pointed-hom}の全同値比較と
系\ref{cor:general-relative-isom-scheme}の有限表示が成り立つ．
全ての係数環上で compact 圏の相対 Serre 関手は $(-2)[3]$ である．

複素体値の一般の $\mathcal O$-marked Abel 圏同型類は十七次元分動く．
正確には，定理\ref{thm:17}の Rosenlicht 商の超越次数十七が，
全ての $\mathcal O$-保存 Abel 同値で割った同型類にもそのまま対応する．
また，そのような同値は二次包絡
\[
 B=T(A_1)/(\ker(A_1\otimes A_1\to A_2))
\]
を次数スカラーを除いた同型まで復元し，三次核の直線も保つ．
従って既存の十七係数から十三係数 ambient への射は，
この範囲で全ての構造対象付き Abel 同値と両立する．
\end{corollary}
\begin{proof}
相対模型の許容性は命題\ref{prop:general17-abel-admissible}である．
命題\ref{prop:general17-universal-pbw}と
\ref{prop:general17-relative-cy-koszul}は
第\ref{sec:relative-koszul-line-cohomology}節の四入力を全て供給し，
命題\ref{prop:general17-universal-saturation-center}は普遍中心条件を供給する．
従って相対 Serre 定理と全 Hom 比較を適用できる．
十七次元という結論は，複素体上で Abel 同値が次数付き同型と正確に一致することと，
定理\ref{thm:17}による．二次包絡は次数二の関係だけから作られるので，
代数同型に関して自然である．次数スカラーは包絡の次数スカラーへ移り，
三次核も元の同型により対応する．この構成は係数変更に両立する．
\end{proof}

この結論は，十七係数の正規形が大域的に一意であることを意味しない．
有限な残余同一視は保持する．また，$\mathcal O$ を忘れた圏，
裸の導来圏，モジュライ全体の成分性・代数性・有限型 atlas，
および一般十七係数族の dg smooth locus は別の問題である．
この節では有限中心性も全ファイバーの tails の正則性も使っていない．

\section{Jacobian 型 Abel 圏の代数的部分モジュライ}
\label{sec:jacobian-algebraic-chart}

本節の主結果は定理\ref{thm:jacobian-algebraic-abel-substack}である。
第\ref{sec:general17-relative-algebra}節の任意係数上の全同値比較を用い、
構造対象だけを標識した Abel 圏の問題の中に、
滑らかな十七次元 Deligne--Mumford stack を構成する。
構成は、十三係数の二次包絡の飽和開、普遍的な低次数中心、
中心三次直線の torsor、全 Abel 同値との比較、という順で進める。
とくに四次元のファイバーは、中心三次完全列の分裂を選ぶ自由度である。
商 stack・étale atlas・アフィン降下の一般論は既知の結果であり、
本節の対象はその具体的な Jacobian 族への適用と、全 Abel 同値との一致である。

\subsection{同変な飽和写像からの étale atlas}

以下の補題は、有限な同型層の方程式に加えて、実際の atlas を作るために用いる。
商 stack の代数性は既知の定理を使い、必要な étale 被覆をここで具体的に構成する。

\begin{lemma}[等次元な同変飽和写像]
\label{lem:equivariant-etale-saturation}
$k$ を標数零の体、$G$ を滑らかな幾何学的整な有限型アフィン群、
$T$ を滑らかな幾何学的整な有限型 $k$-スキーム、$Z$ を整な有限型 $k$-スキームとする。
$G$ は $Z$ に作用し、写像 $j:T\to Z$ による
\[
 a:G\times T\longrightarrow Z,\qquad(g,t)\longmapsto g\cdot j(t)
\]
が支配的で、$\dim Z=\dim G+\dim T$ と仮定する。
$a$ の étale locus を $E$、単位切断 $T\to G\times T$ によるその逆像を $T^e$ と置く。
このとき $T^e$ は非空開部分であり、
\[
 E=G\times T^e,\qquad Y=a(E)
\]
となる。$Y$ は非空な $G$-不変開部分で、$a:E\to Y$ は étale 全射である。
従って
\[
 T^e\longrightarrow[Y/G]
\]
は étale 全射 atlas であり、$[Y/G]$ は滑らかな $\dim T$ 次元の
既約 Deligne--Mumford stack である。
$Z$ が分離的なら、その対角はアフィンかつ有限表示である。
\end{lemma}

\begin{proof}
等次元性から $k(Z)\subset k(G\times T)$ は有限拡大であり、標数零なので分離的である。
源の一般点では局所環の写像はこの体の拡大そのものであり、平坦で、
相対微分加群は零である。射は局所有限表示なので、
étale 判定\cite[Lemma~29.37.15]{StacksAlgebraicChart}により一般点で étale となる。
従ってその étale locus $E$ は非空開部分である。

$G$ の左移動と $Z$ の作用は $a$ と可換する。
étale locus の平坦底変更との両立\cite[Lemma~29.37.17]{StacksAlgebraicChart}を $G\to\operatorname{Spec}k$ に適用し、
$G\times G\times T$ と $G\times Z$ 上の移動同型を用いれば、
$E$ は群スキームの作用について $G$-不変である。
この不変な開部分を単位切断へ引き戻してから再び移動すると
$E=G\times T^e$ を得る。
これは $k$-点だけについての不変性ではなく、開部分スキームの等式である。

étale 射は開写像\cite[Lemma~29.37.13]{StacksAlgebraicChart}なので、$Y=a(E)$ は開部分である。
同変性により $Y$ は $G$-不変であり、$E\to Y$ は étale 全射となる。
$G$ は滑らかなので、作用群oidの二射影は滑らかかつ有限表示であり、
商 $[Y/G]$ は代数的 stack である
\cite[Theorem~97.17.2]{StacksAlgebraicChart}。
$Y\to[Y/G]$ による $T^e\to[Y/G]$ の底変更は
\[
 G\times T^e\longrightarrow Y
\]
そのものである。従って $T^e\to[Y/G]$ は étale 全射である。
これで Deligne--Mumford 性を得る
\cite[Theorem~101.21.6]{StacksAlgebraicChart}。
$T^e$ が滑らかで次元 $\dim T$ であるため、商も同じ次元で滑らかであり、
$Y$ の既約性から商の既約性も従う。

最後に $Z$ が分離的なら $Y$ も分離的である。
商の対角を $Y\times Y$ へ引き戻すと
\[
 G\times Y\longrightarrow Y\times Y,\qquad(g,y)\longmapsto(y,gy)
\]
となる。この写像は $G\times Y\times Y$ 内の閉グラフを経由し、
$G$ がアフィン有限型なのでアフィンかつ有限表示である。
この性質を忠実平坦降下して対角の主張を得る。
対角の proper 性や商の分離性は、この議論からは主張しない。
\end{proof}

\begin{proposition}[Jacobian 周囲の具体的な étale atlas]
\label{prop:jacobian-ambient-etale-chart}
$G=\operatorname{PGL}(V)$、$V=\C^5$ とする。
式\eqref{eq:general17-coefficient-ring}から四つの $g_j$ を除いた基底を
\[
 T_B=\operatorname{Spec}
 \C[a,b,\Delta^{-1}][\alpha,\beta,(m_{ij})]
\]
とし、二次関係を記録する写像を $j_B:T_B\to Q_{113}$ と書く。
$Q_{113}$ は命題\ref{prop:quantum-113-component}の graph 像に被約構造を入れたものとする。
作用写像
\[
 a_B:G\times T_B\longrightarrow Q_{113}
\]
に前補題を適用できる。
従って $T_B$ の非空開 $T_B^e$ と $Q_{113}$ の非空 $G$-不変開 $Y$ が得られ、
\[
 G\times T_B^e\longrightarrow Y,\qquad
 T_B^e\longrightarrow\mathfrak B_{\mathrm J}:=[Y/G]
\]
はそれぞれ étale 全射および étale 全射 atlas である。
$\mathfrak B_{\mathrm J}$ は滑らかな既約十三次元 Deligne--Mumford stack である。
さらに $T_B^e$ は次の有理点 $t_0$ を含む：
\[
 (a,b,\alpha,\beta)=(3,5,2,3),\qquad
 (m_{ij})=
 \begin{pmatrix}1&1&2\\2&3&1\\1&4&5\end{pmatrix}.
\]
\end{proposition}

\begin{proof}
分解 $V^{\otimes2}=\wedge^2V\oplus\operatorname{Sym}^2V$ は線型変換と両立するので、
その graph chart と $Q_{113}$ は $G$-不変である。
次数二の関係部分束は命題\ref{prop:general17-universal-pbw}により任意の底変更に両立するため、
$j_B$ はスキームの射である。
定理\ref{thm:17}の十三係数横断性と、$Q_{113}$ の三十七次元性から、
$a_B$ は支配的である。
$\dim G+\dim T_B=24+13=37$ なので、前補題を適用できる。
この適用では $Q_{113}$ の graph 表示と次元を用いるだけであり、
Pym の形式的・解析的な成分比較を atlas の存在の入力にはしていない。

指定した点での非空性も直接検証できる。
$t_0$ の式は \eqref{eq:rat113-test-F} の $F_*$ であり、
$\Delta$ の七因子の値は
\[
2562,\quad-378,\quad-237100440304,\quad316,\quad5489,\quad-5805,\quad31963501
\]
となる。従って $t_0\in T_B$ である。

二次 bivector の百五十係数を座標とする。
$\operatorname{GL}_5$ の二十五方向と十三係数 slice の微分を合わせた行列を $D$ とすると、
$D$ は $150\times38$ の整数行列である。
$E_{ab}$ が変数 $x_b$ を $x_a$ へ送る規約で、軌道方向の成分は
\[
 (D_{ab}\Pi)_{ij}
 =x_a\partial_b\Pi_{ij}
  -\delta_{ib}\Pi_{aj}-\delta_{jb}\Pi_{ia}
\]
となる。五つの対角方向の和は整数上で零である。
\nolinkurl{computations/verify_jacobian_etale_point.py}は、
\nolinkurl{computations/jacobian_etale_point_certificate.json}で指定した
三十七次小行列式について
\[
 \det D_{37}\equiv127524\pmod{1000003}
\]
を与える。従って $\operatorname{rank}_{\mathbb Q}D=37$ であり、
核はちょうど上記の同時スカラー一方向である。

全 Jacobi・発散方程式の微分を $M$ とすると、同じコードは
$M D=0$ を整数上で確認し、既存の $113$ 次小行列式
\eqref{eq:rat113-minor-certificates}も再現する。
従って $\dim\ker M=150-113=37$ である。
$Q_{113}$ の次元は三十七で、その接空間は $\ker M$ に含まれるので、
$j_B(t_0)$ は $Q_{113}$ の滑らかな点である。
$\operatorname{PGL}_5$ へ移るとスカラー核が除かれ、
$(1,t_0)$ における $d a_B$ は接空間の同型となる。
源・標的がこの点で滑らかかつ同次元であるため、Jacobian 判定により
$a_B$ は $(1,t_0)$ で étale である。
従って $t_0\in T_B^e$ を得る。
ここで有限体を用いたのは整数行列式の非零判定だけであり、
積零とスカラー核は整数上の恒等式として確認している。
\end{proof}

\subsection{周囲の普遍的な低次数中心}

三次式の四つの定数係数を除いた基底を
\[
 P_B=\C[a,b,\Delta^{-1}]
       [\alpha,\beta,(m_{ij})_{1\leq i\leq3,\,0\leq j\leq2}],
 \qquad T_B=\operatorname{Spec}P_B
\]
と書く。$B$ の関係は $g(u,v)$ に依存しない。
命題\ref{prop:general17-universal-pbw}で選んだ中心三次元を
$\mathcal G\in B_3$ とし、$U=P_Bu\oplus P_Bv$ と置く。

\begin{proposition}[中心一次・三次束の普遍性]
\label{prop:general13-universal-centers}
任意の可換 $P_B$-代数 $S$ に対し、
\begin{equation}\label{eq:general13-universal-centers}
 Z(B_S)_1=Su\oplus Sv,
 \qquad
 Z(B_S)_3=S\mathcal G\oplus\operatorname{Sym}^3_S(Su\oplus Sv)
\end{equation}
である。これらは $(B_S)_1,(B_S)_3$ の直和因子であり、
形成は任意の係数変更と可換する。
\end{proposition}

\begin{proof}
まず複素閉点上で計算する。
補題\ref{lem:general17-relative-elliptic-boundary}により
楕円境界の $\sigma,\sigma^2,\sigma^3$ は恒等でない。
従って twisted homogeneous coordinate algebra $H$ の
次数一、二、三の中心は零である。
完全列\eqref{eq:general17-elliptic-quotient}より
\[
 Z(\mathscr S)_1=Z(\mathscr S)_2=0,
 \qquad Z(\mathscr S)_3=k\Omega.
\]
最後の等式では、中心元の $H$ への像が中心であることと、
$\ker(\mathscr S_3\to H_3)=k\Omega$ を用いた。

$B$ の active 濾過の associated graded は
$\mathscr S\otimes k[u,v]$ である。
次数一の中心元の active 次数が正なら、その最高記号は
$Z(\mathscr S)_1$ に属するので零となり、矛盾する。
従って $Z(B)_1=ku\oplus kv$ である。
次数三の中心元では、active 次数三の最高記号は $k\Omega$ に属する。
対応する $\mathcal G$ のスカラー倍を引けば active 次数は二以下となる。
次いで $Z(\mathscr S)_2=Z(\mathscr S)_1=0$ により
active 次数を順に零まで下げられるので、残りは $k[u,v]_3$ に属する。
この議論で必要なのは次数三までの中心であり、全中心の記述ではない。

次に任意係数への普遍性を示す。
PBW 分裂と
$\mathscr S_3/P_0\Omega\simeq H_3$ の有限射影性により、
\[
 E_1=U\subset B_1,
 \qquad E_3=P_B\mathcal G\oplus\operatorname{Sym}^3U\subset B_3
\]
は有限局所自由な直和因子である。
実際、$E_3$ の active 零部分は $P_B[u,v]_3$ 全体であり、
active 三部分では直和因子 $P_B\Omega$ を持つ。
$B_3/E_3$ の残る濾過商は有限射影なので、この商自体も有限射影である。

$E_d$ は中心なので、$d=1,3$ について交換子写像は
\[
 \overline\delta_d:B_d/E_d\longrightarrow B_{d+1}^{\oplus3},
 \qquad \overline c\longmapsto([c,x],[c,y],[c,z])
\]
を誘導する。$u,v$ は中心なので、三つの交換子で十分である。
この写像は前段の計算により全複素閉点で単射である。
$P_B$ は有限型 $\C$-代数の局所化であり Jacobson なので、
各点の近傍で最大階数の小行列式が可逆となる。
従って $\overline\delta_d$ は局所的に分裂単射である。
この単射性は任意の、非平坦な係数変更後にも保たれ、
\eqref{eq:general13-universal-centers}を得る。
\end{proof}

命題\ref{prop:jacobian-ambient-etale-chart}で構成した
$G=\operatorname{PGL}(V)$-不変開集合 $Y$ と、その étale 全射
\[
 G\times T_B^e\longrightarrow Y
\]
を用いる。$Y$ 上の普遍関係部分束を $R\subset V^{\otimes2}\otimes\mathcal O_Y$ とし、
$\mathcal B=T(V\otimes\mathcal O_Y)/(R)$ と置く。
$G$ の元を $\operatorname{GL}(V)$ へ持ち上げる fpqc 被覆上では、
$\mathcal B$ は上の $B$ の係数変更と同型である。
従って各 $\mathcal B_d$ は有限局所自由であり、
命題\ref{prop:general13-universal-centers}と忠実平坦降下により
\[
 \mathcal U=Z(\mathcal B)_1,
 \qquad \mathcal E=Z(\mathcal B)_3,
 \qquad \mathcal W=\operatorname{Sym}^3\mathcal U\subset\mathcal E
\]
はそれぞれ階数二、五、四の部分束となる。
これらの中心の形成は $Y$ 上の任意の基底変更と可換する。
ここでは核の普遍性を fpqc 被覆上の分裂単射により確認しており、
幾何ファイバーの次元だけから非被約な底への結論を出していない。
商 $\mathcal L=\mathcal E/\mathcal W$ は可逆層である。

\subsection{中心三次直線のアフィン四次元 torsor}

\begin{lemma}[射影超平面の補集合と分裂の torsor]
\label{lem:central-cubic-splitting-torsor}
任意のスキーム $Y$ 上の有限局所自由層の完全列
\[
 0\longrightarrow\mathcal W\longrightarrow\mathcal E
   \xrightarrow{q}\mathcal L\longrightarrow0
\]
を考え、$\mathcal L$ の階数を一とする。
$\PP_Y(\mathcal E)$ は一次元部分束を表す規約とし、
\[
 J=\PP_Y(\mathcal E)\setminus\PP_Y(\mathcal W),
 \qquad \mathcal V_4=\mathcal Hom_Y(\mathcal L,\mathcal W)
\]
と置く。$\mathcal W$ の階数が四なら、$J$ は
階数四のベクトル束 $\mathcal V_4$ の加法群の torsor である。
この同定は任意の底変更に自然であり、
$J\to Y$ はアフィン、滑らか、全射で、相対次元四である。
\end{lemma}

\begin{proof}
$Y$-スキーム $S$ に対して $J(S)$ は、
直線部分束 $N\subset\mathcal E_S$ であって、
$q|_N:N\to\mathcal L_S$ が各ファイバーで非零となるものからなる。
可逆層間のこの写像は同型である。
従って $N$ は一意な右逆射
$s:\mathcal L_S\to\mathcal E_S$、$q\circ s=1$ の像となる。
逆に右逆射の像は直線直和因子であり、$J(S)$ の点を与える。
よって $J$ は $q$ の分裂の層を表す。

二つの分裂の差は一意に
$\mathcal Hom_S(\mathcal L_S,\mathcal W_S)$ の切断であり、
$\mathcal V_4$ は分裂の層に自由かつ推移的に作用する。
局所自由層の列は Zariski 局所で分裂するので、
$J$ は局所的に $\mathbb A^4_Y$ に同型である。
特に主張した射の性質が従う。
分裂を一つ選ばない限り、$J$ 自体を零切断付きベクトル束とは同一視しない。
\end{proof}

上で得た中心束の列にこの補題を適用する。
$T_B^e$ 上では Rains の規格化された中心元 $\mathcal G$ が分裂を与えるので、
\begin{equation}\label{eq:jacobian-cubic-chart-pullback}
 J\times_Y T_B^e\simeq T_B^e\times\mathbb A^4,
 \qquad
 (t,g_0,g_1,g_2,g_3)\longmapsto
 [\mathcal G_t+\textstyle\sum_{j=0}^3g_j u^{3-j}v^j].
\end{equation}
これは体値の点の対応ではなく、分裂の関手の同型である。
特に任意の非被約基底も含む。
$J$ 上の普遍直線部分束で $\mathcal B$ を割って得る三次商は、
この被覆上で既存の十七係数族と一致する。
その次数ごとの平坦性と許容性は
命題\ref{prop:general17-universal-pbw}および
命題\ref{prop:general17-abel-admissible}により供給される。

\begin{proposition}[代数的周囲への相対アフィン四次元射]
\label{prop:jacobian-affine-four-morphism}
上の $Y$、$J$ と $G=\operatorname{PGL}(V)$ に対して、
\[
 \mathfrak B_{\mathrm J}=[Y/G],\qquad
 \mathfrak C_{\mathrm J}=[J/G]
\]
と置く。自然な射
\begin{equation}\label{eq:jacobian-affine-four-stack-map}
 \mathfrak C_{\mathrm J}\longrightarrow\mathfrak B_{\mathrm J}
\end{equation}
は階数四のベクトル束の加法群の torsor であり、
スキームで表現可能、アフィン、滑らか、全射で、相対次元四である。
$T_B^e\times\mathbb A^4\to\mathfrak C_{\mathrm J}$ は étale 全射 atlas であり、
$\mathfrak C_{\mathrm J}$ は滑らかな既約十七次元の有限型
Deligne--Mumford stack である。
\end{proposition}

\begin{proof}
まず線型化の重みを確認する。
$\mathcal B_d$ には自然な $\operatorname{GL}(V)$-線型化があり、
中心スカラー $\lambda I$ は $\lambda^d$ として作用する。
従って $\mathcal U,\mathcal W,\mathcal E,\mathcal L$ の中心重みは
それぞれ $1,3,3,3$ である。
$\mathcal U$ や $\mathcal E$ 自体を
$[Y/\operatorname{PGL}(V)]$ 上の通常のベクトル束として降下させてはいない。

一方、$\mathcal V_4=\mathcal Hom(\mathcal L,\mathcal W)$ の中心重みは零である。
$\operatorname{GL}(V)\to G$ が fpqc $\mathbb G_m$-torsor であることから、
その作用写像と線型化は $G$ へ降下する。
$J$ の点は $\mathcal E$ 内の直線なので、中心スカラーは $J$ にも自明に作用し、
$J$ には $G$-作用がある。
前補題の加法作用と torsor 同型は $\operatorname{GL}(V)$-同変であるから
$G$-同変となる。

従って $\mathcal V_4$ は $\mathfrak B_{\mathrm J}$ 上の階数四のベクトル束
$\overline{\mathcal V}_4$ を定め、$[J/G]$ はその加法群の torsor となる。
より具体的に、$S\to\mathfrak B_{\mathrm J}$ を
$G$-torsor $P\to S$ と $G$-同変写像 $P\to Y$ で表示する。
その上の2-ファイバー積は
$P\times_YJ$ の $G$-同変な降下である。
$P\times_YJ\to P$ はアフィンなので、fpqc アフィン降下によって
$S$ 上のアフィンスキームを与える。
同時に $\mathcal V_4$ とその作用も降下し、
このアフィンスキームが $\overline{\mathcal V}_{4,S}$ の torsor であると分かる。
これで表現可能性、アフィン性、滑らかさ、全射性、相対次元を得る。
商 stack の torsor 表示とアフィン降下については
\cite[Section~95.15 and Lemma~35.37.1]{StacksAlgebraicChart}を参照する。

$T_B^e\to\mathfrak B_{\mathrm J}$ の底変更を取ると
\eqref{eq:jacobian-cubic-chart-pullback}になる。
前者は命題\ref{prop:jacobian-ambient-etale-chart}の étale 全射 atlas なので、
$T_B^e\times\mathbb A^4\to\mathfrak C_{\mathrm J}$ も étale 全射である。
これは滑らかな十七次元スキームからの atlas であり、
$\mathfrak C_{\mathrm J}$ の smooth 性、Deligne--Mumford 性、次元十七が従う。
$T_B^e\times\mathbb A^4$ は既約な有限型スキームであり、この atlas は全射なので、
商の既約性と有限型性も従う。
\end{proof}

\begin{remark}[ファイバーと標識の範囲]
周囲 $B$ とその同定を固定した幾何学的2-ファイバーは
$J_B$ であり、中心三次元を一つ選べば $\mathbb A^4$ となる。
周囲の同型類だけを固定して同定を忘れた同型類の集合には、
有限群 $\operatorname{Aut}_{\mathrm{gr}}(B)/\mathbb G_m$ の作用が残り得る。
この有限群が自明であるとは主張しない。

$\mathfrak B_{\mathrm J}$ は周囲の\emph{次数付き}代数を射影的な生成元変換で比較する stack である。
周囲の全ての $\mathcal O$-保存 Abel 圏同値を分類する、という追加主張は含まない。
また、この命題の smooth 性はパラメータ stack の smooth 性であり、
三次商の各 compact 導来圏の dg smooth 性とは異なる。
\end{remark}

\subsection{有限な関係窓による三次族の表示}

\begin{lemma}[中心三次直線の関係窓への埋込み]
\label{lem:jacobian-relation-window-immersion}
命題\ref{prop:jacobian-affine-four-morphism}の $J\to Y$ に対して、
商 $A$ の二次・三次関係を取る写像
\[
 \iota:J\longrightarrow
 X:=\operatorname{Gr}(10,V^{\otimes2})
       \times\operatorname{Gr}(91,V^{\otimes3})
\]
は $G$-同変な局所閉埋込みである。
ここで Grassmannian は指定した階数の部分束を表す規約である。
この記述は任意の係数環上の点について成立する。
\end{lemma}

\begin{proof}
$Y$ 上の二次関係部分束を $R$、その二次代数を $\mathcal B$ とする。
$\mathcal B_3$ は階数三十五の局所自由層なので、
\[
 K_3=\ker(V^{\otimes3}\otimes\mathcal O_Y\to\mathcal B_3)
\]
は階数九十の部分束である。
これは $V\otimes R+R\otimes V$ の像である。
$\mathcal B_3$ 内の直線部分束 $N$ を取ることは、
$K_3$ を含む階数九十一の部分束 $I_3\subset V^{\otimes3}\otimes\mathcal O_Y$ を取ることと同値である。
この対応は $I_3$ を $N$ の逆像、$N$ を $I_3/K_3$ とするものであり、
任意の基底変更に両立する。
従って $\PP_Y(\mathcal B_3)$ は、第二の Grassmannian の $Y$ への引戻しの中で、
$K_3$ を含むという閉条件により表示される。

中心束 $\mathcal E\subset\mathcal B_3$ は部分束なので、
$\PP_Y(\mathcal E)$ は $\PP_Y(\mathcal B_3)$ の閉部分である。
$J$ はそこから $\PP_Y(\mathcal W)$ を除いた開部分である。
$Y$ 自体も第一の Grassmannian の局所閉部分であるから、主張の局所閉埋込みを得る。
全ての操作は線型変換と両立し、中心スカラーは部分束に自明に作用するため、
$\iota$ は $G$-同変である。
\end{proof}

\subsection{代数的な部分モジュライと全 Abel 同値}

\begin{theorem}[Jacobian 型の代数的 Abel 圏部分モジュライ]
\label{thm:jacobian-algebraic-abel-substack}
命題\ref{prop:jacobian-affine-four-morphism}の滑らかな既約十七次元の有限型
Deligne--Mumford stack $\mathfrak C_{\mathrm J}=[J/G]$ から、自然な射
\begin{equation}\label{eq:jacobian-abel-fully-faithful}
 \mathfrak C_{\mathrm J}\longrightarrow
 \mathfrak M^{\mathcal O}_{\mathrm{cub}}
\end{equation}
がある。この射は任意の係数環上で完全忠実である。
その本質像は、fpqc 局所で
$T_B^e\times\mathbb A^4$ 上の十七係数族の引戻しに同値となる対象からなる全部分 stack である。
特にこの全部分 stack は、有限型の代数的 Deligne--Mumford stack として表される。
その対角はアフィン、有限表示かつ非分岐である。

二次包絡への射は、この同定の下で
\[
 \mathfrak C_{\mathrm J}\longrightarrow\mathfrak B_{\mathrm J}
\]
という相対アフィン四次元 torsor に一致する。
この主張は周囲を次数スカラーを除いた次数付き同型で比較するものである。
\end{theorem}

\begin{proof}
$J$ 上には普遍中心直線部分束
$\mathcal N\subset\mathcal B_3|_J$ があり、
\[
 \mathcal A=T(V\otimes\mathcal O_J)/(R,\widetilde{\mathcal N})
\]
を得る。ここで $\widetilde{\mathcal N}$ は三次語の加群における逆像であり、
表示は $\mathcal B/(\mathcal N)$ と同じ意味である。
$G\times T_B^e\to Y$ と、さらに
$\operatorname{GL}(V)\to G$ による fpqc 被覆の上では、
この代数は既存の十七係数族の引戻しに次数付き同型となる。
従って命題\ref{prop:general17-universal-pbw}と
\ref{prop:general17-abel-admissible}による平坦性と許容性を持ち、
$J\to\mathfrak M^{\mathcal O}_{\mathrm{cub}}$ が定まる。
普遍的な $G$ の元を大域的に $\operatorname{GL}(V)$ へ持ち上げたとは仮定していない。

この射の $G$-同変性を構成する。
$g\in G(S)$ は fpqc 局所で $\widetilde g\in\operatorname{GL}(V)(S)$ に持ち上がる。
その持上げは代数の構造移送により、構造対象の同定を備える完全 Abel 同値を与える。
二つの持上げは次数スカラー $\varepsilon_\lambda$ だけ異なる。
右次数付き加群に対する構造移送を
$F_{\varepsilon_\lambda}(M)=M\otimes_{A,\varepsilon_\lambda}A$ と書くと、
\[
 m_d\otimes b\longmapsto\lambda^d m_db
\]
は $F_{\varepsilon_\lambda}(M)\to M$ の自然な次数付き同型である。
この写像は $M=A$ のとき標準の構造対象の同定に一致し、
$\lambda\mu$ に関する合成も保つ。
従って持上げの差は標識付き自然同型で消える。

定理\ref{thm:general-relative-pointed-hom}により、構造対象上で恒等な自然自己同型は自明である。
従ってこれらの標識付き自然同型は一意であり、三重・四重交叉上の整合性も従う。
同型類を射とする現在の $1$-stack の規約では、
$G\times J$ 上の同変同型は合成に関する cocycle 条件を満たす。
商 stack の torsor 表示と stack 化の普遍性
\cite[Section~95.15]{StacksAlgebraicChart}から
\eqref{eq:jacobian-abel-fully-faithful}を得る。
この操作は $G$-torsor が大域的に $\operatorname{GL}_5$-torsorへ持ち上がることや、
降下後の対象が大域的な次数付き代数で表示されることを要求しない。

完全忠実性を示す。
$S$ 上の二点 $x,x':S\to J$ に対し、商 stack の同型層は transporter
\[
 \{g\in G_S:g\cdot x=x'\}
\]
である。前補題の関係窓の復元により、これは
$\underline{\operatorname{Isom}}_{\mathrm{gr}}(A_x,A_{x'})/\mathbb G_m$ に等しい。
この同定は三次直線を関係部分束の商として復元するので、非被約な $S$ にも成立する。
一方、両対象を同時に fpqc 局所化して十七係数模型へ戻せば、
定理\ref{thm:general-relative-pointed-hom}により、この層は全ての
$\mathcal O$-保存 Abel 圏同値の層と一致する。
固定模型の Hom 降下も同定されているため、局所的な一致は元の $S$ 上へ降下する。
任意の $\mathfrak C_{\mathrm J}$ の二対象も fpqc 局所で $J$ 上の対象に戻るので、
全ての対象間の完全忠実性が従う。

本質像の記述では、$T_B^e\times\mathbb A^4$ が
$\mathfrak C_{\mathrm J}$ の étale 全射 atlas であることを用いる。
この atlas の模型へ fpqc 局所的に同値な対象には、
完全忠実性により一意な同型降下データが持ち上がる。
$\mathfrak C_{\mathrm J}$ 自体が stack であるため、このデータは有効である。
よって記述した本質像は全部分 stack となる。

対角のアフィン性と有限表示性は、$J$ が分離有限型であることと
$G$ のアフィン性から、補題\ref{lem:equivariant-etale-saturation}の対角の議論と同様に得る。
非分岐性は Deligne--Mumford 性による。
最後に二次包絡は二次関係部分束だけから復元され、三次核は $I_3/K_3$ である。
従って全同値に自然な二次包絡射は、構成した $[J/G]\to[Y/G]$ と一致する。
\end{proof}

\begin{remark}[大域降下と主張の境界]
この定理は、指定した全部分問題について有限型の atlas を構成する。
全 $\mathfrak M^{\mathcal O}_{\mathrm{cub}}$ の代数性、
その中での開埋込み・局所閉埋込み・成分性は結論していない。
像の条件は\emph{族が fpqc 局所で模型へ戻ること}であり、
各幾何ファイバーの同型類だけについての条件ではない。
一方、補題\ref{lem:jacobian-relation-window-immersion}により、
有限な関係窓の商 $[X/G]$ の中では局所閉部分 stack として表示できる。

$\mathcal O$ は標識付き同値により降下するが、
$\mathcal O(1)$ や次数一の shift が大域的に降下するとは主張しない。
相対 Serre 関手の $(-2)[3]$ という表示は指定した étale atlas 上のものであり、
中心重みや基底線束の問題を無視して全商上へ同じ表示を付けない。
また、パラメータ stack の smooth 性は、各三次圏の dg smooth 性と別である。
幾何学的 stabilizer は有限であるが、自明性・商の分離性・対角の有限性は未証明である。
\end{remark}

\section{相対Kuznetsov成分の構成と降下}
\label{sec:relative-kuznetsov-construction}

この節では第\ref{sec:relative-koszul-line-cohomology}節の
相対入力を満たす三次商 $A=B/(w)$ を扱う。
基底 $P$ 上の dg smoothness は仮定しない。
任意の可換 $P$-代数 $S$ に対し
\[
 \mathcal C_S=D(\operatorname{QGr}(A_S))^c,
 \qquad L_{i,S}=\mathcal O_{A_S}(i)
\]
と置く。以下の圏と関手は、非有界導来圏の標準 dg 増強で取る。
従って評価写像の cone は dg 関手として構成し、
単なる三角圏上での非関手的な cone の選択には依存しない。

\subsection{相対例外対と有限な射影関手}

\begin{theorem}[任意係数上の相対残余圏]
\label{thm:relative-kuznetsov-projector}
上の入力の下で、次が成り立つ。
\begin{enumerate}
\item $(L_{0,S},L_{1,S})$ は $S$-線型相対例外対であり、
\[
 \mathbf R\operatorname{Hom}(L_{i,S},L_{i,S})\simeq S,
 \quad\mathbf R\operatorname{Hom}(L_{1,S},L_{0,S})=0,
 \quad\mathbf R\operatorname{Hom}(L_{0,S},L_{1,S})\simeq (A_S)_1.
\]
最後の加群は階数五の有限射影 $S$-加群である。
\item 相対 dg 部分圏
\[
 \mathcal K_S=\{M\in\mathcal C_S\mid
  \mathbf R\operatorname{Hom}(L_{0,S},M)=
  \mathbf R\operatorname{Hom}(L_{1,S},M)=0\}
\]
は admissible であり、半直交分解
\begin{equation}\label{eq:relative-kuznetsov-sod}
 \mathcal C_S=\langle\mathcal K_S,L_{0,S},L_{1,S}\rangle
\end{equation}
がある。
\item $\mathsf L_L(M)=\operatorname{Cone}
 (\mathbf R\operatorname{Hom}(L,M)\otimes_S^{\mathbf L}L\to M)$
と置くと、
\begin{equation}\label{eq:relative-kuznetsov-left-projector}
 \mathsf P_S=\mathsf L_{L_{0,S}}\circ\mathsf L_{L_{1,S}}
\end{equation}
は包含 $\mathcal K_S\hookrightarrow\mathcal C_S$ の左随伴である。
特に右側の $L_1$ に沿う mutation を先に行う。
\item
\begin{equation}\label{eq:relative-kuznetsov-generator}
 G_S=\mathsf P_S\bigl(L_{2,S}\oplus L_{3,S}\oplus L_{4,S}\bigr)
\end{equation}
は $\mathcal K_S$ の classical generator であり、
\[
 E_{\mathcal K,S}=\mathbf R\operatorname{Hom}(G_S,G_S),
 \qquad\mathcal K_S\simeq\operatorname{Perf}(E_{\mathcal K,S})
\]
を得る。$E_{\mathcal K,S}$ は $S$-perfect である。
\item 任意の係数環の射 $S\to S'$ について、射影関手と上記 generator は
導来係数延長と両立し、標準比較
\[
 E_{\mathcal K,S}\otimes_S^{\mathbf L}S'
 \xrightarrow{\ \sim\ }E_{\mathcal K,S'}
\]
は dg 代数の quasi-isomorphism である。
従って $\mathcal K_S$ は arbitrary derived base change と両立する
実際の相対 proper dg 圏を与える。
\end{enumerate}
\end{theorem}

\begin{proof}
例外対の表示は定理\ref{thm:relative-cubic-universal-line-cohomology}
で差 $j-i=0,-1,1$ を取ったものである。
係数 perfect 性は系\ref{cor:relative-cubic-compact-proper}による。
従って $\mathsf L_L$ の定義に現れる評価写像とその cone は compact 対象の範囲にある。

まず相対例外対象 $L$ に対し、評価三角へ
$\mathbf R\operatorname{Hom}(L,-)$ を適用すると
$\mathbf R\operatorname{Hom}(L,\mathsf L_L(M))=0$ となる。
実際、最初の射は例外性により恒等写像である。
$\mathsf L_{L_1}$ の後で $\mathsf L_{L_0}$ を適用しても、
$\mathbf R\operatorname{Hom}(L_1,L_0)=0$ により前者の直交性は保たれる。
従って $\mathsf P_S(M)\in\mathcal K_S$ である。
また $N\in\mathcal K_S$ なら両評価項が零なので $\mathsf P_S(N)\simeq N$ である。
各評価三角に $\mathbf R\operatorname{Hom}(-,N)$ を適用して
\[
 \mathbf R\operatorname{Hom}(\mathsf P_S(M),N)
 \simeq\mathbf R\operatorname{Hom}(M,N)
\]
を得る。これが左随伴性である。
二つの評価三角を貼り合わせれば、$M\to\mathsf P_S(M)$ の fiber は
$\langle L_0,L_1\rangle$ に属する。
これは半直交性と合わせて分解\eqref{eq:relative-kuznetsov-sod}を与える。

右随伴を具体的に構成する。
定理\ref{thm:koszul-cubic-relative-serre}の相対 Serre 自己同値を
$\mathsf S_{\mathcal C}=(-2)[3]$ と置き、$F_i=\mathsf S_{\mathcal C}(L_i)$ とする。
係数 perfect 双対性から
\[
 \mathcal K_S={}^{\perp}\langle F_0,F_1\rangle,
 \qquad\mathbf R\operatorname{Hom}(F_1,F_0)=0.
\]
相対例外対象 $F$ に対し
\[
 \mathsf R_F(M)=\operatorname{Cone}\bigl(
 M\longrightarrow\mathbf R\operatorname{Hom}(M,F)^\vee
                      \otimes_S^{\mathbf L}F\bigr)[-1]
\]
と置く。ここで $(-)^\vee=\mathbf R\operatorname{Hom}_S(-,S)$ であり、
射は完全双対性による余評価である。
$\mathbf R\operatorname{Hom}(-,F)$ を適用すれば
$\mathbf R\operatorname{Hom}(\mathsf R_F(M),F)=0$ となる。
さらに $\mathbf R\operatorname{Hom}(F_1,F_0)=0$ により
\[
 \mathsf Q_S=\mathsf R_{F_1}\circ\mathsf R_{F_0}
\]
は $\mathcal K_S$ に値を取り、その上で恒等である。
$N\in\mathcal K_S$ に対し二つの余評価三角へ
$\mathbf R\operatorname{Hom}(N,-)$ を適用すると
\[
 \mathbf R\operatorname{Hom}(N,\mathsf Q_S(M))
 \simeq\mathbf R\operatorname{Hom}(N,M)
\]
を得る。従って $\mathsf Q_S$ は包含の右随伴であり、admissibility が従う。

系\ref{cor:relative-cubic-compact-proper}により
$\mathcal C_S=\operatorname{thick}(\bigoplus_{i=0}^4L_i)$ である。
exact な射影関手 $\mathsf P_S$ は有限 cone、shift、直和因子を保ち、
$\mathsf P_S(L_0)=\mathsf P_S(L_1)=0$ である。
任意の $N\in\mathcal K_S$ は $N\simeq\mathsf P_S(N)$ なので
$\mathcal K_S=\operatorname{thick}(G_S)$ を得る。
dg Morita 同値と $E_{\mathcal K,S}$ の係数 perfect 性も従う。

最後に全ての評価・余評価写像は自然であり、
それらの係数 Hom 複体は perfect である。
系\ref{cor:relative-cubic-compact-proper}の任意底変更比較を使うと、
テンソル積、双対、有限 cone は全て導来係数延長と両立する。
従って $\mathsf P_S(M)$ と $G_S$ の底変更を得る。
同じ比較を $G_S$ の自己 Hom に適用すれば、積と両立する
$E_{\mathcal K,S}$ の標準底変更同型を得る。
$G_{S'}$ は前段の議論によって $\mathcal K_{S'}$ の generator であるから、
これは圏全体の Morita 基底変更比較でもある。
\end{proof}

\begin{corollary}[残余圏の相対 Serre 自己同値]
\label{cor:relative-kuznetsov-serre}
$i:\mathcal K_S\hookrightarrow\mathcal C_S$ を包含とする。
相対 Serre 自己同値とその逆は
\[
 \mathsf S_{\mathcal K}=i^!\mathsf S_{\mathcal C}i,
 \qquad \mathsf S_{\mathcal K}^{-1}=i^*\mathsf S_{\mathcal C}^{-1}i
\]
で与えられる。trace、pairing とこれらの関手は任意の導来係数延長と両立する。
\end{corollary}
\begin{proof}
随伴と $\mathcal C_S$ の相対 Serre 双対により
$\mathbf R\operatorname{Hom}_{\mathcal K}(M,N)^\vee
 \simeq\mathbf R\operatorname{Hom}_{\mathcal K}(N,i^!\mathsf S_{\mathcal C}iM)$
を得る。逆 Serre に対して同じ計算を行うと第二の関手を得る。
両者の合成は Yoneda により恒等関手となる。
随伴を上の有限 mutation で構成したため、その底変更も成立する。
\end{proof}

\begin{remark}[可換ファイバー]
$A_S$ が滑らかな相対三次超曲面 $X_S\subset\mathbb P^4_S$ の斉次座標環である場合、
上の構成は通常の
$\operatorname{Ku}(X_S)=\langle\mathcal O_{X_S},\mathcal O_{X_S}(1)\rangle^\perp$
そのものである。従ってこの構成を非可換族の相対 Kuznetsov 成分と呼べる。
非可換一般点における幾何性は、第\ref{sec:kuznetsov-moduli-morphism}節で別に証明する。
\end{remark}

\begin{lemma}[残余生成元の非消滅を与える Euler 数]
\label{lem:relative-kuznetsov-nonzero-euler}
任意の体値ファイバーにおいて、$G_i=\mathsf P(\mathcal O(i))$ と置くと
\[
 [G_2]=[\mathcal O(2)]-5[\mathcal O(1)]+10[\mathcal O],
 \qquad\chi(G_2,G_2)=-9.
\]
特に $G_2$ は全ての体値ファイバーで非零である。
生成元 $G=G_2\oplus G_3\oplus G_4$ 自身については
$\chi(G,G)=-2925$ である。
\end{lemma}
\begin{proof}
非負 twist の Euler 数を
$h_n=\binom{n+4}{4}-\binom{n+1}{4}$（$n\ge0$、$0\le n<3$ では第二項を零とする）
で書くと、$h_0,h_1,h_2,h_3,h_4=1,5,15,34,65$ である。
負の twist は同じ Hilbert 関数の零延長で扱わず、線コホモロジーの双対性から計算する。
二つの評価三角から
\[
 [G_j]=[\mathcal O(j)]-h_{j-1}[\mathcal O(1)]
           +(5h_{j-1}-h_j)[\mathcal O],\qquad j=2,3,4
\]
となる。$\chi(\mathcal O,G_2)=\chi(\mathcal O(1),G_2)=0$ より
\[
 \chi(G_2,G_2)=\chi(\mathcal O(2),G_2)
   =1-5\cdot0+10\cdot(-1)=-9.
\]
最後の $-1$ は線コホモロジーの
$\mathbf R\operatorname{Hom}(\mathcal O(2),\mathcal O)\simeq k[-3]$ による。
同じ計算で
\[
 \bigl(\chi(G_i,G_j)\bigr)_{i,j=2}^4
 =\begin{pmatrix}-9&-36&-90\\-45&-189&-486\\-126&-540&-1404\end{pmatrix}.
\]
全成分の和は $-2925$ である。
これらは普遍的線コホモロジーと評価三角からの厳密な整数計算であり、
特定ファイバーの数値的近似ではない。
\end{proof}

\subsection{一次 twist の対象降下を必要としない圏の降下}

\begin{lemma}[基底上の可逆加群に対する不変性]
\label{lem:kuznetsov-line-twist-invariance}
$\ell$ を可逆 $S$-加群とする。相対例外対象 $L$ と $L\otimes_S\ell$ は
同じ $S$-線型 thick 部分圏を生成し、同じ直交圏を定める。
さらに評価関手と left mutation には標準同型
\[
 \mathbf R\operatorname{Hom}(L\otimes_S\ell,-)
             \otimes_S^{\mathbf L}(L\otimes_S\ell)
 \simeq \mathbf R\operatorname{Hom}(L,-)\otimes_S^{\mathbf L}L,
 \qquad\mathsf L_{L\otimes\ell}\simeq\mathsf L_L
\]
がある。この比較は可逆加群のテンソル積に関して coherent である。
\end{lemma}
\begin{proof}
最初の比較では左辺を
$\mathbf R\operatorname{Hom}(L,-)\otimes_S\ell^\vee\otimes_S L\otimes_S\ell$
と書き、評価 $\ell^\vee\otimes_S\ell\simeq S$ を用いる。
これは $L$ の評価写像と両立するので cone にも降りる。
テンソル積についての coherence は可逆加群の評価の自然性と結合則による。
生成圏の等式は $\ell$ と $\ell^\vee$ による両方向のテンソル積から従う。
\end{proof}

\begin{theorem}[構造対象付き十七次元 chart からの残余圏の射]
\label{thm:jacobian-residual-dg-map}
定理\ref{thm:jacobian-algebraic-abel-substack}の十七次元 chart
$\mathfrak C_J=[J/\operatorname{PGL}_5]$ 上で、上の相対入力を満たす領域を取る。
相対 residual construction は構造対象付き Abel 同値を保つ Morita 同値に自然であり、
射
\begin{equation}\label{eq:jacobian-residual-dg-map}
 \kappa:\mathfrak C_J\longrightarrow\mathfrak{Cat}^{\mathrm{proper}}_{\mathrm{dg}}
\end{equation}
を定める。右辺は small idempotent-complete proper $S$-線型 dg 圏と
Morita 同値からなる高次群oidの fpqc stack 化を表す。
この記号は右辺の代数 stack 性を主張するものではない。
また、$\mathcal O(1)$ 自体の大域降下を仮定しない。
\end{theorem}
\begin{proof}
定理\ref{thm:general-relative-pointed-hom}により、
構造対象付き相対 Abel 同値は fpqc 局所的に graded algebra 同型に由来する。
graded 同型は $L_i$ を対応する $L_i'$ へ送り、導来圏上で評価写像を保つので
$\mathcal K_S$ と $\mathcal K_S'$ の Morita 同値を誘導する。
局所持上げの変更は grading scalar である。
grading scalar $\epsilon_\lambda$ に対して、構造対象を保つ比較は
\[
 F_{\epsilon_\lambda}(M)
       =M\otimes_{A,\epsilon_\lambda}A\longrightarrow M,
 \qquad m_d\otimes b\longmapsto\lambda^d m_db
\]
で与えられ、$\lambda$ の積と両立する。
これを残余圏に制限すれば、重なり上の比較と cocycle の coherence を得る。
同じ内容を $\mathcal O(1)$ の局所持上げで述べれば、異なる持上げは
基底の可逆加群を tensor する違いであり、
補題\ref{lem:kuznetsov-line-twist-invariance}が評価写像と射影関手の比較を与える。
従って残余圏の descent data は一次 twist 対象の自明化に依存しない。
対象・射・比較の基底変更自然性は
定理\ref{thm:relative-kuznetsov-projector}による。
これらを fpqc stack 化すると\eqref{eq:jacobian-residual-dg-map}を得る。
\end{proof}

\begin{remark}[dg 圏への射と三次モジュライへの射]
\eqref{eq:jacobian-residual-dg-map}の終域は dg 圏のモジュライである。
可換三次三様体のモジュライへの射を得るには、
ある非空開上で $\mathcal K_S$ が可換三次族の Kuznetsov 成分であることと、
その族の再構成の代数性・降下を別途示す必要がある。
その証明は第\ref{sec:kuznetsov-moduli-morphism}節で与える。
体値圏の categorical Torelli が与える一意性だけでは、
この存在と相対代数性を代用できない。
\end{remark}

\subsection{全圏と残余圏の smoothness の同値}

\begin{lemma}[係数 perfect 加群と smooth dg 代数]
\label{lem:relative-smooth-perfect-module}
$D$ を $S$-上 dg smooth な dg 代数とする。
右 $D$-加群 $N$ が $S$-perfect なら $D$-perfect である。
\end{lemma}
\begin{proof}
対角 $D$ は $D^{\mathrm{op}}\otimes_S^{\mathbf L}D$ の有限 cell と直和因子から得られる。
これに $N\otimes_D^{\mathbf L}(-)$ を適用すると
$N\simeq N\otimes_D^{\mathbf L}D$ は
$N\otimes_S^{\mathbf L}D$ の有限 cone と直和因子から得られる。
$N$ が $S$-perfect なので後者は右 $D$-perfect であり、結論を得る。
\end{proof}

\begin{lemma}[相対有限 gluing の smoothness]
\label{lem:relative-finite-gluing-smoothness}
$\mathcal D$ を classical generator を持つ $S$-proper dg 圏とし、
$\mathcal D=\langle\mathcal A,\mathcal B\rangle$ を
両成分の包含が左右随伴を持つ半直交分解とする。
このとき
\[
 \mathcal D\text{ は }S\text{-smooth}
 \quad\Longleftrightarrow\quad
 \mathcal A,\mathcal B\text{ はともに }S\text{-smooth}.
\]
\end{lemma}
\begin{proof}
まず smoothness が両側 admissible 成分へ降りることを確認する。
generator を取り $\mathcal D=\operatorname{Perf}(E)$,
$\mathcal A=\operatorname{Perf}(D)$ と書く。
包含を $-\otimes_D^{\mathbf L}P$ で表すと、$P$ は $D$-$E$ 双加群である。
右随伴は compact 対象を保つから、
$Q=\mathbf R\operatorname{Hom}_E(P,E)$ は右 $D$-perfect である。
左随伴の kernel を $T$ と書けば、$T$ は $E$-$D$ 双加群で右 $D$-perfect である。
左随伴性を $D\in\operatorname{Perf}(D)$ に対して評価すると
\[
 P\simeq\mathbf R\operatorname{Hom}_D(T,D)
\]
を $D$-$E$ 双加群として得る。従って $P$ は左 $D$-perfect である。
完全忠実性によって $P\otimes_E^{\mathbf L}Q\simeq D$ である。
$E$ の対角の有限 cell 表示に
$P\otimes_E^{\mathbf L}(-)\otimes_E^{\mathbf L}Q$ を適用する。
自由 $E$-双加群の像は $P\otimes_S^{\mathbf L}Q$ であり、これは
$P$ の左 perfect 性と $Q$ の右 perfect 性から $D$-双加群として perfect である。
従って $D$ の対角は perfect である。$\mathcal B$ も同様である。

逆向きを示す。
半直交分解の dg 模型は、二つの生成元を対角に置いた上三角 gluing
\[
 \Gamma=\begin{pmatrix}D&M\\0&F\end{pmatrix}
\]
と Morita 同値である。
これは二つの充満 dg 部分圏間の消える向きの acyclic Hom を零に置き換えて得られる。
この置換は quasi-fully faithful であり、pretriangulated idempotent 完備化が全圏を生成する。
$D,F$ は smooth なので $D\otimes_S^{\mathbf L}F^{\mathrm{op}}$ も smooth である。
$\mathcal D$ の properness により $M$ は $S$-perfect であり、
前補題から $M$ は $D$-$F$ 双加群として perfect である。

冪等元 $e_D,e_F$ を使うと、対角双加群には三角
\[
 \Gamma e_D\otimes_D^{\mathbf L}M\otimes_F^{\mathbf L}e_F\Gamma
 \longrightarrow
 \bigl(\Gamma e_D\otimes_D^{\mathbf L}D\otimes_D^{\mathbf L}e_D\Gamma\bigr)
 \oplus
 \bigl(\Gamma e_F\otimes_F^{\mathbf L}F\otimes_F^{\mathbf L}e_F\Gamma\bigr)
 \longrightarrow\Gamma
\]
がある。
二項から対角への射は積であり、最初の射は共通の off-diagonal 項を
二つの場所へ逆符号で入れる。
各 corner で計算すればこの三角が分かる。
必要なら $D,F,M$ を相対 flat な dg 模型に置き換えて計算する。
左右の corner による induction は、自由双加群を
$\Gamma$ の自由双加群の直和因子へ送るので perfect 性を保つ。
$D,F,M$ はそれぞれの双加群圏で perfect であるから、三角の最初の二項は
$\Gamma$-双加群として perfect であり、対角 $\Gamma$ も perfect である。
\end{proof}

\begin{corollary}[smoothness の同値]
\label{cor:relative-kuznetsov-smooth-equivalence}
定理\ref{thm:relative-kuznetsov-projector}の状況で
\[
 \mathcal C_S\text{ は }S\text{-dg smooth}
 \quad\Longleftrightarrow\quad
 \mathcal K_S\text{ は }S\text{-dg smooth}.
\]
従って、ある基底上で残余圏を滑らかな相対可換三次族の Kuznetsov 成分と同定できれば、
全非可換 tails 圏の dg smoothness も従う。
\end{corollary}
\begin{proof}
$\langle L_0,L_1\rangle$ は、二頂点と階数五の矢印加群を持つ相対 Kronecker 代数の
perfect 圏である。頂点の対角項と一つの矢印項による長さ一の相対双加群解消があり、
$S$-smooth である。
定理\ref{thm:relative-kuznetsov-projector}の admissibility と
前補題を適用する。
\end{proof}

\begin{remark}[既知の gluing 定理との関係]
体上では上三角 gluing の smoothness は
D.~Orlov, \emph{Smooth and proper noncommutative schemes and gluing of DG categories},
Adv. Math. \textbf{302} (2016), 59--105,
arXiv:1402.7364v5, Theorem~3.25
（同論文の参照 [LS, 3.24]）に述べられている。
同定理は体上で記述されているので、本節では任意係数環に必要な対角三角の証明を示した。
ここに用いた一般的な mutation と gluing の議論自体の新規性は主張しない。
本研究への適用点は、既に構成した任意係数上の proper 模型、有限 generator、
相対 Serre 双対、全構造対象付き同値比較を結び付けることにある。
\end{remark}

\subsection{適用範囲に関する確認事項}

上の全構成は central-CY 相対入力を満たす十七係数 atlas とその任意係数延長に適用される。
滑らかな可換点の全 Abel Artin 模型については、
命題\ref{prop:abel-artin-actual-residual}で別に必要な入力を確認する。
定理\ref{thm:abel-cubic-flat-envelope}\textup{(iii)}の中心商表示、
普遍有限 Koszul 解消と線対象のコホモロジーから compact/proper な有限 dg 模型を得て、
相対 CY orientation を用いずに本節の左射影を構成できる。
$\mathcal O_X(1)$ の持上げの一意性だけで properness を代用せず、
derived torsion と任意の導来底変更をこの模型で確認する。

本節の相対 residual functor、降下、smoothness 比較を入力として、
次節では十七次元 chart の非空開上に可換三次モジュライへの射を構成する。
全 Abel 形式近傍上の残余射と、その射に両立する Poisson 形式近傍への同型は、
第\ref{sec:poisson-cubic-lifts}節で扱う。

\section{Kuznetsov成分による可換三次モジュライへの射}
\label{sec:kuznetsov-moduli-morphism}

$M_{\mathrm{cub}}$ を滑らかな可換三次三様体の粗モジュライ空間とする。
これは滑らかな三次形式の開部分 $H\subset\mathbb P(\operatorname{Sym}^3\mathbb C^5)$
の $\operatorname{PGL}_5$ による幾何学的商であり、分離的な十次元既約代数多様体である。
滑らかな三次の GIT 安定性には \cite[Theorem~1.1]{AllcockCubicGIT} を用いる。
自己同型が自明な三次の非空開部分を $M_{\mathrm{cub}}^\circ$ と書く。
この開上では粗モジュライと通常の三次モジュライ stack は一致する。
この開の存在は \cite[Theorems~1.5--1.6 and Corollary~1.9]{PoonenHypersurfaceAut}
による。$[H/\operatorname{PGL}_5]$ は滑らかなので、$M_{\mathrm{cub}}^\circ$ も滑らかである。

\begin{theorem}[十七次元から十次元への残余射]
\label{thm:kuznetsov-cubic-moduli-morphism}
定理\ref{thm:jacobian-algebraic-abel-substack}の
十七次元 Jacobian 型 Abel 圏モジュライ $\mathfrak C_{\mathrm J}$ には
非空開部分 stack $\mathfrak U$ と支配的な射
\begin{equation}\label{eq:kuznetsov-seventeen-ten}
 \rho:\mathfrak U\longrightarrow M_{\mathrm{cub}}
\end{equation}
が存在する。任意の幾何学的点 $c$ に対し、$\rho(c)=[X_c]$ は
\[
 \mathcal K_c\simeq\operatorname{Ku}(X_c)
\]
を満たす滑らかな可換三次の唯一の同型類である。
同値は dg Morita 同値として取れる。さらに次を満たすように縮小できる。
\begin{enumerate}
\item $\mathfrak U$ の étale 被覆上に滑らかな三次族と相対 Morita 同値が存在する。
相対残余圏と全 compact tails 圏は dg smooth proper であり、任意の導来底変更と両立する。
\item $\rho$ は $M_{\mathrm{cub}}^\circ$ のある非空開部分 $V$ に値を取り、
$\rho:\mathfrak U\to V$ は相対次元七の滑らかな全射である。
従って $\mathfrak U\to V\subset\mathfrak M_{\mathrm{cub}}^{\mathrm{sm}}$
は可換三次のモジュライ stack への実際の射でもある。
\end{enumerate}
特に十七次元の scheme atlas の非空開 $U$ 上にも
$U\to M_{\mathrm{cub}}$ という射を得る。
\end{theorem}

これは全 O-marked Abel モジュライの代数性やその中での
$\mathfrak C_{\mathrm J}$ の開性を主張する定理ではない。
また、式\eqref{eq:kuznetsov-seventeen-ten}を、量子化の入力三次形式を
そのまま出力する写像と同一視しない。
本節では形式的認識、係数 perfect な双加群の効果化、Serre 安定な像の生成性を
順に証明し、最後に有限表示の降下から実際の開集合上の射を得る。
証明の各段階が与える結論を先にまとめる。
\begin{enumerate}[label=\textup{(\roman*)}]
\item 命題\ref{prop:kuznetsov-formal-recognition}は、通常の Artin 基底上の
三次変形と残余圏変形の\emph{同型類}を比較する。
\item 命題\ref{prop:kuznetsov-effective-serre-embedding}は、完備離散付値環上で
係数 perfect な同値核を効果化し、Serre 関手と両立する完全忠実関手を作る。
\item 補題\ref{lem:kuznetsov-graded-thick-simple}と
系\ref{cor:kuznetsov-formal-to-generic}は、三次の特異圏に固有の生成性から
一般ファイバーでの本質的全射性を与える。
\item 補題\ref{lem:kuznetsov-spread-geometric-model}は、その幾何的模型を
有限 étale 被覆上へ広げる。粗モジュライへの射は categorical Torelli によって降下し、
支配性は別に補題\ref{lem:kuznetsov-classical-dominance}で確認する。
\end{enumerate}
従って、形式的変形の認識だけから代数的な一般点の幾何性を結論する段階はない。

\subsection{可換三次成分の全Artin形式認識}

\begin{proposition}[通常のArtin基底上の幾何的認識]
\label{prop:kuznetsov-formal-recognition}
$k/\mathbb C$ を体拡大、$X/k$ を滑らかな三次三様体、
$\mathcal D=\operatorname{Ind}(\operatorname{Ku}(X))$ とする。
特殊ファイバーを同定した変形を考える。
通常の局所 Artin $k$-代数 $T$ 上で、自然写像
\[
 \operatorname{Def}_X(T)\longrightarrow
 \pi_0\operatorname{CatDef}^{c}_{\mathcal D}(T),\qquad
 X_T\longmapsto\operatorname{Ind}(\operatorname{Ku}(X_T))
\]
は全単射である。右辺は compactly generated な線型圏の実際の変形空間の同型類を表す。
これは高次自己同型まで含む変形空間を離散化された左辺と同一視する主張ではない。
\end{proposition}
\begin{proof}
用いる Hochschild 計算と幾何的 Kodaira--Spencer 写像は
\begin{equation}\label{eq:kuznetsov-formal-hh}
 HH^i(\operatorname{Ku}(X))=
 \begin{cases}k&i=0,\\ k^{10}&i=2,\\0&i\ne0,2,\end{cases}
 \qquad
 H^1(X,T_X)\xrightarrow{\sim}HH^2(\operatorname{Ku}(X))
\end{equation}
である。前者は \cite[Theorem~8.9]{KuznetsovHochschild}、
後者は \cite[Theorem~4.4(2)]{JLLZInfinitesimal} である。
後者では単なる次元の一致でなく、三次の変形から成分の変形へ誘導される
実際の写像が同型であることを使う。
これらは普遍的な滑らかな三次族上の代数的コホモロジーとその射として定まり、
複素閉点で同型であることから普遍族上で同型となるため、任意の $k/\mathbb C$ に適用できる。
成分を $\langle\mathcal O(-1),\mathcal O\rangle^\perp$ とする規約との差は
線束 twist による同値である。

$\mathcal D$ は proper な compact 部分圏で生成されるので、
compact 対象間の Ext が有界であり、tamely compactly generated である。
Lurie \cite[Definition~5.3.20, Proposition~5.3.21 and Theorem~5.1.9]{LurieDAGX} により
\[
 \theta:\operatorname{CatDef}^{c}_{\mathcal D}
       \longrightarrow\operatorname{CatDef}^{\wedge}_{\mathcal D}
\]
は $(-1)$-truncated である。従って $\pi_0$ 上単射であり、
各点の $\pi_i$、$i\ge1$ 上同型である。
同じ結論と接複体 $HH(\mathcal D)[2]$ は
\cite[\S3.3]{BKPFormalGenerators} にも明記されている。
Lurie の Theorem~5.3.16 による形式 hull の同定と、標数零の形式変形理論から、
可換 Artin 環への制限は $H^i(\mathfrak g)=HH^{i+1}(\mathcal D)$ を持つ
dg Lie 代数 $\mathfrak g$ により制御される。

$H(\mathfrak g)$ は次数 $-1,1$ にのみ存在する。
移送した極小 $L_\infty$ 構造の $n$ 項演算は次数 $2-n$ を持つ。
奇数次数の $n$ 個の入力に対する出力次数は偶数なので、すべての高次演算は零である。
従って通常の Artin 基底上では
$\pi_0\operatorname{CatDef}^{\wedge}_{\mathcal D}$ は十次元形式円板であり、
$\pi_1=0$ である。
極大イデアルを $\mathfrak m_T$ とすれば、各点で
$\pi_2\simeq HH^0(\mathcal D)\otimes_k\mathfrak m_T$
が残る。これは標数零の対数によって $1+\mathfrak m_T$ と同定される
形式的スカラーであり、上の全単射が高次群oid全体の離散性を意味しない理由である。

一方 $H^0(T_X)=H^2(T_X)=0$、$h^1(T_X)=10$ なので、
$\operatorname{Def}_X$ は滑らかな十次元形式円板である。
実際の幾何的写像の接写像は\eqref{eq:kuznetsov-formal-hh}の同型である。
形式逆写像定理により
$\operatorname{Def}_X\to\pi_0\operatorname{CatDef}^{\wedge}_{\mathcal D}$ は同型である。
幾何的変形は実際の compactly generated な変形を与えるので、
$\theta$ の $\pi_0$ 上の像は hull の全体を含む。
単射性と合わせて主張を得る。
\end{proof}

\subsection{形式的同値からSerre安定な埋込みへ}

\begin{proposition}[係数perfectな効果化とSerre比較]
\label{prop:kuznetsov-effective-serre-embedding}
$R=k[[h]]$、$R_n=R/(h^{n+1})$ とし、$E$ を $R$-perfect な dg $R$-代数とする。
$\mathcal C=\operatorname{Perf}(E)$ が、自然な双加群 pairing を備えた
dg/enriched な相対 Serre 自己同値を持ち、その構造が底変更に両立し、
$\mathcal C\otimes_R k\simeq\operatorname{Ku}(X_0)$ と仮定する。
このとき滑らかな相対三次 $X_R/R$ と完全忠実な $R$-線型 dg 関手
\begin{equation}\label{eq:kuznetsov-effective-embedding}
 \Phi:\operatorname{Perf}(E)\longrightarrow\operatorname{Ku}(X_R)
\end{equation}
が存在する。$\Phi$ の特殊ファイバーは指定した同値であり、
像は target の相対 Serre 関手とその逆で安定である。
ここでは $E$ の dg smoothness は仮定しない。
\end{proposition}
\begin{proof}
命題\ref{prop:kuznetsov-formal-recognition}により
$\operatorname{Perf}(E\otimes_R^{\mathbf L}R_n)\simeq\operatorname{Ku}(X_n)$
となる三次変形 $X_n/R_n$ を得る。
認識の自然性と形式表現性により $X_n$ は互換的に取れる。
同値も互換的に選べる。実際、一段上の同値の還元と既に選んだ同値との差は
特殊ファイバー上恒等な自己同値であり、前証明の $\pi_1=0$ により同型である。
隣接する比較を逐次選び、$\infty$ 圏内で tower の合成として
homotopy coherent な高次の比較を取る。
スカラーの二次自己同型を消滅したものとは扱わない。

三次の偏極は $H^1(\mathcal O)=H^2(\mathcal O)=0$ により互換的に持ち上がる。
大域切断の基底と三次方程式の生成元を互換的に選び、その係数の逆極限を取ると
$X_R\subset\mathbb P^4_R$ を得る。
判別式の特殊値は非零であり $R$ の単元なので、この族は滑らかである。
$\mathcal D=\operatorname{Ku}(X_R)$ と置き、
$\mathcal D\simeq\operatorname{Perf}(D)$ を有限 generator で表示する。
相対 smooth な射影三次の perfect 圏とその admissible 成分なので、$D$ は $R$-smooth proper である。
generator と $D$ は $R_n$ への底変更に両立する。

互換的な同値を $E_n$--$D_n$ 双加群 $P_n$ で表す。
各 $P_n$ は右 $D_n$-perfect であるため $R_n$-perfect である。
Bhatt \cite[Lemma~4.2]{BhattAlgebraization} の $\infty$ 圏としての同値
\[
 \operatorname{Perf}(R)\simeq\varprojlim_n\operatorname{Perf}(R_n)
\]
により、underlying complex は $P\in\operatorname{Perf}(R)$ に効果化される。
作用の効果化も確認する。
$\mathbf R\operatorname{End}_R(P)$ は $R$-perfect、従って導来 $h$-完備であり、
$E_1$-代数として
\[
 \mathbf R\operatorname{End}_R(P)
 \simeq\varprojlim_n\mathbf R\operatorname{End}_{R_n}(P_n).
\]
代数の極限は underlying 加群で計算される。
$\Gamma=E\otimes_R^{\mathbf L}D^{\mathrm{op}}$ からの写像空間は target の極限を保つので、
\[
 \operatorname{Map}_{\operatorname{Alg}_R}
       (\Gamma,\mathbf R\operatorname{End}_R(P))
 \simeq\varprojlim_n
 \operatorname{Map}_{\operatorname{Alg}_{R_n}}
       (\Gamma_n,\mathbf R\operatorname{End}_{R_n}(P_n)).
\]
従って互換的な作用が $E$--$D$ 双加群構造を与える。
この段階で $P$ が $\Gamma$-perfect であるとは主張しない。

$D$ が $R$-smooth なので、補題\ref{lem:relative-smooth-perfect-module}から
$P$ は右 $D$-perfect である。
左作用による射 $E\to\mathbf R\operatorname{End}_D(P)$ の両辺は $R$-perfect であり、
その特殊ファイバーは $P_0$ の同値性により quasi-isomorphism である。
cone に導来 Nakayama を適用して、この射自体が quasi-isomorphism であることを得る。
従って $-\otimes_E^{\mathbf L}P$ は\eqref{eq:kuznetsov-effective-embedding}を与える。

最後に Serre 関手をそれぞれ $S_{\mathcal C},S_{\mathcal D}$ と書く。
完全忠実性と両側の enriched Serre 双対性は、$M,N\in\mathcal C$ に自然な同型
\[
 \mathbf R\operatorname{Hom}_{\mathcal D}(\Phi N,S_{\mathcal D}\Phi M)
 \simeq\mathbf R\operatorname{Hom}_{\mathcal C}(M,N)^\vee
 \simeq
 \mathbf R\operatorname{Hom}_{\mathcal C}(N,S_{\mathcal C}M)
\]
を与える。右辺を $N$ の関手と見て enriched Yoneda を適用すると、
恒等射 $\operatorname{id}_{S_{\mathcal C}M}$ に対応する射
$\alpha_M:\Phi S_{\mathcal C}M\to S_{\mathcal D}\Phi M$ を得る。
この射は source の Serre trace によって特徴付けられ、$M$ に自然である。
同型を dg 双加群の同型として用いているので、比較は homotopy coherent である。
特殊ファイバーの $\Phi_0$ は同値なので $\alpha_M\otimes_R k$ は同型である。
その cone は $D$-perfect、従って $R$-perfect であり、導来 Nakayama により零となる。
$S_{\mathcal C}$ は自己同値なので逆 Serre に対しても安定である。
すべての比較を有限 generator に適用すれば、任意の導来底変更後にも同じ比較が成立する。
一般ファイバーの像の両側 Serre 安定性には、$M=E$ と
$M=S_{\mathcal C}^{-1}E$ の比較だけで十分である。
一般ファイバーの全対象が $R$ 上へ持ち上がるとは仮定しない。
\end{proof}

\subsection{三次成分に固有の生成性}

\begin{lemma}[次数シフトで安定な部分圏の生成]
\label{lem:kuznetsov-graded-thick-simple}
$S=k[x_0,\ldots,x_4]$、$f\in S_3$ とし、$V(f)\subset\mathbb P^4_k$ は滑らかとする。
有限階数の graded matrix factorization の冪等完備化を $\mathcal F$ と書く。
非零の $P\in\mathcal F$ に対し
\[
 \operatorname{thick}\{P(a)\mid a\in\mathbb Z\}=\mathcal F.
\]
従って $\operatorname{Ku}(V(f))$ の非零 thick 部分圏が
Serre 関手とその逆で安定なら、それは全体である。
\end{lemma}
\begin{proof}
$S/(f)$ の特異点集合は原点だけである。
Stevenson \cite[Theorem~6.13]{StevensonSingularityActions} は、
Noetherian かつ局所的に hypersurface な環の特異圏の全 thick 部分圏を
特異点集合の特殊化閉部分集合で分類する。
従って ungraded matrix factorization 圏の冪等完備化では、
どの非零対象も全体を thick-generate する。
affine な定理を使っており、ここに追加の tensor 安定性の仮定はない。

graded 圏への移行を明示する。
忘却関手を $\mathsf F$ とし、任意階数の graded projective 加群による
matrix factorization の大きい homotopy 圏を考える。
ungraded factorization $Q$ に対して、graded 化 $\mathsf G$ を
\[
 (\mathsf GQ_i)_n=Q_i\quad(n\in\mathbb Z)
\]
で定める。$s\in S_e$ は copy $n$ から copy $n+e$ へ $s$ 倍として作用する。
$\mathsf G(S)\simeq\bigoplus_{a\in\mathbb Z}S(a)$ であり、
$\mathsf G$ は直和因子を保つので、projective 加群を graded projective 加群へ送る。
graded factorization を
$P_1\xrightarrow{d_1}P_0\xrightarrow{d_0}P_1(3)$ と書く規約では、
$\mathsf Gd_1$ は copy の次数を保ち、$\mathsf Gd_0$ は三だけ増加させる。
これで二つの合成は $f$ 倍になる。
$\mathsf G$ は homotopy、cone、直和因子を保ち、
copy の添字の平行移動 $\mathsf GQ_i(3)\simeq\mathsf GQ_i$ により shift とも両立する。

graded $P$ について、copy $n$ にある元の元来の内部次数を $j$ とすると、
$a=j-n$ による分解が
\[
 \mathsf G\mathsf F(P)\simeq\bigoplus_{a\in\mathbb Z}P(a)
\]
を与える。$d_0$ は $j,n$ をともに三増加させるので、この分解は微分とも両立する。
特に $P$ は右辺の $a=0$ の直和因子である。
忘却後に $P$ が contractible なら、contracting homotopy の該当する homogeneous 成分を
取ることで graded contracting homotopy を得る。
冪等完備化の対象 $(P,e)$ についても、忘却後の $e$ が null-homotopic なら、
その null-homotopy の該当する homogeneous 成分が graded な $e$ の null-homotopy を与える。
従って冪等完備化の上でも $\mathsf F$ は非零性を保つ。

任意の有限階数 graded $Q$ に対し、ungraded の分類から
$\mathsf FQ\in\operatorname{thick}(\mathsf FP)$ である。
$\mathsf G$ を適用し、上の分解と $Q$ の直和因子表示を使うと
$Q\in\operatorname{Loc}\{P(a)\}$ を得る。
有限階数 factorization は、大きい homotopy 圏で compact である。
実際、各 Hom 複体は有限生成射影加群からの Hom で計算され、直和と可換する。
compact 対象が生成する localizing 圏の compact 部分はその thick 閉包である
\cite[Lemma~2.2]{NeemanCompactLocalization}から、
$Q\in\operatorname{thick}\{P(a)\}$ を得る。

Orlov の同値 \cite[Theorems~2.5 and~3.10--3.11]{OrlovSingularities} と
\cite[\S3.3, Corollary~3.18, Lemma~3.19 and Corollary~4.1]{KuznetsovFractionalCY} の規約で、
次数シフトは rotation $\mathsf R$ に対応する。
Orlov の半直交分解の window は線束 twist と Serre mutation により
$\langle\mathcal O,\mathcal O(1)\rangle^\perp$ の window に移して用いる。
三次三様体では
$\mathsf R^3=[2]$、$S=\mathsf R^{-2}[3]$ である。
後者は $S=\mathsf R[1]$、従って $\mathsf R=S[-1]$ と書ける。
この二つの公式は単なる $S^3=[5]$ より強く、
Serre 安定性を\emph{全ての内部次数シフト}に対する安定性へ移すために用いる。
よって Serre 安定な非零 thick 部分圏には前段の結論を適用できる。
\end{proof}

\begin{corollary}[非零一般ファイバーの幾何性]
\label{cor:kuznetsov-formal-to-generic}
命題\ref{prop:kuznetsov-effective-serre-embedding}の状況で、
$\operatorname{Perf}(E\otimes_R k((h)))\ne0$ とする。
このとき\eqref{eq:kuznetsov-effective-embedding}の一般ファイバーは Morita 同値である。
\end{corollary}
\begin{proof}
完全忠実関手の像は冪等完備な thick 部分圏であり、
Serre 比較の底変更により Serre とその逆で安定である。
非零性と前補題により像は全体である。
\end{proof}

\subsection{可換三次を含む代数的な十八係数の橋渡し族}
\label{subsec:kuznetsov-algebraic-bridge}

\begin{proposition}[完備化を使わない相対量子化族]
\label{prop:kuznetsov-algebraic-bridge}
独立変数 $a,b,\alpha,\beta,m_{ij},g_j,h$ を取り、
\[
 F=a(x^3+y^3+z^3)+bxyz+ux^2+\alpha vy^2+\beta(u+v)z^2
       +\sum_{i=1}^3M_i(u,v)x_i+g(u,v),
\]
$M_i=\sum_{j=0}^2m_{ij}u^{2-j}v^j$、$g=\sum_{j=0}^3g_ju^{3-j}v^j$ と置く。
さらに
\[
 p_h=2-hb,\qquad q_h=2+hb,\qquad r_h=6ha,
 \qquad d(h)=p_hq_h(q_h^3-r_h^3)(p_h^3+r_h^3)
\]
と定める。$d(0)=256$ である。
係数環
\[
 R=\mathbb C[a,b,\alpha,\beta,m_{ij},g_j,h,
                 \{a(27a^3+b^3)d(h)\}^{-1}]
\]
上で $u,v$ を中心変数とし、
\[
 \begin{aligned}
 B_h&=R\langle x,y,z,u,v\rangle/I_h,\\
 I_h&=\bigl([u,-],[v,-],\ [x,y]-h\operatorname{Symm}(F_z),\\
 &\hspace{20mm}[y,z]-h\operatorname{Symm}(F_x),\ [z,x]-h\operatorname{Symm}(F_y)\bigr).
 \end{aligned}
\]
を定める。すると次が成立する。
\begin{enumerate}
\item $B_h$ に中心三次元 $w_h$ が代数的に存在し、
\[
 B_h/(h)=R/(h)[x,y,z,u,v],\qquad w_h\bmod h=F
\]
となる。
\item $B_h$ と $A_h=B_h/(w_h)$ は、
第\ref{sec:relative-koszul-line-cohomology}節の相対入力 (1)--(4) を満たす。
従って任意の可換 $R$-代数上で五つの線対象が compact generator を与え、
その自己 Hom dg 代数は係数 perfect であり、相対 Serre 双対と
任意の導来係数延長比較が成立する。
\item $\operatorname{Spec}R$ は整な有限型十八次元基底であり、
$h=0$ には滑らかな可換三次の点がある。
\item $h\ne0$ 上の族は、係数変換と二つの中心変数の変更により
既存の十七係数族と一致する。
その係数写像は支配的である。
従って $\operatorname{Spec}R$ の任意の非空開から、十七係数族の一般点へ到達する。
\end{enumerate}
\end{proposition}

\begin{proof}
\emph{有限の分母と中心元。}
命題\ref{prop:pair-rains-polynomial}へ
\[
 \begin{gathered}
 q=\frac{q_h}{p_h},\qquad t=\frac{r_h}{p_h},\qquad
 \lambda=\frac b{3a},\\
 (L_1,L_2,L_3)=\frac{(2u,2\alpha v,2\beta(u+v))}{3a},
 \qquad\widetilde M_i=\frac{M_i}{3a}
 \end{gathered}
\]
を代入する。$q=1+\lambda t$ であり、関係は上記 $B_h$ の関係そのものである。
Rains の分母は
\[
 q(q^3-t^3)(q+1)(1+t^3)
 =\frac{4q_h(q_h^3-r_h^3)(p_h^3+r_h^3)}{p_h^8}
\]
なので $R$ 上で可逆である。
active 定数項を零に規格化した Rains 元を $G_t^{\mathrm R}$ と置き、
\[
 w_h=3aG_t^{\mathrm R}+g(u,v)
\]
とする。この有限有理式は $R$ 上で定義され、中心的である。
$q(0)=1,t(0)=0$ なので $w_h\bmod h=F$ となる。

\emph{同一の代数的基底上での全次数の射影性。}
$R_0=\mathbb C[a,b,h,\{a(27a^3+b^3)\}^{-1}]$ とし、
active 最高次部分の代数を $S_h$ と書く。
その関係は固定した滑らかな平面三次
$f_3=a(x^3+y^3+z^3)+bxyz$ に量子化係数 $h$ を付けたものである。
Mori--Smith \cite[Theorem~4.2(1)]{MS} は
$A_{f,\lambda}$ について全ての $\lambda$ を許し、$\lambda=0$ も含む。
従って $R_0$ の全ての幾何学的ファイバーは
三次元 Koszul AS 正則 CY 整域であり、Hilbert 級数は $(1-s)^{-3}$ である。
$h=0$ ではこれは多項式環である。

各 $(S_h)_n$ は有限な語の加群の有限表示商であり、全閉点で
階数 $\binom{n+2}{2}$ を持つ。
$R_0$ は被約な有限型 $\mathbb C$-環なので、表示行列の高すぎる小行列式は
全閉点で零であることから恒等的に零となる。
必要な階数の小行列式が可逆な局所開上で行列を標準形へ直せば、
$(S_h)_n$ は $R_0$-有限射影であり、任意の底変更と両立する。
ここでは次数ごとに基底を縮小していない。
全次数で同じ $R_0$ 上に一定階数が成立することを使っている。

$Q=R[u,v]$ と置く。
active 次数を下げる cyclic potential の項に対する PBW 定理
\cite[Theorem~3.1]{BT}\cite[Proposition~7.3.1(ii)]{EG} を各閉点に適用する。
各 $(S_h)_n$ への語の全射を係数線型に分裂し、その持上げを $B_h$ に取る。
最高 active 次数からの帰納法でこれらは $Q$-加群として $B_h$ を生成する。
有限の関係があれば全閉点の PBW 基底でその係数が零となり、
被約性により関係自体が零となる。
従って
\[
 \operatorname{gr}_{\mathrm{active}}B_h\simeq S_h\otimes_{R_0}Q,
 \qquad B_h\simeq\bigoplus_{n\ge0}(S_h)_n\otimes_{R_0}Q
\]
を得る。後の同型は濾過付き $Q$-加群としての非標準的な分裂である。
特に $B_h$ は $Q$-射影で、全次数成分は有限射影 $R$-加群となる。
この分裂から PBW 同型は非平坦なものも含む任意の係数延長後に保たれる。

\emph{正則三次商。}
各幾何学的ファイバーの $B_h$ は、随伴次数付きが
$S_h\otimes k[u,v]$ であるため整域である。
Rains 元の active 最高記号は各ファイバーで非零である。
実際、命題\ref{prop:pair-rains-polynomial}の三次関係行列の小行列式
$-q^5(q^3-t^3)^2$ は $R$ 上可逆であり、同命題の整列基底表示に現れる
$(q+1)(1+t^3)$ および $(q+1)(q^3-t^3)$ も可逆である。
$3a$ も可逆なので、この最高記号は零にならない。
従って全閉点で $w_h$ は正則である。
有限射影加群の写像
\[
 B_{h,n-3}\xrightarrow{w_h}B_{h,n}
\]
は全閉点で単射で一定階数を持つので、その cokernel は有限射影であり、
各次数の完全列は係数線型に分裂する。
これで $A_{h,n}$ の階数
$\binom{n+4}{4}-\binom{n+1}{4}$ と、任意の底変更後の正則性を得る。
楕円境界の滑らかさや、その translation の有限位数性はこの議論に不要である。

\emph{真正な相対 CY と有限 Koszul 対角。}
命題\ref{prop:general17-relative-cy-koszul}の証明を、この基底上で反復できる。
具体的には
\[
 \Phi_h=xyz-yxz-h\operatorname{Symm}(F-g)
\]
の三変数 potential 双加群複体を $Q$ 上で作る。
各全次数は有限射影 $R$-加群の有限複体であり、
各閉点では PBW 濾過によって三次元 CY の potential 解消となる。
従って同一基底上で分裂完全であり、任意の係数延長後にも完全である。
cyclic Hessian の反転恒等式が、この複体の真正な双加群自己双対
$K_Q^\vee\simeq K_Q(3)[-3]$ を与える。
$B_h$ の $Q$-射影性により二つの中心差
$u_L-u_R,v_L-v_R$ は $B_h\otimes_R B_h^{\mathrm{op}}$ 上で正則列である。
その Koszul 双対と change of rings を使うと
\[
 \mathbf R\underline{\operatorname{Hom}}_{B_h\otimes_RB_h^{\mathrm{op}}}
 (B_h,B_h\otimes_RB_h^{\mathrm{op}})\simeq B_h(5)[-5]
\]
を得る。この同型は相対双加群として自然で、任意の底変更と両立する。

最後に全ファイバーの五次元 Koszul 性を使い、二次関係から定める
有限交わり $W_i$ は $0\le i\le5$ で階数 $\binom5i$、$W_6=0$ となる。
有限行列の一定階数により $W_i$ は有限射影で、基底変更と可換である。
標準 Koszul 双加群複体の各全次数は全閉点で完全な有限射影複体なので、
同じ基底上で普遍的に完全である。
これで相対入力 (1)--(4) が揃う。
全ての次数にわたる正確性を、有限個の次数の開条件から推測したのではなく、
同一基底の全ファイバーに適用できる PBW/CY 定理から証明したことに注意する。

\emph{可換な滑らかな点。}
$h=0$ において
\[
 F_*=x^3+y^3+z^3+ux^2+u^3+v^3
\]
を取る。この係数は $a=1,b=0,\alpha=\beta=0,M_i=0,g=u^3+v^3$ に対応し、
基底 $R$ の条件を満たす。
偏微分の同時零点は $y=z=v=0$ および
$x(3x+2u)=0$, $x^2+3u^2=0$ を満たす。
$x=0$ なら $u=0$、$x=-2u/3$ なら $31u^2/9=0$ なので、
射影空間に同時零点は存在しない。
必要なら $F$ の射影 Jacobian 判別式をさらに反転しても、
この点と整性・支配性は保たれる。

\emph{十七係数一般点への支配性。}
$h\ne0$ で $u'=hu,v'=hv$ と置くと、周囲の関係式は
\[
 a'=ha,\qquad b'=hb,\qquad
 \alpha'=\alpha,\quad\beta'=\beta,\qquad m'_{ij}=h^{-1}m_{ij}
\]
を持つ標準十七係数族の形になる。
中心三次を canonical な未スケールの $G^{\mathrm R}$ で規格化する場合、
$w_h$ を可逆元 $3a$ で割り
\[
 g'_j=\frac{g_j}{3ah^3}
\]
を使う。中心を $3a'G^{\mathrm R}$ に規格化する場合は
$w_h$ を $h$ 倍して $g'_j=h^{-2}g_j$ とする。
いずれも同じ三次直線を記述し、規格化を途中で混在させない。
前者について逆変換は
\[
 a=a'/h,\qquad b=b'/h,\qquad
 m_{ij}=hm'_{ij},\qquad g_j=3a'h^2g'_j
\]
である。
従って $\mathbb C(h)$ 上で係数変換は双有理であり、
十七係数空間の函数体は $\operatorname{Frac}(R)$ に単射となる。
既存の $\Delta$ 開と étale chart の開を引き戻しても非空であり、
これは主張する一般点での一致と支配性を与える。
\end{proof}

\begin{remark}[この橋渡しが解決する点]
ここでは $\mathbb C[[h]]$ 上の次数ごとの Nakayama だけから
有限型の開近傍を推論していない。
固定した滑らかな $f_3$ に対し Mori--Smith の定理が全量子化係数 $h$ で成立すること、
低次 potential の PBW 定理、有限な Rains 公式を組み合わせ、
初めから一つの有限型代数的基底上で任意係数の proper dg 族を作った。
この命題自体は三次商の dg smoothness も残余圏の幾何性も主張しない。
それらを証明するための、滑らかな可換点と十七次元一般点を結ぶ実際の代数的族を供給する。
\end{remark}


\subsection{有限表示への降下と粗モジュライの降下}

\begin{lemma}[幾何的一般点からの開上の幾何的模型]
\label{lem:kuznetsov-spread-geometric-model}
$T$ を滑らかな整有限型 $\mathbb C$-スキーム、$K=\mathbb C(T)$ とし、
$E_T$ を係数 perfect な dg 代数族とする。
ある体拡大 $L/K$ 上で
$\operatorname{Perf}(E_L)\simeq\operatorname{Ku}(Y_L)$ となる滑らかな三次が存在すると仮定する。
すると $T$ の非空開 $T'$ と有限 étale 全射 $V\to T'$ 上に、
滑らかな三次族 $Y_V$ および相対 Morita 同値
$\operatorname{Perf}(E_V)\simeq\operatorname{Ku}(Y_V)$ が存在する。
また、この同値によって一意に定まる射 $T'\to M_{\mathrm{cub}}$ がある。
\end{lemma}
\begin{proof}
$E_L$ は smooth proper なので、補題\ref{lem:dg-smooth-field-descent}により $E_K$ もそうである。
三次の有限個の係数は有限生成 $K$-部分代数 $B\subset L$ へ降り、
判別式を反転して滑らかな三次 $Y_B$ を得る。
その相対成分の generator の自己 Hom 代数を $D_B$ とする。
$D_B$ は $B$-smooth proper である。
$L$ 上の同値 kernel は
$E_K\otimes_K D_B^{\mathrm{op}}$ の底変更上の perfect 加群である。
perfect 加群の連続性 \cite[Lemma~2.10]{ToenVaquieObjects} により、
$B$ を有限生成 $K$-部分代数として拡大すれば、その kernel と逆 kernel、
unit と counit が $B$ 上へ降りる。
tensor 積とこれらの射の cone は $B$-perfect であり、$L$ 上で零となる。
有限個の分母を反転して、二つの cone をともに零にできる。
ここでは両側の smooth proper 性により相対 kernel と Hom の有限性を確保している。

非空な $\operatorname{Spec}B$ の閉点を取ると、有限体拡大 $K'/K$ 上で
三次と同値が存在する。標数零なので $K'/K$ は分離的である。
$T$ の正規化を取り有限 étale になる非空開へ制限する。
さらに一般点から perfect kernel と逆 kernel、unit、counit を広げ、
三次の判別式と係数 perfect な cone の support を除く。
正規化の射は有限なので、除く閉部分の像も閉であり、
$T$ 自体を縮小することで被覆のすべての点上に模型を得る。
必要な相対 smoothness は補題\ref{lem:dg-generic-spreading}により先に確保できる。

この被覆の二重交差 $V\times_{T'}V$ は滑らか、特に被約である。
二つの三次の成分は同じ $E$ と同値なので、各幾何学的点で
categorical Torelli \cite[Theorem~1.1]{BMMSCubicTorelli} により同型である。
従って $M_{\mathrm{cub}}$ への二つの写像は点ごとに一致する。
分離的な target への二写像の等化子は閉部分空間であり、
被約な source の全点を含むので source 全体である。
よって étale 降下から $T'\to M_{\mathrm{cub}}$ を得る。

ここで用いた Torelli の一次資料は複素数体上である。
拡大体上の適用は次のように帰着する。
二つの滑らかな三次と perfect 同値 kernel を有限生成被約 $\mathbb C$-環へ広げ、
同値性を保つ開へ制限する。
すべての複素閉点で二つの粗モジュライ写像は一致し、
Jacobson 性、被約性、分離性から写像自体が一致する。
これを元の体へ戻せば、その代数閉包上で同型類が等しい。
従って解析的な Torelli を未確認のまま任意の体へ拡張してはいない。
\end{proof}

\begin{lemma}[通常の三次族の支配性]
\label{lem:kuznetsov-classical-dominance}
命題\ref{prop:kuznetsov-algebraic-bridge}の $h=0$ の十七係数三次族は
$M_{\mathrm{cub}}$ を支配する。
\end{lemma}
\begin{proof}
三次形式の三十五次元ベクトル空間への写像
$\operatorname{GL}_5\times\mathbb A^{17}\to\operatorname{Sym}^3(\mathbb C^5)$
の微分を考える。二十五本の列は $x_a\partial F/\partial x_b$、
残る十七本は係数についての偏微分である。
\[
 (a,b,\alpha,\beta)=(3,5,2,3),\qquad
 (m_{ij})=\begin{pmatrix}1&1&2\\2&3&1\\1&4&5\end{pmatrix},
 \qquad(g_0,g_1,g_2,g_3)=(1,2,3,4)
\]
において、この $35\times42$ 整数行列の指定した三十五次小行列式は
$752849\pmod{1000003}$ である。
従って標数零で階数三十五となり、写像は支配的である。
行と列の全指定、入力、再生成コードは
\texttt{computations/commutative\_dominance\_certificate.json} と
\texttt{computations/verify\_commutative\_dominance.py} にある。
同じ計算で Euler 恒等式 $\sum_i x_iF_{x_i}=3F$ は整数上で確認する。

滑らかな三次の locus は非空開であり、橋渡し命題の $F_*$ が非空性の証人である。
従って上の支配性を滑らかな locus と $a(27a^3+b^3)\ne0$ へ制限しても保たれる。
幾何学的商を取れば主張を得る。
この行列計算は categorical な幾何性の証明ではなく、通常の三次族の支配性だけを証明する。
\end{proof}

\begin{proof}[定理\ref{thm:kuznetsov-cubic-moduli-morphism}の証明]
橋渡し命題の十七個の係数の函数体を
$K=\mathbb C(a,b,\alpha,\beta,m_{ij},g_j)$ とし、$R=K[[h]]$ に底変更する。
$d(0)\ne0$ なのでこれは許される。
特殊ファイバーは滑らかな三次 $X_0=V(F)$ である。
定理\ref{thm:relative-kuznetsov-projector} と
系\ref{cor:relative-kuznetsov-serre} により、
残余圏は係数 perfect な代数 $E_R$ で表示され、底変更に両立する相対 Serre 自己同値を持つ。
従って命題\ref{prop:kuznetsov-effective-serre-embedding}を適用できる。
補題\ref{lem:relative-kuznetsov-nonzero-euler}の
$\chi(G_2,G_2)=-9$ は一般ファイバーでも成立するので、その残余圏は非零である。
系\ref{cor:kuznetsov-formal-to-generic}から
\begin{equation}\label{eq:kuznetsov-geometric-generic}
 \mathcal K_{K((h))}\simeq\operatorname{Ku}(X_{K((h))})
\end{equation}
を得る。

橋渡し命題の可逆な係数変換により、十七次元 atlas の函数体
$K_{17}$ は $K((h))$ に単射として入り、左辺はその一般点の残余圏の底変更である。
補題\ref{lem:kuznetsov-spread-geometric-model}を atlas に適用すると、
非空開 $U$ 上に étale 局所の三次模型と相対同値、射
$\rho_U:U\to M_{\mathrm{cub}}$ を得る。
残余圏が相対 smooth proper であることと、
系\ref{cor:relative-kuznetsov-smooth-equivalence}から全 compact tails も相対 smooth proper となる。

étale atlas $T\to\mathfrak C_{\mathrm J}$ における $U$ の像を $\mathfrak U$ とする。
étale 射は開なので、これは非空開部分 stack である。
$U\times_{\mathfrak U}U$ は étale 局所的に滑らかで被約である。
残余圏は定理\ref{thm:jacobian-residual-dg-map}により O-marked 同値に沿って降下する。
従って二重交差上の二つの粗モジュライ写像は、Torelli と被約性により一致する。
étale 降下により\eqref{eq:kuznetsov-seventeen-ten}を得る。
同じ議論は、得られた射の各幾何点での特徴付けと一意性を与える。

支配性を確認する。$\rho$ の像の閉包を $Z\subset M_{\mathrm{cub}}$ とする。
式\eqref{eq:kuznetsov-geometric-generic}と Torelli により、
$R$ 上の三次 $X_R$ の一般点は $Z$ に値を取る。
$R$ は整域で $Z$ は閉なので、その特殊点 $[X_0]=[V(F)]$ も $Z$ に入る。
補題\ref{lem:kuznetsov-classical-dominance}により、この特殊点は
通常の三次モジュライの一般点を支配する。
従って $Z=M_{\mathrm{cub}}$ であり $\rho$ は支配的である。

最後に target を自己同型が自明な開 $M_{\mathrm{cub}}^\circ$ に制限する。
支配性により source の逆像は非空である。
source と target は標数零で滑らかであるため、generic smoothness により
さらに非空開へ縮小すれば $\rho$ は滑らかになる。
その像 $V$ は開であり、そこへの射は全射で相対次元 $17-10=7$ である。
$M_{\mathrm{cub}}^\circ$ はモジュライ stack と一致するので、stack への射も得る。
三次族と残余同値については、構成した étale 被覆上の相対模型を用いる。
\end{proof}

\begin{remark}[粗モジュライ、族、同値kernelの区別]
自己同型を持つ三次を含む粗モジュライへの射について、点ごとの Torelli は
被約な二重交差上の写像の一致に使った。
任意基底上の categorical Torelli や全自己同値群の一致を仮定したのではない。
自己同型が自明な開へ縮小すれば、普遍三次族の引戻しが source 上に存在する。
その相対成分との Morita 同値 kernel は étale 局所に構成しており、
source 全体で単一の kernel を選べることや、その選択の一意性は主張しない。
大域的な $\mathcal O(1)$ の選択も不要である。
\end{remark}

\begin{remark}[主張の確実性と既知入力]
本節の証明は、既存の相対 PBW/CY/proper/Serre 定理と代数的 chart を入力とする。
Hochschild 計算、infinitesimal Torelli、Lurie の形式 hull 比較、
Bhatt の perfect complex の効果化、Stevenson の特異圏の分類、
Orlov 比較、categorical Torelli、perfect kernel の連続性は既知の結果である。
それらの適用範囲を明記して、実際の十七係数族の一般点の幾何性とモジュライ射を証明した。
この具体的な組合せと応用の初出は未確認であり、独立の査読を終えたものではない。
一般の proper dg 族における compact kernel の形式効果性や
生成性の開性は使っていない。これらに対する反例と適用範囲の境界は研究ノートに保持する。
\end{remark}


\clearpage
\part{全Artin変形とPoisson幾何}
\section{中心三次対の全Artin形式局所構造と二次包絡}
\label{sec:moduli-artin}
本節は、滑らかな可換三次三様体の近傍において、
有限係数の17次元族と異なる\emph{完備形式局所モジュライ}を扱う。
生成空間の座標標識を残す関手と、その射影座標変更を割った関手を
区別する。構造代数層との比較に用いるのは後者であり、
特殊ファイバーの同一視は保持する。主張は通常の集合値Artin変形関手に関するものであり、
derived stackの同型でも非可換領域全体の大域分類でもない。
既知のKontsevich量子化から中心三次商の族を実際に構成し、
構造代数層の変形関手と比較する。
全ての次数付き代数変形を同次formalityで先に分類することは仮定しない。
構造代数層への移行ではYekutieliの量子化定理の仮定と
相対コホモロジー計算を明示する。

以下$k=\C$、$V=k^5$、$S=\Sym V^*$、
$F\in S_3$を滑らかな三次三様体
$X=V_+(F)\subset\PP(V)$の定義式とする。
$(R,\mathfrak m)$は剰余体$k$の局所Artin代数である。

\begin{lemma}[平坦な正規三次商は自動的に中心商]
\label{lem:artin-normal-central}
$B_R$を$R$上各次数有限自由な連結非負次数付き代数とし、
一次で生成され、特殊ファイバー$B_R/\mathfrak mB_R=S$とする。
$w\in(B_R)_3$が$F$に特殊化するなら$w$は左右正則であり、
以下の条件は同値である：
\[
 B_R/(w)\text{ が各次数で }R\text{-平坦},\qquad
 wB_R=B_Rw,\qquad
 w\in Z(B_R).
\]
成立するとき商のHilbert級数は$(1-t^3)/(1-t)^5$である。
\end{lemma}
\begin{proof}
$w$の左右乗法の各次数成分は、$\mathfrak m$で還元すると
$F:S_{d-3}\hookrightarrow S_d$である。
$F\ne0$と$S$が整域であることから単射となり、
有限自由行列の最大rank小行列式が単元であるため、
元の乗法も分裂単射になる。

$J=(w)$について商が平坦なら
$0\to J_d\to(B_R)_d\to(B_R/J)_d\to0$は
有限自由$R$加群の分裂完全列である。
$J_d/\mathfrak mJ_d=FS_{d-3}$なので、
左右乗法$(B_R)_{d-3}\to J_d$はNakayamaと
先の単射性により同型である。従って$J=wB_R=B_Rw$。
正規性から$wb=\sigma(b)w$で次数付き自己同型$\sigma$が得られる。
特殊ファイバーでは$\bar\sigma=\mathrm{id}$であり、
$b=w$と右消去によって$\sigma(w)=w$となる。
$\sigma\equiv\mathrm{id}\pmod{\mathfrak m^r}$とし、
$\mathfrak m^r/\mathfrak m^{r+1}=I$とおくと、
次の差分は次数零の導分
\[
 D:S\longrightarrow I\otimes_k S,\qquad
 D(F)=0
\]
である。$D=\sum_i\ell_i\partial_{x_i}$、
$\ell_i\in I\otimes S_1$だから
$\sum_i\ell_iF_i=0$。滑らかさから
$F_0,\ldots,F_4$は二次正則列をなし、
Koszul複体に一次係数のsyzygyが存在しない。
よって$D=0$となり、$\mathfrak m$の冪零性から
$\sigma=\mathrm{id}$を得る。
中心性から商平坦性へは、
各次数の分裂単射$w:(B_R)_{d-3}\to(B_R)_d$を使えばよい。
\end{proof}

\begin{corollary}[標識付き対の自己同型]
\label{cor:artin-scalar}
上の中心正規対$(B_R,Rw)$の次数付き自己同型が
特殊ファイバーで恒等なら、次数一の生成元への作用は
すべて同じ単元$\lambda\in1+\mathfrak m$によるスカラー倍である。
従って射影標識を入れた対には非自明な無限小自己同型はない。
\end{corollary}
\begin{proof}
$\phi(w)=cw$、$c\in1+\mathfrak m$と書き、
$c$の一意の三乗根による次数スカラー倍を合成すると
$w$を固定する自己同型に帰着する。
その$\mathfrak m$進差分は前補題の$D(F)=0$と同じ
一次Koszul syzygyで消える。
\end{proof}

\subsection{相対Casimir条件と五つのPl\"ucker方程式}
体積形式$\Omega=dx_0\wedge\cdots\wedge dx_4$を固定する。
$F_R\in R\otimes S_3$を$F$の持上げとし、
$a=\sum_{i<j}a_{ij}\,dx_i\wedge dx_j$を
係数$a_{ij}\in\mathfrak m$の定数二形式とする。
\[
 \iota_{\Pi(F_R,a)}\Omega=dF_R\wedge a
\]
で二次bivectorを定める。
$\{i,j,k,l\}=\{0,\ldots,4\}\setminus\{r\}$、
$i<j<k<l$に対し
\[
 q_r(a)=(-1)^r
 (a_{ij}a_{kl}-a_{ik}a_{jl}+a_{il}a_{jk}),\qquad
 g=\sum_{r=0}^4q_r(a)\,\partial_rF_R
\]
とおく。

\begin{lemma}[全Artin上でのCasimir核]
\label{lem:artin-casimir}
二次bivector$\Pi$が$\iota_{dF_R}\Pi=0$を満たすことと、
一意な$a\in R\otimes\wedge^2V^*$を用いて
$\Pi=\Pi(F_R,a)$と表せることは同値である。
この同型は任意の$R$の基底変換と可換であり、
$\Pi$は発散零である。
\end{lemma}
\begin{proof}
$F$の五つの偏微分は二次正則列だから、そのKoszul複体の
\[
 \wedge^3V\otimes R
 \xrightarrow{\iota_{dF_R}}
 S_{R,2}\otimes\wedge^2V
 \xrightarrow{\iota_{dF_R}}
 S_{R,4}\otimes V
\]
の該当次数は、特殊ファイバーで左端の単射と中央での完全性を持つ。
有限自由加群の複体なので、分裂小行列式の単元性を
Artin基底へ持ち上げて同じ完全性を得る。
$\wedge^3V\simeq\wedge^2V^*$により表示と一意性が従う。
$d(dF_R\wedge a)=0$から発散零である。
\end{proof}

\begin{proposition}[JacobiイデアルとPl\"uckerイデアルの一致]
\label{prop:artin-pfaffian}
$\Pi=\Pi(F_R,a)$に対し、Jacobi方程式の係数イデアルは
\[
 (q_0(a),\ldots,q_4(a))
\]
と一致する。特に集合論的なrank分類だけでなく、
非被約な局所Artin基底上での方程式を同定する。
\end{proposition}
\begin{proof}
$J=\frac12[\Pi,\Pi]$と書くと、五変数の外積計算から
\[
 \iota_{\Pi\wedge\Pi}\Omega=2g\,dF_R,\qquad
 \iota_J\Omega=dg\wedge dF_R
\]
が得られる。これは多項式恒等式である。
符号は$ a=dx_0\wedge dx_1+dx_2\wedge dx_3$
の代入で直接確認できる。
線形写像
\[
 L_F:k^5\longrightarrow\Omega^2_{S,3},\qquad
 (c_r)\longmapsto d\bigl(\textstyle\sum_rc_rF_r\bigr)\wedge dF
\]
は単射である。核の元について、
$F_i$のKoszul完全性から
$d(\sum c_rF_r)$は$dF$の多項式倍となるが、
係数の次数は負なので零である。
Euler式と定数係数のKoszul syzygyの消滅により$c_r=0$。
従って$L_F$の非零$5$次小行列式は
$R$上でも単元であり、$L_{F_R}$は分裂単射となる。
上の$J$の表示の係数から$q_r(a)$を回収でき、
両方のイデアルが一致する。
\end{proof}

\begin{lemma}[同次量子化による中心三次商の構成]
\label{lem:artin-homogeneous-casimir-quantization}
$\Pi\in\mathfrak m\otimes S_2\otimes\wedge^2V$を
$[\Pi,\Pi]=0$を満たす二次bivectorとする。
Kontsevichの多重線形写像を一つ固定すると、
$R\otimes S$上に次数を保つ単位的な星積$B_\Pi$が定まり、
$B_\Pi$は一次生成の二次代数である。
この構成は任意のArtin係数環の準同型と可換である。

twisted Taylor mapの次数三部分
$D^0_\Pi:S_{R,3}\to(B_\Pi)_3$は同型であり、
任意の$F_R\in S_{R,3}$について
\[
 D^0_\Pi(F_R)\in Z(B_\Pi)
 \quad\Longleftrightarrow\quad [\Pi,F_R]=0.
\]
特に$F_R\bmod\mathfrak m=F$および$q_r(a)=0$なら、
$\Pi=\Pi(F_R,a)$、$w_R=D^0_\Pi(F_R)$から
次数別有限自由な中心三次商$A_R=B_\Pi/(w_R)$を得る。
\end{lemma}
\begin{proof}
Kontsevichの写像$U_n$をCattaneo--Felder--Tomassiniの規約で取り、
$\alpha=\Pi/2$として
\[
 f\star g=fg+\sum_{n\ge1}\frac1{n!}
 U_n(\alpha,\ldots,\alpha)(f,g),\qquad
 D^0_\Pi(f)=\sum_{n\ge0}\frac1{n!}
 U_{n+1}(f,\alpha,\ldots,\alpha)
\]
と定める。$\Pi$の係数は$\mathfrak m$に入るので和は有限である。
結合律と単位はformality恒等式から従い、一次交換子は$\Pi$である。
ここで用いる正確な外部入力はCattaneo--Felder--Tomassini
\cite[\S3, Theorem~3.1, Proposition~3.3]{CFT}であり、
その式はテンソルの係数と微分についての普遍的多項式である。
従って実数から複素数へ係数拡大し、冪零Artin係数を代入できる。
\S3で述べられているGL同変性もこの係数拡大後に成立する。

各グラフでは二次bivectorの頂点から二本の辺が出て、
各辺は一回の微分を与える。
従って星積は通常の多項式次数を保ち、
$D^0_\Pi$は三次式を三次式へ送る。
基礎加群は$R\otimes S$なので次数別有限自由であり、
一次生成性は各次数でのNakayamaによる。
一次生成空間からのテンソル代数の核を$K$とすると、
$K_n$は有限自由な直和因子で、その特殊化は多項式環の関係イデアルである。
$K_2$から生成されるイデアルは特殊化後に各$K_n$を生成するから、
$K_n$へのNakayamaにより$K=(K_2)$となる。
基底変換に関する自然性は、上の有限の普遍式から従う。

$D^0_\Pi$の次数三部分
$D^0_\Pi:S_{R,3}\to(B_\Pi)_3$は
還元が恒等の有限自由加群の同型である。
また係数が四次のベクトル場$Y$に対して
\[
 D^1_\Pi(Y):=\sum_{n\ge0}\frac1{n!}
 U_{n+1}(Y,\alpha,\ldots,\alpha)
\]
とおく。その評価写像
\[
 Y\longmapsto
 \bigl(D^1_\Pi(Y)(x_0),\ldots,D^1_\Pi(Y)(x_4)\bigr)
\]
は$S_{R,4}\otimes V\to(B_\Pi)_4^{\oplus5}$の写像であり、
同じく還元が恒等の有限自由加群の同型である。
Cattaneo--Felder--Tomassini
\cite[Corollary~5.2(iv)]{CFT}のtwisted chain identityにより
\[
 D^0_\Pi(F_R)\in Z(B_\Pi)
 \quad\Longleftrightarrow\quad [\Pi,F_R]=0
\]
となる。ここで中心の形成と非平坦基底変換の交換性は
仮定せず、有限自由行列の可逆性によって逆向きを証明する。
命題\ref{prop:artin-pfaffian}により、$q_r(a)=0$なら
$\Pi(F_R,a)$はPoissonであり、定義から$[\Pi(F_R,a),F_R]=0$である。
よって$w_R$は中心元となり、補題\ref{lem:artin-normal-central}から
その正則性と商の次数別平坦性を得る。
この補題は具体的な族の構成であり、
次数付きHochschild変形と同次Poisson変形の全gauge分類を主張していない。
\end{proof}

\subsection{Ore局所化、偏極、および斉次二次包絡}
次数別平坦な斉次変形$A_R$を
$S/(F)$へ特殊化するものとする。
一次元の持上げ$\widetilde x_i$を選び、
標準開集合$U_i=D_+(x_i)\cap X$上で
\[
 \mathscr A_R(U_i)=
 \bigl(A_R[\widetilde x_i^{-1}]\bigr)_0
\]
を取る。$s=\widetilde x_i$、$\mathfrak m^N=0$とすると、
$[A_R,A_R]\subset\mathfrak m A_R$より
$\operatorname{ad}_s(\mathfrak m^rA_R)\subset\mathfrak m^{r+1}A_R$、
従って$\operatorname{ad}_s^N=0$である。
任意の$n\ge1$と$a\in A_R$に対し
\[
 s^{n+N-1}a=
 \sum_{j=0}^{N-1}\binom{n+N-1}{j}
 \operatorname{ad}_s^j(a)s^{n+N-1-j}\in A_Rs^n.
\]
逆側にも同じ展開があり、$\{s^n\}$は左右Ore条件を満たす。
$s$の特殊化は整域$S/(F)$の非零因子であり、
$A_R$の$\mathfrak m$進濾過の商における乗法も単射なので、
$s$は左右正則である。
局所化は$R$加群として$A_R$の乗法写像によるfiltered colimitだから平坦であり、
その各次数成分も直和因子として平坦である。
重なりは二つの一次元をともに可逆化した次数零部分で
記述され、三重重なりのcocycleも局所化の普遍性に従う。
この操作で$X$上に平坦な結合$R$代数層
$\mathscr A_R$が得られる。

\begin{proposition}[偏極とsection algebraの回復]
\label{prop:artin-ore-section}
上のOre局所化層の次数$n$部分を貼り合わせた
$\mathscr E_n$は可逆両側$\mathscr A_R$加群であり、
\[
 \mathscr E_m\otimes_{\mathscr A_R}\mathscr E_n
 \simeq\mathscr E_{m+n}.
\]
$n\ge0$について
\[
 H^i(X,\mathscr E_n)=0\ (i>0),\qquad
 (A_R)_n\xrightarrow{\sim}H^0(X,\mathscr E_n)
\]
が成り立ち、積とも両立し任意のArtin基底変換に可換である。
\end{proposition}
\begin{proof}
局所化の一次元は可逆なので強い$\mathbb Z$次数付けと
両側可逆性が従う。アフィン交叉上では、$\mathscr E_n$の有限な
$\mathfrak m$進濾過の各商は$\OO_X(n)$の有限直和となり、非負次数に限らず
高次コホモロジーを持たない。従って有限標準被覆の\v{C}ech複体で
大域コホモロジーを計算できる。
$R$の小拡大を順に帰納すると、
特殊ファイバーの$H^i(X,\OO_X(n))=0$ ($i>0,n\ge0$)
から高次消滅と$H^0$の有限自由性が得られる。
特殊化で同型な$(A_R)_n\to H^0(X,\mathscr E_n)$は
Nakayamaで同型となる。強次数付けにより積との両立も従う。
\end{proof}

\begin{lemma}[可逆両側偏極の同型類の一意性]
\label{lem:artin-polarization}
$\mathscr A_R$上の二つの平坦な可逆両側加群
$\mathscr M,\mathscr N$が同じ線束$L$に特殊化し、
$H^1(X,\OO_X)=H^0(X,T_X)=0$なら、
$\mathscr M\simeq\mathscr N$である。
$H^0(X,\OO_X)=k$なら各々の自己同型は$R^\times$である。
\end{lemma}
\begin{proof}
小拡大$R'\twoheadrightarrow R$の核$I$で
$\mathfrak mI=0$とする。
特殊ファイバーの局所基底の持上げにより
$\mathscr M,\mathscr N$は右加群として局所自由階数一。
同型の右加群としての持上げの障害は
$I\otimes H^1(X,\OO_X)$に属し消える。
二つの左作用の差は導分に対応し
$I\otimes H^0(X,T_X)=0$なので持上げは両側同型。
自己同型群は
$H^0(X,\OO_X)=k$と小拡大帰納から$R^\times$となる。
\end{proof}

\subsection{構造代数層の変形関手と座標標識の商}
\begin{lemma}[相対三次Poisson変形の形式局所環]
\label{lem:artin-relative-poisson-sheaf}
$\mathcal D_{\mathrm{Pois},X}(R)$を、$\OO_X$への還元同型を持つ
平坦な可換$R$代数層と、還元が零である$R$線形Poisson括弧の
gauge同値類の集合とする。するとこの関手は
\begin{equation}\label{eq:artin-poisson-sheaf-ring}
 \mathcal R_X=
 \frac{k[[t_1,\ldots,t_{10},a_{01},\ldots,a_{34}]]}
 {(q_0(a),\ldots,q_4(a))}
\end{equation}
で表される。ここでは特殊ファイバーの有限自己同型群$\Aut(X)$を割らない。
\end{lemma}
\begin{proof}
まずPoisson括弧を忘れる。滑らかなアフィン開集合上の
冪零な平坦可換変形とその貼合せにより、基礎の代数層は
$X$の滑らかなArtin変形$X_R$の構造層となる。
小拡大の核を$I$とすると、線束$\OO_X(1)$の持上げの障害は
$I\otimes H^2(X,\OO_X)$に、持上げの差は
$I\otimes H^1(X,\OO_X)$に属する。
両者が零なので、$\OO_X(1)$は各$X_R$へ同型を除いて一意に持ち上がる。
三次超曲面の標準完全列から
\[
 H^1(X,\OO_X(n))=H^2(X,\OO_X(n))=0\quad(n\in\mathbb Z)
\]
である。特に$n\ge0$では高次コホモロジーがすべて消えるので、
小拡大帰納により持ち上げた線束の非負次数のsectionは自由で、
任意のArtin基底変更と可換である。
その一次sectionの基底を選ぶと$X_R$は$\PP^4_R$に埋め込まれる。
実際、section環への写像は特殊ファイバーで全射なので
各次数でNakayamaにより全射であり、三次核は階数一である。
その生成元$F_R$の乗法は特殊ファイバーで単射だから分裂単射であり、
同じHilbert関数とNakayamaにより全核は$(F_R)$となる。
従って$X_R=V_+(F_R)$である。

相対随伴公式により
$\wedge^2T_{X_R/R}\simeq\Omega^1_{X_R/R}(2)$である。
相対余法完全列と制限完全列
\[
\begin{split}
0&\longrightarrow\OO_{X_R}(-1)
 \longrightarrow\Omega^1_{\PP^4_R/R}(2)|_{X_R}
 \longrightarrow\Omega^1_{X_R/R}(2)\longrightarrow0,\\
0&\longrightarrow\Omega^1_{\PP^4_R/R}(-1)
 \longrightarrow\Omega^1_{\PP^4_R/R}(2)
 \longrightarrow\Omega^1_{\PP^4_R/R}(2)|_{X_R}\longrightarrow0
\end{split}
\]
を用いる。左端二層の$H^0,H^1$は、超曲面完全列とEuler完全列から
零である。さらにEuler完全列の次数二部分から
\begin{equation}\label{eq:artin-relative-bivectors}
 H^0(X_R,\wedge^2T_{X_R/R})
 \simeq H^0(\PP^4_R,\Omega^1_{\PP^4_R/R}(2))
 \simeq R\otimes_k\wedge^2V^*
\end{equation}
を得る。これらは有限自由加群間の分裂完全列による計算なので、
非平坦なArtin基底変更とも可換である。
$F_R$の尺度と体積形式を固定した上で、この対応は
定数二形式$a$を$\iota_Ea$に送り、対応するbivectorは
$\iota_{\Pi}\Omega=dF_R\wedge a$の射影制限である。
これは偏微分$\partial_iF_R$が単元となる滑らかなchartで
随伴公式の縮約を計算して確認できる。

ここで$J=\frac12[\Pi,\Pi]$、$g=\sum_rq_r(a)\partial_rF_R$とする。
命題\ref{prop:artin-pfaffian}の恒等式をEuler場で縮約すると
\begin{equation}\label{eq:artin-intrinsic-jacobi}
 \iota_E\iota_J\Omega
   =2g\,dF_R-3F_R\,dg
\end{equation}
となる。$F_R=0$に制限しEuler方向を除くと、
内在的Jacobi障害は随伴同型
$\wedge^3T_{X_R/R}\simeq\OO_{X_R}(2)$の下で
$2g$に対応する（全体の符号は縮約の規約による）。
五つの二次式$\partial_rF_R$は$R\otimes S_2$内で
階数五の直和因子を生成する。
実際、特殊ファイバーでの定数係数の関係は
滑らかな$F$の偏微分正則列のKoszul完全性により零であり、
非零最大minorを$R$へ持ち上げればよい。
また$H^0(X_R,\OO_{X_R}(2))=R\otimes S_2$である。
従ってJacobi方程式の係数イデアルは、
集合論的にだけでなく$(q_0,\ldots,q_4)$と一致する。

座標を残すと、$[F_R]$の34係数と$a$の10係数に
この五つのPl\"ucker方程式を課した関手を得る。
滑らかな三次三様体は$H^0(X,T_X)=0$を満たすので、
特殊ファイバーで恒等な射影座標変更の作用は自由であり、
その24次元軌道を分離する形式sliceが存在する。
この作用は分解可能二形式の錐を保つので、$F_R$をsliceへ移す
一意の形式座標変更を$a$にも作用させれば、
残る方程式は同じPl\"ucker方程式である。
これが\eqref{eq:artin-poisson-sheaf-ring}を与える。
特殊ファイバーを恒等に保つ自己同型は、最初の非零な
$\mathfrak m$進差分が$H^0(X,T_X)$に属するため存在しない。
従ってここで必要なのは集合値のprorepresenting環である。
\end{proof}

$\mathcal D_{\mathrm{cp},X}(R)$を、次数別有限自由な二次代数$B_R$と、
その三次部分の階数一の直和因子で中心に含まれる$L_R$の対の同値類の集合とする。
特殊ファイバーの$(B_R,L_R)$と$(S,kF)$の同一視をデータに含め、
同値はこの同一視を保つ次数付き同型とする。
すなわち、一次生成空間の座標を選んだ表示から、
特殊ファイバーで恒等な射影座標変更を割った関手である。
次数スカラー倍による自己同型をrigidifyしても同値類の集合は変わらない。
$\mathcal D_{\mathrm{Ass},X}$は$\OO_X$への還元同型を保つ
平坦結合代数層のArtin変形のgauge同値類の関手とする。

\begin{lemma}[構成した中心商族と代数層変形関手]
\label{lem:artin-constructed-family-comparison}
補題\ref{lem:artin-relative-poisson-sheaf}の三次式の形式slice
$F_{\boldsymbol t}$を固定する。
$\mathcal R_X$のArtin点$(\boldsymbol t,a)$に対し、
補題\ref{lem:artin-homogeneous-casimir-quantization}の
$\Pi(F_{\boldsymbol t},a)$の量子化と中心三次商を取り、
さらにOre局所化を取ることで得る自然変換
\[
 \Psi:\operatorname{Hom}_{\mathrm{loc}}(\mathcal R_X,-)
       \longrightarrow\mathcal D_{\mathrm{Ass},X}
\]
は自然同型である。
\end{lemma}
\begin{proof}
量子化の有限な普遍式、中心元の構成、およびOre局所化の普遍性により、
$\Psi$は任意のArtin基底変換に自然な実際の写像である。
Yekutieli\cite[Corollary~11.2]{YekutieliTwisted}は、滑らかな$X$と
$H^1(X,\OO_X)=H^2(X,\OO_X)=0$の下で、
可換積も動かすPoisson代数層変形と結合代数層変形の
gauge同値類の自然な全単射を与える。
本稿では$k=\C\supset\mathbb R$であり、局所Artin代数は
同定理のパラメータ代数に含まれる。
通常の変形とtwisted変形の同値類の一致は
同論文Corollary~7.5により保証されている。
補題\ref{lem:artin-relative-poisson-sheaf}により、
標的$\mathcal D_{\mathrm{Ass},X}$も完備局所Noetherian環
$\mathcal R_X$で表される。

接写像の単射性を直接示す。
$R=k[\varepsilon]/(\varepsilon^2)$上の点$(\dot{\boldsymbol t},a)$が
$\Psi$によって自明な変形へ移るとする。
一次交換子はgauge不変であり、構成した星積から得る交換子の射影制限は
\[
 \wedge^2V^*\longrightarrow H^0(X,\wedge^2T_X),
 \qquad a\longmapsto\iota_Ea
\]
という\eqref{eq:artin-relative-bivectors}の同型である。
従って$a=0$である。
このとき量子化は通常の積であり、中心三次商とその局所化は
可換な三次変形$F_{\boldsymbol t}=0$そのものである。
これが結合代数層として自明なら、可換代数層としても自明だから、
三次式のsliceの接方向$\dot{\boldsymbol t}$は零である。
よって$d\Psi$は単射であり、両側の接空間は20次元なので同型である。
この議論は、Yekutieliの量子化が選んだ幾何的座標を保つことを仮定しない。

標的の表現環との同型を一つ選ぶと、$\Psi$は$\mathcal R_X$の
連続自己準同型を与え、余接空間上で同型となる。
随伴次数環は次数一で生成されるから、各次数成分への写像は全射である。
それらは同じ有限次元空間の自己写像なので同型となり、
完備性と分離性から元の環準同型も同型である。
従って実際の$\Psi$が全Artin基底で自然同型となる。
\end{proof}

\begin{theorem}[斉次環の全Artin変形の二次cover]
\label{thm:artin-quadratic-cover}
Ore局所化による自然変換
\[
 \Theta:\mathcal D_{\mathrm{cp},X}
       \longrightarrow\mathcal D_{\mathrm{Ass},X}
\]
は自然同型となる。
従って任意の次数別平坦な斉次変形$A_R$は
\[
 V_R=(A_R)_1,\quad
 K_{2,R}=\ker(V_R^{\otimes2}\to(A_R)_2),\quad
 B_R^{\mathrm{quad}}=T_R(V_R)/(K_{2,R})
\]
を用いて一意の中心三次対から回復される。
特に$B_R^{\mathrm{quad}}$は次数別$R$平坦で、
\[
 L_R=\ker\bigl((B_R^{\mathrm{quad}})_3\to(A_R)_3\bigr)
 \subset Z(B_R^{\mathrm{quad}})_3
\]
は階数一の直和因子であり、
$A_R\simeq B_R^{\mathrm{quad}}/(L_R)$。
この構成はArtin基底変換と可換である。
\end{theorem}
\begin{proof}
$\Theta$は上のOre局所化の構成により定義され、
特殊ファイバーを保つ次数付き同型は局所化層のgauge同型を誘導する。
補題\ref{lem:artin-constructed-family-comparison}の$\Psi$は、
具体的な中心三次対の族を経由するので、$\Theta$は全射である。

単射性を示すため、二つの中心対$(B_R,L_R)$、$(B'_R,L'_R)$の
商$A_R$、$A'_R$の局所化層がgauge同型であるとする。
この同型によって同じ代数層上で考えると、
両者の局所化次数一部分は$\OO_X(1)$の可逆両側持上げであり、
補題\ref{lem:artin-polarization}によって両側同型となる。
特殊ファイバーの同型はスカラーを調整して恒等に選べる。
この同型のテンソル冪と命題\ref{prop:artin-ore-section}により、
積と特殊ファイバーの標識を保つ次数付き同型$A_R\simeq A'_R$を得る。
中心三次商は二次関係を変えないので、$B_R$と$B'_R$はそれぞれの
商の二次包絡である。従って上の同型は$B_R\simeq B'_R$へ一意に持ち上がり、
三次核$L_R$を$L'_R$へ送る。これで$\Theta$の単射性を得る。

任意の斉次変形$A_R$についても、そのOre局所化層に
$\Theta$の逆像となる中心対$(B_R,L_R)$を取る。
局所化次数一の両側持上げの一意性とsection回復により、
$A_R$はこの中心対の商と次数付き同型になる。
この比較を行ってから二次包絡を取れば
$B_R^{\mathrm{quad}}\simeq B_R$を得るので、その平坦性は未知のまま仮定していない。
三次核は$L_R$であり、主張する復元が従う。
基底変換は、有限自由な次数一・二の加群間の分裂核と
局所化の普遍性により可換となる。

Ore局所化の一般的な参照先はYekutieli\cite{YekutieliAffine}である。
De~Deken--Lowen\cite[Theorem~4.5]{DeDekenLowen}は
Abel変形とsection $\mathbb Z$-代数変形を比較するが、
同定理だけから次数付き構造や1周期性を推論してはいない。
全Abel変形への移行は定理\ref{thm:abel-cubic-flat-envelope}で、
各線束対象の持上げの一意性を用いて別に行う。
\end{proof}

\begin{theorem}[中心三次対の全Artin形式局所環]
\label{thm:artin-central-pair}
中心三次対の古典的変形関手$\mathcal D_{\mathrm{cp},X}$は
\[
 \operatorname{Spf}\mathcal R_X
 =\operatorname{Spf}\frac{
 k[[t_1,\ldots,t_{10},a_{01},\ldots,a_{34}]]
 }{(q_0(a),\ldots,q_4(a))}
\]
で表され、形式局所次元は17、接空間次元は20である。
一次生成空間の射影座標標識を残す関手は、非標準的に
\[
 \widehat{\mathcal I}_{F}\simeq
 \operatorname{Spf}\frac{
  k[[s_1,\ldots,s_{34},a_{01},\ldots,a_{34}]]
 }{(q_0(a),\ldots,q_4(a))}
\]
で表され、形式局所次元は41、接空間次元は44である。
有限群$\Aut(X)$はこの記述では割らない。
\end{theorem}
\begin{proof}
最初の主張は定理\ref{thm:artin-quadratic-cover}、
YekutieliのCorollary~11.2、および
補題\ref{lem:artin-relative-poisson-sheaf}から従う。
これは中心対の全ての変形を同次formalityで分類することを
前提としない証明である。

射影座標標識を忘れるときの群を
\[
 G_1(R)=\ker\bigl(\PGL_5(R)\longrightarrow\PGL_5(k)\bigr)
\]
とする。標識は自由な一次生成空間の射影基底だから、
その選択に$G_1(R)$が推移的に作用する。
系\ref{cor:artin-scalar}により中心対の無限小自己同型は
次数スカラー倍だけであり、射影基底を変えない。
従って標識を残す関手は$\mathcal D_{\mathrm{cp},X}$上の
$G_1$-torsorである。
補題\ref{lem:artin-constructed-family-comparison}で構成した族の
固定した生成空間$R\otimes V^*$がこのtorsorのsectionを与える。
$G_1$は24次元の滑らかな形式群なので、
標識を残す関手の表現環は$\mathcal R_X[[u_1,\ldots,u_{24}]]$である。
$t_1,\ldots,t_{10},u_1,\ldots,u_{24}$を$s_1,\ldots,s_{34}$と書き換えればよい。

Pl\"ucker環は$\mathrm{Gr}(2,5)$のアフィン錐の原点完備化であり、
その次元は7、接空間次元は10である。
従って$10+7=17$、$10+10=20$、
座標標識を残すと$34+7=41$、$34+10=44$を得る。
\end{proof}

\begin{remark}[一次方向はすべて二次へ延長するわけではない]
\label{rem:artin-obstruction}
$a=\varepsilon(dx_0\wedge dx_1+dx_2\wedge dx_3)$を
$k[\varepsilon]/(\varepsilon^2)$上で考えると
$q_r(a)=0$だが、
$k[\varepsilon]/(\varepsilon^3)$上では$q_4(a)=\varepsilon^2$。
二次の補正$b$を$a+\varepsilon^2b$に加えてもこの項は消えない。
従って接空間20次元と局所次元17の差は実際の二次障害を表し、
全Artin比較を無障害性と混同してはいけない。
\end{remark}

\begin{remark}[明確に主張しない事柄]
\label{rem:artin-limits}
この全Artin同型は偏極付き\emph{古典的}変形関手の同定であり、
derived moduli stack、高次自己同型を保つgroupoid、
構造対象を忘れたKuznetsov成分の変形、
またはそこで構成される可換復元との一致を与えない。
無限小の中心対を得ることと、有限非可換パラメータで
滑らかなtailsを持つことは別問題である。
\end{remark}

\subsection{固定ambientからの忘却の四次元ファイバー}
\begin{proposition}[中心三次直線の射影束]
\label{prop:artin-fixed-ambient-fiber}
ある良い開集合$U$の各点$B$が二次Koszul AS正則CY五次元整域であり、
$\dim Z(B)_1=2$、$\dim Z(B)_3=5$かつ
中心の形成が基底変更と可換であると仮定する。
$E_1=Z(\mathscr B)_1$、$E_3=Z(\mathscr B)_3$とおけば
\[
 0\longrightarrow\Sym^3E_1\longrightarrow E_3
 \longrightarrow L\longrightarrow0,\qquad\rank L=1.
\]
中心三次直線のパラメータ空間は
$\PP_U(E_3)\to U$で各fiberは$\PP^4$。
$\PP(E_3)\setminus\PP(\Sym^3E_1)$は
$\Sym^3E_1\otimes L^{-1}$を作用群とするtorsorである。
同型類を取る場合、$\Gamma_B=\Aut_{\gr}(B)/k^\times$
が有限なら一般fiberの粗い双有理型は
$\PP Z(B)_3/\Gamma_B$となり、その次元は4である。
\end{proposition}
\begin{proof}
中心三次部のベクトル束の射影化で前半を得る。
$B$が整域なので中心三次の非零元は正則、
次数ごとの乗法の核は零で、商の各次数が局所自由となる。
$E_3\to L$に非零に写る直線はこの短完全列のsplittingに対応し、
splittingの差は$\Hom(L,\Sym^3E_1)$に属する。
同型類については、二次包絡と三次核の復元から
$B/(W)\simeq_{\gr}B/(W')$と
$[W'],[W]$が$\Gamma_B$軌道に属することが同値である。
Rosenlichtの有理商を$U$とincidenceへ適用し、
周囲$B$の軌道の上の空間が
$\PGL(V)\times_{\Gamma_B}\PP Z(B)_3$となることから
不変関数体が$k(\PP Z(B)_3)^{\Gamma_B}$に一致する。
有限群商は次元を保つ。
\end{proof}

\begin{remark}[四係数は標識を忘れても消えない]
特に固定した周囲に対し$Z(B)_3=kG\oplus\Sym^3Z(B)_1$なら、
$G+f(u,v)$の四係数が相対アフィンfiberを表す。
有限群$\Gamma_B$による同一視は生じ得るが、
四係数全体を座標変換で消せるわけではない。
原稿中の明示剛体周囲$B_*$についての
$\Aut_{\gr}(B_*)=k^\times$と正則開集合の主張は、
付属の有限整数証明書が現行Git treeに存在しないため、
本稿では独立に再検証した定理として採用しない。
\end{remark}

\section{可換三次点におけるAbel圏の全Artin変形}
\label{sec:abel-formal-cubic}

本節では，特殊ファイバーを固定した無限小変形を扱う．
大域的な無標識Abel圏の復元，有限型atlasの存在，または
任意の変形のAS ambientへの埋込みは結論しない．
比較の一般定理はDe Deken--Lowenの既知結果であり，
ここで行うのは三次三様体に対する仮定の確認と，
構造対象を標識する操作への帰結の明示である．

\subsection{変形の射と特殊ファイバーの固定}

$k=\mathbb C$とし，$X\subset\mathbb P^4_k$を滑らかな三次三様体とする．
$\mathcal C_0=\operatorname{Qcoh}(X)$，$E_i=\mathcal O_X(-i)$と置く．
section $\mathbb Z$-代数は
\[
 a_{ij}=\operatorname{Hom}_{\mathcal C_0}(E_j,E_i)
       =H^0(X,\mathcal O_X(j-i))
\]
である．特に$i>j$では$a_{ij}=0$であり，
$a$は$k[x_0,\ldots,x_4]/(F)$の付随$\mathbb Z$-代数に等しい．

$R$を剰余体$k$の可換局所Artin $k$-代数とする．
$\operatorname{Def}_{\rm ab}(\mathcal C_0;R)$の対象を，
Lowen--Van den Berghの意味で平坦な$R$-線型Abel圏$\mathcal D$と，
その$k$-対象の部分圏$\mathcal D_k$への同値
$\iota:\mathcal C_0\xrightarrow{\sim}\mathcal D_k$とする．
射は$R$-線型圏同値$F:\mathcal D\to\mathcal D'$で
$F\iota\simeq\iota'$を満たすものの自然同型類である．
この両立する自然同型そのものは追加frameとして保持しない．
これは\cite[\S8.2]{LowenVdBDeformation}の1群oidの規約である．
従って，完全な2群oidにあるスカラー2自己同型は，この定義では既に切断されている．
同様に$\operatorname{Def}_{\rm lin}(a;R)$では，Homが$R$-平坦な
線型圏の変形と，還元後の$a$への同値を対象とし，射を線型圏同値の
自然同型類とする．添字を固定した模型に取り換えた後も，
頂点スカラーによる自然同型を射から除いてはならない．

\begin{lemma}[三次三様体の全捩れに対する消滅]
\label{lem:abel-cubic-acyclicity}
任意の整数$m$について
\[
 H^1(X,\mathcal O_X(m))=H^2(X,\mathcal O_X(m))=0.
\]
従って，任意の$i,j\in\mathbb Z$と有限次元$k$-ベクトル空間$M$に対し
\[
 \operatorname{Ext}^q_{\mathcal C_0}(E_j,M\otimes_k E_i)=0
 \qquad(q=1,2)
\]
である．また各$E_i$は例外対象である．
\end{lemma}
\begin{proof}
超曲面の完全列
\[
 0\longrightarrow\mathcal O_{\mathbb P^4}(m-3)
 \xrightarrow{\ F\ }\mathcal O_{\mathbb P^4}(m)
 \longrightarrow\mathcal O_X(m)\longrightarrow0
\]
と，$\mathbb P^4$の線束の中間コホモロジー
$H^q(\mathbb P^4,\mathcal O(r))=0$（$0<q<4$）の消滅から従う．
線束同士のExtは捩れ線束のコホモロジーに等しく，$M$とのテンソル積は
有限直和なのでExtの式を得る．
$H^0(X,\mathcal O_X)=k$であり，
$\omega_X\simeq\mathcal O_X(-2)$とSerre双対性から
$H^3(X,\mathcal O_X)=H^0(X,\mathcal O_X(-2))^*=0$である．
これにより$\operatorname{RHom}(E_i,E_i)\simeq k$を得る．
\end{proof}

\begin{theorem}[既知の変形比較の三次三様体への適用]
\label{thm:abel-cubic-all-artin}
すべての局所Artin $k$-代数$R$について，1群oidの同値
\[
 \operatorname{Def}_{\rm lin}(a;R)
 \simeq\operatorname{Def}_{\rm ab}(\operatorname{Qcoh}(X);R)
 \simeq\operatorname{Def}_{\rm ab}(\operatorname{coh}(X);R)
\]
がある．第一の同値は$\widetilde a\mapsto\operatorname{Qmod}(\widetilde a)$で与えられる．
任意の$\mathcal D$について，各$E_i$には同型を除いて一意な$R$-平坦な持上げ
$\widetilde E_i$があり，逆方向の模型を
\[
 \widetilde a_{ij}=
 \operatorname{Hom}_{\mathcal D}(\widetilde E_j,\widetilde E_i)
 \quad(i\leq j),\qquad \widetilde a_{ij}=0\quad(i>j)
\]
として取れる．その各成分は有限自由で，階数は
\[
 h_{j-i}=\binom{j-i+4}{4}-\binom{j-i+1}{4}
 \quad(j\geq i)
\]
である．ここで下段が上段より大きい二項係数は零とする．
$\widetilde a$は次数一生成であり，局所有限，正，連結，右noetherianである．
$\mathcal D$はlocally noetherianで，そのnoetherian対象の圏は
$\operatorname{qgr}\widetilde a$に同値である．

これらの同値は，任意の局所Artin $k$-代数の射による基底変更と両立する
変形擬関手の同値として成り立つ．ここでAbel圏の基底変更には
新しい係数環の作用を備えた対象の圏を用いる．
\end{theorem}
\begin{proof}
$\mathcal O_X(1)$は豊富であり，$(\mathcal O_X(n))_{n\in\mathbb Z}$は
$\operatorname{Qcoh}(X)$の$\mathbb Z$-generating sequenceをなす．
補題\ref{lem:abel-cubic-acyclicity}は
\cite[Theorem~4.5, pp.~495--496]{DeDekenLowen}のExt消滅仮定を満たす．
同定理により第一の同値とその逆構成を得る．
ここでは同論文\S5のgeometric helix仮定は使用しない．
持上げの存在と同型類の一意性は，対象の障害が$\operatorname{Ext}^2$に，
存在する場合の持上げの同型類の差が$\operatorname{Ext}^1$にあること
\cite[Theorems~6.11--6.12の双対版]{LowenObstruction}からも得られる．

各Homの平坦性と還元の同定は同じ比較定理に含まれる．
有限性の持上げは\cite[Proposition~4.6]{DeDekenLowen}による．
$R$はArtin局所環なので，有限平坦なHomは有限自由であり，その階数は
特殊ファイバーのHilbert関数である．
また，隣接する次数一成分の積から任意の$\widetilde a_{ij}$への写像は
特殊ファイバー上で全射である．余核を$C$とすると$C/\mathfrak m_R C=0$であり，
$\mathfrak m_R$の冪零性から$C=0$となる．従って次数一生成も成り立つ．

noetherianな$\mathbb Z$-代数のtailsはlocally noetherianで，
そのnoetherian対象は有限生成加群からの像である．
よって$\mathcal D\simeq\operatorname{Qmod}(\widetilde a)$のnoetherian対象は
$\operatorname{qgr}\widetilde a$に等しい．
最後に，Indと有限表示対象を取る操作の変形擬関手の同値
\cite[Theorem~8.8, p.~5472]{LowenVdBDeformation}を用いると第二の同値を得る．
冪零核を持つ全射についての基底変更の両立は，引用した変形擬関手の同値に含まれる．
任意の局所Artin射への拡張は次の補題による．
\end{proof}

\begin{lemma}[Artin係数の任意の変更]
\label{lem:abel-artin-base-change}
$b$を定理の$R$-線型$\mathbb Z$-代数の変形とし，$R\to S$を
局所Artin $k$-代数の任意の射とする．$\mathcal D_S$で
$R$-線型Abel圏$\mathcal D$内の$S$-対象の圏を表すと，標準的な同値
\[
 \operatorname{Qmod}(S\otimes_R b)
 \simeq \operatorname{Qmod}(b)_S
\]
がある．この同値は射の合成と両立する．
\end{lemma}
\begin{proof}
局所Artin環の射なので$S$は$R$-加群として有限表示である．
$\mathcal D$が$R$-平坦なら$\mathcal D_S$は$S$-平坦である
\cite[Proposition~4.8(5), p.~5454]{LowenVdBDeformation}．
$b$と$S\otimes_R b$はいずれも右noetherianだから，そのQmodは
加群圏をtails torsionで割ったGabriel商として実現される．
商関手$L:\operatorname{Mod}(b)\to\operatorname{Qmod}(b)$と
その完全忠実な右随伴$\omega$に$S$の作用を移すと，随伴
$L_S\dashv\omega_S$を得る．$L_S$は完全で$\omega_S$は完全忠実なので，
$L_S$もその核によるGabriel商である．

一方，$\operatorname{Mod}(b)_S\simeq\operatorname{Mod}(S\otimes_R b)$
\cite[Proposition~4.4]{LowenVdBDeformation}である．
この同定の下で$L_S$の核は，基礎$b$-加群がtails torsionである
$S\otimes_R b$-加群からなる．各斉次元について，十分高い次数の$b$の作用が
零であることと，十分高い次数の$S\otimes_R b$の作用が零であることは同値である．
従ってこの核は$S\otimes_R b$のtails torsionに等しく，所望の同値を得る．
これらの構成は係数の作用を移す標準的な構成なので，射の合成と両立する．
$\operatorname{coh}$側の基底変更との両立は，有限表示対象を取る操作が
係数変更と可換であること\cite[Proposition~4.6]{LowenVdBDeformation}による．
\end{proof}

\begin{remark}[相対化の範囲]
引用した変形擬関手の定義域は$\operatorname{Rng}^0/k$，すなわち
冪零核を持つ全射である．上の補題は，本節のnoetherianなtails模型では
任意の局所Artin射にも比較が自然であることを別途確認したものである．
これは一般のスキーム基底上でのfpqc降下や，有限型代数スタックの存在までは意味しない．
\end{remark}

\subsection{関係式と正規三次条件への還元}

以下では，定理\ref{thm:abel-cubic-all-artin}から得られた$R$-線型
$\mathbb Z$-代数の変形を$b$と書く．$V_i=b_{i,i+1}$とし，
次数一の成分からなるtensor $\mathbb Z$-代数$T=T_R(V_\bullet)$を
\[
 T_{ij}=V_i\otimes_R\cdots\otimes_R V_{j-1}\quad(i<j),
 \qquad T_{ii}=R,\qquad T_{ij}=0\quad(i>j)
\]
で定める．次数一生成から標準的全射$T\twoheadrightarrow b$があり，
その核を$K$とする．二次関係の族を
$\mathcal R_i=K_{i,i+2}$と書き，二次包絡を
\[
 B=T/(\mathcal R_i\mid i\in\mathbb Z)
\]
と定める．以下の主張は，有限区間への制限による生成元・関係式の有限個化や，
添字の周期性を主張するものではない．

\begin{proposition}[二次関係と一つの三次関係による提示]
\label{prop:abel-quadratic-cubic-presentation}
各$V_i$は自由rank~5，各$\mathcal R_i$は自由rank~10である．
各$i$について$f_i\in K_{i,i+3}$を適切に選ぶと，
\[
 K=(\mathcal R_i,f_i\mid i\in\mathbb Z),
 \qquad b\simeq T/(\mathcal R_i,f_i\mid i\in\mathbb Z)
\]
となる．従って$g_i$を$f_i$の$B_{i,i+3}$への像とすれば，
$J=\ker(B\twoheadrightarrow b)$は$(g_i)_{i\in\mathbb Z}$で生成される．
特に$J_{i,i+3}$はcyclic $R$-加群である．
この段階では，$B$の平坦性や$J_{i,i+3}$の自由rank~1性は主張しない．
\end{proposition}
\begin{proof}
各$T_{ij}$と$b_{ij}$は有限自由であり，$T_{ij}\to b_{ij}$は全射なので，
完全列$0\to K_{ij}\to T_{ij}\to b_{ij}\to0$は$R$-加群として分裂する．
従って$K_{ij}$も有限自由で，その特殊化は特殊ファイバーの関係空間と一致する．
特に$\operatorname{rank}\mathcal R_i=5^2-15=10$である．

各$V_i$の還元を$k^5$と同定し，特殊ファイバーのtensor代数において
三次多項式$F$の順序付き単項式による代表$F_i\in(T\otimes_Rk)_{i,i+3}$を選ぶ．
$F_i$は$K\otimes_Rk$に属するので，これを$f_i\in K_{i,i+3}$へ持ち上げる．
$L\subset K$を$\mathcal R_i$と$f_i$で生成される両側idealとする．
特殊ファイバーの関係idealは可換性の二次関係と各位置の$F_i$で生成されるから，
すべての$i\leq j$について
\[
 L_{ij}\otimes_R k\longrightarrow K_{ij}\otimes_R k
\]
は全射である．従って$(K/L)_{ij}/\mathfrak m_R(K/L)_{ij}=0$であり，
$\mathfrak m_R$の冪零性から$(K/L)_{ij}=0$を得る．

$B$では二次関係の像は零なので，$J$は$g_i$で生成される．
次数三の位置$(i,i+3)$には他の位置の三次生成元を掛けて寄与させることはできず，
$J_{i,i+3}=Rg_i$となる．ただしcyclicであるだけでは$R$-平坦とは限らない．
\end{proof}

\begin{proposition}[次数三の障害idealと有限行列]
\label{prop:abel-degree-three-obstruction}
$I_i=\operatorname{Ann}_R(g_i)$と置くと，$I_i\subset\mathfrak m_R$であり，
非標準的な$R$-加群の同型
\[
 J_i\simeq R/I_i,
 \qquad B_{i,i+3}\simeq R^{34}\oplus R/I_i
\]
がある．特に$B_{i,i+3}$が$R$-平坦であるための必要十分条件は$I_i=0$である．

二次関係を次数三へ延長する写像
\[
 d_i:(V_i\otimes_R\mathcal R_{i+1})
       \oplus(\mathcal R_i\otimes_R V_{i+2})
       \longrightarrow T_{i,i+3},
 \qquad d_i(v\otimes r,r'\otimes v')=vr+r'v'
\]
を考える．各加群の基底を選ぶと，これは$125\times100$行列であり，
その特殊化のrankは$90$である．特殊化後非零となる$90\times90$小行列を
pivotとして消去すれば，$d_i$は可逆な行・列操作により
\[
 \begin{pmatrix}I_{90}&0\\0&D_i\end{pmatrix},
 \qquad D_i\in M_{35\times10}(\mathfrak m_R)
\]
に変形できる．さらに残りの$35$次元自由加群の基底を適切に選ぶと，
\[
 D_i=
 \begin{pmatrix}
  \lambda_{i,1}&\cdots&\lambda_{i,10}\\
  0&\cdots&0\\
  \vdots&&\vdots\\
  0&\cdots&0
 \end{pmatrix},
 \qquad
 I_i=(\lambda_{i,1},\ldots,\lambda_{i,10})
\]
となる．従って$I_i$は高々$10$元で生成される．
このidealは基底やpivotの選択によらず，より内在的には
\[
 I_i=\operatorname{Fitt}_{34}(B_{i,i+3})
\]
である．すなわち，元の$125\times100$行列$d_i$の
$91\times91$小行列式が生成するidealに等しい．
\end{proposition}

\begin{proof}
$b_{i,i+3}$は自由rank~34なので，完全列
\[
 0\longrightarrow J_i\longrightarrow B_{i,i+3}
 \longrightarrow b_{i,i+3}\longrightarrow0
\]
は$R$-加群として分裂する．$J_i=Rg_i$より$J_i\simeq R/I_i$であり，
所望の直和表示を得る．$g_i$の特殊化は非零な三次式$F$なので，
$I_i$は単元を含まず$I_i\subset\mathfrak m_R$である．
$B_{i,i+3}$が平坦であることと$R/I_i$が平坦であることは同値である．
後者は有限平坦加群だから自由であり，その剰余体上の次元は$1$なのでrank~1である．
このとき全射$R\to R/I_i$は同じ長さの$R$-加群間の全射となり，$I_i=0$である．
逆は明らかである．

次に行列表示を示す．$\operatorname{rank}V_i=5$，
$\operatorname{rank}\mathcal R_i=10$であるから，$d_i$の始域は自由rank~100，
終域は自由rank~125である．二次idealの次数三成分は$d_i$の像なので，
\[
 B_{i,i+3}=\operatorname{coker}(d_i).
\]
特殊化後の二次包絡は多項式環$S=k[x_0,\ldots,x_4]$であるから，
$\operatorname{coker}(d_i\otimes_R k)=S_3$は次元$35$を持ち，
$d_i\otimes_R k$のrankは$125-35=90$である．

特殊化後非零となる$90\times90$小行列の行列式は$R$の単元である．
この部分を可逆な行・列操作で単位行列にし，その上下・左右の成分を消去すると，
主張のblock表示を得る．残余行列$D_i$の特殊化はrank~0なので，
すべての成分が$\mathfrak m_R$に属する．

この消去によって取り除いた$90$本の像の直和因子を$U\subset T_{i,i+3}$とし，
$P=T_{i,i+3}/U\simeq R^{35}$と置く．$U$は二次関係の像からなるため，
全射$T_{i,i+3}\to b_{i,i+3}$は全射
\[
 \psi:P\longrightarrow b_{i,i+3}\simeq R^{34}
\]
を誘導する．これは分裂し，その核$H$は自由rank~1である．
$f_i$の$P$への像を$\widetilde g_i$と書くと，
$\widetilde g_i\in H$であり，特殊化後には$S_3$の中の非零元$F$を与える．
従って$\widetilde g_i$は$H$の生成元である．

$D_i$の各列は二次関係の像であるから$\ker\psi=H=R\widetilde g_i$に入る．
よって一意的な$\lambda_{i,a}\in R$を用いて，その第$a$列は
$\lambda_{i,a}\widetilde g_i$と書ける．
$H$は$P$の直和因子なので，$\widetilde g_i$を$P$の基底の最初の元として取り，
これを$\psi$の分裂から得られる$34$個の元で補えば，表示した一行行列を得る．
特に
\[
 \operatorname{im}(D_i)
   =(\lambda_{i,1},\ldots,\lambda_{i,10})\widetilde g_i,
 \qquad
 J_i=H/\operatorname{im}(D_i)
\]
である．従って$I_i=(\lambda_{i,1},\ldots,\lambda_{i,10})$となる．
なお$D_i$の特殊化が零で$\widetilde g_i$がprimitiveであることから，
各$\lambda_{i,a}$も$\mathfrak m_R$に属する．

最後に，この$35$生成元・$10$関係の表示における
$\operatorname{Fitt}_{34}$は$D_i$の$1\times1$小行列式のidealであり，
従って$I_i$に等しい．Fitting idealの表示非依存性と，
元の$125$生成元・$100$関係の表示における定義から，
$d_i$の$125-34=91$次小行列式のidealによる記述も得る．
\end{proof}

\begin{remark}[計算上の抽出と論理的な範囲]
$35\times10$残余行列の各列は，$R^{35}$内の
primitiveな三次関係の持上げ$\widetilde g_i$の倍数である．
$\widetilde g_i$のある成分が単元なら，残余行列の対応する行を
その成分で割るだけで$\lambda_{i,1},\ldots,\lambda_{i,10}$を抽出できる．
残りの行が同じ係数による$\widetilde g_i$の倍数であることは，
$b_{i,i+3}$の自由性と上の核の記述から従う．

この命題は\emph{次数三の平坦性の必要十分条件}を与える．
すべての$i$で$I_i=0$を示しても，それだけで$B$の次数四以降の平坦性，
AS正則性，または全Artin Abel変形が三次incidenceに属することを結論してはならない．
\end{remark}

\begin{proposition}[二次包絡の平坦性から正規正則三次族へ]
\label{prop:abel-flat-envelope-normal-cubic}
前命題の二次包絡$B$について，すべての$B_{ij}$が$R$-平坦であると仮定する．
このとき
\[
 \operatorname{rank}_R B_{ij}=\binom{j-i+4}{4}\quad(i\leq j),
 \qquad
 \operatorname{rank}_R J_{ij}=
 \begin{cases}
  \binom{j-i+1}{4},&j-i\geq3,\\
  0,&j-i<3
 \end{cases}
\]
であり，これらは有限自由である．特に各$J_{i,i+3}$は自由rank~1である．
前命題の$g_i$はその生成元となり，すべての$j\geq i+3$で
\[
 \begin{aligned}
 B_{i+3,j}&\xrightarrow{\ \sim\ }J_{ij},& a&\longmapsto g_i a,\\
 B_{i,j-3}&\xrightarrow{\ \sim\ }J_{ij},& a&\longmapsto a g_{j-3}
 \end{aligned}
\]
を得る．従って$(g_i)$は両側正規正則な次数三の族であり，
$b=B/(g_i\mid i\in\mathbb Z)$である．

より具体的に，$B^{(+3)}_{ij}=B_{i+3,j+3}$と書くと，
$\mathbb Z$-代数の同型
\[
 \sigma:B^{(+3)}\xrightarrow{\sim}B,
 \qquad g_i a=\sigma_{ij}(a)g_j
 \quad(a\in B_{i+3,j+3})
\]
が一意に定まる．相対AS正則性・Koszul性は定理\ref{thm:abel-relative-as-lift}で確認する．
\end{proposition}
\begin{proof}
二次関係空間$\mathcal R_i$の特殊化は可換性の関係空間である．
従って，二次関係で割る操作が特殊化と可換であることから，
$B\otimes_R k$は多項式環$S=k[x_0,\ldots,x_4]$の付随$\mathbb Z$-代数に等しい．
各$B_{ij}$は有限自由加群$T_{ij}$の商なので有限であり，
仮定した平坦性と$R$の局所性から有限自由である．
その階数は$\dim_k S_{j-i}=\binom{j-i+4}{4}$である．

完全列
\[
 0\longrightarrow J_{ij}\longrightarrow B_{ij}
 \longrightarrow b_{ij}\longrightarrow0
\]
は$b_{ij}$の自由性により$R$-加群として分裂する．
従って$J_{ij}$は有限自由で，その特殊化は
$\ker(S_{j-i}\to(S/(F))_{j-i})=F S_{j-i-3}$に等しい．
ここで負次数の$S$は零とする．階数の式と$J_{i,i+3}$の自由rank~1性を得る．
$g_i$の特殊化は$F$であるから，$g_i$は$J_{i,i+3}$の生成元である．

写像$g_i\cdot-:B_{i+3,j}\to J_{ij}$の両辺は同じ階数の有限自由$R$-加群である．
その特殊化は
\[
 S_{j-i-3}\longrightarrow F S_{j-i-3},\qquad a\longmapsto Fa
\]
であり，$F\neq0$および$S$が整域であることから同型である．
従って元の写像も同型である．右から$g_{j-3}$を掛ける写像も同様である．
これにより，両側の積が同じideal成分を与える正規性と，
これらの積写像の単射性である正則性が証明された．

最後に，$a\in B_{i+3,j+3}$に対して，$g_i a\in J_{i,j+3}$を
一意に$\sigma_{ij}(a)g_j$と書ける．両側の積写像が同型なので
$\sigma_{ij}$も$R$-加群の同型である．
$a\in B_{i+3,j+3}$，$c\in B_{j+3,l+3}$について
\[
 g_i(ac)=\sigma_{ij}(a)g_jc
        =\sigma_{ij}(a)\sigma_{jl}(c)g_l
\]
となる．$-\cdot g_l$の単射性から
$\sigma_{il}(ac)=\sigma_{ij}(a)\sigma_{jl}(c)$を得る．
また$g_i\cdot1=1\cdot g_i$なので単位も保たれ，$\sigma$は所望の同型である．
\end{proof}

\begin{remark}[周囲のAbel変形問題と三次incidenceの像]
定理\ref{thm:abel-cubic-all-artin}が記述するのは
$\operatorname{coh}(X)$の\emph{周囲のAbel変形問題}である．
AS正則二次包絡と正規正則三次族を持つ生成列の存在によって
$\mathcal O$-marked非可換三次三様体のクラスを定義した場合，
その形式完備化が周囲のAbel変形問題全体と一致することは，まだ証明していない．
三次incidenceからのessential imageは，現段階では真に小さい可能性を排除していない．

命題\ref{prop:abel-flat-envelope-normal-cubic}により，
二次包絡の次数ごとの平坦性を示せば，正規正則三次核という条件は自動的に従う．
全次数平坦性の下では，定理\ref{thm:abel-relative-as-lift}が相対Koszul--AS条件も与える．
従って残る核心は二次包絡の全次数平坦性自体である．
特に，非正則なArtin基底上での相対AS正則性と，係数環も含めた通常の環の
有限global dimensionを混同してはならない．
\end{remark}

\subsection{多項式環の平坦Artin変形における相対Koszul--AS性}
\label{subsec:abel-relative-as-lifting}

以下の補題は二次包絡の全次数平坦性を仮定する．
二次・三次など有限個の成分の平坦性を，この仮定の代用にはしない．
またArtin基底が非正則な場合，基底込みの環の絶対的大域次元は有限とは限らない．
ここでは頂点上の基底加群の線型解消，AS双対，および基底上平坦な加群に対する
射影次元を分けて述べる．

\begin{theorem}[相対Koszul解消と両側AS双対の持上げ]
\label{thm:abel-relative-as-lift}
$(\Lambda,\mathfrak m,k)$を可換局所Artin $k$-代数とする．
$B$は正の連結$\Lambda$-$\mathbb Z$-代数で，
$B_{ii}=\Lambda$，$B_{ij}=0$ $(j<i)$，各$B_{ij}$は有限自由$\Lambda$加群とする．
$B$は次数一生成であり，その特殊ファイバーは
\[
 \overline B:=B\otimes_\Lambda k
       \simeq\check{\operatorname{Sym}_k(k^5)}
\]
とする．二次表示を仮定してもよいが，それ自体も以下の線型解消から従う．
右表現加群$P_i=e_iB$，左表現加群$Q_i=Be_i$，
頂点$i$に$\Lambda$だけを持つ右・左加群をそれぞれ$S_i,T_i$と書く．
このとき次が成り立つ．
\begin{enumerate}
\item $S_i$は自然な線型射影解消
\begin{equation}\label{eq:abel-relative-right-resolution}
 0\longrightarrow W_i^5\otimes_\Lambda P_{i+5}
 \longrightarrow\cdots\longrightarrow W_i^1\otimes_\Lambda P_{i+1}
 \longrightarrow P_i\longrightarrow S_i\longrightarrow0
\end{equation}
を持ち，$W_i^p$は階数$\binom5p$の有限自由$\Lambda$加群である．
その特殊ファイバーは多項式環のKoszul解消である．
左側にも$Q_{i-p}$を用いる同じ型の線型解消がある．
その第$p$生成加群を${}^\ell W_i^p$と書く．
\item AS双対は
\begin{align}
 \operatorname{Ext}^p_{\mathrm{Gr}B}(S_i,P_j)
  &\simeq
   \begin{cases}
    (W_i^5)^* & p=5,\ j=i+5,\\
    0 & \text{それ以外},
   \end{cases}
   \label{eq:abel-relative-right-as}\\
 \operatorname{Ext}^p_{\mathrm{Gr}B^{\mathrm{op}}}(T_i,Q_j)
  &\simeq
   \begin{cases}
    ({}^\ell W_i^5)^* & p=5,\ j=i-5,\\
    0 & \text{それ以外}
   \end{cases}
   \label{eq:abel-relative-left-as}
\end{align}
である．最上次数の直線は$\Lambda$に同型だが，その自明化は選ばない．
\item 各$P_i$は右noetherian，各$Q_i$は左noetherianである．
従って$B$は通常の$\mathbb Z$-代数の意味で左右noetherianである．
\item 上の解消とAS双対は任意のArtin基底変更に可換である．
また各成分が$\Lambda$上平坦な次数付き右・左$B$加群の射影次元は高々$5$であり，
この上界は$S_i,T_i$によって達成される．
\end{enumerate}
従って，この明示した相対的な意味で次元$5$，AS index $5$の
両側Koszul--AS正則性が持ち上がる．$B$全体の1周期同型の選択は仮定していない．
\end{theorem}
\begin{proof}
\emph{線型解消の構成．}
$Z_i^0=S_i$，$F_i^0=P_i$とし，$Z_i^1=\ker(P_i\to S_i)$と置く．
この短完全列は各頂点成分で$\Lambda$加群として分裂するので，
$Z_i^1$の各成分は有限自由であり，特殊ファイバーを取る操作と可換である．
以下，$Z_i^p$まで構成され，各成分が有限自由，
特殊ファイバーが多項式環の第$p$ syzygyであるとする．
この特殊ファイバーは頂点$i+p$の成分で生成され，
その成分の次元は$\binom5p$である．
頂点$i+p$より小さい成分はNakayamaの補題で零であり，
\[
 W_i^p=(Z_i^p)_{i+p},\qquad
 F_i^p=W_i^p\otimes_\Lambda P_{i+p}
\]
と置ける．$W_i^p$は階数$\binom5p$の有限自由加群である．

生成元をsyzygyへ送る評価写像$F_i^p\to Z_i^p$は，
特殊ファイバーでKoszul解消の全射となる．各頂点$j$における余核は
有限生成$\Lambda$加群なので，再びNakayamaの補題によって零となる．
従ってこの評価写像は全射である．
$Z_i^{p+1}=\ker(F_i^p\to Z_i^p)$と置くと，終項の各成分が自由なので
短完全列は成分ごとに分裂し，$Z_i^{p+1}$も成分ごとに有限自由である．
その特殊ファイバーは正しく次のKoszul syzygyになる．
これで帰納法が進む．微分は$F_i^p\to Z_i^p\hookrightarrow F_i^{p-1}$で定めたので，
任意の行列を持ち上げて$d^2=0$を期待する手順は使っていない．

$p=5$の後に現れる$Z_i^6$は，特殊ファイバーのすべての成分が零である．
各成分の有限生成性とNakayamaの補題により$Z_i^6=0$となる．
これで\eqref{eq:abel-relative-right-resolution}を得た．
同じ構成を左側で行えば左線型解消を得る．
$p=2$で関係加群が次数二で生成されることも示されたので，
次数一生成を仮定すれば二次表示は結論として従う．

\emph{AS双対．}
\eqref{eq:abel-relative-right-resolution}に
$\operatorname{Hom}_{\mathrm{Gr}B}(-,P_j)$を適用すると，
第$p$項が
\[
 (W_i^p)^*\otimes_\Lambda B_{j,i+p}
\]
である有界な有限自由$\Lambda$加群の複体を得る．
$j=i+5$では正性により最後の項以外は零なので，
第$5$コホモロジーは$(W_i^5)^*$である．
$j\ne i+5$では特殊ファイバーの複体は多項式環のAS双対により非輪状である．

局所環上の有界有限自由複体が剰余体上で非輪状なら，その複体は可縮である．
実際，剰余体上で非零の微分成分があれば元の行列に単元成分があるので，
基底変換で$\Lambda\xrightarrow{1}\Lambda$という直和因子を分離して消去する．
有限回の消去後に残る微分は全成分が$\mathfrak m$に入り，
特殊ファイバーで微分が零になる．その特殊ファイバーが非輪状であるため残りの自由加群はすべて零である．
これにより\eqref{eq:abel-relative-right-as}を得る．左側も同様である．
解消の短完全列は成分ごとに分裂し，双対複体は可縮部分と最上次数の自由直線へ分かれるので，
これらの主張は任意の基底変更後にも成立する．

\emph{Noetherian性．}
$\mathfrak m^N=0$となる$N$を取り，$P_i$を
$P_i\supset\mathfrak mP_i\supset\cdots\supset\mathfrak m^NP_i=0$で濾過する．
$\Lambda$平坦性から
\[
 \mathfrak m^rP_i/\mathfrak m^{r+1}P_i
  \simeq (\mathfrak m^r/\mathfrak m^{r+1})
       \otimes_k(P_i\otimes_\Lambda k)
\]
となる．各層は多項式環の表現加群の有限直和であり，右noetherianである．
有限回の拡大もnoetherianなので$P_i$は右noetherianとなる．左側も同様である．

\emph{基底上平坦な加群の射影次元．}
まず$W_i^0=\Lambda$，$W_i^1=V_i$とする．最小次数の核を取った構成から
\[
 W_i^p=\ker\bigl(W_i^{p-1}\otimes_\Lambda V_{i+p-1}
       \longrightarrow W_i^{p-2}\otimes_\Lambda B_{i+p-2,i+p}\bigr)
 \qquad(p\geq2)
\]
となる．各核は構成した短完全列の該当成分であり，
$W_i^q\hookrightarrow V_i\otimes\cdots\otimes V_{i+q-1}$は
$\Lambda$-分裂した埋込みである．
線型解消の第$p$生成加群$W_i^p$は，$p\geq2$について，
$V_i=B_{i,i+1}$，二次関係$R_i\subset V_i\otimes V_{i+1}$を用いて
\[
 W_i^p
  =\bigcap_{r=0}^{p-2}
   V_i\otimes\cdots\otimes V_{i+r-1}
   \otimes R_{i+r}\otimes
   V_{i+r+2}\otimes\cdots\otimes V_{i+p-1}
\]
と表される．交差は$V_i\otimes\cdots\otimes V_{i+p-1}$内で取る．
これは上で構成した微分を最小次数に制限した核の式を帰納的に適用すれば分かる．
各段の短完全列が$\Lambda$分裂することにより，この式の特殊ファイバーも
通常の$\bigwedge^p k^5$となる．特に$W_i^p=0$ $(p>5)$である．

さらに$W_i^0=\Lambda$，$W_i^1=V_i$と置き，双加群のKoszul複体を
\[
 \mathcal K_p
  =\bigoplus_{i\in\mathbb Z}
    Be_i\otimes_\Lambda W_i^p\otimes_\Lambda e_{i+p}B
 \qquad(0\leq p\leq5)
\]
で作る．微分は最初と最後の次数一因子をそれぞれ外側へ掛け，
後者に$(-1)^p$を掛けた和とする．上の交差表示と二次関係から$d^2=0$である．
$\mathcal K_0\to B$は積で与える．固定した両端の頂点$j,\ell$について，
和に現れる$i$は$j\leq i\leq\ell-p$に限られるので，
この成分の増大複体は有界な有限自由$\Lambda$複体である．
その特殊ファイバーは多項式環の通常の双加群Koszul解消であり完全である．
先ほどの有限自由複体の消去を用いれば，元の増大複体も完全である．

この解消は左$B$加群として分裂する．実際，
$e_{i+p}B=\bigoplus_j B_{i+p,j}$は$\Lambda$-自由なので，
各$\mathcal K_p$は左表現加群の直和で射影的である．
最後の$B=\bigoplus_j Be_j$自体も左射影的なので，右端から順に短完全列を分裂できる．
従って任意の右$B$加群$M$に対して$M\otimes_B\mathcal K_\bullet$は$M$を解消する．
各$M_i$が$\Lambda$上平坦なら，Artin環上の平坦加群は射影的なので
（\href{https://stacks.math.columbia.edu/tag/051E}{Stacks Project,
Section~10.101, Lemma~10.101.2}），
\[
 M\otimes_B\mathcal K_p
  =\bigoplus_i M_i\otimes_\Lambda W_i^p\otimes_\Lambda P_{i+p}
\]
は右$B$射影加群である．よって射影次元は高々$5$である．
左側も同様であり，\eqref{eq:abel-relative-right-as}と
\eqref{eq:abel-relative-left-as}の非零な最上次数Extが上界の達成を示す．
無限個の頂点については表現双加群の直和を使っており，
ここから対角双加群が有限個の生成子で生成されると主張してはいない．
\end{proof}

\begin{remark}[絶対的大域次元は持ち上がらない]
\label{rem:abel-relative-not-absolute-gldim}
$\Lambda=k[\epsilon]/(\epsilon^2)$，
$B=\check{\Lambda[x_0,\ldots,x_4]}$は定理の仮定を満たす．
しかし頂点に置いた剰余体$k$は$\Lambda$上で無限射影次元を持つ．
$B$射影加群の各頂点成分は$\Lambda$射影的なので，
この$k$の$B$上の射影次元も無限である．
通常の次数付き環$\Lambda[x_0,\ldots,x_4]$の大域次元も無限である．
本定理で有限射影次元を主張したのは，頂点の$\Lambda$および
各成分が$\Lambda$平坦な加群である．
\end{remark}


\subsection{構造対象の標識と形式hull}
\begin{corollary}[正規三次表示の平坦二次包絡判定]
\label{cor:abel-cubic-flatness-criterion}
定理\ref{thm:abel-cubic-all-artin}のsection $\mathbb Z$-代数変形$b$について，
次は同値である．
\begin{enumerate}[label=\textup{(\roman*)}]
\item 二次包絡$Q(b)$の全ての成分が$R$-平坦である．
\item $b$は，多項式環の標準$\mathbb Z$-代数を特殊ファイバーとする
相対Koszul--AS正則次元$5$，index $5$の平坦な二次代数族の，
正規正則な三次族による商として表示できる．
\end{enumerate}
ここで相対AS条件は定理\ref{thm:abel-relative-as-lift}の明示した意味である．
表示が存在するとき，二次ambientと三次idealはそれぞれ$Q(b)$と
$\ker(Q(b)\to b)$に標準的に同型である．
\end{corollary}
\begin{proof}
\textup{(i)}から\textup{(ii)}は，
命題\ref{prop:abel-flat-envelope-normal-cubic}と
定理\ref{thm:abel-relative-as-lift}による．
逆に$b=B'/(g'_i)$という表示があれば，三次idealは次数一・二の成分を変えない．
従って二次関係で生成される$B'$は$b$の二次包絡に標準的に同型であり，
$B'$の平坦性から\textup{(i)}を得る．三次idealはその商写像の核である．
これらの構成はsection代数の同型と底変更に両立する．
\end{proof}

本系は，標準豊富列を特殊ファイバーに持つ三次表示の存在を，
全次数平坦性という一つの条件へ還元する．その条件が全Abel変形で成立することや，
有限個の次数だけから判定できることを証明したものではない．


\begin{proposition}[構造対象の標識を忘れる局所比較]
\label{prop:abel-cubic-o-marking}
上の変形群oidにおいて，対象に$\mathcal O_X$の平坦な持上げ
$\widetilde{\mathcal O}$を加え，射には
$\alpha:F(\widetilde{\mathcal O})\xrightarrow{\sim}\widetilde{\mathcal O}'$を加える．
特殊ファイバーの還元同型は追加frameとして固定せず，
射は$\alpha$と両立する自然同型を法として取る．
こうして得られる$\mathcal O$-marked 1群oidから
$\operatorname{Def}_{\rm ab}(\mathcal C_0;R)$への忘却は同値である．

ただし，自然同型を保持する2群oidでは同じ結論は成り立たない．
固定した$\mathcal D$内で還元同型を指定した構造対象の持上げ群oidは連結であり，
その自己同型群は$\ker(R^\times\to k^\times)=1+\mathfrak m_R$である．
還元同型を指定しなければ自己同型群は$R^\times$である．
\end{proposition}
\begin{proof}
持上げの存在と同型類の一意性は前定理による．
また$\operatorname{End}_{\mathcal D}(\widetilde E_i)$は有限自由rank~1であり，
単位写像$R\to\operatorname{End}(\widetilde E_i)$は還元後同型だから同型である．
従って，任意の変形同値$F$は$\widetilde{\mathcal O}$を
$\widetilde{\mathcal O}'$へ送る同型$\alpha$を持つ．

同じ$F$に対する二つの$\alpha$は$R^\times$の元だけ異なり，
その差は$F$全体のスカラー自然自己同型で補正できる．
より一般に，$F,G$の間の自然同型$\eta$と，それぞれに指定した
$\alpha,\beta$があれば，$\beta\eta_{\widetilde{\mathcal O}}\alpha^{-1}$はスカラー単元である．
$\eta$をその逆のスカラーで補正すれば，標識と両立する自然同型になる．
従って忘却はHom集合上で全単射であり，対象上本質的全射だから同値である．

持上げ自体の自己同型群については
$\operatorname{End}(\widetilde{\mathcal O})=R$から直ちに従う．
さらに$\operatorname{Aut}_2(\operatorname{Id}_{\mathcal D})=R^\times$である．
実際，任意の自然自己同型は各$\widetilde E_i$上でスカラー$\lambda_i$であり，
隣接Homが忠実な有限自由rank~5加群であることから自然性により
$\lambda_i=\lambda_{i+1}$となる．$\widetilde E_i$は$\mathcal D$を生成するため，
全体への作用もこの共通スカラーである．
構造対象の同定を保つ2自己同型はこのスカラーを1にするので自明である．
従って非切断2群oidの忘却は2慣性を変える．
\end{proof}

\begin{remark}[本研究の三次条件と形式hullとの関係]
定理\ref{thm:abel-cubic-all-artin}と命題\ref{prop:abel-cubic-o-marking}は
homogeneous formalityを仮定しない．また，次数一生成と二次・三次関係による提示は証明されたが，
二次包絡の全次数平坦性は，この比較定理には含まれない．
全次数平坦性が得られれば，定理\ref{thm:abel-relative-as-lift}によって相対AS性も従う．
正規正則三次族の存在は，二次包絡の全次数平坦性の下では
命題\ref{prop:abel-flat-envelope-normal-cubic}から従う．
そのため，すべてのAbel変形が現行の正規三次incidenceに入ることは別途証明する必要がある．

Hochschild制御\cite[Theorem~3.1]{LowenVdBHochschild}から，
Abel変形の一次空間は$HH^2(X)$である．既存の計算では
$HH^2(X)=H^1(X,T_X)\oplus H^0(X,\wedge^2T_X)$は20次元である．
本節の比較により，構造対象を標識してもこの一次空間は変わらない．
一方，10変数とGrassmannian錐の積という完全な形式hullを主張するには，
既存原稿のformality入力と，その模型がここで定めた変形問題を表すことの
比較がなお必要である．この二つの仮定を本節の定理によって消したことにはしない．

さらに，$\mathcal O_X(n)$の全列が変形先へ持ち上がることは，
その列を一つ移す自己同値の存在を意味しない．
通常の次数付き代数への1周期化，Serre関手による反標準偏極の認識，
ならびに特殊ファイバーを固定しない大域的な無標識復元は別の問題である．
\end{remark}

\section{構造対象付きAbel圏stackの可換点における形式完備化}
\label{sec:abel-formal-descent}

本節では，第\ref{sec:abel-formal-cubic}節の標準相対模型と
全fpqc stackの形式近傍を比較する．
任意の平坦被覆上の線対象を元のArtin基底に関して持ち上げ，
三重交叉のscalar cocycleを補正してから，Homと合成だけを降下する．
最後にframed自己同型を直接計算し，形式変形群oidそのものを表示する．
ここで形式近傍は特殊ファイバーへの同一視を指定したものとし，
有限自己同型群を含む無標識の完備化と区別する．
到達点は系\ref{cor:abel-full-stack-discrete-formal}の群oid値関手の表示である．
その証明は，定理\ref{thm:abel-stackified-relative-formal}による
標準模型の本質的全射性と，命題\ref{prop:abel-framed-artin-automorphisms}による
自己同型の消滅を別々に示してから組み合わせる．

\subsection{任意の平坦被覆上での線対象の持上げと復元}
\label{subsec:abel-flat-cover-reconstruction}

$k=\mathbb C$、$R$ を剰余体 $k$ の可換局所 Artin $k$-代数とする。
$\mathfrak m=\mathfrak m_R$ とし、$S$ を任意の平坦な可換 $R$-代数とする。
$S$ は noetherian または coherent とは仮定しない。
\[
 I=\mathfrak mS,\qquad Q=S/I
\]
と書く。$X/k$ は固定した滑らかな三次三様体である。
この小節で既存の Abel 圏変形定理を適用する係数全射は、終始
\[
 R\longrightarrow k
\]
である。非 coherent であり得る $S\to Q$ に、coherent な係数環を仮定する
定理を直接適用することはしない。

\begin{lemma}[平坦な線型圏の Gabriel 商]
\label{lem:abel-flat-quotient-over-artin}
$R$ を可換 coherent 環、$c$ を Hom が $R$-平坦な小さい $R$-線型圏とする。
$\pi:\operatorname{Mod}(c)\to\mathcal D$ を $R$-線型 Gabriel 商とし、
その完全忠実な右随伴を $\omega$ とする。
このとき $\mathcal D$ は Lowen--Van den Bergh の意味で $R$-flat である。
従って上記の平坦な $R\to S$ と、全 Hom が $S$-flat な tails 模型 $c_S$ について、
$\operatorname{Qmod}(c_S)$ は $R$-flat である。
\end{lemma}

\begin{proof}
\cite[Proposition~3.7]{LowenVdBDeformation} により
$\operatorname{Mod}(c)$ は $R$-flat である。
$J$ を $\mathcal D$ の injective 対象とすると、$\pi$ が exact なので
$\omega J$ は $\operatorname{Mod}(c)$ の injective 対象である。
同論文 Proposition~3.4 により $\omega J$ は $R$-coflat である。
任意の有限表示 $R$-加群 $M$ に対し、内部 $\operatorname{Hom}_R(M,-)$ は
有限直和間の写像の kernel で計算される。
$\pi$ は exact なので、この kernel と可換し、
\[
 \operatorname{Hom}_R(M,J)
 \simeq\pi\operatorname{Hom}_R(M,\omega J)
\]
となる。従って $J$ は $R$-coflat である。
再び Proposition~3.4 により $\mathcal D$ は $R$-flat である。
最後の主張は $S$-flat な Hom が $R$ 上も flat であることから従う。
ここでは $S$ の coherent 性も、$R$ 上の有限表示性も使わない。
\end{proof}

\begin{lemma}[元の Artin 係数から平坦被覆の係数への平坦性]
\label{lem:abel-flat-coefficient-upgrade}
$\mathcal D$ を $S$-線型 Grothendieck 圏、$E\in\mathcal D$ を
$R$-flat な対象とする。
$E_0=Q\otimes_S E=k\otimes_R E$ が
$\mathcal D_k$ において $Q$-flat であれば、$E$ は $S$-flat である。
ここで任意の係数環 $T$ に関する対象の flatness は、
全ての $T$-加群に対する $-\otimes_T E$ の完全性で定義する。
coherent な係数環では、これは有限表示加群で試す通常の定義と一致する。
\end{lemma}

\begin{proof}
Grothendieck 圏では任意直和が exact なので、係数加群の自由分解により
内部 Tor が定義できる。有限表示加群上の $R$-flatness は、
filtered colimit の完全性から全ての $R$-加群上の完全性を意味する。
$k$ の自由 $R$-分解 $P_\bullet\to k$ を取る。
$S$ が $R$-flat なので $S\otimes_RP_\bullet$ は
$Q=S\otimes_Rk$ の自由 $S$-分解であり、
\[
 \operatorname{Tor}^S_q(Q,E)
 =\operatorname{Tor}^R_q(k,E)=0\qquad(q>0),
 \quad Q\otimes_SE=E_0.
\]
任意の $Q$-加群 $M$ について、自由 $Q$-分解を取り、
その項が $\operatorname{Tor}^S_{>0}(-,E)$ に関して acyclic であることを使う。
通常の二重分解による係数変更から
\[
 M\otimes_S^{\mathbf L}E
 \simeq M\otimes_Q^{\mathbf L}E_0
\]
を得る。右辺は $E_0$ の $Q$-flatness により次数零に集中する。
この計算は係数加群の自由分解とその tensor だけからなり、
$\mathcal D$ に十分な射影対象があるとは仮定していない。

$I^N=0$ となる $N$ を取る。
任意の $S$-加群 $M$ は有限 filtration
\[
 M\supset IM\supset\cdots\supset I^NM=0
\]
を持ち、その商は $Q$-加群である。
前段の消滅と Tor の長完全列を帰納的に適用すると
$\operatorname{Tor}^S_q(M,E)=0$ for $q>0$ となる。
従って $E$ は $S$-flat である。
\end{proof}

\begin{lemma}[平坦被覆上の線対象と section 復元]
\label{lem:abel-flat-cover-reconstruction}
$S$-線型 Grothendieck 圏 $\mathcal D$ は $R$-flat であり、
$Q$-線型同値
\[
 \iota:\operatorname{Qcoh}(X_Q)\xrightarrow{\sim}\mathcal D_k
 \quad\bigl(X_Q=X\times_k\operatorname{Spec}Q\bigr)
\]
が指定されているとする。
$E_{i,0}=\iota(\mathcal O_{X_Q}(-i))$ と置く。
このとき以下が成立する。
\begin{enumerate}
\item 各 $E_{i,0}$ は $\mathcal D$ 内に $S$-flat な持上げ $E_i$ を持つ。
還元同型を付けた $S$-flat な持上げの群oidは連結である。
各 $E_i$ は有限表示対象であり、
\[
 \operatorname{End}_{\mathcal D}(E_i)=S,
 \qquad \operatorname{Aut}(E_i,E_{i,0})=1+I
\]
である。最初の同型は scalar action による標準写像である。
\item section $\mathbb Z$-代数
\[
 b_{ij}=\operatorname{Hom}_{\mathcal D}(E_j,E_i)
\]
は $i>j$ で零であり、$i\leq j$ では有限自由 $S$-加群である。
その階数は $h_{j-i}=\dim_k H^0(X,\mathcal O_X(j-i))$ であり、
$b$ は正、連結、次数一生成である。
\item 任意の $S$-加群 $M$ と任意の $i,j$ に対し、標準写像
\begin{equation}\label{eq:abel-flat-cover-universal-hom}
 M\otimes_S b_{ij}\xrightarrow{\sim}
 \operatorname{Hom}_{\mathcal D}(E_j,M\otimes_S E_i)
\end{equation}
は同型であり、
$\operatorname{Ext}^1_{\mathcal D}(E_j,M\otimes_S E_i)=0$ である。
\item section functor は $S$-線型同値
\[
 \mathcal D\simeq\operatorname{Qmod}(b)
\]
を与える。任意の可換 $S$-代数 $T$ に対し、
\[
 \mathcal D_T\simeq\operatorname{Qmod}(T\otimes_S b),
 \qquad
 \operatorname{Hom}_{\mathcal D_T}(T\otimes_SE_j,T\otimes_SE_i)
 =T\otimes_S b_{ij}
\]
が成立する。これらは係数変更の合成、乗法、section functor と両立する。
この部分では $T$ の平坦性も noetherian 性も不要である。
\end{enumerate}
\end{lemma}

\begin{proof}
まず $\operatorname{Qcoh}(X_Q)$ を $k$-線型 Grothendieck 圏として見る。
その列 $(\mathcal O_{X_Q}(n))_{n\in\mathbb Z}$ は
De Deken--Lowen の意味で $\mathbb Z$-generating である。
この事実には noetherian 性は必要ない。
例えば $X$ の通常の section 代数 $a$ と
補題\ref{lem:abel-arbitrary-coefficient-tails} を用いると
\[
 \operatorname{Qmod}(Q\otimes_ka)
 \simeq\operatorname{Qcoh}(X)_Q
 \simeq\operatorname{Qcoh}(X_Q)
\]
となり、標準対象が上記の列に対応する。
第二の同値は affine 射 $X_Q\to X$ の準連接加群の記述による。

有限 affine 標準被覆の \v Cech 複体を使うと、任意の $Q$-加群 $M$ について
\begin{equation}\label{eq:abel-flat-cover-special-ext}
 \operatorname{Ext}^q_{\operatorname{Qcoh}(X_Q)}
 (\mathcal O_{X_Q}(-j),M\otimes_Q\mathcal O_{X_Q}(-i))
 \simeq M\otimes_kH^q(X,\mathcal O_X(j-i)).
\end{equation}
これは $k$ 上の複体を $M$ と tensor する計算であり、
$M$ の $Q$-flatness を仮定しない。
三次三様体の中間コホモロジー消滅により、$q=1,2$ の右辺は零である。
特に有限次元 $k$-ベクトル空間を係数とする必要な消滅も成立する。
体 $k$ 上では全ての対象は係数に関して flat である。

従って coherent な全射 $R\to k$ に
\cite[Theorem~4.5, pp.~495--496]{DeDekenLowen} を適用できる。
同定理は Grothendieck 圏と $\mathbb Z$-generating 列を扱い、
Hom の有限次元性、圏の locally noetherian 性、または locally coherent 性を仮定しない。
これにより $R$-flat な持上げ $E_i$ が得られ、
その非負方向の Hom が $R$-flat な線型変形 $b$ をなし、
\[
 k\otimes_R b_{ij}\simeq Q\otimes_k a_{ij},\qquad
 \mathcal D\simeq\operatorname{Qmod}(b)
\]
となる。持上げの群oidの連結性は、
\cite[Theorems~6.11--6.12 の双対版]{LowenObstruction}
と自己 Ext の一次・二次消滅からも従う。
ここで $\operatorname{End}(E_{i,0})=Q$ であり、
$E_{i,0}$ を $k$ 上の例外対象と呼ぶ必要はない。

$S$ の scalar action は section 代数 $b$ の中心に作用する。
$b_{ij}$ は $R$-flat であり、その $Q$ 還元は $Q^{h_{j-i}}$ である。
還元の基底を持ち上げて $S^{h_{j-i}}\to b_{ij}$ を作る。
$S^{h_{j-i}}$ と $b_{ij}$ はともに $R$-flat であり、
任意の $R$-flat 加群 $V$ には
\[
 \mathfrak m^rV/\mathfrak m^{r+1}V
 \simeq(\mathfrak m^r/\mathfrak m^{r+1})
       \otimes_k(V/\mathfrak mV)
\]
が成り立つ。
従って上の写像は有限 $\mathfrak m$-filtration の全ての商で同型であり、
元の写像も同型である。
この議論では加群が $R$ 上有限生成であるとは仮定していない。
従って $b_{ij}$ は主張した有限自由加群となる。
$i=j$ では還元基底を $1$ と選べるので、同型は標準の scalar action
$S\to\operatorname{End}(E_i)$ そのものである。

負方向については、$E_i$ の $\mathfrak m$-filtration の商が
\[
 \mathfrak m^rE_i/\mathfrak m^{r+1}E_i
 \simeq(\mathfrak m^r/\mathfrak m^{r+1})\otimes_k E_{i,0}
\]
であることを使う。還元と包含の随伴により、$E_j$ からこの商への Hom は
$E_{j,0}$ からの特殊ファイバーの Hom に等しい。
$i>j$ のときそれは零なので、有限 filtration について帰納すると
$\operatorname{Hom}(E_j,E_i)=0$ を得る。
次数一成分の積による生成性も、還元上の生成性と cokernel への冪零 Nakayama によって従う。
従って $b$ は $S$ 上正、連結、次数一生成である。
この有限生成性により通常の tails torsion は localizing であり、
DDL の $\operatorname{Qmod}$ は通常の Gabriel 商と一致する。

ここで $E_i$ の $S$-flatness を直接確認できる。
section functor
$H:\mathcal D\to\operatorname{Mod}(b)$、
$H(N)_i=\operatorname{Hom}_{\mathcal D}(E_i,N)$ は、
$S$-線型性により precomposition と postcomposition の scalar action が一致するので
$S$-線型である。
その左随伴 $\pi$ も $S$-線型であり、DDL によって exact な Gabriel quotient となる。
$E_i=\pi(b(-,i))$ であり、representable $b(-,i)$ は各成分が $S$-自由なので
$S$-flat な対象である。
$\pi$ は直和と cokernel を保つため任意の係数 tensor と可換し、exact なので
$E_i$ も $S$-flat である。
逆に $S$-flat な対象は $R\to S$ の平坦性により $R$-flat なので、
持上げの群oidの連結性は $S$-flat な持上げ全体について成り立つ。
従ってこの主要な復元の証明は、
補題\ref{lem:abel-flat-coefficient-upgrade} の一般的な内部 Tor 計算を使わずにも完結する。

次に任意係数の Hom を証明する。
$N_0\in\mathcal D_k$ に対し、$E_j$ の $R$-flatness と
\cite[Proposition~4.7 の双対版]{LowenVdBDeformation} から
\begin{equation}\label{eq:abel-flat-cover-ext-reduction}
 \operatorname{Ext}^q_{\mathcal D}(E_j,N_0)
 \simeq\operatorname{Ext}^q_{\mathcal D_k}(E_{j,0},N_0)
 \quad(q\geq0)
\end{equation}
となる。すなわち、係数全射 $R\to k$ の Tor を用いる change-of-rings
spectral sequence が、$E_j$ の flatness により退化する。
$q=0$ は左随伴 $k\otimes_R-$ の通常の随伴そのものであり、
$q=1$ だけなら補題\ref{lem:abel-flat-object-ext-one} の
拡大の還元と引戻しによる Yoneda の証明も使える。

$M$ が $Q$-加群の場合、
\eqref{eq:abel-flat-cover-special-ext} と
\eqref{eq:abel-flat-cover-ext-reduction} により
\eqref{eq:abel-flat-cover-universal-hom} と必要な $\operatorname{Ext}^1$ 消滅を得る。
一般の $S$-加群 $M$ については有限 $I$-filtration を使う。
$E_i$ と $b_{ij}$ は $S$-flat なので、係数の短完全列は
$-\otimes_SE_i$ と $-\otimes_Sb_{ij}$ の下で短完全列のままである。
Hom--Ext の長完全列により $\operatorname{Ext}^1$ の消滅は拡大について閉じており、
Hom の標準写像も短完全列の比較から同型となる。
有限 filtration の長さについて帰納して任意の $M$ に結論を得る。

有限表示性も noetherian 性なしに確認できる。
$E_{i,0}$ は直線束なので、任意の準連接層を第二変数とする全ての Ext は、
有限 affine 被覆の \v Cech 複体で計算でき、filtered colimit と可換する。
式 \eqref{eq:abel-flat-cover-ext-reduction} により、
同じことが $\mathcal D_k$ に値を持つ filtered diagram について成立する。
一般の filtered diagram $N_\lambda$ は
$\mathfrak m^rN_\lambda/\mathfrak m^{r+1}N_\lambda$ により有限に分解する。
Grothendieck 圏の filtered colimit は exact なので、
これらの filtration とその商を取る操作と可換する。
長完全 Ext 列の比較を filtration の長さについて帰納すると、
全ての $\operatorname{Ext}^q_{\mathcal D}(E_i,-)$ が filtered colimit と可換する。
特に $q=0$ により $E_i$ は有限表示である。

持上げの還元 frame を保つ自己同型は、$\operatorname{End}(E_i)=S$ のうち
$Q$ で $1$ となる単元である。従ってその群は $1+I$ である。

最後に、既に確認した $S$-線型 section 同値と、任意の可換 $S$-代数 $T$ に対して
補題\ref{lem:abel-arbitrary-coefficient-tails} を適用すれば
$\mathcal D_T\simeq\operatorname{Qmod}(T\otimes_Sb)$ を得る。
係数作用の自由拡張と忘却の随伴により
\begin{align*}
 \operatorname{Hom}_{\mathcal D_T}(T\otimes_SE_j,T\otimes_SE_i)
 &\simeq\operatorname{Hom}_{\mathcal D}(E_j,T\otimes_SE_i)\\
 &\simeq T\otimes_Sb_{ij}.
\end{align*}
これは標準の乗法写像であるため、合成と係数変更の合成に両立する。
有限表示性も自由拡張・忘却の随伴と忘却の filtered colimit 保存から保たれる。

特に $S\to T$ が平坦な場合、$T\otimes_SE_i$ は $T$-flat、従って $R$-flat であり、
$T_0=T/\mathfrak mT$ 上の $\mathcal O_{X_{T_0}}(-i)$ を持ち上げる。
$T\otimes_Sb$ の Hom は $T$-自由なので、
補題\ref{lem:abel-flat-quotient-over-artin} により $\mathcal D_T$ は $R$-flat である。
この場合の Hom の係数変更は、任意係数の議論とは独立にも確認できる。
元の係数全射 $R\to k$ に対する DDL の定理を $\mathcal D_T$ に適用すると、
\[
 T\otimes_S b_{ij}\longrightarrow
 \operatorname{Hom}_{\mathcal D_T}(T\otimes_SE_j,T\otimes_SE_i)
\]
の両辺は $R$-flat であり、還元は
$T_0\otimes_k a_{ij}$ の恒等写像となる。
有限 $\mathfrak m$-filtration の各商を比較すれば同型である。
従って fpqc 被覆の多重交叉で必要な Hom の係数変更は、
$E_i$ の有限表示性を前提とせずに成立する。
\end{proof}

\begin{remark}[section 復元の相対性と圏同値との整合性]
\label{rem:abel-flat-cover-reconstruction-natural}
上の同値は選んだ線対象とその全 Hom に関して自然である。
実際、係数作用の忘却を $U:\mathcal D_T\to\mathcal D$ とすると、
任意の $N\in\mathcal D_T$ に対して
\[
 H_T(N)_i
 =\operatorname{Hom}_{\mathcal D_T}(T\otimes_SE_i,N)
 \simeq\operatorname{Hom}_{\mathcal D}(E_i,UN)=H(UN)_i
\]
となる。右辺には $N$ の係数作用から $T$ が作用し、この同型はその作用と両立する。
従って $H_T$ は $H$ を係数作用付き対象へ延長した functor であり、
Gabriel 商の係数変更と合わせて、復元同値の相対性を与える。

さらに $F:\mathcal D\to\mathcal D'$ を $S$-線型完全同値とし、
$\theta_i:F(E_i)\xrightarrow{\sim}E_i'$ を指定する。
これらは section 代数の同型 $\phi:b\xrightarrow{\sim}b'$ を定める。
全ての $N$ について、完全忠実性と $\theta_i$ による同型
\[
 \operatorname{Hom}_{\mathcal D}(E_i,N)
 \xrightarrow{\sim}\operatorname{Hom}_{\mathcal D'}(E_i',FN)
\]
は、$\phi$ による module 構造の移送と両立する。
従って $H'F$ は $H$ の $\phi$ による移送と自然同型である。
section functor $H'$ の完全忠実性から、$F$ は
$\phi$ から復元される $\operatorname{Qmod}$ の同値と自然同型になる。
これらの同型は合成と係数変更に両立する。
この議論で線対象が射影対象であるとは仮定しておらず、
線対象の有限表示による射の持上げも必要ない。
\end{remark}

\begin{remark}[実際の降下に供給する内容と残る段階]
この補題は、fpqc 被覆上の局所圏を、固定した特殊線対象の平坦持上げと
section $\mathbb Z$-代数で表示する。
被覆やその多重交叉の係数環が非 noetherian であっても適用できる。
局所対象を別の局所対象へ移す $S$-線型同値は平坦持上げを平坦持上げへ送り、
還元上の frame を保つ比較同型の差は正確に $1+\mathfrak mS$ である。
従って、残る降下はこの群の cocycle と section 代数の加群降下に帰着する。

ただし持上げの群oidが連結であることだけから、比較同型が三重交叉で
cocycle 条件を満たすとは結論しない。
その差の $2$-cocycle を別途消す必要がある。
またこの補題だけで、任意の Grothendieck 圏の $2$-降下の有効性を主張するものではない。
局所圏の特殊ファイバーへの同値 $\iota$ と被覆上の実際の transition functor は、
固定した係数族の Hom 降下を用いて先に実現しておく。
\end{remark}

\begin{remark}[一次文献の正確な適用範囲]
\begin{enumerate}
\item De Deken--Lowen, \S4.1, p.~493 は $R\to k$ を coherent 可換環の冪零全射、
$k$ を有限表示 $R$-加群と仮定する。
Theorem~4.5, pp.~495--496 はこの設定の Grothendieck 圏に適用され、
圏の noetherian 性や Hom の有限性は追加仮定でない。
従って本補題では $R\to k$ を固定したまま特殊圏を
$\operatorname{Qcoh}(X_Q)$ に取り換えられる。
\item 同論文 Proposition~4.6 の「connected」「locally finite」を
$k$ 上の $Q\otimes_ka$ にそのまま適用していない。
一般にその対角 End は $Q$ であり、$k$ ではない。
$S$ 上の有限自由性と connected 性は本補題内で別に証明した。
\item Lowen--Van den Bergh, \S4, pp.~5451--5454 の一般の係数変更命題には
coherent 性と係数加群の有限表示性がある。
任意の $R\to S$ に Proposition~4.8 を直接適用して平坦性を結論していない。
Gabriel 商の平坦性は補題\ref{lem:abel-flat-quotient-over-artin} の直接証明で供給する。
\item Lowen, \S6.4, Theorems~6.11--6.12 は、coherent な係数環の冪零変形に対する
対象と射の障害理論である。本文で coflat 版を述べ、\S5.2 で flat 版への双対化を明記する。
本補題で使うのは Artin 係数 $R\to k$ に対するその flat 版である。
\end{enumerate}
\end{remark}

\subsection{fpqc stack 化後の Artin 形式近傍}

この小節では、特殊ファイバーとの同一視を指定した形式変形を扱う。
$x=(\operatorname{Qcoh}(X),\mathcal O_X)$ とし、
$\widehat{\mathfrak M}^{\mathcal O}_{\mathrm{cub},X}(R)$ の対象は
$\mathfrak M^{\mathcal O}_{\mathrm{cub}}(R)$ の対象 $\xi$ と同型
$\iota:\xi_{\mathbb C}\xrightarrow{\sim}x$ の組とする。
射は $\iota$ を保つものとする。
この定義では、$\xi$ 自身が $R$ 上の大域的な標準係数模型を持つことを仮定しない。

\begin{lemma}[還元上恒等な scalar の Amitsur 複体]
\label{lem:abel-unipotent-amitsur}
$R$ を局所 Artin $\mathbb C$-代数、$\mathfrak m$ をその極大イデアルとし、
$R\to S$ を忠実平坦な可換環準同型とする。
$S^{[q]}=S^{\otimes_R(q+1)}$ と置く。
還元上恒等な単元の余単体群
\[
 U(S^{[q]})=
 \ker\bigl((S^{[q]})^\times\longrightarrow
          (S^{[q]}/\mathfrak mS^{[q]})^\times\bigr)
 =1+\mathfrak mS^{[q]}
\]
の \v{C}ech cohomology は正次数で零である。
特に、$U(S^{[2]})$ に値を持つ descent の二次 cocycle は
$U(S^{[1]})$ に値を持つ一次 cochain の coboundary であり、
一次 cocycle も零次 cochain の coboundary である。
$S$ または $S/\mathfrak mS$ の noetherian 性、有限型性、Artin 性は必要ない。
\end{lemma}

\begin{proof}
$\mathfrak m^{N+1}=0$ を満たす $N$ を取る。
任意の可換 $R$-代数 $T$ 上の有限多項式
\[
 \log(1+z)=\sum_{r=1}^N\frac{(-1)^{r+1}}r z^r,
 \qquad
 \exp(z)=\sum_{r=0}^N\frac1{r!}z^r
\]
は、$1+\mathfrak mT$ と加法群 $\mathfrak mT$ の互いに逆な同型である。
可換性と標数零を用いており、これらの同型は全ての係数写像に自然である。
$S^{[q]}$ は $R$ 上平坦なので、
\[
 \mathfrak m\otimes_R S^{[q]}\xrightarrow{\sim}
 \mathfrak mS^{[q]}.
\]
従って、上の乗法的余単体群の複体は、対数によって
\[
 \mathfrak m\longrightarrow\mathfrak m\otimes_R S
 \longrightarrow\mathfrak m\otimes_R S^{[1]}
 \longrightarrow\mathfrak m\otimes_R S^{[2]}
 \longrightarrow\cdots
\]
という拡大 Amitsur 複体に移る。
この複体の完全性は忠実平坦降下の基本計算
\cite[Lemma~35.3.6, Tag~023M]{StacksAbel}である。
具体的には $S$ に係数変更すると、環準同型
$S\to S\otimes_R S$ が積による retraction を持つため縮約を持ち、
元の完全性は忠実平坦性から従う。
指数を取り戻せば主張を得る。
\end{proof}

\begin{theorem}[stack 化後の全 Artin 対象の大域模型]
\label{thm:abel-stackified-relative-formal}
$X$ を滑らかな複素三次三様体とする。
定義\ref{def:abel-cubic-pointed}の全ての構造対象付き相対完全同値を射とする
fpqc stack $\mathfrak M^{\mathcal O}_{\mathrm{cub}}$ について、
任意の局所 Artin $\mathbb C$-代数 $R$ 上の自然な関手
\begin{equation}\label{eq:abel-full-stack-formal-equivalence}
 \widehat{\mathfrak M}^{\mathcal O,\mathrm{can}}_{\mathrm{cub},X}(R)
 \longrightarrow
 \widehat{\mathfrak M}^{\mathcal O}_{\mathrm{cub},X}(R)
\end{equation}
は群oidの同値である。
従って、定理\ref{thm:abel-canonical-relative-formal}の比較は
\[
 \widehat{\mathfrak M}^{\mathcal O}_{\mathrm{cub},X}(R)
 \simeq
 \operatorname{Def}^{\mathcal O}_{\mathrm{ab}}
       (\operatorname{Qcoh}(X);R)
\]
へ延びる。
これは局所 Artin 環の任意の射に自然であり、
得られる大域模型の任意の係数環への標準延長は、もとの相対対象と同値である。
\end{theorem}

\begin{proof}
\emph{局所代表と実際の遷移同値。}
$\mathfrak m=\mathfrak m_R$ とし、
$(\xi,\iota)\in\widehat{\mathfrak M}^{\mathcal O}_{\mathrm{cub},X}(R)$ を取る。
stack 化の局所表示により、あるアフィン忠実平坦被覆 $R\to S$ 上で
$\xi_S$ は許容な標準 tails 係数族 $\mathcal D$ により表せる。
ここで $S$ を Artin 環に取れるとは仮定しない。
$S_0=S/\mathfrak mS$ と置く。

重なり上の二つの標準係数模型の間に stack の射があるとき、
補題\ref{lem:abel-canonical-family-equivalences}の Hom 降下により、
その射は実際の構造対象付き相対完全同値によって表される。
同じ理由で、$\iota$ を $S_0$ 上へ制限した射は、実際の同値
\[
 \mathcal D_{S_0}\simeq
 \operatorname{Qcoh}(X\times_{\mathbb C}\operatorname{Spec}S_0)
\]
によって表せる。
従って、特殊ファイバー上の被覆を新たに選んで $R$ 上へ持ち上げる必要はない。

二つの marked 同値の間に同型があれば、その marked 自然同型は一意である。
実際、許容な列上の End は係数環であり、隣接 Hom の忠実性と生成性から
自然自己同型は一つの scalar で、構造対象の標識がこれを $1$ にする。
よって、重なり上の遷移同値を $F_{12}$ と書くと、三重の重なり上には
一意な marked 自然同型
\[
 a_{123}:F_{23}F_{12}\xRightarrow{\sim}F_{13}
\]
があり、四重の重なり上の pentagon はこの一意性から成立する。
これは実際の遷移データを得る議論であり、圏の一般的な有効 $2$-降下を
仮定する議論ではない。

\emph{各線対象の framed descent。}
局所 tails 模型の Hom は $S$-flat であり、$S$ は $R$-flat なので、
補題\ref{lem:abel-flat-quotient-over-artin}により $\mathcal D$ は $R$-flat である。
補題\ref{lem:abel-flat-cover-reconstruction}を、元の Artin 環
$R\twoheadrightarrow\mathbb C$ に対して適用する。
その特殊圏は $\mathbb C$-線型圏
$\operatorname{Qcoh}(X_{S_0})$ であり、$S_0$ を体と見なしていない。
各 $i\in\mathbb Z$ について、$\mathcal O_{X_{S_0}}(-i)$ の平坦な持上げ
$E_{i,S}\in\mathcal D$ と、その還元の同定を取れる。
$i=0$ では、もとから与えられている構造対象 $E_{0,S}=\mathcal O_S$ と
その還元同定を用いる。
同補題は任意の平坦な $R$-係数環に適用できるから、
$S^{[q]}$ 上でもこれらの持上げの同型類は一意であり、
\[
 \operatorname{End}(E_{i,S^{[q]}})=S^{[q]},\qquad
 \operatorname{Aut}(E_{i,S^{[q]}},\text{還元同定})
       =1+\mathfrak mS^{[q]}.
\]
したがって二重の重なり上で、還元の同定を保つ同型
\[
 \phi_{i,12}:F_{12}(E_{i,1})\xrightarrow{\sim}E_{i,2}
\]
を選べる。
$\phi_{0,12}$ は構造対象に既に指定された同型とし、変更しない。

三重の重なり上で descent のずれを
\[
 \phi_{i,23}F_{23}(\phi_{i,12})
 =c_{i,123}\,\phi_{i,13}a_{123,E_{i,1}}
\]
によって定める。
左右の写像は還元後に一致するので
$c_{i,123}\in1+\mathfrak mS^{[2]}$ である。
遷移同値は係数線型で scalar を保つため、pentagon は
$c_i$ が可換群値の二次 \v{C}ech cocycle であることを意味する。
補題\ref{lem:abel-unipotent-amitsur}により
$c_i=\delta u_i$ となる $u_i\in1+\mathfrak mS^{[1]}$ が取れる。
$\phi_{i,12}$ を $u_{i,12}^{-1}\phi_{i,12}$ に変更すると、
厳密な descent の等式を得る。
$i=0$ では最初から $c_0=1$ なので $u_0=1$ とする。
全ての添字に同じ被覆 $R\to S$ を使い、この補正を各添字について行う。
無限積の群が quasi-coherent であることや、
無限個の線束を一つの有限表示被覆で自明化できることは用いない。

この計算は、特殊ファイバーで trivialization を指定した線対象の lift gerbe の
band が $\mathbb G_m$ 全体ではなく $1+\mathfrak m\mathcal O$ であることを利用している。
可換 band の gerbe に対する通常の cocycle 補正と同じ計算である
\cite[Lemma~21.11.1]{StacksAbel}。
ここでは必要な二次 cocycle を固定した被覆上で直接消しており、
一般の $\mathbb G_m$-gerbe の分類は使用していない。

\emph{section 代数の通常の加群降下。}
\[
 b^S_{ij}=\operatorname{Hom}_{\mathcal D}(E_{j,S},E_{i,S})
\]
と置く。局所再構成補題により、これは標準 Hilbert 関数の階数を持つ
有限自由 $S$-加群であり、特に重なりを定める平坦な係数変更と可換する。
補正後の $\phi_i$ と $F_{12}$ を用いて
\[
 f\longmapsto\phi_{i,12}\,F_{12}(f)\,\phi_{j,12}^{-1}
\]
とすれば、$b^S_{ij}$ に通常の $S/R$-加群降下データが入る。
この降下データは単位と合成を保つ。
忠実平坦な加群降下と有限射影性の降下
\cite[Proposition~35.3.9, Tag~023N;
 Proposition~10.83.3, Tag~058S]{StacksAbel}
により、有限射影 $R$-加群 $b_{ij}$ とその積を得る。
$R$ は局所環なので各 $b_{ij}$ は有限自由である。
結合律と単位律は $S$ で確認でき、これらは一つの $R$-線型
$\mathbb Z$-代数 $b$ を定める。

還元の同定を保つように $\phi_i$ を選んだので、還元後の降下データは
\[
 \operatorname{Hom}_X(\mathcal O_X(-j),\mathcal O_X(-i))
       \otimes_{\mathbb C}S_0
\]
の標準降下データである。
従って $b\otimes_R\mathbb C\simeq\check A_X$ である。
負次数の成分の消滅と次数一生成も、局所再構成補題および忠実平坦性から従う。
ここで降下したのは通常の加群、積、単位のみである。
Grothendieck 圏全体の降下の有効性は仮定していない。

\emph{大域模型と局所代表の比較。}
$b$ は $\check A_X$ の平坦な $R$-線型変形なので、
De Deken--Lowen の変形比較
\cite[Theorem~4.5]{DeDekenLowen}
により $\mathcal D_R=\operatorname{Qmod}(b)$ は
$\operatorname{Qcoh}(X)$ の平坦 Abel 変形である。
その次数零の標準対象を構造対象とする。
定理\ref{thm:abel-cubic-flat-envelope}と
系\ref{cor:abel-universal-relative-cubic}により、これは任意の係数変更後の
相対三次条件・飽和・復元を備える許容な標準模型である。
この段階で、持上げた新しい局所列の section 代数が予め
相対三次条件を満たすことを仮定してはいない。
大域の平坦 Abel Artin 変形を構成してから、その既存の定理を適用している。

局所再構成と係数変更の Gabriel 比較から
\[
 (\mathcal D_R)_S
 \simeq\operatorname{Qmod}(S\otimes_R b)
 \simeq\operatorname{Qmod}(b^S)
 \simeq\mathcal D
\]
を得る。
ここで再構成と実際の遷移同値との両立を直接確認する。
section 関手を
\[
 H_S(M)_i=\operatorname{Hom}_{\mathcal D}(E_{i,S},M)
\]
と書くと、再構成は $\pi H_S:\mathcal D\to\operatorname{Qmod}(b^S)$
という同値である。
$\phi_i$ による前合成と $F_{12}$ の完全忠実性は、
$H_2F_{12}$ と $H_1$ の代数移送を、通常の加群の段階で自然に同定する。
この同定は Hom の合成から定義した $b^S$ の descent 写像と一致する。
従って $\pi$ を適用した後にも、再構成は actual な $F_{12}$ と可換する。
構造対象の descent は変更していないから、これらは marked 自然同型である。

任意の $S$-代数 $T$ と $N\in\mathcal D_T$ に対する自由係数作用の随伴
\[
 \operatorname{Hom}_{\mathcal D_T}(T\otimes_S E_{i,S},N)
 \simeq\operatorname{Hom}_{\mathcal D}(E_{i,S},U(N))
\]
は、右辺をその自然な $T$-作用付き加群と見れば、section 関手と係数作用が
両立することを示す。
$\pi$ の係数変更には補題\ref{lem:abel-arbitrary-coefficient-tails}を適用できる。
よってこの再構成の相対同値は標準係数延長そのものであり、
全てのさらなる係数変更に両立する。
これは補題\ref{lem:abel-canonical-family-equivalences}の延長の一意性とも一致する。
重なり上での同型は marked $2$-自己同型の消滅により coherent であり、
$\mathfrak M^{\mathcal O}_{\mathrm{cub}}$ の射の降下を使うと
$(\mathcal D_R)_\bullet\simeq\xi$ を得る。
還元の同定との両立も、各 $E_i$ とその Hom の還元の構成から従う。
これで \eqref{eq:abel-full-stack-formal-equivalence} の本質的全射性を得る。

完全忠実性は、固定した二つの標準係数族の間の marked 同値と自然同型の
降下、すなわち補題\ref{lem:abel-canonical-family-equivalences}である。
従って構成に使った被覆、持上げ、cocycle の補正の選択にかかわらず、
元の $\xi$ とその射が同じ canonical 部分群oidによって回復される。
ここで比較関手そのものは標準模型から stack 化への自然な関手であり、
これらの選択を用いて定義したものではない。
よって各 Artin 環で示した同値は、Artin 環の任意の射に自然である。
最後に定理\ref{thm:abel-canonical-relative-formal}を合成する。
\end{proof}

\begin{corollary}[全 stack の framed 形式局所環]
\label{cor:abel-full-stack-formal-ring}
特殊ファイバー $X$ との同一視を保った
$\widehat{\mathfrak M}^{\mathcal O}_{\mathrm{cub},X}$ の同型類関手は、
系\ref{cor:abel-cubic-formal-ring}の完備局所環 $\mathcal R_X$ で表される。
特に Krull 次元は $17$、接空間次元は $20$ である。
\end{corollary}
\begin{proof}
前定理と系\ref{cor:abel-cubic-formal-ring}を合成する。
\end{proof}

\begin{remark}[主張の境界]
この定理は、選んだ滑らかな可換三次点における Artin 形式近傍について、
fpqc stack 化が新たな非標準形式を付け加えないことを証明する。
全ての基底上の任意の Grothendieck 圏に対する有効 $2$-降下定理、
非可換な全ての点での同じ局所完全性、代数化、有限型 atlas、
対角の代数性は結論に含めない。
また、特殊ファイバーへの同一視を忘れる場合は、その自己同型の作用を残す必要があり、
この framed 形式局所環だけをそのまま無標識の stack と同一視しない。
\end{remark}

\subsection{framed Artin 変形の自己同型の消滅}

ここでは，特殊ファイバーとの同一視を保つ射を，
構造対象を保つ自然同型で割った $1$-群oid の規約を使う．
全 Abel Artin 比較と section 復元を用いれば，
Hochschild コホモロジーによる一般の射の制御定理を追加せずに，
この群oid の自己同型が自明であることを直接示せる．

\begin{lemma}[三次 section 代数の導分]
\label{lem:abel-cubic-section-derivations}
$k=\mathbb C$，$A=\operatorname{Sym}_k(V)/(F)$ とし，
$\dim_kV=5$，$F\in\operatorname{Sym}^3_kV$ は滑らかな射影三次三様体を定めるとする．
$a=\check A$ をその付随 $\mathbb Z$-代数とし，
$J$ を有限次元 $k$-ベクトル空間とする．
任意の $k$-線型導分
\[
 d_{ij}:a_{ij}\longrightarrow J\otimes_k a_{ij}
\]
で，単位上零かつ $d(xy)=d(x)y+xd(y)$ を満たすものは，
ある族 $\lambda_i\in J$，$\lambda_0=0$ により
\[
 d_{ij}(x)=(\lambda_i-\lambda_j)x
\]
と表される．
\end{lemma}
\begin{proof}
次数一生成により $d$ は
$D_i=d_{i,i+1}\in J\otimes_k\operatorname{End}_k(V)$ により決まる．
二次関係は $\bigwedge^2V$ であるから，$v,w\in V$ について
\[
 (D_i-D_{i+1})(v)w-(D_i-D_{i+1})(w)v=0
 \quad\text{in }J\otimes_k\operatorname{Sym}^2V.
\]
$J$ の基底について成分ごとに考えると，この条件を満たす
$T\in\operatorname{End}_k(V)$ は scalar である．
実際，基底 $x_p,x_q$ を $v,w$ に代入すると，
$T(x_p)x_q-T(x_q)x_p=0$ の係数比較から非対角成分が零となり，
全ての対角成分が等しくなる．
従って $D_i=M+s_i\operatorname{Id}_V$ と書ける．

$\lambda_0=0$ とし，両方向の有限な再帰により
$\lambda_i-\lambda_{i+1}=s_i$ を解く．
導分 $x\mapsto(\lambda_i-\lambda_j)x$ を差し引くと，
全ての隣接成分が同じ $M$ に等しくなる．
この導分は各区間で $\operatorname{Sym}V$ の通常の次数零導分 $D_M$ を誘導する．
三次関係 $F$ を保つことから
\[
 D_M(F)\in J\otimes_k kF
\]
である．

再び $J$ の各成分について，$D_M(F)=cF$ とする．
Euler の等式 $\sum_px_p\partial_pF=3F$ により
$M'=M-(c/3)\operatorname{Id}_V$ は
\[
 \sum_p M'(x_p)\,\partial_pF=0
\]
を満たす．
滑らかさと Euler の等式から，五つの二次式
$\partial_0F,\ldots,\partial_4F$ の共通零点はアフィン空間の原点だけである．
従ってこれらは五変数多項式環の斉次パラメータ系であり，
多項式環の Cohen--Macaulay 性から正則列である．
その Koszul 解消において一次 syzygy の生成元は，係数が次数二の
Koszul syzygy である．係数が次数一の syzygy は存在しないので
$M'(x_p)=0$ が全ての $p$ について成り立つ．
よって $M$ は scalar である．

残った一様 scalar $c'\in J$ も，
$\lambda_i=-ic'$ と置けば
$\lambda_i-\lambda_{i+1}=c'$ により除ける．
これは $\lambda_0=0$ を保つ．
以上の二つの族を加えれば主張を得る．
\end{proof}

\begin{proposition}[framed Artin 自己同型の消滅]
\label{prop:abel-framed-artin-automorphisms}
$X/\mathbb C$ を滑らかな三次三様体とする．
任意の局所 Artin $\mathbb C$-代数 $R$ と，
$\operatorname{Def}^{\mathcal O}_{\rm ab}(\operatorname{Qcoh}(X);R)$ の対象
$(\mathcal D,\mathcal O_R,\iota)$ について，この framed $1$-群oid における
自己同型群は自明である．
\end{proposition}
\begin{proof}
$R$ の長さについて帰納する．$R=\mathbb C$ では frame を保つという条件が，
自己同値を恒等の自然同型類に限定する．
一般の場合，$0\ne J\subset\mathfrak m_R$，
$\mathfrak m_RJ=0$ を満たすイデアルを選び，$R'=R/J$ と置く．
構造対象付き自己同値 $F:\mathcal D\to\mathcal D$ が与えられたとする．
帰納法により，その $R'$-係数変更には marked 自然同型
$\eta':F_{R'}\xRightarrow{\sim}\operatorname{Id}_{\mathcal D_{R'}}$ がある．

全 Artin 比較の標準線対象を $E_i$，$E_0=\mathcal O_R$ とし，
$b_{ij}=\operatorname{Hom}_{\mathcal D}(E_j,E_i)$ と置く．
各 $b_{ij}$ は有限自由であり，$b\otimes_R\mathbb C=\check A_X$ である．
$F$ は係数テンソル積を保つため $F(E_i)$ も平坦であり，
特殊ファイバー上で $E_i$ と同定される．
持上げの同型類の一意性により同型
$\phi_i:F(E_i)\xrightarrow{\sim}E_i$ を選べる．
$i=0$ では $F$ に指定された構造対象の標識を使う．

任意の $i$ で $\phi_i$ の還元と $\eta'_{E_{i,R'}}$ の差は
$\operatorname{End}(E_{i,R'})=R'$ の単元である．
$R^\times\to(R')^\times$ は全射なので，$\phi_i$ を scalar で補正し，
その還元を $\eta'_{E_{i,R'}}$ に一致させる．
$i=0$ では $\eta'$ が marked なので既に一致しており，変更しない．

この選択は section 代数の自己同型
\[
 \alpha_{ij}(x)=\phi_iF(x)\phi_j^{-1}
\]
を与え，$\alpha\equiv1\pmod J$ である．
$\delta=\alpha-1$ と置くと，$J^2=0$ により積に関する導分則が成り立つ．
$\mathfrak m_RJ=0$ と $R$-線型性により $\delta$ は
\[
 d:\check A_X\longrightarrow J\otimes_{\mathbb C}\check A_X
\]
へ因子化する．ここでは $b$ の有限自由性により
$Jb_{ij}\simeq J\otimes_{\mathbb C}(b_{ij}/\mathfrak m_Rb_{ij})$
を用いる．
前補題により
$\delta_{ij}(x)=(\lambda_i-\lambda_j)x$，$\lambda_0=0$ と書ける．

$u_i=1+\lambda_i\in R^\times$ と置けば
$u_ixu_j^{-1}=x+(\lambda_i-\lambda_j)x$ なので，
$\phi_i$ を $u_i^{-1}\phi_i$ に取り換えると section 上の自己同型は恒等になる．
$u_0=1$ であるから構造対象の標識は保たれる．

最後に，この列上の同型を全圏へ延長する．
$H(M)_i=\operatorname{Hom}_{\mathcal D}(E_i,M)$ とすると，
\[
 \operatorname{Hom}(E_i,FM)
 \xrightarrow{\ \circ\phi_i\ }
 \operatorname{Hom}(FE_i,FM)
 \xrightarrow{\sim}\operatorname{Hom}(E_i,M)
\]
は $b$-加群の同型 $HF\simeq H$ を与える．
section 上の作用を恒等にしたことが $b$-線型性を保証する．
tails 商後の $\pi H$ は同値なので，これから
$F\xRightarrow{\sim}\operatorname{Id}_{\mathcal D}$ を得る．
これは $E_0$ 上で指定された標識に一致する．
従って $F$ は marked 自然同型で割った自己同型群では恒等であり，帰納法が閉じる．
\end{proof}

\begin{corollary}[離散な framed 形式近傍]
\label{cor:abel-full-stack-discrete-formal}
定理\ref{thm:abel-stackified-relative-formal}の同値の下で，
framed 形式変形群oid は同型類の離散群oidと自然に同値である．
従って系\ref{cor:abel-cubic-formal-ring}の非標準的な座標同定を選ぶと，
局所 Artin $\mathbb C$-代数上の群oid値関手として
\[
 \widehat{\mathfrak M}^{\mathcal O}_{\mathrm{cub},X}
 \simeq \operatorname{Spf}\mathcal R_X
\]
を得る．右辺は連続な局所 $\mathbb C$-代数準同型の集合を離散群oidとして扱う．
ここで
\[
 \mathcal R_X=
 \frac{\mathbb C[[t_1,\ldots,t_{10},a_{ij}\ (0\le i<j\le4)]]}
 {(a_{ij}a_{kl}-a_{ik}a_{jl}+a_{il}a_{jk}\mid i<j<k<l)}.
\]
従ってこの形式近傍は、十次元形式円板と
$\operatorname{Gr}(2,5)$ のアフィン錐の頂点完備化の積であり、
Krull 次元は $17$、接空間次元は $20$ である。
\end{corollary}
\begin{proof}
自己同型群が自明な群oidでは，同型な二対象の間の射は一意であり，
同型類の離散群oidへの標準関手は同値である．
これは基底変更に自然である．
前命題，全 stack 化後の比較，同型類関手の既存の prorepresentability を合成する．
\end{proof}

\begin{remark}[規約と範囲]
消えたものは，特殊ファイバーへの同一視を保つ自己同値の自然同型類である．
標識を外した特殊ファイバーの有限自己同型群はここに含まれず，
非標識の $2$-群oid における scalar $2$-自己同型についても何も変更していない．
また，この形式表示は代数化や有限型 atlas の存在を意味しない．
定理\ref{thm:jacobian-algebraic-abel-substack}の Jacobian 部分 stack は、
二次包絡の一次中心が階数二となる非可換族から構成されるため、
この可換点そのものを含まない。
両者の次元が $17$ であることから、Jacobian 部分 stack が
全問題の成分または可換点の代数的近傍であるとは結論できない。
導分の計算と全 Artin section 比較による直接証明であり，
$HH^1$ による一般の自己同型制御を未確認のまま引用するものではない．
この直接的適用の初出は未確認である．
\end{remark}

\begin{remark}[一次方向の障害と完全交叉性]
\label{rem:abel-formal-quadratic-obstruction}
この表示の非可換座標を二形式
$a=\sum_{i<j}a_{ij}e_i\wedge e_j$と書くと，方程式は$a\wedge a=0$である．
全ての方程式が二次なので，双対数環上では十個の$a_{ij}$を独立に選べる．
しかし例えば$t=0$，$a_{01}=a_{23}=\varepsilon$，他の$a_{ij}=0$は
$\C[\varepsilon]/(\varepsilon^2)$上の変形を定める一方，
$\C[\varepsilon]/(\varepsilon^3)$上へは持ち上がらない．
実際，任意の持上げの$(0,1,2,3)$-Pfaffianは
$\varepsilon^2$となり，二次の座標補正では消せない．

形式環の埋込み次元は$20$，Krull次元は$17$である．
五つの二次Pfaffianはその次数二部分で線型独立なので，
最小関係数は$5$であり，余次元$3$に等しくない．
従ってこの形式近傍は完全交叉ではない．
これらは既に同定したGrassmannian錐の形式環の具体的な帰結であり，
新しい独立な変形分類を主張するものではない．
\end{remark}

% Classical Artin deformation functors; all categorical equivalences are marked.
% The projective frame removes outer automorphisms, not scalar 2-automorphisms.
\section{可換点におけるambient変形関手の比較}
\label{sec:formal-ambient-comparison}

本節では、非可換射影空間のモジュライを定義する際に現れる
Abel圏、$\mathbb Z$-代数、規格化CY代数の三つの表示を、
可換な$\PP^4$の全Artin変形について比較する。
接空間の一致だけでは比較を結論しない。
有限Beilinson窓の形式sliceを構成し、二つのminimal hullの一意性を用いて、
実際の$\mathsf{Qmod}$自然変換が同型であることを示す。
Poisson座標の選択と、この自然変換自体とは区別する。

$k=\C$とし、$\mathsf{Art}_k$を剰余体$k$の局所Artin $k$-代数の圏とする。
本節の$\mathcal D_{\mathrm{ab},Y}(R)$は
$\mathsf{Qcoh}(Y)$の平坦Abel $R$-変形の同値類であり、
特殊ファイバーとの同値をデータに含める。
変形の間の同値はこの標識と自然同型を除いて可換とする。
$\mathcal D_{\mathbb Z,Y}$は$(\OO_Y(i))_{i\in\mathbb Z}$の
section $\mathbb Z$-代数の線型変形の同値類とする。
いずれも集合値関手であり、自己同値間の自然同型まで保持した
$2$-groupoidを表す記号ではない。

\subsection{豊富列の持上げと構造代数層}

\begin{proposition}[Abel圏、$\mathbb Z$-代数、Poisson変形の同値類]
\label{prop:formal-abel-poisson}
$Y=\PP^4$または滑らかな三次三様体$X\subset\PP^4$とする。
特殊ファイバーの標識を保つ全Artin変形について、自然な全単射
\begin{equation}\label{eq:formal-abel-poisson-chain}
 \mathcal D_{\mathbb Z,Y}
 \simeq \mathcal D_{\mathrm{ab},Y}
 \simeq \mathcal D_{\mathrm{Ass},Y}
 \simeq \mathcal D_{\mathrm{Pois},Y}
\end{equation}
がある。最後の二関手では構造層の結合代数変形、
および可換積も変形させるPoisson変形のgauge同値類を取る。
最後の同定にはYekutieliの量子化を用いる。

各$\OO_Y(i)$は任意のAbel変形に平坦対象として同型を除いて一意に持ち上がる。
$m\le n$なら、その持上げの間のHomは有限自由であり、
任意のArtin基底変換と可換である。
ただし、この一意性は持上げを結ぶ同型の一意性でも、
添字を一つ動かす自己同値の持上げの一意性でもない。
\end{proposition}
\begin{proof}
両方の$Y$について
\[
 H^1(Y,\OO_Y(j))=H^2(Y,\OO_Y(j))=0
 \qquad(j\in\mathbb Z)
\]
である。射影空間では線束のコホモロジー、三次の場合は超曲面完全列で従う。
従って小拡大の核$I$に対する対象$\OO_Y(i)$の持上げの障害群
$I\otimes\Ext^2(\OO_Y(i),\OO_Y(i))$と差の群
$I\otimes\Ext^1(\OO_Y(i),\OO_Y(i))$は消える。
豊富列は$\mathbb Z$-生成列であり、$m\le n$に対し
\[
 \Ext^{1,2}(\OO_Y(m),M\otimes_k\OO_Y(n))=0
 \quad(M\text{ は有限次元 }k\text{-空間})
\]
が成立する。De~Deken--Lowen\cite[Theorem~4.5]{DeDekenLowen}を
適用でき、その逆構成が持上げのHomの平坦性と基底変換を与える。
これで最初の同定を得る。

Lowen\cite[Theorems~3.14, 3.26]{LowenRinged}により、
有限アフィン被覆とその交叉を用いたalgebroid変形がAbel変形を表す。
ここでは$Y$が準コンパクトかつ分離的なので同定理の被覆条件を満たす。
このalgebroidは局所的な$\OO_Y$の平坦持上げから構成される。
$H^2(Y,\OO_Y)=0$により大域的持上げを選べ、
$H^1(Y,\OO_Y)=0$によりその同型類は一意である。
その局所自己準同型層を取れば結合代数層を得る。
これはtwisted変形から通常の変形へ移る
Yekutieli\cite[Corollary~7.5]{YekutieliTwisted}の消滅条件でもある。
したがって同値類について第二の同定を得る。
同論文Corollary~11.2を、滑らかな$Y$、$k=\C\supset\mathbb R$、
$H^1(Y,\OO_Y)=H^2(Y,\OO_Y)=0$に適用して最後の同定を得る。
そのパラメータ代数は局所Artin代数を含み、係数環の準同型に関して自然である。

持上げた$\OO_Y(i)$の自己準同型環は階数一の自由$R$加群で、
単位写像$R\to\End(\widetilde{\OO_Y(i)})$が同型になる。
従って還元が恒等な自己同型$1+\mathfrak m_R$は残る。
また、対象を個別に持ち上げることは、それらを結ぶ自己同値を
整合的に選ぶことを含まない。
\end{proof}

\begin{remark}[選択対象とsmooth proper性]
同様の持上げ論法を一般の選択対象$E$へ使う際に必要なのは
$\Ext^2(E,E)=\Ext^1(E,E)=0$である。
smooth proper性だけではこの消滅は得られない。
例えば滑らかな正次元射影多様体の点のskyscraper層は
通常$\Ext^1$を持ち、点を動かす持上げがある。
本節で$\OO_Y(i)$を使えるのは上記の消滅による。
三次三様体上のこれらの対象を、完全な例外生成列と見なすこともしない。
\end{remark}

\subsection{射影Poisson錐とgauge作用}

$W=H^0(\PP^4,\OO(1))$、$P=\PP(W^*)$とおき、
\[
 \mathcal P(P)=
 \{\pi\in H^0(P,\wedge^2T_P)\mid[\pi,\pi]=0\}
\]
をJacobi方程式で定義したアフィンスキームとする。
その原点完備化を$\widehat{\mathcal P}_0$と書く。
また
\[
 G_1(R)=\ker\bigl(\PGL(W)(R)\longrightarrow\PGL(W)(k)\bigr)
\]
とする。

\begin{proposition}[全Artin上の射影Poisson表示]
\label{prop:formal-projective-poisson-quotient}
自然な同定
\begin{equation}\label{eq:formal-projective-poisson-quotient}
 \mathcal D_{\mathrm{ab},P}(R)
 \simeq\widehat{\mathcal P}_0(R)/G_1(R)
\end{equation}
があり、その接空間は$H^0(P,\wedge^2T_P)$、次元$126$である。
商写像$\widehat{\mathcal P}_0\to\mathcal D_{\mathrm{ab},P}$は
形式的に滑らかで接空間同型を誘導する、すなわちminimal hullである。
\end{proposition}
\begin{proof}
$H^1(P,T_P)=0$だから、$P$の標識付き可換Artin変形は自明である。
これは小拡大ごとに局所自明化の差の$T_P$-値cocycleを消して確認できる。
自明化の差は$G_1(R)$であり、Poisson括弧は
$\mathfrak m_R\otimes H^0(P,\wedge^2T_P)$のJacobi零点となる。
命題\ref{prop:formal-abel-poisson}により
\eqref{eq:formal-projective-poisson-quotient}を得る。

Euler完全列の外積から
\[
 H^0(P,\wedge^2T_P)
 \simeq
 \frac{\Sym^2W\otimes\wedge^2W^*}
 {E\wedge\End(W^*)},\qquad
 \dim=150-(25-1)=126
\]
を得る。ここで$E$はEuler場であり、分母の核は$kE$である。
Jacobi式は二次なので、原点の接空間はこの全空間である。
$G_1(k[\varepsilon]/(\varepsilon^2))$の作用はこの接空間上で自明である。
さらに、全射$R'\to R$に対して$G_1(R')\to G_1(R)$は全射である。
与えられたPoisson代表の還元を一致させるgaugeを$R'$へ持ち上げれば、
商写像の形式的滑らかさが従う。
\end{proof}

\begin{lemma}[発散零の持上げはスキームとして一意]
\label{lem:formal-unimodular-projective}
射影化は、発散零の二次Poisson構造のスキームと
$\mathcal P(P)$の間の同型を与える。
特にこれは非被約Artin基底上でも成立する。
\end{lemma}
\begin{proof}
体積形式を固定し、発散作用素を$\Delta$とする。
線型ベクトル場$Z$について
\[
 \Delta(E\wedge Z)=5Z-\operatorname{tr}(Z)E
\]
である。二次bivectorの発散は跡零の線型場なので、
Euler方向を加減することにより、一意の発散零の持上げ$\Pi$を得る。
この操作は係数環によらない線型分裂である。
射影Jacobiが零なら$J=[\Pi,\Pi]$は$E\wedge Q$と書ける。
$\Delta\Pi=0$から$\Delta J=0$であり、同次BV恒等式は
\[
 0=\Delta(E\wedge Q)=5Q-E\wedge\Delta Q
\]
を与える。両辺に$E$を外積すると$5J=0$、従って$J=0$となる。
これは任意の$k$-代数上の恒等式であるから、Jacobi零点スキームを同定する。
\end{proof}

\begin{remark}[接空間で見えないgauge同一視]
\label{rem:formal-gauge-counterexample}
標識付きPoisson錐そのものを
$\mathcal D_{\mathrm{ab},P}$と同一視することはできない。
$R=k[\varepsilon]/(\varepsilon^3)$とし、射影ベクトル場
\[
 u=x_0\partial_{x_0},\quad v=x_1\partial_{x_1},
 \quad\xi=x_2\partial_{x_0},\quad P_0=u\wedge v
\]
を取る。$[u,v]=0$、$[\xi,P_0]=\xi\wedge v\ne0$である。
二つの異なる点
\[
 \pi=\varepsilon P_0,
 \qquad \pi'=\varepsilon P_0+\varepsilon^2[\xi,P_0]
\]
はともにPoissonであり、$\pi'=\exp(\varepsilon\operatorname{ad}\xi)\pi$だから
同じgauge類を与える。一次の接空間はどちらも126次元でも、
この同一視は二次で初めて現れる。
原点の自己同型方向は$\mathfrak{pgl}_5$である。
従って商groupoidを考えるなら、原点の接複体の次数$-1,0$部分には
$[\mathfrak{pgl}_5\xrightarrow{0}H^0(P,\wedge^2T_P)]$が現れる。
Abel圏の自然自己同型まで保持すると、さらにスカラーの$2$-自己同型を扱う必要がある。
\end{remark}

\subsection{Beilinson窓の射影標識と実際のCY比較}

$\mathcal D^{\mathrm{pf}}_{\mathrm{ab},P}$を、
Abel変形に$\Hom(\widetilde\OO,\widetilde\OO(1))$の
標準基底へ還元する\emph{射影基底}を付けた同値類の関手とする。
このHomは命題\ref{prop:formal-abel-poisson}により階数五の自由加群である。
持上げ対象の同型を変更するとHomの基底は共通スカラー倍だけ変わるので、
射影基底という条件はその選択に依存しない。

\begin{lemma}[有限窓の形式slice]
\label{lem:formal-beilinson-frame}
$\mathcal D^{\mathrm{pf}}_{\mathrm{ab},P}$は完備局所Noetherian環
$\mathcal R_{\mathrm{pf}}$でprorepresentされる。
射影基底を忘れる写像
$\mathcal D^{\mathrm{pf}}_{\mathrm{ab},P}\to\mathcal D_{\mathrm{ab},P}$は
minimal hullである。
\end{lemma}
\begin{proof}
$\OO,\ldots,\OO(4)$の有限線型圏を$\mathfrak b$とする。
そのHom空間は
\[
 H_{ij}=\begin{cases}\Sym^{j-i}W&i\le j,\\0&i>j,\end{cases}
 \qquad 0\le i,j\le4.
\]
標準列は射影空間の幾何的helixをなすので、
De~Deken--Lowen\cite[Theorem~5.15]{DeDekenLowen}は
この有限窓の変形と$\mathcal D_{\mathrm{ab},P}$を比較する。
ここで三次三様体へ同じ有限窓の主張を適用しているのではない。

固定した各$H_{ij}$上の合成写像を変数とし、単位と結合律を課した
有限型アフィンスキームを$M$、通常の対称積を$\mu_0$とする。
基底変更群は
\[
 K=\prod_{0\le i<j\le4}\GL(H_{ij}).
\]
射影基底を固定する部分群$K^{\mathrm{pf}}$は、
$H_{01}=W$への作用がスカラーである元からなる。
対象ごとのスカラー変更
\[
 T=(\mathbb G_m)^5/\mathbb G_m,
 \qquad (\lambda_i)\colon H_{ij}\longmapsto
       (\lambda_j/\lambda_i)H_{ij}
\]
は$M$に自明に作用する中央部分群である。
$H=K^{\mathrm{pf}}/T$は滑らかな群であり、
形式変形の同値類は$\widehat M_{\mu_0}/\widehat H_1$となる。
ここでは関手の自然同型が与える対象スカラーを既に除いている。

$\mu_0$における$H$の無限小固定群が零であることを確認する。
有限圏$\mathfrak b$は隣接Hom
$V_i=H_{i,i+1}=W$、$0\le i\le3$で生成され、
隣接二段の関係は$\wedge^2W\subset V_{i+1}\otimes V_i$で生成される。
圏の導分の$V_i$上の成分を$D_i$とすると、関係の保存から
\[
 (D_{i+1}-D_i)u\mathbin{\odot}v
 -(D_{i+1}-D_i)v\mathbin{\odot}u=0
 \quad(u,v\in W)
\]
を得る。基底ベクトルを代入すると、各$D_{i+1}-D_i$はスカラーである。
射影基底条件で$D_0$もスカラーだから、すべての$D_i$はスカラーとなる。
それらは$T$のLie代数から来る導分であり、$\Lie(H)$の固定部分は零である。

$r=\dim H$とし、合成写像のアフィン空間上の$r$個の線型関数を、
軌道写像の微分上で独立となるように選ぶ。
それらを$\mu_0$での値に固定した$M$内の形式sliceを$S$とする。
形式陰関数定理により、各Artin点$\mu$には一意な
$h\in\widehat H_1$が存在して$h\mu\in S$となる。
実際、$h$の$r$個の形式座標に関するJacobianが特殊ファイバーで可逆であり、
小拡大に沿う帰納で解が存在し一意となる。
従って
\[
 \widehat H_1\times S\simeq\widehat M_{\mu_0},
 \qquad\mathcal D^{\mathrm{pf}}_{\mathrm{ab},P}\simeq S.
\]
この$S$の完備局所環が所要の$\mathcal R_{\mathrm{pf}}$である。

基底を忘れる写像は形式的に滑らかである。
小拡大でAbel変形を持ち上げた後、自由Hom加群の射影基底を持ち上げればよい。
一次では基底変更の24方向はすべて$\mu_0$の射影座標導分から来る。
任意の一次積$\mu_0+\varepsilon\mu_1$に対し
$1+\varepsilon D$は、$D$が$\mu_0$の導分ならその積の自己同型である。
従ってこれら24方向は同値類では消え、忘却の接写像は同型となる。
\end{proof}

Pym\cite[Definition~2.8]{Pym}の規格化Koszul CY変形芽の
可換点完備化を$\widehat{\mathcal Q}_{\CY}$と書く。
生成空間$W$は固定し、これは関係部分空間の形式スキームである。
Artin点は相対的なCY変形を意味し、Artin基底上の絶対大域次元を主張しない。

\begin{theorem}[可換$\PP^4$での全Artin ambient比較]
\label{thm:formal-ambient-cy-comparison}
次数付き代数$B_R$からその$\mathbb Z$-代数と$\mathsf{Qmod}$を取る自然変換
\begin{equation}\label{eq:formal-cy-framed-map}
 \Psi:\widehat{\mathcal Q}_{\CY}
 \longrightarrow\mathcal D^{\mathrm{pf}}_{\mathrm{ab},P}
\end{equation}
は、集合値Artin関手の同型である。
特にすべての$R\in\mathsf{Art}_k$について、実際の$\mathsf{Qmod}$写像は
\begin{equation}\label{eq:formal-ambient-cy-quotient}
 \widehat{\mathcal Q}_{\CY}(R)/G_1(R)
 \xrightarrow{\ \sim\ }
 \mathcal D_{\mathrm{ab},P}(R)
 \simeq\mathcal D_{\mathbb Z,P}(R)
\end{equation}
を与える。
従って可換点の形式近傍では、規格化CY表示と
一般$\mathbb Z$-代数表示の間に追加の古典的Artin変形方向はない。
これは特殊ファイバーの標識を保つ局所主張である。
\end{theorem}
\begin{proof}
まず$B_R$に対応する$\mathbb Z$-代数$\check B_R$は、
各Homが有限自由な$\check{\Sym W}$の線型変形である。
De~Deken--Lowen\cite[Proposition~4.4, Theorem~4.5]{DeDekenLowen}により
$\mathsf{Qmod}(\check B_R)$はAbel変形となり、
同定理の逆構成は指定対象の間のHomを$\check B_R$と同定する。
特に$\Hom(\widetilde\OO,\widetilde\OO(1))=(B_R)_1=W\otimes R$なので、
生成元の基底から射影基底を得る。
これで$\Psi$は基底変換と可換な実際の自然変換として定義された。

命題\ref{prop:formal-projective-poisson-quotient}と
補題\ref{lem:formal-beilinson-frame}により、
$\widehat{\mathcal P}_0$と
$\operatorname{Spf}\mathcal R_{\mathrm{pf}}$は同じ集合値関手のminimal hullである。
従ってその完備局所環は非標準的に同型となる
\cite[Tag~06T5]{StacksFormal}。
ここで必要な一意性は、二つの形式的に滑らかな写像を
Artin商ごとに持ち上げて環準同型を双方向に作り、
接空間同型を持つ合成自己準同型が同型となることからも直接従う。

一方、Pym\cite[Proposition~2.9]{Pym}と
補題\ref{lem:formal-unimodular-projective}により、
$\widehat{\mathcal Q}_{\CY}$の完備局所環も
$\widehat{\mathcal P}_0$のそれと同型である。
この段階ではPymの解析同型の座標同変性も、
その同型とYekutieliの量子化との一致も用いていない。

$d\Psi$は、発散零二次bivectorの射影化である。
この一次の照合では、$R=k[\varepsilon]/(\varepsilon^2)$上で
交換子が$\varepsilon B_R$に入ることを用いる。
第\ref{sec:moduli-artin}節と同じ有限の交換子展開により
各標準一次元の冪をOre局所化できる。
その次数零部分の一次交換子は、通常の商の微分則により
二次bivectorの射影化に等しい。
従って誘導された構造代数層の第一Poisson括弧はこの射影bivectorである。
命題\ref{prop:formal-abel-poisson}の量子化も第一括弧を保つ。
従って補題\ref{lem:formal-unimodular-projective}により
$d\Psi$は126次元空間の同型である。

既に抽象的に同型と分かっている二つの完備局所Noetherian環を同一視すると、
$\Psi$は余接空間上で同型な連続自己準同型となる。
随伴次数環は余接空間で生成されるから各次数の写像は全射であり、
各次数は同じ有限次元空間なので同型である。
完備性と分離性から環準同型自体が同型となり、
\eqref{eq:formal-cy-framed-map}を得る。
接空間同型だけから結論したのではなく、minimal hullの比較で
両環の抽象的同型を先に確立した点が必要である。

$\Psi$は生成元の射影座標変更と射影基底変更に関して同変である。
すべての射影基底は$G_1(R)$の作用で移り合うので、
それを忘れた商が\eqref{eq:formal-ambient-cy-quotient}となる。
\end{proof}

\begin{corollary}[37次元成分と中心incidenceの形式像]
\label{cor:formal-113-component}
命題\ref{prop:quantum-113-component}の$Q_{113}$の可換頂点完備化は、
$\mathcal D^{\mathrm{pf}}_{\mathrm{ab},P}$の形式局所模型の
37次元既約成分を与え、その一般点において全局所環は被約かつ正則である。
命題\ref{prop:formal-abel-poisson}を通して、これは射影基底付き
$\mathbb Z$-代数変形の形式局所模型の成分でもある。

$Q_{113}$の可換頂点の近傍を、普遍二次代数の次数三・四部分が
局所自由となるように取る。
その上の$\PP(\mathscr B_3)$内の中心直線の閉incidenceにおいて、
既存の支配的成分の閉包を$I_{113}$とする。
可換ambient上の\emph{全fiberに沿って}これを形式完備化し、
実際の$\mathsf{Qmod}$写像を合成すると、像の形式閉包は上記の成分全体となる。
従ってこの対応成分内での像の余次元は零である。
この完備化のArtin点は正則中心三次商を与えるので、
正規三次超曲面型のモジュライへの写像として解釈できる。
ここでは一つの指定した三次式$F$での対の完備化だけを
取っているのではなく、商のsmooth proper性を追加した部分関手への
同じ主張もここからは結論しない。
\end{corollary}
\begin{proof}
$Q_{113}$はgraph座標で正次数の斉次整域のスペクトルである。
頂点完備化における非零元の最初の斉次項を取れば、
その積が非零であることが分かるので完備化も整域である。
全Poissonスキームは有限型$k$-スキームであるから、その局所環はexcellentであり、
完備化への写像はregularである\cite[Tag~07QS]{StacksFormal}。
Pymの形式比較を通すと、命題\ref{prop:quantum-113-component}の
成分性と一般正則性がこの形式完備化へ移る。
さらに定理\ref{thm:formal-ambient-cy-comparison}の
実際の同型$\Psi$で主張した形式模型の成分となる。

次数三・四の普遍成分は有限表示加群であり、
可換点での完備化後は平坦な二次変形の対応次数成分として自由である。
完備化の忠実平坦性によって可換点の局所環上でも自由となるため、
局所自由となる近傍を取れる。
選んだ近傍では、中心条件は
$\mathscr L\otimes\mathscr B_1\to\mathscr B_4$という
交換子写像の零点として$\PP(\mathscr B_3)$内で閉に定義される。
ここで$\mathscr L$はtautological直線束である。
その支配的既約成分の閉包を取るので、
$f:I_{113}\to Q_{113}$は、この近傍へ制限して射影的かつ支配的である。
$Q_{113}$は被約なので
$\OO_{Q_{113}}\to f_*\OO_{I_{113}}$は単射である。
可換点の局所化と完備化は平坦であり、
射影的写像の大域切断はこの平坦基底変換と可換である
\cite[Tag~02KH]{StacksFormal}。
従って完備化後もこの単射を保つ。
形式関数定理\cite[Tag~02OC]{StacksFormal}により、
同じ結論を全特殊fiberに沿う形式完備化で表せる。
これが像の形式閉包が全成分であることを意味する。
最後に$\Psi$を合成してもこの性質は変わらない。

正則性も全Artin点で確認できる。
特殊ファイバーの中心直線は非零の三次式$F\in\Sym^3W$である。
任意の局所Artin持上げ$R$で直線の生成元$g$を選ぶと、
各次数の左右乗法は、還元すると多項式環での
$F:S_{n-3}\hookrightarrow S_n$となる。
$F\ne0$なのでこれは単射であり、最大rankの小行列式が単元に持ち上がるため、
元の有限自由$R$加群間の乗法も分裂単射である。
incidence条件から$g$は中心なので、両側イデアルは$gB_R=B_Rg$となり、
その商は各次数で自由、Hilbert級数$(1-t^3)/(1-t)^5$を持つ。
従ってこの形式完備化が与える対は実際に正則正規三次商の条件を満たす。
この段階では$F$の滑らかさを仮定していない。
この議論は具体的な量子incidenceと$\mathsf{Qmod}$を使っており、
独立に選んだPoisson座標どうしの可換性を必要としない。
\end{proof}

\begin{remark}[形式同型の意味と大域化]
定理\ref{thm:formal-ambient-cy-comparison}の$\Psi$は
$\mathsf{Qmod}$から定まる自然変換である。
その証明に使ったPoisson座標やminimal hull間の同型は選択に依存し得る。
また、同定理は任意の非可換点での形式近傍、全AS正則
$\mathbb Z$-代数の有限型モジュライ、あるいはそれらの大域成分を同定しない。
相対的な周期同型の選択をすべて保持したstackの同型も主張しない。
集合値の同型から高次自己同型を保つ同値を推論してはいけない。
\end{remark}

\subsection{三次三様体からのambient写像}

\begin{corollary}[可換三次点での実際のambient変形写像]
\label{cor:formal-cubic-ambient-map}
滑らかな三次三様体$X=V_+(F)\subset P$について、
定理\ref{thm:artin-central-pair}および
定理\ref{thm:artin-quadratic-cover}の同次formality入力の下で、
二次包絡は自然変換
\begin{equation}\label{eq:formal-cubic-ambient-map}
 \mathcal D_{\mathrm{ab},X}
 \longrightarrow\mathcal D_{\mathrm{ab},P}
 \simeq\widehat{\mathcal Q}_{\CY}/G_1
\end{equation}
を与える。
左辺は$\mathcal R_X$、すなわち
\eqref{eq:artin-poisson-sheaf-ring}の17次元・接空間20次元の環で表される。
原点における接写像は
\[
 H^1(X,T_X)\oplus H^0(X,\wedge^2T_X)
 \longrightarrow H^0(P,\wedge^2T_P),
 \qquad(\dot F,a)\longmapsto\Pi(F,a)_{\mathrm{proj}}.
\]
その核は10次元の可換三次変形方向であり、像は10次元である。
\end{corollary}
\begin{proof}
命題\ref{prop:formal-abel-poisson}と
定理\ref{thm:artin-quadratic-cover}により、
Abel変形は中心三次商の変形で表示される。
そのsection $\mathbb Z$-代数はDe~Deken--Lowenの逆構成によるものと同一視される。
従って二次包絡の内在的復元はこの表示から得た$B_R$の
$\mathbb Z$-代数に一致し、平坦なambient変形を与える。
これに$\mathsf{Qmod}$を取り、
定理\ref{thm:formal-ambient-cy-comparison}を用いれば
\eqref{eq:formal-cubic-ambient-map}を得る。
規格化CY側への同定は、存在するambient圏に同定理の逆写像を適用する意味であり、
任意の次数付き表示が既に規格化CYであると仮定していない。

一次では、$F+\varepsilon\dot F$だけを動かしても二次包絡は可換多項式環である。
Poisson方向の一次交換子は
$\iota_{\Pi(F,a)}\Omega=dF\wedge a$なので上記の式となる。
その射影化を$X$へ制限すると
$a\mapsto\iota_Ea$という
\eqref{eq:artin-relative-bivectors}の同型を得るから、
$a$方向の写像は単射である。
\end{proof}

\begin{remark}[Poisson座標での式と比較の範囲]
相対Poisson計算自体は、全Artin上の自然写像
\[
 \mathcal D_{\mathrm{Pois},X}\longrightarrow
 \widehat{\mathcal P}_0/G_1,
 \qquad (F_R,a)\longmapsto\Pi(F_R,a)_{\mathrm{proj}}
\]
を与える。この写像と\eqref{eq:formal-cubic-ambient-map}を、
独立に選んだ二つの量子化同定によって作った正方形の中で
同じものと主張するには、量子化とCasimir商・局所化の両立を照合する必要がある。
本節はその未確認の両立を証明の前提にしていない。
実際のambient写像の存在、全Artin ambient比較、および接写像の式は上で得られている。
また、ここでの接写像の10次元の核は、非零一般点における
ambient忘却の4次元ファイバーとは異なる。
原点のJacobi方程式が二次で始まることを無視して両者を比較してはいけない。
\end{remark}

\section{Poisson構造付き三次への持上げと形式近傍}
\label{sec:poisson-cubic-lifts}

定理\ref{thm:kuznetsov-cubic-moduli-morphism}の射を
$\rho:\mathfrak U\to V$ とし、$V$ は自己同型が自明な滑らかな三次の
モジュライの非空開とする。本節では Poisson 構造の尺度を記憶する。
構造を非零スカラー倍で同一視した射影 Poisson モジュライとは区別する。
中心となるのは、可換三次点の全Abel Artin近傍において、
実際の残余操作から得られる射とPoisson構造の忘却射を比較することである。
系\ref{cor:abel-poisson-formal-lift}では、両者に両立する形式同型を構成し、
残余射の形式ファイバーを$\operatorname{Gr}(2,5)$のアフィン錐の完備化と同定する。

有限型の開上でも、相対有理座標を選べば支配的なPoisson持上げが得られる。
こちらは相対次元七とPoisson錐の有理性を使う存在結果であり、
形式同型の代数化や交換子からの逆量子化とは別の構成である。
最後に、中心三次の対称化逆像がPoisson超曲面を定めない具体例を示し、
量子中心と古典Casimirを比較する際に必要な補正を明らかにする。

\subsection{相対Poisson錐とその有理座標}

\begin{proposition}[相対Poisson錐の構造]
\label{prop:relative-poisson-cubic-cone}
$f:X\to V$ を普遍的な滑らかな三次族とする。
相対 Poisson 構造の関手は、$V$ 上アフィン有限表示な錐
$\mathcal P\to V$ により表される。その零切断の補集合を
$\mathcal P^\times$ と書く。
\begin{enumerate}
\item $V$ の Zariski 開被覆上で
\[
 \mathcal P\simeq
 V\times\operatorname{Spec}
 \frac{\mathbb C[a_{ij}\mid0\le i<j\le4]}
 {(a_{ij}a_{kl}-a_{ik}a_{jl}+a_{il}a_{jk}\mid i<j<k<l)}.
\]
従って $\mathcal P^\times\to V$ は滑らかで相対次元七である。
その射影化は相対次元六の Grassmannian 束である。
\item ある非空開 $V_0\subset V$ 上で、$\mathcal P^\times$ は開部分
\[
 \mathcal W\simeq V_0\times\mathbb A^6\times\mathbb G_m
\]
を持つ。
\end{enumerate}
\end{proposition}
\begin{proof}
まず族の射影座標について確認する。普遍三次は、
射影座標 torsor から降下した Brauer--Severi $\mathbb P^4$ 束
$p:\mathcal B\to V$ の有効 Cartier 因子 $D$ として存在する。
線束
\[
 \mathcal H=\omega_{\mathcal B/V}\otimes\mathcal O_{\mathcal B}(2D)
\]
の各幾何学的ファイバーへの制限は $\mathcal O_{\mathbb P^4}(1)$ である。
これは $-5+2\cdot3=1$ による。
従って $p_*\mathcal H$ は階数五の局所自由層であり、評価写像は
$\mathcal B$ をその射影束と同定する。
この主張は Brauer--Severi 束を自明化する étale 被覆上で通常の
$\mathcal O(1)$ の計算に帰着し、射自体は評価によって大域的に定まる。
特にこの三次族の ambient に Brauer 障害は残らない。

$V$ を Zariski 局所化し、この階数五の束と三次方程式の基底直線束を
自明化すると、$X=V(F)\subset\mathbb P^4_V$ と書ける。
体積形式と $F$ の尺度もこの表示に従って固定する。
相対余法完全列、超曲面制限完全列、Euler 完全列から
\[
 \mathcal E=f_*\wedge^2T_{X/V}\simeq\mathcal O_V^{10},
 \qquad
 \mathcal F=f_*\wedge^3T_{X/V}\simeq\mathcal O_V^{15}
\]
を得る。高次直像は零であり、これらの同定は任意の係数変更に両立する。
具体的には
$\wedge^2T_{X/V}\simeq\Omega^1_{X/V}(2)$、
$H^0(\Omega^1_{\mathbb P^4/V}(2))\simeq\wedge^2\mathcal O_V^5$
を用いる。基底の直線束をまだ自明化していない場合は、対応する基底twistを残す。

相対 Schouten 括弧は $\operatorname{Sym}^2\mathcal E\to\mathcal F$ を与える。
その二次写像 $\pi\mapsto[\pi,\pi]$ の零点schemeを $\mathcal P$ と定義する。
これにより族と底変更に関する表現性が得られる。

方程式を確認する。$\epsilon_{ij}=x_jdx_i-x_idx_j$ を対応する
$\Omega^1(2)$ の基底とし、$\pi=\sum_{i<j}a_{ij}\epsilon_{ij}$ と書く。
Mayanskiy \cite[\S2.2, Lemma~2 and Theorem~1]{Mayanskiy} の計算は
\[
 \begin{aligned}
 \frac14[\pi,\pi]
 &=q_{1234}\partial_0F-q_{0234}\partial_1F
  +q_{0134}\partial_2F-q_{0124}\partial_3F
  +q_{0123}\partial_4F,\\
 q_{ijkl}&=a_{ij}a_{kl}-a_{ik}a_{jl}+a_{il}a_{jk}.
 \end{aligned}
\]
を与える。これは係数の多項式恒等式なので相対族でも成立する。
五つの偏微分の張る写像
$\mathcal O_V^5\to\mathcal O_V^{15}$ は各幾何学的ファイバーで単射である。
もし定数係数の非零な線型関係があれば、線型座標変更の後で
$F$ が一つの変数に依存しなくなり、その三次超曲面が特異になるためである。
従ってこの写像は局所的に分裂単射である。
ゆえに Schouten 方程式の係数イデアルは、五つの $q_{ijkl}$ のイデアルと一致する。
非被約基底への変更後も左逆が残るため、集合論的な一致にとどまらない。
これは第\ref{sec:moduli-artin}節で用いた相対計算の
通常の有限型基底上での形でもある。

非零 Pl\"ucker 錐は Grassmannian 上の tautological 直線束の零切断の補集合であり、
滑らかで次元 $6+1=7$ である。
さらに $a_{01}\ne0$ の chart では、二形式は一意に
\[
 a=\lambda
 (e_0+u_2e_2+u_3e_3+u_4e_4)
 \wedge
 (e_1+v_2e_2+v_3e_3+v_4e_4)
\]
と表される。ここで $\lambda=a_{01}\in\mathbb G_m$ であり、
残り六つの係数は任意である。これが後半の chart を与える。
\end{proof}

\subsection{全Abel Artin変形上の残余射とPoisson形式近傍}
\label{subsec:poisson-formal-residual}

以下では滑らかな複素三次三様体 $X$ を固定し、特殊ファイバーの同定を
持つ通常の局所 Artin $\mathbb C$-代数上の変形を考える。
構造対象付き Abel 変形の群oidは、標識を保つ自然同型で射を割った規約とする。
第\ref{sec:abel-formal-descent}節によりこれは同型類の離散群oidと同値である。

\begin{proposition}[全Abel形式近傍上の実際の残余射]
\label{prop:abel-artin-actual-residual}
任意の構造対象付き Abel Artin 変形 $\mathcal D_R$ に対し、
proper な $R$-線型 dg 圏 $\mathcal K_R$ を自然に構成できる。
特殊ファイバーは $\operatorname{Ku}(X)$ であり、その形成は任意の Artin
係数変更に両立する。従って自然写像
\[
 \begin{aligned}
 r:\operatorname{Def}^{\mathcal O}_{\mathrm{ab}}(\operatorname{Qcoh}(X))
 &\longrightarrow
 \pi_0\operatorname{CatDef}^{c}_{\operatorname{Ind}\operatorname{Ku}(X)}\\
 &\xrightarrow{\ \sim\ }\operatorname{Def}_X.
 \end{aligned}
\]
がある。最後の同型は命題\ref{prop:kuznetsov-formal-recognition}による。
可換三次変形に制限した $r$ は、その三次変形自身を返す。
\end{proposition}

\begin{proof}
定理\ref{thm:abel-cubic-flat-envelope}\textup{(iii)}により、
$\mathcal D_R$ を次数別有限自由な通常の二次代数 $B_R$ と
中心正則三次元 $w_R$ によって
\[
 A_R=B_R/(w_R),\qquad
 \mathcal D_R\simeq\operatorname{QGr}(A_R)
\]
と表示できる。この表示は全 Abel Artin 変形について存在する。
一般の非周期的 $\mathbb Z$-模型をそのまま中央 CY 模型と仮定するのではなく、
上の既存比較定理を用いて通常の中心商模型を選ぶ。
本証明では $B_R$ の相対 CY orientation は必要としない。

まず compact generation を確認する。
系\ref{cor:abel-universal-relative-cubic}により、$B_R$ には任意係数変更後にも
完全な長さ五の有限二次 Koszul 双加群解消がある。
これは補題\ref{lem:relative-koszul-torsion-model}の入力\textup{(1)--(2)}を満たす。
同補題は CY 同型を要求しないので、任意の可換 $R$-代数 $S$ に対し
\[
 T_{B,S}=\bigoplus_{i=0}^{4}\mathcal O_{B_S}(i)
\]
は $D(\operatorname{QGr}B_S)$ の compact generator である。
有限二次双加群解消の普遍完全性は、各内部次数で有限自由な複体の還元が
多項式環の Koszul 対角解消であることからも確認できる。
その augmentation の余核も係数自由であるため、有限 Artin 基底上の
分裂持上げにより各内部次数の完全部分は分裂し、任意の $S$ への変更を保つ。

商 $B_S\to A_S$ に沿う restriction は
\[
 j_*:D(\operatorname{QGr}A_S)\longrightarrow D(\operatorname{QGr}B_S)
\]
を定め、完全、保守的で、任意直和を保つ。
実際、$A_S$-加群の $A_S$-torsion と $B_S$-torsion は一致し、
次数付き restriction は完全かつ保守的なので Abel 商へ降りる。
その導来関手も cohomology 上の保守性から保守的である。
非有界 Gabriel 商を使う点も確認する。
補題\ref{lem:relative-cubic-torsion-comparison}の証明で使うのは
有限 Koszul 解消と中心正則元であり、CY 同型ではない。
その証明を適用すると、$A_S$-injective $J$ に対する $B_S$ 上の高次 torsion は、
$\operatorname{Tor}_1^{B_S}(B_S/B_{S,\ge n},A_S)$ の遷移が
$n$ を三増やすと零になることにより消える。
従って $A_S$ の derived torsion は $B_S$ の有限幅の模型で計算でき、
有限 cohomological dimension と直和保存を持つ。
これにより $D(\operatorname{QGr}A_S)$ は
$D(\operatorname{Gr}A_S)$ の cohomologically torsion な部分圏による商として扱える。
中心正則三次による長さ一の解消から
$-\otimes_{B_S}^{\mathbf L}A_S$ は torsion cohomology を torsion cohomology へ
送るため、左随伴 $\mathbf Lj^*$ が tails の非有界導来圏上で定まる。
従って
\[
 T_{A,S}=\mathbf Lj^*T_{B,S}
         =\bigoplus_{i=0}^{4}\mathcal O_{A_S}(i)
\]
は compact である。$T_{A,S}$ の全 shift に直交する $N$ に対し、随伴により
$j_*N$ は $T_{B,S}$ の全 shift に直交するので零である。
保守性から $N=0$ となり、$T_{A,S}$ は compact generator である。
この議論は特殊ファイバーでの生成性だけから一般の生成性を推論していない。

命題\ref{prop:artin-ore-universal-cohomology}は、すべての線対象間の
RHom の係数 perfect 性と任意の導来底変更を与える。
従って $\mathcal C_S=D(\operatorname{QGr}A_S)^c$ は proper であり、
\[
 L_i=\mathcal O_{A_S}(i),\quad
 \mathbf R\operatorname{Hom}(L_0,L_0)=
 \mathbf R\operatorname{Hom}(L_1,L_1)=S,\quad
 \mathbf R\operatorname{Hom}(L_1,L_0)=0.
\]
dg evaluation による有限 mutation
\[
 P_S=\mathbb L_{L_0}\mathbb L_{L_1}
\]
を用い、$\mathcal K_S=\langle L_0,L_1\rangle^\perp$ と置く。
定理\ref{thm:relative-kuznetsov-projector}の左射影の証明は例外性と
propernessだけを用いるので、この場合にも適用できる。
特に
\[
 G_S=P_S(L_2\oplus L_3\oplus L_4),\qquad
 \mathcal K_S=\operatorname{thick}(G_S)
             \simeq\operatorname{Perf}(E_S),\quad
 E_S=\mathbf R\operatorname{End}(G_S).
\]
有限 mutation と線対象の Hom 比較から $E_S$ は $S$-perfect で、
任意の $S\to S'$ に対して
$E_S\otimes_S^{\mathbf L}S'\simeq E_{S'}$ が積と両立する。
ここでは右随伴、相対 Serre 同型、$\mathcal C_S$ 自体の smoothness を
仮定していない。

模型からの独立性と自然性を確認する。
各 Artin 基底上で $L_1$ は $\mathcal O_X(1)$ の平坦持上げである。
定理\ref{thm:abel-cubic-all-artin}に使った持上げの一意性により、
特殊ファイバーの同定を保つ完全 Abel 同値は $L_1$ を同型を除いて保つ。
この比較の選択差は $\operatorname{End}(L_1)^\times=R^\times$ であり、
評価項 $\mathbf R\operatorname{Hom}(L_1,M)\otimes_R L_1$ 上で相殺する。
$L_0$ は指定した構造対象である。
従って両 mutation とその直交部分圏は全ての framed O-marked 同値に自然である。
任意の Artin 係数変更にも同じ一意性と標準の評価写像が適用できる。
以上により、特定の通常次数付き表示やその周期構造の選択を
モジュライのデータに追加せずに $\mathcal K_R$ が定まる。

特殊ファイバーでは通常の例外対を用いるので
$\mathcal K_k=\operatorname{Ku}(X)$ である。
実際、
\[
 \operatorname{Ind}(\mathcal K_R)=\operatorname{Mod}(E_R),\qquad
 \operatorname{Mod}(E_R)\otimes_{\operatorname{Mod}(R)}
                 \operatorname{Mod}(k)
 \simeq\operatorname{Mod}(E_R\otimes_R^{\mathbf L}k)
 \simeq\operatorname{Ind}\operatorname{Ku}(X)
\]
という指定した Morita framing がある。
従ってこれは実際の compactly generated 変形であり、
命題\ref{prop:kuznetsov-formal-recognition}を適用できる。
同命題の tame 条件は固定特殊圏の条件であり、変形源の smoothness は要求しない。
可換三次族の場合には構成そのものがその族の相対 Ku だから、最後の主張も従う。
\end{proof}

\begin{corollary}[残余射に両立するPoisson形式近傍への同型]
\label{cor:abel-poisson-formal-lift}
固定した $X$ の上で、Poisson 構造を零に特殊化する変形の形式モジュライを
$\widehat{\mathfrak P}_{(X,0)}$ とする。非標準的な同型
\[
 \widehat{\mathfrak M}^{\mathcal O}_{\mathrm{cub},X}
 \xrightarrow{\ \sim\ }
 \widehat{\mathfrak P}_{(X,0)}
\]
を、前命題の残余射 $r$ と Poisson 構造を忘れる射に両立するように選べる。
この同型は可換部分上で恒等であり、対応する $r$ の形式ファイバーは
$\operatorname{Gr}(2,5)$ のアフィン Pl\"ucker 錐の頂点完備化である。
ファイバーの次元は七、頂点接空間の次元は十である。
\end{corollary}
\begin{proof}
\[
 A_0=\mathbb C[[a_{ij}\mid0\le i<j\le4]]/
       (\operatorname{Pfaff}_4(a)),\qquad
 R_X=A_0[[s_1,\ldots,s_{10}]]
\]
と置く。補題\ref{lem:artin-relative-poisson-sheaf}により
$\widehat{\mathfrak P}_{(X,0)}$ は $R_X$ で表され、
Poisson 構造を忘れる射は $s$-座標への射である。
補題\ref{lem:artin-constructed-family-comparison}と全 Abel 比較による
量子化同型 $Q:\widehat{\mathfrak P}_{(X,0)}\to
\widehat{\mathfrak M}^{\mathcal O}_{\mathrm{cub},X}$ を一つ選ぶ。
構成により $Q$ は $a=0$ 上で通常の可換三次族を返す。

合成 $rQ$ を十次元幾何的形式座標で
$t_i=f_i(s,a)$ と書けば、前命題から $f_i(s,0)=s_i$ である。
従って $s$ に関するJacobianは原点で恒等行列である。
完備係数環 $A_0$ 上の形式逆写像定理により、
\[
 \Theta:(s,a)\longmapsto(f(s,a),a)
\]
は $\operatorname{Spf}R_X$ の自己同型である。
$A_0$ が特異でも、この定理は $s$ の線型係数行列の可逆性から
極大イデアル次数について帰納して証明できる。
求める同型は $\Theta Q^{-1}$ である。
$a$ を固定しているためPl\"ucker方程式は変わらず、
Poisson忘却との両立と $a=0$ 上での恒等性も従う。
\end{proof}

\begin{remark}[形式的同型の範囲]
量子化比較と形式slice、Poisson係数の自明化を選んでいるため、上の同型を
選択に依存しない脱量子化写像とは呼ばない。
また元の量子化比較 $Q^{-1}$ が残余射とそのまま両立するとは証明しておらず、
そのために幾何的座標を $f(s,a)$ へ変更した。
形式的な無限級数の存在は代数的射の存在や収束を保証しない。
実際の十七次元開集合上で作る非標準的Poisson構造の選択と、
この特定の形式同型の代数化・一致は別の主張である。
\end{remark}

\subsection{相対有理座標による開集合上の持上げ}

\begin{proposition}[相対有理座標による非標準的Poisson持上げ]
\label{thm:noncanonical-poisson-cubic-lift}
$\mathfrak U$ と $V$ をさらに非空開へ縮小すると、
\[
 \widetilde\rho:\mathfrak U\longrightarrow\mathcal P^\times,
 \qquad p\circ\widetilde\rho=\rho
\]
となる射が存在する。しかも選択を行えば、次を同時に満たせる。
\begin{enumerate}
\item $\widetilde\rho$ は支配的であり、源と標的の有理函数体の拡大は有限分離である。
\item ある十七次元 scheme の étale atlas $T\to\mathfrak U$ に対し、
$T\to\mathcal P^\times$ は étale である。
\end{enumerate}
ここでは $\widetilde\rho$ の representability を主張しない。
源の有限 stabilizer が残る可能性があり、函数体上の有限性は
射全体の有限性や次数一の双有理性を意味しない。
\end{proposition}
\begin{proof}
命題\ref{prop:relative-poisson-cubic-cone}の chart $\mathcal W\simeq V_0\times\mathbb A^6\times\mathbb G_m$ を取る。
$\rho$ が支配的なので、源を $\rho^{-1}(V_0)$ へ制限してよい。
単に非零持上げの存在を得るには、この chart の一定の非零点を選び、
$\rho$ と合成すれば足りる。以下ではさらに支配的な持上げを構成する。

$K=\mathbb C(V)$、$L=\mathbb C(\mathfrak U)$ とする。
後者は商表示 $\mathfrak C_{\mathrm J}=[J/\operatorname{PGL}_5]$ における
不変有理函数体 $\mathbb C(J)^{\operatorname{PGL}_5}$ である。
系\ref{cor:central-incidence-rational-dimensions}の Rosenlicht 商の次元計算と
$\rho$ の支配性から
\[
 \operatorname{trdeg}_K L=17-10=7.
\]
有限生成体拡大なので、超越基底 $t_1,\ldots,t_7\in L$ を取ると
$L/K(t_1,\ldots,t_7)$ は有限であり、標数零なので分離的である。

各 $t_i$ を $J$ 上の不変有理関数として表す。
有限個の極を除いた共通の定義域は非空不変開に取れるため、
商stackの非空開上の通常の関数へ降下する。
さらに $t_7$ が単元となるように縮め、命題\ref{prop:relative-poisson-cubic-cone}の chart で
\[
 (u_2,u_3,u_4,v_2,v_3,v_4,\lambda)
 =(t_1,t_2,t_3,t_4,t_5,t_6,t_7)
\]
と置く。これが $\rho$ 上の射 $\widetilde\rho$ を与える。
函数体の準同型は $K(t_1,\ldots,t_7)\hookrightarrow L$ なので、
この射は支配的で、主張した有限分離拡大を誘導する。

命題\ref{prop:jacobian-affine-four-morphism}の十七次元 étale atlas を
$\mathfrak U$ に制限し、その支配的な既約成分を $T$ とする。
$\mathbb C(T)/L$ は有限分離拡大である。
従って合成 $T\to\mathcal P^\times$ は支配的かつ一般点で étale であり、
非空 étale locus $T^e\subset T$ を持つ。
その $\mathfrak U$ への像は非空開であり、
$T^e$ はその像の étale atlas となる。
これに制限すれば二番目の主張を得る。
\end{proof}

\begin{remark}[構成の意味と未解決部分]
この持上げの存在は、$\rho$の相対次元と相対Poisson錐の有理性から従う。
座標、非零 Pl\"ucker chart、七つの有理関数の選択に依存し、
非可換交換子や半古典極限から指定された Poisson 構造を抽出してはいない。
選択を固定した後には商stack上の実際の射であるから、
全 $\mathcal O$ 保存 Abel 同値と底変更に両立する。
未解決なのは補助的な選択からの独立性、非可換構造の量子化としての意味、
および系\ref{cor:abel-poisson-formal-lift}の形式同型との両立である。
とくに七次元という次元の一致と上の有限分離拡大から、
$\rho$ の一般fiberが Poisson 錐と双有理であるとは結論しない。
\end{remark}

\subsection{量子中心と古典三次の同定に必要な補正}

\begin{proposition}[中心三次の対称化逆像は一般にPoisson超曲面を与えない]
\label{prop:poisson-naive-symmetrization-fails}
標数零体 $k$ と不定元 $h$ を取り、$K=k(h)$ とする。
$u,v$ を中心とする五変数二次代数
\[
 B_h=K\langle x,y,z,u,v\rangle/
 \bigl([u,-],[v,-],[x,y]-hz^2,[y,z]-hx^2,[z,x]-hy^2\bigr)
\]
を考える。$f=(x^3+y^3+z^3)/3$、
$T=xyz+xzy+yxz+yzx+zxy+zyx$ と置く。
中心三次
\[
 w_h=f-\frac{h^2}{9}T+u^3+v^3
\]
の完全対称化写像 $\operatorname{sym}:\operatorname{Sym}^3K^5\to(B_h)_3$
による逆像は
\[
 H_h=f-\frac{2h^2}{3}xyz+u^3+v^3
\]
である。しかし、二次関係の graph から復元する Poisson 括弧
\[
 \{x,y\}=hz^2,\qquad\{y,z\}=hx^2,\qquad\{z,x\}=hy^2,
 \qquad\{u,-\}=\{v,-\}=0
\]
について、$H_h$ は Casimir でなく、$(H_h)$ は Poisson ideal でもない。
従って、中心三次直線を完全対称化の逆写像で可換三次直線へ送るだけでは、
Poisson 構造付き三次への写像を定義できない。
\end{proposition}
\begin{proof}
$w_h$ の中心性は定理\ref{thm:pair-full-fermat-center}の
$L=e=m=0$ の場合である。
全次数三の対称化は、$h=0$ で通常の対称代数への同型となる。
相対次数成分の平坦性と有限表示から、その行列式は零関数でなく、
函数体 $K$ 上では同型である。
$\operatorname{sym}(xyz)=T/6$ なので、表示した $H_h$ が一意な逆像である。
一方、$f$ と $u,v$ は上の Poisson 括弧の中心に属し、
\[
 \{x,xyz\}=h x(z^3-y^3),
 \qquad
 \{x,H_h\}=-\frac{2h^3}{3}x(z^3-y^3)\ne0.
\]
この次数四の多項式が $(H_h)$ に入るなら、
$\{x,H_h\}=\ell H_h$ となる一次式
$\ell=a_xx+a_yy+a_zz+a_uu+a_vv$ が存在する。
右辺の $u^4,u^3v,u^3x,u^3y,u^3z$ の係数は
それぞれ $a_u,a_v,a_x,a_y,a_z$ であるが、左辺は $u$ を含まない。
従って $\ell=0$ となり矛盾する。
射影化後にも失敗することを確認する。
$u\ne0$ の chart で
$X=x/u,\ Y=y/u,\ Z=z/u,\ V=v/u$ と置くと、
\[
 \begin{aligned}
 \overline H=H_h/u^3
 &=\frac{X^3+Y^3+Z^3}{3}-\frac{2h^2}{3}XYZ+1+V^3,\\
 \{X,\overline H\}&=-\frac{2h^3}{3}X(Z^3-Y^3).
 \end{aligned}
\]
$u$ が Poisson 中心なので、右の括弧はこの chart の括弧である。
$\overline H$ は $V$ に関して monic 三次であるが、括弧は $V$ を含まない
非零多項式である。従って $(\overline H)$ に属さず、
射影超曲面 $V_+(H_h)$ もこの ambient 括弧の Poisson 部分多様体ではない。
なお $V_+(H_h)$ は一般の $h$ で滑らかである。
実際 $h=0$ では五変数 Fermat 三次なので、その判別式は零多項式ではない。
この反例の失敗は可換三次の特異性によるものではない。
\end{proof}

\begin{remark}[この障害の範囲]
二次関係の graph から ambient の Poisson 構造を復元する操作自体は
$\operatorname{GL}_5$-同変である。
問題は、その Poisson ambient の古典 Casimir と量子中心三次との同定にある。
上の命題は、適切な量子補正を施した中心同定や、別の内在的な逆量子化の
不存在を主張するものではない。
また、Kuznetsov 成分から復元した可換三次は、ここでの素朴な $V(H_h)$ と
同じであるとは証明していない。
\end{remark}

\begin{remark}[形式量子化文献の適用範囲]
Yekutieli \cite[Corollary~11.2]{YekutieliTwisted} は
$H^1(\mathcal O_X)=H^2(\mathcal O_X)=0$ の下で、
完備局所 parameter algebra 上の Poisson 変形と結合的変形の
gauge 同値類を比較し、一次括弧と係数変更を保つ。
可換積自体が変わり得る比較であり、
固定した層上の bracket と star product の比較である
同論文 Corollary~11.3 はさらに $H^1(T_X)=H^1(\mathcal D_X)=0$ を要求する。
三次三様体では $h^1(T_X)=10$ なので、この固定積版を直接は用いない。
Pym \cite[Theorem~2.3 and the following paragraph]{Pym} の
解析芽の同型は Artin 近似で得られ、元の Kontsevich 形式同型と一致するとは限らない。
同論文 Proposition~2.9 と Theorem~2.10 は
規格化 Koszul CY 変形芽と unimodular Poisson 錐の局所比較を与える。
これらの定理の結論には、一般の非可換点での標準的な代数的逆量子化、
量子中心三次と古典 Casimir の同変な対応、
指定した残余射 $\rho$ との両立は含まれない。
\end{remark}

\section{Point scheme、Poisson零点と古典構造の復元}
\label{sec:point-scheme-poisson}

本節では、埋込み点スキームとそのシフト対応が保持する古典的な幾何を調べる。
二つの中心一次変数を持つJacobian模型では、点スキームは楕円曲線$E$と
長さ二十四の完全交叉$D$の非交和となる。
一般の境界三次に対して、このデータから五次元の三次Casimir空間と
ambient Poisson構造の射影類を復元し、シフト対応を加えると
二次包絡とPoisson構造の尺度も復元できる。
一方、有限Hesse族の$E\sqcup D$は、同じ射影空間のどの滑らかな三次上でも
Poisson零点スキームにはならない。
この非実現性は、点スキームからの復元と零点スキームとの同一視の差を示す。

証明はまず一次持上げの障害を計算し、古典零点の完全交叉表示を得た後、
有限係数での分解・復元・非実現性へ進む。
最後にシフトの一次速度をmodular residueと比較し、
これらの構成が構造対象を保つAbel同値に降下する範囲を明らかにする。

\paragraph{先行研究との関係.}
Vancliff \cite[\S6, Proposition~6.17 and Theorem~6.18]{VancliffQuantumSpaces}は、
多項式環の平坦な二次変形でpoint correspondenceが自己同型のgraphをなし、
非零パラメータで埋込みpoint schemeが一定であるという仮定の下で、
Poisson零点schemeへの包含を示している。同文献のExample~6.12では等号が失敗する。
ここで扱うHesse曲線は二次補正で動くため、その一定性の仮定を満たさない。

またLe~Bruyn--Smith--Van~den~Bergh
\cite[Theorem~4.2.2 and Proposition~4.3.7(1)]{LeBruynSmithVdBCentral}
は、一つの中心一次変数による一般Sklyanin拡大のpoint varietyが
楕円曲線と八点からなり、中心超平面の外ではshiftが恒等であることを示している。
以下は二つの中心一次変数と中心三次商を持つ模型について、
対応のscheme構造と長さ二十四の有限成分を明示するものである。
勾配空間の一意性にはJacobian Torelliを用い、modular residueの規約は
Gualtieri--Pymに従う。各々の適用箇所で、既知の結果と本模型への追加部分を述べる。


\subsection{点加群の一次持ち上げと Poisson 零点}

$k$ を標数零の体，$R=k[\epsilon]/(\epsilon^2)$ とする．
$V$ の基底を $x_0,\ldots,x_N$ とし，$B_R$ を
\[
 B_R=R\langle x_0,\ldots,x_N\rangle/
 \bigl(x_ix_j-x_jx_i-\epsilon\widetilde q_{ij}\bigr)_{i<j}
\]
で与えられる $R$-平坦な二次次数付き代数とする．
ここで $\widetilde q_{ij}$ の可換化は二次形式 $q_{ij}$ であり，
$q_{ji}=-q_{ij}$，$q_{ii}=0$ とする．
$B_R/\epsilon B_R=\operatorname{Sym}V$ と同一視する．
二次双ベクトル
\[
 \Pi=\sum_{i<j}q_{ij}\,\partial_{x_i}\wedge\partial_{x_j}
\]
は Poisson 双ベクトルであると仮定する．
この最後の仮定は形式的な量子化の一次切断であれば満たされるが，
双対数上の結合的変形だけから Jacobi 恒等式を結論してはならない．

$w_\epsilon\in Z(B_R)_3$ を次数三の中心元とし，
$F=w_\epsilon\bmod\epsilon$ と置く．
$X=V_+(F)\subset\mathbb P(V^*)$ は滑らかであるとする．
$A_R=B_R/(w_\epsilon)$ と置く．
中心性より $\Pi^\#dF=0$ であるから，$\Pi$ は $X$ 上の
Poisson 双ベクトル $\pi_X$ を誘導する．
以下の議論では $A_R$ の平坦性を別に用いないが，ここでは
$F$ が非零であることと $B_R$ の平坦性から $A_R$ も $R$-平坦である．

\begin{proposition}[点加群の一次持ち上げ]
\label{prop:first-order-point-poisson-zero}
$p\in X(k)$ とする．可換三次環 $A_0=\operatorname{Sym}V/(F)$ の
点加群 $M(p)$ が $A_R$ 上の $R$-平坦な点加群に持ち上がるための
必要十分条件は
\[
 \pi_X(p)=0
\]
である．ここで点加群とは次数零で生成され，各次数成分が
$R$ 上階数一の自由加群である次数付き右加群を意味する．

さらに任意の $k$-スキーム $T$ と $p:T\to X$ について，
$p$ に対応する点加群族が $T$ 上 Zariski 局所的に
$T[\epsilon]$ 上へ持ち上がることと，$p$ が零スキーム
$Z(\pi_X)\hookrightarrow X$ を経由することは同値である．
従って、局所的に一次持ち上げを持つ点の部分関手は $Z(\pi_X)$ により表現される．
ここで表現するのは持ち上げ可能な点の関手であり、持ち上げの選択を含む関手ではない．
\end{proposition}

\begin{proof}
まず $k$-点の場合を示す．$p$ の非零ベクトル表示を同じ文字で
$p=(p_0,\ldots,p_N)$ と書く．持ち上げの各次数の基底 $m_n$ を
選べば作用は
\[
 m_nx_i=a_{n,i}m_{n+1},\qquad
 a_{n,i}=p_i+\epsilon b_{n,i}
 \quad(n\geq0)
\]
と書ける．基底の剰余を通常の点加群の基底に同定している．
二次関係の $\epsilon$ 係数は
\begin{equation}\label{eq:first-order-point-wedge}
 q_{ij}(p)=p_i(b_{n+1,j}-b_{n,j})
                -p_j(b_{n+1,i}-b_{n,i})
\end{equation}
となる．従って $c_n=b_{n+1}-b_n$ と置くと
$\Pi(p)=p\wedge c_n$ である．
射影空間の接空間は Euler 方向 $kp$ を除いて得られるため，
この等式は射影化された双ベクトルが $p$ で消えることを意味する．
$X\hookrightarrow\mathbb P(V^*)$ の接束の単射から
$\bigwedge^2T_X\hookrightarrow\bigwedge^2T_{\mathbb P(V^*)}|_X$ も
単射であるので，これは $\pi_X(p)=0$ と同値である．
以上で必要性が従う．

逆に $\pi_X(p)=0$ と仮定する．あるベクトル $c$ が存在して
\begin{equation}\label{eq:first-order-point-c}
 q_{ij}(p)=p_ic_j-p_jc_i
\end{equation}
となる．任意のベクトル $b$ に対して
\begin{equation}\label{eq:first-order-point-action}
 a_n=p+\epsilon(b+nc)
\end{equation}
と置けば，式\eqref{eq:first-order-point-wedge}により
全ての二次関係を満たす．残る問題は中心三次関係である．

完全対称化により
\[
 w_\epsilon=\operatorname{sym}(F)+\epsilon\widetilde w_1,
 \qquad G=\widetilde w_1\bmod\epsilon\in\operatorname{Sym}^3V
\]
と書く．$B_R$ の平坦性により $G$ はこの等式から一意に定まる．
式\eqref{eq:first-order-point-action}の作用を三つ続けて
適用し，各単項式の三つの位置について平均を取ると
\begin{equation}\label{eq:first-order-central-evaluation}
 m_nw_\epsilon=
 \epsilon\bigl(dF_p(b)+(n+1)dF_p(c)+G(p)\bigr)m_{n+3}.
\end{equation}
実際，$b$ の寄与は通常の微分 $dF_p(b)$ であり，
$c$ の係数として現れる位置 $n,n+1,n+2$ の平均は $n+1$ である．
零次の項 $F(p)$ は $p\in X$ により消える．

中心性 $[x_i,w_\epsilon]=0$ の一次項から
$\Pi^\#dF=0$ を得る．一方，Euler の公式は
$dF_p(p)=3F(p)=0$ を与える．
式\eqref{eq:first-order-point-c}をこの中心性に代入すれば
$dF_p(c)p=0$ となり，$p\neq0$ から
\[
 dF_p(c)=0
\]
を得る．$X$ は滑らかなので $dF_p\neq0$ である．従って
$dF_p(b)=-G(p)$ を満たす $b$ を選べる．
式\eqref{eq:first-order-central-evaluation}は全ての $n$ で消え，
$A_R$ の作用が定義される．$p_\ell\neq0$ となる $\ell$ を選ぶと
$a_{n,\ell}$ は全ての $n$ で $R$ の単元である．従ってこの加群は
$m_0$ で生成され，各次数階数一の自由加群であり，必要な持ち上げを与える．

族についても議論は局所的である．$p_\ell=1$ のチャートでは
$c_\ell=0$，$c_i=q_{\ell i}(p)$ と置ける．
射影双ベクトルの零スキームの方程式は
\begin{equation}\label{eq:projective-poisson-zero-chart}
 q_{ij}(p)-p_iq_{\ell j}(p)+p_jq_{\ell i}(p)=0
 \qquad(i,j\neq\ell)
\end{equation}
である．これらの等式は任意の係数環上で
式\eqref{eq:first-order-point-c}と同値である．
必要性には次数 $0,1,2$ の線束を局所的に自明化すれば十分であり，
同じ式から $p^*\pi_X=0$ を得る．
十分性では $X$ の滑らかさより偏微分たちは単位イデアルを生成するので，
$T$ を Zariski 局所化して $dF_p(b)=-G(p)$ を解ける．
その後の構成は同じである．この議論は非被約な $T$ にも有効であり，
得られるものは零スキームの被約台だけではない．
\end{proof}

\begin{remark}
この命題は一次の持ち上げの存在を特徴付ける．
$k[[\hbar]]$ 上への持ち上げ，有限の非零 $\hbar$ における点加群の存在，
または固定された非可換代数の point scheme と $Z(\pi_X)$ の同型は
結論していない．高次には追加の障害があり得る．
また，持ち上げ自体が一意であるとも主張していない．
\end{remark}

\begin{remark}[一般次数]
同じ議論は中心元の次数 $d$ に対しても成立する．
式\eqref{eq:first-order-central-evaluation}の $(n+1)$ を
$n+(d-1)/2$ に置き換えればよい．
中心性と $F(p)=0$ により $dF_p(c)=0$ となるので，
滑らかな中心超曲面に対する必要十分条件は同じである．
\end{remark}

\begin{remark}[滑らかさの必要性]
滑らかさを外すと，Poisson 零点であっても中心関係自体が障害となる．
例えば $B_R=R[x,y]$，$w_\epsilon=x^3+\epsilon y^3$ とし，
$p=[0:1]\in V_+(x^3)$ を取る．Poisson 双ベクトルは恒等的に零であるが，
任意の点加群の持ち上げで $x$ の作用係数は $\epsilon$ の倍数，
$y$ の作用係数は単元である．従って $w_\epsilon$ の作用は非零の
$\epsilon$ 倍となり，$A_R$ 上の点加群には持ち上がらない．
\end{remark}


\subsection{滑らかな平面三次を基底とするPoisson零点scheme}
\label{subsec:poisson-zero-geometry}

以下では、零点は接双ベクトルの切断としてのscheme論的零点を意味する。
滑らかな三次三様体上の非零Poisson構造が二次元線型系から生じることは
Mayanskiy \cite[\S2.2, Lemma~2 and Theorem~1]{Mayanskiy}による。
Pencilの基底と特異ファイバーの特異点が零点を与える一般的な描像は
Pym \cite[\S3.2.2]{PymSurvey}にある。
次の命題では、この描像を三次三様体についてscheme論的に具体化する。
以下の完全交叉表示は、後に有限係数の点スキームと二次イデアルを比較するために用いる。

\begin{proposition}[平面三次と長さ二十四の完全交叉]
\label{prop:poisson-zero-elliptic-polar}
$X=V_+(F)\subset\mathbb P^4_{x,y,z,u,v}$を滑らかな三次三様体とする。
同型$\wedge^2T_X\simeq\Omega_X^1(2)$のもとで
$v\,du-u\,dv$に対応する非零Poisson構造を$\pi$とし、
\[
 E=X\cap V_+(u,v),\qquad
 D=V_+(F,F_x,F_y,F_z)
\]
と置く。$E$が滑らかな平面三次ならば、$D$は$E$と交わらない
零次元完全交叉であり、
\[
 \operatorname{length}(D)=3\cdot2^3=24,
 \qquad Z(\pi)=E\sqcup D
\]
がschemeとして成立する。$D$の被約性はこの仮定だけからは従わない。
\end{proposition}

\begin{proof}
まず$E$の近傍を調べる。$\mathcal O_X(1)$を局所的に自明化し、
$u,v$を局所関数とみなす。$E\subset X$は滑らかな余次元二の
完全交叉なので、$du,dv$は$E$上で一次独立である。
従って各点の近傍で$\Omega_X^1$の基底の一部に延長できる。
この基底に関する$v\,du-u\,dv$の係数イデアルは$(u,v)$である。
ゆえに$E$の近傍で$Z(\pi)$は$E$そのものであり、
この成分に余分な非被約構造や埋込み成分はない。

次に$u\ne0$上で$u=1$とし、$t=v/u$、$f=F/u^3$と書く。
上の一形式は可逆関数倍を除いて$-dt$になる。
$f_t$が可逆な開では、関係$df=0$から
\[
 dt=-\frac{f_x\,dx+f_y\,dy+f_z\,dz}{f_t}
\]
を得るので、零点イデアルは$(f_x,f_y,f_z)$である。
一方、$f_x,f_y,f_z$のいずれかが可逆な開では$dt$は余接束の基底の
一部となり、両方の零点イデアルは単位イデアルである。
$X$の滑らかさにより、これらの開は$u\ne0$を覆う。
$v\ne0$上でも同様である。従って$X\setminus E$上の零点schemeは$D$と一致する。
また$F_x,F_y,F_z$の$u=v=0$への制限は、
滑らかな平面三次$F(x,y,z,0,0)$の三つの偏微分であるから、$D\cap E=\varnothing$である。

$D$が有限であることも確認する。Gauss写像
\[
 \gamma:X\longrightarrow(\mathbb P^4)^*,\qquad
 p\longmapsto[F_x(p):F_y(p):F_z(p):F_u(p):F_v(p)]
\]
は$\gamma^*\mathcal O(1)=\mathcal O_X(2)$を満たす。
この線束はampleなので、$\gamma$のファイバーは正次元の完備曲線を含まない。
$\gamma$はproperであるから、その像への有限射である。
Pencilの直線を$\ell=\mathbb P\langle u,v\rangle\subset(\mathbb P^4)^*$と置くと、
$D=\gamma^{-1}(\ell)$である。Pencilの一般元は$X\setminus E$上でBertiniにより
滑らかであり、$E$上では$du,dv$の一次独立性からすべての元が滑らかである。
従って$\ell$は双対多様体$X^\vee=\gamma(X)$に含まれず、
$\ell\cap X^\vee$は有限である。ゆえに$D$も有限である。

多項式環$S=\mathbb C[x,y,z,u,v]$の同次イデアル
$(F,F_x,F_y,F_z)$は高さ四であり、四つの生成元は正則列をなす。
従って$D$は型$(3,2,2,2)$の完全交叉で、長さは二十四である。
その同次座標環は次元一のCohen--Macaulay環なので、このイデアルは飽和している。
先に得た局所計算と$D\cap E=\varnothing$を合わせると所望のscheme論的分解を得る。
\end{proof}

\begin{proposition}[被約な二十四点を持つ非空開]
\label{prop:poisson-zero-lefschetz-open}
上の組$(X,\langle u,v\rangle)$のパラメータ空間には、$D$が二十四個の
相異なる被約点からなり、pencilの各特異ファイバーが一個の通常二重点を持つ
非空開が存在する。
\end{proposition}

\begin{proof}
被約性、臨界点の非退化性、臨界値の相異性はいずれも開条件である。
非空性を明示的に確認するため
\[
 F=\frac{x^3+y^3+z^3}{3}-u^2(x+4y+16z)+u^3+v^3
\]
を取る。基底曲線は滑らかなFermat平面三次であり、
\[
 F_x=x^2-u^2,\qquad F_y=y^2-4u^2,\qquad F_z=z^2-16u^2.
\]
$D$上では$u\ne0$なので、$u=1$とすると
\[
 (x,y,z)=(\varepsilon_1,2\varepsilon_2,4\varepsilon_3),\qquad
 v^3=\frac23(\varepsilon_1+8\varepsilon_2+64\varepsilon_3)-1,
 \quad \varepsilon_i\in\{1,-1\}.
\]
八つの符号和は$\pm55,\pm57,\pm71,\pm73$である。
従って右辺は相異なる非零数で、二十四個の点と二十四個の臨界値$v$を得る。
方程式$(F_x,F_y,F_z,F)$の$(x,y,z,v)$に関するJacobian行列式は
$24xyzv^2\ne0$なので、すべての点は被約である。
また$F_v=3v^2$を用いて$X$上で$v$を局所関数として解くと、
臨界点でのそのHessianは
$-(3v^2)^{-1}\operatorname{diag}(2x,2y,2z)$であり、非退化である。

$X$自体も滑らかである。特異点があれば$F_v=0$から$v=0$となる。
$u=0$ならすべての座標が零になる。$u\ne0$なら$u=1$に正規化でき、
上の三つの二次方程式と$F_u=0$から
$\varepsilon_1+8\varepsilon_2+64\varepsilon_3=3/2$を得るが、これは不可能である。
以上で必要な非空例を得た。
\end{proof}

\begin{proposition}[零点曲線による射影Poisson構造の復元]
\label{prop:poisson-zero-projective-recovery}
$X$を固定し、命題\ref{prop:poisson-zero-elliptic-polar}の仮定を満たす
Poisson構造を考える。$Z(\pi)$は射影類$[\pi]$を一意に定める。
より具体的には、その唯一の一次元成分$E$から
\[
 U=H^0(X,\mathcal I_{E/X}(1))=\langle u,v\rangle
\]
を復元し、Pl\"ucker直線$\wedge^2U$から$[\pi]$を得る。
一方、$c\in\mathbb C^\times$に対して$Z(c\pi)=Z(\pi)$なので、
零点schemeだけでは非零スカラーの尺度は復元できない。
\end{proposition}

\begin{proof}
滑らかな平面三次$E$はその基底平面を張る。
従って$E$上で消える線型形式はちょうど$u,v$の張る二次元空間である。
命題\ref{prop:relative-poisson-cubic-cone}の同定により、
この二次元空間の外積が射影Poisson構造を定める。
最後の主張は切断を可逆な定数倍にしても係数イデアルが変わらないことによる。
\end{proof}


\subsection{二つの中心一次元を持つ代数の点スキーム}
\label{subsec:point-central-two}

以下では標数零の代数閉体上で、点加群を次数零から生成され各次数の次元が一である
右次数付き加群とする。族では各次数を可逆層とし、次数零を基底の構造層に同定する。
その点列を $(p_0,p_1,\ldots)$ と書く。

\begin{proposition}[楕円曲線と定常点の分離]
\label{prop:central-two-point-scheme}
$V=\langle x,y,z,u,v\rangle$、$h\ne0$ とし、
$F\in k[x,y,z,u,v]_3$、$f_3=F|_{u=v=0}$ とする。
$f_3$ は滑らかな平面三次を定めると仮定する。
$u,v$ を中心とし、
\[
 [y,z]=h\operatorname{sym}(F_x),\qquad
 [z,x]=h\operatorname{sym}(F_y),\qquad
 [x,y]=h\operatorname{sym}(F_z)
\]
で二次代数 $B$ を定める。
$S=B/(u,v)$ の二次関係の対応が、滑らかな平面三次 $E$ 上の
固定点を持たない自己同型 $\sigma:E\to E$ のグラフであると仮定する。
また中心三次 $w\in B_3$ の $S$ における像は、$E$ の全点加群を
消す三次の元とする。$A=B/(w)$ と置く。

$w$のtensor代数における斉次代表を一つ選び、その可換化を
$w_{\mathrm{ab}}\in\operatorname{Sym}^3V$と書く。
この多項式自体は代表に依存し得るが、二つの代表の差の可換化は
$(F_x,F_y,F_z)$に属する。従って次のスキームは代表に依存しない：
\[
 C=V(F_x,F_y,F_z),\qquad
 D_w=V(F_x,F_y,F_z,w_{\mathrm{ab}})\subset\mathbb P^4.
\]
このとき点加群の関手は
\[
 \operatorname{Pt}(B)=E\amalg C,\qquad
 \operatorname{Pt}(A)=E\amalg D_w
\]
で表される。次数シフトは $E$ 上で $\sigma$、残りの成分上で恒等である。
$C$ は平面 $V(u,v)$ と交わらない次数八の $(2,2,2)$ 完全交叉曲線である。
$w_{\mathrm{ab}}$が$C$のどの既約成分上でも恒等的に零でない場合、$D_w$は長さ二十四の
$(2,2,2,3)$ 完全交叉零次元スキームとなる。
\end{proposition}

\begin{proof}
二次関係の双線形化の零点スキームを
$\Gamma_B\subset\mathbb P^4\times\mathbb P^4$ とし、二つの射影座標を $p,q$ と書く。
中心性は、全ての一次座標 $z_i$ について
\[
 u(p)z_i(q)-z_i(p)u(q)=0,\qquad
 v(p)z_i(q)-z_i(p)v(q)=0
\]
を与える。
$u(p)=1$のchartで第二因子の座標$q_j=1$を局所的に選ぶと、
中心関係から$1=p_j u(q)$となるので$u(q)$は単元である。
これは任意の係数環上で成り立つ。
両方を一に規格化すると $p=q$ となる。
$v(p)\ne0$ の chart も同じである。
この対角上で三つの残りの関係は $hF_x=hF_y=hF_z=0$ となるので、
この部分は $C$ 上の恒等写像のグラフである。

次に $p$ と $q$ が異なる chart を取る。
ある二座標の小行列式
$z_i(p)z_j(q)-z_j(p)z_i(q)$ が可逆になる。
この二座標を使う中心性の二本の式から、スキームとして
\[
 u(p)=u(q)=v(p)=v(q)=0
\]
が従う。残りは $S$ の二次対応そのものであり、仮定から
$\Gamma_\sigma\subset E\times E$ となる。
$\sigma$ に固定点がないため、この小行列式chartは $\Gamma_\sigma$ 全体を覆う。

これらのchartが $\Gamma_B$ を覆うことも確認する。
幾何点で $u(p)=v(p)=0$ なら中心性から $u(q)=v(q)=0$ であり、
従って $(p,q)\in\Gamma_\sigma$ である。
一方、$C$ は $V(u,v)$ と交わらない。
もし交点があれば、そこでは $(f_3)_x=(f_3)_y=(f_3)_z=0$ となり、
Eulerの公式と $f_3$ の滑らかさに反する。
従って上の二つのグラフは互いに交わらず、chart上の計算により
非被約構造も含めて
\[
 \Gamma_B=\Gamma_\sigma\amalg\Delta_C
\]
となる。対応が自己同型のグラフなので、二点以上の二次関係を満たす点列、および無限点列は
初期点から一意に定まり、$\operatorname{Pt}(B)$ の表示を得る。

中心三次を加えると、全ての連続三点で $w$ の三重線型化が零であることを要求する。
$E$上では仮定により全ての幾何点で零であり、$E$が滑らかで被約なので切断として恒等的に零である。
定常点列では、どの三点で評価しても $w_{\mathrm{ab}}(p)$ となる。
よって $\operatorname{Pt}(A)=E\amalg D_w$ であり、族についても成立する。

三つの二次式で切られる $C$ の各成分の次元は一以上である。
二次元以上の成分は射影次元定理から平面 $V(u,v)$ と交わるため、存在しない。
従って三つの二次式のイデアルの高さは三で、正則列となる。
次数は $2^3=8$ である。この完全交叉はCohen--Macaulayで埋込み素点を持たないため、
どの既約成分も含まない三次切断は非零因子となる。従って正則列は長さ四となり、
$D_w$ の長さは $2^3\cdot3=24$ である。
\end{proof}

\begin{remark}[一般点での二十四個の相異なる点]
全混合係数を動かす族では、$C$ が滑らかで $D_w$ が被約となる条件は非空開である。
非空性は
\[
 F=\frac{x^3+y^3+z^3}{3}
   -q_1(u,v)x-q_2(u,v)y-q_3(u,v)z+g(u,v)
\]
で、三つの二元二次式 $q_i$ の六つの零点を互いに異なるものに取れば分かる。
$C$ は $x^2=q_1,\ y^2=q_2,\ z^2=q_3$ であり、
$\mathbb P^1_{u:v}$ 上の三つの二重被覆のファイバー積として滑らかである。
例えば $q_1=uv,\ q_2=u^2-v^2,\ q_3=u^2-4v^2$ を使える。
中心三次の $g(u,v)$ を動かす線型系には基点がなく、
$\mathbb P^1$ への射は有限なので、標数零の一般切断は二十四個の相異なる点となる。
この説明は、特別な分離型の点スキームも被約だと主張するものではない。
\end{remark}


\subsection{点対応からの二次包絡とPoisson尺度の復元}

\begin{proposition}[点対応の双線型イデアルによる二次包絡の復元]
\label{prop:point-correspondence-quadratic-recovery}
命題\ref{prop:central-two-point-scheme}の仮定に加え、$D_w$が
型$(2,2,2,3)$の零次元完全交叉であるとする。
$R_B\subset V\otimes V$を$B$の二次関係空間とし、
\[
 \Gamma_A=\Gamma_\sigma\amalg\Delta_{D_w}
 \subset\mathbb P(V^*)\times\mathbb P(V^*)
\]
を全埋込みpoint scheme上のshiftのgraphとする。
これは三点までの切断point schemeそのものとは区別する。このとき
\[
 \boxed{H^0(\mathcal I_{\Gamma_A}(1,1))=R_B.}
\]
従って埋込みpoint schemeとshiftの組は二次包絡$B$を復元する。
さらに、二次関係を
$[x_i,x_j]=\operatorname{sym}(q_{ij})$と一意に書けるgraph開では、
\[
 \Pi_{\mathrm{fin}}=\sum_{i<j}q_{ij}\,
       \partial_{x_i}\wedge\partial_{x_j}
\]
をスカラー倍の曖昧さなしに復元する。
本Jacobian族では$\Pi_{\mathrm{fin}}=h\Pi_F$でありPoisson構造である。
この構成に有限自己同型の対数を選ぶ操作は含まれない。
\end{proposition}

\begin{proof}
標数零なので$V\otimes V$を対称部分と反対称部分に分ける。
対角$D_w$上の消滅は、対称部分が
$Q=I(D_w)_2=\langle F_x,F_y,F_z\rangle$に属することと同値である。
従って候補空間は
$\operatorname{Alt}^2V\oplus\operatorname{sym}(Q)$で、次元は$10+3=13$である。
$u,v$を含む反対称tensorの空間は七次元であり、
$\Gamma_\sigma$上でも恒等的に零である。
これで割って境界平面へ制限すると、候補は
\[
 \operatorname{Alt}^2(V/U)\oplus
 \operatorname{sym}\langle(f_3)_x,(f_3)_y,(f_3)_z\rangle
 \subset (V/U)\otimes(V/U)
\]
という六次元空間に単射する。
偏微分の独立性は滑らかな平面三次がconeでないことから従う。

一方、$E$上で$L=\mathcal O_E(1)$、$M=\sigma^*L$と置くと、
双線型形式の制限は乗法写像
\[
 H^0(E,L)\otimes H^0(E,M)\longrightarrow H^0(E,L\otimes M)
\]
である。この写像は全射である。
$L\not\simeq M$なら、$L$の基点を持たない二次元の切断空間を取り、
base-point-free pencil trickを適用する。
次数零の非自明線束$M\otimes L^{-1}$の$H^1$は零なので全射となる。
$L\simeq M$なら、平面三次の二次イデアルが零であることから
$\operatorname{Sym}^2H^0(E,L)\simeq H^0(E,L^2)$となり、やはり全射である。
右辺の次元は六なので、$\Gamma_\sigma$の全双線型イデアルの次元は三である。

上の六次元候補には$S=B/(u,v)$の三つの独立な二次関係が含まれる。
従ってその核は正確に三次元である。
元の候補に戻した核は$7+3=10$次元で、既に十次元の$R_B$を含むから等しい。
ここで$R_B$の独立性は、反対称部分への射影が可逆であることからも分かる。
これで二次包絡の復元を得た。

graph条件の下では、各反対称tensor
$x_i\otimes x_j-x_j\otimes x_i$を持つ関係は一意である。
その対称部分を$-\operatorname{sym}(q_{ij})$と書くことにより
$q_{ij}$が一意に定まり、表示した双ベクトルを得る。
本族では定義から$q_{yz}=hF_x$と巡回式であり、$u,v$はCasimirだから、
これは実際に$h\Pi_F$である。
\end{proof}

\begin{remark}[先行結果と尺度の意味]
一中心変数の楕円型拡大で、点対応の双線型イデアルから
二次関係を復元する結果は\cite[Proposition~4.2.5]{LeBruynSmithVdBCentral}にある。
上の命題では中心三次商を取った後にも、有限成分の二次イデアルが
変わらないことを使って二中心変数の包絡を復元する。
点scheme単独からの$[\Pi_F]$復元と異なり、shiftを付けた点対応は
有限係数の$h\Pi_F$の尺度も保持する。
特定の量子化経路で半古典構造$\Pi_F$と比較するには、
指定したパラメータ$h$で割る必要がある。
この二つの意味の尺度を混同しない。
\end{remark}

\subsection{埋込み点スキームからの古典Casimir空間の復元}

\label{subsec:point-gradient-reconstruction}

有限成分の二次イデアルから、古典的な偏微分の空間を抽出する。
点スキームの抽象同型型だけではなく、その $\mathbf P(V^*)$ への埋込みを用いる。
ここではシフト対応を忘れた後に残るデータを記述する。

\begin{proposition}[古典Casimir三次空間の復元]
\label{prop:point-gradient-reconstruction}
上の分解 $\operatorname{Pt}(A)=E\amalg D_w$ が成り立ち、
$D_w$ が次数 $(2,2,2,3)$ の完全交叉であるとする。
$U=\langle u,v\rangle\subset V$ とし、$F$ の $u=v=0$ への制限が
\[
 f_3=\frac{x^3+y^3+z^3}{3}+\lambda xyz,
 \qquad \lambda(\lambda^3-8)\ne0
\]
であると仮定する。埋込み点スキームから
\[
 U=I(E)_1,\qquad W=\operatorname{ann}(U)\subset V^*,
 \qquad Q=I(D_w)_2
\]
を復元できる。さらに
\[
 R_Q:=\{G\in\operatorname{Sym}^3V\mid
              \partial_aG\in Q\text{ for every }a\in W\}
       =kF+\operatorname{Sym}^3U.
\]
特に $R_Q$ は五次元である。これによって、$u,v$ をCasimirとする
古典的な斉次二次Jacobian Poisson bivectorの射影類も復元できる。
\end{proposition}
\begin{proof}
$E$ は平面 $V(u,v)$ を張るので $I(E)_1=U$ である。
$D_w$ の斉次イデアルは完全交叉の飽和イデアルであり、
三次の生成元は次数二の部分に寄与しないので
\[
 Q=\langle F_x,F_y,F_z\rangle.
\]
これら三つの二次形式は、$u=v=0$ への制限の
$x^2,y^2,z^2$ の各係数から一次独立である。

$G\in R_Q$ に対し、一意な定数行列 $M=(m_{ij})\in\operatorname{Mat}_3(k)$ があり
\[
 (G_x,G_y,G_z)^{\mathsf t}=M(F_x,F_y,F_z)^{\mathsf t}
\]
となる。偏微分の可換性により $M\operatorname{Hess}_{x,y,z}(F)$ は対称である。
$u=v=0$ に制限し、$x,y,z$ の係数を比較すれば
\[
 H=\operatorname{Hess}(f_3)=
 \begin{pmatrix}
  2x&\lambda z&\lambda y\\
  \lambda z&2y&\lambda x\\
  \lambda y&\lambda x&2z
 \end{pmatrix}
\]
について $MH$ が対称でなければならない。
対角成分の差には
\[
 \lambda(m_{11}-m_{22})=
 \lambda(m_{11}-m_{33})=
 \lambda(m_{22}-m_{33})=0
\]
が課される。非対角成分には二組の循環関係
\[
 \lambda m_{13}=2m_{21},\quad
 \lambda m_{21}=2m_{32},\quad
 \lambda m_{32}=2m_{13},
\]
\[
 \lambda m_{23}=2m_{12},\quad
 \lambda m_{12}=2m_{31},\quad
 \lambda m_{31}=2m_{23}
\]
が課される。従って $\lambda^3\ne8$ なら全非対角成分は零となり、
$\lambda\ne0$ なら全対角成分が等しい。ゆえに $M=cI_3$ である。
以上より $\partial_x(G-cF)=\partial_y(G-cF)=\partial_z(G-cF)=0$ となり、
標数零より $G-cF\in\operatorname{Sym}^3U$ を得る。逆包含は直接明らかである。

$U$ の基底 $u,v$ と定数体積形式 $0\ne\Omega\in\bigwedge^5V$ を選び、
\[
 \iota_{\Pi_F}\Omega=dF\wedge du\wedge dv
\]
で斉次二次bivectorを定める。
$F$ を $cF+g(u,v)$ に替えると $\Pi_F$ は $c\Pi_F$ になる。
$U$ の基底と体積形式の変更も非零定数倍を与えるだけであるから、
$[\Pi_F]$ は復元された $U,R_Q$ だけに依存する。
\end{proof}

\begin{remark}[復元の範囲]
この命題は有限量子係数での点スキームから古典データを復元する主張であり、
点スキーム自体をPoisson零スキームと同一視する主張ではない。
$R_Q/\operatorname{Sym}^3U$ は一次元であるが、そこから $R_Q$ の特定の三次形式を
選ぶことはしていない。したがって特定の滑らかなcubic threefold、
中心三次関係の古典的な代表、Poisson bivectorの尺度までは復元していない。
この条件は本Jacobian族の中では非空開条件であり、例えば $\lambda=1$ で成立する。
$\lambda=0$ または $\lambda^3=8$ では上の九元連立一次方程式の階数は六に落ちる。
ただし、その場合でも $F$ のmixed項が追加条件を与え、
$R_Q=kF+\operatorname{Sym}^3U$ が成立する可能性はある。
\end{remark}


\begin{corollary}[三次Casimir全体の内在的な復元]
\label{cor:point-full-casimir-space}
命題\ref{prop:point-gradient-reconstruction}の条件の下で、復元した斉次二次
Jacobian Poisson括弧の次数三の中心は正確に
\[
 Z_{\mathrm{Pois}}(\operatorname{Sym}V)_3
 =kF\oplus\operatorname{Sym}^3U=R_Q
\]
である。ここで中心とはaffine括弧に関するCasimirを意味し、
射影的なPoisson不変超曲面一般を意味しない。
\end{corollary}
\begin{proof}
三つの偏微分$F_x,F_y,F_z$は正則列なので、共通の非定数因子を持たない。
三次同次式$G$がCasimirなら
$F_iG_j-F_jG_i=0$、$i,j\in\{x,y,z\}$である。
従って分数体上で$G_i=rF_i$となる$r$がある。
$r=a/b$を既約表示すると、$bG_i=aF_i$と多項式環のUFD性から
$b$は三つの$F_i$をすべて割る。よって$b$は単元で、$r$は多項式である。
両辺の斉次次数が二だから$r$は定数であり、
$G-cF\in k[u,v]_3$となる。逆包含は明らかである。
このCasimirの記述自体には$\lambda(\lambda^3-8)\ne0$を使わず、
その仮定は$R_Q$との同定にだけ必要である。
\end{proof}

\begin{remark}[Jacobian Torelliとの関係]
可積分な勾配空間からの一意性は既知のJacobian Torelliの特殊例である。
Ueda--Yoshinaga \cite[\S3, Lemma~3]{UedaYoshinagaJacobian}によれば、
異なる二つの滑らかな超曲面が同じJacobi idealを持てば、
その定義式はSebastiani--Thom型である。
従って上の復元命題は、$f_3$が滑らかで非Sebastiani--Thom型という
座標に依存しない仮定の下でも成立する。
実際、$\nabla g=M\nabla f_3$なら、一般の$t$で$M+tI$は可逆であり、
$g+tf_3$と$f_3$のJacobi idealは等しい。
この共通の勾配空間は基点を持たないから$g+tf_3$も滑らかである。
同補題により$g+tf_3$は$f_3$の定数倍となり、$M$はスカラーである。
平面三次の場合、除外される型は$j=0$に対応する。
Hwang \cite[Definition~2.1 and Theorems~1.3, 1.5--1.6]{HwangSymmetrizer}
のsymmetrizerは、このHessianの対称条件を一般的に扱う。
本稿の追加部分は、埋込みpoint schemeの有限成分から勾配空間を取り出し、
二つの中心変数を保ってPoisson直線へ戻すことであり、
Jacobian Torelli自体の新規性は主張しない。
\end{remark}

\subsection{Hesse曲線の有限補正と零点データの実現判定}
\begin{proposition}[境界点曲線のHessian補正]
\label{prop:point-hesse-hessian-correction}
$f_3=(x^3+y^3+z^3)/3+\lambda xyz$の対称化Jacobian関係について、
$h\ne0$とする。二次対応の行列式を定める三次式は、非零定数倍を除いて
\[
 e_h=A(x^3+y^3+z^3)+Bxyz,
 \quad A=1-\lambda^2h^2/4,
 \quad B=3\lambda+(1+\lambda^3/4)h^2
\]
であり、
\[
 e_h=3\left(f_3+\frac{h^2}{24}\det\operatorname{Hess}(f_3)\right).
\]
この曲線が滑らかで二次対応が自己同型のgraphとなる開では、
境界point schemeは$E_h=V_+(u,v,e_h)$である。
\end{proposition}
\begin{proof}
$p=1-\lambda h/2$、$q=-1-\lambda h/2$、$r=-h$と置く。
関係は$pyz+qzy+rx^2$とその巡回置換である。
右点加群の対応行列は
\[
 M_h=\begin{pmatrix}rx&qz&py\\pz&ry&qx\\qy&px&rz\end{pmatrix},
 \qquad \det M_h=(p^3+q^3+r^3)xyz-pqr(x^3+y^3+z^3).
\]
代入すると$\det M_h=-he_h$となる。一方、
\[
 \det\operatorname{Hess}(f_3)
 =(8+2\lambda^3)xyz-2\lambda^2(x^3+y^3+z^3)
\]
なので二つ目の式も従う。
滑らかな行列式曲線では行列の階数は二であり、
第二射影についても同じ計算をすると対応が自己同型のgraphとなる。
\end{proof}

この補正は有限係数の多項式恒等式であり、一次の近似だけの主張ではない。
古典的な零曲線$V_+(f_3)$とのずれは二次から始まる。
量子中心の単純な対称化逆像の補正とは係数も異なる。




\begin{proposition}[零点schemeの線型実現判定]
\label{prop:poisson-zero-linear-realization}
$P=V_+(u,v)\subset\mathbb P^4$上の滑らかな平面三次
$E=V_P(g)$と、$E$と交わらない長さ二十四の零次元scheme
$D\subset\mathbb P^4$を与える。
$E\sqcup D$がある滑らかな三次三様体上のPoisson構造の零点schemeとなるための
必要十分条件は、次の条件を満たす三次式$F$と$c\in\mathbb C^\times$が存在することである：
\[
 \begin{gathered}
 F\in H^0(\mathbb P^4,\mathcal I_D(3)),\qquad
 F|_P=cg,\\
 F_x,F_y,F_z\in H^0(\mathbb P^4,\mathcal I_D(2)),\qquad
 V_+(F)\text{ が滑らか}.
 \end{gathered}
\]
最初の三条件は$(F,c)$に関して線型であり、$c\ne0$と滑らかさは開条件である。
\end{proposition}

\begin{proof}
十分性を先に示す。$X=V_+(F)$には$\langle u,v\rangle$に対応する非零Poisson構造がある。
命題\ref{prop:poisson-zero-elliptic-polar}により、その零点schemeは
$E\sqcup D'$、$D'=V_+(F,F_x,F_y,F_z)$であり、$D'$の長さは二十四である。
与えられた線型条件は閉埋込み$D\subset D'$を与える。
両者の長さが等しいので$D=D'$となる。

必要性では、零Poisson構造なら零点schemeが三次元となるので除外する。
非零構造にはMayanskiy \cite[Theorem~1]{Mayanskiy}の分類を適用する。
最初からpencilの基底平面が$P$であるとは仮定しない。
その基底平面を$P'$とする。$P'\cap X$は零点schemeに含まれる。
$P'\subset X$なら零点schemeが二次元成分を持つので不可能である。
従って$P'\cap X$は平面三次であり、$E\sqcup D$の唯一の一次元成分が
滑らかな既約三次$E$であることから、その台は$E$である。
ゆえに$P'=\langle E\rangle=P$であり、三次という次数から$P\cap X=E$がschemeとして成立する。
命題\ref{prop:poisson-zero-elliptic-polar}を適用すれば、残りの必要条件を得る。
\end{proof}

\begin{proposition}[Hesse型データのPoisson零点schemeとしての非実現性]
\label{prop:hesse-point-zero-nonrealization}
$\lambda\in\mathbb C$、$1+\lambda^3\ne0$、$\hbar\in\mathbb C^\times$とし、
\[
 A=1-\frac{\lambda^2\hbar^2}{4},\qquad
 B=3\lambda+\left(1+\frac{\lambda^3}{4}\right)\hbar^2,
 \qquad e_\hbar=A(x^3+y^3+z^3)+Bxyz
\]
と置く。$E_\hbar=V_+(u,v,e_\hbar)$は滑らかと仮定する。
三つの二次式$Q_x,Q_y,Q_z$は、$P=V_+(u,v)$への制限が
\[
 Q_x|_P=x^2+\lambda yz,\qquad
 Q_y|_P=y^2+\lambda zx,\qquad
 Q_z|_P=z^2+\lambda xy
\]
であるとする。三次式$W$に対して
$D_\hbar=V_+(Q_x,Q_y,Q_z,W)$が零次元完全交叉ならば、
同じ$\mathbb P^4$に埋め込まれたどの滑らかな三次三様体$X'$と
そのPoisson構造$\pi'$に対しても
\[
 Z(\pi')\ne E_\hbar\sqcup D_\hbar
\]
である。
\end{proposition}

\begin{proof}
$1+\lambda^3\ne0$により、$P$上の三つの二次式
$x^2+\lambda yz,y^2+\lambda zx,z^2+\lambda xy$は共通零点を持たない。
実際、一つの座標が零ならすべて零となり、三つとも非零なら三式の積から
$1+\lambda^3=0$となる。従って$D_\hbar$は$P$と交わらない。
完全交叉のイデアルは飽和しているから、
\[
 H^0(\mathbb P^4,\mathcal I_{D_\hbar}(2))
 =\langle Q_x,Q_y,Q_z\rangle.
\]

もし所望の$(X',\pi')$が存在したならば、
命題\ref{prop:poisson-zero-linear-realization}の必要性により
$X'=V_+(F')$の方程式は$F'|_P=ce_\hbar$、$c\ne0$を満たし、さらに
\[
 F'_x|_P=c(3Ax^2+Byz)
 \in\langle x^2+\lambda yz,y^2+\lambda zx,z^2+\lambda xy\rangle
\]
となる。平方単項式の係数を比較すると、この線型結合の係数は
$(3Ac,0,0)$でなければならない。$yz$係数も比較して$B=3\lambda A$を得る。
しかし
\[
 B-3\lambda A=\hbar^2(1+\lambda^3)\ne0
\]
であり、矛盾する。
\end{proof}


\subsection{Fermat族の有限量子係数と平坦極限}

\label{subsec:fermat-point-flat-limit}

\begin{proposition}[Fermat点スキームの明示式]
\label{prop:fermat-point-scheme}
$k=\mathbb C$ とし、$u,v$ を中心とする
\[
 B_h=k\langle x,y,z,u,v\rangle/
 ([x,y]-hz^2,[y,z]-hx^2,[z,x]-hy^2,[u,-],[v,-])
\]
を考える。
\[
 \begin{gathered}
 f=\frac{x^3+y^3+z^3}{3},\qquad
 T=xyz+xzy+yxz+yzx+zxy+zyx,\\
 w_h=f-\frac{h^2}{9}T+u^3+v^3,\qquad A_h=B_h/(w_h)
 \end{gathered}
\]
と置く。中心性は定理\ref{thm:pair-full-fermat-center}による。
$h\ne0$、$h^6+27\ne0$のとき
\[
 \operatorname{Pt}(A_h)=E_h\amalg D_h,
\]
\[
 E_h=V(u,v,x^3+y^3+z^3+h^2xyz),
\]
\[
 D_h=V\left(x^2,y^2,z^2,
           u^3+v^3-\frac{2h^2}{3}xyz\right).
\]
$E_h$ は滑らかな平面三次である。
$D_h$ は長さ二十四で、被約台は
$[0:0:0:1:\zeta]$、$\zeta^3=-1$ の三点であり、
各点の局所長は八である。
\end{proposition}

\begin{proof}
境界代数の点対応行列は
\[
 M(x,y,z)=
 \begin{pmatrix}
 -hx&-z&y\\
 z&-hy&-x\\
 -y&x&-hz
 \end{pmatrix},
 \qquad
 \det M=-h(x^3+y^3+z^3+h^2xyz).
\]
$h^6+27\ne0$ なら表示した三次は滑らかであり、
行列はその曲線上で常に階数二となる。
左右の射影について同じ議論をすれば、対応が自己同型のグラフであることが分かる。
固定点は $x^2=y^2=z^2=0$ を要求するので存在しない。

三次 $w_h$ が境界点加群を消すことは、既存の Sklyanin 中心商表示からも従うが、
次の恒等式で直接確認できる。
$p=(x,y,z)$ とし、行列の最初の二行の外積に対応して
\[
 q=(xz+hy^2,\ yz-hx^2,\ z^2+h^2xy),\qquad
 r=\sigma_{\rm pol}(q)
\]
と置く。ここで $\sigma_{\rm pol}(X,Y,Z)=(XZ+hY^2,YZ-hX^2,Z^2+h^2XY)$ である。
$e_h=x^3+y^3+z^3+h^2xyz$ と書けば
\[
\begin{aligned}
 &3\sum_i p_iq_ir_i-h^2\sum_{\{i,j,k\}=\{1,2,3\}}p_iq_jr_k\\
 &\qquad=
 3e_h\left(
 2h^4x^2y^2+h^3x^3z-h^3y^3z+h^2xyz^2+z^4
 \right).
\end{aligned}
\]
左辺は $9(w_h|_{u=v=0})^{\rm ml}(p,q,r)$ である。
非空chartで零となるので、滑らかな既約曲線全体でこの切断は零である。
前命題を適用し、$w_h$ の可換像を二次式 $x^2,y^2,z^2$ で簡約すると
表示した $D_h$ を得る。

$D_h$ は $u\ne0$ に含まれ、そのchartの座標環は
\[
 k[x,y,z,v]/(x^2,y^2,z^2,v^3+1-\tfrac{2h^2}{3}xyz).
\]
$x^iy^jz^kv^\ell$、$0\le i,j,k\le1,\ 0\le\ell\le2$ が基底であり、長さは二十四。
被約台は三点で、$v$ の各単根の近傍では $v$ を一意に消去できるので
局所長は八となる。
\end{proof}

\begin{corollary}[全点関手の特殊ファイバーと平坦閉包の区別]
\label{cor:fermat-point-flat-closure}
$h=0$ を含む十分小さいZariski開基底上で、上の式が定める
$E_h\amalg D_h$ は平坦な射影族であり、一般点の点スキームのスキーム論的閉包である。
その特殊ファイバーは
\[
 E_0=V(u,v,x^3+y^3+z^3),\qquad
 D_0=V(x^2,y^2,z^2,u^3+v^3)
\]
の非交和である。これは$X_0$上で
$\{y,z\}=x^2$、$\{z,x\}=y^2$、$\{x,y\}=z^2$、$u,v$がCasimirとなる
\emph{非零の半古典Poisson構造}$\pi_0$の零点schemeに等しい。
$\pi_0$は交換子を$h$で割ってから特殊化することで得る。
一方、$A_0$ 自体の点スキームは滑らかな可換三次
\[
 X_0=V\left((x^3+y^3+z^3)/3+u^3+v^3\right)
\]
全体である。
従ってこの平坦閉包を、全点加群関手の特殊ファイバーと同一視しない。
\end{corollary}

\begin{proof}
$E_h$ は平面内の三次超曲面の平坦族であり、$h=0$ で滑らかである。
$D_h$ の上の座標環は、同じ二十四単項式により $k[h]$ 上自由である。
両者は基底全体で交わらない。
従って非交和は平坦で、その構造層に $h$-torsion がないため、
一般ファイバーのスキーム論的閉包に等しい。
命題\ref{prop:poisson-zero-elliptic-polar}により、この特殊ファイバーは
$Z(\pi_0)$である。$h=0$の代数は通常の可換三次の斉次環なので、
全点加群の表示は$X_0$である。
\end{proof}


\begin{corollary}[Fermat point schemeの有限係数での非実現性]
\label{cor:fermat-point-zero-nonrealization}
$h\ne0$、$h^6+27\ne0$の上のFermat模型で、
埋込みscheme$E_h\amalg D_h$は、同じ$\mathbb P^4$の
どの滑らかな三次三様体上のPoisson零点schemeにもならない。
\end{corollary}
\begin{proof}
$D_h$は表示した型$(2,2,2,3)$の完全交叉なので、
命題\ref{prop:hesse-point-zero-nonrealization}を$\lambda=0$で適用できる。
\end{proof}

\begin{remark}[境界曲線の抽象同型型も変わる]
Artebani--Dolgachev \cite[\S2, equations~(6)--(8)]{AD}のWeierstrass係数に
$(u_0,u_1)=(1,h^2/6)$を代入すると、通常の規格化で
\[
 j(E_h)=\frac{h^6(216-h^6)^3}{(27+h^6)^3}.
\]
実際、同文献の$A=12u_1(u_0^3-u_1^3)$、
$B=2(u_0^6-20u_0^3u_1^3-8u_1^6)$に対して
$j=1728\cdot4A^3/(4A^3+27B^2)$を計算すれば上式となる。
従って一般の$E_h$は$j(E_0)=0$の古典境界と抽象的にも同型ではない。
ただし、これだけで任意の別三次上の零点としての実現を否定するのではなく、
その強い非実現性には前の二次イデアル比較を使う。
\end{remark}

\subsection{シフト自己同型の一次速度}

\label{subsec:point-shift-first-order}

\begin{proposition}[Fermat境界に残る一次の量子尺度]
\label{prop:fermat-shift-velocity}
Fermat族の $E_h$ は $h$ の一次まで定数族であり、
点加群のシフト自己同型の一次速度は、$z=1$ のchartで
\[
 \xi=y^2\frac{\partial}{\partial x}
      -x^2\frac{\partial}{\partial y}
\]
である。このベクトル場は $E_0$ 上の大域的な非消滅ベクトル場へ延びる。
$h$ を $ch$ に置き換えると速度は $c\xi$ となる。
\end{proposition}
\begin{proof}
上の点対応式を $z=1$ で規格化すると
\[
 \sigma_h(x,y)=
 \left(\frac{x+hy^2}{1+h^2xy},
       \frac{y-hx^2}{1+h^2xy}\right).
\]
従って $h=0$ における微分は表示した $\xi$ である。
$\xi(x^3+y^3+1)=0$ なので曲線に接する。
滑らかな平面三次のresidue一形式を
$\omega=dx/y^2=-dy/x^2$ と規格化すると $\iota_\xi\omega=1$ である。
従ってこのchart上のベクトル場は、大域的な非消滅一形式の双対として
曲線全体へ正則かつ非消滅に延びる。
尺度に関する主張はchain ruleによる。
\end{proof}

\begin{remark}[有限自己同型と無限小速度]
この命題は指定した量子化経路とパラメータ $h$ の下での無限小の主張である。
有限の自己同型から代数的・一意な対数を取れるとは主張しない。
また $E_h,D_h$ は $h$ の符号を変えても同じであるが、
二次関係を反転すると対応の二つの射影が交換されるので
$\sigma_{-h}=\sigma_h^{-1}$ である。
したがって裸の点スキームが忘れる情報の一部は、
点スキームとシフト自己同型の組に残っている。
\end{remark}


\subsection{Modular residueによる尺度の回収}


零点schemeが失う尺度を補う自然な微分データとして、canonical Poisson moduleの
modular residueがある。これは既知の構成であり、
Gualtieri--Pym \cite[Definition~4.16 and Theorem~7.2]{GualtieriPym}による。
本稿では同論文の規約
\[
 \operatorname{Res}^{0}_{\mathrm{mod}}(\pi)
 =-\operatorname{Tr}\bigl(D\pi|_{Z(\pi)}\bigr)
\]
を採用する。右辺は$\pi$の零点に沿うfirst jetの縮約であり、
零点scheme自身の接ベクトル場になる。

\begin{proposition}[平面三次とmodular residueによる復元]
\label{prop:poisson-modular-scale-recovery}
命題\ref{prop:poisson-zero-elliptic-polar}の仮定のもとで、
$\eta_\pi=\operatorname{Res}^{0}_{\mathrm{mod}}(\pi)|_E$は
$E$上の至る所非零な大域接ベクトル場である。
固定した$X$上では組$(E\subset X,\eta_\pi)$が$\pi$を一意に定める。
\end{proposition}

\begin{proof}
$E$の点の近くで$(u,v,z)$を局所座標とし、体積形式を選ぶと、
ある可逆関数$a$により
\[
 \pi=a(u,v,z)(u\partial_u+v\partial_v)\wedge\partial_z
\]
と書ける。従って$E$上で
\[
 D\pi=a|_E\bigl(du\otimes(\partial_u\wedge\partial_z)
                  +dv\otimes(\partial_v\wedge\partial_z)\bigr),
 \qquad \eta_\pi=-2a|_E\partial_z.
\]
この式は$\eta_\pi$が非零であることを示す。
縮約による定義は座標と体積形式に依存せず、局所計算は貼り合わされる。
局所標準形で$a=1$の場合は
\cite[Example~7.4]{GualtieriPym}の計算そのものである。
命題\ref{prop:poisson-zero-projective-recovery}により$E$は$[\pi]$を決め、
さらに$\eta_{c\pi}=c\eta_\pi$と$\eta_\pi\ne0$から$c$が決まる。
\end{proof}

\begin{proposition}[Jacobian模型の一般の速度比較]
\label{prop:jacobian-modular-shift-comparison}
$X=V_+(F)$と$E=X\cap V_+(u,v)$はともに滑らかとする。
半古典括弧を、homogeneous coordinate ring上で
\[
 \{y,z\}=F_x,\qquad \{z,x\}=F_y,\qquad \{x,y\}=F_z,
 \qquad \{u,-\}=\{v,-\}=0
\]
と規格化する。この規格化を持つ二次変形において、$E_0=E$を持つ相対平面曲線族
$E_\hbar$上のpoint correspondenceが$\sigma_\hbar$のgraphをなし、$\sigma_0=\mathrm{id}_E$で、
その一次変形が定義されると仮定する。このとき
\[
 \left.\frac{d\sigma_\hbar}{d\hbar}\right|_{\hbar=0}
 =\frac12\operatorname{Res}^0_{\mathrm{mod}}(\pi)|_E.
\]
ここで$\pi$は$X$上に誘導されたPoisson構造である。
\end{proposition}

\begin{proof}
$z=1$で$f=F(x,y,1,0,0)$と置く。
$p'=p+\hbar\delta$をpoint relationの一次項に代入する。
$yz'-zy'=\hbar F_x+O(\hbar^2)$と
$zx'-xz'=\hbar F_y+O(\hbar^2)$から
$\delta_y=-f_x$、$\delta_x=f_y$を得る。
従ってshiftの速度は$\xi_f=f_y\partial_x-f_x\partial_y$である。
この計算は二点の差の一次項なので、曲線の変形の局所自明化の選択に依存しない。

一方、degree-zero局所化の括弧を計算すると
\[
 \{x/z,u/z\}=uF_y/z^3,\qquad
 \{x/z,v/z\}=vF_y/z^3,\qquad
 \{u/z,v/z\}=0.
\]
従って$z=1$、$F_y\ne0$では$(x,u,v)$を$X$の局所座標として
$\pi=F_y(u\partial_x\wedge\partial_u+v\partial_x\wedge\partial_v)$である。
そのfirst jetを縮約すると
$\operatorname{Res}^0_{\mathrm{mod}}(\pi)|_E=2f_y\partial_x=2\xi_f$を得る。
他の座標chartでも同様である。
\end{proof}


\subsection{構造対象付きAbel同値に対する降下の範囲}


この部分では全相対 Hom 比較定理
\ref{thm:general-relative-pointed-hom}を用いる．
その定理が適用される中心二変数 Jacobian 開では，
\[
 \underline{\operatorname{Isom}}_{\mathrm{gr}}(A,A')/\mathbb G_m
 \simeq
 \underline{\operatorname{Eq}}_{\mathcal O}
       (\operatorname{QGr}A,\operatorname{QGr}A')
\]
である．次数付き代数の同型は各長さの切断点加群関手，
それらの切断写像，および次数移動による対応を自然に保つ．
次数スカラー $x_i\mapsto\lambda x_i$ は次数 $n$ の基底を
$\lambda^n$ 倍することで吸収されるので，これらの点加群関手上には
自明に作用する．従って点加群の対応は全ての
$\mathcal O$-保存 Abel 同値に対してこの開上で降下する．
この帰結のために大域的な $\mathcal O(1)$ の選択を追加する必要はない．
ただし一般の Abel 圏において内在的な条件だけで点対象を特徴付けた，
あるいは任意の五次元 AS 正則代数の点対応が有限段階で自己同型のグラフに
安定化すると主張しているのではない．


\begin{proposition}[復元の相対性とAbel同値への自然性]
\label{prop:point-recovery-relative-descent}
既存Jacobian族の中で、命題\ref{prop:point-gradient-reconstruction}の仮定、
特に$D_w$が相対零次元完全交叉となる条件と
$\lambda(\lambda^3-8)\ne0$を満たす開を考える。
埋込みpoint schemeからの$U,Q,R_Q$と$[\Pi]$の構成は代数的であり、
任意の底変更と両立する。構成は$\operatorname{GL}(V)$-同変であり、
定理\ref{thm:general-relative-pointed-hom}が適用される開では
全ての$\mathcal O$-保存Abel同値に自然である。
shiftを付けた復元についても
$H^0(\mathcal I_{\Gamma_A}(1,1))=R_B$は相対的に成立し、
$\Pi_{\mathrm{fin}}$は中心重み零の実際の双ベクトルとして降下する。
\end{proposition}
\begin{proof}
$E$は相対平面三次であり、その一次イデアルは階数二の部分束$U$である。
相対完全交叉$D$のKoszul解消と底変更から、
次数二イデアルは階数三の部分束$Q=\langle F_x,F_y,F_z\rangle$である。
$W=\operatorname{ann}(U)$とすると
\[
 R_Q=\ker\left(\operatorname{Sym}^3V\longrightarrow
 W^*\otimes(\operatorname{Sym}^2V/Q)\right).
\]
Hesseの証明に用いた二つの三項循環関係は、$2,\lambda,\lambda^3-8$が
単元である任意の係数環でそのまま有効である。
したがって$R_Q=\mathcal O F\oplus\operatorname{Sym}^3U$は階数五の部分束となり、
任意の底変更と両立する。
商直線$R_Q/\operatorname{Sym}^3U$からの双ベクトルの射影類も同様である。
全ての操作は線型形式の空間に関して座標自由なので同変であり、
前述の次数スカラーの自明性と全Hom比較によりAbel同値に対して降下する。
中心重みを持つ$U$や$R_Q$そのものが商上の通常のベクトル束になるとは
主張せず、対応する射影データとその同変構成を降下する。

shift付きの主張では、楕円曲線上の$L\otimes\sigma^*L$は次数六なので
$H^1$が零で、その切断束は階数六で任意の底変更と両立する。
上の乗法写像は全ての幾何ファイバーで全射なので相対的にも全射である。
対角$D$で消える階数十三の候補束は既知の関係部分束$R_B$を含み、
商は階数三である。この商から境界切断束への写像は全ファイバーで単射なので、
三次小行列式により局所分裂単射となる。従って非平坦底変更の後にも
核は正確に$R_B$である。
最後に$\Pi_{\mathrm{fin}}\in\operatorname{Sym}^2V\otimes\bigwedge^2V^*$の
$\operatorname{GL}(V)$の中心重みは$2-2=0$なので、
これは射影類だけでなく実際の双ベクトルとして$\operatorname{PGL}(V)$に降下する。
\end{proof}

ここでいう開条件は本Jacobian locus内の条件である。
任意の$(U,Q)$の全Grassmannian上で核の階数五が開になるという主張ではない。
非空性は、Hesse係数の非空開、完全交叉条件、および
命題\ref{prop:jacobian-affine-four-morphism}の既約Jacobian atlasの非空開を
共通部分に取ることで得る。
係数環\eqref{eq:general17-coefficient-ring}では$\lambda=b/(3a)$、$h=3a$なので、追加する非ST条件は
$b(b^3-216a^3)\ne0$である。境界の中心商表示は
式\eqref{eq:general17-elliptic-quotient}から得られ、固定点は古典境界の勾配の
同時零点を要求するので存在しない。
$C$が$V(u,v)$と交わらないためbinary cubicの線型系は$C$上で基点を持たず、
その一般元を中心三次へ加えれば真の切断となる。
点scheme単独からはambient Poisson直線と三次Casimir空間を復元し、
shiftも含めれば命題\ref{prop:point-correspondence-quadratic-recovery}により
尺度を持つ$\Pi_{\mathrm{fin}}$まで復元できる。
その相対版と全$\mathcal O$-保存同値への自然性は直前の命題で証明した。
量子中心直線に対応する特定の古典Casimir直線と、
Kuznetsov残余射$\rho$の出力三次との一致は依然として別の問題である。


\section{今後の問題}\label{sec:status}

\paragraph{量子中心と残余三次の比較.}
主定理\ref{thm:intro-points}は、埋込み点データから五次元の三次 Casimir 空間を復元する。
その中から中心商 $B/(w)$ に対応する直線を選び、
主定理\ref{thm:intro-residual}の $X_c$ を得るためには、量子中心の補正を記述する必要がある。
命題\ref{prop:poisson-naive-symmetrization-fails}と
命題\ref{prop:hesse-point-zero-nonrealization}は、補正が満たすべき制約を与える。
中心三次の具体式と相対 Casimir 空間を比較することが一つの入口である。
この選択が全 $\mathcal O$-保存同値と係数変更に自然となるかも問題の一部である。

\paragraph{七次元ファイバーの幾何.}
可換点の形式ファイバーは Pl\"ucker 錐であるが、
一般の代数的ファイバーの双有理型は定まっていない。
残余圏と五矢 Kronecker 圏の接合双加群を通じて、
残余射が失う七方向を理解したい。
同時に、二次包絡射 $\tau$ の四次元 torsor 上で $\rho$ がどう変化するかを調べる必要がある。
相対超越基底から作る非標準的な持上げは、この双有理分類を与えない。

\paragraph{全 Abel モジュライ内の位置.}
$\mathfrak C_{\mathrm J}$ の有限型代数性と完全忠実性は主定理\ref{thm:intro-atlas}で得た。
一般の非可換点における全 Abel 変形を有限関係窓と比較できれば、
全モジュライ内での開性・局所閉性・成分性を検討できる。
この比較では接空間に加えて小拡大の持上げと障害を扱う必要がある。
対角の有限性、分離性、$\mathcal O$ 自体の内在的な認識も独立の問題として残る。

\appendix
\clearpage
\part{補助構成と証明}
\section{正則正規超曲面対とAS正則$\mathbb Z$-代数の周期性}
\label{sec:moduli-ambient-periodicity}
\label{sec:periodic-pair-route}
本節では，通常の次数付き表示をあらかじめ選ばず，
正則正規超曲面とそのambientを同時に周期化する．
標準5次元AS正則$\mathbb Z$-代数の正規三次商については，
三次イデアルを保つ1周期同型が存在し，通常の次数付き正規三次商として表示できる．
その証明では，正規族による有限頂点orbit代数を構成し，
Reyes--Rogalskiの一般化AS正則性とtwisted Calabi--Yau性の同値，
Nakayama双加群と片側Extの比較，中心への作用を使う
\cite[Theorem~5.15, Lemmas~5.9 and~4.5]{RRGeneralizedAS}．
これらの既知結果は通常の有限頂点次数付き代数に適用する．
抽象的な非可換cubicすべての超曲面表示や，族全体での周期構造の降下を
この定理だけから結論することはしない．

\subsection{添字，次数，AS条件}

$\mathcal B$を体$k$上の局所有限な正の連結$\mathbb Z$-代数とする．
すなわち$\mathcal B_{ii}=k$，$\mathcal B_{ij}=0$（$j<i$）で，積は
\begin{equation}
 \mathcal B_{ij}\otimes\mathcal B_{j\ell}\longrightarrow\mathcal B_{i\ell}
\end{equation}
の順とする．右$\mathcal B$-加群$M$は$M_i$と
$M_i\otimes\mathcal B_{ij}\to M_j$からなる．
$P_i=e_i\mathcal B$，$S_i=P_i/(P_i)_{>i}$と書く．
左単純加群は${}_iS$と表す．

$\mathcal B(r)_{ij}=\mathcal B_{i+r,j+r}$と置き，
添字を保つ同型$\mathcal B\simeq\mathcal B(r)$があるとき$r$周期的という．
通常の連結次数付き代数$C$に付随する$\check C_{ij}=C_{j-i}$は1周期的である．
逆に1周期同型を選ぶと，それを用いて$\mathcal B_{0m}$に積を定め，
$\mathcal B\simeq\check C$と表示できる．この表示は周期同型の選択を含む．
正規族またはイデアルを持つ場合には，周期同型がそれも保存することを別に要求する．

本節で必要なAS仮定は次の明示的なものである．
\begin{enumerate}[label=\textup{(AS\arabic*)}]
\item 右および左$\mathcal B$-加群圏のglobal dimensionは$h<\infty$である．
\item 各$S_i$は有限生成射影加群による有限解消を持ち，ある固定整数$N>0$について
\begin{equation}\label{eq:pair-route-as-right}
 \Ext^q_{\mathcal B}(S_i,P_j)
 \simeq
 \begin{cases}
  k,&q=h,\ j=i+N,\\
  0,&\text{それ以外}
 \end{cases}
\end{equation}
である．左右を交換した条件も仮定する．
\end{enumerate}
標準5次元Koszul AS正則型では$h=N=5$であり，
最小右解消の末項は$P_{i+5}$である．
後述する正規超曲面対の周期性の証明にはKoszul性自体は不要で，uniformなAS添字$N$を使う．

通常の次数付き代数$D$の加群のシフトは$M(s)_n=M_{n+s}$とする．
従って$k(N)$は内部次数$-N$に集中する．
内部次数を記録するExtを
\begin{equation}
 \underline{\Ext}^{q}_{D}(M,D)_s
  =\Ext^q_{\operatorname{Gr}D}(M,D(s))
\end{equation}
で定義する．以下の$\operatorname{Gr}D$は全整数次数の右加群圏である．
通常の代数としての$D$の次数は非負である．
本節では下付き$D$または$\mathcal B$のExtを右加群圏で取る．
外部文献が左加群を標準規約とする場合の$\Ext_{D^{\rm op}}$は，
本節では$\Ext_D$と書く．ただし双加群Extの下付き$D^e$は
対角双加群を左$D^e$-加群とした通常の規約に従う．

次数$d>0$の左右正則正規族$0\ne g_i\in\mathcal B_{i,i+d}$を仮定する．
正規性は
\begin{equation}\label{eq:moduli-ambient-normal-family}
 \mathcal B_{ij}g_j=g_i\mathcal B_{i+d,j+d}
\end{equation}
である（正規族の定義については
Chan--Nyman\cite[Definition~8.6]{PChanNymanNormalFamily}も参照）．
ここでは中心性を仮定しない．左右正則性により一意な線型同型
\begin{equation}\label{eq:moduli-ambient-sigma}
 \sigma_{ij}:\mathcal B_{ij}\longrightarrow\mathcal B_{i+d,j+d},
 \qquad ag_j=g_i\sigma_{ij}(a)
\end{equation}
が定まる．$a\in\mathcal B_{ij}$，$b\in\mathcal B_{j\ell}$について
\begin{equation}
 (ab)g_\ell
  =a(g_j\sigma(b))
  =g_i\sigma(a)\sigma(b)
\end{equation}
となるので，$g_i$の正則性による消去から$\sigma(ab)=\sigma(a)\sigma(b)$を得る．
$a=g_i$を代入すれば$\sigma(g_i)=g_{i+d}$となる．
従って$\sigma$は$g$も保つ$d$周期同型である．
負の冪は逆同型を用いて定義する．


\subsection{正規族を仮定しないKoszul ambientの周期}
\begin{lemma}[Koszul--Frobenius双対性からの周期]
\label{lem:moduli-ambient-koszul-period}
$\mathcal B$は二次Koszul AS正則$\mathbb Z$-代数で，
左右のAS条件を満たし，各単純加群は有限生成射影加群による有限の最小線型解消を持つとする．
AS次元とGorenstein parameterがともに$N$なら，
正規族の存在を仮定せず$\mathcal B\simeq\mathcal B(N)$となる．
\end{lemma}
\begin{proof}
Koszul双対 $E=\mathcal B^!$ を考える。
有限線型解消と左右AS双対性から，$E$ はindex $N$ の
局所Frobenius $\mathbb Z$-代数となる
（Bondal--Polishchuk~\cite[Proposition~4.1, \S5]{PBondalPolishchukHelices}）。
具体的には，最高次数の同定
$\lambda_i:E_{i,i+N}\xrightarrow{\sim}k$ を選ぶと，
\begin{equation}
 E_{ij}\otimes E_{j,i+N}\longrightarrow E_{i,i+N}
 \xrightarrow{\lambda_i}k
\end{equation}
が各相補次数で完全対合である。
この対合から，任意の $a\in E_{ij}$ に対して
\begin{equation}
 \lambda_i(ab)=\lambda_j(b\nu_{ij}(a))
 \quad(b\in E_{j,i+N})
\end{equation}
となる線型同型 $\nu_{ij}:E_{ij}\to E_{i+N,j+N}$ が定まる。
$a\in E_{ij},c\in E_{j\ell},b\in E_{\ell,i+N}$ に対して
\begin{equation}
 \lambda_i((ac)b)
 =\lambda_j((cb)\nu(a))
 =\lambda_\ell(b\nu(a)\nu(c))
\end{equation}
であるから，対合の完全性により
$\nu(ac)=\nu(a)\nu(c)$ となる。
合計の長さが$N$を超える場合には積がともに零なので，同じ結論となる。
従って $\nu:E\xrightarrow{\sim}E(N)$ は $\mathbb Z$-代数同型である。
次数一部分で逆双対を取り，二次関係を保つことを用いると，
\begin{equation}\label{eq:moduli-ambient-rho-unmarked}
 \rho_0:\mathcal B\xrightarrow{\sim}\mathcal B(N)
\end{equation}
が得られる。
\end{proof}
この直接証明はambientのみの周期を与える．
指定した正規族を保つ$N$周期同型を得るには，以下のorbit構成を使う．

\subsection{中心元を持つ有限頂点orbit代数}

各整数$i$を一意に$i=\bar i+dq_i$，$0\leq\bar i<d$と書く．
通常の非負次数付き代数$D$を
\begin{equation}\label{eq:pair-route-orbit-space}
 D_m=\bigoplus_{r=0}^{d-1}\mathcal B_{r,r+m}\quad(m\geq0)
\end{equation}
として定める．成分$\mathcal B_{r,r+m}$は
$e_rD_me_{\overline{r+m}}$である．
$a\in\mathcal B_{r,r+m}$，$b\in\mathcal B_{s,s+n}$に対し，
$s\not\equiv r+m\pmod d$なら$a\cdot b=0$とし，
そうでなければ
\begin{equation}\label{eq:pair-route-orbit-product}
 a\cdot b=a\,\sigma^{(r+m-s)/d}(b)
   \in\mathcal B_{r,r+m+n}
\end{equation}
と置く．$\sigma$が代数同型なのでこの積は結合的である．
具体的には三つの合成可能な成分の積では，第二・第三因子を
同じ開始添字まで$\sigma$で移してから$\mathcal B$内で掛けるため，
二つの括弧付けは同じ三因子積になる．
$1_D=\sum_{r=0}^{d-1}e_r$，$D_0=k^d$であり，$D$は局所有限である．

\begin{lemma}[orbit内の中心元]\label{lem:pair-route-orbit-center}
\begin{equation}\label{eq:pair-route-central-z}
 z=\sum_{r=0}^{d-1}g_r\in D_d
\end{equation}
は中心正則元である．
\end{lemma}
\begin{proof}
$a\in\mathcal B_{r,r+m}\subset D_m$について
\eqref{eq:pair-route-orbit-product}を用いると
\begin{equation}
 a\cdot z=ag_{r+m},\qquad
 z\cdot a=g_r\sigma(a)
\end{equation}
となる．両者は\eqref{eq:moduli-ambient-sigma}により一致する．
各斉次成分・各開始頂点で$g_r$の左右乗法と$\sigma$は単射なので，
$z$の左右乗法も単射である．
\end{proof}

この$D$は，単に任意の既知の1周期構造で
$\mathcal B$を通常の連結次数付き代数に読み替えたものではない．
\emph{正規族自身から構成した}$\sigma$で畳み込むため，
一般の正規族が中心元$z$になる．ここが本証明の要点である．

\subsection{AS正則性の移行}

\begin{lemma}[graded module categoryの分解]
\label{lem:pair-route-sector-equivalence}
自然な完全圏同値
\begin{equation}\label{eq:pair-route-category-product}
 \operatorname{Gr}D\simeq
   \prod_{c=0}^{d-1}\operatorname{Mod}\mathcal B
\end{equation}
が存在する．特に$D$のgraded global dimensionは$h$である．
\end{lemma}
\begin{proof}
$M\in\operatorname{Gr}D$を$M_ne_r$に分解する．
次数$m$の矢印の作用は$(n,r)$を$(n+m,\overline{r+m})$へ送るので，
$r-n\bmod d$は保存される．従って
\begin{equation}
 M=\bigoplus_{c\in\mathbb Z/d\mathbb Z}M^{(c)},\qquad
 M^{(c)}_n=M_ne_{\overline{c+n}}
\end{equation}
は$D$-部分加群への分解である．異なるsector間の次数零準同型は零である．

$c$の代表を$0\leq c<d$に選ぶ．右$\mathcal B$-加群$L$から
$c$-sectorの$D$-加群$\Gamma_cL$を
\begin{equation}
 (\Gamma_cL)_n=L_{c+n}
 \quad\text{（右頂点$\overline{c+n}$に置く）}
\end{equation}
で定める．$a\in\mathcal B_{\overline{c+n},\overline{c+n}+m}\subset D_m$
の作用は，$q=(c+n-\overline{c+n})/d$として
$\sigma^q(a)\in\mathcal B_{c+n,c+n+m}$の作用で定義する．
\eqref{eq:pair-route-orbit-product}から結合律が従う．

逆に$c$-sectorの$M$に対して$L_j=M_{j-c}e_{\bar j}$とし，
$b\in\mathcal B_{i,j}$を
$\sigma^{-q_i}(b)\in e_{\bar i}D_{j-i}e_{\bar j}$として作用させる．
二つの構成は互いに逆であり，各成分での完全性を保つ．
有限個のsectorの直積はglobal dimensionの最大値を持つので最後の主張を得る．
\end{proof}

\begin{lemma}[通常のAS条件と頂点置換]
\label{lem:pair-route-orbit-as}
右単純$D$-加群$T_r=e_rD/e_rD_{\geq1}$について
\begin{equation}\label{eq:pair-route-orbit-ext}
 \underline{\Ext}^{q}_{D}(T_r,D)
 \simeq
 \begin{cases}
 {}_{\overline{r+N}}T(N),&q=h,\\
 0,&q\ne h
 \end{cases}
\end{equation}
である．ここで${}_sT$は左頂点$s$の単純加群である．
従って$D$は一般化AS正則次元$h$であり，graded twisted Calabi--Yau代数である．
通常の左右global dimensionも$h$となる．
\end{lemma}
\begin{proof}
$T_r$は$r$-sectorに属し，前補題の同値で$S_r$に対応する．
$D(s)=\bigoplus_{j=0}^{d-1}e_jD(s)$のうち$r$-sectorに属する項は，
$j\equiv r-s\pmod d$を満たす唯一の項である．
この項は$\Gamma_r$の逆同値で$P_{r-s}$に対応する．実際，その添字$t$の成分は
\begin{equation}
 (e_jD)_{t-r+s}e_{\bar t}
 =\mathcal B_{j,j+t-r+s}
 \simeq\mathcal B_{r-s,t},
\end{equation}
最後の同型は$\sigma$によるものである．
従って
\begin{equation}
 \Ext^q_{\operatorname{Gr}D}(T_r,D(s))
 \simeq\Ext^q_{\mathcal B}(S_r,P_{r-s}).
\end{equation}
\eqref{eq:pair-route-as-right}により，これは$q=h,s=-N$の場合に限り一次元である．
そのときの左頂点は$j=\overline{r+N}$である．
左正次数部分は内部次数を変えるので作用が零となり，
\eqref{eq:pair-route-orbit-ext}を得る．
各単純加群の有限生成射影解消も$\Gamma_r$から移る．

前補題のgraded global dimensionとこのExt条件は
Reyes--Rogalskiの一般化AS条件（右加群版）である．
しかも$D/D_{\geq1}=k^d$は$k$上separableなので，
\cite[Theorem~5.15]{RRGeneralizedAS}によりgraded twisted Calabi--Yau性と
通常の左右global dimensionについての主張が従う．
$D_0=k^d$のglobal dimensionが$0$であるだけで
$D$のglobal dimensionを結論しているのではなく，
前補題で全graded module categoryを比較していることに注意する．
\end{proof}

\subsection{Nakayama自己同型とそのlift}

\begin{lemma}[uniform indexとNakayama自己同型]
\label{lem:pair-route-nakayama}
$D$は次数を保つNakayama自己同型$\mu$を持ち，
\begin{equation}\label{eq:pair-route-mu-vertices}
 \mu(e_r)=e_{\overline{r+N}},\qquad \mu(z)=z
\end{equation}
とできる．本節の右AS規約では頂点置換は$+N$であり，$-N$ではない．
\end{lemma}
\begin{proof}
graded Nakayama双加群を
\begin{equation}
 U=\underline{\Ext}^{h}_{D^e}(D,D^e)
\end{equation}
とする．前補題により$U$はgraded invertibleである．
$S=D_0=k^d$にはtrace対合による$S^*\simeq S$という双加群同型がある．
\cite[Lemma~5.9(1)]{RRGeneralizedAS}を$X=S$に適用すると，
\begin{equation}\label{eq:pair-route-U-top}
 U\otimes_D S\simeq
 \underline{\Ext}^{h}_{D}(S,D)
\end{equation}
を得る．右辺は右$D$-加群$S$からのExtであり，
その右$e_r$成分は\eqref{eq:pair-route-orbit-ext}の
左単純加群${}_{\overline{r+N}}T(N)$である．
すなわち$Ue_r\twoheadrightarrow U\otimes_D {}_rT$の終項は
${}_{\overline{r+N}}T(N)$となる．

$U\otimes_D-$はgraded左加群圏の同値なので，
$De_r\twoheadrightarrow{}_rT$という射影被覆を射影被覆に移す．
従って，graded射影被覆の一意性から
\begin{equation}\label{eq:pair-route-U-free}
 Ue_r\simeq De_{\overline{r+N}}(N),\qquad
 {}_DU\simeq {}_DD(N)
\end{equation}
となる．各直和因子をこの順序で対応させる左加群同型を選ぶ．
$U$の右作用をこれによって$D(N)$へ移すと，
$\End_D(D)\simeq D^{\rm op}$から次数を保つ代数準同型
$\mu:D\to D$を得る．$U$がinvertibleなので右作用写像
$D^{\rm op}\to\End_D(U)$は同型であり，$\mu$は自己同型である．
従って$U\simeq {}^1D^\mu(N)$となる．
\eqref{eq:pair-route-U-free}で$Ue_r$が$De_{\overline{r+N}}(N)$に対応するため，
$\mu(e_r)=e_{\overline{r+N}}$である．

twisted Calabi--Yau代数のNakayama自己同型は中心を点ごとに固定する
\cite[Lemma~4.5]{RRGeneralizedAS}．
この事実は，$c\in Z(D)$に対して
$c\otimes1-1\otimes c^{\rm op}\in Z(D^e)$が対角双加群$D$に零作用するため，
その作用が$U=\underline{\Ext}^{h}_{D^e}(D,D^e)$でも零であることから従う．
$U\simeq{}^1D^\mu(N)$では$c=\mu(c)$を意味する．
補題\ref{lem:pair-route-orbit-center}を適用して$\mu(z)=z$を得る．
\end{proof}

\begin{lemma}[正規族を保つ$N$周期同型]
\label{lem:pair-route-rho}
族$(g_i)$を保つ$N$周期同型
\begin{equation}\label{eq:moduli-ambient-rho}
 \rho:\mathcal B\xrightarrow{\sim}\mathcal B(N)
\end{equation}
が存在する．
さらに$\rho$と$\sigma$は添字シフト込みで可換にできる．
\end{lemma}
\begin{proof}
各$i\leq j$に対して，orbit成分との線型同型
\begin{equation}\label{eq:pair-route-Fij}
 F_{ij}:\mathcal B_{ij}\xrightarrow{\sim}
       e_{\bar i}D_{j-i}e_{\bar j},\qquad
 F_{ij}(a)=\sigma^{-q_i}(a)
\end{equation}
を用いる．orbit積の定義から，これらは合成と両立する：
\begin{equation}
 F_{ij}(a)\cdot F_{j\ell}(b)=F_{i\ell}(ab).
\end{equation}
実際，左辺の第二因子に掛ける$\sigma$の冪は$q_j-q_i$であり，
全体は$\sigma^{-q_i}(a)\sigma^{-q_i}(b)$となる．

\eqref{eq:pair-route-mu-vertices}により
\begin{equation}\label{eq:pair-route-rho-formula}
 \rho_{ij}=F_{i+N,j+N}^{-1}\,\mu\,F_{ij}
\end{equation}
は定義できる．$\mu$の次数保存性と乗法保存性から
$\rho:\mathcal B\to\mathcal B(N)$は$\mathbb Z$-代数同型である．
また
\begin{equation}
 F_{i,i+d}(g_i)=e_{\bar i}z
\end{equation}
だから，$\mu(z)=z$と頂点置換から$\rho(g_i)=g_{i+N}$を得る．

最後に
$F_{i+d,j+d}\sigma_{ij}=F_{ij}$が成り立つので，
\eqref{eq:pair-route-rho-formula}から
\begin{equation}
 \rho_{i+d,j+d}\sigma_{ij}
 =\sigma_{i+N,j+N}\rho_{ij}
\end{equation}
が従う．従ってこの選択では二つの周期同型は可換である．
\end{proof}

\begin{theorem}[正規超曲面対の最大公約数周期性]
\label{thm:pair-route-gcd}
\label{thm:moduli-ambient-gcd}
\textup{(AS1)--(AS2)}を満たす$\mathcal B$と，次数$d$の左右正則正規族$(g_i)$に対して，
組$(\mathcal B,(g_i))$は$\gcd(N,d)$周期的である．
その族付きの最小正周期，イデアル付きの最小正周期，ambientのみの最小正周期は，
いずれも$\gcd(N,d)$を割る．
従って$I=(g_i)$，$\mathcal A=\mathcal B/I$もその周期同型を持つ．
特に標準5次元AS正則ambientと正規三次超曲面対は，
商イデアルを保つ1周期同型を持つ．
\end{theorem}
\begin{proof}
前補題の$\rho$と\eqref{eq:moduli-ambient-sigma}の$\sigma$は，それぞれ$N$と$d$だけ添字を動かし，
いずれも$g_i$を移動後の添字の$g$へ送る．
$ud+vN=\gcd(N,d)$となる整数$u,v$を選び，
$\theta=\sigma^u\rho^v$とする．これは所望の周期同型であり，$I$を保つので商にも降りる．
周期の集合は合成と逆同型により$\mathbb Z$の部分群となるので，
最小正周期についての主張も従う．
$N=5,d=3$の場合には
\begin{equation}\label{eq:pair-route-one-period}
 \theta=\rho^{-1}\sigma^2,\qquad 2\cdot3-5=1
\end{equation}
とすればよい．成分を明記すると
$\theta_{ij}=\rho^{-1}_{i+1,j+1}\sigma_{i+3,j+3}\sigma_{ij}$であり，
$\theta(g_i)=g_{i+1}$が成り立つ．
\end{proof}

\subsection{通常の次数付き商への移行と中心性の限界}

\begin{corollary}[三次三様体対の次数付き表示]
\label{cor:pair-route-graded-quotient}
\label{cor:moduli-cubic-ambient-periodic}
$\mathcal B$を標準$(1,5,10,10,5,1)$型の5次元Koszul AS正則
$\mathbb Z$-代数とし，次数3の正則正規族を持つとする．
このとき通常の連結次数付き代数$C$と
次数3の正則正規元$G\in C_3$を選び，
\begin{equation}
 (\mathcal B,I)\simeq(\check C,\check{(G)}),\qquad
 \mathcal A\simeq\check{C/(G)}
\end{equation}
と表示できる．$C$は標準5次元AS正則である．
特に$\dim\mathcal B_{i,i+m}=\binom{m+4}{4}$なら$H_C(t)=(1-t)^{-5}$である．
この表示は1周期同型の選択に依存し，一意な次数付き代数$C$を与えるとは主張しない．
$G$が中心的であることも一般にはこの構成からは従わない．
\end{corollary}
\begin{proof}
$C_m=\mathcal B_{0m}$，$a*b=a\theta^m(b)$（$a\in C_m$）と置く．
$\theta$の乗法保存性から$*$は結合的であり，
$\theta^i:C_{j-i}\to\mathcal B_{ij}$が$\check C\simeq\mathcal B$を与える．
$G=g_0$とすれば$\theta^i(G)=g_i$なので，この同型は三次イデアルを保つ．
左右正則性は同型から従う．次数付き加群圏と$\mathcal B$-加群圏の同値により
AS条件・global dimensionも移る．

正規自己同型を明記するため，
$\tau=\theta^{-d}\sigma$を$\mathcal B_{0m}=C_m$に制限する．
$\sigma$と$\theta$は可換なので$\tau$は$C$の次数付き自己同型である．
\eqref{eq:moduli-ambient-sigma}から
\begin{equation}\label{eq:pair-route-normal-tau}
 a*G=G*\tau(a)
\end{equation}
を得る．$N=5,d=3$の選択\eqref{eq:pair-route-one-period}では
\begin{equation}
 \tau=\rho^3\sigma^{-5}
\end{equation}
という添字を動かさない自己同型であり，これが恒等であるとは証明されていない．
従って中心性を結論しない．
\end{proof}

\begin{remark}[Zhang twistとの整合性]
\label{ex:pair-route-zhang}
体の標数を$0$とし，$S=k[x_0,\ldots,x_4]$，
\begin{equation}
 F=x_0^2x_1+x_1^3+x_2^3+x_3^3+x_4^3
\end{equation}
と置く．$V(F)\subset\PP^4$は滑らかである．実際，勾配が零なら
$x_2=x_3=x_4=0$，$x_0x_1=0$，$x_0^2+3x_1^2=0$から全座標が零になる．
$\varphi(x_0)=-x_0$，$\varphi(x_j)=x_j$（$j>0$）とし，
Zhang twist $C=S^\varphi$の積を
$a*b=a\varphi^{\deg a}(b)$とする．$\varphi(F)=F$だから
\begin{equation}
 a*F=F*\varphi(a)
\end{equation}
となり，$F$は正則正規であるが中心的ではない．
例えば$x_0*F=x_0F$に対して$F*x_0=-Fx_0$である．
一方
\begin{equation}
 \check C_{ij}\longrightarrow\check S_{ij},\qquad
 a\longmapsto\varphi^i(a)
\end{equation}
は$\mathbb Z$-代数同型であり，$F$が生成する三次イデアルも保つ．
従って同じ$\mathbb Z$-代数対が，中心三次商の表示と非中心的な正規三次商の表示を持つ．
本節の1周期化はこの例と両立し，特定の次数付き表示の中心性を不変量とする主張ではない．
\end{remark}

\subsection{低次数からの判定と適用範囲}
\begin{proposition}[低次数の核からの正規性判定]\label{prop:moduli-cubic-normal-criterion}
$\mathcal B$ は次数一生成で
$\dim\mathcal B_{i,i+1}=5$ とする。
斉次両側イデアル $I\subset\mathcal B$ が
\begin{equation}
 I_{i,i+3}=k g_i,\qquad \dim I_{i,i+4}=5
\end{equation}
を満たし，各 $g_i$ が左右正則なら，$(g_i)$ は次数3の正則正規族である。
\end{proposition}
\begin{proof}
$I_{i,i+4}$ の二つの部分空間
\begin{equation}
 g_i\mathcal B_{i+3,i+4},
 \qquad \mathcal B_{i,i+1}g_{i+1}
\end{equation}
は，正則性によってともに5次元であり，よって一致する。
これは長さ一の正規関係である。
次数一生成と結合律による帰納法で任意の $j\geq i$ に対する
\eqref{eq:moduli-ambient-normal-family}を得る。
\end{proof}

標準多項式型の行Hilbert級数$(1-t)^{-5}$と
次数3の正則元による商の級数$(1-t^3)/(1-t)^5$が既知なら、
次数3と4の差から
$\dim I_{i,i+3}=35-34=1$、$\dim I_{i,i+4}=70-65=5$を計算できる。
ただしHilbert級数だけでAS正則性や正則性を導いたわけではない。
一般に$N=n+2,d=3$なら$n=2,3$では対の1周期性が従う一方、
$n=4$では対の最小正周期は3の約数としか言えず、3周期の非1周期例の存在は主張しない。


\begin{remark}[適用範囲と既知結果との関係]
\label{rem:moduli-ambient-pair-scope}
本稿の中心三次対$(B_F,[G_F])$は初めから通常の次数付き代数と中心正則元からなるため，
その付随$\mathbb Z$-代数対は明らかに1周期的である．
系\ref{cor:moduli-cubic-ambient-periodic}は，この表示を初めから選ばなくても，
標準5次元AS正則ambientの正則正規三次超曲面対が1周期化できることを示す．
一方，復元される通常の次数付き表示の一意性，次数3の関係元の中心化，
ambientのCalabi--Yau化は別問題である．

本節は体上の正規元族を扱う．相対三次核が非自明な直線束で与えられる場合には，
$g_i$や同時の$d$周期同型を大域的に選べるとは限らない．
従って，これだけで族全体の1周期構造，その一意性，あるいは底への降下は結論しない．
抽象的な非可換cubicすべてがこの超曲面表示を持つことも主張しない．

正規族からのorbit構成と既知の一般化AS正則性・twisted Calabi--Yau双対性を
組み合わせた証明を上に与えたが，この組合せの初出・新規性の網羅的確認は未了である．
Koszul--Frobenius双対性自体，およびReyes--Rogalskiの各入力定理は既知の結果である．
\end{remark}

\section{周期構造を忘れる射とZhang同一視}
\label{sec:period-forgetting}
本節では体を代数閉体$k$、標数を0とする。
次数付き代数の同型と、その$\mathbb Z$-代数の同型は異なる。
この差を単なる$\Aut_{\mathrm{gr}}$の作用と同一視せず、
周期構造の共役問題として記述する。
その後、本稿の一般族では追加同一視が有限個であることを、
二因子stabilizerの整数行列証明書から示す。

\subsection{厳密なファイバーと固定同型類上の群oid}
連結な$\mathbb Z$-代数$\mathcal B$について
\[
 \operatorname{Per}_1(\mathcal B)
 =\{\theta:\mathcal B\xrightarrow{\sim}\mathcal B(1)
       \mid\theta\text{ は添字を1だけ動かす代数同型}\}
\]
と置く。$\mathcal B(1)_{ij}=\mathcal B_{i+1,j+1}$という記法であり、
$\theta$は各$\mathcal B_{ij}$を$\mathcal B_{i+1,j+1}$へ送る。
添字を固定する自己同型群を$G=\Aut_0(\mathcal B)$とする。
非空なら$\operatorname{Per}_1(\mathcal B)$は$G$のtorsorである。

\begin{proposition}[周期構造による圧縮と忘却]\label{prop:period-fiber}
周期構造$\theta$を選ぶと、
\begin{equation}\label{eq:period-compression-product}
 B^\theta_n=\mathcal B_{0n},\qquad
 a*_{\theta}b=a\,\theta^m(b)
 \quad(a\in\mathcal B_{0m},\ b\in\mathcal B_{0n})
\end{equation}
は連結次数付き代数となり、その付随$\mathbb Z$-代数は$\mathcal B$と同型である。
この構成は、次数付き代数の群oidと1周期構造付き$\mathbb Z$-代数の群oidの同値を与える。

周期構造を忘れる射の、\emph{指定された}$\mathcal B$への2ファイバーは
$\operatorname{Per}_1(\mathcal B)$を対象とし、恒等射だけを持つ群oidである。
一方、$\mathcal B$の同型類だけを固定した部分群oidは共役作用のaction群oidである。
そのfpqc局所形式も含むstackを表す記法では
\begin{equation}\label{eq:period-residual-quotient}
 [\operatorname{Per}_1(\mathcal B)/\Aut_0(\mathcal B)]
\end{equation}
となり、作用は共役作用である。
\end{proposition}
\begin{proof}
結合性は$\theta^{m+n}=\theta^m\theta^n$から従う。
$B^\theta_{j-i}=\mathcal B_{0,j-i}\to\mathcal B_{ij}$を$\theta^i$で与えると
付随$\mathbb Z$-代数の同型になる。
同型$g$が周期構造を保つ条件は$g\theta=\theta'g$であり、
これは圧縮した次数付き代数の同型条件である。
この圧縮はSierra\cite[Proposition~3.3]{SierraGAlgebras}のprincipal mapの構成である。
2ファイバーではambientへの同型もデータに含むため、
その同型を使ってambientを$\mathcal B$と同一視した後に許される射は恒等だけになる。
同型の指定を外したときに初めて共役による同一視が生じる。
\end{proof}

ここで2ファイバーが「恒等射だけを持つ」とは、パラメータ集合が零次元という意味ではない。
$G$自体は一般に無限個の次数成分にわたる自己同型を含み、
\eqref{eq:period-residual-quotient}を有限型代数的stackと呼ぶ根拠はまだない。

次数付き代数$B$の付随代数を$\check B$と書く。
Sierra\cite[Theorem~1.2, Corollary~4.4]{SierraGAlgebras}により
\begin{equation}\label{eq:sierra-zhang-comparison}
 \check B\simeq\check B'
 \quad\Longleftrightarrow\quad
 B'\text{ は }B\text{ のあるZhang twisting systemによる捩れと次数付き同型}
\end{equation}
である。単一の自己同型の冪による捩れはこの特別な場合であり、
全てのtwisting systemを最初からその形だとは仮定しない。
対$(\mathcal B,\mathcal I)$では、$G$と$\operatorname{Per}_1$を
それぞれ$\mathcal I$を保つ自己同型と周期構造に置き換える。

\subsection{二因子stabilizerからの有限次元化}
$B=T(V)/(R)$、$\dim V=N<\infty$とし、
\begin{equation}\label{eq:two-factor-stabilizer}
 H_R=\{(g,h)\in\PGL(V)\times\PGL(V)
       \mid(g\otimes h)R=R\}
\end{equation}
と置く。これは有限型の閉代数群である。
射影行列の持上げのスカラーは部分空間$R$を変えないので、条件はwell-definedである。

\begin{lemma}[中心二変数による一意延長]\label{lem:two-factor-injective}
本稿のPBW開部分の$B_F$について、$H_{R_F}$の両射影は単射である。
\end{lemma}
\begin{proof}
$(1,h)\in H_{R_F}$とし、$h\in\GL(V)$を持ち上げる。
中心元$u,v$と任意の$x\in V$の交換関係に$1\otimes h$を作用させると
\begin{equation}\label{eq:one-factor-central-equation}
 u h(x)=x h(u),\qquad u h(v)=v h(u)
\end{equation}
を得る。補題\ref{lem:pbw}の自由性と濾過から、次数2にはベクトル空間直和分解
\[
 (B_F)_2=k[u,v]_2\oplus(k[u,v]_1\otimes W)\oplus\widetilde{(S_{f_3})_2}
\]
がある。最後の項は任意に選んだactive次数2の持上げでよい。
$u h(v)=v h(u)$の$uW,vW$成分を比較すると、$h(u),h(v)$の$W$成分は零になる。
さらに$u^2,uv,v^2$を比較すると
$h(u)=\lambda u$、$h(v)=\lambda v$となる。
\eqref{eq:one-factor-central-equation}と$u$の正則性から$h(x)=\lambda x$であり、
$h$は射影的に恒等である。$(h,1)$も左右を交換した同じ議論で従う。
標数0の代数群は滑らかなので、これは群スキームの核も自明であることを意味する。
\end{proof}

\begin{proposition}[有限次元群による周期の比較]\label{prop:finite-dimensional-period-group}
$H_R$の両射影が単射とする。
このとき有限型の代数群$K_R\subset\PGL(V)$とその自己同型$\alpha$があり、
以下が成り立つ。
\begin{enumerate}
\item $\Aut_0(\check B)$を各次数間のスカラー変更で割った幾何点の群は$K_R(k)$である。
\item 同じ$\mathbb Z$-代数を持つ次数付き代数の同型類は、$K_R$上の作用
\begin{equation}\label{eq:twisted-conjugacy-action}
   k\cdot a=k a\alpha(k)^{-1}
\end{equation}
の軌道と一対一に対応する。
\item $\Aut_{\mathrm{gr}}(B)/\mathbb G_m=K_R^\alpha$である。
周期の無限小変更を添字付き代数の自己同型で同一視した空間は
\begin{equation}\label{eq:period-tangent-cokernel}
 \operatorname{coker}(1-d\alpha:\Lie K_R\to\Lie K_R)
\end{equation}
である。この空間は命題\ref{prop:period-fiber}の厳密な2ファイバーの接空間とは異なる。
\end{enumerate}
\end{proposition}
\begin{proof}
$H_1,H_2\subset\PGL(V)$を二つの射影像とする。
代数群の準同型の像は閉部分群であり、標数0で核が自明な準同型はその像への同型なので、
$H_R$は同型$\phi:H_1\xrightarrow{\sim}H_2$のgraphである。
添字を固定する$\check B$の自己同型は、各隣接次数の線型同型$g_i:V\to V$で
\[
  (g_i\otimes g_{i+1})R=R\qquad(i\in\mathbb Z)
\]
を満たすものに等しい。従ってその射影類は$g_0$から両方向に一意に延びる。
$g_0$が$n$歩まで両方向に延びる条件を$K_n$とすると、
$K_0=H_1\cap H_2$から始めて
\[
 K_{n+1}=\{g\in K_n:\phi(g),\phi^{-1}(g)\in K_n\}
\]
で定まる閉部分群の降鎖を得る。
Noether性により降鎖は有限段で安定する。
その安定値$K_R$では$\phi$が自己同型となり、これが両無限列の全体を表す。
特にこの有限次元化は、固定した$R$での主張である。
全ての$R$に一様な安定段数があるとは主張しない。

各$g_i$のスカラー変更は自由に選べる。
これらの群を$D\simeq\prod_{i\in\mathbb Z}\mathbb G_m$と書くと、
頂点のスカラー$(\lambda_i)$による$\lambda_{i+1}\lambda_i^{-1}$と同じ群である。
従って$D$は$\Aut_0(\check B)$の中心部分群で、幾何点では商が$K_R$になる。
標準周期構造を$\theta$とし、$\alpha=\operatorname{Ad}\theta$が誘導する
$K_R$の自己同型を取る。
他の周期構造を$a\theta$と書けば、共役作用は\eqref{eq:twisted-conjugacy-action}になる。
$D$上で$\alpha$は添字の平行移動であり、差分準同型
$d\mapsto d\alpha(d)^{-1}$は、初期値を一つ選ぶ逐次解法により全射である。
従ってスカラー変更から追加の同型類は生じず、その固定部分は
次数スカラー$\mathbb G_m$だけである。
これで軌道の対応を得る。一般の底上の無限積の降下や商stackの代数性を
ここから主張しているわけではない。

$\theta$と可換な自己同型では$g_i$が一定なので、
射影次数付き自己同型群は$K_R^\alpha$になる。
$a=1+\epsilon X$、$k=1+\epsilon Y$を
\eqref{eq:twisted-conjugacy-action}へ代入すると
$X\mapsto X+(1-d\alpha)Y$となり、接空間の主張を得る。
スカラー差分の逐次解法は局所Artin環の単元にも適用できる。
\end{proof}

\begin{corollary}[連続した追加同一視がないための条件]\label{cor:period-no-continuous-drop}
前命題の仮定の下で$\Aut_{\mathrm{gr}}(B)/\mathbb G_m$が有限なら、
\eqref{eq:period-tangent-cokernel}は零である。
同じ$\mathbb Z$-代数を持つ次数付き代数のうち、射影次数付き自己同型群が有限なものの
同型類は有限個しかない。さらに$H_R$そのものが有限なら、
自己同型群について追加仮定を置かず、全ての次数付き表示の同型類が有限個である。
\end{corollary}
\begin{proof}
有限性より$\ker(1-d\alpha)=\Lie(K_R^\alpha)=0$である。
$\Lie K_R$は有限次元なので余核も零になる。
捩れ共役作用で有限stabilizerを持つ軌道は$\dim K_R$次元であり、$K_R$の開部分になる。
異なる開軌道は同じ既約成分に入れないため、その数は有限である。
最後の場合は$K_R$自体が有限である。
\end{proof}

有限な$\Aut_{\mathrm{gr}}/\mathbb G_m$だけから結論するのではなく、
命題\ref{prop:finite-dimensional-period-group}による有限次元化を先に用いている。
また$K_R$が正次元の可能性が残る場合、同型類の次元とstackの次元はなお区別する。

\subsection{一般族の二因子stabilizerの整数証明書}
\begin{theorem}[一般ambientの有限な二因子stabilizer]\label{thm:generic-two-factor-finite}
本稿の相対Jacobian族には、$H_{R_F}$が有限である非空Zariski開部分がある。
従って一般ambientの次数付き同型類13次元、一般中心三次対の次数付き同型類17次元は、
添字付き$\mathbb Z$-代数または対の同型による追加同一視で連続的に減少しない。
正確には、既存の一般有理商上で、同じ$\mathbb Z$-同型類を与える点の集合は有限である。
\end{theorem}
\begin{proof}
式\eqref{eq:F0}の$F_0$を使う。二元三次部分は二次関係に寄与しない。
$V^{\otimes2}$を順序付き語$x_ix_j$、$\Sym^2V$を15個の可換単項式で座標表示する。
二次関係$R_F$は$\wedge^2V$から$\Sym^2V$へのgraphなので、
関係を表す$25\times10$行列$\widetilde R$と、
商への射影を表す$15\times25$行列$Q$を取れる。
対称化の分母を除くため、関係の列と商射影をそれぞれ2倍し、両行列を整数行列にする。
その場合も$Q\widetilde R=0$である。

二因子の無限小作用は次の$150\times50$整数行列で表される：
\begin{equation}\label{eq:two-factor-lie-matrix}
 M(X,Y)=\operatorname{vec}
       \bigl(Q(X\otimes1+1\otimes Y)\widetilde R\bigr).
\end{equation}
列は左因子$E_{ab}$の25個、次に右因子$E_{ab}$の25個の順とし、
行は10個の関係それぞれに対する15個の商座標とする。
$(I,0)$と$(0,I)$は整数上で独立な核ベクトルなので$\rank_{\mathbb Q}M\le48$である。
\nolinkurl{computations/verify_two_factor_stabilizer.py}は全行列をこの式から生成し、
\nolinkurl{computations/two_factor_stabilizer_certificate.json}の行・列で指定した小行列式について
\begin{equation}\label{eq:two-factor-nonzero-minor}
 \det M_{48}\equiv815096\pmod{1000003}
\end{equation}
を得る。従って$\rank_{\mathbb Q}M\ge48$であり、階数はちょうど48である。
有限体での非零を整数行列式の非零証明に使い、核ベクトルの等式は整数上で検証している。
よって$\PGL(V)\times\PGL(V)$におけるstabilizerのLie代数は零であり、
標数0の有限型代数群として$H_{R_{F_0}}$は有限である。
同じ非零小行列式の条件はZariski開条件であり、
既存のPBW・AS正則な非空開部分との共通部分も非空である。

この開部分で補題\ref{lem:two-factor-injective}と
系\ref{cor:period-no-continuous-drop}を適用する。
対$(\check B,\check I)$の自己同型群は$\Aut_0(\check B)$の部分群であり、
全ての頂点スカラーは$\check I$を保つ。
従ってスカラーを除いた対の自己同型群も有限な$K_R$の部分群となる。
イデアルを保つ周期構造についても同じ議論が適用でき、対の場合の有限性も従う。
最後に定理\ref{thm:17}および第\ref{sec:ambient-incidence}節の一般有理商の次元を用いる。
\end{proof}

この定理は、有限個の同一視の個数を求めたものではない。
また、13または17を全$\mathbb Z$-代数モジュライstackの次元と呼ぶには、
そのstackの定義・頂点スカラーの扱い・代数性の比較が必要である。
全モジュライの既約成分性もこの有限性だけからは従わない。

\subsection{一般点での中心条件とCY条件の形式的剛性}
\begin{lemma}[単位に還元する次数付き自己同型]\label{lem:scalar-automorphism-rigidity}
$B$を次数一生成の連結次数付き代数とし、$\Aut_{\mathrm{gr}}(B)/\mathbb G_m$は有限とする。
局所Artin $k$-代数$T$上の次数一生成平坦変形$B_T$について、
特殊ファイバーで恒等になる次数付き$T$-代数自己同型は、
ある$c\in1+\mathfrak m_T$による次数スカラー$\xi_c$である。
\end{lemma}
\begin{proof}
$\Lie\Aut_{\mathrm{gr}}(B)=kI$である。
$\mathfrak m_T$の冪について帰納し、自己同型を既に得た次数スカラーで割って
$\beta\equiv1\pmod{\mathfrak m_T^r}$とする。
$\mathfrak m_T^r/\mathfrak m_T^{r+1}$の各係数は、
関係保存の一次式により$B$の次数0導分を与えるので全てスカラーである。
これを$c$の次の係数へ吸収する。有限回で$\beta=\xi_c$となる。
\end{proof}

\begin{corollary}[中心三次の近傍では正規条件は中心条件になる]
\label{cor:normal-central-artin-rigidity}
前補題の下で$G\in Z(B)_3$を正則元とする。
その正則正規な持上げ$G_T\in(B_T)_3$が存在し、
正規自己同型$\sigma_T$が$G_Tb=\sigma_T(b)G_T$で定まるとする。
このとき$G_T$は中心元である。
\end{corollary}
\begin{proof}
正則性と$G$の中心性により$\sigma_T$の特殊ファイバーは恒等である。
前補題から$\sigma_T=\xi_c$、$c\equiv1\pmod{\mathfrak m_T}$となる。
$b=G_T$を正規性の式に代入し、$G_T$を消去すると$\sigma_T(G_T)=G_T$を得る。
従って$c^3=1$であり、3が可逆で$c\equiv1$なので$c=1$である。
\end{proof}

\begin{proposition}[添字付きnormal incidenceの形式的中心化]
\label{prop:indexed-normal-central-rigidity}
$B=T(V)/(R)$は次数一生成の連結二次代数で、
$H_R\subset\PGL(V)^2$は有限とする。
さらに、中心正則元$u\in B_1$と$G\in\Z(B)_3$が存在するとする。
局所Artin $k$-代数$T$上の任意の次数付き平坦二次変形$B_T$に対し、
次の二集合は、中央三次直線を全ての添字で繰り返す写像によって自然に一致する：
\begin{enumerate}
\item $kG$に還元する中心正則三次直線$L\subset(B_T)_3$；
\item $\check B_T$のnormal incidenceで、特殊ファイバーが
      $L_i=kG\subset B_3$となる直線束族$(L_i)_{i\in\mathbb Z}$。
\end{enumerate}
従って、ここでは次数付き表示を選んだcentral部分問題と、
添字付き代数の\emph{全normal incidence}の形式近傍が一致する。
ただし、底にあるambient変形は次数付き$B_T$から来るものとする。
\end{proposition}
\begin{proof}
各$L_i$の生成元$G_i$を$G$に還元するように選ぶ。
長さ四のnormal条件と左右の正則性から、一意な線型同型
$\phi_i:(B_T)_1\to(B_T)_1$があり、
\[
 G_i x=\phi_i(x)G_{i+1}\qquad(x\in(B_T)_1)
\]
を満たす。$G$の中心性と正則性より$\phi_i\equiv1\pmod{\mathfrak m_T}$である。
二次関係$r=\sum x\otimes y\in R_T$に対し、
\[
 0=G_i\sum xy
   =\sum\phi_i(x)\phi_{i+1}(y)G_{i+2}
\]
である。右の$G_{i+2}$を消去すると
$(\phi_i\otimes\phi_{i+1})R_T\subset R_T$を得る。
両辺は同階数の部分束で、特殊ファイバーでは同じなので、これは等号である。
従って射影対$([\phi_i],[\phi_{i+1}])$は
二因子stabilizerの$T$-点で、特殊ファイバーで恒等となる。

そのような$T$-点は恒等に限る。
実際、$\mathfrak m_T$の冪に関する最初の非零項は
$\Lie H_R\otimes_k(\mathfrak m_T^r/\mathfrak m_T^{r+1})$に入り、
$H_R$が有限で標数0だから零である。有限回の帰納で結論を得る。
この議論に相対stabilizer自体の平坦性は不要である。
よって$\phi_i=c_i1$、$c_i\in1+\mathfrak m_T$である。
$s_0=1$、$s_{i+1}=s_ic_i$を両方向に解き、$G_i$を$s_iG_i$で置き換えると
\begin{equation}\label{eq:indexed-central-recursion}
 G_i x=xG_{i+1}
\end{equation}
となる。

次数一生成により$G_i b=bG_{i+3}$が全ての$b\in(B_T)_3$に成り立つ。
$b=G_i$として左の$G_i$を消去すれば$G_{i+3}=G_i$となる。
そこで$H=(G_0+G_1+G_2)/3$と置くと、
\eqref{eq:indexed-central-recursion}と3周期性から$Hx=xH$であり、$H$は中心元である。
$D_i=G_i-H$は3周期的で、$D_0+D_1+D_2=0$を満たす。
$u$の任意の次数一の持上げ$u_T$について
\[
 D_i u_T=u_TD_{i+1}
\]
である。$D_i\in\mathfrak m_T^r(B_T)_3$とし、その
$\mathfrak m_T^r/\mathfrak m_T^{r+1}$での像を$d_i$とする。
特殊ファイバーの$u$は中心正則だから$u(d_i-d_{i+1})=0$より$d_i=d_{i+1}$となる。
その和が零で3が可逆なので全て$d_i=0$である。
$D_i$は初めから$\mathfrak m_T$に入るため、帰納により$D_i=0$を得る。
従って$L_i=TH$は全添字で同じ中心直線である。

逆に中心正則直線を繰り返せばnormal incidenceの条件を満たす。
逆写像は$(L_i)\mapsto L_0$であるため、上の選択には依存せず、Artin基底変換とも可換である。
\end{proof}

\begin{corollary}[固定ambient上の全ファイバーの四次元形式近傍]
\label{cor:indexed-normal-four-dimensional-fibre}
本稿の一般開部分で$\dim\Z(B)_3=5$とし、$G\in\Z(B)_3$を正則元とする。
指定したambient $\check B$上の全normal incidenceの$\check{(G)}$における形式完備化は
\[
 \widehat{\PP(\Z(B)_3)}_{[G]}
 \simeq\operatorname{Spf}k[[t_1,t_2,t_3,t_4]]
\]
である。これは定理\ref{thm:moduli-incidence-equivalence}の厳密な2ファイバーについての主張であり、
ambient自己同型による商の形式近傍を表すものではない。
\end{corollary}
\begin{proof}
前命題を定数変形$B_T=B\otimes_kT$に適用する。
この場合$\Z(B_T)_3=\Z(B)_3\otimes_kT$であり、
正則な特殊ファイバーを持つ元の乗法は各次数で部分束への単射となる。
従って$[G]$の全Artin近傍では正則性の追加方程式はなく、
中心直線の射影空間の形式完備化を得る。
\end{proof}

\begin{proposition}[相対Koszul--Frobenius条件下のNakayama剛性]
\label{prop:relative-nakayama-rigidity}
$B$を本稿の一般開部分にある5次元Koszul CY代数とし、$B_T$を局所Artin底上の平坦二次変形とする。
以下の相対条件を仮定する：Koszul双対$E_T$は次数0から5にある有限自由なFrobenius代数で、
$(E_T)_5\simeq T$、Frobenius積の対は完全であり、相対Koszul双加群複体は
Nakayama自己同型$\mu_T$を与える。その次数一部分はFrobenius自己同型$\nu_T$の双対
\begin{equation}\label{eq:relative-koszul-nakayama-dual}
 \mu_T|_{(B_T)_1}=(\nu_T|_{(E_T)_1})^*
\end{equation}
であるとする（次元5では通常の符号$(-1)^{5+1}$は正である）。
このとき$\mu_T=1$であり、$B_T$は相対CYである。
\end{proposition}
\begin{proof}
$\mu_T$は特殊ファイバーで恒等になるので、
補題\ref{lem:scalar-automorphism-rigidity}より$\mu_T=\xi_c$となる。
\eqref{eq:relative-koszul-nakayama-dual}から$\nu_T$は$(E_T)_1$に$c$倍で作用する。
Koszul双対は次数一生成なので、最高次数には$c^5$倍で作用する。
Frobenius汎関数$\varepsilon:E_T\to T$を取ると、
Nakayamaの定義に単位を代入して$\varepsilon\nu_T=\varepsilon$を得る。
最高次数で$\varepsilon$は同型だから$c^5=1$である。
$c\equiv1$と標数0から$c=1$となる。
\end{proof}

体上のKoszul双対によるNakayama比較とhomological determinantの恒等式は
Reyes--Rogalski--Zhang\cite[Theorem~6.3 and its proof]{RRZHomological}を参照できる。
全斉次成分が有限自由な平坦族については、
補題\ref{lem:p3-relative-bar}と注意\ref{rem:relative-bar-general-dimension}により、
二次双対の有限自由性と完全Frobenius対は幾何学的ファイバーから従う。
上の命題でなお明示的な入力として残すのは、相対双加群複体によるNakayamaの構成と
式\eqref{eq:relative-koszul-nakayama-dual}の具体的同定である。

Pym\cite[\S2.2, Definition~2.8]{Pym}による可換点近傍のCY規格化は、
Nakayamaの適切な根による捩れから得られる。
その自己同型による捩れの公式は
Reyes--Rogalski--Zhang\cite[Theorem~0.3]{RRZHomological}にある。
本節の有限同一視の定理は任意のtwisting systemを含むため、
Pymの議論を全$\mathbb Z$-代数の大域的な一意性へ拡張して使う必要がない。
残る比較は、一般の$\mathbb Z$-代数の\emph{族}が相対1周期構造を持つことと、
その降下・有限型の局所atlasを構成することである。

\section{非可換射影三空間の四頂点superpotential表示}
\label{sec:projective-three-potentials}

非可換$\PP^3$の標準的な座標$\mathbb Z$-代数は4周期的である．
この事実は四頂点quiverによる有限個の係数への還元を可能にするが，
任意の周期同型を選んだ表示と，対象自体のモジュライとは区別する必要がある．
本節ではKoszul双対のFrobenius traceから周期を規格化し，
まず体上の対象と自己同型についてこの差を処理する．
さらに，全成分が局所自由で幾何学的ファイバーが標準AS正則である族では，
相対bar複体からFrobenius双対性を導き，頂点スカラーを除いたfpqc stackの比較まで行う．
得られるpotentialは四次の符号付きsuperpotentialであり，
関係は長さ二の経路による微分から作る．通常の一階Jacobian代数ではない．

基礎体は$k=\C$とする．対象$C$は正連結・局所有限な二次Koszul
$\mathbb Z$-代数で，左右AS正則性，AS次元とGorenstein parameterがともに4，
各単純加群の有限最小線型解消の型が$(1,4,6,4,1)$であるとする．
従って$V_i=C_{i,i+1}$の次元は4，二次関係$R_i\subset V_i\otimes V_{i+1}$の
次元は6，行Hilbert級数は$(1-t)^{-4}$である．
以下で\emph{標準4次元データ}というときはこの対象域を指す．
必要なら左右noetherian性や第\ref{sec:categorical-moduli-comparison}節の
saturationを加えた部分問題へ制限できる．
単に四生成元・六関係という数値だけを仮定するものではない．
すべての同型は整数添字を保つものとし，添字の平行移動は同一視しない．

\subsection{Frobenius周期による規格化}

前向き二次双対を$E=C^!$と書き，$E_{i,i+1}=V_i^*$の規約を用いる．
Koszul--Frobenius双対性により
\begin{equation}\label{eq:p3-local-frobenius}
 \dim E_{i,i+p}=\binom4p\quad(0\le p\le4),\qquad
 E_{i,i+p}=0\quad(p>4)
\end{equation}
であり，相補次数の積は完全対合である
\cite[Proposition~4.1, \S5]{PBondalPolishchukHelices}．
この既知入力から周期同型を得る証明は
補題\ref{lem:moduli-ambient-koszul-period}の$N=4$の場合である．
以下ではその選択依存性も用いるので，必要な式を明記する．

各$i$で非零trace $\lambda_i:E_{i,i+4}\to k$を選ぶ．
完全対合により
\begin{equation}\label{eq:p3-frobenius-period}
 \lambda_i(ab)=\lambda_j\bigl(b\nu_{ij}(a)\bigr),
 \qquad a\in E_{ij},\quad b\in E_{j,i+4}
\end{equation}
が$\nu:E\xrightarrow{\sim}E(4)$を定める．
結合律と完全対合から$\nu(ac)=\nu(a)\nu(c)$となるので，これは代数同型である．
次数一部分で逆双対を取ると
$\tau^\lambda:C\xrightarrow{\sim}C(4)$を得る．
これを\emph{Frobenius周期}と呼ぶ．

\begin{lemma}[Frobenius周期の選択独立性]
\label{lem:p3-canonical-period}
二つのtrace族が与えるFrobenius周期は，頂点スカラー自己同型$h$により
\[
 \tau^{\lambda'}=h(4)\tau^\lambda h^{-1}
\]
と比較できる．また$C\simeq C'$という任意の添字保存同型は，
終域の頂点スカラー自己同型で補正すれば，指定された二つのFrobenius周期と可換にできる．
\end{lemma}
\begin{proof}
$\lambda'_i=t_i\lambda_i$とすると
\[
 \nu^{\lambda'}_{ij}=(t_i/t_j)\nu^\lambda_{ij}.
\]
$s_{i+4}/s_i=t_i$となる$s_i\in k^\times$を，四つの剰余類について
それぞれ一つの初期値から両方向へ再帰的に選ぶ．
$E$上の頂点スカラー$s_i/s_j$による共役は上のtrace変更を実現する．
逆双対を取れば，$C_{ij}$上の$s_j/s_i$倍による自己同型$h$が所望の式を満たす．
この構成に無限積は現れない．

同型$f:C\to C'$を与えると，双対の同型を通じてtraceを移送できる．
式\eqref{eq:p3-frobenius-period}の特徴付けにより，移送されたFrobenius周期は
$f(4)\tau^\lambda f^{-1}$である．
終域で選ばれたtraceとの差を，既に証明した頂点スカラー共役で補正すればよい．
\end{proof}

この補題の内容は\texttt{asz4}のFrobenius標準化の議論
\cite[定理\texttt{canonical-four}]{ASZ4Canonical}と同じであり，
本節では以後のCY表示と群oid比較に必要な証明を再掲した．
選択を含まない一つの具体的な線型写像$\tau$が得られた，という意味ではない．

\subsection{四頂点への畳込みと4-Calabi--Yau性}

一般の4周期同型$\tau:C\to C(4)$を選び，それに沿って各生成空間の基底を延長する．
すると$V_{i+4}=V_i$，$R_{i+4}=R_i$と厳密に表示できる．
四頂点$i\in\mathbb Z/4$と，$i$から$i+1$への四次元矢印空間$V_i$を持つ
quiverを$Q$とし，$S=k^4$，$V=\bigoplus_iV_i$と置く．
畳込み代数は
\begin{equation}\label{eq:p3-rolled-algebra}
 D(C,\tau)=T_S(V)/(R_0,R_1,R_2,R_3).
\end{equation}
長さ$m$の始点$i$の成分は$C_{i,i+m}$であり，
終点は$i+m\bmod4$である．
逆にこの四頂点の二次代数を整数列上に展開すれば，厳密な4周期同型を備えた$C$に戻る．
従ってこの操作は，周期を\emph{指定した}対象と，頂点を固定する次数付きquiver代数との同値である．

\begin{theorem}[Frobenius周期のCY orbit]
\label{thm:p3-canonical-cy-orbit}
$C$を標準4次元データとし，$\tau$をFrobenius周期とする．
このとき$D(C,\tau)$は二次Koszul 4-Calabi--Yau代数である．
異なるtraceから作った代数は頂点を保つ次数付き代数として同型である．
任意の4周期を使った畳込みについては，このuntwistedなCY結論を無条件に主張しない．
\end{theorem}
\begin{proof}
整数列上の線型解消を畳み込むと，$D_0=S$の各単純加群に線型解消を得る．
従って$D$はKoszulである．その二次双対は$E$を$\nu$で畳み込んだ有限次元代数
$\overline E$で記述できる．通常の$S$-双加群双対で矢印を逆向きに置く規約では
この前向き表示の反対代数を取るが，以下の対合とgraded対称性は反対代数でも保たれる．

式\eqref{eq:p3-frobenius-period}で$j=i+4$，$b=1$とすると
\[
 \lambda_{i+4}(\nu(a))=\lambda_i(a).
\]
従ってtraceは畳込みと両立する．$\overline E$上で
$\lambda=\sum_{i=0}^3\lambda_i$を次数4に置き，その他の次数では零とすると，
\eqref{eq:p3-frobenius-period}は
$\lambda(ab)=\lambda(ba)$という通常の対称性になる．
相補次数の完全性により，これはFrobenius汎関数である．

4次元のCY符号を得るには，頂点ごとにさらに符号を入れる．
\begin{equation}\label{eq:p3-graded-trace}
 \varepsilon|_{e_i\overline E_4e_i}=(-1)^i\lambda_i.
\end{equation}
これは$i\bmod4$で定義される．
$a\in\overline E_p$，$b\in\overline E_{4-p}$が積を持つとき，
始点は$p$だけ変わるので
\begin{equation}\label{eq:p3-graded-symmetry}
 \varepsilon(ab)=(-1)^p\varepsilon(ba)
                 =(-1)^{p(4-p)}\varepsilon(ba).
\end{equation}
よって$\overline E$はgraded対称Frobenius代数である．

この汎関数を四重テンソルへ双対化すると，次数4の符号付き巡回テンソル$W$を得る．
完全対合により，$W$の二階微分の空間は$D$の二次関係空間と一致し，
各階の微分の空間は二次双対の各次数の双対に一致する．
従って$W$から作る複体は$D$のKoszul双加群解消そのものであり，完全である．
Bocklandt--Schedler--Wemyssの定理
\cite[Theorem~6.2, Lemmas~6.4--6.5 and equation~(6.2)]{BSWPotentials}
により，この符号付き自己双対解消は4-CY性を与える．
これはKoszul性をFrobenius性だけから推論したものではなく，
最初から$C$のKoszul性を用いた議論である．

最後の同型独立性は補題\ref{lem:p3-canonical-period}から従う．
一般の周期を使うと，$\nu$と選んだ周期の差が畳込み後のNakayamaの捩れに残る．
その場合はtwisted CYとtwisted superpotentialの対応
\cite[Theorem~6.2の証明直後のtwisted版への拡張]{BSWPotentials}
を用いることになる．
\end{proof}

\begin{remark}[通常の対称性とCY符号]
前証明で通常の対称traceだけを使うと，4次元の符号を見落とす．
同じ現象は，Koszul双対の通常のNakayamaが恒等であるとき，
元代数のNakayamaには次数符号が現れる，と記述することもできる．
四頂点cycleでは$a\mapsto(-1)^{\deg a}a$は
$\sum_i(-1)^ie_i$による内的共役なので，CY性の障害にはならない．
式\eqref{eq:p3-graded-trace}はその内的補正をtrace側で実行している．
\end{remark}

\subsection{256係数と二階微分による関係}

以下，四つの矢印空間$V_i\simeq k^4$を固定する．
長さ四の閉経路の空間は
\[
 \bigoplus_{i=0}^3V_i\otimes V_{i+1}\otimes V_{i+2}\otimes V_{i+3}.
\]
最初の矢印を最後へ送る巡回回転を$c$と書く．
$c^4=1$であり，次数4の\emph{符号付きsuperpotential}は$cW=-W$を満たす元である．
これはBSWの符号$(-1)^{4-1}$に対応する
\cite[\S1, equation~(1.3)]{BSWPotentials}．

任意の
\begin{equation}\label{eq:p3-potential-tensor}
 w\in\mathcal W:=V_0\otimes V_1\otimes V_2\otimes V_3,
 \qquad \dim\mathcal W=4^4=256
\end{equation}
に対し
\begin{equation}\label{eq:p3-signed-potential}
 W(w)=w-cw+c^2w-c^3w
\end{equation}
と置く．始点0の成分を取ることが逆写像になるので，
これは$\mathcal W$と符号付きsuperpotential空間の線型同型である．
四つの成分は異なる始点の直和因子にある．
従って$\mathcal W$自体に，一頂点のテンソルで課すような追加の巡回不変条件はない．

経路$p$による左微分は，経路$q$に対し
\[
 \partial_p(q)=
 \begin{cases}r&q=pr,\\0&\text{それ以外}\end{cases}
\]
と定め，線型に延長する．矢印空間の双対による縮約として基底に依存せず定義できる．
二次関係と代数を
\begin{equation}\label{eq:p3-derivation-quotient}
 R_i(w)=\operatorname{span}\{\partial_pW(w):
       p\text{は}i+2\text{から}i\text{への長さ二の経路}\},
 \qquad D_w=T_S(V)/(R_i(w))
\end{equation}
と置く．ここで$R_i(w)\subset V_i\otimes V_{i+1}$である．
すなわち$D_w=\mathcal D(W(w),2)$というBSWのderivation quotientである．
一階微分なら三次関係となるため，本節の二次AS正則代数には対応しない．
また$W(w)$を通常の符号なしcyclic wordの商で扱うと，
$cW=-W$からその像は零になる．本節では経路代数内の符号付きテンソルとして扱う．

\begin{definition}[許容potentialの集合]
\label{def:p3-admissible-potentials}
$\mathcal U_4(k)\subset\PP(\mathcal W)(k)$を，$D_w$がKoszul 4-CYで，
各頂点の線型解消の型が$(1,4,6,4,1)$となる$[w]$の集合とする．
同値に，BSWの微分空間を$W_j$とすると，各始点におけるその次元が
$1,4,6,4,1$となり，複体
\begin{equation}\label{eq:p3-superpotential-complex}
 0\longrightarrow D_w\otimes_S W_4\otimes_S D_w
 \longrightarrow\cdots\longrightarrow
 D_w\otimes_S W_0\otimes_S D_w
 \longrightarrow D_w\longrightarrow0
\end{equation}
が完全であることを要求する．微分はBSWの式(6.1)に従う．
この同値性は同論文Theorem~6.2による．
\end{definition}

この定義で展開した$C_w$は標準4次元データである．
実際，線型双加群解消を各頂点の単純加群へ適用すると，
指定された左右の線型解消と，自己双対性によるAS双対性を得る．
そのEuler関係から各行のHilbert級数は$(1-t)^{-4}$となる．
逆に定理\ref{thm:p3-canonical-cy-orbit}は，すべての標準4次元データを
この形で与える．$w$の非零スカラー倍は関係空間を変えないため，
代数のモジュライには$w$自体ではなく$[w]\in\PP^{255}$を用いる．

\begin{remark}[有限係数表示と正則locus]
必要条件$\dim R_i(w)=6$は$16\times16$の縮約行列のrankが6という条件である．
rankが6以下という閉条件と6以上という開条件を同時に課すので，
全$\mathbb A^{256}$の開部分を取るだけの問題ではない．
有限個の微分空間の次元とFrobenius対合は有限行列で検査できる．
しかし\eqref{eq:p3-superpotential-complex}の完全性は全内部次数にわたる条件であり，
有限次元のFrobenius性だけからは従わない．
本節は$\mathcal U_4$全体を有限型の開または局所閉schemeとして構成したとは主張しない．
そのためには完全性・平坦性の有限判定や適切な有界性を別途証明する必要がある．
既にPBW性などを認定できるchartでは，その追加定理を用いて具体的な代数的locusを作れる．
\end{remark}

\begin{remark}[可換例と正則性を満たさないpotential]
\label{ex:p3-potential-examples}
各$V_i$の基底を$x_0^{(i)},\ldots,x_3^{(i)}$とする．
交代テンソル
\[
 w_{\mathrm{alt}}=
 \sum_{\sigma\in\mathfrak S_4}\operatorname{sgn}(\sigma)
 x_{\sigma(0)}^{(0)}x_{\sigma(1)}^{(1)}
 x_{\sigma(2)}^{(2)}x_{\sigma(3)}^{(3)}
\]
から\eqref{eq:p3-signed-potential}を作ると，二階微分の空間は
\[
 R_i=\operatorname{span}
 \{x_a^{(i)}x_b^{(i+1)}-x_b^{(i)}x_a^{(i+1)}:a<b\}.
\]
従って$C_{w_{\mathrm{alt}}}=\check{k[x_0,x_1,x_2,x_3]}$である．
四cycleの巡回置換の符号が負なので，この例でも$cW=-W$が必要である．

一方，
\[
 w_{\mathrm{bad}}=\sum_{a,b=0}^3
 x_a^{(0)}x_b^{(1)}x_a^{(2)}x_b^{(3)}
\]
を取ると，$V_2\otimes V_3$を縮約するだけで
$V_0\otimes V_1$の全16基底が得られる．
同じことがすべての始点で成立するので$R_i(w_{\mathrm{bad}})=V_i\otimes V_{i+1}$である．
よって長さ二以上の成分は零となり，標準Hilbert関数もAS正則性も満たさない．
この例は256係数のpotential空間全体を非可換$\PP^3$のパラメータ空間と
呼べないことを直接示す．
\end{remark}


\subsection{Directed Beilinson四窓と三次syzygyによるpotential復元}
\label{subsec:p3-directed-potential}

非可換$\PP^2$のBeilinson関係式と三頂点三次potentialの比較
\cite[\S7, 特に(7.27)--(7.28)]{AOUNCPlanes}は，
非可換$\PP^3$の場合に次の形で一般化できる．
ただし四次では二つの連続する二次関係空間があり，追加する矢の双対は
\emph{三次syzygy}から作る．

$V_0,V_1,V_2$をrank $4$のベクトル束とし，
$R_0\subset V_0\otimes V_1$，$R_1\subset V_1\otimes V_2$
をrank $6$の直和因子とする．
\begin{equation}\label{eq:p3-directed-K}
 K=(R_0\otimes V_2)\cap(V_0\otimes R_1)
 \subset V_0\otimes V_1\otimes V_2.
\end{equation}
相対的な交わりは二つの包含の差の核で定義する．
次の条件を課したfpqc部分関手を$\mathcal D_4$と書く：
(i) $K$がrank $4$の直和因子で底変更に両立すること；
(ii) 自然な縮約
\begin{equation}\label{eq:p3-directed-spanning}
 K\otimes V_2^*\longrightarrow R_0,\qquad
 V_0^*\otimes K\longrightarrow R_1
\end{equation}
がともに全射であること．
全体のパラメータは
$\Gr(6,V_0\otimes V_1)\times\Gr(6,V_1\otimes V_2)$
であるが，一般点の$K$は零になり得るので，
$\mathcal D_4$はこの積全体ではない．

\begin{proposition}[Directed relationと四次potentialの群oid同値]
\label{prop:p3-directed-potential}
$\mathcal W=V_0\otimes V_1\otimes V_2\otimes V_3$とする．
その非零テンソルの部分関手$\mathcal U^{\mathrm{dir}}$を以下で定める：
$w:V_3^*\hookrightarrow V_0\otimes V_1\otimes V_2$の像が
$K(R_0(w),R_1(w))$と一致するrank $4$の直和因子で，
二階縮約で得る$R_0(w),R_1(w)$がrank $6$の直和因子である．
このとき
\begin{equation}\label{eq:p3-dir-pot-groupoid}
 \left[\mathcal U^{\mathrm{dir}}/
 \prod_{i=0}^{3}\GL(V_i)\right]
 \simeq
 \left[\mathcal D_4/\prod_{i=0}^{2}\GL(V_i)\right].
\end{equation}
両側で展開した四頂点二次代数にKoszul・AS正則性を課した
全部分群oidにも，この同値は制限される．
ここで左辺は非射影化したテンソルである．
射影potentialを使う場合はスカラー慣性のrigidificationを
\emph{両側で}揃える必要がある．
\end{proposition}
\begin{proof}
$K$がrank $4$の直和因子なら$V_3=K^*$と置き，
恒等テンソル
\[
 w=\operatorname{coev}_K
 \in K\otimes K^*
 \subset V_0\otimes V_1\otimes V_2\otimes V_3
\]
を取る．局所基底では$w=\sum_{\alpha=1}^4k_\alpha c_\alpha$である．
$K\subset R_0\otimes V_2$と$K\subset V_0\otimes R_1$から，
二階縮約の像はそれぞれ$R_0,R_1$に含まれ，
式\eqref{eq:p3-directed-spanning}によって一致する．
残りの二箇所の関係空間は反対側の二対二flatteningの
transposeであり，同じrank $6$を持つ．

逆に$\mathcal U^{\mathrm{dir}}$の$w$は
$V_3^*\simeq K$を与えるので，この同型で
$w$は$\operatorname{coev}_K$に戻る．
固定$(R_0,R_1)$上のその同型の選択空間は$\GL(V_3)$-torsorである．
これはfpqc局所で自明であり，射と自己同型も同じであるから，
両群oidの同値を得る．
AS・Koszul条件は両表示が生成する\emph{同じ}四頂点代数に課す．
\end{proof}

四次のBSW型非捩れsuperpotentialは
\[
 W=w-cw+c^2w-c^3w,\qquad cW=-W
\]
であり，関係は長さ二の\emph{左}経路微分から得る．
符号なしの通常のcyclic path algebraに$W$を射影すると零になるので，
これは通常のDWZ型\emph{一階Jacobian代数}ではない．
他方，符号なしの巡回和$W_+=w+cw+c^2w+c^3w$も
四つの始点ごとに得られる二階縮約の関係\emph{空間}は同じである．
したがって「ordinary potential」の意味を，
BSWの符号付き高階微分か，符号なし巡回和の二階微分かで指定すべきである．

\begin{remark}[可換例，境界とmoduliの意義]
$V_i=V$，$R_0=R_1=\wedge^2V$なら
$K=\wedge^3V$（rank $4$）で，
両方の式\eqref{eq:p3-directed-spanning}の像は$\wedge^2V$である．
双対の体積形式を選べば$w$は完全交代四重テンソルに対応する．
全Grassmannianの一般の二関係空間では$K=0$となり得るので，
このpotential表示は\emph{非自明な三次syzygyのrank層}上の表示である．
境界で$\dim K>4$なら$V_3=K^*$という指定は破れ，
$\Gr(4,K)$のようなsyzygy選択を保持するincidence compactificationと，
関係式だけのGITコンパクト化は一般に一致しない．
従ってpotentialは分類のパラメータ数を直ちに減らす方法ではないが，
四つの二次関係を一つのテンソルにまとめ，
4-Calabi--Yau双対性・Koszul複体・Poisson量子化の条件を
同時に表すという意味で有用である．
大域AS locusの有限型性とGIT境界の同定は依然として未証明である．
\end{remark}

\subsection{軌道と自己同型の正確な比較}

$G=\prod_{i=0}^3\GL(V_i)$，$H=\prod_{i=0}^3\PGL(V_i)$とする．
$G$は$\mathcal W$にテンソル積で作用する．各矢印空間のスカラー
$Z_G=(k^\times)^4$は$w$を四つのスカラーの積だけ倍にするため，
$\PP(\mathcal W)$には自明に作用する．
従って射影potentialの有効作用は$H$を通る．
頂点はラベル付きのままであり，四頂点の回転による追加の商は取らない．

\begin{lemma}[potential直線の一意性]
\label{lem:p3-potential-unique}
許容代数$D_w$から，符号付きsuperpotentialの直線$kW(w)$を一意に復元できる．
従って，頂点を保つ次数付き同型$D_w\simeq D_{w'}$は，
$G$の元が$[w]$を$[w']$へ送ることと同値である．
\end{lemma}
\begin{proof}
Koszul双対の各頂点の最高次数は一次元である．
その二つのgraded対称Frobenius traceの比を$t_i\in k^\times$とする．
隣接する次数1と次数3の完全対合に
\eqref{eq:p3-graded-symmetry}を適用すると$t_i=t_{i+1}$となる．
従って全traceの比は一つの共通スカラーである．
traceを双対化して得る$W$も同じ共通スカラーを除いて一意である．
代数同型は二次双対とこのtrace直線を保つので，potential直線を保つ．
逆にpotential直線を保つ基底変更は二階微分の空間を保つため代数同型を与える．
\end{proof}

標準4次元データの群oidから，すべての頂点スカラー
\[
 \rho_\lambda(a_{ij})=\lambda_i\lambda_j^{-1}a_{ij},
 \qquad(\lambda_i)\in(k^\times)^{\mathbb Z}
\]
を射において同一視した群oidを$\mathcal Z_4^{\mathrm{pt}}(k)$と書く．
共通スカラーはもともと恒等に作用するので，$\lambda_0=1$としてよい．
saturationを課す場合，これは第\ref{sec:categorical-moduli-comparison}節の
一点の同定を保持する豊富列付き圏の群oidである．

\begin{theorem}[体上のpotential軌道による比較]
\label{thm:p3-potential-moduli}
次が成り立つ．
\begin{enumerate}[label=\textup{(\roman*)}]
\item 標準4次元データの添字保存同型類は
$\mathcal U_4(k)/G(k)=\mathcal U_4(k)/H(k)$と一対一に対応する．
\item 点の作用群oidとして
\begin{equation}\label{eq:p3-pointed-groupoid}
 [\mathcal U_4(k)/H(k)]\simeq\mathcal Z_4^{\mathrm{pt}}(k)
\end{equation}
が成り立つ．これはスキーム上の族のstack同値を表す式ではない．
\item $\tau$をFrobenius周期，$D=D(C,\tau)$とすると
\begin{equation}\label{eq:p3-automorphism-comparison}
 \frac{\Aut_{S,\mathrm{gr}}(D)}{(k^\times)^4}
 \simeq
 \frac{\Aut_{\mathbb Z}(C)}{\rho((k^\times)^{\mathbb Z})}.
\end{equation}
左辺の四つのスカラーは四つの矢印空間のスカラーである．
\end{enumerate}
\end{theorem}
\begin{proof}
定理\ref{thm:p3-canonical-cy-orbit}で各$C$から許容potentialを得る．
逆に$D_w$を展開すると標準4次元データ$C_w$を得る．
この展開の厳密な4周期はFrobenius周期であることも確認できる．
実際，$D_w$のgraded対称traceの頂点$i$成分を$(-1)^i$倍すると，
通常の対称traceになり，\eqref{eq:p3-frobenius-period}が厳密4移動を与える．
補題\ref{lem:p3-canonical-period}と補題\ref{lem:p3-potential-unique}から
同型類の全単射を得る．

自己同型とHom集合について，周期の選択を忘れる差を確認する．
$\Aut_{S,\mathrm{gr}}(D)$は$\Aut_{\mathbb Z}(C,\tau)$と同じである．
任意の$f\in\Aut_{\mathbb Z}(C)$に対し，移送された
$f(4)\tau f^{-1}$もFrobenius周期である．
補題\ref{lem:p3-canonical-period}により$f$を頂点スカラーで補正すれば
$\tau$と可換にできる．従って\eqref{eq:p3-automorphism-comparison}の右辺への写像は全射である．

その核は，周期と可換な頂点スカラーである．
$\rho_\lambda(4)=\rho_\lambda$は
\[
 \lambda_{i+4}/\lambda_i=\kappa\qquad(i\in\mathbb Z)
\]
となる一つの$\kappa\in k^\times$が存在することと同値である．
従って隣接矢印に作用する比$\lambda_i/\lambda_{i+1}$は4周期的である．
逆に任意の四つの非零比から$\lambda_0=1$を初期値としてこのような頂点スカラーを作れる．
よって核はちょうど$(k^\times)^4$であり，自己同型群の式を得る．

二つの異なる対象間の同型についても同じ補正を行える．
また二つの周期保存同型が頂点スカラーだけ異なるとき，
その差は今確認した四つの矢印スカラーである．
従って各Hom集合について全単射であり，対象については本質的全射である．
補題\ref{lem:p3-potential-unique}によるpotentialの作用群oidの表示と合わせて
\eqref{eq:p3-pointed-groupoid}を得る．
\end{proof}

\begin{remark}[スカラー商のlevel]
$[\mathcal U_4(k)/G(k)]$には四次元のスカラー自己同型が残るので，
$[\mathcal U_4(k)/H(k)]$とは同じ群oidではない．
式\eqref{eq:p3-pointed-groupoid}では，終域でも対応する頂点スカラーを除いたことが重要である．
全対象の同定を保持する厳密な$\mathbb Z$-代数stackとの同値ではない．
以下では，両側でこのスカラーの扱いを揃えてから相対比較を行う．
\end{remark}

\subsection{相対bar複体とpointed stackの比較}

体上で周期を規格化できるという事実だけでは，族の比較は得られない．
ここでは本稿の代数stackの定義に含まれる\emph{全成分の局所自由性}を使う．
$\mathfrak P_3^{\mathrm{AS}}$を，正連結・次数一生成の相対二次
$\mathbb Z$-代数$C/\OO_S$で，全$C_{ij}$が階数
$\binom{j-i+3}{3}$の有限局所自由層（$i\le j$），すべての幾何学的ファイバーが
標準4次元データとなるもののfpqc stackとする．
射は厳密な添字保存同型である．これは
定義\ref{def:moduli-two-stacks}のambientの定義を4次元へ置き換えたものである．

\begin{lemma}[相対Koszul双対は全底変更と両立する]
\label{lem:p3-relative-bar}
$C\in\mathfrak P_3^{\mathrm{AS}}(S)$とする．
その相対二次双対$E=C^!$の各成分は有限局所自由で，
\eqref{eq:p3-local-frobenius}の階数を持ち，任意の底変更と両立する．
特に$L_i=E_{i,i+4}$は直線束であり，積は完全対合
\begin{equation}\label{eq:p3-relative-frobenius}
 E_{ij}\xrightarrow{\sim}
 \mathcal Hom_{\OO_S}(E_{j,i+4},L_i)
 \qquad(i\le j\le i+4)
\end{equation}
を与える．基底の被約性やnoetherian性は必要ない．
\end{lemma}
\begin{proof}
始点$i$と内部長$n>0$を固定する．相対normalized bar複体の該当部分は
\begin{equation}\label{eq:p3-bar-complex}
 K^{i,n}_p=
 \bigoplus_{i=i_0<i_1<\cdots<i_p=i+n}
 C_{i_0i_1}\otimes_{\OO_S}\cdots\otimes_{\OO_S}C_{i_{p-1}i_p}
 \quad(1\le p\le n)
\end{equation}
であり，微分は隣接する二因子の積の交代和である．
これは有限個の有限局所自由層からなる有界複体であり，
底変更後も同じbar複体となる．$n=0$では次数0に$\OO_S$を置く．
幾何学的ファイバーのKoszul性により，そのhomologyは
$1\le n\le4$ならhomological degree $n$だけにあり，次元は$\binom4n$，
$n>4$なら全て零である．

有界な局所自由複体の全幾何学的ファイバーが一つの次数だけに
同じ次元のhomologyを持てば，複体は局所的にその次数のベクトル束と同値になる．
ここでは次の初等的な議論を使える．一点の近傍で項を自由にし，
その点で可逆な微分の小行列を順に消去すると，contractibleな二項直和因子を除ける．
残る複体の微分はその点で零なので，各項のその点での階数はhomologyの次元である．
従って一つの次数以外の項は消え，残る項は所定の階数の自由層となる．
この操作は開近傍で行え，被約性を使わない
\cite[Tags~0BDK, 0BCD]{StacksPerfect}．
従って$T_{i,i+n}:=H_n(K^{i,n})$は所定の階数のベクトル束であり，
他のhomologyは零である．この局所記述により，homologyは非平坦なものを含む
任意の底変更と可換になる．

bar複体の経路を切断するcoproductを双対化すると，
$E_{i,i+n}=T_{i,i+n}^*$に積を得る．これが相対二次双対であることも確認できる．
bar複体の最高項は$C_{i,i+1}\otimes\cdots\otimes C_{i+n-1,i+n}$であり，
その微分の核は各隣接位置に二次関係を置いた部分空間の共通部分である．
上で得た局所分裂を双対化すれば，その双対は二次関係の直交空間で
テンソル代数を割った商そのものである．
従って$E=C^!$であり，この同定と積も底変更に自然である．

最後に，\eqref{eq:p3-relative-frobenius}は同階数のベクトル束間の写像である．
各幾何学的ファイバーではAS双対性により同型なので，その行列式は局所的に単元となる．
従って相対的にも完全対合である．
\end{proof}

\begin{remark}[相対bar補題の次元に依存しない部分]
\label{rem:relative-bar-general-dimension}
補題\ref{lem:p3-relative-bar}の証明で$4$を$N$に置き換えると，
全成分が有限局所自由で，幾何学的ファイバーが標準$N$次元Koszul
AS正則である族について同じ結論を得る．
すなわち，二次双対の階数は$\binom Np$，最高次数は直線束で，
相補次数の積は任意の底変更と両立する完全対合である．
従って，本稿の5次元ambientの族にもこのFrobenius構造の存在は適用できる．
相対Nakayama双加群の具体的同定や，正規三次イデアルを保つ周期の降下は，
この補題単独の結論とは区別する．
\end{remark}

\begin{definition}[相対許容potential関手]
\label{def:p3-relative-potential-functor}
$U_4^{\mathrm{AS}}$を$\PP(\mathcal W)$のfpqc部分関手として次で定める．
$S$上のpotential直線$\mathscr L\subset\mathcal W\otimes\OO_S$について，
\eqref{eq:p3-signed-potential}から得る二階縮約の像$R_i$がrank6の部分束であり，
その関係で作った四頂点代数を展開した$C_w$が
$\mathfrak P_3^{\mathrm{AS}}(S)$に属することを要求する．
ここで部分束とは余商も局所自由であることを含む．
\end{definition}

この条件はfpqc局所的であり，底変更と両立する．
縮約の像については，定義に含めた部分束条件により核と余商が局所分裂することを使う．
体上では$U_4^{\mathrm{AS}}(k)=\mathcal U_4(k)$である．
実際，符号付き巡回性と二階微分の関係により，$W$の各頂点成分は
二次Koszul複体の長さ四の共通核に入る．
この核は一次元であり$W$の成分は非零なので，その生成元となる．
従って双対traceはgraded対称で，完全対合を持ち，
定理\ref{thm:p3-canonical-cy-orbit}と同じBSWの議論が適用できる．

また$\mathfrak P_3^{\mathrm{pt}}$を，
$\mathfrak P_3^{\mathrm{AS}}$の同型層を頂点スカラー
$\mathbb T_0=\{(\lambda_i)\in(\mathbb G_m)^{\mathbb Z}:\lambda_0=1\}$で割り，
その後fpqc stack化したものとする．
この定義は定義\ref{def:categorical-three-moduli}のpointed levelと同じである．
なお圏の族との同定には，同節で明示した相対section復元とsaturationの条件を用いる．

\begin{theorem}[相対potential表示とpointedモジュライ]
\label{thm:p3-potential-stack}
fpqc stackの同値
\begin{equation}\label{eq:p3-potential-stack-equivalence}
 [\,U_4^{\mathrm{AS}}/\prod_{i=0}^3\PGL(V_i)\,]
 \xrightarrow{\ \sim\ }\mathfrak P_3^{\mathrm{pt}}
\end{equation}
が成り立つ．左辺は上のfpqc部分関手の商stackとして定義する．
この定理は$U_4^{\mathrm{AS}}$のschemeとしての表現可能性，開性，有限型性を主張しない．
\end{theorem}
\begin{proof}
まず，固定した四つの矢印空間上のpotentialから$C_w$を作ることは底変更と可換である．
$\prod_i\GL(V_i)$の基底変更は周期を保つ$\mathbb Z$-代数同型を誘導する．
その中心$(\mathbb G_m)^4$は，整数列上で
$\lambda_{i+1}=\lambda_i/a_i$を解くことで頂点スカラーと書ける．
従って終域のpointed化の後は$\prod_i\PGL(V_i)$の作用に因子化し，
\eqref{eq:p3-potential-stack-equivalence}の射を得る．
$\PGL$の切断をfpqc局所で$\GL$へ持ち上げる際の二つの選択の差は中心であり，
pointed化の後に恒等となるので，この射は選択に依存せず降下する．

本質的全射性はfpqc局所で示せばよい．
$C\in\mathfrak P_3^{\mathrm{AS}}(S)$に対し，
補題\ref{lem:p3-relative-bar}により最高次数の直線束$E_{i,i+4}$と完全Frobenius対合を得る．
全ての最高次数直線束のframe束の積を取り，必要な四つの$V_i$も自明化すれば，
同時にtraceと基底を選べるfpqc被覆を得る
（命題\ref{prop:categorical-relative-descent}の証明と同じ議論である）．
この被覆上で\eqref{eq:p3-frobenius-period}は相対的な$\nu$と$\tau$を与える．
四頂点へ畳み，\eqref{eq:p3-graded-trace}を双対化すると
potential直線$\mathscr L\subset\mathcal W\otimes\OO_S$を得る．
完全対合により二階微分の像はもとのrank6の関係部分束に一致する．
その展開はもとの$C$なので，この直線は$U_4^{\mathrm{AS}}$に属する．

逆に$U_4^{\mathrm{AS}}$の点から作った$C_w$では，
$W$の各頂点成分が相対bar複体の最高次核に入る．
この核は補題\ref{lem:p3-relative-bar}による直線束であり，
potential直線からの写像は各ファイバーで同型なので相対的にも同型である．
従って，potential直線を局所的に自明化して得るgraded対称traceを
頂点符号$(-1)^i$で補正すると，展開の厳密4周期がFrobenius周期になる．

次に完全忠実性を同型層について局所的に確認する．
二つのこのような$C,C'$の厳密な同型$f$を与える．
最高次数直線束のtraceを局所的に自明化すれば，移送したtraceと指定したtraceの比は
全整数$i$に対する単元$t_i\in\OO_S^\times$となる．
補題\ref{lem:p3-canonical-period}の再帰式
$s_{i+4}/s_i=t_i$は単元の有限積によって相対的にも解ける．
従って$f$を頂点スカラーで補正すれば周期を保存し，四つの可逆な矢印写像から定まる．
potential直線の一意性も相対的に成り立つ．
実際，二つのtraceの比$t_i$は隣接する完全対合から$t_i=t_{i+1}$を満たすので，
共通単元に限る．従って四つの矢印写像はpotential直線を保つ．

二つの周期保存同型の差が頂点スカラーであるための条件は，
体上と同じ計算により，四つの矢印に4周期的なスカラーとして作用することである．
これはちょうど$\prod_i\GL(V_i)$の中心であり，商は$\prod_i\PGL(V_i)$である．
よって同型のfpqc層について完全忠実性を得る．
両側をstack化しているので，局所本質的全射性とこの完全忠実性からstack同値が従う．
\end{proof}

\begin{remark}[この商表示が含む形式と残る幾何]
左辺の$\PGL$-torsorは大域的な$\GL$-torsorへ持ち上がるとは限らない．
定理はその場合も両側のfpqc stack化で扱っており，
すべての対象に四つの大域的ベクトル束や，大域的な周期同型を指定しているのではない．
従ってBrauer型の捩れを，根拠なく無いものとして扱ってはいない．
右辺を厳密な$\mathfrak P_3^{\mathrm{AS}}$へ戻すと頂点スカラーの慣性が復活するので，
同じ$\PGL$商表示にはならない．

本定理によって，定義したpointed fpqc stackの比較自体は得られた．
残る代数幾何の問題は，$U_4^{\mathrm{AS}}$の全次数の平坦性とKoszul条件を
有限の係数条件で認識できる範囲，有界性，代数性，および既約成分の決定である．
$256-60$などの形式的な引算を，AS正則モジュライの次元とは解釈しない．
\end{remark}

\subsection{一般次元と五頂点ambientへの帰結}

\begin{corollary}[標準$N$次元のpotential商]
\label{cor:standard-n-potential-stack}
$N\ge2$とする．全成分が有限局所自由で，幾何学的ファイバーが
二次Koszul・左右AS正則，AS次元とparameterが$N$，
線型解消の型が$\bigl(\binom N0,\ldots,\binom NN\bigr)$となる
$\mathbb Z$-代数を対象とする．
各$V_i$が$N$次元である$N$頂点cycleを取り，
\[
 \mathcal W_N=\bigotimes_{i=0}^{N-1}V_i,\qquad
 W_N(w)=\sum_{i=0}^{N-1}(-1)^{(N-1)i}c^i w
\]
と置く．関係は長さ$N-2$の経路による微分で与える．
$U_N^{\mathrm{AS}}\subset\PP(\mathcal W_N)$を，微分像が
rank $\binom N2$の部分束で，この代数を展開したものが
上記の相対標準AS条件を満たすfpqc部分関手とすると，
\[
 [\,U_N^{\mathrm{AS}}/\PGL_N^N\,]
 \simeq\mathfrak P_{N-1}^{\mathrm{pt}}
\]
が成り立つ．特に$N=5$では
\begin{equation}\label{eq:p4-five-potential-stack}
 [\,U_5^{\mathrm{AS}}/\PGL_5^5\,]
 \simeq\mathfrak P_4^{\mathrm{pt}},
 \qquad U_5^{\mathrm{AS}}\subset\PP^{3124}.
\end{equation}
ここでも部分関手全体のscheme表現可能性や有限型性は主張しない．
\end{corollary}
\begin{proof}
注意\ref{rem:relative-bar-general-dimension}が相対Frobenius双対を与える．
traceを選ぶと$\nu:E\to E(N)$が得られ，
Frobenius周期で畳み込んだ双対のtraceは通常対称となる．
頂点$i$のtraceを
\[
 \varepsilon_i=(-1)^{(N-1)i}\lambda_i
\]
とする．$N(N-1)$が偶数なのでこの符号は$i\bmod N$で定義される．
次数$p$と$N-p$の相補対合では
$(-1)^{(N-1)p}=(-1)^{p(N-p)}$となり，graded対称性を得る．
従ってBSWのKoszul双加群複体は$N$-CY解消であり，
そのpotentialは表示した$W_N$となる．

trace変更の補正式は$s_{i+N}/s_i=t_i$であり，
$N$個の剰余類ごとに単元の有限積で解ける．
周期を保存する頂点スカラーは，各矢空間に独立な$N$個のスカラーとして作用する．
相対bar，potential直線の一意性，frame束のfpqc自明化およびHomの商の証明は
定理\ref{thm:p3-potential-stack}と同一であり，
$\GL_N^N$の中心を除くと主張の$\PGL_N^N$商になる．
$N=5$では符号は全て正，係数の数は$5^5=3125$，
二次関係を与える微分の長さは3である．
\end{proof}

式\eqref{eq:p4-five-potential-stack}はambient側だけの結論である．
三次イデアルを指定した対について，Frobenius周期がそのイデアルを
相対的に保つことまでは証明していない．
従って次節の正規三次incidenceを，すべての基底で自動的に
五つの直線束だけへ置き換える根拠にはしない．

\subsection{可換点の完備化における二つのrank条件}
\label{subsec:p3-formal-rank}

前述の大域的な完全性の問題に対し，可換点の形式近傍では有限個の
determinantal equationsだけで十分である．以下のrankは体上の点のrankではなく，
すべての$7\times7$ minorの消失と，少なくとも一つの$6\times6$ minorの可逆性を
要求するscheme論的なrank条件である．

テンソル$w$の二つのflatteningを
\[
 M_{01|23}(w):(V_2\otimes V_3)^*\longrightarrow V_0\otimes V_1,
 \qquad
 M_{12|30}(w):(V_3\otimes V_0)^*\longrightarrow V_1\otimes V_2
\]
とし，いずれもrank $6$となる局所閉部分schemeを
\[
 Z_6\subset\PP(\mathcal W)
\]
と書く．残る二つのflatteningはこれらの転置と符号で異なるだけなので，
この条件は全ての$R_i(w)$のrankが$6$であることと同値である．
$[w_{\mathrm{alt}}]\in Z_6$であり，その点で非零な二つの$6$次minorを
固定した開chartだけを用いても，以下の完備局所環は変わらない．

\begin{theorem}[rank条件による可換点の形式模型]
\label{thm:p3-formal-rank-six}
$R$を剰余体$k=\C$の局所Artin $k$-代数とし，
$[w_R]\in Z_6(R)$の剰余点が$[w_{\mathrm{alt}}]$であるとする．
このとき$D_{w_R}$の長さ$m$の各始点成分は
rank $\binom{m+3}{3}$の自由$R$-加群であり，
各頂点加群$R e_i$は型$(1,4,6,4,1)$の相対線型Koszul解消を持つ．
さらにpotential双加群複体は相対的に完全かつ自己双対である．
従って$D_{w_R}$は$R$上相対的にhomologically smoothな4-CY代数であり，
相対AS双対性も満たす．

逆に，$[w_{\mathrm{alt}}]$を持ち上げる射影potentialで，
$D_{w_R}$が指定されたHilbert ranksを持つ平坦族となるものは$Z_6(R)$に属する．
従って
\[
 \widehat{(Z_6)}_{[w_{\mathrm{alt}}]}
\]
は，可換点を特殊化として持つ，矢印空間を固定した射影potential付き
標準4次元データの形式変形関手を表す．
\end{theorem}
\begin{proof}
$\Lambda_R=R^4$，$A=D_{w_R}$とし，$\mathcal W_p$を
$W(w_R)$の$(4-p)$階微分の像とする．$R$は局所環なので$w_R$の生成元を選べる．
この選択を単元倍しても，微分空間も代数も変わらない．

各始点の$\mathcal W_4$は，特殊ファイバーで非零な係数を持つ
$W(w_R)$の成分が張る直線部分加群であり，rank $1$の直和因子である．
各始点の$\mathcal W_3$はrank $4$の自由加群からの縮約写像の像である．
特殊ファイバーである$4$次minorが非零なので，この写像はsplit injectionであり，
像はrank $4$の直和因子である．$\mathcal W_2$は仮定によりrank $6$の
直和因子である．同様に，特殊ファイバーで全射である三階・四階縮約から
$\mathcal W_1=V_R$，$\mathcal W_0=\Lambda_R$を得る．
従って各始点のranksは順に$1,4,6,4,1$である．

符号付き巡回性から，$\mathcal W_p$の一つの端の因子を切り離すと
$V_R\otimes_{\Lambda_R}\mathcal W_{p-1}$および
$\mathcal W_{p-1}\otimes_{\Lambda_R}V_R$へ入る．
隣接二因子を積にすると二次関係で消えるので，左Koszul微分は二乗零となる．
これはBSWのpotential複体の構成
\cite[\S6, equation~(6.1) and Lemma~6.5]{BSWPotentials}
と同じテンソル恒等式であり，$R$上でも成立する．

内部長$m$に関する帰納法で，$A_m$の自由性と左Koszul複体の完全性を同時に示す．
$m=0$では$A_0=\Lambda_R$である．$m>0$とし，$A_r$が$r<m$について
所定のrankの自由加群であると仮定する．正のhomological degreesだけを残した複体
\begin{equation}\label{eq:p3-positive-koszul-induction}
 F^{(m)}_p=(A\otimes_{\Lambda_R}\mathcal W_p)_m
          =A_{m-p}\otimes_{\Lambda_R}\mathcal W_p,
 \qquad 1\le p\le4
\end{equation}
を考える．負の添字の項は零とする．各項は帰納仮定により有限自由である．
特殊ファイバーは多項式環の四頂点表示のKoszul複体から
homological degree $0$の項を除いたものである．従ってhomologyはdegree $1$だけにあり，
各始点でその次元は$\binom{m+3}{3}$である．

perfect複体の一つの次数への集中判定
\cite[Tags~0BDK, 0BCD]{StacksPerfect}により，
$F^{(m)}_\bullet$もdegree $1$以外のhomologyを持たず，
$H_1(F^{(m)}_\bullet)$は所定のrankの自由$R$-加群である．
ここで$H_1(F^{(m)}_\bullet)=A_m$は二次presentationから直接従う．
実際，$A_{m-1}\otimes_{\Lambda_R}V_R$を
$A_{m-2}\otimes_{\Lambda_R}\mathcal W_2$の像で割ると，
長さ$m$の全ての位置に二次関係を課した商，すなわち$A_m$になる．
この同定により$A_m$の自由性を得るとともに，
$F^{(m)}_1\to A_m$を付加した複体の完全性も得る．帰納法が完了する．
特に，この議論は未知の$A_m$の平坦性を帰納段階の仮定に含めていない．

次に双加群potential複体を考える．全ての$A_m$が自由になったので，
その内部長$m$の各項
\[
 \bigoplus_{a+b+p=m}
 A_a\otimes_{\Lambda_R}\mathcal W_p\otimes_{\Lambda_R}A_b
\]
およびaugmentationの終域$A_m$は有限自由である．
特殊ファイバーは多項式環の完全な双加群Koszul解消であるから，
同じperfect複体の判定により，相対双加群複体も全内部次数で完全である．

potentialによる$\mathcal W_p$と$\mathcal W_{4-p}$の対合は完全である．
実際，関連する縮約行列をrank $\binom4p$の可逆小行列でsplitし，
その像と核の商を用いれば完全対合を得る．
従ってBSWの符号付き自己双対性の証明
\cite[Lemma~6.4 and the proof of Theorem~6.2]{BSWPotentials}
は$R$上でも成立する．各$\mathcal W_p$は$\Lambda_R$-双加群として有限射影なので，
これは相対enveloping algebra $A\otimes_R A^{\mathrm{op}}$上の
有限射影自己双対解消である．よって相対4-CY性を得る．
同じ解消を左右の頂点加群に用いると相対AS双対性を得る．
ここで$R$自体の正則性を仮定しておらず，
$k$上の絶対的global dimensionについての主張ではない．

逆向きには，各始点の長さ二の成分がrank $10$の自由加群であることから，
その二次関係空間はrank $6$の直和因子となる．
二階縮約行列の像はこの関係空間そのものなので，
各行列は基底変更により$\operatorname{diag}(I_6,0)$となる．
従ってscheme論的rank $6$条件を満たす．
両構成は底変更に自然であり，形式関手の主張を得る．
\end{proof}

\begin{proposition}[可換potential点の接空間と軌道]
\label{prop:p3-rank-six-tangent}
$\widetilde Z_6\subset\mathcal W$を上記二つのrank $6$条件のaffine locusとする．
このとき
\[
 \dim T_{w_{\mathrm{alt}}}\widetilde Z_6=91,
 \qquad
 \dim T_{[w_{\mathrm{alt}}]}Z_6=90.
\]
$G=\prod_{i=0}^3\GL(V_i)$のaffine orbitの次元は$46$，
射影orbitの次元は$45$である．従って射影potentialの変形を
恒等に還元する基底変更で割った一次変形の空間は$45$次元である．
\end{proposition}
\begin{proof}
指定した基底で全ての$V_i$を四次元空間$V$と同一視する．
交代テンソルの第一flatteningの核と余核はともに$\operatorname{Sym}^2V$型の
$10$次元空間である．rank $6$ determinantal locusの接条件は
核から余核への誘導写像が零となることなので，
第一条件の$100$個の独立線型汎関数は
\[
 U_{01|23}=
 (V^{*\otimes4})^{\langle(01),(23)\rangle}
 \cong\operatorname{Sym}^2(V^*)\otimes\operatorname{Sym}^2(V^*)
\]
を張る．第二条件のrow spaceは同様に
$U_{12|30}=(V^{*\otimes4})^{\langle(12),(30)\rangle}$である．
これら四つのtranspositionは$\mathfrak S_4$を生成するため，
\[
 U_{01|23}\cap U_{12|30}
   =(V^{*\otimes4})^{\mathfrak S_4}
   =\operatorname{Sym}^4(V^*),
 \qquad \dim\operatorname{Sym}^4(V^*)=35.
\]
従って全線型条件のrankは$100+100-35=165$であり，
affine接空間の次元は$256-165=91$である．
非零tensorのスカラー方向を除くと射影接空間の次元は$90$となる．

最後にLie安定化群を計算する．$X_i\in\mathfrak{gl}(V)$について
\[
 \sum_{i=0}^3
 (1\otimes\cdots\otimes X_i\otimes\cdots\otimes1)w_{\mathrm{alt}}
 \in k w_{\mathrm{alt}}
\]
を仮定する．二つの位置に同じ基底添字$b$が現れ，添字$a\ne b$が欠ける係数を
比較すると，任意の$i,j$について$(X_i)_{ba}=(X_j)_{ba}$を得る．
次に四つの添字が全て異なる係数を，二位置の添字を交換して比較すると，
$(X_i)_{aa}-(X_j)_{aa}$は$a$によらない．
従って$X_i=H+c_iI$であり，$H\in\mathfrak{sl}(V)$と取れる．
逆にこの形は全て射影potentialを保つ．
よって射影Lie安定化群の次元は$15+4=19$，射影orbitは$64-19=45$次元である．
affine potential自体を保つ条件はさらに$\sum_i c_i=0$であるため，
affine安定化群は$18$次元，orbitは$46$次元となる．
一次変形を基底変更で割った空間の次元は$90-45=45$である．
\end{proof}

\begin{remark}[形式模型の範囲]
定理\ref{thm:p3-formal-rank-six}は，全ての局所Artin基底上での同定である．
それは$Z_6$の可換点における完備局所環を有限個のminorで記述するが，
同じrank条件を満たす全ての体上のpotentialの正則性や，
$Z_6$の一つのZariski開近傍全体での全次数の平坦性を証明したものではない．
また，接空間の次元計算だけから形式関手の同型を推論したのではなく，
先の全Artin基底上の完全性の証明が形式同定を与えている．
\end{remark}

\begin{corollary}[五頂点ambientの可換点の形式方程式]
\label{cor:p4-formal-rank-ten}
$N=5$の係数空間$\PP(\mathcal W_5)=\PP^{3124}$において，
五つのcyclic flattening
\[
 (V_{i+2}\otimes V_{i+3}\otimes V_{i+4})^*
   \longrightarrow V_i\otimes V_{i+1},
 \qquad i\in\mathbb Z/5,
\]
がscheme論的にrank $10$となる局所閉schemeを$Z_{10}^{(5)}$とする．
可換五変数環を与える交代テンソル$w_{\mathrm{alt},5}$において，
$\widehat{(Z_{10}^{(5)})}_{[w_{\mathrm{alt},5}]}$は
矢印空間を固定した射影potential付き標準5次元データの
全Artin形式変形関手を表す．
\end{corollary}
\begin{proof}
各flatteningは$25\times125$行列である．
rank条件は長さ二の微分像$\mathcal W_2$をrank $10$の直和因子にする．
符号なし巡回性により，長さ三の微分像を与える行列は
これらの転置なので，$\mathcal W_3$もrank $10$となる．
可換特殊ファイバーのunit minorにより，
$\mathcal W_4,\mathcal W_5,\mathcal W_1,\mathcal W_0$の各頂点rankは
それぞれ$5,1,5,1$である．
従って微分空間の全ranksは$(1,5,10,10,5,1)$となる．
定理\ref{thm:p3-formal-rank-six}の内部長帰納法を，
$1\le p\le5$としてそのまま適用でき，
全次数の平坦性，Koszul解消，相対5-CY性を得る．
逆に平坦な標準Hilbert型なら，長さ二の商のrank $15$から
関係像はrank $10$の部分束であり，上の五つのrank条件を満たす．
これらの構成は底変更に自然である．
\end{proof}

この系は有限個のminorによる形式方程式を与えるが，
$Z_{10}^{(5)}$全体のAS性や既約成分を決定したものではない．
五頂点の接空間行列の直接計算も，本稿の四頂点の有限検算とは別課題である．

\section{非可換三次三様体側との比較}
\label{sec:cubic-p3-comparison}

非可換$\PP^3$と非可換三次三様体は、射影的な次元はともに3であるが、
座標代数のホモロジーは異なる。
本研究で三次商から復元するambientは非可換$\PP^4$であり、
二次包絡が与えるモジュライの射も三次側から$\PP^4$側への射である。
そのambient側のpointed stackには
系\ref{cor:standard-n-potential-stack}の五頂点potential商表示が適用できる。
以下の比較は、三次側から非可換$\PP^3$への自然な射を構成するものではない。

\subsection{AS正則代数とAS Gorenstein超曲面商}

体の標数を0とする。$B$は左右noetherianな連結次数付き代数で、
次数一生成、二次Koszul AS正則、global dimensionとGorenstein parameterがともに5、
$H_B(t)=(1-t)^{-5}$であるとする。
$G\in B_3$は非零の左右正則正規元とし、$A=B/(G)$と置く。
正規性とは$GB=BG$であり、中心性はここでは仮定しない。

\begin{proposition}[三次商のホモロジー的型]
\label{prop:cubic-p3-gorenstein-type}
$A$はAS Gorensteinで、左右のinjective dimensionは4、
Gorenstein parameterは2である。一方、$k$の左右のprojective dimensionと
$A$の左右のglobal dimensionは無限である。
また$A$は二次代数ではなく、従って通常のKoszul代数でもない。
\end{proposition}
\begin{proof}
正則正規元によるReesの補題はAS Gorenstein性を保存し、
injective dimensionを一つ下げる
\cite[arXiv版v1, Lemma~4.4]{CSBGCCKM}。
内部次数も記録したchange-of-ringsの同定は、単純加群上では
正規自己同型による捩れを除いて
\[
 \underline{\operatorname{Ext}}^{i}_{A}(k,A)
 \simeq\underline{\operatorname{Ext}}^{i+1}_{B}(k,B)(-3)
\]
となる。従って非零の項は$i=4$での$k(5-3)=k(2)$だけであり、
左右ともGorenstein parameterは2である。

正則性と$0\to B(-3)\xrightarrow{\,G\,}B\to A\to0$から
\begin{equation}\label{eq:cubic-p3-hilbert}
 H_A(t)=\frac{1-t^3}{(1-t)^5},\qquad
 H_A(t)^{-1}=\frac{(1-t)^4}{1+t+t^2}.
\end{equation}
仮に$k$が有限projective dimensionを持てば、noetherian性により
有限生成自由加群による有限の最小次数付き解消を持つ。
そのHilbert級数の交代和から$H_A(t)^{-1}$は多項式となる。
しかし\eqref{eq:cubic-p3-hilbert}の分母は原始三乗根で零になり、
分子はそこで零にならないため矛盾する。左加群にも同じ証明が使える。

$G$の次数は3なので$A_1=B_1$、$A_2=B_2$である。
従って$A$の二次包絡は$B$に等しく、
$\ker(B_3\to A_3)=kG\ne0$である。
これは既存の二次関係からは従わない三次関係であり、$A$は二次ではない。
\end{proof}

比較を表にまとめる。最後の行は、適切なnoetherian条件の下での
Orlovの定理\cite[Theorem~2.5]{OrlovSingularities}による。
$\operatorname{pd}_A k=\infty$なので$k$はperfectでなく、
$D_{\mathrm{sg}}^{\mathrm{gr}}(A)$は非零である。
\begin{center}
\small
\begin{tabular}{p{.23\textwidth}p{.32\textwidth}p{.35\textwidth}}
\hline
 & 非可換$\PP^3$の標準座標代数$C$ & 非可換三次三様体の座標代数$A$\\
\hline
一次生成元の数 & $4$ & $5$\\
基本関係 & 二次関係 & ambientの二次関係と正規三次関係\\
Hilbert級数 & $(1-t)^{-4}$ & $(1-t^3)(1-t)^{-5}$\\
ホモロジー的次元 & AS正則、$\operatorname{gldim}C=4$ & AS Gorenstein、$\operatorname{injdim}A=4$、$\operatorname{gldim}A=\infty$\\
Gorenstein parameter & $4$ & $2$\\
射影導来圏の標準分解 & 四つの線対象による完全例外列 & 二つの線対象と非自明な特異圏成分\\
\hline
\end{tabular}
\end{center}
AS Gorensteinのinjective dimensionが4であることを、
4次元AS正則性や4-Calabi--Yau性と読み替えることはできない。
また$A$のglobal dimensionが無限でも、$\operatorname{qgr}A$が滑らかであることとは両立する。

\subsection{三次側の有限potentialデータ}

Bocklandt--Schedler--Wemyssの定理によれば、Koszul twisted CY代数は、
twisted superpotentialとその高階導分から復元される。
逆向きには付随する双加群複体の完全性が必要である
\cite[Theorem~6.2とその証明直後のtwisted版]{BSWPotentials}。
従ってCY規格化を選べるchartでは、比較すべき表示は
\begin{equation}\label{eq:cubic-p3-potential-degrees}
 C=\mathcal D(\omega_4,2),\qquad
 B=\mathcal D(\omega_5,3),\qquad A=B/(G)
\end{equation}
である。$\mathcal D(\omega,r)$は長さ$r$の部分語による導分を関係とする代数を表す。
これは通常の一階のcyclic derivativeによるJacobi代数とは異なる。
同文献の符号規約では、長さ$m$のsuperpotentialは巡回移動で$(-1)^{m-1}$倍となる。
従って四次では反巡回、五次では巡回である。
一般のAS正則表示ではNakayamaの捩れも記録し、CY性を無条件に仮定しない。

三次側に必要な追加データは、$B_3$内の直線$L=kG$である。
中心三次chartでは
\[
 L\subset\ker\bigl(B_3\longrightarrow
        \operatorname{Hom}(B_1,B_4),\quad g\longmapsto[x\mapsto gx-xg]\bigr)
\]
を課す。全normal問題では中心条件を正規条件に置き換える。
$G$は$B_3$の元であり、自由代数内の三次liftを変更しても同じ直線になり得る。
従ってpotentialの係数空間に任意の三次式を付加するだけではモジュライ問題を表せない。

添字付き表示では、体上で次の選択肢がある。
\begin{enumerate}
\item \textbf{五頂点のambient表示。}
三次イデアルを保つ5周期同型を選び、$i\in\mathbb Z/5$、
五次元の矢空間$V_i:i\to i+1$を取る。
必要なNakayamaの捩れも記録した五次superpotentialの三階導分で二次関係を与え、
さらに長さ三の直線$L_i\subset B_{i,i+3}$と
\begin{equation}\label{eq:cubic-p3-normal-incidence}
 L_iV_{i+3}=V_iL_{i+1}\subset B_{i,i+4}
\end{equation}
を課す。五頂点の向き付けられたcycle上で長さ3のpathは閉じないので、
$L_i$は通常のquiver potentialではなくnormal incidenceのデータである。
\item \textbf{一頂点の次数付き表示。}
系\ref{cor:moduli-cubic-ambient-periodic}の正規超曲面対の周期性定理を、AS index $5$と三次正規族へ適用すると、
$\gcd(5,3)=1$によりイデアルを保つ1周期同型を選べる。
従って五つのloop、五次のtwisted potential、正規三次直線による
\eqref{eq:cubic-p3-potential-degrees}の表示を得る。
この選択は一意ではなく、異なる表示を比較するにはZhang twisting systemと周期の忘却を扱う。
\end{enumerate}
正規族自身の3周期で折り畳む方法もある。その場合は三頂点上で正規三次族が中心的な閉路になる一方、
ambientのNakayama置換は$+5\equiv+2\pmod3$となり、五次potentialに頂点の捩れが残る。
これらは体上の表示であり、族での周期構造や直線束の降下は別に証明する必要がある。

有限tensorの表示を得ても、全てのtensorが求めるAS正則代数を与えるわけではない。
五次のpotentialには導分空間の型$(1,5,10,10,5,1)$とKoszul複体の完全性を課し、
三次直線には正則性と\eqref{eq:cubic-p3-normal-incidence}を課す。
射影三次圏のsmooth proper性も独立の条件である。
従って研究上の自然な有限表示の候補は、\emph{ambientのpotential空間上のnormal incidence}である。

\subsection{例外列の障害と残余圏を保持する模型}

滑らかな可換三次三様体$X\subset\mathbb P^4$では
\[
 D^b(X)=\langle\operatorname{Ku}(X),\mathcal O_X,\mathcal O_X(1)\rangle,
 \qquad S_{\operatorname{Ku}(X)}^3\simeq[5]
\]
である\cite[Corollary~4.1]{KuznetsovFractionalCY}。
従って残余圏のfractional CY次元は$5/3$であり、
非可換$\mathbb P^3$の四対象完全例外列と同じ構造ではない。

\begin{proposition}[普通のtilting代数に対する可換点の障害]
\label{prop:cubic-p3-no-ordinary-tilting}
$X$上の三つの連続する線束$\mathcal O_X,\mathcal O_X(1),\mathcal O_X(2)$は例外列ではない。
さらに$D^b(X)$は完全例外列を持たず、通常の有限次元代数のperfect導来圏とも同値でない。
特に、四頂点の通常のtilting代数によりこの可換点を表すことはできない。
\end{proposition}
\begin{proof}
$\omega_X=\mathcal O_X(-2)$なのでSerre双対により
$\operatorname{Ext}^3(\mathcal O_X(2),\mathcal O_X)=k$である。
従って三つ目の線束を入れた時点で半直交性が破れる。

Hochschild homologyの次元は体拡大で不変なので、有限生成部分体へ降下して
$\mathbb C$へ埋め込むことにより、以下は$k=\mathbb C$で計算してよい。
次に$H$を超平面類とすると
$c(T_X)=(1+H)^5/(1+3H)$より$c_3(T_X)=-2H^3$、$\chi_{\mathrm{top}}(X)=-6$である。
Lefschetz定理から$b_0=b_2=b_4=b_6=1$、$b_1=b_5=0$となるので$b_3=10$である。
$H^0(X,\omega_X)=0$とHodge分解から$h^{2,1}=h^{1,2}=5$を得る。
HKR同型と半直交分解に関するHochschild homologyの加法性
\cite[Theorem~8.3, Corollary~7.5]{KuznetsovHochschild}により
\[
 HH_1(X)\simeq k^5,\qquad HH_{-1}(X)\simeq k^5.
\]
例外対象一つのHochschild homologyは次数0にしかないため、完全例外列は存在しない。
また通常の代数$\Lambda$のbar複体は非負homological次数にあり、$HH_{-1}(\Lambda)=0$である。
Hochschild homologyの導来Morita不変性から
$D^b(X)\simeq\operatorname{Perf}(\Lambda)$も不可能である。
実際、そのような三角圏同値があれば$\Lambda$の像は自己Extが次数0に集中する
生成対象となり、そのdg endomorphism代数は通常の$\Lambda$と擬同型の鎖で結ばれるため、
Hochschild homologyの比較に必要な導来Morita同値が得られる。
\end{proof}

この障害は有限頂点の\emph{dgまたは$A_\infty$模型}を排除しない。
例えば$T=\bigoplus_{i=0}^{4}\mathcal O_X(i)$はcompact generatorである。
実際、埋込み$j:X\hookrightarrow\mathbb P^4$と$M\in D_{\mathrm{qc}}(X)$について
$\operatorname{RHom}(T,M)=0$ならBeilinson生成と随伴から$j_*M=0$、従って$M=0$となる。
よって$D^b(X)\simeq\operatorname{Perf}(\operatorname{RHom}_X(T,T))$であるが、
このdg代数には上で検出した高次Extがあり、その次数0部分だけでは足りない。

非可換商$A=B/(G)$についても、Orlovの定理は適切なmutationの後に
二つの線対象と$D^{\mathrm{gr}}_{\mathrm{sg}}(A)$による分解を与える。
さらに$B$がnoetherian AS正則、$G$が正則正規なら、後者は
twisted matrix factorizationのhomotopy圏と同値になる
\cite[公刊版, Theorem~5.8]{CSBGCCKM}。
同文献は左加群を使うので、本稿の右加群の規約では反対環へ適用する。
従って圏のモジュライまで比較する際には、ambientのpotentialと三次直線に加え、
この残余圏と二つの線対象への貼り合わせを保持するのが適切である。
一般正規商では正規自己同型の捩れが残るため、
可換点の式$S^3=[5]$をそのまま一般の非可換商へ移すことはしない。

\section{対角混合族の完全な同型条件と反対代数}
\label{sec:pair-diagonal-classification}
$D=x_1^3+x_2^3+x_3^3$、$T$を六つの混合語の和とし、
$s=(u,v)^t$, $q_i\in k[u,v]_2$, $f\in k[u,v]_3$に対して
\begin{gather*}
 B(h,q)=\frac{k[u,v]\langle x_1,x_2,x_3\rangle}
 {([x_j,x_k]-h(x_i^2+q_i))_{(i,j,k)\text{ cyclic}}},\\
 W_{h,q,f}=D/3-h^2T/9+\sum_iq_ix_i+f,
 \quad A(h,q,f)=B(h,q)/(W_{h,q,f})
\end{gather*}
と置く。以下$h\ne0$とする。境界の滑らかさは同型判定には不要である。

\begin{lemma}
$Z(A)_1=ku\oplus kv$であり、$\Sym^2(A_1)\to A_2$は同型である。
\end{lemma}
\begin{proof}
二次関係は$\wedge^2A_1$から$\Sym^2A_1$へのgraphで、三次商は次数二を変えない。
従って対称二次式は$A_2$の基底である。
一次元$\sum a_ix_i+bu+cv$の交換子を取ると、独立な$x_i^2$の係数と$h\ne0$から
$a_i=0$を得る。
\end{proof}
\begin{theorem}[対角族の同型の必要十分条件]
\label{thm:pair-diagonal-isomorphism}
$h,h'\ne0$について$A(h',q',f')\simeq_{\gr} A(h,q,f)$であることは、
$p\in S_3$, $N\in\GL_2$, $\lambda_i\ne0$, $\kappa\ne0$が存在して
\begin{gather*}
 \lambda_1^3=\lambda_2^3=\lambda_3^3=\kappa,\qquad
 h'=\operatorname{sgn}(p)h\lambda_1\lambda_2\lambda_3/\kappa,\\
 q_i'(Ns)=\lambda_i^2q_{p(i)}(s),\qquad f'(Ns)=\kappa f(s)
\end{gather*}
を満たすことと同値である。全ての同型は
$s'\mapsto Ns$, $x_i'\mapsto\lambda_i x_{p(i)}$であり、特に$h^6$は不変量である。
\end{theorem}
\begin{proof}
二次包絡の復元により、同型は周囲へ持ち上がる。
中心一次部分から$s'\mapsto Ns$, $x'\mapsto Mx+\ell(s)$であり、
関係式を対称二次式として比較すると
\[
 h\sum_a\operatorname{cof}(M)_{ia}(x_a^2+q_a(s))
 =h'\bigl((Mx)_i+\ell_i(s)\bigr)^2+h'q_i'(Ns).
\]
$x_as_b$の係数から$\ell_i=0$、$x_ax_b$ ($a\ne b$)の係数から
$M_{ia}M_{ib}=0$となる。可逆性により$M$はmonomialである。
$x_{p(i)}^2$の係数を比較して
$h'\lambda_i^2=\operatorname{sgn}(p)h\lambda_j\lambda_k$を得て、
$\lambda_i$を掛ければ三つの立方は等しい。
中央二次部分の比較が$q_i$の条件を与える。
$L=\prod\lambda_i$と置けば$L^3=\kappa^3$, $(h')^2L=\kappa h^2$なので、
$W'$の像は$\kappa(D/3-h^2T/9+\sum q_ix_i)+f'(Ns)$となる。
最高active三次部分の比較から比例係数は$\kappa$であり、$f$の条件を得る。
この最高次元は$h^6=-27$でも非零である。
実際、三変数代数の可換化は$k[x,y,z]/(x^2,y^2,z^2)$で、
$\Omega_h$の像は$-2h^2xyz/3\ne0$である。
逆に表示条件を満たす可逆線形写像は全ての関係を保つので同型を与える。
$h'$の式の六乗が最後の不変量を与える。
\end{proof}

\begin{proposition}[軌道次元と自己同型]
$q_1q_2q_3$の六零点が相異なる範囲で、周囲の対角族の同型類は6次元、
中心三次対の同型類は10次元である。全37次元周囲族と41次元対の中で、
それぞれの$\PGL_5$飽和像の余次元は7となる。
さらに六点集合の$\PGL_2$ stabilizerが自明で$f\ne0$なら
$\Aut_{\gr}A=k^\times$、$A\not\simeq_{\gr} A^{\rm op}$である。
\end{proposition}
\begin{proof}
連続な変換パラメータは$(c,N)\in\mathbb G_m\times\GL_2$の5次元であり、
全生成元の同時スカラーという1次元核を持つ。
六点のstabilizerは有限で、各候補の射影$N$を固定すると$q_i$の条件は
$N$のスカラーと$c$の比を有限個に制限する。
従って軌道は4次元で、周囲の10係数・対の14係数からそれぞれ4を引く。
一般の$\PGL_5$軌道は24次元なので飽和像は30次元・34次元であり、余次元7を得る。
stabilizerが自明なら$N=tI$, $p=1$で、$\lambda_i^2=t^2$と立方の条件から
全ての$\lambda_i$は同じ$\lambda=\pm t$である。
$f\ne0$が$\lambda^3=t^3$を強制し、自己同型は同時スカラーのみとなる。
反対代数は$A(-h,q,f)$であるが、同じ条件は$h'=h$を強制して矛盾する。
\end{proof}
\begin{corollary}[有理係数の非自己反対例と標識付き導来圏]
\label{cor:pair-not-opposite}
$q_1=uv$, $q_2=(u-v)(u-2v)$, $q_3=(u-4v)(u-9v)$とする。
六点集合のstabilizerは自明である。
$h=\pm1$について、両方のtailsが正則となる$f$の非空開集合が存在し、
その中の$f\ne0$では$A$と$A^{\rm op}$の間に構造対象をline objectのshiftへ送る
導来同値は存在しない。有理係数のそのような$f$はZariski稠密に存在する。
\end{corollary}
\begin{proof}
$\{0,\infty,1,2,4,9\}$の最初の三点の像を選ぶ120通りの有理射影行列を列挙し、
全六点を保つものを照合すると恒等変換だけである。
\texttt{verify\_central\_pair\_orbits.py}はこの厳密列挙を行う。
$h=\pm1$の境界の位数は補題\ref{CSBGlem:order-six}により6であり、
定理\ref{CSBGthm:global-open}を両周囲に適用する。
全ての$f$で相対PBW・中心正則商の自由性があり、最高次中心元は非零正則なので全仮定が揃う。
二つの非空開集合は交わり、$f=0$を除いても非空である。
Serre関手・偶数Veroneseによる復元定理と前命題から導来同値の不存在を得る。
有理点の稠密性は$\mathbb Q^4\subset\mathbb C^4$のZariski稠密性による。
構造対象を忘れた全ての導来同値についての主張ではない。
\end{proof}

\section{本質的混合項の量子中心と相対Picardの制約}
\label{sec:mixed-quantum-center}
分離型を越える量子化では中心変数一次の補正が最初の課題となる。
以下の中心元表示はEtingof--Ginzburg \cite[\S9]{EG} のRains型計算および
第\ref{sec:pair-explicit-centers}節に証明した全混合中心式のFermat規格化であり、新しい中心公式とは主張しない。
第\ref{thm:pair-full-fermat-center}定理の自由代数恒等式は、
原点近傍だけでなく全有限パラメータの中心性を証明する。

\subsection{全15混合方向の中心三次元}
$f=(x^3+y^3+z^3)/3$、$U=\langle u,v\rangle$とし、
$A,B,C\in U$、$D,E,H\in\Sym^2U$、$g\in\Sym^3U$を任意に取る。
\begin{equation}\label{eq:full-mixed-fermat}
 F=f+Axy+Byz+Czx+Dx+Ey+Hz+g
\end{equation}
とおき、$u,v$を中心変数として
\begin{equation}\label{eq:full-mixed-brackets}
 [x,y]=hF_z,\quad[y,z]=hF_x,\quad[z,x]=hF_y
\end{equation}
の右辺を完全対称化して読む。$T=xyz+xzy+yxz+yzx+zxy+zyx$、
$\operatorname{sym}(xy)=(xy+yx)/2$とする。

\begin{remark}[Fermat規格化の全混合中心元]\label{rem:full-mixed-center}
定理\ref{thm:pair-full-fermat-center}の中心元は、上の係数では
\begin{align}\label{eq:full-mixed-center}
\Omega_h={}&\frac{x^3+y^3+z^3}{3}-\frac{h^2}{9}T\notag\\
&+A\left(\operatorname{sym}(xy)+\frac{h^4}{18}z^2\right)
 +B\left(\operatorname{sym}(yz)+\frac{h^4}{18}x^2\right)
 +C\left(\operatorname{sym}(zx)+\frac{h^4}{18}y^2\right)\notag\\
&+\left(D+\frac{h^2B^2}{6}-\frac{h^4AC}{18}\right)x
 +\left(E+\frac{h^2C^2}{6}-\frac{h^4AB}{18}\right)y\notag\\
&+\left(H+\frac{h^2A^2}{6}-\frac{h^4BC}{18}\right)z+g(u,v)
\end{align}
となり、$h=0$で$F$に特殊化する。
この式は第\ref{thm:pair-full-fermat-center}定理の$e_1=B,e_2=C,e_3=A$への代入であり、
そこに全係数を保持する関係イデアル恒等式を収録した。
四次補正を省けない反例も同節にある。
\end{remark}

\subsection{一次中心持上げの一意性}
\begin{proposition}\label{prop:mixed-center-jet-unique}
$S_h=\CC\langle x,y,z\rangle/
([x,y]-hz^2,[y,z]-hx^2,[z,x]-hy^2)$とし、
$Z(S_h)_2=0$を仮定する。
$R=\CC[u,v]/(u,v)^2$上の\eqref{eq:full-mixed-brackets}の
量子化族$\mathcal B_R$が次数ごとに$R$-自由なら、
固定された$S_h$の次数三の中心元$\Omega_{h,0}$への
次数三の中心持上げ$\Omega_R\in Z(\mathcal B_R)_3$は、
存在すれば一意である。
\end{proposition}
\begin{proof}
二つの持上げの差$\Delta$は$(u,v)(\mathcal B_R)_3$に属し、
次数別自由性と平方零条件により
$\Delta=uQ+vR'$、$Q,R'\in(S_h)_2$と一意に書ける。
$\Delta$の中心性を各生成元$a=x,y,z$との交換子に適用すると
$u[Q,a]+v[R',a]=0$を得る。
$u,v$が独立な平方零係数なので、
$Q,R'\in Z(S_h)_2=0$。従って$\Delta=0$。
\end{proof}
\begin{remark}[一次混合6方向の明示的持上げ]
\eqref{eq:full-mixed-center}を$(u,v)^2=0$に制限すると
\[
\begin{aligned}
 \Omega_h^{(1)}
 ={}&\frac{x^3+y^3+z^3}{3}-\frac{h^2}{9}T
 +A\left(\operatorname{sym}(xy)+\frac{h^4}{18}z^2\right)\\
 &+B\left(\operatorname{sym}(yz)+\frac{h^4}{18}x^2\right)
 +C\left(\operatorname{sym}(zx)+\frac{h^4}{18}y^2\right).
\end{aligned}
\]
これが全ての一次混合6方向の持上げの表示である。
命題\ref{prop:mixed-center-jet-unique}の一意性自体は、
式\eqref{eq:full-mixed-center}を仮定しない独立した証明である。
$h=0$では$Z(S_0)_2\ne0$なので適用しない。
\end{remark}

\subsection{相対Picardに対するスカラー方向の障害}
\begin{proposition}\label{prop:relative-scalar-picard}
$X\subset\PP^4$を滑らかなcubic threefoldとし、
$T_0\subset\operatorname{Gr}(2,H^0(\PP^4,\OO(1)))$を、
指定する平面$\Lambda_U=V(U)$との切断$C_U=X\cap\Lambda_U$が
滑らかな楕円曲線になる非空開部分とする。
$\mathscr C\to T_0$を楕円曲線族、
$\mathscr P=\operatorname{Pic}^0(\mathscr C/T_0)$を相対Picardとする。
非零Poisson錐の$\mathbb G_m$-束
$\mathscr L^\times\to T_0$から$\mathscr P$への
$T_0$-上の任意の有理写像は、稠密開部分において
$\mathscr L^\times\to T_0$を通じて有理的に因子化する。
すなわちPoissonスカラー方向には定数である。
幾何学的一般ファイバーで連結な有限\emph{étale}被覆を許しても同様である。
\end{proposition}
\begin{proof}
$K=\CC(T_0)$とする。$\mathscr L^\times$の一般ファイバーは
$\mathbb G_{m,K}$である。
$\overline K$へ移すと、与えられた有理写像は
$\mathbb G_m\dashrightarrow E$、$E$は種数1曲線となる。
固有性によって$\PP^1\to E$に延長されるが、
Riemann--Hurwitz公式から定数である。
従ってもとの写像は一般ファイバーに沿って定数であり、
$T_0$上の有理sectionを通じて因子化する。
また代数閉体上の$\mathbb G_m$の連結有限étale被覆は
$t\mapsto t^n$型で、その源も$\mathbb G_m$である。
同じ議論が適用できる。
\end{proof}
\begin{remark}[復元への含意と限界]
元の$X$とPoisson pencil $U$を幾何側でも固定したまま、
非零スカラーを楕円Picardの非自明な点へ送る大域有理模型は存在しない。
しかし量子化後に別の可換cubic $Y$や平面が動く比較、
分岐被覆を介した多価対応、形式局所同値の不存在は結論しない。
一般の非零パラメータでのrotation付き可換復元に必要な比較は未完成である。
有限Artin底での残余圏の構成と可換三次としての形式認識は、
命題\ref{prop:abel-artin-actual-residual}と
命題\ref{prop:kuznetsov-formal-recognition}で証明する。
これらは、rotationや追加の線対象を保つ相対MF比較を主張するものではない。
\end{remark}

% Part of the present nc cubic threefold moduli manuscript.
% Restored from the 2026-11-01 central-pair paper V, sections 19--21.
\section{全混合族の低次数中心、偶数Veroneseと剛体ambient}
\label{sec:moduli-rigid-ambient}
本節では、可換点の全Artin比較とは別に、有限非可換パラメータ
において一つのambientを固定したまま三次方程式を動かす問題を調べる。
全混合中心三次商の低次数中心と偶数 Veronese による復元を示し、
明示周囲 $B_*$ の自己同型を本節の座標変換計算で決定する。
境界の位数六は補題\ref{CSBGlem:order-six}、
正則な中心商の非空開集合は定理\ref{CSBGthm:global-open}と
系\ref{rem:moduli-rigid-generic-smoothness}で扱う。
本節の導来同値の比較には、定理\ref{thm:moduli-marked-derived}に
明示した tails の有限大域次元を仮定する。

\subsection{三次商による二次中心の非増大}
$U=\Span(u,v)$、$R=k[u,v]$、$k=\C$とおく。
$F\in k[x,y,z,u,v]_3$は
$F(x,y,z,0,0)$が滑らかな平面三次式$f_3(x,y,z)$であると仮定する。
$h\in k^\times$に対して
\[
 B(h,F)=\frac{R\langle x,y,z\rangle}
 {([x,y]-h\Symm F_z,\ [y,z]-h\Symm F_x,\
    [z,x]-h\Symm F_y)}
\]
とし、次数三の中心正則元$W$で割った
$A=B/(W)$を考える。ここで$B$は一次関係を持たないとする。

\begin{lemma}[中心三次商では二次中心が増えない]
\label{lem:moduli-no-new-quadratic-center}
上の仮定で、自然な同一視$B_2\simeq A_2$の下で
$Z(A)_2=Z(B)_2$となる。
\end{lemma}
\begin{proof}
$Q\in B_2$の像が$A$で中心なら、
$[Q,x_i]=c_iW$、$c_i\in k$、
$(x_1,x_2,x_3)=(x,y,z)$と書ける。
$u,v,W$が中心なので
\[
 0=[Q,[x_i,x_j]]
 =h[Q,\Symm F_{x_k}]
 =hW\sum_\ell c_\ell F_{x_kx_\ell}
\]
である（$i,j,k$は巡回）。
$hW$を消去すると、右側の一次式が$B_1$で零。
一次関係を持たないから可換多項式として零であり、
$u=v=0$で
\[
  \sum_\ell c_\ell(f_3)_{x_kx_\ell}=0
  \quad(k=1,2,3).
\]
すなわち$\sum_\ell c_\ell(f_3)_{x_\ell}$は
定数の斉次二次式なので零になる。
$c\ne0$なら座標を取り換えれば$f_3$は一変数を含まず、
対応する射影点で偏微分がすべて消えて滑らかさに矛盾する。
従って$c=0$で、$Q$は既に$B$の中心である。
\end{proof}

\begin{lemma}[全混合族の次数一・二の中心]
\label{lem:moduli-low-center}
$S=B/(u,v)$とし、中心三次元$\gamma\in S_3$で
\[
 S/(\gamma)\simeq B(E,\mathcal L,\tau)
\]
となるとする。$E$は滑らかな楕円曲線、
$\deg\mathcal L=3$、$\tau$は平行移動で$\tau^2\ne1$と仮定する。
さらに相対PBWにより
$\gr_{x,y,z} B\simeq S\otimes_k k[u,v]$とする。
このとき
\[
 Z(B)_1=Z(A)_1=U,\qquad
 Z(B)_2=Z(A)_2=\Sym^2U.
\]
\end{lemma}
\begin{proof}
twisted homogeneous coordinate ringは
$k(E)[t,t^{-1};\tau]$に埋め込まれる。
そこに次数$m=1,2$の中心元$at^m$があれば
一次切断$s_1,s_2$との可換性から
$\tau^m(s_1/s_2)=s_1/s_2$を得る。
次数三の完全線形系の一次切断の比は$k(E)$を生成するため
$\tau^m=1$となり仮定に反する。
従って$S/(\gamma)$には低次数中心がなく、
$S\to S/(\gamma)$は次数一・二で単射だから
$Z(S)_1=Z(S)_2=0$。
$B$の中心元の最高相対次数を取ると、
正の相対次数一・二は消え、中心は$k[u,v]$に属する。
三次商でも次数一の関係は増えず、
二次は補題\ref{lem:moduli-no-new-quadratic-center}を使う。
\end{proof}

\subsection{偶数Veroneseからの標識付き代数の回復}
\begin{lemma}[中心二変数の平方根と奇数部分]
\label{lem:moduli-even-veronese}
$A$はconnectedな次数付き$k$代数で、
二つの一次中心元$u,v$が多項式環$S=k[u,v]$を生成し、
$A$が$S$上自由、
$Z(A)_2=S_2$とする。
$R=A^{(2)}=\bigoplus_{n\ge0}A_{2n}$とおくと、
$R$の次数付き同型類は$A$の次数付き同型類を決定する。
詳しくは、同じ仮定を満たす$A'$について、
任意の$\phi:R\xrightarrow{\sim}(A')^{(2)}$は
正確に二つの次数付き同型$A\to A'$へ持ち上がり、
その差は奇数次数に$-1$を掛ける自己同型である。
\end{lemma}
\begin{proof}
$z\in Z(A)_1$について$uz,vz\in S_2$。
$u(vz)=v(uz)$と$S$の一意分解性から
$uz=u\ell$、$\ell\in S_1$であり、
$u$の正則性から$z=\ell$となる。
$z\in Z(R)$と奇数次数の$a\in A$に対しては
$ua\in R$なので
$u[z,a]=[z,ua]=0$より$[z,a]=0$。
従って$Z(R)=Z(A)^{(2)}$。
その次数一中心の生成する内在的部分環は
\[
 D=k[u^2,uv,v^2]
   \simeq k[\alpha,\beta,\gamma]/(\alpha\gamma-\beta^2).
\]
$\Proj D$は二次曲線で$\OO_{\Proj D}(1)=\OO_{\PP^1}(2)$。
任意の$D\simeq D'$は二次Veronese表現
$\GL_2\to\Aut_{\gr}D$により二通りだけ
$S\simeq S'$へ持ち上がり、その差は全一次元の符号反転。

$u,v$を局所化して$R$の次数を二倍すると
\[
 A[u^{-1}]
  \simeq R[(u^2)^{-1}]
    \otimes_{D[(u^2)^{-1}]}S[u^{-1}]
\]
であり、$v$の場合も同様である。
$A$は$S$上自由なので
\[
 A[u^{-1}]\cap A[v^{-1}]=A
 \quad\text{in }A[(uv)^{-1}].
\]
これが$R$と$S$だけからの復元式である。
同型を持ち上げた後、偶数次数と$S_1$上で恒等なら、
奇数元$a$について$u\phi(a)=\phi(ua)=ua$から
$\phi(a)=a$。残る一つは全奇数次数の符号反転である。
\end{proof}

\begin{theorem}[構造対象を保つ導来同値：条件付き全混合復元]
\label{thm:moduli-marked-derived}
$A=B/(W)$と$A'=B'/(W')$は上の全混合族の中心正則三次商で、
補題\ref{lem:moduli-low-center}の仮定と
それぞれの$k[u,v]$上の自由性を満たすとする。
両方の$\qgr$は有限大域次元を持つと仮定する。
次は同値である：
\begin{enumerate}
\item $A\simeq_{\gr}A'$;
\item $\OO_A$を$\OO_{A'}$に送る$k$線形Abel圏同値
 $\qgr A\simeq\qgr A'$が存在する;
\item ある$m,r\in\mathbb Z$について
 $\OO_A$を$\OO_{A'}(m)[r]$に送る$k$線形三角同値
 $D^b(\qgr A)\simeq D^b(\qgr A')$が存在する。
\end{enumerate}
これらの条件は、周囲の二次代数と三次中心直線
$(B,kW)$の同型類も決定する。
\end{theorem}
\begin{proof}
$(1)\Rightarrow(2)\Rightarrow(3)$は標準的である。
$(3)$の同値にshiftとtwistを合成して構造対象を一致させる。
$B$はCY五次元、$W$は中心正則三次元だから
balanced dualizing complexは$A(-2)[4]$である。
tailsの有限大域次元を使ってSerre関手は
$\mathsf S_A\simeq(-2)[3]$となる。
これは\href{https://www.math.bgu.ac.il/~amyekut/publications/nc-serre/paper.pdf}
{Yekutieli--ZhangのSerre双対}の適用を入力とする。
同値はSerre関手と可換なので
$\mathsf H_A=\mathsf S_A^{-1}[3]\simeq(2)$を保ち、
その軌道環
\[
 \bigoplus_{n\ge0}\Hom(\OO_A,\mathsf H_A^n\OO_A)
 \simeq A^{(2)}
\]
を積とともに同定する。
補題\ref{lem:moduli-low-center}と
\ref{lem:moduli-even-veronese}により$A\simeq A'$。
最後に二次関係の核から$B$を、三次の核から$kW$を回復する
定理\ref{thm:moduli-iso}を使う。
\end{proof}
\begin{remark}
\label{rem:moduli-marking}
ここで除いたのは次数一shiftの\emph{指定}であり、
構造対象の条件自体は残る。任意の導来自己同値が
構造対象をline objectのshiftへ送るとは主張しない。
有限大域次元を持たないtailsへの適用も主張しない。
\end{remark}

\subsection{全混合ポテンシャルの有限同型判定}
\begin{proposition}[中心blockと三次ポテンシャルの共変性]
\label{prop:moduli-mixed-block-isomorphism}
上の低次数中心を持つ$B(h,F)$、$B(h',F')$について
任意の次数付き同型$B(h',F')\to B(h,F)$は
\[
 s' \longmapsto Ns,\qquad x'\longmapsto Lx+Ms,
 \quad
 L\in\GL_3(k),\ N\in\GL_2(k),\ M\in\mathrm{Mat}_{3\times2}(k)
\]
の形を持つ。これは同型である必要十分条件が
ある$r\in k[u,v]_3$について
\[
 h'F'(Lx+Ms,Ns)=h\det(L)F(x,s)+r(s)
\]
となることに等しい。
中心商の次数付き同型にはさらに
$\psi_{L,M,N}(W')=\kappa W$ ($\kappa\in k^\times$)
を要し、これで必要十分である。
\end{proposition}
\begin{proof}
低次数中心が$U=\Span(u,v)$なので同型はblock形となる。
$([y,z],[z,x],[x,y])^t$の変換は
$\det(L)L^{-t}$であり、関係式の保存条件は
\[
 h'\nabla_{x'}F'(Lx+Ms,Ns)
 =h\det(L)L^{-t}\nabla_xF(x,s).
\]
二次関係空間が交代二次式から対称二次式へのgraphであるため、
これは可換二次多項式の恒等式と同値。
chain ruleで勾配の式を積分すると、差は$k[u,v]_3$に入る。
最後の中心直線保存は二次包絡と三次核の回復に等しい。
\end{proof}

\subsection{射影自己同型のない明示ambient}
次の例は元の中心三次対原稿第21節にあったものであり、
正則な中心商の開集合の存在は、中心商論文の有限位数境界に関する定理から従う。

\begin{proposition}[明示的な剛体ambient：中心の入力を指定]
\label{prop:moduli-rigid-example}
\[
 F_*=\frac{x^3+y^3+z^3}{3}
     +uxy+vyz+(u+v)zx
\]
に対する$h=1$の代数
\[
 B_*=\frac{k[u,v]\langle x,y,z\rangle}
 {([x,y]-z^2-vy-(u+v)x,\
   [y,z]-x^2-uy-(u+v)z,\
   [z,x]-y^2-ux-vz)}
\]
を考える。$S_*=B_* /(u,v)$の中心三次商が、
平面三次$E_*:x^3+y^3+z^3+xyz=0$に付随する
位数6の平行移動を持つTHCRであることを入力とする。
このとき
\[
 Z(B_*)_1=ku\oplus kv,\qquad
 Z(B_*)_3=kG_*\oplus k[u,v]_3,
 \qquad \Aut_{\gr}(B_*)=k^\times.
\]
ここで
\[
\begin{aligned}
 \Omega_1&=\frac{x^3+y^3+z^3}{3}
       -\frac19\sum_{\sigma\in S_3}\sigma(xyz),\\
 G_*={}&\Omega_1+
       u\Symm(xy)+v\Symm(yz)+(u+v)\Symm(zx)\\
 &+\frac1{18}\bigl(vx^2+(u+v)y^2+uz^2\bigr)\\
 &+\frac1{18}\bigl((-u^2-uv+3v^2)x
       +(3u^2+5uv+3v^2)y+(3u^2-uv-v^2)z\bigr).
\end{aligned}
\]
従って$B_*/(G_*+f)$ ($f\in k[u,v]_3$)は、
一つの周囲の中で相互に次数付き非同型な四次元代数族を与える。
さらに、$f$の非空Zariski開部分ではtailsが正則で大域次元3となる。
\end{proposition}
\begin{proof}
Fermat型の混合中心の普遍公式
\href{https://github.com/uedakazushi/nc-cubic/blob/main/papers/support/papers/sections/fermat_mixed_universal_center.tex}
{（nc-cubic論文A）}へ
$A=u,B=v,C=u+v,D=E=H=g=0,h=1$を代入すれば
上記の$G_*$とその中心性が得られる。
指定されたTHCRでは平行移動の位数が6なので
次数一・二の中心は零、次数三では$k\Omega_1$である。
相対PBWと中心元の三次還元から
$Z(B_*)_1=U$、$Z(B_*)_3=kG_*\oplus k[u,v]_3$となる。

任意の次数付き自己同型は中心一次空間を保つから、
$s=(u,v)^t$、$x=(x,y,z)^t$について
$s\mapsto Ns,\ x\mapsto Mx+\ell(s)$の形である。
命題\ref{prop:moduli-mixed-block-isomorphism}の勾配比較より
\[
 F_*(Mx+\ell(s),Ns)=\det(M)F_*(x,s)+r(s)
\]
となる。最高$x$-次数で
$\sum_i(Mx)_i^3=\det(M)\sum_ix_i^3$。
Fermat三次のHessianは非零定数倍の$xyz$なので、
$M$は座標三直線の合併を保ち、monomialである：
$(Mx)_i=\lambda_i x_{p(i)}$。
$x_{p(i)}^2s$項を比較すると$\ell=0$。
$\lambda_i^3=\det(M)$と
$\det(M)=\operatorname{sgn}(p)\prod_i\lambda_i$から
$p$は偶置換である。
同時スカラー倍で$\lambda_1=1$に規格化すると
$\lambda_i^3=1$、$\lambda_1\lambda_2\lambda_3=1$。
混合係数を
$(r_1,r_2,r_3)=(v,u+v,u)$とすれば
$r_i(Ns)=\lambda_i r_{p(i)}(s)$。
これらの一次式の唯一の関係の係数は
$c=(-1,1,-1)$である。よって
$\lambda_i=\mu c_{p(i)}/c_i$。
$p$が三cycleなら三つの比に$+1$と$-1$が混じり、
$\lambda_i^3=1$に矛盾する。
従って$p=1$、$\lambda_i=1$、$N=1$となり、
自己同型は同時スカラー倍だけ。
最後に次数二包絡・三次核の復元により、
$B_*/(G_*+f)\simeq_{\gr}B_*/(G_*+f')$なら
$[G_*+f]=[G_*+f']$。
$G_*$の係数を1と規格化したので$f=f'$である。
\end{proof}

\begin{corollary}[剛体周囲の中の正則な四係数非同型族]
\label{rem:moduli-rigid-generic-smoothness}
ある非零多項式$P(f)$が存在し、$P(f)\ne0$では
$\gldim\qgr B_*/(G_*+f)=3$である。
同時に可換な$V(F_*+f)\subset\PP^4$も滑らかとなる$f$は非空開部分をなし、
その中の有理係数点はZariski稠密である。
\end{corollary}
\begin{proof}
周囲はnoetherian Koszul CY5、相対associated gradedは$S_1\otimes k[u,v]$である。
第\ref{sec:pair-explicit-centers}節の中心式を代入すると$G_*$を得て、
全ての$f$で商は$k[u,v]$上自由である。
最高次中心元は非零正則、境界は位数6の滑らかな楕円THCRなので、
定理\ref{CSBGthm:global-open}の全仮定を満たす。
そこでは全係数chartの代数が$D_u[b,c,d]$となること、
中心持上げにより可換有限型中心上有限となることを確認してからBertiniを適用する。

可換な滑らかさについては$u=v=0$上に勾配の共通射影零点はない。
$u=1,t=v/u$上のactive勾配は
$x^2+y+(1+t)z$, $y^2+x+tz$, $z^2+(1+t)x+ty$であり、
$x,y,z$全次数を先に比較すると互いに素な先頭項$x^2,y^2,z^2$を持つ
monic Gr\"obner基底である。従って臨界スキームは$k[t]$上rank8の有限自由環を持ち、次元1である。
各点が全三様体の特異点となる条件は$(F_*+f)(P)=\partial_t(F_*+f)(P)=0$で、
$f(1,t)$の値と微分は四係数に対する独立な二条件である。
従って特異点incidenceの次元は高々3である。$v=1$も同様。
全incidenceは$\PP^4\times\A^4$内で閉で、その係数空間への像はproper性により閉である。
ゆえに滑らかな係数は非空開をなし、tails正則開部分とも交わる。
有理点の稠密性から最後の主張が従う。
これは可換三様体との残余圏同値を主張するものではない。
\end{proof}

\section{有限位数楕円境界を持つ四係数族の中心上有限な相対模型}
\label{sec:four-parameter-central-model}

この節では、固定した周囲の代数 $B$ における
中心三次斉次元 $G$ の商だけを動かす。
係数環と族を
\[
 P=\mathbb C[a,b,c,d],\qquad
 R=P[u,v],\qquad
 \mathcal A=B\otimes_{\mathbb C}P/
 (G+au^3+bu^2v+cuv^2+dv^3)
\]
と書く。$a,b,c,d$ の通常の次数は零、$u,v,x,y,z$ の通常の次数は一とする。
以下で用いる入力は、既存の PBW 構成により
\begin{equation}\label{eq:central-model-filtered-input}
 \operatorname{gr}_{x,y,z}\mathcal A
 \simeq H\otimes_{\mathbb C}P[u,v],\qquad
 H=B(E,\mathcal L,\tau)
\end{equation}
となり、$E$ が滑らかな楕円曲線、$\mathcal L$ が豊富な直線束、
$\tau$ が有限位数の平行移動であることである。
各 $F_j\mathcal A$ は $R$ 上有限自由で、閉点への基底変更と filtration が両立する。
また $\mathcal A$ は通常の次数一で生成され、$\mathcal A_0=P$ とする。
具体的な $B_*$ については、命題\ref{prop:moduli-rigid-example}の
PBW条件と楕円境界の同定を入力とする。
境界の平行移動の位数六は補題\ref{CSBGlem:order-six}で確認される。

式 \eqref{eq:central-model-filtered-input} の各斉次部分は $R$ 上自由であるから、
その基底を filtration に沿って持ち上げると
\begin{equation}\label{eq:central-model-R-free}
 \mathcal A\text{ は次数付き }P[u,v]\text{-加群として自由}
\end{equation}
である。
特に任意の可換 $P$-代数 $S$ について、
$\mathcal A_S=\mathcal A\otimes_P S$ は $S[u,v]$ 上自由であり、
$u,v$ およびその正冪は普遍的な中心正則列となる。

\subsection{中心元の一様な持上げ}

補題\ref{CSBGlem:global-center-lifting}の中心持上げを，
$x,y,z$に関するfiltrationへ適用する。

\begin{theorem}[全四係数族の有限な中心模型]
\label{thm:central-model-finite}
上述の入力の下で、有限型次数付き可換 $P$-代数
\[
 P[u,v]\subset C\subset Z(\mathcal A),\qquad C_0=P
\]
であって $\mathcal A$ が $C$ 上有限となるものが存在する。
任意の noetherian 可換 $P$-代数 $S$ に対し、
$\mathcal A_S$ は $C_S=C\otimes_P S$ 上有限であり、左右 noetherian である。
\end{theorem}

\begin{proof}
補題\ref{CSBGlem:global-center-lifting}を $R_0=P[u,v]$、$U=\mathcal A$ に適用する。
そこで持ち上げる中心元 $z\in Z(H)_j$ は、通常の次数も $j$ である。
中心に持ち上げた元を通常の次数で分解すると、各斉次成分も中心に入る。
通常の次数 $j$ の成分を取れば、最高 filtration symbol は依然として $z$ である。
従って中心生成元の持上げを通常の次数で斉次に選べる。
$P[u,v]$ と有限個のこれらの持上げが生成する代数を $C$ とする。
同補題の filtration 帰納法は $\mathcal A$ の $C$ 上の有限性を与える。

有限生成加群であることは任意の基底変更後にも保たれる。
$C_S$ は noetherian 環 $S$ 上有限型の可換代数なので noetherian である。
$\mathcal A_S$ の左イデアルまたは右イデアルの昇鎖は $C_S$-部分加群の昇鎖であり、停止する。
\end{proof}

\begin{remark}[全中心および非被約基底の扱い]
\label{rem:central-model-not-full-center}
ここでは $C=Z(\mathcal A)$ を主張していない。
また、任意の $P\to S$ の後に $C_S\to\mathcal A_S$ が単射であるとも主張しない。
必要なのは、中心に像を持つ環準同型
$C_S\to Z(\mathcal A_S)$ と、その作用に関する有限加群性だけである。
中心持上げに Rains の定理と被約性を使うのは $P[u,v]$ 上で一度だけであり、
非被約な $S$ 上でそれらを再適用していない。
\end{remark}

\begin{corollary}[剛体族の普遍的な二次中心]
\label{cor:central-model-universal-center-two}
命題\ref{prop:moduli-rigid-example}の$B_*$とその四係数族について，
任意の可換$P$-代数$S$に対し
\[
 Z(\mathcal A_S)_2=S u^2\oplus S uv\oplus S v^2
\]
である．従ってこの族は定義\ref{def:abel-two-center-family}の
普遍的な中心二変数条件を全係数基底上で満たす．
この主張では$S$のnoetherian性を仮定しない．
\end{corollary}
\begin{proof}
三次商は二次関係を変えないので，$(\mathcal A)_2$は
$(B_*)_2\otimes P$であり，$Pu^2\oplus Puv\oplus Pv^2$は
その定数な直和因子である．従って商は自由rank~12である．
$(\mathcal A)_1$は自由rank~5，$(\mathcal A)_3$は自由rank~34だから，
命題\ref{prop:abel-center-rank-open}の有限自由加群間の写像
\[
 (\mathcal A)_2/\operatorname{Sym}^2(Pu\oplus Pv)
 \longrightarrow\operatorname{Hom}_P((\mathcal A)_1,(\mathcal A)_3)
\]
を使える．各閉点$s$では，三次元$G_*+f_s$の最高symbolは非零正則な$\gamma$であり，
補題\ref{lem:moduli-low-center}と\ref{lem:moduli-no-new-quadratic-center}から
$Z(\mathcal A_s)_2=k(s)u^2\oplus k(s)uv\oplus k(s)v^2$である．
ここで固定した主変数の三次式の滑らかさと境界translationの位数六を使う．
従って全閉点がrank~12の開部分に入る．$P$はJacobson環なのでその開部分は全体である．
写像は局所的に分裂単射であり，任意の$S$への係数変更後も単射である．
全次数の有限自由性と$S[u,v]$-自由性も既に示したので，定義の他の条件も満たす．
\end{proof}

\subsection{射影スキーム上の有限代数層と普遍的な section 復元}

以下、$\ell>0$ を十分割り切れる偶数とし、$D=C^{(\ell)}$ が $D_1$ で生成されるように取る。
\[
 Y=\operatorname{Proj}D,\qquad
 \mathscr A=\widetilde{\mathcal A^{(\ell)}}
\]
と置く。$Y\to\operatorname{Spec}P$ は射影的で、$\mathscr A$ は有限な
$\mathcal O_Y$-代数である。任意の可換 $P$-代数 $S$ について
$Y_S=\operatorname{Proj}(D\otimes_P S)$ と書き、$\mathscr A_S$ を基底変更した代数層とする。
中心準同型に核が生じる場合も、この模型をその support とともに用いる。

\begin{proposition}[相対 tails の有限代数層模型]
\label{prop:central-model-qgr}
任意の可換 $P$-代数 $S$ について、基底変更に自然な同値
\[
 \operatorname{QGr}(\mathcal A_S)
 \simeq \operatorname{QCoh}(Y_S,\mathscr A_S)
\]
がある。$S$ が noetherian なら、これは同値
\[
 \operatorname{qgr}(\mathcal A_S)
 \simeq \operatorname{coh}(Y_S,\mathscr A_S)
\]
に制限される。
構造対象 $\mathcal O=\pi\mathcal A_S$ は $\mathscr A_S$ に対応し、
$\mathcal O(m)$ は
$\widetilde{\mathcal A_S(m)^{(\ell)}}$ に対応する。
大きい加群圏は $S$-線型 Grothendieck 圏であり、
$S$ が noetherian なら locally noetherian である。
これらの同値は任意の基底変更の合成および fppf 降下に両立する。
\end{proposition}

\begin{proof}
標準開 $D_+(c)$、$c\in D_1=C_\ell$ を取る。
$\mathcal A_S$ は次数一生成なので、$c$ を次数一元の長さ $\ell$ の積の和で書ける。
中心局所化 $T=\mathcal A_S[c^{-1}]$ において、その式を最初または最後の因子で切れば
\[
 1\in T_1T_{-1},\qquad 1\in T_{-1}T_1
\]
となる。
従って $T$ は強く $\mathbb Z$-次数付きであり、次数零部分を取る関手は
次数付き $T$-加群と $T_0$-加群の同値を与える。
実際 $T_n$ は可逆な $T_0$-両側加群であり、
$N\mapsto N\otimes_{T_0}T$ が逆関手である。

一方、$\mathcal A_S$ が $C_S$ 上有限なので、$\mathcal A_S$-torsion と
$C_{S,+}$-torsion は一致する。
例えば $mC_{S,\geq n}=0$ なら、有限個の斉次 $C_S$-加群生成元の次数の最大値を $e$ として
$m\mathcal A_{S,\geq n+e}=0$ となる。
逆方向は $C_S$ の像が中心にあることから直ちに従う。
標準開上の上記の同値を貼り合わせると tails の同値を得る。
これは通常の Serre の Proj 対応の証明を有限な代数層に適用したものであり、
中心局所化と次数零部分だけで構成される。
逆向きには、quasi-coherent な $\mathscr A_S$-加群 $N$ から
\[
 \bigoplus_{n\in\mathbb Z}
 \Gamma(Y_S,N\otimes_{\mathcal O_{Y_S}}\mathcal O_{Y_S}(n))
\]
という次数付き $\mathcal A_S^{(\ell)}$-加群を取る。
この加群の層化は $N$ を復元する
\cite[Lemma~28.29.5, Tag~0AG5]{StacksProj}。
さらに次数を $\ell$ 倍にして $\mathcal A_S$ へ誘導すれば、
上記の強次数付き局所化の逆関手を各標準開上で得る。
有限中心被覆上で層化の核が torsion と一致することから、
これらは商圏上で互いに逆である。
ここまでの構成は $S$ の noetherian 性を使わない。
また中心局所化と次数零部分は任意の係数拡大と可換する。
したがって基底変更とその合成に自然である。
coherent 部分の対応には $S$ の noetherian 性を用いる。

有限代数層の quasi-coherent 加群は通常の加群とその作用写像の降下として
fppf 降下を持つ。大きい tails 商は Grothendieck 圏であり、
$S$ が noetherian なら $\mathcal A_S$ も noetherian なので
この商圏は locally noetherian である。
固定した二つの模型の間の同値および自然変換の降下には，
補題\ref{lem:abel-canonical-family-equivalences}を用いる。
\end{proof}

\begin{theorem}[任意基底変更後の全次数 saturation]
\label{thm:central-model-universal-sections}
任意の可換 $P$-代数 $S$ と任意の整数 $m$ に対し、標準写像
\[
 (\mathcal A_S)_m\longrightarrow
 \operatorname{Hom}_{\operatorname{QGr}(\mathcal A_S)}
 (\mathcal O,\mathcal O(m))
\]
は同型である。これらの同型は任意の可換な基底変更、乗法および合成と両立する。
特に負次数の右辺も零である。
非noetherianな$S$については大きい$\operatorname{QGr}$を用いる。
\end{theorem}

\begin{proof}
まず$S$がnoetherianである場合を示す。
$D_S=D\otimes_P S$、$I=D_{S,+}$ とし、次数付き有限 $D_S$-加群
$M_m=\mathcal A_S(m)^{(\ell)}$ を考える。
式 \eqref{eq:central-model-R-free} により、$u^\ell,v^\ell\in (D_S)_1$
は $M_m$ 上の正則列である。
商を取ることと Veronese 部分を取ることは可換なので、
この主張は $\mathcal A_S$ が $S[u,v]$ 上自由であることから従う。

$u^\ell$ の単射性から $H^0_I(M_m)=0$ である。
また $M_m/u^\ell M_m$ 上では $v^\ell$ が単射なので
$H^0_I(M_m/u^\ell M_m)=0$ となる。
local cohomology の長完全列により、$u^\ell$ は $H^1_I(M_m)$ 上でも単射である。
しかし $H^1_I(M_m)$ は $I$-torsion であり、各元は $u^\ell$ のある冪で消える。
従って $H^1_I(M_m)=0$ である。

斉次生成元による増大\v{C}ech複体とその次数成分を比較して得られる，
標準次数付き可換環のProjに対するlocal cohomologyの完全列
\[
 0\longrightarrow H^0_I(M_m)\longrightarrow M_m
 \longrightarrow\bigoplus_{n\in\mathbb Z}
 H^0(Y_S,\widetilde M_m(n))\longrightarrow H^1_I(M_m)
 \longrightarrow0
\]
の次数零部分を取る。
$\widetilde M_m=\mathcal O(m)$、$\mathcal O=\mathscr A_S$ なので、
\[
 H^0(Y_S,\widetilde M_m)
 =\operatorname{Hom}_{\mathscr A_S}(\mathscr A_S,\widetilde M_m)
\]
であり、所望の同型を得る。
ここで用いた写像は、元を左乗法による加群準同型へ送る標準写像そのものである。
この段階では、同じ証明が全てのnoetherianな$S$に適用できる。

一般の$S$を、その有限型$P$-部分代数$S_\lambda$のfiltered colimitとして書く。
$\mathcal D=\operatorname{QGr}(\mathcal A)$とし、$E_m=\mathcal O(m)$と置く。
$\mathcal D$はlocally noetherianであり、各$E_m$は有限表示対象である。
補題\ref{lem:abel-arbitrary-coefficient-tails}により
$\operatorname{QGr}(\mathcal A_S)\simeq\mathcal D_S$であり、
自由係数作用と忘却の随伴から
\begin{align*}
 \operatorname{Hom}_{\mathcal D_S}(S\otimes_P E_0,S\otimes_P E_m)
 &\simeq\operatorname{Hom}_{\mathcal D}(E_0,S\otimes_P E_m)\\
 &\simeq\varinjlim_\lambda
       \operatorname{Hom}_{\mathcal D}(E_0,S_\lambda\otimes_P E_m)\\
 &\simeq\varinjlim_\lambda(S_\lambda\otimes_P\mathcal A_m)
  =S\otimes_P\mathcal A_m
\end{align*}
となる。各$S_\lambda$はnoetherianなので、最後から二番目の同型には既に証明した場合を使える。
これは標準section写像を計算したものであり、基底変更前後の正方形は可換である。
従って同型は任意の基底変更、乗法および合成と両立する。
\end{proof}

\subsection{偶数列に関する補足}

次の各主張は、上の saturation と区別して記載する。
相対Serre双対は第\ref{sec:relative-serre-projective}節で別に構成する。

\begin{proposition}[偶数列の豊富性、復元、関手の決定]
\label{prop:central-model-even-sequence}
noetherian な可換 $P$-代数 $S$ を固定し、
上記の模型において $E_i=\mathcal O(-2i)$、$i\geq0$ とする。
\begin{enumerate}
 \item 任意の coherent 対象と任意の下限 $i_0$ に対し、
 その対象は $i\geq i_0$ を満たす $E_i$ の有限直和の商である。
 また任意の coherent 対象間の全射 $M\to N$ に対し、
 $\operatorname{Hom}(E_i,M)\to\operatorname{Hom}(E_i,N)$ は $i\gg0$ で全射である。
 \item Veronese 関手は基底変更に自然な同値
 $\operatorname{QGr}(\mathcal A_S)\simeq\operatorname{QGr}(\mathcal A_S^{(2)})$ を与える。
 偶数列の section 代数は前定理によって $\mathcal A_S^{(2)}$ となる。
 \item 同じ始域 $\operatorname{QGr}(\mathcal A_S)$ と終域を持つ
 二つの完全同値の偶数列上の制限の間に、全ての Hom と合成に両立する同型が与えられたなら、
 その同型は圏全体の自然同型へ一意的に延びる。
\end{enumerate}
(2) は任意の可換 $P$-代数 $S$ について成り立ち、noetherian 性を必要としない。
\end{proposition}

\begin{proof}
$\ell$ を偶数に選んだので、次数捩れ $(\ell)$ は
有限代数層模型において $\mathcal O_{Y_S}(1)$ との tensor に一致する。
通常のSerreの生成性と消滅\cite[Lemma~30.14.2]{StacksProj}を、coherent な $\mathscr A_S$-加群の台となる
$\mathcal O_{Y_S}$-加群に適用する。
例えばその大きな twist の $\mathcal O_{Y_S}$-加群としての有限生成切断は、
$\mathscr A_S$-加群としての $\mathscr A_S$ の有限直和からの全射を与える。
有限個の剰余類 $2i\bmod\ell$ について同じ議論を行えば (1) を得る。

$\ell$ は二で割れるので、$\mathcal A_S$ と $\mathcal A_S^{(2)}$ は同じ中心標準開上で
同じ次数零局所化代数を持つ。
前命題の局所化と貼合せによる構成は、この二つを同じ $\mathscr A_S$-加群圏へ送る。
これが (2) の同値であり、基底変更にも自然である。
この同値は$\pi\mathcal A_S(-2i)$を$\pi\mathcal A_S^{(2)}(-i)$へ送る。
従って商関手$\operatorname{Gr}(\mathcal A_S^{(2)})\to
\operatorname{QGr}(\mathcal A_S^{(2)})$の右随伴の非負次数部分は，
この同値の下で$\bigoplus_{i\ge0}\operatorname{Hom}(E_i,-)$に対応する。

(1) により coherent 対象 $M$ は、$E_i$ の有限直和を項とする表示
$L_1\to L_0\to M\to0$ を持つ。
二つの完全同値を $F,G$ と書く。
偶数列上の与えられた同型を有限直和に延ばし、
表示の余核を取ると $\alpha_M:F(M)\to G(M)$ が得られる。
固定した全射 $L_0\to M$ に対して、余核上の写像は
$\alpha_M F(L_0\to M)=G(L_0\to M)\alpha_{L_0}$ という等式で一意に決まり、
核を覆う $L_1$ の選択には依存しない。
二つの全射 $L_0\to M$、$L'_0\to M$ を比較するときは、
共通の全射 $L_0\oplus L'_0\to M$ と、その coherent な核の有限表示を取る。
この共通表示から得る写像は、各直和因子について上記の等式を満たす。
$F(L_0)\to F(M)$ と $F(L'_0)\to F(M)$ の全射性により、
元の二つの表示から得る写像はともにこの写像に等しい。

射 $h:M\to N$ に対しては、全射 $L_0\to M$ と $P_0\to N$ を取り、
$(h\circ(L_0\to M),P_0\to N):L_0\oplus P_0\to N$ という全射を用いる。
これから構成した $\alpha_N$ の等式を $L_0$ に制限し、
$F(L_0)\to F(M)$ が全射であることを使うと
$\alpha_NF(h)=G(h)\alpha_M$ が得られる。
この議論は $E_i$ が射影的であることや、表示への任意の射の持上げを仮定しない。
従って coherent 部分圏で表示によらない自然同型を得る。
余核表示はその一意性も与える。

最後に、locally noetherian Grothendieck 圏の各対象は
noetherian 部分対象の filtered colimit であり、
完全同値は noetherian 対象と全ての colimit を保つ。
従って coherent 部分圏上の自然同型は圏全体へ一意的に延びる。
\end{proof}

\subsection{局所 perfect 対象について得られる相対 RHom}

$\operatorname{Perf}(Y_S,\mathscr A_S)$ を、中心標準開上で
有限生成射影 $\mathscr A_S$-加群の有界複体に同型な対象の dg 圏とする。
命題\ref{prop:central-model-qgr}と
\cite[Definitions~3.1, 3.13, Theorem~3.14]{BDGMinors}により、
任意の可換 $P$-代数 $S$ について
\[
 D(\operatorname{QGr}(\mathcal A_S))^c
 \simeq \operatorname{Perf}(Y_S,\mathscr A_S)
\]
である。ここで上添字 $c$ は非有界導来圏の compact 対象を表す。
この同定に必要なのは $Y_S$ の quasi-compact 性と separated 性であり、
$S$ の noetherian 性や fiber の正則性ではない。
全係数基底上では，これを「有界coherentかつ底$S$上有限Tor振幅」という
内在的な部分圏と同一視しない。後の命題\ref{prop:central-model-intrinsic-perfect-open}は，
一般ファイバーの正則性の下で，非空基底開上に限ってその同一性を証明する。

\begin{proposition}[局所 perfect 圏の proper 性と基底変更]
\label{prop:central-model-perfect-rhom}
任意の noetherian $P$-代数 $S$ について、全ての $\mathcal O(m)$ は
$\operatorname{Perf}(Y_S,\mathscr A_S)$ に入る。
この圏の二対象 $K,L$ の相対 RHom は $S$-perfect である。
任意の noetherian $S$-代数 $S'$ に対して、標準写像
\[
 \operatorname{RHom}_{\mathscr A_S}(K,L)\otimes_S^{\mathbf L}S'
 \longrightarrow
 \operatorname{RHom}_{\mathscr A_{S'}}(K_{S'},L_{S'})
\]
は同型である。ここで右辺への移行は導来基底変更である。
正確には、中心 affine chart の代数 $\Lambda$ を使って
$K_{S'}=K\otimes_{\Lambda}^{\mathbf L}(\Lambda\otimes_SS')$
と定め、これらを貼り合わせる。
これは代数の係数拡大であり、$Y_S\to\operatorname{Spec}S$ の平坦性を仮定しないまま
台の $\mathcal O_{Y_S}$-複体の通常の導来逆像と同一視しない。
\end{proposition}

\begin{proof}
中心標準開上で $\mathscr A_S$ は
$\Lambda=(\mathcal A_S[c^{-1}])_0$ で表される。
この代数は $S$ 上平坦である。
実際、次数零局所化は有限自由 $S$-加群
$(\mathcal A_S)_{j\deg c}$ の filtered colimit である。
強次数付き性から $(\mathcal A_S[c^{-1}])_m$ は有限生成射影 $\Lambda$-加群であり、
従って全ての $\mathcal O(m)$ は局所 perfect である。

$K,L$ を局所 perfect とする。
$\mathbf R\mathcal Hom_{\mathscr A_S}(K,L)$ は有界 coherent な
$\mathcal O_{Y_S}$-複体である。
上記の各標準開で有限射影複体を用いて計算する。
有限射影右$\Lambda$-加群間のHomは，$S$-加群として$\Lambda$の有限直和の直和因子であり，
従ってこの局所Hom複体は$S$-平坦加群の有界複体となる。
このHomを自然な右$\Lambda$-加群とみなす必要はない。
有限個の開を選べるので、底に関する Tor 振幅は一様に有限である。
$Y_S\to\operatorname{Spec}S$ は射影的なので、
proper pushforward の perfect 性
\cite[Lemma 36.27.1, Tag~08EV]{StacksPerfect}
をこの $\mathcal O_{Y_S}$-複体に適用できる。
従ってその大域 RHom は $S$-perfect である。

基底変更は、中心標準開とその交叉による有限\v{C}ech 複体で確認する。
各交叉は affine であり、局所 RHom は有限射影 $\Lambda$-複体で計算できる。
$\Lambda$ の $S$-平坦性と有限射影性により、この局所計算は導来基底変更と可換する。
有限\v{C}ech totalization も導来 tensor と可換するので、表示した同型を得る。
この議論は $Y_S$ 自体が $S$ 上平坦であるとは仮定していない。
\end{proof}

\subsection{非空開上の相対対角と内在的なperfect部分圏}

全係数基底上では局所perfect圏だけを用いた。
ここでは一般の幾何学的ファイバーのtails大域次元を入力として，
有限な中心標準開被覆の各相対対角をperfectにする非空基底開を得る。
$u\ne0$と$v\ne0$の二chartだけでは$u=v=0$の境界を覆わないため，
射影中心模型の全標準開被覆を用いる。

\begin{lemma}[中心 affine 開への Ext の局所化]
\label{lem:central-affine-ext-localization}
$k$ を体、$Y$ を noetherian 射影 $k$-スキーム、
$\mathscr A$ を有限な $\mathcal O_Y$-代数とする。
豊富な可逆層 $\mathcal O_Y(1)$ の切断 $c$ に対し
$U=Y\setminus V(c)$ が affine であるとし、$j:U\hookrightarrow Y$、
$\Lambda=\Gamma(U,\mathscr A)$ と置く。
coherent な右 $\mathscr A$-加群 $M,N$ について、自然な同型
\[
 \operatorname{Ext}^i_\Lambda(M|_U,N|_U)
 \simeq\varinjlim_n
 \operatorname{Ext}^i_{\mathscr A}(M,N\otimes\mathcal O_Y(n))
 \qquad(i\geq0)
\]
がある。右辺の遷移写像は $c$ による乗法である。
さらに $\operatorname{coh}(Y,\mathscr A)$ の大域次元が高々 $d$ なら、
$\operatorname{gldim}\Lambda\leq d$ である。
\end{lemma}

\begin{proof}
まず coherent 対象 $M$ に対して、Serre の生成性を繰り返すと、
各項が $\mathscr A\otimes\mathcal O_Y(-n)$ の有限直和である
有界上方の分解 $L^\bullet\to M$ を取れる。
核は常に coherent なので、この操作を任意の長さまで継続できる。
$L^\bullet$ の各項は局所的に有限自由な $\mathscr A$-加群である。
$Y$ の有限 affine 開被覆を選ぶと、
$\mathcal Hom_{\mathscr A}(L^\bullet,N)$ の有限\v{C}ech totalization は
$\operatorname{RHom}_{\mathscr A}(M,N)$ を計算する。
sheaf Hom の次数および\v{C}ech の次数はともに非負で、後者は一様に有界である。
従って固定した total degree に関わる項は有限個である。
局所有限自由性、affine 上の切断、有限\v{C}ech 複体は filtered colimit と可換するので、
$\operatorname{Ext}^i_{\mathscr A}(M,-)$ は quasi-coherent 加群の filtered colimit と可換する。

可逆層を局所的に自明化して通常の環の局所化を使えば
\[
 j_*j^*N=\varinjlim_n N\otimes\mathcal O_Y(n)
\]
となる。
$Y$ は separated で $U$ は affine なので $j$ は affine である。
従って quasi-coherent 加群に対する $j_*$ は完全であり、
$j^*$ との導来随伴と直前の colimit との可換性から表示した Ext 同型を得る。
ここでcoherent対象間のExtはcoherent圏でもquasi-coherent圏でも同じである。
実際，locally noetherian Grothendieck圏ではnoetherian対象への任意の全射は，
あるnoetherian部分対象上で既に全射となる。
従って\cite[Lemma~13.17.4, Tag~0FCL]{StacksAbel}をnoetherian対象のSerre部分圏に
適用でき，必要な有界上方導来圏の比較を得る。

任意の有限生成 $\Lambda$-加群 $Q$ は coherent な $\mathscr A$-加群へ延長できる。
実際、$Q$ の有限個の生成元は
$\Gamma(U,Q)=\Gamma(Y,j_*Q)$ の切断であり、
対応する $\mathscr A^{\oplus r}\to j_*Q$ の像は coherent で、$U$ 上では $Q$ に一致する。
従って大域次元の仮定と Ext の同型により、
全ての有限生成 $\Lambda$-加群 $Q,Q'$ に対して
$\operatorname{Ext}^{i}_\Lambda(Q,Q')=0$、$i>d$ となる。
$\Lambda$ は affine noetherian 可換環上有限なので左右 noetherian である。
有限生成 $Q$ の有限自由分解における $d$ 次 syzygy は有限生成であり、
それに対する有限自由全射の核も有限生成である。
その短完全列は上記の Ext 消滅によって分裂するので、$\operatorname{pd}_\Lambda Q\leq d$ である。
とくに全ての巡回加群 $\Lambda/I$ の射影次元が高々 $d$ である。
Baer の判定を dimension shifting とともに使えば、全ての加群の入射次元が高々 $d$ となり、
$\operatorname{gldim}\Lambda\leq d$ を得る。
\end{proof}

\begin{lemma}[相対対角と有限 Tor 振幅]
\label{lem:relative-diagonal-finite-tor}
$S$ を可換 noetherian 環、$\Lambda$ を $S$ 上平坦な左右 noetherian 代数とする。
$\Lambda$ が対角両側加群として
$\Lambda^{\mathrm{op}}\otimes_S\Lambda$ 上 perfect であると仮定する。
このとき、$K\in D^b_{\mathrm{fg}}(\Lambda)$ に対し
\[
 K\in\operatorname{Perf}(\Lambda)
 \quad\Longleftrightarrow\quad
 K\text{ は }S\text{ 上有限 Tor 振幅を持つ}.
\]
\end{lemma}

\begin{proof}
左辺から右辺は $\Lambda$ の $S$-平坦性から従う。
逆に $K$ の $S$-Tor 振幅を $[a,b]$ とする。
対角 $\Lambda$ の有界有限射影
$\Lambda^{\mathrm{op}}\otimes_S\Lambda$-分解を $Q^\bullet$ とする。
任意の左 $\Lambda$-加群 $N$ に対し
\[
 K\otimes_\Lambda^{\mathbf L}N
 \simeq K\otimes_\Lambda^{\mathbf L}Q^\bullet
              \otimes_\Lambda^{\mathbf L}N
\]
を使う。
$Q^j$ は有限個の外積 $\Lambda\otimes_S\Lambda$ の直和因子なので、
各項の derived tensor は有限個の
$K\otimes_S^{\mathbf L}N$ の直和因子である。
$Q^\bullet$ の長さだけ Tor 振幅を広げれば、全ての $N$ に共通する有界区間を得る。
従って $K$ は $\Lambda$ 上有限 flat dimension を持つ。
$\Lambda$ は noetherian で $K$ の cohomology は有界有限生成なので、
有限生成自由加群の有界上方分解を十分低い次数で切ると、最後の syzygy は有限表示かつ flat になる。
これは射影的であり、$K$ は perfect である。
\end{proof}

この補題は有限次元 $S$-加群性を仮定しない。
$K\otimes_S\Lambda$ は一般には有限生成射影 $\Lambda$-加群ではないので、
それを誤って有限射影と扱うことは避ける。
証明でまず得るのは有限 flat dimension であり、
coherent 性と noetherian 性によって最後に有限射影分解へ移る。

\begin{lemma}[中心上有限な affine 族の一般点からの拡張]
\label{lem:relative-diagonal-generic-spread}
$P$ を noetherian 整域、$K=\operatorname{Frac}P$ を標数零の体とする。
$D$ を有限型可換 $P$-代数、$\Lambda$ を有限な $D$-代数とし、
$\Lambda$ は $P$ 上平坦であるとする。
一般 fiber の対角 $\Lambda_K$ が
$\Lambda_K^{\mathrm{op}}\otimes_K\Lambda_K$ 上 perfect であると仮定する。
この仮定は、例えば $\Lambda_K$ または $\Lambda_{\overline K}$ が有限大域次元を持てば満たされる。
このとき
ある $0\ne p\in P$ に対し $\Lambda_p$ の相対対角は
$\Lambda_p^{\mathrm{op}}\otimes_{P_p}\Lambda_p$ 上 perfect である。
その有限射影分解は任意の $P_p$-基底変更へ移せる。
\end{lemma}

\begin{proof}
有限大域次元から表示した generic perfect 性が従うという補足を先に確認する。
Rennemo--Segal--Van den Bergh\cite[Proposition~2.1, Lemmas~3.3--3.6, Corollary~3.9]{RSVdB}の証明の
steps (3)--(5) は、標数零の体上で有限型可換代数上有限な代数について、
有限大域次元から homological smoothness を導く。
ここでは彼らの Lemma 3.4（反対代数）、Lemma 3.5（PI 代数の tensor）、
Lemma 3.6（代数閉包からの降下）が使われる。
幾何学的一般 fiber $\Lambda_{\overline K}$ の有限大域次元だけを入力とする場合は、
まず代数閉体上の steps (4) によりその対角を perfect とし、
Lemma 3.6 により $K$ 上へ降下すればよい。

包絡代数 $\Lambda^{\mathrm{op}}\otimes_P\Lambda$ は
有限型可換 $P$-代数 $D\otimes_PD$ 上有限であり、
対角 $\Lambda$ はその有限生成加群である。
これは同論文 Corollary 3.9 の適用範囲である。
必要な spreading の議論を明示するため、包絡代数を $\Gamma$ と書き、
$\Lambda$ の有限生成自由 $\Gamma$-加群による分解を取る。
$\Gamma$ は noetherian で、$\Lambda_K$ は perfect なので、
十分高い syzygy $M$ は有限表示であり、$M_K$ は有限生成射影 $\Gamma_K$-加群である。
有限自由加群を経由する分裂
\[
 M_K\xrightarrow{i_K}\Gamma_K^{\oplus r}
       \xrightarrow{p_K}M_K,\qquad p_Ki_K=1
\]
を取る。
$M$ の有限表示性により Hom は中心局所化と可換するので、
分母を払えば、ある $0\ne s\in P$ の後で $i_K,p_K$ は $i_s,p_s$ に降りる。
$p_si_s-1$ は一般点で零であり、$M_s$ の有限個の生成元でその値の分母をさらに払うと、
ある $0\ne t\in P$ の後で $p_{st}i_{st}=1$ が成り立つ。
従って $M_{st}$ は有限生成射影であり、相対対角は $\Gamma_{st}$ 上 perfect となる。
この分裂の分母払いは既知の generic projectivity の議論であり、新規性は主張しない。
$\Lambda$ は $P$-平坦なので、対角の有界有限射影分解に任意の基底変更を適用しても、
導来基底変更した対角は通常の $\Lambda\otimes_PS$ と一致する。
包絡代数も基底変更した包絡代数に一致するため、所望の分解を得る。
\end{proof}

\begin{proposition}[非空開上の内在的な有限 Tor 部分圏]
\label{prop:central-model-intrinsic-perfect-open}
定理\ref{thm:central-model-finite}の仮定を満たす四係数族を考え、
$K=\operatorname{Frac}P$ とする。
追加の入力として
\begin{equation}\label{eq:geometric-generic-tails-input}
 \operatorname{gldim}\operatorname{qgr}\mathcal A_{\overline K}<\infty
\end{equation}
を仮定する。
ある非零 $p\in P$ が存在し、任意の noetherian $P_p$-代数 $S$ について
\[
 \operatorname{Perf}(Y_S,\mathscr A_S)
 =\{K'\in D^b_{\mathrm{coh}}(Y_S,\mathscr A_S):
          K'\text{ は }S\text{ 上有限 Tor 振幅を持つ}\}
\]
となる。
この部分圏は代数の係数拡大による導来基底変更に保たれ、
二対象間の相対 RHom は $S$-perfect であり、任意の noetherian 基底変更と可換する。
\end{proposition}

\begin{proof}
$D=C^{(\ell)}$ の有限個の次数一元 $c_i$ により、$Y$ の中心標準開被覆を取る。
その代数を
\[
 D_i=(D[c_i^{-1}])_0,\qquad
 \Lambda_i=(\mathcal A[c_i^{-1}])_0
\]
と書く。
$D_i$ は有限型可換 $P$-代数であり、$\Lambda_i$ は $D_i$ 上有限で $P$-平坦である。
後者の平坦性は次数零局所化が有限自由 $P$-加群の filtered colimit であることによる。

仮定 \eqref{eq:geometric-generic-tails-input} と中心模型の同値、
補題 \ref{lem:central-affine-ext-localization} を
$Y_{\overline K}$ の各中心 affine 開に適用する。
各 $\Lambda_{i,\overline K}$ は有限大域次元を持つ。
従って前補題の議論によって $\Lambda_{i,K}$ の対角は perfect である。
さらに同補題により、ある非零 $p_i\in P$ の後で $\Lambda_{i,p_i}$ の相対対角は perfect となる。
$p=\prod_i p_i$ とすれば有限個の chart で同時にこれが成立する。
この段階で重要なのは $u,v$ の二 chart ではなく、$Y$ 全体を覆う有限中心被覆を用いたことである。

任意の noetherian $P_p$-代数 $S$ に対し、
$\Lambda_{i,S}$ は $S$-平坦で左右 noetherian であり、相対対角の有限射影分解を持つ。
補題 \ref{lem:relative-diagonal-finite-tor} によって、各 chart で
有界 coherent かつ有限 $S$-Tor 振幅を持つ対象は $\Lambda_{i,S}$-perfect に一致する。
内在的な Tor 振幅の定義は $S$-加群作用から与えられるので、この chart への制限と可換する。
被覆が有限であるため、局所的な有限 Tor 振幅は共通の有限区間で抑えられる。
従って表示した大域的な部分圏の同一性を得る。
基底変更と相対 RHom に関する主張は、有限中心模型における
命題\ref{prop:central-model-perfect-rhom}から従う。
\end{proof}

\paragraph{具体的な族への適用入力と残る問題.}
$B_*$ の固定四係数族については、既存の開な正則性の定理
\ref{CSBGthm:global-open} が、PBW/Koszul 解消と AS 条件の体拡大への両立の下で
任意の代数閉拡大体の fiber の正則性を与える。
これを非零多項式の開集合の geometric generic point へ適用すれば
式 \eqref{eq:geometric-generic-tails-input} の値は三となる。
ここで $K$ 自身への tails 大域次元の降下を省略して仮定したわけではない。
局所化の補題は最初に $\overline K$ 上で適用し、
affine 対角の perfect 性を RSV の Lemma 3.6 により $K$ へ降下している。

条件\textup{(C1),(C2)}は上の全四係数族で検証された。
この命題を条件\textup{(C1)--(C4)}の比較へ用いる場合、さらに得られるのは
非空基底開上の (C3) の内在的部分圏の同定と、その部分圏上の
perfect な相対 RHom および基底変更である。
これは(C4)のRHom部分を満たす。
相対Serre関手の存在と$(-2)[3]$という同定は，
第\ref{sec:relative-serre-projective}節で全四係数基底上に証明する。
そこでcompact対象を使えば，Hom比較には本節の非空開への制限が不要となる。
大域 dg smooth 性には、局所対角に加えて大域的な生成と Morita 比較が必要である。
第\ref{sec:dg-smooth-torsion-families}節では、左右の正則性の下で
外積 classical generator と対角 kernel による比較を与え、
有限位数境界上の相対 smooth proper 開を別に構成する。


\section{四係数族の相対Serre双対と構造対象付きAbel圏}
\label{sec:relative-serre-projective}

この節では、有限な中心模型を持つ固定周囲代数の四係数族について、
体上の双対両側加群の計算を可換係数環上へ移す。
体を次数零部分とする既存の非可換 Serre 双対定理を、
次数零部分が多項式環である場合に直接適用することはしない。
必要な相対双対性は通常の射影空間の Grothendieck 双対性から構成する。

\subsection{仮定と体上の入力}

\[
 P=\mathbb C[a,b,c,d],\qquad
 \mathcal A=B_*\otimes_{\mathbb C}P/
 (G+au^3+bu^2v+cuv^2+dv^3)
\]
とし、有限中心模型の節で構成した
\[
 C\longrightarrow Z(\mathcal A),\quad
 D=C^{(\ell)},\quad
 Y=\operatorname{Proj}D,\quad
 \mathscr A=\widetilde{\mathcal A^{(\ell)}}
\]
を用いる。$\ell$ は十分割り切れる正の整数で、$D$ は次数一で生成される。
有限個の次数一生成元を選び、閉埋込
\[
 j:Y\hookrightarrow Z=\mathbb P^N_P
\]
を固定する。$q:Z\to\operatorname{Spec}P$ を射影とする。
各整数 $m$ に対し
\[
 \mathscr E_m=\widetilde{\mathcal A(m)^{(\ell)}}
\]
と書く。中心局所化の強い次数付き性により、
$\mathscr E_m$ は可逆な $\mathscr A$-両側加群であり、
$\mathscr E_m\otimes_{\mathscr A}\mathscr E_n\simeq\mathscr E_{m+n}$ である。
ここで $\mathscr E_m$ の $m$ は元の次数であり、
射影空間または $Y$ の通常の twist と混同しない。

以下の三つの入力を明記する。
\begin{enumerate}
\item $\mathcal A_n$ は有限自由 $P$-加群である。
  従って中心局所化の次数成分として得られる
  $\mathscr A$ と全ての $\mathscr E_m$ は $P$ 上平坦であり、
  それらの形成は任意の可換基底変更と可換する。
\item 任意の $s\in\operatorname{Spec}P$ に対し、$k=\kappa(s)$ と書くと
  \begin{equation}\label{eq:relative-serre-field-cohom-input}
  R\Gamma(Y_k,\mathscr A_k)\simeq k[0].
  \end{equation}
  これは $\mathcal A_k$ の AS-Gorenstein 次元四、指数二の局所双対性から従う。
  実際、tail の局所コホモロジー完全列で
  $H^1(Y_k,\mathscr A_k)=H^2(Y_k,\mathscr A_k)=0$、
  $H^3(Y_k,\mathscr A_k)\simeq(\mathcal A_k)_{-2}^*=0$ となり、
  saturation により $H^0(Y_k,\mathscr A_k)=(\mathcal A_k)_0=k$ となる。
\item 体上の双対両側加群の補題を次の形で用いる。
  \begin{equation}\label{eq:relative-serre-field-bimodule-input}
  R\mathcal Hom_{\mathcal O_{Z_k}}
  (j_{k,*}\mathscr A_k,\omega_{Z_k/k}[N])
  \simeq j_{k,*}\mathscr E_{-2,k}[3]
  \end{equation}
  は $\mathscr A_k$-\emph{両側加群}の導来圏における同型である。
  この入力は全ての residue field、特に
  $k=\operatorname{Frac}P$ に対して必要である。
  その証明は有限中心多項式環上の局所双対性と
  balanced dualizing complex を比較する補題
  \ref{lem:finite-central-field-dualizing-bimodule} に置く。
  右加群としての双対対象の同定だけでは、この入力を代替できない。
\end{enumerate}

三入力のうち、有限自由性は第\ref{sec:four-parameter-central-model}節で示した。
体上の二つの双対入力は、下の系\ref{cor:four-parameter-field-dualizing}から従う。
従って以下の相対化に新たなファイバー正則性の仮定は不要である。
既存のPBW・AS正則・CY構成への依存は同系に明記する。

\subsection{体上の有限中心模型における双加群双対化}

次の補題では、通常の次数シフトと中心の Veronese 次数を区別する。
特に $\mathscr E_{-2}$ は通常の次数を $-2$ だけ移す双加群であり、
中心の射影空間から来る $\mathcal O(-2)$ とは限らない。

\begin{lemma}[有限中心模型の双加群双対化]
\label{lem:finite-central-field-dualizing-bimodule}
$k$ を体、$A$ を連結、局所有限、左右 noetherian な
非負次数付き $k$-代数とする。
非負次数付きの有限型可換 $k$-代数 $C$、$C_0=k$ から
$Z(A)$ への次数付き準同型があり、
$A$ は $C$ 上有限であるとする。
$C^{(\ell)}$ が次数一で生成されるように $\ell>0$ を選び、
\[
 Y=\operatorname{Proj}C^{(\ell)},\qquad
 \mathscr A=\widetilde{A^{(\ell)}},\qquad
 \mathscr E_t=\widetilde{A(t)^{(\ell)}}
\]
と置く。
$C^{(\ell)}$ の次数一生成元による閉埋込み
$j:Y\hookrightarrow Z=\mathbb P^N_k$ を選ぶ。
さらに $A$ の balanced dualizing complex が
\begin{equation}\label{eq:field-balanced-dualizing-input}
 R_A\simeq A(-2)[4]
 \quad\text{in }D(\operatorname{Gr}A^e)
\end{equation}
を満たすとする。
このとき、$j_*\mathscr A$ の双加群の導来圏において
\begin{equation}\label{eq:field-groth-dualizing-bimodule}
 \mathbb D_{\mathscr A/k}
 :=\mathbf R\mathcal Hom_{\mathcal O_Z}
       (j_*\mathscr A,\omega_Z[N])
 \simeq j_*\mathscr E_{-2}[3]
\end{equation}
である。
有限大域次元、$Y$ の正則性、および $\mathscr A$ の正則性は仮定しない。
体 $k$ の代数閉性も不要である。

右辺に移した双対性の trace は、
\begin{equation}\label{eq:field-groth-dualizing-trace}
 \mathbf R\Gamma(Z,\mathbb D_{\mathscr A/k})
 \xrightarrow{\,\operatorname{ev}_1\,}
 \mathbf R\Gamma(Z,\omega_Z[N])
 \xrightarrow{\,\operatorname{Tr}_{Z/k}\,}k
\end{equation}
で与えられる。
balanced dualizing complex の規格化を選べば、
この trace と有限中心写像の graded trace は一致する。
規格化を指定しない場合の同型は、その規格化の共通スカラーの自由度を持つ。
\end{lemma}

\begin{proof}
$C_\ell$ の上記生成元に対し、
\[
 T=k[y_0,\ldots,y_N],\qquad \deg y_i=\ell
\]
と置く。
$C$ は $C^{(\ell)}$ 上有限であり、$T\to C^{(\ell)}$ は全射である。
従って $T\to C\to A$ は中心に像を持ち、$A$ は $T$ 上有限である。
$C\to A$ または $T\to A$ の単射性は必要ない。
$T$ の balanced dualizing complex は
\[
 R_T=T(-(N+1)\ell)[N+1]
\]
である。
ここで
\[
 S_A=\mathbf R\underline{\operatorname{Hom}}_T(A,R_T)
\]
には、左作用と右作用をともに残す。
具体的には、$R_T$ を表す有界な次数付き $T$-injective 複体 $I$ を取り、
$\underline{\operatorname{Hom}}_T(A,I)$ に
\[
 (a\varphi b)(x)=\varphi(bxa)
\]
で $A$-双加群構造を入れる。
coinduction の随伴と中心性により、この複体は左および右の
$\mathbf R\underline{\operatorname{Hom}}_T(A,R_T)$ を同時に表す。
従って \cite[Proposition~5.2]{YekutieliBalanced} により
$S_A$ は $A$ の dualizing complex である。

この段階で balanced 性を確認しておく。
$\mathfrak n=T_{>0}$、$\mathfrak m=A_{>0}$ とする。
$A$ の $T$ 上の有限性により、$A$-加群に対する
$\mathfrak m$-torsion と $\mathfrak n$-torsion は一致する。
実際、$A/\mathfrak n^q A$ は有限次元なので、
ある十分高い次数以上の $A$ の成分は $\mathfrak n^qA$ に含まれる。
逆向きの包含は次数から従う。
さらに $A$ の有限個の斉次 $T$-生成元を取ると、各項について
\[
 \Gamma_{\mathfrak m}
    \underline{\operatorname{Hom}}_T(A,I^r)
 =\underline{\operatorname{Hom}}_T
    (A,\Gamma_{\mathfrak n}I^r)
\]
となる。
右辺で像の各元を殺す冪が取れるとき、有限個の生成元に対する冪の最大を取れば、
準同型全体を殺す一つの冪を選べるからである。
coinduction により左辺の torsion を取る前の各項は $A$-injective であり、
\cite[Proposition~4.6]{YekutieliBalanced} により
$\Gamma_{\mathfrak n}I^r$ は次数付き $T$-injective である。
$T'=\underline{\operatorname{Hom}}_k(T,k)$ と書くと、
$T$ の balanced 性から
$\Gamma_{\mathfrak n}I\simeq T'$ を得る。
従って、左側の torsion について双加群として
\[
 \mathbf R\Gamma_{\mathfrak m}(S_A)
 \simeq\mathbf R\underline{\operatorname{Hom}}_T(A,T')
 \simeq\underline{\operatorname{Hom}}_k(A,k)=A'
\]
となる。
$T'$ は次数付き $T$-injective なので、最後の式に高次 Ext はない。
右側の torsion も同じ議論で計算できる。
よって $S_A$ は balanced であり、balanced dualizing complex の
双加群としての一意性
\cite[Definition~4.1 の直後、p.~57]{YekutieliBalanced} により
\begin{equation}\label{eq:field-central-polynomial-duality}
 R_A\simeq
 \mathbf R\underline{\operatorname{Hom}}_T(A,R_T)
 \quad\text{in }D(\operatorname{Gr}A^e)
\end{equation}
を得る。
以上の議論は、片側の Serre 対象の一意性だけから双加群の同型を推測するものではない。

次に $T$ の次数を $\ell$ で割り、標準次数付き多項式環 $T_0$ と見なす。
$\operatorname{Gr}T$ は次数の $\ell$ を法とした剰余類で直積分解する。
$R_T$ は剰余類零に属するので、上の $I$ もこの剰余類だけに取れる。
すると複体の各項について
\[
 \bigl(\underline{\operatorname{Hom}}_T(A,I)\bigr)^{(\ell)}
 =\underline{\operatorname{Hom}}_{T_0}
      (A^{(\ell)},I^{(\ell)})
\]
である。
左辺の次数が $\ell$ の倍数である準同型は、$A$ の剰余類零以外を零へ送る。
また $A^{(\ell)}$ の左右の作用は剰余類零を保つ。
従って、この同型は両側の作用と両立し、
\begin{equation}\label{eq:field-veronese-duality}
 (R_A)^{(\ell)}\simeq
 \mathbf R\underline{\operatorname{Hom}}_{T_0}
 \bigl(A^{(\ell)},T_0(-N-1)[N+1]\bigr)
\end{equation}
を双加群として得る。

この同型を $Z=\operatorname{Proj}T_0$ 上で層化する。
導来 Hom と層化の比較はここでも双加群として行える。
実際、$A^{(\ell)}$ は $T_0$ 上有限なので Hom は標準開集合への局所化と可換であり、
noetherian 可換環上の injective 加群の局所化は injective である。
次数付き版は、斉次イデアルの有限表示によって分母を払い、
次数付き Baer 判定法を適用する同じ証明で従う。
標準開集合の局所化環は次数一の可逆元を持つので、
次数零を取る操作は次数付き加群圏と次数零環の加群圏との同値であり、injective 性を保つ。
従って上で選んだ coinduced 複体の局所化が、
両側の作用を保持したまま層の導来 Hom を計算する。
また $\widetilde{T_0(-N-1)}=\omega_Z$ であるから、
\[
 \mathbf R\mathcal Hom_{\mathcal O_Z}
       (j_*\mathscr A,\omega_Z[N])
 \simeq \widetilde{(R_A)^{(\ell)}}[-1]
 \simeq j_*\mathscr E_{-2}[3].
\]
ここで最後にだけ \eqref{eq:field-balanced-dualizing-input} を用いた。

trace について、有限中心写像の dualizing trace は
$T\to A$ の双対、すなわち $\varphi\mapsto\varphi(1)$ である。
中心性により左右の trace は一致する
\cite[Remark~5.7]{YekutieliBalanced}。
剰余類零の取り出しと局所化は単位元およびこの評価写像を保ち、
多項式環の residue の規格化は射影空間の標準 trace を与える。
従って、比較によって運ばれる trace は
\eqref{eq:field-groth-dualizing-trace} である。
既に \eqref{eq:field-balanced-dualizing-input} の同型が選ばれている場合には、
balanced duality の規格化に合わせてそのスカラーを選ぶ。
$A_0=k$ なので、$A(-2)$ の次数零の双加群自己同型の自由度は $k^\times$ である。
\end{proof}

\begin{corollary}[四係数族の全ての体値における untwisted 双対複体]
\label{cor:four-parameter-field-dualizing}
命題\ref{prop:moduli-rigid-example}の $B_*$ について、
補題\ref{lem:pbw}と\ref{lem:CY}の PBW、Koszul AS 正則性、
中山自己同型が恒等であるという入力を用いる。
$P\to k$ を任意の体値とし、対応する二変数三次式を $f$ とする。
\[
 B_k=B_*\otimes_{\mathbb C}k,\qquad
 W_k=G_*+f,\qquad A_k=B_k/(W_k)
\]
と置くと、次数付き双加群として
\[
 R_{B_k}\simeq B_k(-5)[5],\qquad
 R_{A_k}\simeq A_k(-2)[4]
\]
が balanced dualizing complex である。
従って補題\ref{lem:finite-central-field-dualizing-bimodule}は、
四係数族の全ての体値に適用できる。
\end{corollary}

\begin{proof}
まず noetherian 性の体拡大について説明する。
定理\ref{thm:central-model-finite}を四つの係数が零である
$P\to k$ に適用すると、$A_{0,k}=B_k/(G_*)$ は左右 noetherian である。
$G_*$ の正則性は体拡大で保たれ、
$G_*$-adic 随伴次数付き環は $A_{0,k}[t]$ である。
斉次イデアルの随伴次数付きイデアルの有限個の生成元を持ち上げると、
$\deg G_*=3>0$ により、固定した通常の次数では持上げの帰納法が有限回で終わる。
従って $B_k$ の斉次左・右イデアルは有限生成である。
任意の左・右イデアルに通常の次数 filtration を適用すれば、$B_k$ は左右 noetherian である。
これは任意の体拡大が一般に noetherian 性を保つと主張する議論ではない。

既存の PBW filtration の等式と有限な Koszul 解消は、平坦な体拡大と可換である。
有限な次数付き双加群射影解消とその CY 同型も体拡大で保たれる。
よって $B_k$ は Koszul AS 正則で、中山自己同型は恒等であり、
AS index は Koszul 次元の $5$ である。
AS 正則代数の balanced dualizing complex と中山自己同型の関係
\cite[Lemma~3.5 の証明、Theorem~6.3 の証明]{RRZHomological}
から $R_{B_k}\simeq B_k(-5)[5]$ を双加群として得る。

$W_k$ の最高 $x,y,z$-filtration symbol は、
$\mathbb C$ 上の正則な中心三次元を体拡大したものであるから、
$W_k$ は中心正則な次数三の元である。
中心性を保った二項分解
\[
 0\longrightarrow B_k(-3)
 \xrightarrow{\,W_k\,}B_k\longrightarrow A_k\longrightarrow0
\]
により、左右の相対 Hom はともに、自然な両側の作用を持つ
$A_k(3-5)[5-1]=A_k(-2)[4]$ を与える。
ここで通常の normal 元の商に現れる normalizing automorphism は、
$W_k$ が中心元であるため恒等である。
従って \cite[Proposition~5.2, Corollary~5.5]{YekutieliBalanced} を
全射 $B_k\to A_k$ に適用すれば、これが balanced であることも従う。
この計算によって Nakayama の捩れが新たに導入されることはない。
\end{proof}

\begin{remark}[入力の範囲]
この系で使用する $B_*$ の PBW・AS 正則・CY 性は、
補題\ref{lem:pbw}、\ref{lem:CY}とその引用入力に依存する。
この系の独立した証明内容は、体拡大との両立、中心三次商の双加群計算、
および balanced 性の確認である。
有限中心模型の存在は定理\ref{thm:central-model-finite}に依存する。
\end{remark}

\subsection{相対双対両側加群の構成と基底変更}

\begin{lemma}\label{lem:relative-serre-ambient-dual}
$\mathscr A$ を $j$ により $Z$ 上の有限代数層とみなし、
\[
 \mathscr D_P=
 R\mathcal Hom_{\mathcal O_Z}
 (j_*\mathscr A,\omega_{Z/P}[N]),
 \qquad \omega_{Z/P}=\mathcal O_Z(-N-1)
\]
と置く。この複体は自然に $\mathscr A$-両側加群の複体であり、
任意の可換 $P$-代数 $S$ に対して自然な同型
\begin{equation}\label{eq:relative-serre-dual-basechange}
 \mathscr D_P\otimes_P^{\mathbf L}S\simeq
 \mathscr D_S:={}
 R\mathcal Hom_{\mathcal O_{Z_S}}
 (j_{S,*}\mathscr A_S,\omega_{Z_S/S}[N])
 \quad\text{in }D(\mathcal O_{Z_S})
\end{equation}
を持つ。これは両側作用、$1$ における評価写像、および基底変更の合成と両立する。
この段階の導来基底変更は、まず underlying な
$\mathcal O_Z$-perfect 複体の導来基底変更として読む。
両側作用との両立は、全ての左乗法と右乗法についての自然性を意味する。
以下で cohomology が一つの $P$-flat 両側加群に集中すると示した後には、
これは通常の両側加群の係数基底変更そのものになる。
$\mathscr A^{\mathrm{op}}\otimes_{\mathcal O_Z}\mathscr A$ が
$P$ 上平坦であるという仮定は用いない。
\end{lemma}

\begin{proof}
$P$ は有限 Krull 次元の正則 noetherian 環であり、$Z$ も正則である。
従って coherent $\mathcal O_Z$-加群 $j_*\mathscr A$ は
$\mathcal O_Z$ 上 perfect である。
各 affine open 上で有限長の有限射影 $\mathcal O_Z$-分解を取れる。
その各項は $P$ 上平坦であり、分解される加群 $j_*\mathscr A$ 自身も $P$ 上平坦である。
従ってこの分解は任意の $P\to S$ の後にも分解である。
これは、増大した分解が有界な平坦加群の完全複体であり、
その cycle が端からの帰納法で平坦となることから従う。

分解の各項は有限射影なので、$\omega$ への Hom は任意の係数基底変更と可換する。
よって \eqref{eq:relative-serre-dual-basechange} を得る。
この写像は分解の選択を用いず、導来 Hom の通常の基底変更写像として定義できる。
したがって局所計算は貼り合わさり、基底変更の合成にも両立する。

両側作用は、$\mathscr A$ の左乗法と右乗法を Hom の第一変数へ作用させて得る。
例えば通常の Hom の段階では
$(a\varphi b)(x)=\varphi(bxa)$ である。
導来 Hom の基底変更写像の自然性はこの二つの作用を保つ。
同様に、単位写像 $\mathcal O_Z\to j_*\mathscr A$ による反変性が与える
\begin{equation}\label{eq:relative-serre-evaluation-one}
 \epsilon_S:\mathscr D_S\longrightarrow\omega_{Z_S/S}[N]
\end{equation}
とも両立する。
両側加群の複体は、$\omega[N]$ の injective 分解 $I^\bullet$ に対する
$\mathcal Hom_{\mathcal O_Z}(j_*\mathscr A,I^\bullet)$ としても表せる。
その全項は $Y$ の定義イデアルで消えるため、$Y$ 上の両側加群複体として扱える。
この説明では、閉埋込 $j$ の導来 direct image が fully faithful であるとは仮定していない。
\end{proof}

\begin{lemma}[相対振幅零の判定]\label{lem:relative-serre-amplitude-zero}
$f:Z\to\operatorname{Spec}P$ は上の射影空間とし、
$K$ を $Z$ 上の perfect 複体とする。
全ての $s\in\operatorname{Spec}P$ について
$K\otimes_P^{\mathbf L}\kappa(s)$ が次数零のみに cohomology を持つなら、
$K\simeq W[0]$ となる coherent 層 $W$ が存在し、$W$ は $P$ 上平坦である。
\end{lemma}

\begin{proof}
$z\in Z$、$s=f(z)$ を固定し、局所 noetherian 環の局所準同型
$R=P_s\to B=\mathcal O_{Z,z}$ を考える。
$B$ は $R$ 上平坦である。
$K_z$ を有限個の有限自由 $B$-加群の複体 $F^\bullet$ で表す。
この複体の $R$ の剰余体への導来 tensor は、通常の tensor で計算できる。
仮定により $F^\bullet\otimes_R\kappa(s)$ は次数零以外で完全である。

最高の非零次数 $b>0$ がある場合、
$\operatorname{coker}(F^{b-1}\to F^b)\otimes_R\kappa(s)=0$ である。
cokernel は有限 $B$-加群であり、$\mathfrak m_RB\subset\mathfrak m_B$ なので
Nakayama の補題から $F^{b-1}\to F^b$ は全射である。
$F^b$ は自由であるからこの全射は分裂し、可縮な二項複体を除ける。
これを繰り返し、$F^i=0$ for $i>0$ としてよい。

残った有限複体の全項は $R$ 上平坦であり、剰余体後に次数零より左で完全である。
\cite[Lemma~10.99.5, Tag~\texttt{00MI}]{StacksRelativeDuality}により、
元の複体も次数零より左で完全で、終端 cokernel は $R$ 上平坦となる。
これを全ての $z$ に適用すると、$K$ の非零 cohomology は $H^0(K)=W$ のみであり、
$W$ は $P$ 上平坦である。
\end{proof}

\begin{proposition}[双対両側加群の相対的同定]
\label{prop:relative-serre-dual-bimodule}
上記の体上の入力の下で、$\mathscr A$-両側加群の導来圏における同型
\begin{equation}\label{eq:relative-serre-global-dual}
 \Theta_P:\mathscr E_{-2}[3]\xrightarrow{\ \sim\ }\mathscr D_P
\end{equation}
を選べる。一度この同型を選ぶと、任意の可換 $P$-代数 $S$ に対し
\begin{equation}\label{eq:relative-serre-all-base-dual}
 \Theta_S:\mathscr E_{-2,S}[3]\xrightarrow{\ \sim\ }\mathscr D_S
\end{equation}
が基底変更に両立して定まる。
\end{proposition}

\begin{proof}
補題 \ref{lem:relative-serre-ambient-dual} と
入力 \eqref{eq:relative-serre-field-bimodule-input} により、
$\mathscr D_P[-3]$ の全ての導来 fibre は次数零に集中する。
この複体は $\mathcal O_Z$ 上 perfect なので、
補題 \ref{lem:relative-serre-amplitude-zero} より
\[
 \mathscr D_P\simeq\mathscr W[3]
\]
となる coherent な $P$-flat 層 $\mathscr W$ が得られる。
この同型は、両側加群複体の標準 truncation として取れる。
従って $\mathscr W$ は $\mathscr A$-両側加群であり、
$\mathscr W\otimes_P\kappa(s)\simeq\mathscr E_{-2,\kappa(s)}$ は
両側加群として成立する。
実際、$\mathscr W$ の $P$-flatness により補題の underlying な導来基底変更は
通常の $\mathscr W\otimes_P\kappa(s)$ に集中し、
左右の乗法についての自然性がこの cohomology 同型の両側線型性を与える。

まず右加群としての大域同型を作る。
\[
 \mathscr V=
 \mathcal Hom_{\mathscr A}(\mathscr E_{-2},\mathscr W)
 \simeq\mathscr W\otimes_{\mathscr A}\mathscr E_2
\]
と置く。
$\mathscr E_2$ は左 $\mathscr A$-加群として局所有限射影であるから、
$\mathscr V$ は局所的に $\mathscr W$ の有限直和の直和因子となる。
従って $\mathscr V$ は coherent かつ $P$-flat であり、
形成は任意の基底変更と可換する。
その fibre は $\mathscr A_{\kappa(s)}$ と同型なので、
\eqref{eq:relative-serre-field-cohom-input} により
\[
 R\Gamma(Y_{\kappa(s)},\mathscr V_{\kappa(s)})\simeq\kappa(s)[0].
\]

$Y$ は $P$ 上 proper で、$\mathscr V$ は coherent かつ $P$-flat である。
\cite[Lemma~30.22.1, Tag~\texttt{07VK}]{StacksRelativeDuality}により
$R\Gamma(Y,\mathscr V)$ は perfect $P$-複体で、任意の基底変更と可換する。
その全ての導来 fibre は一次元で次数零に集中するため、
\[
 R\Gamma(Y,\mathscr V)\simeq L[0]
\]
となる階数一の有限射影 $P$-加群 $L$ が存在する。
この最後の主張は有限自由複体の局所最小化、または有限射影複体の fibre による
Tor 振幅判定から従う。

Hom の評価は右 $\mathscr A$-加群写像
\begin{equation}\label{eq:relative-serre-line-evaluation}
 e:L\otimes_P\mathscr E_{-2}\longrightarrow\mathscr W
\end{equation}
を与える。
全ての fibre で、これは
$\operatorname{Hom}_{\mathscr A_k}(\mathscr E_{-2,k},\mathscr W_k)
 \otimes_k\mathscr E_{-2,k}\to\mathscr W_k$
と一致する。
右加群として $\mathscr W_k\simeq\mathscr E_{-2,k}$ であり、
\[
 \operatorname{End}_{\mathscr A_k}(\mathscr E_{-2,k})
 \simeq H^0(Y_k,\mathscr A_k)=k
\]
なので、各 fibre の評価は同型である。
まず有限生成 cokernel に局所 Nakayama を適用すると $e$ は全射である。
次に右辺が $P$-flat であるため、$\ker e$ の fibre は fibre 写像の核に一致する。
これも零なので、もう一度局所 Nakayama を適用して $\ker e=0$ を得る。

多項式環 $P$ の Picard 群は零である。
例えば $P$ が UFD であることから、その可逆イデアルは主イデアルとなる。
従って $L\simeq P$ を一度選び、右加群同型
\[
 e_0:\mathscr E_{-2}\xrightarrow{\sim}\mathscr W
\]
を得る。

この同型が左作用も保つことを示す。
$\eta$ を $\operatorname{Spec}P$ の generic point、$K=\operatorname{Frac}P$ とする。
体上の両側加群入力により、ある両側加群同型
$\beta_\eta:\mathscr E_{-2,K}\to\mathscr W_K$ が存在する。
$\beta_\eta^{-1}e_{0,\eta}$ は右加群の自己同型であり、
上の End 計算により $K^\times$ の scalar multiplication である。
したがって $e_{0,\eta}$ は左作用も保つ。

二つの左作用を比較する写像
\[
 \Delta:\mathscr A\otimes_{\mathcal O_Y}\mathscr E_{-2}
 \longrightarrow\mathscr W,\qquad
 a\otimes x\longmapsto e_0(ax)-a e_0(x)
\]
は generic fibre で零である。
$\mathscr W$ は整域 $P$ 上平坦なので、各 affine open の加群は
$P$-torsion-free である。
従って $\mathscr W\to\mathscr W\otimes_P K$ は単射であり、$\Delta=0$ となる。
よって $e_0$ は両側加群同型である。
これに shift $[3]$ を施して \eqref{eq:relative-serre-global-dual} を得る。

$\mathscr W$ は $P$-flat なので、任意の $P\to S$ に対して
補題 \ref{lem:relative-serre-ambient-dual} の ambient な基底変更を cohomology に取ると
\[
 \mathscr W\otimes_P S\xrightarrow{\sim}H^{-3}(\mathscr D_S)
\]
を得て、$\mathscr D_S$ の他の cohomology は零である。
左右の乗法との自然性により、この写像は通常の
$\mathscr A_S$-両側加群の同型である。
従って $e_0\otimes_P S$ と両側加群複体の標準 truncation を合成すると、
\eqref{eq:relative-serre-all-base-dual} が
両側加群の導来圏における真正な同型として得られる。
この構成で、enveloping algebra に沿う導来係数 tensor の平坦性は必要ない。
全ての写像は一つの $e_0$ と自然な ambient 比較の cohomology から定まるため、
係数環の変更の合成に両立する。
\end{proof}

\begin{remark}[正規化と平坦性の使い分け]
同型 $L\simeq P$ の選択により、$\Theta_P$ には scalar の正規化の自由度がある。
これを各 fibre で別々に選び直すことはしない。
一度選んだ $\Theta_P$ のみから全ての基底上の trace と双対性を定める。
また、この証明では $Y\to\operatorname{Spec}P$ の平坦性を使わない。
使うのは $Z=\mathbb P^N_P$ の平坦性と、$\mathscr A$、
$\mathscr E_m$、$\mathscr W$ の係数環上の平坦性である。
\end{remark}

\subsection{局所 perfect 対象の相対 Serre 双対性}

可換 noetherian $P$-代数 $S$ を固定する。
$\operatorname{Perf}(Y_S,\mathscr A_S)$ は、中心の affine open 被覆上で
有限生成射影右 $\mathscr A_S$-加群の有界複体となる対象の圏とする。
$K(m)=K\otimes_{\mathscr A_S}\mathscr E_{m,S}$ と書く。
ここでは、有限 $S$-Tor 振幅を持つ coherent 複体全体がこの圏に等しいとは主張しない。

\begin{lemma}[内部 Hom の双対性]\label{lem:relative-serre-internal-hom}
$K,L\in\operatorname{Perf}(Y_S,\mathscr A_S)$ に対し、
自然な同型
\begin{equation}\label{eq:relative-serre-internal-duality}
 R\mathcal Hom_{\mathcal O_{Z_S}}
 \bigl(j_{S,*}R\mathcal Hom_{\mathscr A_S}(K,L),\omega_{Z_S/S}[N]\bigr)
 \simeq
 j_{S,*}R\mathcal Hom_{\mathscr A_S}
 \bigl(L,K\otimes_{\mathscr A_S}^{\mathbf L}\mathscr D_S\bigr)
\end{equation}
がある。これは変数 $K,L$、係数基底変更、および評価による pairing に自然である。
\end{lemma}

\begin{proof}
局所的に、可換環 $T$ 上有限な代数 $\Lambda$ とその右 perfect 複体を考える。
$K^*=R\operatorname{Hom}_{\Lambda}(K,\Lambda)$ は左 perfect 複体であり、
\[
 R\operatorname{Hom}_{\Lambda}(K,L)
 \simeq L\otimes_{\Lambda}^{\mathbf L}K^*.
\]
係数の双対複体を $J$、$D_\Lambda=R\operatorname{Hom}_T(\Lambda,J)$ と書くと、
tensor--Hom 随伴と有限射影加群の双対性により
\begin{align*}
 R\operatorname{Hom}_T
  (R\operatorname{Hom}_{\Lambda}(K,L),J)
 &\simeq
 R\operatorname{Hom}_{\Lambda}
  (L,R\operatorname{Hom}_T(K^*,J))\\
 &\simeq
 R\operatorname{Hom}_{\Lambda}
  (L,K\otimes_{\Lambda}^{\mathbf L}D_\Lambda).
\end{align*}
第二の同型は $K=\Lambda$ のとき定義そのものである。
有限直和、shift、cone、直和因子を取ることに両立するため、
任意の有限射影複体に成立する。
従って $\Lambda$ の有限 global dimension は必要ない。
導来随伴を具体的に確かめるには、$J$ の $T$-injective 分解 $I^\bullet$ を取ればよい。
$\operatorname{Hom}_T(\Lambda,I^q)$ は、右 $\Lambda$-加群の忘却関手が exact であることと
その右随伴から、右 $\Lambda$-injective である。
有限射影左加群 $K^*$ に対する $\operatorname{Hom}_T(K^*,I^q)$ もその有限直和の
直和因子なので、同じ計算が導来 Hom を計算する。

この計算で $D_\Lambda$ の右作用は、
$\operatorname{Hom}_T(\Lambda,J)$ の引数への左乗法から来る。
従って最後の右 $\Lambda$-加群構造は
$K\otimes_\Lambda D_\Lambda$ の正しい右作用であり、
右加群としての同型だけからこの式を得ているのではない。
全ての写像は評価と随伴の単位・余単位で定義される。
そのため局所計算は貼り合わさり、
\eqref{eq:relative-serre-internal-duality} を与える。

基底変更については、$\mathscr A_S$ が $S$-flat であることから
局所 perfect 複体は有界な $S$-flat 複体で表せる。
有限射影間の Hom は $S$-加群として
$\mathscr A_S$ の有限直和の直和因子である。
また $j_{S,*}\mathscr A_S$ は補題
\ref{lem:relative-serre-ambient-dual} の分解を基底変更した
$\mathcal O_{Z_S}$ 上の有限 perfect 分解を持つ。
従って $j_{S,*}R\mathcal Hom_{\mathscr A_S}(K,L)$ も
$\mathcal O_{Z_S}$ 上 perfect である。
実際、局所有限射影表示の各 Hom の項は
$j_{S,*}\mathscr A_S$ の有限直和の $\mathcal O_{Z_S}$-直和因子である。
したがって上の有限射影 Hom と ambient perfect dual の基底変更写像は同型である。
評価と tensor--Hom 随伴はこれらの基底変更写像に両立する。
なお、用いている基底変更は係数環に関する導来 tensor、
または $\mathscr A_S$ に沿う導来 extension of scalars である。
非平坦であり得る $Y_S\to\operatorname{Spec}S$ に沿う
通常の $\mathcal O_{Y_S}$-導来 pullback との無条件の同一視はしていない。
一方、$Z_S\to\operatorname{Spec}S$ は flat なので、
$Z_S$ 上の複体については係数の導来 tensor が
$Z_{S'}\to Z_S$ に沿う導来 pullback に一致する。
\end{proof}

\begin{theorem}[四係数族全体の相対Serre公式]
\label{thm:relative-serre-projective}
上記の入力の下で、全ての可換 noetherian $P$-代数 $S$ と
$K,L\in\operatorname{Perf}(Y_S,\mathscr A_S)$ に対し、
\begin{equation}\label{eq:relative-serre-formula}
 R\operatorname{Hom}_S
 \bigl(R\operatorname{Hom}_{\mathscr A_S}(K,L),S\bigr)
 \simeq
 R\operatorname{Hom}_{\mathscr A_S}(L,K(-2)[3])
\end{equation}
という自然な同型がある。
両辺は perfect $S$-複体であり、この同型は任意の noetherian な
$S\to S'$ による導来基底変更と両立する。
従って $K\mapsto K(-2)[3]$ はこの perfect 圏の相対 Serre 自己同値である。
\end{theorem}

\begin{proof}
$M=j_{S,*}R\mathcal Hom_{\mathscr A_S}(K,L)$ と置く。
前補題の証明により、$M$ は $\mathcal O_{Z_S}$ 上 perfect である。
$q_S:Z_S\to\operatorname{Spec}S$ は proper、flat、有限表示なので、
\cite[Lemma~36.30.4(1), Tag~\texttt{0B91}]{StacksRelativeDuality}により
$R\Gamma(Z_S,M)$ は perfect $S$-複体であり、
形成は任意の係数基底変更と可換する。

$q_S:Z_S=\mathbb P^N_S\to\operatorname{Spec}S$ は smooth projective であり、
その相対双対複体は $\omega_{Z_S/S}[N]$ である。
従って通常の Grothendieck 双対性は
\[
 R\operatorname{Hom}_S(R\Gamma(Z_S,M),S)
 \simeq
 R\Gamma\bigl(Z_S,
 R\mathcal Hom_{\mathcal O_{Z_S}}(M,\omega_{Z_S/S}[N])\bigr)
\]
を与える。
射影空間上の双対複体と trace については
\cite[Lemma~48.15.1, Tag~\texttt{0A9W}]{StacksRelativeDuality}、
相対 pairing による上の同型については Remark 48.12.6 を用いる。
補題 \ref{lem:relative-serre-internal-hom} と
$\Theta_S$ を代入すると \eqref{eq:relative-serre-formula} を得る。
右辺が perfect であることは、
$K(-2)[3]$ も局所 perfect であることから同じ議論で従う。

最後に、同型を与える trace と基底変更の整合性を明示する。
構造対象の trace は
\begin{equation}\label{eq:relative-serre-trace-structure}
 R\Gamma(Y_S,\mathscr E_{-2,S}[3])
 \xrightarrow{R\Gamma(\Theta_S)}
 R\Gamma(Z_S,\mathscr D_S)
 \xrightarrow{R\Gamma(\epsilon_S)}
 R\Gamma(Z_S,\omega_{Z_S/S}[N])
 \xrightarrow{\operatorname{Tr}_{q_S}} S
\end{equation}
と定義する。
一般の $K$ に対する trace は、既に構成した双対同型の $L=K$ の場合と
単位 $\operatorname{id}_K:S\to R\operatorname{Hom}_{\mathscr A_S}(K,K)$ を用いて
\[
 \tau_K:R\operatorname{Hom}_{\mathscr A_S}(K,K(-2)[3])
 \xrightarrow{\sim}
 R\operatorname{Hom}_S
  (R\operatorname{Hom}_{\mathscr A_S}(K,K),S)
 \xrightarrow{(\operatorname{id}_K)^*}S
\]
と定義する。$K=\mathscr A_S$ の場合、これは
\eqref{eq:relative-serre-trace-structure} と一致する。
この定義は非可換代数の射影加群から両側加群への通常の matrix trace を
無条件に定義することを避け、導来随伴の評価だけを用いる。
これと導来合成写像により
\[
 R\operatorname{Hom}_{\mathscr A_S}(L,K(-2)[3])
 \otimes_S^{\mathbf L}
 R\operatorname{Hom}_{\mathscr A_S}(K,L)
 \longrightarrow S
\]
を得る。
内部 Hom の補題の評価随伴を追えば、この pairing の転置は
ちょうど \eqref{eq:relative-serre-formula} の同型である。
具体的には双対同型の第二変数についての自然性により、
$f:K\to L$ と $g:L\to K(-2)[3]$ の評価は
$\tau_K(g\circ f)$ に一致する。複体の場合の符号は導来合成の Koszul 規約に従う。
したがって、単に全ての fibre で非退化な pairing が存在することを用いているのではない。

$\Theta_S$ は一つの $\Theta_P$ の基底変更であり、
$\epsilon_S$ は単位の反変性で定まる。
また $q_S$ は flat なので、全ての $S\to S'$ は $q_S$ と Tor independent である。
射影空間の trace は任意の基底変更に両立する
（\cite[Lemma~48.7.1, Remark~48.12.5]{StacksRelativeDuality}）。
Hom の導来基底変更は上で既に確認されており、
perfect $S$-複体の $S$-双対も任意の基底変更と可換する。
従って pairing とその転置は基底変更に両立する。
評価、随伴、trace のいずれも自然変換であるため、
変数 $K,L$ の射と合成、および基底変更の合成に関する整合性も成立する。
\end{proof}

\subsection{内在的な perfect 圏と既存の有限 Tor 条件}

中心模型から定義した局所 perfect 圏は、実際には Abel 圏から内在的に回復できる。
Burban--Drozd--Gavran, \emph{Minors and resolutions of non-commutative schemes},
\cite[Definition~3.13, Theorem~3.14]{BDGMinors}により、quasi-compact separated な
スキーム上の quasi-coherent 代数層に対して
\[
 D(\operatorname{QCoh}(Y_S,\mathscr A_S))^c
 =\operatorname{Perf}(Y_S,\mathscr A_S)
\]
となる。左辺は非有界導来圏の compact 対象である。
Abel 圏の exact equivalence は非有界導来圏へ成分ごとに延長し、
任意直和を保つので compact 対象を保つ。
従って、この perfect 圏は選んだ有限中心模型から独立である。

この既知の特徴付けは、有限 $S$-Tor 振幅を持つ coherent 対象全体と
局所 perfect 圏との一致を意味しない。
例えば体上では全ての coherent 対象が係数体に関して Tor 振幅零であるが、
代数層が有限 global dimension を持たなければ全てが perfect となるわけではない。
元の公理で有限係数 Tor 圏が指定されている場合、
compact 圏を用いる公理への変更と、その元の公理の証明とは区別する。

\begin{lemma}[compact 版の相対比較で必要な保存性]
\label{lem:relative-serre-compact-intrinsic}
相対 Hom 定理の条件 \textup{(C3)} で、$\mathcal P_U$ を
標準 dg enhancement の compact 対象の full dg 部分圏
$D(\mathcal C_U)^c$ に置き換える。
全ての $E_i$ の compact 性と、指定した基底変更による compact 性の保存を仮定する。
この条件と既存の \textup{(C1),(C2),(C4)} の下で、
同定対象 $\mathcal P_U$ の保存性、相対 Serre との可換性、
および構造対象付き相対 Hom 定理の証明は同じ結論を与える。
有限中心模型では、この compact 性と基底変更による保存性は成立する。
\end{lemma}

\begin{proof}
exact $S$-linear Abel 圏同値 $F:\mathcal C_U\to\mathcal C'_U$ は
複体へ項ごとに作用し、acyclic 複体を保存する。
圏同値は全ての直和を保つため、誘導される非有界導来同値は compact 対象を保つ。
標準の injective enhancement を用いるときも、
acyclic 複体 $M$ と K-injective 複体 $I$ に対して
\[
 \operatorname{Hom}^{\bullet}(M,F(I))
 \simeq\operatorname{Hom}^{\bullet}(F^{-1}(M),I)
\]
なので、$F$ は K-injective 複体を保つ。
従って $S$-複体としての相対 RHom とその合成を保つ dg 同値を得る。
相対 Serre 双対とその trace を含む表現の dg Yoneda 一意性により
$F\mathsf S\simeq\mathsf S'F$ が定まり、同値の合成と指定した基底変更に整合する。

有限中心模型では、$\mathscr E_m$ は局所有限射影の右加群である。
Burban--Drozd--Gavran の compact 性の特徴付けにより、全て compact である。
係数基底変更の extension of scalars は、局所的な有限射影加群を有限射影加群へ送る。
有界性と直和因子も保存するため、局所 perfect 対象、従って compact 対象を保存する。
その相対 RHom と基底変更は既に上で構成したものに一致する。

相対 Hom 定理が後に用いるのは、Serre 軌道に全ての $E_i$ が属すること、
これらの Hom を双対性と同値が保つこと、そして \textup{(C2)} の豊富列復元だけである。
これらは全て compact 部分圏上で成立するので、同じ証明が適用できる。
\end{proof}

\begin{remark}[Serre 表示の変更に必要な scalar の主張]
\label{rem:relative-serre-scalar-on-lines}
この比較において、perfect 圏全体の関手の全ての自然自己同型が
scalar であるとは主張しない。
必要なのは $T=(-2)[3]$ の二つの表示の比較を line 対象へ制限したときの主張である。
$\alpha:T\Rightarrow T$ を可逆な自然変換とすると、
$T(E_i)=E_{i+1}[3]$ と $\operatorname{End}(E_{i+1})=S$ により
$\alpha_{E_i}$ はある $\lambda_i\in S^\times$ による scalar multiplication となる。
$\operatorname{Hom}(E_{i+1},E_i)=(\mathcal A_S)_2$ 上の自然性から
\[
 (\lambda_{i+1}-\lambda_i)(\mathcal A_S)_2=0
\]
を得る。$(\mathcal A_S)_2$ は正の階数の有限局所自由加群なので忠実であり、
$\lambda_{i+1}=\lambda_i$ となる。
従って比較は全ての $E_i$ 上で一つの $\lambda\in S^\times$ による。
Serre の反復では、軌道環の次数 $n$ に $\lambda^{n}$ または
逆の表示規約では $\lambda^{-n}$ が掛かる。
いずれも相対同型層で既に割っている次数スカラーである。
この議論は $S$ の被約性を仮定しない。
Abel 圏の自然変換を全圏上で決める段階では、
別途 \textup{(C2)} の豊富列生成性を用いる。
\end{remark}

\subsection{剛体性の群スキームへの昇格}

\begin{lemma}[周囲の次数付き自己同型群スキーム]
\label{lem:rigid-ambient-group-scheme}
$\operatorname{Aut}_{\mathrm{gr}}(B_*)(\mathbb C)=\mathbb C^\times$
という既存の入力の下で、群スキームとして
\[
 \underline{\operatorname{Aut}}_{\mathrm{gr}}(B_*)=\mathbb G_m
\]
である。従って任意の可換 $\mathbb C$-代数 $S$ 上でも、
全ての次数付き $S$-線型自己同型は次数スカラーである。
\end{lemma}

\begin{proof}
$V=(B_*)_1\simeq\mathbb C^5$、
$R_2=\ker(V\otimes V\to(B_*)_2)$ とする。
$B_*$ は二次代数なので、その次数付き自己同型関手は
$R_2$ を保つ $\operatorname{GL}(V)$ の閉部分群スキームで表される。
特にこれは $\mathbb C$ 上有限型である。
任意の基底上で $g(R_2)\subset R_2$ と書いた閉条件は等号を意味する。
実際、$R_2$ とその補空間に関する $g\otimes g$ の行列は上三角ブロックであり、
全体が可逆なので対角ブロックも可逆となる。

標数零の有限型群スキームは Cartier の定理により smooth、従って被約である
（\cite[Lemma~39.8.2, Tag~\texttt{047N}]{StacksRelativeDuality}）。
次数スカラーの $\mathbb G_m$ はこの群の閉部分群であり、
入力により全ての複素閉点を含む。
両者を $\operatorname{GL}(V)$ の閉部分スキームとしてみると、
Hilbert の零点定理と被約性から同じ定義イデアルを持つ。
従って群スキームとして等しい。
\end{proof}

\begin{corollary}[任意の係数環上の三次商の同型]
\label{cor:rigid-family-all-coeff-isom}
$S$ を任意の可換 $\mathbb C$-代数とし、$f,g\in S[u,v]_3$ とする。
$A_f=B_{*,S}/(G_*+f)$、$A_g=B_{*,S}/(G_*+g)$ に対し
\[
 \operatorname{Isom}_{\mathrm{gr},S}(A_f,A_g)
 =\begin{cases}
   S^\times & f=g,\\
   \varnothing & f\ne g
  \end{cases}
\]
であり、$S^\times$ は次数スカラーとして作用する。
\end{corollary}

\begin{proof}
三次関係は次数二以下を変えないので、$A_f,A_g$ の二次包絡はともに $B_{*,S}$ である。
従って次数付き同型は $B_{*,S}$ の自己同型へ一意に延長し、前補題から次数スカラー $\lambda$ である。
三次核の次数三部分は、それぞれ $S(G_*+f)$、$S(G_*+g)$ である。
これらの直線が対応するので、ある $r\in S$ により
\[
 \lambda^3(G_*+f)=r(G_*+g)
\]
と書ける。
主変数 filtration の最高次 symbol は両方とも固定した $\gamma$ であり、
この成分を比較すると $r=\lambda^3$ となる。
$\lambda$ は単元なので $f=g$ である。
逆に $f=g$ なら全ての次数スカラーが同型を与える。
この議論では $S$ の被約性も noetherian 性も用いない。
\end{proof}

\subsection{compact 版比較と stack 化後の射}

\begin{corollary}[全四係数族の相対同型層]
\label{cor:central-model-relative-hom}
$T$を有限型$\mathbb C$-スキーム，$f,g:T\to\mathbb A^4$を二つの係数写像とする。
対応する族の間に，底変更と合成に自然なfppf層の同型
\[
 \underline{\operatorname{Eq}}_{\mathcal O,T}(A_f,A_g)
 \simeq
 \underline{\operatorname{Isom}}_{\mathrm{gr},T}(A_f,A_g)/\mathbb G_m
\]
がある。左辺は構造対象を保つ全ての相対完全同値の自然同型類を含む。
同定対象上で恒等な2自己同型は自明である。
ファイバーの有限大域次元または正則性は仮定しない。
\end{corollary}
\begin{proof}
第\ref{sec:four-parameter-central-model}節は条件\textup{(C1),(C2)}と
普遍的中心二変数条件を与える。
条件\textup{(C3)}の代わりに，
補題\ref{lem:relative-serre-compact-intrinsic}のcompact部分圏を使う。
その上の\textup{(C4)}は定理\ref{thm:relative-serre-projective}で証明した。
従って定理\ref{thm:abel-relative-pointed-hom}の証明と
注意\ref{rem:relative-serre-scalar-on-lines}を適用できる。
元の有限係数Tor振幅部分圏による\textup{(C3)}の全基底上での成立を用いたのではない。
\end{proof}

ここで定義\ref{def:abel-cubic-pointed}の相対模型としての許容性も確認する。
固定した$B_*$のKoszul解消とその両側AS双対複体は，
体上の有限自由構成を係数環へtensorするため任意の底変更後も完全である。
三次核は固定した非零正則symbol$\gamma$を持ち，
商の各次数は普遍的に自由なので，$G_*+f$による乗法は任意係数上で左右正則である。
全shiftのsection飽和は定理\ref{thm:central-model-universal-sections}により
非noetherian係数も含めて成立する。
Gabriel商の標準対象列は生成族であり，そのsection代数はこの飽和によって
元の座標代数に一致する。商関手とその右随伴による復元は任意係数変更に自然である。
従って四係数族は，指定した全係数変更を備える
$\mathfrak M^{\mathcal O}_{\mathrm{cub}}$の相対模型を与える。
非noetherianな底では大きい$\operatorname{QGr}$を使い，局所noetherian性は要求しない。

\begin{proposition}[四係数空間からのモジュライ射の monomorphism]
\label{prop:rigid-four-family-monomorphism}
上の入力と compact 版相対比較の下で、四係数族が与える射
\[
 \mathbb A^4\longrightarrow\mathfrak M^{\mathcal O}_{\mathrm{cub}}
\]
は、有限型 $\mathbb C$-スキームの site に制限すると stack の monomorphism である。
すなわち、任意の有限型基底$T$と二つの係数写像$f,g:T\to\mathbb A^4$に対し、
その像の間の射の層は閉等化子スキーム$\operatorname{Eq}_T(f,g)$が表す層である。
特に$T$全体上の射は$f=g$の場合に一意に存在し、それ以外には存在しない。
この等化子は係数関数の差で定義され，非被約な構造も保持する。
\end{proposition}

\begin{proof}
アフィン上で前系と compact 版比較を組み合わせると、
次数スカラー商は二つの係数写像の等化子の層に一致する。
この同定は係数変更と両立するので一般の $T$ へ貼り合わさる。

ここでは stack 化によって新しい射が加わらないことも確認する必要がある。
固定した二つの標準 Gabriel 係数族の相対同値・自然同型は、
補題 \ref{lem:abel-canonical-family-equivalences} により
\emph{任意の fpqc 係数被覆}について降下する。
従って、有限型 site の外に出る fpqc 被覆上で表された局所的な同値も、
既に $T$ 上の同値へ降下する。

また構造対象上で恒等な自然自己同型は自明である。
任意係数における saturation により各標準対象の自己射環は係数環であり、
隣接 Hom は正階数の自由加群なので、その上のスカラーは全て等しい。
構造対象上で恒等なら全スカラーが一となり、生成族上の自然性から圏全体で恒等となる。
したがって標識付き Hom 群oid の同型類は既に fpqc 層であり、
その stack 化はこの二つの大域模型の間の射を増やさない。
これでモジュライ stack 上でも同じ等化子の記述を得る。
\end{proof}

この monomorphism は、モジュライ問題における完全忠実性の主張である。
開埋込み、モジュライの一つの成分の尽くし、有限型 atlas、
対角の代数性または任意の基底上での比較定理を追加で意味するものではない。

\section{有限位数楕円境界の中心商に必要な正則性の証明}
\label{CSBGsec:pair-finite-order-background}
この節では、固定した AS 正則な周囲の中心三次商について、
tails 大域次元が3となる係数範囲を構成する。
二つの中心 chart の正則性を楕円境界に沿って貼り合わせ、
有限位数の場合は Rains の中心持上げと有限中心代数の Bertini 定理から非空開集合を得る。
以下で$B$は標数零の体上の左右noetherian connected graded AS正則代数、
$W\in Z(B)$は正次数の正則斉次元、$A=B/(W)$とする。
$\mathcal P_A=\operatorname{thick}\{\pi A(n):n\in\mathbb Z\}$、
$\mathcal Q_A=\mathrm D^b(\qgr A)/\mathcal P_A$と置く。
有限graded自由行列因子化と次数付き特異圏の比較には
\cite[arXiv v1, Theorems~4.7, 5.5]{CSBGCCKM}を用いる。
右加群には反対環を適用し、中心元のnormalizing automorphismは恒等である。
従って有限syzygy後のcoherent対象を二周期因子化で扱える。
以下の補題の引用入力と仮定を保ち、一般の非正則底へ拡張しない。
\begin{lemma}[有限射影因子化の縮約と分母払い]
\label{CSBGlem:gluing-contraction}
$2$が可逆な左右noetherian環$S$と中心正則元$w$を取り、
有限生成射影加群の行列因子化を$d^2=w$と書く。
$S/(w)$上のcokernelが有限射影次元なら、この因子化は縮約可能である。
さらに中心正則元$u$で局所化した後に縮約可能なら、
ある$N$について$u^N\mathrm{id}=dK+Kd$となる。
次数付きの場合も、因子化の次数shiftを含めて同じ主張が成立する。
\end{lemma}
\begin{proof}
$\bar S=S/(w)$上の二周期分解では、cokernelの有限射影次元により
十分高い syzygy が射影的である。周期性から全ての周期 syzygy が射影的となる。
従って分解は分裂し、
$\bar S$上で$d\bar H+\bar Hd=1$となる奇写像を取れる。
有限射影性からこれを$S$上の$H$へ持ち上げると
$dH+Hd=1+wT$となる。
左辺は$d$と可換なので、$w$の正則性により$[d,T]=0$。
ゆえに$H'=H-\frac12dT$は$dH'+H'd=1$を満たす。
局所化後の$H'$の有限個の成分の分母を$u^N$で払い、
$S\hookrightarrow S[u^{-1}]$の単射性を使えば最後の等式を得る。
\end{proof}

\begin{lemma}[有限大域次元からtwistsによる生成へ]
\label{CSBGlem:global-finite-to-perfect}
$A$は次数1生成で、$\operatorname{qgr}A$が有限大域次元を持つなら、
$\mathrm D^b(\operatorname{qgr}A)=\mathcal P_A$である。
\end{lemma}
\begin{proof}
大域次元を$e$とし、graded自由表示を繰り返して
$0\to K\to L_e\to\cdots\to L_0\to M\to0$を取る。
各$L_i$は$\pi A(n)$の有限直和である。
最後のYoneda拡大は$\operatorname{Ext}^{e+1}(M,K)=0$なので、
$M$は有限複体$[L_e\to\cdots\to L_0]$の直和因子である。
line object自体が射影的であるとは仮定していない。
\end{proof}

\begin{proposition}[中心零化元による十分条件]
\label{CSBGprop:global-annihilators}
$u,v\in Z(B)_1$が$A$上正則列をなし、
$\mathrm D^b(\operatorname{qgr}A/(u,v))$が構造対象のtwistsで厚く生成されるとする。
正次数の斉次中心元$p_1,\ldots,p_r\in Z(A)$が
次数付き特異圏の全対象を零化し、ある$N$について
$(u,v)^N\subset(p_1,\ldots,p_r)A$なら、$\mathcal Q_A=0$である。
\end{proposition}
\begin{proof}
graded lift $M$のKoszul複体$K(p_1,\ldots,p_r;M)$のcohomologyは
$(p_1,\ldots,p_r)$に殺され、従って$(u,v)$の冪に殺される。
有限filtrationの各商は$A/(u,v)$加群である。
正則列のKoszul分解から$\pi(A/(u,v))(n)\in\mathcal P_A$なので、
境界の生成性によりこのKoszul複体は$\mathcal P_A$に入る。
一方、$\mathrm D^b(\operatorname{gr}A)\to\mathcal Q_A$は
次数付きperfect複体を殺すから次数付き特異圏を経由する。
従って各$p_i$は$\mathcal Q_A$で零作用し、反復coneは
$M$のshiftとtwistの直和へ分裂する。ゆえに$M$自身も零となる。
\end{proof}

\begin{proposition}[二つの中心chartの貼合せ]
\label{CSBGprop:global-two-charts}
$u,v\in Z(B)_1$が$B$上正則で$A$上正則列をなし、
$A[u^{-1}]_0,A[v^{-1}]_0$が有限大域次元を持つとする。
$\operatorname{qgr}A/(u,v)$の有界導来圏が構造対象のtwistsで厚く生成されるなら、
$\mathcal Q_A=0$である。
\end{proposition}
\begin{proof}
各 coherent 対象の graded lift に有限回の syzygy を取ると、
行列因子化の cokernel を得る。途中の有限 graded 自由部分は
$\mathcal Q_A$ で消えるので、元の対象はこの cokernel の shift と同型である。
中心一次元を逆にした次数零部で因子化を考えると、
chart の有限大域次元により cokernel の射影次元は有限となり、因子化は縮約する。
補題\ref{CSBGlem:gluing-contraction}の分母払いから、対象ごとに選んだ
$u^r,v^s$が$\mathcal Q_A$で零作用する。
二重Koszul coneのcohomologyは$(u^r,v^s)$に殺され、
$(u,v)^{r+s-1}\subset(u^r,v^s)$から境界の有限filtrationを持つ。
前命題と同じ議論でconeはperfectである。
中心作用は自然変換なので、最初の零coneの直和分解とも可換し、
二重coneは元の対象の四つのshift・twistの直和となる。
従って元の対象は零である。一様な$r,s$は要求しない。
\end{proof}

\begin{lemma}[perfect性とSerre双対性からの正確な次元]
\label{CSBGlem:gluing-serre-dimension}
$A$は左右noetherian connected graded AS--Gorensteinで、
balanced dualizing complexが$A_\nu(-\ell)[d+1]$（$d\ge0$）で、
$\pi A\ne0$とする。
$\mathrm D^b(\operatorname{qgr}A)=\mathcal P_A$なら
$\operatorname{gldim}\operatorname{qgr}A=d$である。
\end{lemma}
\begin{proof}
\cite[Theorem~4.2, Corollary~4.3]{YZ}は左右の$\chi$と
有限cohomological dimensionを与える。
特に$\pi A(n)$は有限射影次元を持つ。
Chen\cite[Theorem~3.10]{CSBGChenGeneralizedSerre}の一般化Serre双対性は
有限射影次元の対象に適用でき、その関手は
$S(M)=M_\nu(-\ell)[d]$となる。
仮定により全coherent対象が適用範囲に入るので、
\[
 \operatorname{Ext}^i(M,N)^*
 \simeq\operatorname{Ext}^{d-i}(N,M_\nu(-\ell)).
\]
$i>d$では右辺が負次数ゆえ零。
また$M=\pi A$、$N=M_\nu(-\ell)$と取ると
$\operatorname{Ext}^d(M,N)^*\simeq\operatorname{Hom}(N,N)\ne0$なので等号となる。
\end{proof}

\begin{lemma}[超越的中心値の有限大域次元]
\label{CSBGlem:global-transcendental}
$k_0\subset\CC$、$R_0$は$k_0$-代数、$z\in Z(R_0)$は$k_0$上超越的とする。
$R=R_0\otimes_{k_0}\CC$が左右noetherianで$\operatorname{gldim}R\le n$なら、
$k_0$上超越的な全ての$\lambda\in\CC$について
$\operatorname{gldim}R/(z-\lambda)\le n$である。
\end{lemma}
\begin{proof}
$\lambda$を含む$\CC/k_0$の超越基底$T$を選び、$K=k_0(T\setminus\{\lambda\})$とする。
体拡大は自由かつ忠実平坦で、$K$-線形な$\CC\to K$の$1\mapsto1$という分裂がある。
従って大域次元の上界とnoetherian性は$R_K=R_0\otimes K$へ降下する。
$z$の冪は$K$上でも独立なので、中心局所化
$R_\eta=R_K\otimes_{K[z]}K(z)$も大域次元高々$n$。
$z\mapsto\lambda$で$K(z)=K(\lambda)$とみなすと
$R_\eta\otimes_{K(z)}\CC=R/(z-\lambda)$である。
最後の体拡大は代数的分離拡大で、右辺は$R$の商ゆえnoetherian。
\cite[Lemmas~3.2, 3.3]{RSVdB}の仮定が揃い、上界が保たれる。
任意の超越的体拡大が有限大域次元を保つと仮定したわけではない。
\end{proof}

\begin{theorem}[固定した良いambientでのvery-generalな正則性]
\label{CSBGthm:global-very-general}
$B$は中心一次元$u,v$を持つ左右noetherian AS正則五次元代数とし、
$G\in Z(B)_3$、$\operatorname{gr}_{x,y,z}B=S\otimes\CC[u,v]$とする。
$S$は整域、$\operatorname{gr}G=\gamma\in Z(S)_3$は非零正則元とする。
全ての$f\in\CC[u,v]_3$について$B/(G+f)$は$\CC[u,v]$上自由であり、
$\operatorname{qgr}(S/(\gamma))$は滑らかな楕円曲線のcoherent圏と仮定する。
有限表示の定義体$k_0$に対し、$a,d$が$k_0(b,c)$上代数的独立なら
\[
 \operatorname{gldim}\operatorname{qgr}
 B/(G+au^3+bu^2v+cuv^2+dv^3)=3.
\]
\end{theorem}
\begin{proof}
$D_u=B[u^{-1}]_0$はnoetherianで有限大域次元を持つ。
実際、中心局所化は大域次元を増やさず、
$B[u^{-1}]=D_u[u,u^{-1}]$から制限・誘導と評価$u=1$で上界を降下できる。
$t=v/u$とすると商の$u$-chartは
$D_u/(G/u^3+bt+ct^2+dt^3+a)$。
括弧内の$a$を除く中心元は$k_0(b,c,d)$上定義され、
最高主変数成分が非零正則なので、その冪は独立である。
前補題でchartの有限大域次元を得る。
$v$-chartも$d$を中心値として同様に扱える。
命題\ref{CSBGprop:global-two-charts}から有界 coherent 導来圏は twists で厚く生成される。
ここで $W=G+au^3+bu^2v+cuv^2+dv^3$ の最高成分は正則元 $\gamma$ なので、
$W$ 自身も正則である。正則中心商の Rees の補題
\cite[arXiv v1, Lemma~4.4(3)]{CSBGCCKM}により $A=B/(W)$ は
AS--Gorenstein 四次元であり、\cite[Corollary~4.3 とその直前]{YZ}から
balanced dualizing complex は $A_\nu(-\ell)[4]$ の形を持つ。
$A$ の $\CC[u,v]$-自由性から $\pi A\ne0$ でもある。
従って補題\ref{CSBGlem:gluing-serre-dimension}を $d=3$ として適用できる。
\end{proof}

\begin{lemma}[被約な係数環上の中心持上げ]
\label{CSBGlem:global-center-lifting}
$k$ を代数閉標数零体、$R_0$ を被約有限型 $k$-代数とし、
$U$ を非負の増大する exhaustive filtration $F_\bullet$ を持つ $R_0$-代数とする。
$R_0$ は中心に入り、$F_jU$ は有限自由、
$\operatorname{gr}U=H\otimes R_0$ が閉点への基底変更と両立すると仮定する。
$H=B(E,\mathcal L,\tau)$は滑らかな楕円曲線$E$，豊富な直線束$\mathcal L$，
有限位数の平行移動$\tau$によるtwisted homogeneous coordinate ringとする。
このとき $Z(H)$ の各斉次元は $Z(U)$ に持ち上がる。
さらに $U$ は有限型可換中心部分代数上有限である。
\end{lemma}

\begin{proof}
$H$ の有限個の斉次代数生成元を $U$ に持ち上げて $q_i$ と書く。
最高次 symbol を消す次数帰納法により、$q_i$ は $U$ を $R_0$-代数として生成する。
各整数 $j\geq 0$ に対し、有限自由 $R_0$-加群間の交換子写像
\[
 M_j:F_jU\longrightarrow\bigoplus_i F_{j+\deg q_i}U,
 \qquad w\longmapsto ([w,q_i])_i
\]
を取る。その核は $F_jU\cap Z(U)$ である。
Rains の定理 \cite[Theorem 2.13]{Rains}
は、標数零で有限位数の平行移動によりねじられた楕円代数の
\emph{任意の filtered deformation} の中心について
$\operatorname{gr}Z(U_s)=Z(H)$ を与える。
従って全ての閉点 $s$ で $\dim\ker(M_j\otimes k(s))$
は $\sum_{i\leq j}\dim Z(H)_i$ に等しい。
ここは単なる「associated graded が中心上有限」という一般論ではない。
楕円代数および有限位数平行移動という Rains の仮定を用いている。

$M_j$ の rank を $r_j$ とする。
$(r_j+1)$ 次 minor は全閉点で消える。
$R_0$ が被約かつ Jacobson なので、これらの minor は環の元として零である。
$r_j$ 次 minor の主開集合は $\operatorname{Spec}R_0$ を覆う。
各主開上で可逆な minor を用いて行列を変形すると、
$M_j$ は単位行列と零行列の直和になる。
従って $\ker M_j$ は有限射影で、その形成は任意の基底変更と可換する。

最高次 symbol は写像
\[
 \ker M_j\longrightarrow Z(H)_j\otimes R_0
\]
を与える。
この写像は全閉点で全射であるから、有限生成 cokernel に Nakayama を適用して全射である。
右辺は自由なので、各 $z\in Z(H)_j$ は $Z(U)\cap F_jU$ に持ち上がる。

最後に、有限位数 $r$ の平行移動に対し
$\mathcal L_r=\mathcal L\otimes\tau^*\mathcal L\otimes\cdots
\otimes(\tau^{r-1})^*\mathcal L$ は巡回置換の線型化を持ち、
$E/\langle\tau\rangle$ 上の豊富な直線束へ降下する。
その section ring は、次数を $r$ 倍にして $Z(H)$ に入る。
$H$ の各剰余次数部分は、有限射 $E\to E/\langle\tau\rangle$ の下で
coherent 層の section module と表されるので、$H$ はこの中心 section ring 上有限である。
従って $H$ が有限加群となる有限型斉次中心部分代数 $Z_0\subset Z(H)$ を取れる。
$Z_0$ の有限個の代数生成元と $H$ の有限個の $Z_0$-加群生成元を持ち上げる。
最高次 symbol を順次消して filtration 次数を下げる帰納法により、
持ち上げた中心元と $R_0$ が生成する可換部分代数上で $U$ は有限となる。
\end{proof}


\begin{theorem}[有限位数境界を持つ固定ambientの開な正則範囲]
\label{CSBGthm:global-open}
$k$ を代数閉標数零体とし、定理\ref{CSBGthm:global-very-general}の
良い固定ambientの仮定を $\CC$ の代わりに $k$ 上で満たす $B$ を取る。
さらに $H=S/(\gamma)=B(E,\mathcal L,\tau)$ のtranslationが有限位数とする。
ある非零多項式$P\in k[a,b,c,d]$が存在して、$P\ne0$上では
中心三次商のtails大域次元が3となる。
周囲のPBW/Koszul解消とnoetherian AS条件が基底変更と両立する構成では、
この十分条件は任意の代数閉拡大体上でも成立する。
特に $k=\CC$ で周囲の表示が有理数係数なら、有理数四係数の正則な商がZariski稠密に存在する。
\end{theorem}
\begin{proof}
$D_u=B[u^{-1}]_0$、$t=v/u$、$G_u=G/u^3$とする。
全係数を動かす$u$-chartの全代数は
\[
 \mathcal U=D_u[a,b,c,d]/(G_u+a+bt+ct^2+dt^3)
 \simeq D_u[b,c,d]
\]
なので左右noetherianかつ有限大域次元である。
$R=k[t,a,b,c,d]$をfiltration次数零とすると
$\operatorname{gr}\mathcal U=H\otimes R$で、各filtered pieceは有限自由。
前補題から$\mathcal U$は有限型可換中心部分代数$C_u$上有限である。
この$C_u$は係数環$R_0=k[a,b,c,d]$上も有限型である。
ここで\cite[Proposition~2.1, proof steps (6)--(8)]{RSVdB}を
$(R,S,A)=(R_0,C_u,\mathcal U)$へ適用する。
この命題は、有限型可換中心代数上の加群有限性と全代数の有限大域次元を仮定し、
$C_u$ の正規性や $\mathcal U$ が NCR であることを要求しない。
その証明は、ある非零 $P_u\in R_0$ を逆にした $R'=R_0[P_u^{-1}]$ 上で、
$\mathcal U'=\mathcal U\otimes_{R_0}R'$ の $R'$-平坦性と有限射影解消
\[
 0\longrightarrow Q_m\longrightarrow\cdots\longrightarrow Q_0
 \longrightarrow\mathcal U'\longrightarrow0,
 \qquad Q_i\in\operatorname{proj}
       ((\mathcal U')^{\mathrm{op}}\otimes_{R'}\mathcal U')
\]
を与える。
各 $Q_i$ と $\mathcal U'$ は $R'$-平坦であり、短完全列を順に見れば
全 syzygy も $R'$-平坦である。従って任意の体値 $R'\to L$ へのテンソル後も
これは対角 $\mathcal U'_L$ の有限射影解消である。
この解消は片側 $\mathcal U'_L$-加群として分裂するので、任意の
$\mathcal U'_L$-加群と片側でテンソルしても完全であり、
$L$ が体であることから長さ高々 $m$ の射影解消を得る。
ゆえに全ての幾何学的 fiber で有限大域次元が成立し、体拡大も許される。
$v$-chartでも$d$を消去して$P_v$を得る。
$P=P_uP_v$とし、二chartの貼合せとSerre双対性で結論する。
$k=\CC$ の場合、$\mathbb Q^4$は複素四次元affine空間でもZariski稠密なので、
非空開集合には有理数係数点が稠密に存在する。
この証明は Rains の標数零の中心持上げ、被約 Jacobson 環上の一定階数、
および有限中心代数に対する RSV の一般正則性だけを用い、
very-general 点の存在や $k$ の非可算性に依存しない。
反対環にも同じ議論を適用し、二つの非零多項式を掛ければ左右同時の結論を得る。
\end{proof}

\begin{lemma}[位数6の具体的な楕円曲線計算]\label{CSBGlem:order-six}
$E_1=V(x^3+y^3+z^3+xyz)$で$O=[1:-1:0]$を原点とすると、
$p=[1:-1:-1]$は位数6であり、$2p=[1:0:-1]$である。
\end{lemma}
\begin{proof}
$F=x^3+y^3+z^3+xyz$ と置く。
偏微分が全て零なら、いずれかの座標が零の場合は全座標が零となる。
全座標が非零の場合も、$3x^2=-yz$ とその巡回式を掛けると
$27x^2y^2z^2=-x^2y^2z^2$ となり矛盾する。従って $E_1$ は滑らかである。
$O$ の接線は $z=3(x+y)$ で、
\[
 F(x,y,3(x+y))=28(x+y)^3
\]
だから $O$ は変曲点である。
$p$ の接線 $2x+y+z=0$ については
\[
 F\bigl(-(y+z)/2,y,z\bigr)=\frac78(y+z)(y-z)^2.
\]
従って接線の第三交点は $r=[0:1:-1]$ であり、$p$ での交点重複度は正確に2である。
$r,O,q=[1:0:-1]$ は直線 $x+y+z=0$ 上にあるので、
楕円曲線の弦接線法から $2p=-r=q$ となる。
さらに $q$ の接線 $y=3(x+z)$ について
$F(x,3(x+z),z)=28(x+z)^3$ だから、$q\ne O$ は位数3の変曲点である。
従って $p$ の位数は3または6であるが、$p$ は変曲点でないので位数3ではなく、位数6となる。
\end{proof}
$h=-1$の場合は$x,y$の交換が$S_{-1}$と$S_1$の次数付き同型を与えるので、境界translationの位数も同じである。

\section{有限位数境界上の相対 dg smooth proper 族}
\label{sec:dg-smooth-torsion-families}

本節では、構成済みの十七次元 Jacobian atlas の中で、
十六次元の局所閉部分族の非空開上に相対 dg smooth proper 性を証明する。
このような部分族の合併は atlas に稠密である。
十七次元 atlas 自体の非空開は、この結果には含めない。
五線 compact generator の係数上 perfect 性を既存の入力とし、
有限中心の射影模型に対する大域対角の議論を加える。

\subsection{係数上 perfect な dg 代数の一般点からの広がり}
\label{sec:dg-generic-spreading}

以下では dg 代数は結合的かつ単位的であり、テンソル積と enveloping algebra は
すべて導来の意味で用いる。\(R\) 上の dg 代数 \(E\) に対し
\(E^e=E^{\mathrm{op}}\otimes_R^{\mathbf L}E\) と書く。
\(E\) が \(R\)-proper であるとは \(E\in\operatorname{Perf}(R)\)、
\(R\)-smooth であるとは対角加群 \(E\in\operatorname{Perf}(E^e)\)
であることをいう。

\begin{lemma}[一般点からの smoothness の広がり]
\label{lem:dg-generic-spreading}
\(R\) を整域である Noether 環、\(K=\operatorname{Frac}(R)\) とし、
\(E\) を \(E\in\operatorname{Perf}(R)\) を満たす \(R\)-dg 代数とする。
\(E_K=E\otimes_R^{\mathbf L}K\) が \(K\)-smooth ならば、
ある \(0\ne f\in R\) に対し \(E_f\) は \(R_f\)-smooth proper である。
したがって任意の \(R_f\)-代数 \(S\) に対し \(E_f\otimes_{R_f}^{\mathbf L}S\)
は \(S\)-smooth proper である。
\end{lemma}

\begin{proof}
\(D=E^e\) とおく。\(E\) が係数上 perfect なので \(D\) も
\(R\)-perfect である。\(R\) の非零元による局所化の filtered system に対し、
perfect dg 加群は有限段階で定義できる：
\[
 \mathop{\mathrm{colim}}_{0\ne g\in R}\operatorname{Perf}(D_g)
 \simeq \operatorname{Perf}(D_K).
\]
これは To\"en--Vaqui\'e\cite[Lemma~2.10]{ToenVaquieObjects}、または
To\"en\cite[事実 F4]{ToenDefinitionRings}による。この連続性に \(D\) の homotopical finite presentation
は必要ない。有限 cell 加群とその retract の有限個のデータを降下させる性質である。

仮定より \(E_K\in\operatorname{Perf}(D_K)\) であるから、
ある \(g\ne0\) と \(P\in\operatorname{Perf}(D_g)\) を選んで
\(P_K\simeq E_K\) とできる。この同型を \(\alpha_K\) とする。
source \(P\) の perfectness により、任意の \(D_g\)-加群 \(M\) について
\[
 \mathbf R\operatorname{Hom}_{D_g}(P,M)\otimes_{R_g}^{\mathbf L}K
 \simeq
 \mathbf R\operatorname{Hom}_{D_K}(P_K,M_K)
\]
が成り立つ。自由加群 \(D_g[n]\) に対して明らかであり、
有限直和・有限 cone・retract に関して保たれるためである。
ここでは target \(M\) の \(D_g\)-perfectness を仮定していない。
したがって \(M=E_g\) として \(\alpha_K\) の有限個の分母を払い、
さらに局所化し、記号を変えずに書けば射
\(\alpha:P\longrightarrow E_g\) を得る。

その cone \(C\) は \(R_g\)-perfect である。
実際、\(D_g\) は \(R_g\)-perfect なので制限関手は
\(\operatorname{Perf}(D_g)\) を \(\operatorname{Perf}(R_g)\) に送る。
よって \(P\) と \(E_g\) はともに係数上 perfect である。
また \(C_K=0\) である。
係数上 perfect な複体のコホモロジーは有限個の有限生成加群であり、
各 \(H^i(C)\) は \(K\) への局所化後に消える。
有限個の生成元を殺す非零元を掛け合わせると、さらにある局所化上で
\(C=0\) となる。したがってその上で \(E\simeq P\) は
\(E^e\)-perfect、すなわち smooth である。
properness および smoothness の任意の導来基底変換による保存は、
perfect 複体および perfect 対角加群の基底変換から従う。
\end{proof}

\begin{lemma}[体拡大からの降下]
\label{lem:dg-smooth-field-descent}
\(k\subset L\) を体拡大、\(E\) を \(k\)-proper な dg 代数とする。
\(E_L\) が \(L\)-smooth ならば \(E\) は \(k\)-smooth である。
\end{lemma}

\begin{proof}
\(D=E^{\mathrm{op}}\otimes_k E\) とする。
\(L\) をその有限生成 \(k\)-部分代数の filtered colimit とみなし、
上記と同じ perfect 加群の連続性を適用する。
ある有限生成整域 \(A\subset L\) と \(P\in\operatorname{Perf}(D_A)\)
および同型 \(P_L\simeq E_L\) が存在する。
source \(P\) の compactness によって、\(A\) を必要なら大きくして、
この同型を射 \(\alpha:P\to E_A\) として定義できる。
cone \(C\) は \(A\)-perfect であり \(C\otimes_A L=0\) を満たす。
\(\operatorname{Frac}(A)\subset L\) の忠実平坦性から
\(C_{\operatorname{Frac}(A)}=0\)、したがって非零元を一つ逆にすると
\(C=0\) とできる。こうして \(E_A\in\operatorname{Perf}(D_A)\)
となる有限生成 \(k\)-整域 \(A\subset L\) が得られた。

\(\operatorname{Spec}(A)\) の閉点を取り、その剰余体を \(k'\) とする。
Zariski の補題により \([k':k]<\infty\) である。
\(E_{k'}\) は \(D_{k'}\)-perfect であり、
\(D_{k'}\) は \(D\)-加群として有限自由なので、
制限した \(E_{k'}\) は \(D\)-perfect である。
\(\lambda(1)=1\) を満たす \(k\)-線型写像 \(\lambda:k'\to k\) を選ぶと
\[
 E\longrightarrow E\otimes_k k'
   \xrightarrow{\,1\otimes\lambda\,}E
\]
の合成は恒等写像である。いずれも \(D\)-線型なので、
\(E\) は \(D\)-perfect 加群の retract、したがって \(D\)-perfect である。
\end{proof}

\subsection{大域対角を扱うための補題}

\begin{proposition}[左右の正則性から大域 dg smooth 性へ]
\label{prop:seed-finite-central-global-smooth}
$k$ を完全体、$Y$ を有限型分離 $k$-スキーム、$\mathscr A$ を
coherent な $\mathcal O_Y$-代数とする。
左右それぞれについて
\[
 \mathrm D^b(\operatorname{coh}\mathscr A)
   =\operatorname{Perf}(Y,\mathscr A),\qquad
 \mathrm D^b(\operatorname{coh}\mathscr A^{\mathrm{op}})
   =\operatorname{Perf}(Y,\mathscr A^{\mathrm{op}})
\]
と仮定する。すると $\operatorname{Perf}_{\mathrm{dg}}(Y,\mathscr A)$ は
$k$ 上 dg smooth である。さらに $Y$ が $k$ 上 proper なら、この dg 圏は proper である。
\end{proposition}

\begin{proof}
$Y$ は noetherian J-2 スキームである。
Elagin--Lunts--Schn\"urer\cite[Theorem~4.15]{ELSSmoothAlgebras}により
$\mathrm D^b(\operatorname{coh}\mathscr A)$ は classical generator $G$ を持つ。
仮定により $G$ は $\mathscr A$-perfect である。
以下、$\mathscr A$-加群は右加群とする。
双対
\[
 G^\vee=\mathbf R\mathcal Hom_{\mathscr A}(G,\mathscr A)
\]
は右 $\mathscr A^{\mathrm{op}}$-perfect であり、perfect 加群の双対は反同値なので、
第二の仮定により $G^\vee$ は
$\mathrm D^b(\operatorname{coh}\mathscr A^{\mathrm{op}})$ の classical generator となる。

$Y\times_kY$ 上で
\[
 \mathscr B=\mathscr A^{\mathrm{op}}\boxtimes_k\mathscr A,
 \qquad K=G^\vee\boxtimes_kG
\]
と置く。Elagin--Lunts--Schn\"urer\cite[Theorem~4.18]{ELSSmoothAlgebras}は、完全体上の
noetherian J-2 スキームと coherent 代数に対して、
\emph{任意に選んだ}二つの classical generator の外積が
外積代数の有界 coherent 導来圏を厚く生成することを述べている。
従って
\[
 \mathrm D^b(\operatorname{coh}\mathscr B)=\operatorname{thick}(K).
\]
分離性により対角 $\delta:Y\hookrightarrow Y\times_kY$ は閉埋込みであり、
$\delta_*\mathscr A$ は coherent な右 $\mathscr B$-加群である。
右作用は、局所的に
$m\cdot(a^{\mathrm{op}}\otimes b)=amb$ と与える。
ゆえに $\delta_*\mathscr A\in\operatorname{thick}(K)$ である。

$E=\mathbf R\operatorname{Hom}_{\mathscr A}(G,G)$ と置く。
有限 affine 被覆の \v Cech 模型と perfect 加群の tensor--Hom 恒等式から
\[
 \mathbf R\operatorname{Hom}_{\mathscr B}(K,K)
       \simeq E^{\mathrm{op}}\otimes_k E,
 \qquad
 \mathbf R\operatorname{Hom}_{\mathscr B}(K,\delta_*\mathscr A)
       \simeq E
\]
を得る。後者は $E^{\mathrm{op}}\otimes_kE$-加群の同型であり、右辺は対角作用を持つ。
次数零では $f\cdot(e_1^{\mathrm{op}}\otimes e_2)=e_1fe_2$ であり、
一般の次数では通常の Koszul 符号を入れる。
この比較は局所的に $G$ を有限生成射影 $\mathscr A$-加群の有界複体で表して確かめ、
\v Cech 複体で大域化できる。
$\mathscr A$ が $\mathcal O_Y$ 上平坦であるとは仮定していない。
外積の基底が体 $k$ であることと、$\mathscr A$-perfect 性を用いる。

第一の同型により、$\mathbf R\operatorname{Hom}_{\mathscr B}(K,-)$ は
$\operatorname{thick}(K)$ を
$\operatorname{Perf}(E^{\mathrm{op}}\otimes_kE)$ に移す。
第二の同型と $\delta_*\mathscr A\in\operatorname{thick}(K)$ から、
対角 $E$ は $E^{\mathrm{op}}\otimes_kE$ 上 perfect である。
Morita 不変性により、これが所望の dg smooth 性である。

$Y$ が proper なら、二つの perfect $\mathscr A$-複体の導来内部 Hom は
有界 coherent $\mathcal O_Y$-複体である。
その大域切断は有界かつ各次数有限次元なので、proper 性も従う。
\end{proof}

\begin{remark}[局所対角と大域対角の役割]
有限中心の affine chart の対角が perfect であることだけを、
大域 dg smooth 性の証明と同一視していない。
上の証明で大域化を担う入力は、外積 classical generator の生成定理と
$\delta_*\mathscr A$ に対する dg Morita 比較である。
この命題の一般形の初出は未確認であり、
主張の新規性は既知の Elagin--Lunts--Schn\"urer の生成定理と区別する。
\end{remark}

\subsection{有限位数境界の点は二係数開集合に稠密である}

$T_0=\operatorname{Spec}\mathbb C[a,b,\Delta^{-1}]$ とし、
現行原稿の規約
\[
 p=2-b,\qquad q=-(2+b),\qquad r=-6a
\]
を用いる。$p-q=4$ なので
\[
 \mathbb A^2_{a,b}\xrightarrow{\sim}
 \{[p:q:r]\in\mathbb P^2\mid p-q\ne0\},
 \qquad
 a=-\frac{2r}{3(p-q)},\quad
 b=-\frac{2(p+q)}{p-q}.
\]
従って $T_0$ はこの射影平面の非空開集合と同一視できる。

\begin{lemma}[torsion 境界パラメータの稠密性]
\label{lem:seed-torsion-boundary-dense}
$T_0(\mathbb C)$ において、楕円境界の自己同型 $\sigma$ が
有限位数を持つ点の集合は Zariski 稠密である。
\end{lemma}

\begin{proof}
$pqr$ が可逆な開集合上で
\[
 \mu=\frac{p^3+q^3+r^3}{pqr}
\]
を取る。$\Delta\ne0$ の条件は $\mu^3\ne27$ を含む。
$\mu$ のファイバーは、滑らかな Hesse 三次曲線
\[
 E_\mu:\quad X^3+Y^3+Z^3-\mu XYZ=0
\]
の開部分である。
Sklyanin の同定により、$O=[1:-1:0]$ を原点とすると、
そのパラメータにおける $\sigma$ は
$P=[p:q:r]\in E_\mu$ による平行移動である。
ここで関係の規約は
$pyz+qzy+rx^2=0$ とその巡回式であり、
Walton--Wang--Yakimov\cite[Definition--Lemma~2.13]{WWYSklyaninPoisson}と一致する。

任意の非空開集合 $W\subset T_0$ を取る。
ある閉点 $w\in W$ と $\mu_0=\mu(w)$ を選ぶと、
$W\cap E_{\mu_0}$ は滑らかな既約楕円曲線 $E_{\mu_0}$ の非空開集合である。
複素楕円曲線の torsion 点は無限集合であり、
既約曲線の真の閉集合は有限なので Zariski 稠密である。
従って $W\cap E_{\mu_0}$ に torsion 点 $P$ が存在する。
その平行移動 $\sigma$ は有限位数である。
これで任意の非空 $W$ との交わりが得られた。
\end{proof}

\begin{remark}
torsion の位数を一つに固定していない。
任意の指定された有限個の位数を除いても、
各楕円曲線上には無限個の torsion 点が残るので同じ証明が成立する。
特に $\Delta$ の小位数除外条件は障害にならない。
一方、位数6の既存例 $B_*$ 自身を使うことはできない。
その最高次パラメータは $a=1/3,b=0$、従って
$(p,q,r)=(2,-2,-2)$ であり、$q^3-r^3=0$ のため $T_0$ の外にある。
\end{remark}

\subsection{固定した有限位数 ambient の相対開}

\begin{proposition}[有限位数固定 ambient 上の相対 smooth proper 開]
\label{prop:seed-fixed-ambient-relative-smooth}
$t\in T_B(\C)$ を境界 translation が有限位数となる固定 ambient 点とし、
$A_{t,g}=B_t/(\mathcal G_t+\sum_{j=0}^3g_ju^{3-j}v^j)$ と置く。
$P=\mathbb C[g_0,g_1,g_2,g_3]$ とする。
ある非零元 $h\in P$ が存在し、任意の可換 $P_h$-代数 $S$ に対して、
\[
 \mathrm D(\operatorname{QGr}A_{t,S})^c
\]
は $S$ 上 dg smooth proper となる。
この同定と dg Morita 模型は任意の基底変更と両立する。

\end{proposition}

\begin{proof}
系\ref{cor:relative-cubic-compact-proper}と
命題\ref{prop:general17-relative-cy-koszul}により、
\[
 T=\bigoplus_{i=0}^{4}\mathcal O(i),\qquad
 E=\mathbf R\operatorname{Hom}(T,T)
\]
を取ると、$T$ は compact generator、$E$ は $P$-perfect であり、
全ての可換 $P$-代数 $S$ について
\[
 E\otimes_P^{\mathbf L}S
   \simeq\mathbf R\operatorname{Hom}(T_S,T_S)
\]
が成り立つ。この比較は乗法を保つ canonical な dg 比較である。

$K=\operatorname{Frac}P$ とする。
固定した有限位数境界についての全四係数族は、
定理\ref{thm:central-model-finite} により
有限型次数付き可換 $P$-代数 $C$ 上有限である。
従って $Y=\operatorname{Proj}C^{(\ell)}$ と
$\mathscr A=\widetilde{A^{(\ell)}}$ は
$K$ と $\overline K$ への基底変更後にも使える。
$C_S\to\mathscr A_S$ の単射性や全中心の基底変更可換性は不要であり、
中心作用と有限加群性だけを用いる。

定理\ref{CSBGthm:global-open}を左右に適用して得た二つの非零多項式は、
generic point $\operatorname{Spec}\overline K\to\mathbb A^4$ ではともに可逆である。
同定理の幾何学的基底変更版から、
$A_{t,\overline K}$ の左右 tails 大域次元はともに3である。
左右の tails 大域次元が有限であるから、
補題\ref{CSBGlem:global-finite-to-perfect}により有界 coherent 導来圏は
全 twist で厚く生成される。各 twist は局所 $\mathscr A$-perfect なので、
左右とも $\mathrm D^b(\operatorname{coh}\mathscr A)=\operatorname{Perf}(Y,\mathscr A)$ となる。
命題\ref{prop:seed-finite-central-global-smooth}を
$(Y_{\overline K},\mathscr A_{\overline K})$ に適用して
$E_{\overline K}$ の dg smooth 性を得る。
補題\ref{lem:dg-smooth-field-descent}で $E_K$ へ降下する。

補題\ref{lem:dg-generic-spreading}を整域 $P$ と $P$-perfect dg 代数 $E$ に適用すれば、
ある $0\ne h\in P$ について
\[
 E_h\in\operatorname{Perf}
    (E_h^{\mathrm{op}}\otimes_{P_h}^{\mathbf L}E_h)
\]
となる。任意の $P_h\to S$ の導来基底変更はこの perfect 対角を保つ。
$E_h$ 自身の $P_h$-perfect 性も基底変更で保たれるので、
$E\otimes_P^{\mathbf L}S$ は $S$ 上 smooth proper である。
普遍 generator とその endomorphism の基底変更比較により、
これを $\mathrm D(\operatorname{QGr}A_{t,S})^c$ へ移す。
\end{proof}

\begin{remark}
$t\in T_B^e$ の場合、この相対四次元開は新しい十七次元 atlas の中で
ambient $t$ を固定した fiber 内の非空開集合である。
十七次元 atlas 自体の非空開集合とは区別する。
\end{remark}

\subsection{十五次元および十六次元の相対 smooth proper 族}

\begin{proposition}[有限位数境界を固定した全混合変形]
\label{prop:seed-torsion-base-relative-smooth}
$Z\subset T_0$ を integral affine な局所閉部分スキームとし、
その幾何学的一般点における境界 translation が有限位数を持つと仮定する。
十一個の mixed 係数と四個の cubic 係数を全て動かす族
\[
 X_Z=Z\times\mathbb A^{11}_{\alpha,\beta,m_{ij}}
           \times\mathbb A^4_{g_0,g_1,g_2,g_3}
\]
には、非空主開集合 $U_Z$ があり、その上の compact tails は
相対 dg smooth proper である。
この主張は $U_Z$ からの任意の可換係数環への基底変更に両立する。
従って $Z$ が有限位数境界の複素点なら、
十五次元 affine 空間の非空開上の相対 dg smooth proper 族が得られる。
\end{proposition}

\begin{proof}
$R_0=\Gamma(Z,\mathcal O_Z)$、
$R=R_0[\alpha,\beta,m_{ij},g_0,g_1,g_2,g_3]$ と置く。
この $R$ は有限型 $\mathbb C$-整域である。
普遍五線 generator の dg 自己準同型代数 $E_R$ は $R$-perfect であり、
任意の基底変更と両立する。
\[
 F=\operatorname{Frac}(R_0)(\alpha,\beta,m_{ij}),\qquad
 K=F(g_0,g_1,g_2,g_3)=\operatorname{Frac}(R)
\]
と置き、$\overline K$ の中で $F$ の代数閉包 $\overline F$ を選ぶ。
$\overline F$ 上の ambient は固定され、
その境界 translation は仮定により有限位数である。

定理\ref{CSBGthm:global-open}は任意の代数閉標数零体上で成立する。
実際、その有限中心入力は Rains\cite[Theorem~2.13]{Rains}、
開性入力は Rennemo--Segal--Van den Bergh\cite[Proposition~2.1]{RSVdB}であり、
両者ともこの体の仮定で成立する。
二つの中心 chart、行列因子化、境界の twists 生成、Serre 双対性の議論も
\(\mathbb C\) に固有の位相的性質を使わない。
従って $B_{\overline F}$ とその反対環にそれぞれ適用すると、
四係数空間上の二つの非零正則性多項式を得る。

$g_0,g_1,g_2,g_3$ は $\overline F$ 上も代数的独立なので、
この二つの多項式は幾何学的一般点で消えない。
従って $A_{\overline K}$ の左右 tails 大域次元は3である。
被約 Jacobson 係数環の交換子行列を使う
補題\ref{CSBGlem:global-center-lifting}と定理\ref{thm:central-model-finite}の証明も
この体上で成立する。従って有限中心射影模型が存在する。
前命題と同じ左右の twist 生成により、その有界 coherent 導来圏は局所 perfect 圏に一致するので、
命題\ref{prop:seed-finite-central-global-smooth}から
$E_{\overline K}$ は smooth である。
補題\ref{lem:dg-smooth-field-descent}で $E_K$ へ降下し、
補題\ref{lem:dg-generic-spreading}で $R_h$ 上へ延ばす。
$U_Z=D(h)$ とすれば、相対 smooth proper 性と任意の底変更への両立が従う。
全 $X_Z$ 上で中心を一様に選ぶことは、この証明の仮定ではない。
\end{proof}

\begin{theorem}[十六次元 smooth proper 族の稠密な集合]
\label{thm:seed-sixteen-dimensional-dense-families}
十七次元 atlas $T_B^e\times\mathbb A^4$ の任意の非空開集合 $W$ に対し、
smooth integral affine 曲線 $Z\subset T_0$ と、
\[
 V\subset (Z\times\mathbb A^{15})\cap W
\]
なる非空開部分スキームが存在して、
$V$ 上の compact tails は相対 dg smooth proper となる。
$V$ は十六次元であり、任意の可換係数環への基底変更に両立する。
従ってこの atlas では、十六次元の相対 smooth proper 族の像の合併が稠密である。
\end{theorem}

\begin{proof}
Hesse elliptic scheme を
\[
 \mathcal E\longrightarrow M
   =\operatorname{Spec}\mathbb C[\mu,(\mu^3-27)^{-1}],
 \qquad X^3+Y^3+Z^3-\mu XYZ=0,
 \qquad O=[1:-1:0]
\]
とする。
$pqr\ne0$ 上では $\mu=(p^3+q^3+r^3)/(pqr)$ によって
$\mu$ が一意に決まるので、
$T_0$ は $\mathcal E$ の非空開部分スキームと同一視できる。
正整数 $n$ に対し、乗法写像 $[n]:\mathcal E\to\mathcal E$ の
相対微分は $n$ 倍なので可逆であり、滑らかな $M$-曲線間の Jacobian 判定から
この射は \emph{\'etale} である。各幾何学的ファイバーの乗法写像は有限である
\cite[Lemmas~39.9.8,39.9.9]{StacksEllipticTorsion}。
従って $[n]$ は proper かつ quasi-finite、ゆえに有限であり、
その kernel $\mathcal E[n]$ は $M$ 上有限 \emph{\'etale} である。
従って $\mathcal E[n]\cap T_0$ の各非空連結成分は
smooth integral affine 曲線である。
affine 性は、$\mathcal E[n]\cap T_0$ が affine な $T_0$ の閉部分スキームであり、
smooth 曲線の連結成分が開かつ閉であることから従う。

$W$ の $T_0$ への射影像は非空開集合である。
補題\ref{lem:seed-torsion-boundary-dense}により、そこに torsion 点 $P$ を選べる。
$nP=0$ となる $n$ を取り、
$P$ を含む $\mathcal E[n]\cap T_0$ の連結成分を $Z$ とする。
その全ての点で $\sigma^n=1$ なので、
$Z$ は命題\ref{prop:seed-torsion-base-relative-smooth}の仮定を満たす。

$P$ は $W$ の射影像に入るので、
$W\cap(Z\times\mathbb A^{15})$ は空でない。
これは integral な十六次元スキーム $Z\times\mathbb A^{15}$ の開部分である。
命題\ref{prop:seed-torsion-base-relative-smooth}が与える非空主開 $U_Z$ と交わらせ、
\[
 V=U_Z\cap W
\]
とすれば、非空性、次元、相対 smooth proper 性を得る。
\end{proof}

\begin{remark}[十七次元の開性との区別]
前定理では、各 $Z$ 上で有限位数の上界 $n$ を固定している。
得られる十六次元族の合併は稠密であるが、
十七次元 atlas の generic point を含むとは述べていない。
十七次元の相対 dg smooth proper 非空開の構成には、
一般の無限位数境界を扱う別の議論が必要である。
これは第\ref{sec:kuznetsov-moduli-morphism}節で残余圏の幾何性から実現する。
\end{remark}


\begin{corollary}[smooth proper 点の稠密性]
\label{thm:seed-smooth-points-dense}
$T_B^e\times\A^4$ および $\mathfrak C_{\mathrm J}$ の中で、
compact tails が dg smooth proper となる複素点の集合は Zariski 稠密である。
\end{corollary}
\begin{proof}
前定理で得た非空な有限型 $\C$-スキーム $V$ は複素点を持つ。
従って atlas の任意の非空開集合はそのような点を含む。
atlas の全射性から、$\mathfrak C_{\mathrm J}$ の非空開部分の逆像も非空であり、
同じ主張が stack 上で従う。
\end{proof}

\begin{remark}[結論の範囲と既知入力]
外積 generator の定理、perfect dg 加群の filtered colimit、
楕円曲線の torsion と既存の有限中心 Bertini は既知の入力である。
本節の具体的な Jacobian 族への適用と、その相対十六次元族の初出は未確認である。
$V$ は atlas 内の局所閉部分スキームであり、その像が独立に閉部分 stack を
定めることや、異なる torsion 曲線の像が互いに異なることまでは主張しない。
十七次元の一般点の dg smooth 性とその非空開は
定理\ref{thm:kuznetsov-cubic-moduli-morphism}で別に証明する。
全 smooth locus の開性と、全 Abel モジュライ内での成分性は依然として主張しない。
\end{remark}

\section{17係数正規形の正確な行列証明書}\label{app:rank}
係数を $\mathbb Q$ 上で計算する．単項式基底は
\[
 x^iy^jz^ku^\ell v^{3-i-j-k-\ell},\quad
 i,j,k,\ell\ge0,\quad i+j+k+\ell\le3
\]
を辞書式の添字順で並べる．$i+j+k>0$ の31単項式だけを行として取り出す．
19個の群方向は $M_{ij}$（$1\le i,j\le3$）９個，
$N_{ij}$（$1\le i,j\le2$）４個，
$C_{ij}$（$1\le i\le3,1\le j\le2$）６個の順に
\eqref{eq:Lie} の係数微分を並べる．このうち最初の１列を除く．
13個の正規形方向は
\[
 x^3+y^3+z^3,\ xyz,\ vy^2,\ (u+v)z^2,\quad
 x(u^2,uv,v^2),\ y(u^2,uv,v^2),\ z(u^2,uv,v^2)
\]
の順である．得られる $31\times31$ 行列の行列式は
\[
 -5538243375000=-2^3\cdot3^5\cdot5^6\cdot7^2\cdot61^2.
\]
付属の \nolinkurl{computations/independent/verify_generic17_slice.py} は \eqref{eq:F0} に対して
これを厳密整数計算で再現する．中心三次方向
$u^3,u^2v,uv^2,v^3$ を最後に付け足すと，
$35\times35$ 行列式も同じ非零値となる．
したがって binary cubic に入る量子的な正規化補正の詳細を仮定せずに
横断性を結論できる．


\begin{thebibliography}{99}
\bibitem{BehrendNoohiModuli}
K.~Behrend and B.~Noohi,
\emph{Moduli of non-commutative polarized schemes},
Math. Ann. \textbf{371} (2018), 1375--1408.
\href{https://doi.org/10.1007/s00208-017-1624-1}{doi:10.1007/s00208-017-1624-1}.
本稿では\S\S2.3,2.5の乗法空間・scalar gerbe・有限切断を参照する。

\bibitem{AOUNCPlanes}
T.~Abdelgadir, S.~Okawa and K.~Ueda,
\emph{Compact moduli of noncommutative projective planes},
\href{https://arxiv.org/abs/1411.7770}{arXiv:1411.7770} (2014).
参照箇所：\S\S3--5, Theorem~7.3, Appendix~A, Corollaries~A.7 and~A.10.

\bibitem{SierraGAlgebras}
S.~J.~Sierra,
\emph{$G$-algebras, twistings, and equivalences of graded categories},
Algebr. Represent. Theory \textbf{14} (2011), 377--390.
\href{https://doi.org/10.1007/s10468-009-9193-y}{doi:10.1007/s10468-009-9193-y};
\href{https://arxiv.org/abs/math/0608791}{arXiv:math/0608791}.
参照箇所：Proposition~3.3, Theorem~3.6, Corollary~4.4.

\bibitem{RRZHomological}
M.~Reyes, D.~Rogalski and J.~J.~Zhang,
\emph{Skew Calabi--Yau algebras and homological identities},
Adv. Math. \textbf{264} (2014), 308--354.
\href{https://doi.org/10.1016/j.aim.2014.07.010}{doi:10.1016/j.aim.2014.07.010};
\href{https://arxiv.org/abs/1302.0437}{arXiv:1302.0437}.
参照箇所：Theorems~0.3,6.3と後者の証明。

\bibitem{LowenRinged}
W.~Lowen,
\emph{Algebroid prestacks and deformations of ringed spaces},
Trans. Amer. Math. Soc. \textbf{360} (2008), no.~3, 1631--1660.
\href{https://arxiv.org/abs/math/0511197}{arXiv:math/0511197v2}.
Theorems~3.14 and~3.26.

\bibitem{StacksFormal}
The Stacks Project Authors,
\emph{The Stacks Project}.
\href{https://stacks.math.columbia.edu/tag/06T5}{Tag~06T5}
（minimal versal formal objectの一意性）、
\href{https://stacks.math.columbia.edu/tag/07QS}{Tag~07QS}
（excellent rings）、
\href{https://stacks.math.columbia.edu/tag/02KH}{Tag~02KH}
（flat base change）、
\href{https://stacks.math.columbia.edu/tag/02OC}{Tag~02OC}
（formal functions）。2026年11月14日閲覧。

\bibitem{RRGeneralizedAS}
M.~L.~Reyes and D.~Rogalski,
\emph{Graded twisted Calabi--Yau algebras are generalized Artin--Schelter regular},
Nagoya Math. J. \textbf{245} (2022), 100--153.
\href{https://doi.org/10.1017/nmj.2020.32}{doi:10.1017/nmj.2020.32};
\url{https://arxiv.org/abs/1807.10249}.
Theorem~5.15, Lemmas~5.9(1) and~4.5.

\bibitem{YekutieliTwisted}
A.~Yekutieli,
\emph{Twisted deformation quantization of algebraic varieties},
Adv. Math. \textbf{268} (2015), 241--305.
\href{https://doi.org/10.1016/j.aim.2014.09.003}{doi:10.1016/j.aim.2014.09.003};
\href{https://arxiv.org/abs/0905.0488v5}{arXiv:0905.0488v5}.
Corollary~11.2.

\bibitem{YekutieliAffine}
A.~Yekutieli,
\emph{Deformations of affine varieties and the Deligne crossed groupoid},
\url{https://arxiv.org/abs/1209.6243} (2012).

\bibitem{DeDekenLowen}
O.~De~Deken and W.~Lowen,
\emph{Abelian and derived deformations in the presence of
$\mathbb Z$-generating geometric helices},
J. Noncommut. Geom. \textbf{5} (2011), 477--505.
\href{https://doi.org/10.4171/JNCG/83}{doi:10.4171/JNCG/83};
\url{https://ems.press/content/serial-article-files/30497}.
Theorem~4.5 (pp.~495--496), Proposition~4.6 (p.~496), Theorem~5.15.

\bibitem{LowenVdBDeformation}
W.~Lowen and M.~Van den Bergh,
\emph{Deformation theory of abelian categories},
Trans. Amer. Math. Soc. \textbf{358} (2006), 5441--5483.
\href{https://doi.org/10.1090/S0002-9947-06-03871-2}{doi:10.1090/S0002-9947-06-03871-2};
\url{https://arxiv.org/abs/math/0405226}.
出版版Propositions~4.4,4.6,4.8, Definition~8.2,
Theorem~8.8 (p.~5472), Theorem~8.15 (p.~5476).

\bibitem{LowenObstruction}
W.~Lowen,
\emph{Obstruction Theory for Objects in Abelian and Derived Categories},
Comm. Algebra \textbf{33} (2005), 3195--3223.
\href{https://doi.org/10.1081/AGB-200066155}{doi:10.1081/AGB-200066155}.
Theorems~6.11--6.12 (pp.~3218--3219); flat版はcoflat版の双対．

\bibitem{LowenVdBHochschild}
W.~Lowen and M.~Van den Bergh,
\emph{Hochschild cohomology of abelian categories and ringed spaces},
Adv. Math. \textbf{198} (2005), 172--221.
\href{https://doi.org/10.1016/j.aim.2004.11.010}{doi:10.1016/j.aim.2004.11.010}.
Theorem~3.1 (pp.~182--184), Remark~3.5 (p.~184), Theorem~6.1 (p.~200).

\bibitem{BarmeierWangSheaves}
S.~Barmeier and Z.~Wang,
\emph{Deformations of categories of coherent sheaves via quivers with relations},
Algebraic Geometry \textbf{11} (2024), 1--36.
\href{https://doi.org/10.14231/AG-2024-001}{doi:10.14231/AG-2024-001};
\url{https://content.algebraicgeometry.nl/2024-1/2024-1-001.pdf}.
Theorem~1.2 = Theorem~3.11; tiltingを仮定するTheorem~3.18とは区別する．

\bibitem{StacksAbel}
The Stacks Project Authors, \emph{The Stacks Project}.
\href{https://stacks.math.columbia.edu/tag/0FY3}{Tag~0FY3, Lemma~57.3.5}
（双対同型を備えたSerre関手の一意性），
\href{https://stacks.math.columbia.edu/tag/03PK}{Tag~03PK, Lemma~59.28.1}
（Kummer列），
\href{https://stacks.math.columbia.edu/tag/0CJZ}{Tag~0CJZ, Section~21.11}
（gerbeの降下障害），
\href{https://stacks.math.columbia.edu/tag/058G}{Tag~058G, Theorem~10.81.4}
（Lazardの定理），
\href{https://stacks.math.columbia.edu/tag/023M}{Tag~023M, Lemma~35.3.6}
（忠実平坦な拡大Amitsur複体の完全性），
\href{https://stacks.math.columbia.edu/tag/023N}{Tag~023N, Proposition~35.3.9}，
\href{https://stacks.math.columbia.edu/tag/039Y}{Tag~039Y, Remark~35.3.11}
（通常の加群の忠実平坦降下とcosimplicial表示），
\href{https://stacks.math.columbia.edu/tag/058S}{Tag~058S, Proposition~10.83.3}
（有限射影加群の忠実平坦降下），
\href{https://stacks.math.columbia.edu/tag/0FCL}{Tag~0FCL, Lemma~13.17.4}
（Serre部分圏の有界上方導来圏との比較）．2026年11月19日閲覧．

\bibitem{CFT}
A.~S.~Cattaneo, G.~Felder and L.~Tomassini,
\emph{From local to global deformation quantization of Poisson manifolds},
Duke Math. J. \textbf{115} (2002), 329--352.
\url{https://arxiv.org/abs/math/0012228}.

\bibitem{PolishchukProj}
A.~Polishchuk,
\emph{Noncommutative Proj and coherent algebras},
Math. Res. Lett. \textbf{12} (2005), 63--74.
Theorem~2.4 and the following Remarks~2;
\url{https://arxiv.org/abs/math/0212182}, v4 (2004).

\bibitem{PoissonP4Incidence}
\texttt{uedakazushi/poisson-p4},
\emph{Poisson部分多様体から見る射影四空間の成分分類},
現行研究原稿、\nolinkurl{papers/poisson_subvariety_incidence_ja.tex},
三次延長定理と三次incidenceのファイバー命題。
参照commit \texttt{f118dfb}。
\url{https://github.com/uedakazushi/poisson-p4/blob/f118dfbfc5bb9b7a394d1985b4800b6228eebaf6/papers/poisson_subvariety_incidence_ja.tex}.
公刊済みの外部文献ではなく、関連プロジェクトの証明を持つ研究入力である。

\bibitem{PairBW} S.~Barmeier and Z.~Wang,
\emph{Deformations of path algebras of quivers with relations},
\url{https://arxiv.org/abs/2002.10001}, v4 (2021), Theorem 10.3 and Lemma 10.9.
\bibitem{RainsMagma} E.~Rains, central-element computations in MAGMA form,
\url{https://math.mit.edu/~etingof/delpezzocenter}; E6 data.
\bibitem{AS}
M.~Artin and W.~F.~Schelter,
\textit{Graded algebras of global dimension 3},
Adv. Math. \textbf{66} (1987), 171--216.

\bibitem{AZ}
M.~Artin and J.~J.~Zhang,
\textit{Noncommutative projective schemes},
Adv. Math. \textbf{109} (1994), 228--287.
\href{https://doi.org/10.1006/aima.1994.1087}{doi:10.1006/aima.1994.1087}.

\bibitem{ATV}
M.~Artin, J.~Tate and M.~Van den Bergh,
\textit{Some algebras associated to automorphisms of elliptic curves},
in \textit{The Grothendieck Festschrift}, Vol.~I,
Birkh\"auser (1990), 33--85.
\href{https://doi.org/10.1007/978-0-8176-4574-8_3}{doi:10.1007/978-0-8176-4574-8\_3}.

\bibitem{AD}
M.~Artebani and I.~Dolgachev,
\textit{The Hesse pencil of plane cubic curves},
Enseign. Math. \textbf{55} (2009), 235--273,
\href{https://doi.org/10.4171/LEM/55-3-3}{doi:10.4171/LEM/55-3-3}.

\bibitem{BT}
R.~Berger and R.~Taillefer,
\textit{Poincar\'e--Birkhoff--Witt deformations of Calabi--Yau algebras},
J. Noncommut. Geom. \textbf{1} (2007), 241--270,
\href{https://doi.org/10.4171/JNCG/6}{doi:10.4171/JNCG/6},
\href{https://arxiv.org/abs/math/0610112}{arXiv:math/0610112}.

\bibitem{CKS}
A.~Chirvasitu, R.~Kanda and S.~P.~Smith,
\textit{New Artin--Schelter regular and Calabi--Yau algebras via normal extensions},
\href{https://arxiv.org/abs/1706.05754}{arXiv:1706.05754v3},
\href{https://doi.org/10.1090/tran/7672}{doi:10.1090/tran/7672}.

\bibitem{EG}
P.~Etingof and V.~Ginzburg,
\textit{Noncommutative del Pezzo surfaces and Calabi--Yau algebras},
J. Eur. Math. Soc. \textbf{12} (2010), 1371--1416,
\href{https://doi.org/10.4171/JEMS/235}{doi:10.4171/JEMS/235},
\href{https://arxiv.org/abs/0709.3593}{arXiv:0709.3593v4}.

\bibitem{LPT}
F.~Loray, J.~V.~Pereira and F.~Touzet,
\textit{Foliations with trivial canonical bundle on Fano 3-folds},
Math. Nachr. \textbf{286} (2013), 921--940,
\href{https://arxiv.org/abs/1112.3871}{arXiv:1112.3871}.

\bibitem{Mayanskiy}
E.~Mayanskiy,
\textit{Poisson cohomology of two Fano threefolds},
J. Algebra \textbf{424} (2015), 21--45,
\href{https://doi.org/10.1016/j.jalgebra.2014.08.049}{doi:10.1016/j.jalgebra.2014.08.049},
\href{https://arxiv.org/abs/1303.5916}{arXiv:1303.5916}.

\bibitem{MS}
I.~Mori and S.~P.~Smith,
\textit{The classification of 3-Calabi--Yau algebras with 3 generators and 3 quadratic relations},
Math. Z. (2017),
\href{https://doi.org/10.1007/s00209-016-1824-5}{doi:10.1007/s00209-016-1824-5},
\href{https://arxiv.org/abs/1502.07403}{arXiv:1502.07403v2}.

\bibitem{Pym}
B.~Pym,
\textit{Quantum deformations of projective three-space},
Adv. Math. \textbf{281} (2015), 1216--1241,
\href{https://doi.org/10.1016/j.aim.2015.06.005}{doi:10.1016/j.aim.2015.06.005},
\href{https://arxiv.org/abs/1403.6444}{arXiv:1403.6444}.

\bibitem{Rains}
E.~M.~Rains,
\textit{Filtered deformations of elliptic algebras},
\href{https://arxiv.org/abs/2107.13540v2}{arXiv:2107.13540v2} (2022),
\S2, Theorem~2.13 and its preceding discussion.

\bibitem{RSVdB}
J.~V.~Rennemo, E.~Segal and M.~Van den Bergh,
\textit{A non-commutative Bertini theorem},
J. Noncommut. Geom. \textbf{13} (2019), 609--616,
\href{https://doi.org/10.4171/JNCG/334}{doi:10.4171/JNCG/334},
\href{https://arxiv.org/abs/1705.01366v2}{arXiv:1705.01366v2} (2020),
Proposition~2.1, Lemmas~3.3--3.6 and Corollary~3.9.

\bibitem{Rosenlicht}
M.~Rosenlicht,
\textit{Some basic theorems on algebraic groups},
Amer. J. Math. \textbf{78} (1956), 401--443.

\bibitem{YZ}
A.~Yekutieli and J.~J.~Zhang,
\textit{Serre duality for noncommutative projective schemes},
Proc. Amer. Math. Soc. \textbf{125} (1997), 697--707,
\href{https://doi.org/10.1090/S0002-9939-97-03782-9}{doi:10.1090/S0002-9939-97-03782-9}.

\bibitem{ZSL}
G.-S.~Zhou, Y.~Shen and D.-M.~Lu,
\textit{Skew Calabi--Yau property of normal extensions},
Manuscripta Math. \textbf{161} (2020), 125--140,
\href{https://doi.org/10.1007/s00229-018-1064-6}{doi:10.1007/s00229-018-1064-6},
\href{https://arxiv.org/abs/1709.09292}{arXiv:1709.09292}.

\bibitem{Zhang}
J.~J.~Zhang,
\textit{Twisted graded algebras and equivalences of graded categories},
Proc. London Math. Soc. \textbf{72} (1996), 281--311,
\href{https://doi.org/10.1112/plms/s3-72.2.281}{doi:10.1112/plms/s3-72.2.281}.
\bibitem{PBondalPolishchukHelices}
A.~I.~Bondal and A.~E.~Polishchuk,
\emph{Homological properties of associative algebras: the method of helices},
Russian Acad. Sci. Izv. Math. \textbf{42} (1994), no.~2, 219--260.
Proposition~4.1 and \S5 (Koszul--Frobenius duality and periodicity).
\href{https://doi.org/10.1070/IM1994v042n02ABEH001536}
{doi:10.1070/IM1994v042n02ABEH001536}.

\bibitem{PChanNymanNormalFamily}
D.~Chan and A.~Nyman,
\emph{Noncommutative linear systems and noncommutative elliptic curves},
\href{https://arxiv.org/abs/2211.12465}{arXiv:2211.12465} (2022).
Definition~8.6 (normal families of elements in indexed algebras).

\bibitem{CSBGCCKM}
T.~Cassidy, A.~Conner, E.~Kirkman and W.~F.~Moore,
\emph{Periodic free resolutions from twisted matrix factorizations},
J. Algebra 455 (2016), 137--163;
\href{https://doi.org/10.1016/j.jalgebra.2016.01.037}{doi:10.1016/j.jalgebra.2016.01.037};
\url{https://arxiv.org/abs/1307.8415}.
arXiv v1のLemma~4.4（Reesの補題）と公刊版のTheorem~5.8
（twisted matrix factorizationと特異圏）を版を区別して用いる。
有限位数の議論にはTheorems~4.7 and~5.5を用いる。

\bibitem{BSWPotentials}
R.~Bocklandt, T.~Schedler and M.~Wemyss,
\emph{Superpotentials and higher order derivations},
J. Pure Appl. Algebra \textbf{214} (2010), no.~9, 1501--1522.
\href{https://doi.org/10.1016/j.jpaa.2009.07.013}{doi:10.1016/j.jpaa.2009.07.013};
\href{https://arxiv.org/abs/0802.0162v3}{arXiv:0802.0162v3}.
Equation~(1.3), Theorem~6.2, Lemmas~6.4--6.5,
およびTheorem~6.2証明直後のtwisted版への拡張。

\bibitem{StacksPerfect}
The Stacks Project Authors, \emph{The Stacks Project}.
\href{https://stacks.math.columbia.edu/tag/0BDK}{Tag~0BDK, Lemma~36.31.3},
\href{https://stacks.math.columbia.edu/tag/0BCD}{Tag~0BCD, Lemma~15.77.7}.
perfect複体のファイバーでの単一次数集中と局所自由性。
\href{https://stacks.math.columbia.edu/tag/08EV}{Tag~08EV, Lemma~36.27.1}
（底に関して有限Tor振幅を持つcoherent複体のproper pushforward）を併用する。
2026年11月19日閲覧。

\bibitem{StacksProj}
The Stacks Project Authors, \emph{The Stacks Project}.
\href{https://stacks.math.columbia.edu/tag/01YR}{Section~30.14, Tag~01YR},
特にLemma~30.14.2（coherent層の生成性・消滅）と
Proposition~30.14.4（標準次数付き環のProj対応）。2026年11月19日閲覧。

\bibitem{MoriNymanZ}
I.~Mori and A.~Nyman,
\emph{A Categorical Characterization of Quantum Projective $\mathbb Z$-spaces},
Appl. Categ. Structures \textbf{33} (2025), article~14.
\href{https://doi.org/10.1007/s10485-025-09806-2}{doi:10.1007/s10485-025-09806-2};
\url{https://arxiv.org/abs/2307.15253}.
Theorem~4.19, Lemma~5.24, Theorems~6.4 and~6.11。
著者版：\url{https://faculty.nymana.wwu.edu/regV8.pdf}。

\bibitem{ASZ4Canonical}
\texttt{uedakazushi/asz4},
\emph{四次元二次AS正則$\mathbb Z$-代数の分類：
4周期表示・1周期構造の忘却・全分類への有限atlas},
\nolinkurl{papers/asz4_classification.tex}, 定理\texttt{canonical-four}。
参照commit \texttt{91b4504}。
\url{https://github.com/uedakazushi/asz4/blob/91b45047739187fc5acb6f906b25e48f49a42c2d/papers/asz4_classification.tex}.
関連プロジェクトの研究入力であり、公刊済み文献として扱わない。

\bibitem{BILNComponents}
V.~Bhatoy, C.~Ingalls, F.~LaRoche and R.~Nookala,
\emph{Some components of the moduli space of Koszul Artin--Schelter regular algebras of dimension four},
\href{https://arxiv.org/abs/2511.08390v1}{arXiv:2511.08390v1}
(2025年11月11日), \S2.3, Proposition~5.1 and Theorem~5.8。
本稿では定義と同値関係を比較し、その成分判定を証明の入力にしない。

\bibitem{BIDiscrete}
V.~Bhatoy and C.~Ingalls,
\emph{Discrete Invariants of Koszul Artin--Schelter Regular Algebras of Dimension four},
\href{https://arxiv.org/abs/2602.13178v1}{arXiv:2602.13178v1}
(2026年2月13日), Introduction。

\bibitem{OrlovLandauGinzburg}
D.~Orlov,
\emph{Triangulated categories of singularities and D-branes in Landau--Ginzburg models},
\href{https://arxiv.org/abs/math/0302304v2}{arXiv:math/0302304v2}
(2004), \S3.1 and Theorem~3.9。

\bibitem{OrlovSingularities}
D.~Orlov,
\emph{Derived Categories of Coherent Sheaves and Triangulated Categories of Singularities},
in \emph{Algebra, arithmetic, and geometry: in honor of Yu.~I.~Manin},
Vol.~II, Progress in Mathematics \textbf{270}, Birkh\"auser, 2009, 503--531.
\href{https://doi.org/10.1007/978-0-8176-4747-6_16}{doi:10.1007/978-0-8176-4747-6\_16};
\href{https://arxiv.org/abs/math/0503632v3}{arXiv:math/0503632v3},
Theorem~2.5。

\bibitem{KuznetsovFractionalCY}
A.~Kuznetsov,
\emph{Calabi--Yau and fractional Calabi--Yau categories},
J. reine angew. Math. \textbf{753} (2019), 239--267.
\href{https://doi.org/10.1515/crelle-2017-0004}{doi:10.1515/crelle-2017-0004};
\href{https://arxiv.org/abs/1509.07657v2}{arXiv:1509.07657v2},
Corollary~4.1。

\bibitem{KuznetsovHochschild}
A.~Kuznetsov,
\emph{Hochschild homology and semiorthogonal decompositions},
\href{https://arxiv.org/abs/0904.4330v1}{arXiv:0904.4330v1} (2009).
Theorems~8.3, 8.9, Corollary~7.5 and Proposition~8.4。

\bibitem{CSBGChenGeneralizedSerre}
X.-W.~Chen, \emph{Generalized Serre duality},
\url{https://arxiv.org/abs/math/0610258}, Theorem~3.10.

\bibitem{PPQuadratic}
A.~Polishchuk and L.~Positselski,
\emph{Quadratic Algebras}, University Lecture Series \textbf{37},
American Mathematical Society, Providence, RI, 2005, xii+159 pp.
\href{https://doi.org/10.1090/ulect/037}{doi:10.1090/ulect/037}.
Chapter~6, Proposition~7.2 and Theorem~7.3, pp.~129--130。
Noetherian局所基底上のKoszul変形原理。本稿の非周期的Artin版には直接証明を与える。


\bibitem{YekutieliBalanced}
A.~Yekutieli,
\emph{Dualizing complexes over noncommutative graded algebras},
J. Algebra \textbf{153} (1992), 41--84.
\href{https://doi.org/10.1016/0021-8693(92)90148-F}{doi:10.1016/0021-8693(92)90148-F};
\url{https://www.math.bgu.ac.il/~amyekut/publications/dual-com/dual-com.pdf}.
Definition~4.1と双加群としての一意性（p.~57），Proposition~4.6（p.~59），
Proposition~5.2，Corollary~5.5，Remark~5.7（pp.~65--68）。

\bibitem{BDGMinors}
I.~Burban, Yu.~Drozd and V.~Gavran,
\emph{Minors and resolutions of non-commutative schemes},
\url{https://arxiv.org/abs/1501.06023}, v2 (2017),
Definitions~3.1,3.13 and Theorem~3.14。
quasi-compact separatedな非可換スキームの局所perfect対象とcompact対象の一致。

\bibitem{StacksRelativeDuality}
The Stacks Project Authors, \emph{The Stacks Project}.
\href{https://stacks.math.columbia.edu/tag/00MI}{Lemma~10.99.5, Tag~00MI}，
\href{https://stacks.math.columbia.edu/tag/07VK}{Lemma~30.22.1, Tag~07VK}，
\href{https://stacks.math.columbia.edu/tag/0B91}{Lemma~36.30.4, Tag~0B91}，
\href{https://stacks.math.columbia.edu/tag/0A9W}{Lemma~48.15.1, Tag~0A9W}，
\href{https://stacks.math.columbia.edu/tag/0B6J}{Lemma~48.7.1, Tag~0B6J}，
\href{https://stacks.math.columbia.edu/tag/0B6S}{Remark~48.12.5, Tag~0B6S}，
\href{https://stacks.math.columbia.edu/tag/0E4H}{Remark~48.12.6}，
\href{https://stacks.math.columbia.edu/tag/047N}{Lemma~39.8.2, Tag~047N}。
振幅とflatness，proper base change，射影空間の双対・trace，Cartierのsmoothness。
2026年11月19日閲覧。

\bibitem{KellerDgCategories}
B.~Keller, \emph{On differential graded categories},
\url{https://arxiv.org/abs/math/0601185v5} (2006).
\S2.1（可換係数環の規約），Theorems~3.4,3.8(c), \S3.8。

\bibitem{KeelerAmpleFilters}
D.~S.~Keeler, \emph{Ample filters of invertible sheaves},
J. Algebra \textbf{259} (2003), no.~1, 243--283;
Corrigendum, J. Algebra \textbf{507} (2018), 592--598.
\url{https://arxiv.org/abs/math/0108068v5}（訂正を含む2018年6月2日版）。
Theorem~1.3, Proposition~2.7, Theorem~7.14。

\bibitem{MizunoZDuality}
Y.~Mizuno, \emph{Dualizing complexes over $\mathbb Z$-algebras},
\url{https://arxiv.org/abs/2509.13073v3}（2026年9月4日版）。
\S2の係数体の規約とTheorem~5.8の有限global dimension仮定を参照。

\bibitem{BrownSridharSerre}
M.~K.~Brown and P.~Sridhar, \emph{Serre duality for dg-algebras},
\url{https://arxiv.org/abs/2410.07204v2}（2025年8月15日版）。
\S2の係数体の規約を参照。
\bibitem{BrownSotoSridharBalanced}
M.~K.~Brown, A.~J.~Soto Levins and P.~Sridhar,
\emph{Existence of balanced dualizing dg-modules},
\url{https://arxiv.org/abs/2506.02398v1}（2025年6月3日版）。
Notation~1.4, Setup~2.1, Corollary~1.3。

\bibitem{StacksAlgebraicChart}
The Stacks Project Authors, \emph{The Stacks Project}.
\href{https://stacks.math.columbia.edu/tag/03WT}{Lemma~29.37.13, Tag~03WT}，
\href{https://stacks.math.columbia.edu/tag/02GU}{Lemma~29.37.15, Tag~02GU}，
\href{https://stacks.math.columbia.edu/tag/0476}{Lemma~29.37.17, Tag~0476}，
\href{https://stacks.math.columbia.edu/tag/0245}{Lemma~35.37.1, Tag~0245}，
\href{https://stacks.math.columbia.edu/tag/04UV}{Section~95.15, Tag~04UV}，
\href{https://stacks.math.columbia.edu/tag/06FI}{Theorem~97.17.2, Tag~06FI}，
\href{https://stacks.math.columbia.edu/tag/06N3}{Theorem~101.21.6, Tag~06N3}。
étale locus とその底変更、群作用商・torsor の表示、DM atlas、アフィン降下。
2026年11月19日閲覧。

\bibitem{ELSSmoothAlgebras}
A.~Elagin, V.~A.~Lunts and O.~M.~Schn\"urer,
\emph{Smoothness of Derived Categories of Algebras},
Moscow Math.~J. \textbf{20} (2020), no.~2, 277--309.
\url{https://arxiv.org/abs/1810.07626}, v2 (2019), Theorems~4.15,4.18.
\url{https://www.mathjournals.org/mmj/2020-020-002/2020-020-002-003.html}.
生成定理は著者版26--27頁。

\bibitem{ToenVaquieObjects}
B.~To\"en and M.~Vaqui\'e,
\emph{Moduli of objects in dg-categories},
Ann.~Sci.~\'Ec.~Norm.~Sup\'er.~(4) \textbf{40} (2007), 387--444.
\href{https://doi.org/10.1016/j.ansens.2007.05.001}{doi:10.1016/j.ansens.2007.05.001}.
\url{https://www.numdam.org/article/ASENS_2007_4_40_3_387_0.pdf}.
Lemma~2.10（405--407頁）、Corollary~2.13、Proposition~2.14。

\bibitem{ToenDefinitionRings}
B.~To\"en,
\emph{Anneaux de d\'efinition des dg-alg\`ebres propres et lisses},
著者版（2007年5月）, 事実 F4（2頁）.
\url{https://www.math.univ-toulouse.fr/~btoen/defdgalg.pdf};
\url{https://arxiv.org/abs/math/0611546}.
F4 の perfect 加群の連続性は任意の出発 dg 代数に適用される。
F5 の pseudoperfect 加群についての hfp 仮定とは区別する。

\bibitem{WWYSklyaninPoisson}
C.~Walton, X.~Wang and M.~Yakimov,
\emph{Poisson geometry of PI 3-dimensional Sklyanin algebras},
Proc.~Lond.~Math.~Soc. (3) \textbf{118} (2019), 1471--1500.
\url{https://arxiv.org/abs/1704.04975}, v3 (2018),
Definition--Lemma~2.13（9頁）と Proposition~2.22。
平行移動点の規約については原典\cite{ATV}も参照。

\bibitem{StacksEllipticTorsion}
The Stacks Project Authors, \emph{The Stacks Project}.
\href{https://stacks.math.columbia.edu/tag/0BFG}{Tag~0BFG, Lemma~39.9.8},
\href{https://stacks.math.columbia.edu/tag/0BFH}{Tag~0BFH, Lemma~39.9.9}.
体上の乗法写像の有限性と、標数と互いに素な場合の \'etale 性。
Hesse 族上の主張には本文の相対微分と properness の議論を加える。
2026年11月19日閲覧。

\bibitem{JLLZInfinitesimal}
A.~Jacovskis, X.~Lin, Z.~Liu and S.~Zhang,
\emph{Infinitesimal categorical Torelli theorems for Fano threefolds},
\href{https://arxiv.org/abs/2203.08187v3}{arXiv:2203.08187v3} (2023),
Theorem~4.4(2)。実際の幾何的 Kodaira--Spencer 写像を使用する。

\bibitem{LurieDAGX}
J.~Lurie, \emph{Derived Algebraic Geometry X: Formal Moduli Problems},
November 5, 2011,
\url{https://www.math.ias.edu/~lurie/papers/DAG-X.pdf},
Theorems~5.1.9, 5.3.16 and Proposition~5.3.21。

\bibitem{BKPFormalGenerators}
A.~Blanc, L.~Katzarkov and P.~Pandit,
\emph{Generators in formal deformations of categories},
Compos. Math. \textbf{154} (2018), 2055--2089,
\href{https://doi.org/10.1112/S0010437X18007303}{doi:10.1112/S0010437X18007303};
\href{https://arxiv.org/abs/1705.00655v1}{arXiv:1705.00655v1}, \S3.3。

\bibitem{BhattAlgebraization}
B.~Bhatt, \emph{Algebraization and Tannaka duality},
\href{https://arxiv.org/abs/1404.7483v1}{arXiv:1404.7483v1} (2014),
Lemma~4.2。
通常の係数環上の perfect complex の効果化として用い、
任意の proper dg 代数上の compact 加群へ拡張しない。

\bibitem{StevensonSingularityActions}
G.~Stevenson, \emph{Subcategories of singularity categories via tensor actions},
Compos. Math. \textbf{150} (2014), 229--272,
\href{https://arxiv.org/abs/1105.4698v3}{arXiv:1105.4698v3},
Theorem~3.1(3), Notation~3.2, Lemma~3.10 and Theorem~6.13。

\bibitem{NeemanCompactLocalization}
A.~Neeman, \emph{The connection between the K-theory localization theorem of Thomason,
Trobaugh and Yao and the smashing subcategories of Bousfield and Ravenel},
Ann. Sci. \'Ecole Norm. Sup. (4) \textbf{25} (1992), 547--566,
\href{https://doi.org/10.24033/asens.1659}{doi:10.24033/asens.1659}, Lemma~2.2;
\url{https://www.sas.rochester.edu/mth/sites/doug-ravenel/otherpapers/neeman-loc.pdf}。

\bibitem{BMMSCubicTorelli}
M.~Bernardara, E.~Macr\`i, S.~Mehrotra and P.~Stellari,
\emph{A categorical invariant for cubic threefolds},
\href{https://arxiv.org/abs/0903.4414v2}{arXiv:0903.4414v2} (2011), Theorem~1.1。
複素数体上の定理を用い、拡大体への適用は本文で有限表示により帰着する。

\bibitem{PoonenHypersurfaceAut}
B.~Poonen, \emph{Varieties without extra automorphisms III: hypersurfaces},
Finite Fields Appl. \textbf{11} (2005), 230--268,
\url{https://math.mit.edu/~poonen/papers/projaut.pdf},
Theorems~1.5--1.6 and Corollary~1.9。

\bibitem{AllcockCubicGIT}
D.~Allcock, \emph{The moduli space of cubic threefolds},
J. Algebraic Geom. \textbf{12} (2003), 201--223,
\href{https://doi.org/10.1090/S1056-3911-02-00313-2}{doi:10.1090/S1056-3911-02-00313-2};
\url{https://web.ma.utexas.edu/users/allcock/research/3foldgit.pdf}, Theorem~1.1。


\bibitem{LeBruynSmithVdBCentral}
L.~Le~Bruyn, S.~P.~Smith and M.~Van~den~Bergh,
\textit{Central extensions of three dimensional Artin--Schelter regular algebras},
Math. Z. \textbf{222} (1996), 171--212,
\href{https://doi.org/10.1007/PL00004532}{doi:10.1007/PL00004532}.
\url{https://sites.math.washington.edu/~smith/Research/Central-extns.pdf}.

\bibitem{VancliffQuantumSpaces}
M.~Vancliff,
\textit{Non-commutative Spaces for Graded Quantum Groups and Graded Clifford Algebras},
in R.~Ab{\l}amowicz and B.~Fauser (eds.),
\textit{Clifford Algebras and their Applications in Mathematical Physics, I},
Progress in Physics \textbf{18}, Birkh\"auser Boston (2000), 303--320,
\href{https://doi.org/10.1007/978-1-4612-1368-0_16}{doi:10.1007/978-1-4612-1368-0\_16}.
\url{https://bpb-us-e1.wpmucdn.com/websites.uta.edu/dist/0/5123/files/2022/10/cliff.pdf}.

\bibitem{PymSurvey}
B.~Pym,
\textit{Constructions and classifications of projective Poisson varieties},
Lett. Math. Phys. \textbf{108} (2018), 573--632,
\href{https://doi.org/10.1007/s11005-017-0984-5}{doi:10.1007/s11005-017-0984-5},
\href{https://arxiv.org/abs/1701.08852}{arXiv:1701.08852}.

\bibitem{GualtieriPym}
M.~Gualtieri and B.~Pym,
\textit{Poisson modules and degeneracy loci},
Proc. Lond. Math. Soc. (3) \textbf{107} (2013), no.~3, 627--654,
\href{https://doi.org/10.1112/plms/pds090}{doi:10.1112/plms/pds090},
\href{https://arxiv.org/abs/1203.4293}{arXiv:1203.4293}.

\bibitem{UedaYoshinagaJacobian}
K.~Ueda and M.~Yoshinaga,
\textit{Logarithmic vector fields along smooth divisors in projective spaces},
Hokkaido Math. J. \textbf{38} (2009), 409--415,
\href{https://doi.org/10.14492/hokmj/1258553970}{doi:10.14492/hokmj/1258553970}.
\url{https://www.math.sci.hokudai.ac.jp/hmj/page/38-3/pdf/HMJ38_0409-0415.pdf}.

\bibitem{HwangSymmetrizer}
J.-M.~Hwang,
\textit{Symmetrizer group of a projective hypersurface},
J. Math. Soc. Japan \textbf{78} (2026), no.~1, 55--62,
\href{https://doi.org/10.2969/jmsj/93469346}{doi:10.2969/jmsj/93469346},
\href{https://arxiv.org/abs/2501.08633}{arXiv:2501.08633}.

\bibitem{LinZhangDeformingEquivalences}
X.~Lin and S.~Zhang,
\emph{Categorical Torelli and Fourier--Mukai partners of hypersurfaces},
\href{https://arxiv.org/abs/2611.05915v1}{arXiv:2611.05915v1} (2026),
Theorem~1.7 and \S2.

\end{thebibliography}
\end{document}
